\documentclass{article}
\usepackage[paperwidth=16cm,paperheight=14cm,margin=0cm]{geometry}
\usepackage[utf8]{inputenc}\usepackage[T1]{fontenc}
\usepackage{mathpazo}\usepackage{amsmath,amssymb}
\usepackage[most]{tcolorbox}
\pagestyle{empty}\setlength{\parindent}{0pt}

\definecolor{cari}{HTML}{1F6FA8}
\definecolor{clit}{HTML}{2E7D6B}
\definecolor{cfon}{HTML}{6A4C93}
\definecolor{cgeo}{HTML}{C0503A}
\definecolor{csta}{HTML}{8A5A2B}
\definecolor{cgra}{HTML}{C77F0A}
\begin{document}
\begin{tcolorbox}[enhanced,width=16cm,height=14cm,colback=white,colframe=cari,boxrule=2pt,arc=0mm,
left=6mm,right=6mm,top=0mm,bottom=4mm,sharp corners,
overlay={\fill[cari] (frame.north west) rectangle ([yshift=-1.25cm]frame.north east);
\node[anchor=west,text=white,font=\sffamily\bfseries\LARGE] at ([xshift=5mm,yshift=-0.62cm]frame.north west) {Question 1};
\node[anchor=east,text=white,font=\sffamily\itshape\large] at ([xshift=-5mm,yshift=-0.62cm]frame.north east) {Arithmétique et puissances};}]
\vspace*{1.75cm}
\begin{minipage}[t][6.9cm][c]{\linewidth}\LARGE\centering
\begin{minipage}{0.95\linewidth}\raggedright\linespread{1.25}\selectfont Lequel de ces nombres est premier~?\end{minipage}
\end{minipage}

\vfill
\begin{tcolorbox}[enhanced,nobeforeafter,width=7.1cm,height=2.1cm,valign=center,colback=cari!6,colframe=cari,arc=2.5mm,boxrule=0.9pt,left=3mm]{\Large\bfseries\color{cari}a)}\quad\begin{minipage}[c]{5.3cm}\raggedright\Large $79$\end{minipage}\end{tcolorbox}\hfill\begin{tcolorbox}[enhanced,nobeforeafter,width=7.1cm,height=2.1cm,valign=center,colback=cari!6,colframe=cari,arc=2.5mm,boxrule=0.9pt,left=3mm]{\Large\bfseries\color{cari}b)}\quad\begin{minipage}[c]{5.3cm}\raggedright\Large $81$\end{minipage}\end{tcolorbox}\par\vspace{2.5mm}
\begin{tcolorbox}[enhanced,nobeforeafter,width=7.1cm,height=2.1cm,valign=center,colback=cari!6,colframe=cari,arc=2.5mm,boxrule=0.9pt,left=3mm]{\Large\bfseries\color{cari}c)}\quad\begin{minipage}[c]{5.3cm}\raggedright\Large $85$\end{minipage}\end{tcolorbox}\hfill\begin{tcolorbox}[enhanced,nobeforeafter,width=7.1cm,height=2.1cm,valign=center,colback=cari!6,colframe=cari,arc=2.5mm,boxrule=0.9pt,left=3mm]{\Large\bfseries\color{cari}d)}\quad\begin{minipage}[c]{5.3cm}\raggedright\Large $77$\end{minipage}\end{tcolorbox}
\end{tcolorbox}
\newpage
\begin{tcolorbox}[enhanced,width=16cm,height=14cm,colback=white,colframe=cari,boxrule=2pt,arc=0mm,
left=6mm,right=6mm,top=0mm,bottom=4mm,sharp corners,
overlay={\fill[cari] (frame.north west) rectangle ([yshift=-1.25cm]frame.north east);
\node[anchor=west,text=white,font=\sffamily\bfseries\LARGE] at ([xshift=5mm,yshift=-0.62cm]frame.north west) {Question 2};
\node[anchor=east,text=white,font=\sffamily\itshape\large] at ([xshift=-5mm,yshift=-0.62cm]frame.north east) {Arithmétique et puissances};}]
\vspace*{1.75cm}
\begin{minipage}[t][6.9cm][c]{\linewidth}\LARGE\centering
\begin{minipage}{0.95\linewidth}\raggedright\linespread{1.25}\selectfont Lequel de ces nombres est premier~?\end{minipage}
\end{minipage}

\vfill
\begin{tcolorbox}[enhanced,nobeforeafter,width=7.1cm,height=2.1cm,valign=center,colback=cari!6,colframe=cari,arc=2.5mm,boxrule=0.9pt,left=3mm]{\Large\bfseries\color{cari}a)}\quad\begin{minipage}[c]{5.3cm}\raggedright\Large $63$\end{minipage}\end{tcolorbox}\hfill\begin{tcolorbox}[enhanced,nobeforeafter,width=7.1cm,height=2.1cm,valign=center,colback=cari!6,colframe=cari,arc=2.5mm,boxrule=0.9pt,left=3mm]{\Large\bfseries\color{cari}b)}\quad\begin{minipage}[c]{5.3cm}\raggedright\Large $59$\end{minipage}\end{tcolorbox}\par\vspace{2.5mm}
\begin{tcolorbox}[enhanced,nobeforeafter,width=7.1cm,height=2.1cm,valign=center,colback=cari!6,colframe=cari,arc=2.5mm,boxrule=0.9pt,left=3mm]{\Large\bfseries\color{cari}c)}\quad\begin{minipage}[c]{5.3cm}\raggedright\Large $57$\end{minipage}\end{tcolorbox}\hfill\begin{tcolorbox}[enhanced,nobeforeafter,width=7.1cm,height=2.1cm,valign=center,colback=cari!6,colframe=cari,arc=2.5mm,boxrule=0.9pt,left=3mm]{\Large\bfseries\color{cari}d)}\quad\begin{minipage}[c]{5.3cm}\raggedright\Large $65$\end{minipage}\end{tcolorbox}
\end{tcolorbox}
\newpage
\begin{tcolorbox}[enhanced,width=16cm,height=14cm,colback=white,colframe=cari,boxrule=2pt,arc=0mm,
left=6mm,right=6mm,top=0mm,bottom=4mm,sharp corners,
overlay={\fill[cari] (frame.north west) rectangle ([yshift=-1.25cm]frame.north east);
\node[anchor=west,text=white,font=\sffamily\bfseries\LARGE] at ([xshift=5mm,yshift=-0.62cm]frame.north west) {Question 3};
\node[anchor=east,text=white,font=\sffamily\itshape\large] at ([xshift=-5mm,yshift=-0.62cm]frame.north east) {Arithmétique et puissances};}]
\vspace*{1.75cm}
\begin{minipage}[t][6.9cm][c]{\linewidth}\LARGE\centering
\begin{minipage}{0.95\linewidth}\raggedright\linespread{1.25}\selectfont Quelle est la décomposition de $120$ en produit de facteurs premiers~?\end{minipage}
\end{minipage}

\vfill
\begin{tcolorbox}[enhanced,nobeforeafter,width=7.1cm,height=2.1cm,valign=center,colback=cari!6,colframe=cari,arc=2.5mm,boxrule=0.9pt,left=3mm]{\Large\bfseries\color{cari}a)}\quad\begin{minipage}[c]{5.3cm}\raggedright\Large $2^{4}\times3\times5$\end{minipage}\end{tcolorbox}\hfill\begin{tcolorbox}[enhanced,nobeforeafter,width=7.1cm,height=2.1cm,valign=center,colback=cari!6,colframe=cari,arc=2.5mm,boxrule=0.9pt,left=3mm]{\Large\bfseries\color{cari}b)}\quad\begin{minipage}[c]{5.3cm}\raggedright\Large $2^{3}\times3\times5^{2}$\end{minipage}\end{tcolorbox}\par\vspace{2.5mm}
\begin{tcolorbox}[enhanced,nobeforeafter,width=7.1cm,height=2.1cm,valign=center,colback=cari!6,colframe=cari,arc=2.5mm,boxrule=0.9pt,left=3mm]{\Large\bfseries\color{cari}c)}\quad\begin{minipage}[c]{5.3cm}\raggedright\Large $24\times5$\end{minipage}\end{tcolorbox}\hfill\begin{tcolorbox}[enhanced,nobeforeafter,width=7.1cm,height=2.1cm,valign=center,colback=cari!6,colframe=cari,arc=2.5mm,boxrule=0.9pt,left=3mm]{\Large\bfseries\color{cari}d)}\quad\begin{minipage}[c]{5.3cm}\raggedright\Large $2^{3}\times3\times5$\end{minipage}\end{tcolorbox}
\end{tcolorbox}
\newpage
\begin{tcolorbox}[enhanced,width=16cm,height=14cm,colback=white,colframe=cari,boxrule=2pt,arc=0mm,
left=6mm,right=6mm,top=0mm,bottom=4mm,sharp corners,
overlay={\fill[cari] (frame.north west) rectangle ([yshift=-1.25cm]frame.north east);
\node[anchor=west,text=white,font=\sffamily\bfseries\LARGE] at ([xshift=5mm,yshift=-0.62cm]frame.north west) {Question 4};
\node[anchor=east,text=white,font=\sffamily\itshape\large] at ([xshift=-5mm,yshift=-0.62cm]frame.north east) {Arithmétique et puissances};}]
\vspace*{1.75cm}
\begin{minipage}[t][6.9cm][c]{\linewidth}\LARGE\centering
\begin{minipage}{0.95\linewidth}\raggedright\linespread{1.25}\selectfont Quelle est la décomposition de $98$ en produit de facteurs premiers~?\end{minipage}
\end{minipage}

\vfill
\begin{tcolorbox}[enhanced,nobeforeafter,width=7.1cm,height=2.1cm,valign=center,colback=cari!6,colframe=cari,arc=2.5mm,boxrule=0.9pt,left=3mm]{\Large\bfseries\color{cari}a)}\quad\begin{minipage}[c]{5.3cm}\raggedright\Large $2^{2}\times7^{2}$\end{minipage}\end{tcolorbox}\hfill\begin{tcolorbox}[enhanced,nobeforeafter,width=7.1cm,height=2.1cm,valign=center,colback=cari!6,colframe=cari,arc=2.5mm,boxrule=0.9pt,left=3mm]{\Large\bfseries\color{cari}b)}\quad\begin{minipage}[c]{5.3cm}\raggedright\Large $14\times7$\end{minipage}\end{tcolorbox}\par\vspace{2.5mm}
\begin{tcolorbox}[enhanced,nobeforeafter,width=7.1cm,height=2.1cm,valign=center,colback=cari!6,colframe=cari,arc=2.5mm,boxrule=0.9pt,left=3mm]{\Large\bfseries\color{cari}c)}\quad\begin{minipage}[c]{5.3cm}\raggedright\Large $2\times7^{3}$\end{minipage}\end{tcolorbox}\hfill\begin{tcolorbox}[enhanced,nobeforeafter,width=7.1cm,height=2.1cm,valign=center,colback=cari!6,colframe=cari,arc=2.5mm,boxrule=0.9pt,left=3mm]{\Large\bfseries\color{cari}d)}\quad\begin{minipage}[c]{5.3cm}\raggedright\Large $2\times7^{2}$\end{minipage}\end{tcolorbox}
\end{tcolorbox}
\newpage
\begin{tcolorbox}[enhanced,width=16cm,height=14cm,colback=white,colframe=cari,boxrule=2pt,arc=0mm,
left=6mm,right=6mm,top=0mm,bottom=4mm,sharp corners,
overlay={\fill[cari] (frame.north west) rectangle ([yshift=-1.25cm]frame.north east);
\node[anchor=west,text=white,font=\sffamily\bfseries\LARGE] at ([xshift=5mm,yshift=-0.62cm]frame.north west) {Question 5};
\node[anchor=east,text=white,font=\sffamily\itshape\large] at ([xshift=-5mm,yshift=-0.62cm]frame.north east) {Arithmétique et puissances};}]
\vspace*{1.75cm}
\begin{minipage}[t][6.9cm][c]{\linewidth}\LARGE\centering
\begin{minipage}{0.95\linewidth}\raggedright\linespread{1.25}\selectfont Quelle est la forme irréductible de $\dfrac{126}{180}$~?\end{minipage}
\end{minipage}

\vfill
\begin{tcolorbox}[enhanced,nobeforeafter,width=7.1cm,height=2.1cm,valign=center,colback=cari!6,colframe=cari,arc=2.5mm,boxrule=0.9pt,left=3mm]{\Large\bfseries\color{cari}a)}\quad\begin{minipage}[c]{5.3cm}\raggedright\Large $\dfrac{10}{7}$\end{minipage}\end{tcolorbox}\hfill\begin{tcolorbox}[enhanced,nobeforeafter,width=7.1cm,height=2.1cm,valign=center,colback=cari!6,colframe=cari,arc=2.5mm,boxrule=0.9pt,left=3mm]{\Large\bfseries\color{cari}b)}\quad\begin{minipage}[c]{5.3cm}\raggedright\Large $\dfrac{63}{90}$\end{minipage}\end{tcolorbox}\par\vspace{2.5mm}
\begin{tcolorbox}[enhanced,nobeforeafter,width=7.1cm,height=2.1cm,valign=center,colback=cari!6,colframe=cari,arc=2.5mm,boxrule=0.9pt,left=3mm]{\Large\bfseries\color{cari}c)}\quad\begin{minipage}[c]{5.3cm}\raggedright\Large $\dfrac{42}{60}$\end{minipage}\end{tcolorbox}\hfill\begin{tcolorbox}[enhanced,nobeforeafter,width=7.1cm,height=2.1cm,valign=center,colback=cari!6,colframe=cari,arc=2.5mm,boxrule=0.9pt,left=3mm]{\Large\bfseries\color{cari}d)}\quad\begin{minipage}[c]{5.3cm}\raggedright\Large $\dfrac{7}{10}$\end{minipage}\end{tcolorbox}
\end{tcolorbox}
\newpage
\begin{tcolorbox}[enhanced,width=16cm,height=14cm,colback=white,colframe=cari,boxrule=2pt,arc=0mm,
left=6mm,right=6mm,top=0mm,bottom=4mm,sharp corners,
overlay={\fill[cari] (frame.north west) rectangle ([yshift=-1.25cm]frame.north east);
\node[anchor=west,text=white,font=\sffamily\bfseries\LARGE] at ([xshift=5mm,yshift=-0.62cm]frame.north west) {Question 6};
\node[anchor=east,text=white,font=\sffamily\itshape\large] at ([xshift=-5mm,yshift=-0.62cm]frame.north east) {Arithmétique et puissances};}]
\vspace*{1.75cm}
\begin{minipage}[t][6.9cm][c]{\linewidth}\LARGE\centering
\begin{minipage}{0.95\linewidth}\raggedright\linespread{1.25}\selectfont Quelle est la forme irréductible de $\dfrac{90}{252}$~?\end{minipage}
\end{minipage}

\vfill
\begin{tcolorbox}[enhanced,nobeforeafter,width=7.1cm,height=2.1cm,valign=center,colback=cari!6,colframe=cari,arc=2.5mm,boxrule=0.9pt,left=3mm]{\Large\bfseries\color{cari}a)}\quad\begin{minipage}[c]{5.3cm}\raggedright\Large $\dfrac{30}{84}$\end{minipage}\end{tcolorbox}\hfill\begin{tcolorbox}[enhanced,nobeforeafter,width=7.1cm,height=2.1cm,valign=center,colback=cari!6,colframe=cari,arc=2.5mm,boxrule=0.9pt,left=3mm]{\Large\bfseries\color{cari}b)}\quad\begin{minipage}[c]{5.3cm}\raggedright\Large $\dfrac{45}{126}$\end{minipage}\end{tcolorbox}\par\vspace{2.5mm}
\begin{tcolorbox}[enhanced,nobeforeafter,width=7.1cm,height=2.1cm,valign=center,colback=cari!6,colframe=cari,arc=2.5mm,boxrule=0.9pt,left=3mm]{\Large\bfseries\color{cari}c)}\quad\begin{minipage}[c]{5.3cm}\raggedright\Large $\dfrac{5}{14}$\end{minipage}\end{tcolorbox}\hfill\begin{tcolorbox}[enhanced,nobeforeafter,width=7.1cm,height=2.1cm,valign=center,colback=cari!6,colframe=cari,arc=2.5mm,boxrule=0.9pt,left=3mm]{\Large\bfseries\color{cari}d)}\quad\begin{minipage}[c]{5.3cm}\raggedright\Large $\dfrac{14}{5}$\end{minipage}\end{tcolorbox}
\end{tcolorbox}
\newpage
\begin{tcolorbox}[enhanced,width=16cm,height=14cm,colback=white,colframe=cari,boxrule=2pt,arc=0mm,
left=6mm,right=6mm,top=0mm,bottom=4mm,sharp corners,
overlay={\fill[cari] (frame.north west) rectangle ([yshift=-1.25cm]frame.north east);
\node[anchor=west,text=white,font=\sffamily\bfseries\LARGE] at ([xshift=5mm,yshift=-0.62cm]frame.north west) {Question 7};
\node[anchor=east,text=white,font=\sffamily\itshape\large] at ([xshift=-5mm,yshift=-0.62cm]frame.north east) {Arithmétique et puissances};}]
\vspace*{1.75cm}
\begin{minipage}[t][6.9cm][c]{\linewidth}\LARGE\centering
\begin{minipage}{0.95\linewidth}\raggedright\linespread{1.25}\selectfont Quel est le plus grand diviseur commun à $84$ et $126$~?\end{minipage}
\end{minipage}

\vfill
\begin{tcolorbox}[enhanced,nobeforeafter,width=7.1cm,height=2.1cm,valign=center,colback=cari!6,colframe=cari,arc=2.5mm,boxrule=0.9pt,left=3mm]{\Large\bfseries\color{cari}a)}\quad\begin{minipage}[c]{5.3cm}\raggedright\Large $84$\end{minipage}\end{tcolorbox}\hfill\begin{tcolorbox}[enhanced,nobeforeafter,width=7.1cm,height=2.1cm,valign=center,colback=cari!6,colframe=cari,arc=2.5mm,boxrule=0.9pt,left=3mm]{\Large\bfseries\color{cari}b)}\quad\begin{minipage}[c]{5.3cm}\raggedright\Large $252$\end{minipage}\end{tcolorbox}\par\vspace{2.5mm}
\begin{tcolorbox}[enhanced,nobeforeafter,width=7.1cm,height=2.1cm,valign=center,colback=cari!6,colframe=cari,arc=2.5mm,boxrule=0.9pt,left=3mm]{\Large\bfseries\color{cari}c)}\quad\begin{minipage}[c]{5.3cm}\raggedright\Large $21$\end{minipage}\end{tcolorbox}\hfill\begin{tcolorbox}[enhanced,nobeforeafter,width=7.1cm,height=2.1cm,valign=center,colback=cari!6,colframe=cari,arc=2.5mm,boxrule=0.9pt,left=3mm]{\Large\bfseries\color{cari}d)}\quad\begin{minipage}[c]{5.3cm}\raggedright\Large $42$\end{minipage}\end{tcolorbox}
\end{tcolorbox}
\newpage
\begin{tcolorbox}[enhanced,width=16cm,height=14cm,colback=white,colframe=cari,boxrule=2pt,arc=0mm,
left=6mm,right=6mm,top=0mm,bottom=4mm,sharp corners,
overlay={\fill[cari] (frame.north west) rectangle ([yshift=-1.25cm]frame.north east);
\node[anchor=west,text=white,font=\sffamily\bfseries\LARGE] at ([xshift=5mm,yshift=-0.62cm]frame.north west) {Question 8};
\node[anchor=east,text=white,font=\sffamily\itshape\large] at ([xshift=-5mm,yshift=-0.62cm]frame.north east) {Arithmétique et puissances};}]
\vspace*{1.75cm}
\begin{minipage}[t][6.9cm][c]{\linewidth}\LARGE\centering
\begin{minipage}{0.95\linewidth}\raggedright\linespread{1.25}\selectfont Quel est le plus grand diviseur commun à $60$ et $90$~?\end{minipage}
\end{minipage}

\vfill
\begin{tcolorbox}[enhanced,nobeforeafter,width=7.1cm,height=2.1cm,valign=center,colback=cari!6,colframe=cari,arc=2.5mm,boxrule=0.9pt,left=3mm]{\Large\bfseries\color{cari}a)}\quad\begin{minipage}[c]{5.3cm}\raggedright\Large $60$\end{minipage}\end{tcolorbox}\hfill\begin{tcolorbox}[enhanced,nobeforeafter,width=7.1cm,height=2.1cm,valign=center,colback=cari!6,colframe=cari,arc=2.5mm,boxrule=0.9pt,left=3mm]{\Large\bfseries\color{cari}b)}\quad\begin{minipage}[c]{5.3cm}\raggedright\Large $15$\end{minipage}\end{tcolorbox}\par\vspace{2.5mm}
\begin{tcolorbox}[enhanced,nobeforeafter,width=7.1cm,height=2.1cm,valign=center,colback=cari!6,colframe=cari,arc=2.5mm,boxrule=0.9pt,left=3mm]{\Large\bfseries\color{cari}c)}\quad\begin{minipage}[c]{5.3cm}\raggedright\Large $180$\end{minipage}\end{tcolorbox}\hfill\begin{tcolorbox}[enhanced,nobeforeafter,width=7.1cm,height=2.1cm,valign=center,colback=cari!6,colframe=cari,arc=2.5mm,boxrule=0.9pt,left=3mm]{\Large\bfseries\color{cari}d)}\quad\begin{minipage}[c]{5.3cm}\raggedright\Large $30$\end{minipage}\end{tcolorbox}
\end{tcolorbox}
\newpage
\begin{tcolorbox}[enhanced,width=16cm,height=14cm,colback=white,colframe=cari,boxrule=2pt,arc=0mm,
left=6mm,right=6mm,top=0mm,bottom=4mm,sharp corners,
overlay={\fill[cari] (frame.north west) rectangle ([yshift=-1.25cm]frame.north east);
\node[anchor=west,text=white,font=\sffamily\bfseries\LARGE] at ([xshift=5mm,yshift=-0.62cm]frame.north west) {Question 9};
\node[anchor=east,text=white,font=\sffamily\itshape\large] at ([xshift=-5mm,yshift=-0.62cm]frame.north east) {Arithmétique et puissances};}]
\vspace*{1.75cm}
\begin{minipage}[t][6.9cm][c]{\linewidth}\Large\centering
\begin{minipage}{0.95\linewidth}\raggedright\linespread{1.25}\selectfont Un fleuriste a $60$ roses et $84$ tulipes. Il veut faire le plus grand nombre de bouquets identiques en utilisant toutes les fleurs. Combien de bouquets peut-il faire~?\end{minipage}
\end{minipage}

\vfill
\begin{tcolorbox}[enhanced,nobeforeafter,width=7.1cm,height=2.1cm,valign=center,colback=cari!6,colframe=cari,arc=2.5mm,boxrule=0.9pt,left=3mm]{\Large\bfseries\color{cari}a)}\quad\begin{minipage}[c]{5.3cm}\raggedright\Large $6$\end{minipage}\end{tcolorbox}\hfill\begin{tcolorbox}[enhanced,nobeforeafter,width=7.1cm,height=2.1cm,valign=center,colback=cari!6,colframe=cari,arc=2.5mm,boxrule=0.9pt,left=3mm]{\Large\bfseries\color{cari}b)}\quad\begin{minipage}[c]{5.3cm}\raggedright\Large $144$\end{minipage}\end{tcolorbox}\par\vspace{2.5mm}
\begin{tcolorbox}[enhanced,nobeforeafter,width=7.1cm,height=2.1cm,valign=center,colback=cari!6,colframe=cari,arc=2.5mm,boxrule=0.9pt,left=3mm]{\Large\bfseries\color{cari}c)}\quad\begin{minipage}[c]{5.3cm}\raggedright\Large $12$\end{minipage}\end{tcolorbox}\hfill\begin{tcolorbox}[enhanced,nobeforeafter,width=7.1cm,height=2.1cm,valign=center,colback=cari!6,colframe=cari,arc=2.5mm,boxrule=0.9pt,left=3mm]{\Large\bfseries\color{cari}d)}\quad\begin{minipage}[c]{5.3cm}\raggedright\Large $24$\end{minipage}\end{tcolorbox}
\end{tcolorbox}
\newpage
\begin{tcolorbox}[enhanced,width=16cm,height=14cm,colback=white,colframe=cari,boxrule=2pt,arc=0mm,
left=6mm,right=6mm,top=0mm,bottom=4mm,sharp corners,
overlay={\fill[cari] (frame.north west) rectangle ([yshift=-1.25cm]frame.north east);
\node[anchor=west,text=white,font=\sffamily\bfseries\LARGE] at ([xshift=5mm,yshift=-0.62cm]frame.north west) {Question 10};
\node[anchor=east,text=white,font=\sffamily\itshape\large] at ([xshift=-5mm,yshift=-0.62cm]frame.north east) {Arithmétique et puissances};}]
\vspace*{1.75cm}
\begin{minipage}[t][6.9cm][c]{\linewidth}\Large\centering
\begin{minipage}{0.95\linewidth}\raggedright\linespread{1.25}\selectfont Un fleuriste a $48$ roses et $72$ tulipes. Il veut faire le plus grand nombre de bouquets identiques en utilisant toutes les fleurs. Combien de bouquets peut-il faire~?\end{minipage}
\end{minipage}

\vfill
\begin{tcolorbox}[enhanced,nobeforeafter,width=7.1cm,height=2.1cm,valign=center,colback=cari!6,colframe=cari,arc=2.5mm,boxrule=0.9pt,left=3mm]{\Large\bfseries\color{cari}a)}\quad\begin{minipage}[c]{5.3cm}\raggedright\Large $120$\end{minipage}\end{tcolorbox}\hfill\begin{tcolorbox}[enhanced,nobeforeafter,width=7.1cm,height=2.1cm,valign=center,colback=cari!6,colframe=cari,arc=2.5mm,boxrule=0.9pt,left=3mm]{\Large\bfseries\color{cari}b)}\quad\begin{minipage}[c]{5.3cm}\raggedright\Large $48$\end{minipage}\end{tcolorbox}\par\vspace{2.5mm}
\begin{tcolorbox}[enhanced,nobeforeafter,width=7.1cm,height=2.1cm,valign=center,colback=cari!6,colframe=cari,arc=2.5mm,boxrule=0.9pt,left=3mm]{\Large\bfseries\color{cari}c)}\quad\begin{minipage}[c]{5.3cm}\raggedright\Large $12$\end{minipage}\end{tcolorbox}\hfill\begin{tcolorbox}[enhanced,nobeforeafter,width=7.1cm,height=2.1cm,valign=center,colback=cari!6,colframe=cari,arc=2.5mm,boxrule=0.9pt,left=3mm]{\Large\bfseries\color{cari}d)}\quad\begin{minipage}[c]{5.3cm}\raggedright\Large $24$\end{minipage}\end{tcolorbox}
\end{tcolorbox}
\newpage
\begin{tcolorbox}[enhanced,width=16cm,height=14cm,colback=white,colframe=cari,boxrule=2pt,arc=0mm,
left=6mm,right=6mm,top=0mm,bottom=4mm,sharp corners,
overlay={\fill[cari] (frame.north west) rectangle ([yshift=-1.25cm]frame.north east);
\node[anchor=west,text=white,font=\sffamily\bfseries\LARGE] at ([xshift=5mm,yshift=-0.62cm]frame.north west) {Question 11};
\node[anchor=east,text=white,font=\sffamily\itshape\large] at ([xshift=-5mm,yshift=-0.62cm]frame.north east) {Arithmétique et puissances};}]
\vspace*{1.75cm}
\begin{minipage}[t][6.9cm][c]{\linewidth}\Large\centering
\begin{minipage}{0.95\linewidth}\raggedright\linespread{1.25}\selectfont Un phare clignote toutes les $6$ secondes et un autre toutes les $8$ secondes. Ils viennent de clignoter ensemble. Dans combien de secondes clignoteront-ils de nouveau ensemble~?\end{minipage}
\end{minipage}

\vfill
\begin{tcolorbox}[enhanced,nobeforeafter,width=7.1cm,height=2.1cm,valign=center,colback=cari!6,colframe=cari,arc=2.5mm,boxrule=0.9pt,left=3mm]{\Large\bfseries\color{cari}a)}\quad\begin{minipage}[c]{5.3cm}\raggedright\Large $24~\text{s}$\end{minipage}\end{tcolorbox}\hfill\begin{tcolorbox}[enhanced,nobeforeafter,width=7.1cm,height=2.1cm,valign=center,colback=cari!6,colframe=cari,arc=2.5mm,boxrule=0.9pt,left=3mm]{\Large\bfseries\color{cari}b)}\quad\begin{minipage}[c]{5.3cm}\raggedright\Large $14~\text{s}$\end{minipage}\end{tcolorbox}\par\vspace{2.5mm}
\begin{tcolorbox}[enhanced,nobeforeafter,width=7.1cm,height=2.1cm,valign=center,colback=cari!6,colframe=cari,arc=2.5mm,boxrule=0.9pt,left=3mm]{\Large\bfseries\color{cari}c)}\quad\begin{minipage}[c]{5.3cm}\raggedright\Large $2~\text{s}$\end{minipage}\end{tcolorbox}\hfill\begin{tcolorbox}[enhanced,nobeforeafter,width=7.1cm,height=2.1cm,valign=center,colback=cari!6,colframe=cari,arc=2.5mm,boxrule=0.9pt,left=3mm]{\Large\bfseries\color{cari}d)}\quad\begin{minipage}[c]{5.3cm}\raggedright\Large $48~\text{s}$\end{minipage}\end{tcolorbox}
\end{tcolorbox}
\newpage
\begin{tcolorbox}[enhanced,width=16cm,height=14cm,colback=white,colframe=cari,boxrule=2pt,arc=0mm,
left=6mm,right=6mm,top=0mm,bottom=4mm,sharp corners,
overlay={\fill[cari] (frame.north west) rectangle ([yshift=-1.25cm]frame.north east);
\node[anchor=west,text=white,font=\sffamily\bfseries\LARGE] at ([xshift=5mm,yshift=-0.62cm]frame.north west) {Question 12};
\node[anchor=east,text=white,font=\sffamily\itshape\large] at ([xshift=-5mm,yshift=-0.62cm]frame.north east) {Arithmétique et puissances};}]
\vspace*{1.75cm}
\begin{minipage}[t][6.9cm][c]{\linewidth}\Large\centering
\begin{minipage}{0.95\linewidth}\raggedright\linespread{1.25}\selectfont Un phare clignote toutes les $4$ secondes et un autre toutes les $10$ secondes. Ils viennent de clignoter ensemble. Dans combien de secondes clignoteront-ils de nouveau ensemble~?\end{minipage}
\end{minipage}

\vfill
\begin{tcolorbox}[enhanced,nobeforeafter,width=7.1cm,height=2.1cm,valign=center,colback=cari!6,colframe=cari,arc=2.5mm,boxrule=0.9pt,left=3mm]{\Large\bfseries\color{cari}a)}\quad\begin{minipage}[c]{5.3cm}\raggedright\Large $40~\text{s}$\end{minipage}\end{tcolorbox}\hfill\begin{tcolorbox}[enhanced,nobeforeafter,width=7.1cm,height=2.1cm,valign=center,colback=cari!6,colframe=cari,arc=2.5mm,boxrule=0.9pt,left=3mm]{\Large\bfseries\color{cari}b)}\quad\begin{minipage}[c]{5.3cm}\raggedright\Large $20~\text{s}$\end{minipage}\end{tcolorbox}\par\vspace{2.5mm}
\begin{tcolorbox}[enhanced,nobeforeafter,width=7.1cm,height=2.1cm,valign=center,colback=cari!6,colframe=cari,arc=2.5mm,boxrule=0.9pt,left=3mm]{\Large\bfseries\color{cari}c)}\quad\begin{minipage}[c]{5.3cm}\raggedright\Large $2~\text{s}$\end{minipage}\end{tcolorbox}\hfill\begin{tcolorbox}[enhanced,nobeforeafter,width=7.1cm,height=2.1cm,valign=center,colback=cari!6,colframe=cari,arc=2.5mm,boxrule=0.9pt,left=3mm]{\Large\bfseries\color{cari}d)}\quad\begin{minipage}[c]{5.3cm}\raggedright\Large $14~\text{s}$\end{minipage}\end{tcolorbox}
\end{tcolorbox}
\newpage
\begin{tcolorbox}[enhanced,width=16cm,height=14cm,colback=white,colframe=cari,boxrule=2pt,arc=0mm,
left=6mm,right=6mm,top=0mm,bottom=4mm,sharp corners,
overlay={\fill[cari] (frame.north west) rectangle ([yshift=-1.25cm]frame.north east);
\node[anchor=west,text=white,font=\sffamily\bfseries\LARGE] at ([xshift=5mm,yshift=-0.62cm]frame.north west) {Question 13};
\node[anchor=east,text=white,font=\sffamily\itshape\large] at ([xshift=-5mm,yshift=-0.62cm]frame.north east) {Arithmétique et puissances};}]
\vspace*{1.75cm}
\begin{minipage}[t][6.9cm][c]{\linewidth}\LARGE\centering
\begin{minipage}{0.95\linewidth}\raggedright\linespread{1.25}\selectfont Que vaut $(3\times10^{5})\times(2\times10^{-3})$, en notation scientifique~?\end{minipage}
\end{minipage}

\vfill
\begin{tcolorbox}[enhanced,nobeforeafter,width=7.1cm,height=2.1cm,valign=center,colback=cari!6,colframe=cari,arc=2.5mm,boxrule=0.9pt,left=3mm]{\Large\bfseries\color{cari}a)}\quad\begin{minipage}[c]{5.3cm}\raggedright\Large $60\times10^{1}$\end{minipage}\end{tcolorbox}\hfill\begin{tcolorbox}[enhanced,nobeforeafter,width=7.1cm,height=2.1cm,valign=center,colback=cari!6,colframe=cari,arc=2.5mm,boxrule=0.9pt,left=3mm]{\Large\bfseries\color{cari}b)}\quad\begin{minipage}[c]{5.3cm}\raggedright\Large $6\times10^{1}$\end{minipage}\end{tcolorbox}\par\vspace{2.5mm}
\begin{tcolorbox}[enhanced,nobeforeafter,width=7.1cm,height=2.1cm,valign=center,colback=cari!6,colframe=cari,arc=2.5mm,boxrule=0.9pt,left=3mm]{\Large\bfseries\color{cari}c)}\quad\begin{minipage}[c]{5.3cm}\raggedright\Large $6\times10^{2}$\end{minipage}\end{tcolorbox}\hfill\begin{tcolorbox}[enhanced,nobeforeafter,width=7.1cm,height=2.1cm,valign=center,colback=cari!6,colframe=cari,arc=2.5mm,boxrule=0.9pt,left=3mm]{\Large\bfseries\color{cari}d)}\quad\begin{minipage}[c]{5.3cm}\raggedright\Large $6\times10^{-15}$\end{minipage}\end{tcolorbox}
\end{tcolorbox}
\newpage
\begin{tcolorbox}[enhanced,width=16cm,height=14cm,colback=white,colframe=cari,boxrule=2pt,arc=0mm,
left=6mm,right=6mm,top=0mm,bottom=4mm,sharp corners,
overlay={\fill[cari] (frame.north west) rectangle ([yshift=-1.25cm]frame.north east);
\node[anchor=west,text=white,font=\sffamily\bfseries\LARGE] at ([xshift=5mm,yshift=-0.62cm]frame.north west) {Question 14};
\node[anchor=east,text=white,font=\sffamily\itshape\large] at ([xshift=-5mm,yshift=-0.62cm]frame.north east) {Arithmétique et puissances};}]
\vspace*{1.75cm}
\begin{minipage}[t][6.9cm][c]{\linewidth}\LARGE\centering
\begin{minipage}{0.95\linewidth}\raggedright\linespread{1.25}\selectfont Que vaut $(4\times10^{-2})\times(2\times10^{6})$, en notation scientifique~?\end{minipage}
\end{minipage}

\vfill
\begin{tcolorbox}[enhanced,nobeforeafter,width=7.1cm,height=2.1cm,valign=center,colback=cari!6,colframe=cari,arc=2.5mm,boxrule=0.9pt,left=3mm]{\Large\bfseries\color{cari}a)}\quad\begin{minipage}[c]{5.3cm}\raggedright\Large $80\times10^{3}$\end{minipage}\end{tcolorbox}\hfill\begin{tcolorbox}[enhanced,nobeforeafter,width=7.1cm,height=2.1cm,valign=center,colback=cari!6,colframe=cari,arc=2.5mm,boxrule=0.9pt,left=3mm]{\Large\bfseries\color{cari}b)}\quad\begin{minipage}[c]{5.3cm}\raggedright\Large $8\times10^{-12}$\end{minipage}\end{tcolorbox}\par\vspace{2.5mm}
\begin{tcolorbox}[enhanced,nobeforeafter,width=7.1cm,height=2.1cm,valign=center,colback=cari!6,colframe=cari,arc=2.5mm,boxrule=0.9pt,left=3mm]{\Large\bfseries\color{cari}c)}\quad\begin{minipage}[c]{5.3cm}\raggedright\Large $8\times10^{4}$\end{minipage}\end{tcolorbox}\hfill\begin{tcolorbox}[enhanced,nobeforeafter,width=7.1cm,height=2.1cm,valign=center,colback=cari!6,colframe=cari,arc=2.5mm,boxrule=0.9pt,left=3mm]{\Large\bfseries\color{cari}d)}\quad\begin{minipage}[c]{5.3cm}\raggedright\Large $8\times10^{3}$\end{minipage}\end{tcolorbox}
\end{tcolorbox}
\newpage
\begin{tcolorbox}[enhanced,width=16cm,height=14cm,colback=white,colframe=cari,boxrule=2pt,arc=0mm,
left=6mm,right=6mm,top=0mm,bottom=4mm,sharp corners,
overlay={\fill[cari] (frame.north west) rectangle ([yshift=-1.25cm]frame.north east);
\node[anchor=west,text=white,font=\sffamily\bfseries\LARGE] at ([xshift=5mm,yshift=-0.62cm]frame.north west) {Question 15};
\node[anchor=east,text=white,font=\sffamily\itshape\large] at ([xshift=-5mm,yshift=-0.62cm]frame.north east) {Arithmétique et puissances};}]
\vspace*{1.75cm}
\begin{minipage}[t][6.9cm][c]{\linewidth}\LARGE\centering
\begin{minipage}{0.95\linewidth}\raggedright\linespread{1.25}\selectfont Quelle est l'écriture décimale de $5^{-2}$~?\end{minipage}
\end{minipage}

\vfill
\begin{tcolorbox}[enhanced,nobeforeafter,width=7.1cm,height=2.1cm,valign=center,colback=cari!6,colframe=cari,arc=2.5mm,boxrule=0.9pt,left=3mm]{\Large\bfseries\color{cari}a)}\quad\begin{minipage}[c]{5.3cm}\raggedright\Large $-25$\end{minipage}\end{tcolorbox}\hfill\begin{tcolorbox}[enhanced,nobeforeafter,width=7.1cm,height=2.1cm,valign=center,colback=cari!6,colframe=cari,arc=2.5mm,boxrule=0.9pt,left=3mm]{\Large\bfseries\color{cari}b)}\quad\begin{minipage}[c]{5.3cm}\raggedright\Large $0{,}1$\end{minipage}\end{tcolorbox}\par\vspace{2.5mm}
\begin{tcolorbox}[enhanced,nobeforeafter,width=7.1cm,height=2.1cm,valign=center,colback=cari!6,colframe=cari,arc=2.5mm,boxrule=0.9pt,left=3mm]{\Large\bfseries\color{cari}c)}\quad\begin{minipage}[c]{5.3cm}\raggedright\Large $-10$\end{minipage}\end{tcolorbox}\hfill\begin{tcolorbox}[enhanced,nobeforeafter,width=7.1cm,height=2.1cm,valign=center,colback=cari!6,colframe=cari,arc=2.5mm,boxrule=0.9pt,left=3mm]{\Large\bfseries\color{cari}d)}\quad\begin{minipage}[c]{5.3cm}\raggedright\Large $0{,}04$\end{minipage}\end{tcolorbox}
\end{tcolorbox}
\newpage
\begin{tcolorbox}[enhanced,width=16cm,height=14cm,colback=white,colframe=cari,boxrule=2pt,arc=0mm,
left=6mm,right=6mm,top=0mm,bottom=4mm,sharp corners,
overlay={\fill[cari] (frame.north west) rectangle ([yshift=-1.25cm]frame.north east);
\node[anchor=west,text=white,font=\sffamily\bfseries\LARGE] at ([xshift=5mm,yshift=-0.62cm]frame.north west) {Question 16};
\node[anchor=east,text=white,font=\sffamily\itshape\large] at ([xshift=-5mm,yshift=-0.62cm]frame.north east) {Arithmétique et puissances};}]
\vspace*{1.75cm}
\begin{minipage}[t][6.9cm][c]{\linewidth}\LARGE\centering
\begin{minipage}{0.95\linewidth}\raggedright\linespread{1.25}\selectfont Quelle est l'écriture décimale de $2^{-3}$~?\end{minipage}
\end{minipage}

\vfill
\begin{tcolorbox}[enhanced,nobeforeafter,width=7.1cm,height=2.1cm,valign=center,colback=cari!6,colframe=cari,arc=2.5mm,boxrule=0.9pt,left=3mm]{\Large\bfseries\color{cari}a)}\quad\begin{minipage}[c]{5.3cm}\raggedright\Large $-8$\end{minipage}\end{tcolorbox}\hfill\begin{tcolorbox}[enhanced,nobeforeafter,width=7.1cm,height=2.1cm,valign=center,colback=cari!6,colframe=cari,arc=2.5mm,boxrule=0.9pt,left=3mm]{\Large\bfseries\color{cari}b)}\quad\begin{minipage}[c]{5.3cm}\raggedright\Large $\dfrac{1}{6}$\end{minipage}\end{tcolorbox}\par\vspace{2.5mm}
\begin{tcolorbox}[enhanced,nobeforeafter,width=7.1cm,height=2.1cm,valign=center,colback=cari!6,colframe=cari,arc=2.5mm,boxrule=0.9pt,left=3mm]{\Large\bfseries\color{cari}c)}\quad\begin{minipage}[c]{5.3cm}\raggedright\Large $0{,}125$\end{minipage}\end{tcolorbox}\hfill\begin{tcolorbox}[enhanced,nobeforeafter,width=7.1cm,height=2.1cm,valign=center,colback=cari!6,colframe=cari,arc=2.5mm,boxrule=0.9pt,left=3mm]{\Large\bfseries\color{cari}d)}\quad\begin{minipage}[c]{5.3cm}\raggedright\Large $-6$\end{minipage}\end{tcolorbox}
\end{tcolorbox}
\newpage
\begin{tcolorbox}[enhanced,width=16cm,height=14cm,colback=white,colframe=clit,boxrule=2pt,arc=0mm,
left=6mm,right=6mm,top=0mm,bottom=4mm,sharp corners,
overlay={\fill[clit] (frame.north west) rectangle ([yshift=-1.25cm]frame.north east);
\node[anchor=west,text=white,font=\sffamily\bfseries\LARGE] at ([xshift=5mm,yshift=-0.62cm]frame.north west) {Question 17};
\node[anchor=east,text=white,font=\sffamily\itshape\large] at ([xshift=-5mm,yshift=-0.62cm]frame.north east) {Calcul littéral et équations};}]
\vspace*{1.75cm}
\begin{minipage}[t][6.9cm][c]{\linewidth}\LARGE\centering
\begin{minipage}{0.95\linewidth}\raggedright\linespread{1.25}\selectfont Développer et réduire $(2x+3)(x-4)$.\end{minipage}
\end{minipage}

\vfill
\begin{tcolorbox}[enhanced,nobeforeafter,width=7.1cm,height=2.1cm,valign=center,colback=clit!6,colframe=clit,arc=2.5mm,boxrule=0.9pt,left=3mm]{\Large\bfseries\color{clit}a)}\quad\begin{minipage}[c]{5.3cm}\raggedright\Large $2x^2-11x-12$\end{minipage}\end{tcolorbox}\hfill\begin{tcolorbox}[enhanced,nobeforeafter,width=7.1cm,height=2.1cm,valign=center,colback=clit!6,colframe=clit,arc=2.5mm,boxrule=0.9pt,left=3mm]{\Large\bfseries\color{clit}b)}\quad\begin{minipage}[c]{5.3cm}\raggedright\Large $2x^2-5x-12$\end{minipage}\end{tcolorbox}\par\vspace{2.5mm}
\begin{tcolorbox}[enhanced,nobeforeafter,width=7.1cm,height=2.1cm,valign=center,colback=clit!6,colframe=clit,arc=2.5mm,boxrule=0.9pt,left=3mm]{\Large\bfseries\color{clit}c)}\quad\begin{minipage}[c]{5.3cm}\raggedright\Large $2x^2-12$\end{minipage}\end{tcolorbox}\hfill\begin{tcolorbox}[enhanced,nobeforeafter,width=7.1cm,height=2.1cm,valign=center,colback=clit!6,colframe=clit,arc=2.5mm,boxrule=0.9pt,left=3mm]{\Large\bfseries\color{clit}d)}\quad\begin{minipage}[c]{5.3cm}\raggedright\Large $3x^2-5x-12$\end{minipage}\end{tcolorbox}
\end{tcolorbox}
\newpage
\begin{tcolorbox}[enhanced,width=16cm,height=14cm,colback=white,colframe=clit,boxrule=2pt,arc=0mm,
left=6mm,right=6mm,top=0mm,bottom=4mm,sharp corners,
overlay={\fill[clit] (frame.north west) rectangle ([yshift=-1.25cm]frame.north east);
\node[anchor=west,text=white,font=\sffamily\bfseries\LARGE] at ([xshift=5mm,yshift=-0.62cm]frame.north west) {Question 18};
\node[anchor=east,text=white,font=\sffamily\itshape\large] at ([xshift=-5mm,yshift=-0.62cm]frame.north east) {Calcul littéral et équations};}]
\vspace*{1.75cm}
\begin{minipage}[t][6.9cm][c]{\linewidth}\LARGE\centering
\begin{minipage}{0.95\linewidth}\raggedright\linespread{1.25}\selectfont Développer et réduire $(3x-1)(2x+5)$.\end{minipage}
\end{minipage}

\vfill
\begin{tcolorbox}[enhanced,nobeforeafter,width=7.1cm,height=2.1cm,valign=center,colback=clit!6,colframe=clit,arc=2.5mm,boxrule=0.9pt,left=3mm]{\Large\bfseries\color{clit}a)}\quad\begin{minipage}[c]{5.3cm}\raggedright\Large $6x^2+17x-5$\end{minipage}\end{tcolorbox}\hfill\begin{tcolorbox}[enhanced,nobeforeafter,width=7.1cm,height=2.1cm,valign=center,colback=clit!6,colframe=clit,arc=2.5mm,boxrule=0.9pt,left=3mm]{\Large\bfseries\color{clit}b)}\quad\begin{minipage}[c]{5.3cm}\raggedright\Large $6x^2-5$\end{minipage}\end{tcolorbox}\par\vspace{2.5mm}
\begin{tcolorbox}[enhanced,nobeforeafter,width=7.1cm,height=2.1cm,valign=center,colback=clit!6,colframe=clit,arc=2.5mm,boxrule=0.9pt,left=3mm]{\Large\bfseries\color{clit}c)}\quad\begin{minipage}[c]{5.3cm}\raggedright\Large $5x^2+13x-5$\end{minipage}\end{tcolorbox}\hfill\begin{tcolorbox}[enhanced,nobeforeafter,width=7.1cm,height=2.1cm,valign=center,colback=clit!6,colframe=clit,arc=2.5mm,boxrule=0.9pt,left=3mm]{\Large\bfseries\color{clit}d)}\quad\begin{minipage}[c]{5.3cm}\raggedright\Large $6x^2+13x-5$\end{minipage}\end{tcolorbox}
\end{tcolorbox}
\newpage
\begin{tcolorbox}[enhanced,width=16cm,height=14cm,colback=white,colframe=clit,boxrule=2pt,arc=0mm,
left=6mm,right=6mm,top=0mm,bottom=4mm,sharp corners,
overlay={\fill[clit] (frame.north west) rectangle ([yshift=-1.25cm]frame.north east);
\node[anchor=west,text=white,font=\sffamily\bfseries\LARGE] at ([xshift=5mm,yshift=-0.62cm]frame.north west) {Question 19};
\node[anchor=east,text=white,font=\sffamily\itshape\large] at ([xshift=-5mm,yshift=-0.62cm]frame.north east) {Calcul littéral et équations};}]
\vspace*{1.75cm}
\begin{minipage}[t][6.9cm][c]{\linewidth}\LARGE\centering
\begin{minipage}{0.95\linewidth}\raggedright\linespread{1.25}\selectfont Factoriser $x^2-36$.\end{minipage}
\end{minipage}

\vfill
\begin{tcolorbox}[enhanced,nobeforeafter,width=7.1cm,height=2.1cm,valign=center,colback=clit!6,colframe=clit,arc=2.5mm,boxrule=0.9pt,left=3mm]{\Large\bfseries\color{clit}a)}\quad\begin{minipage}[c]{5.3cm}\raggedright\Large $(x-36)(x+36)$\end{minipage}\end{tcolorbox}\hfill\begin{tcolorbox}[enhanced,nobeforeafter,width=7.1cm,height=2.1cm,valign=center,colback=clit!6,colframe=clit,arc=2.5mm,boxrule=0.9pt,left=3mm]{\Large\bfseries\color{clit}b)}\quad\begin{minipage}[c]{5.3cm}\raggedright\Large $(x-6)^2$\end{minipage}\end{tcolorbox}\par\vspace{2.5mm}
\begin{tcolorbox}[enhanced,nobeforeafter,width=7.1cm,height=2.1cm,valign=center,colback=clit!6,colframe=clit,arc=2.5mm,boxrule=0.9pt,left=3mm]{\Large\bfseries\color{clit}c)}\quad\begin{minipage}[c]{5.3cm}\raggedright\Large $(x+6)^2$\end{minipage}\end{tcolorbox}\hfill\begin{tcolorbox}[enhanced,nobeforeafter,width=7.1cm,height=2.1cm,valign=center,colback=clit!6,colframe=clit,arc=2.5mm,boxrule=0.9pt,left=3mm]{\Large\bfseries\color{clit}d)}\quad\begin{minipage}[c]{5.3cm}\raggedright\Large $(x-6)(x+6)$\end{minipage}\end{tcolorbox}
\end{tcolorbox}
\newpage
\begin{tcolorbox}[enhanced,width=16cm,height=14cm,colback=white,colframe=clit,boxrule=2pt,arc=0mm,
left=6mm,right=6mm,top=0mm,bottom=4mm,sharp corners,
overlay={\fill[clit] (frame.north west) rectangle ([yshift=-1.25cm]frame.north east);
\node[anchor=west,text=white,font=\sffamily\bfseries\LARGE] at ([xshift=5mm,yshift=-0.62cm]frame.north west) {Question 20};
\node[anchor=east,text=white,font=\sffamily\itshape\large] at ([xshift=-5mm,yshift=-0.62cm]frame.north east) {Calcul littéral et équations};}]
\vspace*{1.75cm}
\begin{minipage}[t][6.9cm][c]{\linewidth}\LARGE\centering
\begin{minipage}{0.95\linewidth}\raggedright\linespread{1.25}\selectfont Factoriser $9x^2-4$.\end{minipage}
\end{minipage}

\vfill
\begin{tcolorbox}[enhanced,nobeforeafter,width=7.1cm,height=2.1cm,valign=center,colback=clit!6,colframe=clit,arc=2.5mm,boxrule=0.9pt,left=3mm]{\Large\bfseries\color{clit}a)}\quad\begin{minipage}[c]{5.3cm}\raggedright\Large $(3x+2)^2$\end{minipage}\end{tcolorbox}\hfill\begin{tcolorbox}[enhanced,nobeforeafter,width=7.1cm,height=2.1cm,valign=center,colback=clit!6,colframe=clit,arc=2.5mm,boxrule=0.9pt,left=3mm]{\Large\bfseries\color{clit}b)}\quad\begin{minipage}[c]{5.3cm}\raggedright\Large $(3x-2)(3x+2)$\end{minipage}\end{tcolorbox}\par\vspace{2.5mm}
\begin{tcolorbox}[enhanced,nobeforeafter,width=7.1cm,height=2.1cm,valign=center,colback=clit!6,colframe=clit,arc=2.5mm,boxrule=0.9pt,left=3mm]{\Large\bfseries\color{clit}c)}\quad\begin{minipage}[c]{5.3cm}\raggedright\Large $(3x-2)^2$\end{minipage}\end{tcolorbox}\hfill\begin{tcolorbox}[enhanced,nobeforeafter,width=7.1cm,height=2.1cm,valign=center,colback=clit!6,colframe=clit,arc=2.5mm,boxrule=0.9pt,left=3mm]{\Large\bfseries\color{clit}d)}\quad\begin{minipage}[c]{5.3cm}\raggedright\Large $(3x-4)(3x+4)$\end{minipage}\end{tcolorbox}
\end{tcolorbox}
\newpage
\begin{tcolorbox}[enhanced,width=16cm,height=14cm,colback=white,colframe=clit,boxrule=2pt,arc=0mm,
left=6mm,right=6mm,top=0mm,bottom=4mm,sharp corners,
overlay={\fill[clit] (frame.north west) rectangle ([yshift=-1.25cm]frame.north east);
\node[anchor=west,text=white,font=\sffamily\bfseries\LARGE] at ([xshift=5mm,yshift=-0.62cm]frame.north west) {Question 21};
\node[anchor=east,text=white,font=\sffamily\itshape\large] at ([xshift=-5mm,yshift=-0.62cm]frame.north east) {Calcul littéral et équations};}]
\vspace*{1.75cm}
\begin{minipage}[t][6.9cm][c]{\linewidth}\LARGE\centering
\begin{minipage}{0.95\linewidth}\raggedright\linespread{1.25}\selectfont Quelles sont les solutions de l'équation $(x-3)(2x+5)=0$~?\end{minipage}
\end{minipage}

\vfill
\begin{tcolorbox}[enhanced,nobeforeafter,width=7.1cm,height=2.1cm,valign=center,colback=clit!6,colframe=clit,arc=2.5mm,boxrule=0.9pt,left=3mm]{\Large\bfseries\color{clit}a)}\quad\begin{minipage}[c]{5.3cm}\raggedright\Large $\{-\dfrac{5}{2}\,;\,3\}$\end{minipage}\end{tcolorbox}\hfill\begin{tcolorbox}[enhanced,nobeforeafter,width=7.1cm,height=2.1cm,valign=center,colback=clit!6,colframe=clit,arc=2.5mm,boxrule=0.9pt,left=3mm]{\Large\bfseries\color{clit}b)}\quad\begin{minipage}[c]{5.3cm}\raggedright\Large $\{3\}$\end{minipage}\end{tcolorbox}\par\vspace{2.5mm}
\begin{tcolorbox}[enhanced,nobeforeafter,width=7.1cm,height=2.1cm,valign=center,colback=clit!6,colframe=clit,arc=2.5mm,boxrule=0.9pt,left=3mm]{\Large\bfseries\color{clit}c)}\quad\begin{minipage}[c]{5.3cm}\raggedright\Large $\{-3\,;\,5\}$\end{minipage}\end{tcolorbox}\hfill\begin{tcolorbox}[enhanced,nobeforeafter,width=7.1cm,height=2.1cm,valign=center,colback=clit!6,colframe=clit,arc=2.5mm,boxrule=0.9pt,left=3mm]{\Large\bfseries\color{clit}d)}\quad\begin{minipage}[c]{5.3cm}\raggedright\Large $\{-3\,;\,\dfrac{5}{2}\}$\end{minipage}\end{tcolorbox}
\end{tcolorbox}
\newpage
\begin{tcolorbox}[enhanced,width=16cm,height=14cm,colback=white,colframe=clit,boxrule=2pt,arc=0mm,
left=6mm,right=6mm,top=0mm,bottom=4mm,sharp corners,
overlay={\fill[clit] (frame.north west) rectangle ([yshift=-1.25cm]frame.north east);
\node[anchor=west,text=white,font=\sffamily\bfseries\LARGE] at ([xshift=5mm,yshift=-0.62cm]frame.north west) {Question 22};
\node[anchor=east,text=white,font=\sffamily\itshape\large] at ([xshift=-5mm,yshift=-0.62cm]frame.north east) {Calcul littéral et équations};}]
\vspace*{1.75cm}
\begin{minipage}[t][6.9cm][c]{\linewidth}\LARGE\centering
\begin{minipage}{0.95\linewidth}\raggedright\linespread{1.25}\selectfont Quelles sont les solutions de l'équation $(x+4)(x-1)=0$~?\end{minipage}
\end{minipage}

\vfill
\begin{tcolorbox}[enhanced,nobeforeafter,width=7.1cm,height=2.1cm,valign=center,colback=clit!6,colframe=clit,arc=2.5mm,boxrule=0.9pt,left=3mm]{\Large\bfseries\color{clit}a)}\quad\begin{minipage}[c]{5.3cm}\raggedright\Large $\{-\dfrac{1}{4}\,;\,1\}$\end{minipage}\end{tcolorbox}\hfill\begin{tcolorbox}[enhanced,nobeforeafter,width=7.1cm,height=2.1cm,valign=center,colback=clit!6,colframe=clit,arc=2.5mm,boxrule=0.9pt,left=3mm]{\Large\bfseries\color{clit}b)}\quad\begin{minipage}[c]{5.3cm}\raggedright\Large $\{-4\}$\end{minipage}\end{tcolorbox}\par\vspace{2.5mm}
\begin{tcolorbox}[enhanced,nobeforeafter,width=7.1cm,height=2.1cm,valign=center,colback=clit!6,colframe=clit,arc=2.5mm,boxrule=0.9pt,left=3mm]{\Large\bfseries\color{clit}c)}\quad\begin{minipage}[c]{5.3cm}\raggedright\Large $\{-1\,;\,4\}$\end{minipage}\end{tcolorbox}\hfill\begin{tcolorbox}[enhanced,nobeforeafter,width=7.1cm,height=2.1cm,valign=center,colback=clit!6,colframe=clit,arc=2.5mm,boxrule=0.9pt,left=3mm]{\Large\bfseries\color{clit}d)}\quad\begin{minipage}[c]{5.3cm}\raggedright\Large $\{-4\,;\,1\}$\end{minipage}\end{tcolorbox}
\end{tcolorbox}
\newpage
\begin{tcolorbox}[enhanced,width=16cm,height=14cm,colback=white,colframe=clit,boxrule=2pt,arc=0mm,
left=6mm,right=6mm,top=0mm,bottom=4mm,sharp corners,
overlay={\fill[clit] (frame.north west) rectangle ([yshift=-1.25cm]frame.north east);
\node[anchor=west,text=white,font=\sffamily\bfseries\LARGE] at ([xshift=5mm,yshift=-0.62cm]frame.north west) {Question 23};
\node[anchor=east,text=white,font=\sffamily\itshape\large] at ([xshift=-5mm,yshift=-0.62cm]frame.north east) {Calcul littéral et équations};}]
\vspace*{1.75cm}
\begin{minipage}[t][6.9cm][c]{\linewidth}\LARGE\centering
\begin{minipage}{0.95\linewidth}\raggedright\linespread{1.25}\selectfont Factoriser $(x+1)(2x-3)+(x+1)(x+5)$.\end{minipage}
\end{minipage}

\vfill
\begin{tcolorbox}[enhanced,nobeforeafter,width=7.1cm,height=2.1cm,valign=center,colback=clit!6,colframe=clit,arc=2.5mm,boxrule=0.9pt,left=3mm]{\Large\bfseries\color{clit}a)}\quad\begin{minipage}[c]{5.3cm}\raggedright\Large $(x+1)(x-8)$\end{minipage}\end{tcolorbox}\hfill\begin{tcolorbox}[enhanced,nobeforeafter,width=7.1cm,height=2.1cm,valign=center,colback=clit!6,colframe=clit,arc=2.5mm,boxrule=0.9pt,left=3mm]{\Large\bfseries\color{clit}b)}\quad\begin{minipage}[c]{5.3cm}\raggedright\Large $(x+1)(3x+2)$\end{minipage}\end{tcolorbox}\par\vspace{2.5mm}
\begin{tcolorbox}[enhanced,nobeforeafter,width=7.1cm,height=2.1cm,valign=center,colback=clit!6,colframe=clit,arc=2.5mm,boxrule=0.9pt,left=3mm]{\Large\bfseries\color{clit}c)}\quad\begin{minipage}[c]{5.3cm}\raggedright\Large $(x+1)(2x-15)$\end{minipage}\end{tcolorbox}\hfill\begin{tcolorbox}[enhanced,nobeforeafter,width=7.1cm,height=2.1cm,valign=center,colback=clit!6,colframe=clit,arc=2.5mm,boxrule=0.9pt,left=3mm]{\Large\bfseries\color{clit}d)}\quad\begin{minipage}[c]{5.3cm}\raggedright\Large $(x+1)^2(3x+2)$\end{minipage}\end{tcolorbox}
\end{tcolorbox}
\newpage
\begin{tcolorbox}[enhanced,width=16cm,height=14cm,colback=white,colframe=clit,boxrule=2pt,arc=0mm,
left=6mm,right=6mm,top=0mm,bottom=4mm,sharp corners,
overlay={\fill[clit] (frame.north west) rectangle ([yshift=-1.25cm]frame.north east);
\node[anchor=west,text=white,font=\sffamily\bfseries\LARGE] at ([xshift=5mm,yshift=-0.62cm]frame.north west) {Question 24};
\node[anchor=east,text=white,font=\sffamily\itshape\large] at ([xshift=-5mm,yshift=-0.62cm]frame.north east) {Calcul littéral et équations};}]
\vspace*{1.75cm}
\begin{minipage}[t][6.9cm][c]{\linewidth}\LARGE\centering
\begin{minipage}{0.95\linewidth}\raggedright\linespread{1.25}\selectfont Factoriser $(x-2)(3x+1)+(x-2)(x+4)$.\end{minipage}
\end{minipage}

\vfill
\begin{tcolorbox}[enhanced,nobeforeafter,width=7.1cm,height=2.1cm,valign=center,colback=clit!6,colframe=clit,arc=2.5mm,boxrule=0.9pt,left=3mm]{\Large\bfseries\color{clit}a)}\quad\begin{minipage}[c]{5.3cm}\raggedright\Large $(x-2)(3x+4)$\end{minipage}\end{tcolorbox}\hfill\begin{tcolorbox}[enhanced,nobeforeafter,width=7.1cm,height=2.1cm,valign=center,colback=clit!6,colframe=clit,arc=2.5mm,boxrule=0.9pt,left=3mm]{\Large\bfseries\color{clit}b)}\quad\begin{minipage}[c]{5.3cm}\raggedright\Large $(x-2)(4x+5)$\end{minipage}\end{tcolorbox}\par\vspace{2.5mm}
\begin{tcolorbox}[enhanced,nobeforeafter,width=7.1cm,height=2.1cm,valign=center,colback=clit!6,colframe=clit,arc=2.5mm,boxrule=0.9pt,left=3mm]{\Large\bfseries\color{clit}c)}\quad\begin{minipage}[c]{5.3cm}\raggedright\Large $(x-2)(2x-3)$\end{minipage}\end{tcolorbox}\hfill\begin{tcolorbox}[enhanced,nobeforeafter,width=7.1cm,height=2.1cm,valign=center,colback=clit!6,colframe=clit,arc=2.5mm,boxrule=0.9pt,left=3mm]{\Large\bfseries\color{clit}d)}\quad\begin{minipage}[c]{5.3cm}\raggedright\Large $(x-2)^2(4x+5)$\end{minipage}\end{tcolorbox}
\end{tcolorbox}
\newpage
\begin{tcolorbox}[enhanced,width=16cm,height=14cm,colback=white,colframe=clit,boxrule=2pt,arc=0mm,
left=6mm,right=6mm,top=0mm,bottom=4mm,sharp corners,
overlay={\fill[clit] (frame.north west) rectangle ([yshift=-1.25cm]frame.north east);
\node[anchor=west,text=white,font=\sffamily\bfseries\LARGE] at ([xshift=5mm,yshift=-0.62cm]frame.north west) {Question 25};
\node[anchor=east,text=white,font=\sffamily\itshape\large] at ([xshift=-5mm,yshift=-0.62cm]frame.north east) {Calcul littéral et équations};}]
\vspace*{1.75cm}
\begin{minipage}[t][6.9cm][c]{\linewidth}\LARGE\centering
\begin{minipage}{0.95\linewidth}\raggedright\linespread{1.25}\selectfont Quelles sont les solutions de l'inéquation $3x-5>7$~?\end{minipage}
\end{minipage}

\vfill
\begin{tcolorbox}[enhanced,nobeforeafter,width=7.1cm,height=2.1cm,valign=center,colback=clit!6,colframe=clit,arc=2.5mm,boxrule=0.9pt,left=3mm]{\Large\bfseries\color{clit}a)}\quad\begin{minipage}[c]{5.3cm}\raggedright\Large $x>\dfrac{2}{3}$\end{minipage}\end{tcolorbox}\hfill\begin{tcolorbox}[enhanced,nobeforeafter,width=7.1cm,height=2.1cm,valign=center,colback=clit!6,colframe=clit,arc=2.5mm,boxrule=0.9pt,left=3mm]{\Large\bfseries\color{clit}b)}\quad\begin{minipage}[c]{5.3cm}\raggedright\Large $x<4$\end{minipage}\end{tcolorbox}\par\vspace{2.5mm}
\begin{tcolorbox}[enhanced,nobeforeafter,width=7.1cm,height=2.1cm,valign=center,colback=clit!6,colframe=clit,arc=2.5mm,boxrule=0.9pt,left=3mm]{\Large\bfseries\color{clit}c)}\quad\begin{minipage}[c]{5.3cm}\raggedright\Large $x>-4$\end{minipage}\end{tcolorbox}\hfill\begin{tcolorbox}[enhanced,nobeforeafter,width=7.1cm,height=2.1cm,valign=center,colback=clit!6,colframe=clit,arc=2.5mm,boxrule=0.9pt,left=3mm]{\Large\bfseries\color{clit}d)}\quad\begin{minipage}[c]{5.3cm}\raggedright\Large $x>4$\end{minipage}\end{tcolorbox}
\end{tcolorbox}
\newpage
\begin{tcolorbox}[enhanced,width=16cm,height=14cm,colback=white,colframe=clit,boxrule=2pt,arc=0mm,
left=6mm,right=6mm,top=0mm,bottom=4mm,sharp corners,
overlay={\fill[clit] (frame.north west) rectangle ([yshift=-1.25cm]frame.north east);
\node[anchor=west,text=white,font=\sffamily\bfseries\LARGE] at ([xshift=5mm,yshift=-0.62cm]frame.north west) {Question 26};
\node[anchor=east,text=white,font=\sffamily\itshape\large] at ([xshift=-5mm,yshift=-0.62cm]frame.north east) {Calcul littéral et équations};}]
\vspace*{1.75cm}
\begin{minipage}[t][6.9cm][c]{\linewidth}\LARGE\centering
\begin{minipage}{0.95\linewidth}\raggedright\linespread{1.25}\selectfont Quelles sont les solutions de l'inéquation $-2x+4<10$~?\end{minipage}
\end{minipage}

\vfill
\begin{tcolorbox}[enhanced,nobeforeafter,width=7.1cm,height=2.1cm,valign=center,colback=clit!6,colframe=clit,arc=2.5mm,boxrule=0.9pt,left=3mm]{\Large\bfseries\color{clit}a)}\quad\begin{minipage}[c]{5.3cm}\raggedright\Large $x<-3$\end{minipage}\end{tcolorbox}\hfill\begin{tcolorbox}[enhanced,nobeforeafter,width=7.1cm,height=2.1cm,valign=center,colback=clit!6,colframe=clit,arc=2.5mm,boxrule=0.9pt,left=3mm]{\Large\bfseries\color{clit}b)}\quad\begin{minipage}[c]{5.3cm}\raggedright\Large $x>-3$\end{minipage}\end{tcolorbox}\par\vspace{2.5mm}
\begin{tcolorbox}[enhanced,nobeforeafter,width=7.1cm,height=2.1cm,valign=center,colback=clit!6,colframe=clit,arc=2.5mm,boxrule=0.9pt,left=3mm]{\Large\bfseries\color{clit}c)}\quad\begin{minipage}[c]{5.3cm}\raggedright\Large $x>-7$\end{minipage}\end{tcolorbox}\hfill\begin{tcolorbox}[enhanced,nobeforeafter,width=7.1cm,height=2.1cm,valign=center,colback=clit!6,colframe=clit,arc=2.5mm,boxrule=0.9pt,left=3mm]{\Large\bfseries\color{clit}d)}\quad\begin{minipage}[c]{5.3cm}\raggedright\Large $x>3$\end{minipage}\end{tcolorbox}
\end{tcolorbox}
\newpage
\begin{tcolorbox}[enhanced,width=16cm,height=14cm,colback=white,colframe=clit,boxrule=2pt,arc=0mm,
left=6mm,right=6mm,top=0mm,bottom=4mm,sharp corners,
overlay={\fill[clit] (frame.north west) rectangle ([yshift=-1.25cm]frame.north east);
\node[anchor=west,text=white,font=\sffamily\bfseries\LARGE] at ([xshift=5mm,yshift=-0.62cm]frame.north west) {Question 27};
\node[anchor=east,text=white,font=\sffamily\itshape\large] at ([xshift=-5mm,yshift=-0.62cm]frame.north east) {Calcul littéral et équations};}]
\vspace*{1.75cm}
\begin{minipage}[t][6.9cm][c]{\linewidth}\LARGE\centering
\begin{minipage}{0.95\linewidth}\raggedright\linespread{1.25}\selectfont Quelles sont les solutions de l'équation $x^2=16$~?\end{minipage}
\end{minipage}

\vfill
\begin{tcolorbox}[enhanced,nobeforeafter,width=7.1cm,height=2.1cm,valign=center,colback=clit!6,colframe=clit,arc=2.5mm,boxrule=0.9pt,left=3mm]{\Large\bfseries\color{clit}a)}\quad\begin{minipage}[c]{5.3cm}\raggedright\Large $\{-8\,;\,8\}$\end{minipage}\end{tcolorbox}\hfill\begin{tcolorbox}[enhanced,nobeforeafter,width=7.1cm,height=2.1cm,valign=center,colback=clit!6,colframe=clit,arc=2.5mm,boxrule=0.9pt,left=3mm]{\Large\bfseries\color{clit}b)}\quad\begin{minipage}[c]{5.3cm}\raggedright\Large $\{-4\,;\,4\}$\end{minipage}\end{tcolorbox}\par\vspace{2.5mm}
\begin{tcolorbox}[enhanced,nobeforeafter,width=7.1cm,height=2.1cm,valign=center,colback=clit!6,colframe=clit,arc=2.5mm,boxrule=0.9pt,left=3mm]{\Large\bfseries\color{clit}c)}\quad\begin{minipage}[c]{5.3cm}\raggedright\Large aucune solution\end{minipage}\end{tcolorbox}\hfill\begin{tcolorbox}[enhanced,nobeforeafter,width=7.1cm,height=2.1cm,valign=center,colback=clit!6,colframe=clit,arc=2.5mm,boxrule=0.9pt,left=3mm]{\Large\bfseries\color{clit}d)}\quad\begin{minipage}[c]{5.3cm}\raggedright\Large $\{4\}$\end{minipage}\end{tcolorbox}
\end{tcolorbox}
\newpage
\begin{tcolorbox}[enhanced,width=16cm,height=14cm,colback=white,colframe=clit,boxrule=2pt,arc=0mm,
left=6mm,right=6mm,top=0mm,bottom=4mm,sharp corners,
overlay={\fill[clit] (frame.north west) rectangle ([yshift=-1.25cm]frame.north east);
\node[anchor=west,text=white,font=\sffamily\bfseries\LARGE] at ([xshift=5mm,yshift=-0.62cm]frame.north west) {Question 28};
\node[anchor=east,text=white,font=\sffamily\itshape\large] at ([xshift=-5mm,yshift=-0.62cm]frame.north east) {Calcul littéral et équations};}]
\vspace*{1.75cm}
\begin{minipage}[t][6.9cm][c]{\linewidth}\LARGE\centering
\begin{minipage}{0.95\linewidth}\raggedright\linespread{1.25}\selectfont Quelles sont les solutions de l'équation $x^2=36$~?\end{minipage}
\end{minipage}

\vfill
\begin{tcolorbox}[enhanced,nobeforeafter,width=7.1cm,height=2.1cm,valign=center,colback=clit!6,colframe=clit,arc=2.5mm,boxrule=0.9pt,left=3mm]{\Large\bfseries\color{clit}a)}\quad\begin{minipage}[c]{5.3cm}\raggedright\Large $\{-6\,;\,6\}$\end{minipage}\end{tcolorbox}\hfill\begin{tcolorbox}[enhanced,nobeforeafter,width=7.1cm,height=2.1cm,valign=center,colback=clit!6,colframe=clit,arc=2.5mm,boxrule=0.9pt,left=3mm]{\Large\bfseries\color{clit}b)}\quad\begin{minipage}[c]{5.3cm}\raggedright\Large $\{6\}$\end{minipage}\end{tcolorbox}\par\vspace{2.5mm}
\begin{tcolorbox}[enhanced,nobeforeafter,width=7.1cm,height=2.1cm,valign=center,colback=clit!6,colframe=clit,arc=2.5mm,boxrule=0.9pt,left=3mm]{\Large\bfseries\color{clit}c)}\quad\begin{minipage}[c]{5.3cm}\raggedright\Large $\{-18\,;\,18\}$\end{minipage}\end{tcolorbox}\hfill\begin{tcolorbox}[enhanced,nobeforeafter,width=7.1cm,height=2.1cm,valign=center,colback=clit!6,colframe=clit,arc=2.5mm,boxrule=0.9pt,left=3mm]{\Large\bfseries\color{clit}d)}\quad\begin{minipage}[c]{5.3cm}\raggedright\Large aucune solution\end{minipage}\end{tcolorbox}
\end{tcolorbox}
\newpage
\begin{tcolorbox}[enhanced,width=16cm,height=14cm,colback=white,colframe=clit,boxrule=2pt,arc=0mm,
left=6mm,right=6mm,top=0mm,bottom=4mm,sharp corners,
overlay={\fill[clit] (frame.north west) rectangle ([yshift=-1.25cm]frame.north east);
\node[anchor=west,text=white,font=\sffamily\bfseries\LARGE] at ([xshift=5mm,yshift=-0.62cm]frame.north west) {Question 29};
\node[anchor=east,text=white,font=\sffamily\itshape\large] at ([xshift=-5mm,yshift=-0.62cm]frame.north east) {Calcul littéral et équations};}]
\vspace*{1.75cm}
\begin{minipage}[t][6.9cm][c]{\linewidth}\LARGE\centering
\begin{minipage}{0.95\linewidth}\raggedright\linespread{1.25}\selectfont Le nombre $-2$ est-il solution de l'équation $x^2+3x+2=0$~?\end{minipage}
\end{minipage}

\vfill
\begin{tcolorbox}[enhanced,nobeforeafter,width=7.1cm,height=2.1cm,valign=center,colback=clit!6,colframe=clit,arc=2.5mm,boxrule=0.9pt,left=3mm]{\Large\bfseries\color{clit}a)}\quad\begin{minipage}[c]{5.3cm}\raggedright\Large non\end{minipage}\end{tcolorbox}\hfill\begin{tcolorbox}[enhanced,nobeforeafter,width=7.1cm,height=2.1cm,valign=center,colback=clit!6,colframe=clit,arc=2.5mm,boxrule=0.9pt,left=3mm]{\Large\bfseries\color{clit}b)}\quad\begin{minipage}[c]{5.3cm}\raggedright\Large oui\end{minipage}\end{tcolorbox}\par\vspace{2.5mm}
\begin{tcolorbox}[enhanced,nobeforeafter,width=7.1cm,height=2.1cm,valign=center,colback=clit!6,colframe=clit,arc=2.5mm,boxrule=0.9pt,left=3mm]{\Large\bfseries\color{clit}c)}\quad\begin{minipage}[c]{5.3cm}\raggedright\large seulement si $x$ est positif\end{minipage}\end{tcolorbox}\hfill\begin{tcolorbox}[enhanced,nobeforeafter,width=7.1cm,height=2.1cm,valign=center,colback=clit!6,colframe=clit,arc=2.5mm,boxrule=0.9pt,left=3mm]{\Large\bfseries\color{clit}d)}\quad\begin{minipage}[c]{5.3cm}\raggedright\Large on ne peut pas savoir\end{minipage}\end{tcolorbox}
\end{tcolorbox}
\newpage
\begin{tcolorbox}[enhanced,width=16cm,height=14cm,colback=white,colframe=clit,boxrule=2pt,arc=0mm,
left=6mm,right=6mm,top=0mm,bottom=4mm,sharp corners,
overlay={\fill[clit] (frame.north west) rectangle ([yshift=-1.25cm]frame.north east);
\node[anchor=west,text=white,font=\sffamily\bfseries\LARGE] at ([xshift=5mm,yshift=-0.62cm]frame.north west) {Question 30};
\node[anchor=east,text=white,font=\sffamily\itshape\large] at ([xshift=-5mm,yshift=-0.62cm]frame.north east) {Calcul littéral et équations};}]
\vspace*{1.75cm}
\begin{minipage}[t][6.9cm][c]{\linewidth}\LARGE\centering
\begin{minipage}{0.95\linewidth}\raggedright\linespread{1.25}\selectfont Le nombre $3$ est-il solution de l'équation $x^2-5x+6=0$~?\end{minipage}
\end{minipage}

\vfill
\begin{tcolorbox}[enhanced,nobeforeafter,width=7.1cm,height=2.1cm,valign=center,colback=clit!6,colframe=clit,arc=2.5mm,boxrule=0.9pt,left=3mm]{\Large\bfseries\color{clit}a)}\quad\begin{minipage}[c]{5.3cm}\raggedright\Large non\end{minipage}\end{tcolorbox}\hfill\begin{tcolorbox}[enhanced,nobeforeafter,width=7.1cm,height=2.1cm,valign=center,colback=clit!6,colframe=clit,arc=2.5mm,boxrule=0.9pt,left=3mm]{\Large\bfseries\color{clit}b)}\quad\begin{minipage}[c]{5.3cm}\raggedright\large seulement si $x$ est positif\end{minipage}\end{tcolorbox}\par\vspace{2.5mm}
\begin{tcolorbox}[enhanced,nobeforeafter,width=7.1cm,height=2.1cm,valign=center,colback=clit!6,colframe=clit,arc=2.5mm,boxrule=0.9pt,left=3mm]{\Large\bfseries\color{clit}c)}\quad\begin{minipage}[c]{5.3cm}\raggedright\Large oui\end{minipage}\end{tcolorbox}\hfill\begin{tcolorbox}[enhanced,nobeforeafter,width=7.1cm,height=2.1cm,valign=center,colback=clit!6,colframe=clit,arc=2.5mm,boxrule=0.9pt,left=3mm]{\Large\bfseries\color{clit}d)}\quad\begin{minipage}[c]{5.3cm}\raggedright\Large on ne peut pas savoir\end{minipage}\end{tcolorbox}
\end{tcolorbox}
\newpage
\begin{tcolorbox}[enhanced,width=16cm,height=14cm,colback=white,colframe=clit,boxrule=2pt,arc=0mm,
left=6mm,right=6mm,top=0mm,bottom=4mm,sharp corners,
overlay={\fill[clit] (frame.north west) rectangle ([yshift=-1.25cm]frame.north east);
\node[anchor=west,text=white,font=\sffamily\bfseries\LARGE] at ([xshift=5mm,yshift=-0.62cm]frame.north west) {Question 31};
\node[anchor=east,text=white,font=\sffamily\itshape\large] at ([xshift=-5mm,yshift=-0.62cm]frame.north east) {Calcul littéral et équations};}]
\vspace*{1.75cm}
\begin{minipage}[t][6.9cm][c]{\linewidth}\Large\centering
\begin{minipage}{0.95\linewidth}\raggedright\linespread{1.25}\selectfont Programme de calcul~: choisir un nombre ; ajouter $4$ ; multiplier le résultat par le nombre choisi ; soustraire $2$. Quel résultat obtient-on en choisissant $3$~?\end{minipage}
\end{minipage}

\vfill
\begin{tcolorbox}[enhanced,nobeforeafter,width=7.1cm,height=2.1cm,valign=center,colback=clit!6,colframe=clit,arc=2.5mm,boxrule=0.9pt,left=3mm]{\Large\bfseries\color{clit}a)}\quad\begin{minipage}[c]{5.3cm}\raggedright\Large $19$\end{minipage}\end{tcolorbox}\hfill\begin{tcolorbox}[enhanced,nobeforeafter,width=7.1cm,height=2.1cm,valign=center,colback=clit!6,colframe=clit,arc=2.5mm,boxrule=0.9pt,left=3mm]{\Large\bfseries\color{clit}b)}\quad\begin{minipage}[c]{5.3cm}\raggedright\Large $5$\end{minipage}\end{tcolorbox}\par\vspace{2.5mm}
\begin{tcolorbox}[enhanced,nobeforeafter,width=7.1cm,height=2.1cm,valign=center,colback=clit!6,colframe=clit,arc=2.5mm,boxrule=0.9pt,left=3mm]{\Large\bfseries\color{clit}c)}\quad\begin{minipage}[c]{5.3cm}\raggedright\Large $13$\end{minipage}\end{tcolorbox}\hfill\begin{tcolorbox}[enhanced,nobeforeafter,width=7.1cm,height=2.1cm,valign=center,colback=clit!6,colframe=clit,arc=2.5mm,boxrule=0.9pt,left=3mm]{\Large\bfseries\color{clit}d)}\quad\begin{minipage}[c]{5.3cm}\raggedright\Large $23$\end{minipage}\end{tcolorbox}
\end{tcolorbox}
\newpage
\begin{tcolorbox}[enhanced,width=16cm,height=14cm,colback=white,colframe=clit,boxrule=2pt,arc=0mm,
left=6mm,right=6mm,top=0mm,bottom=4mm,sharp corners,
overlay={\fill[clit] (frame.north west) rectangle ([yshift=-1.25cm]frame.north east);
\node[anchor=west,text=white,font=\sffamily\bfseries\LARGE] at ([xshift=5mm,yshift=-0.62cm]frame.north west) {Question 32};
\node[anchor=east,text=white,font=\sffamily\itshape\large] at ([xshift=-5mm,yshift=-0.62cm]frame.north east) {Calcul littéral et équations};}]
\vspace*{1.75cm}
\begin{minipage}[t][6.9cm][c]{\linewidth}\Large\centering
\begin{minipage}{0.95\linewidth}\raggedright\linespread{1.25}\selectfont Programme de calcul~: choisir un nombre ; ajouter $5$ ; multiplier le résultat par le nombre choisi ; soustraire $3$. Quel résultat obtient-on en choisissant $-2$~?\end{minipage}
\end{minipage}

\vfill
\begin{tcolorbox}[enhanced,nobeforeafter,width=7.1cm,height=2.1cm,valign=center,colback=clit!6,colframe=clit,arc=2.5mm,boxrule=0.9pt,left=3mm]{\Large\bfseries\color{clit}a)}\quad\begin{minipage}[c]{5.3cm}\raggedright\Large $-3$\end{minipage}\end{tcolorbox}\hfill\begin{tcolorbox}[enhanced,nobeforeafter,width=7.1cm,height=2.1cm,valign=center,colback=clit!6,colframe=clit,arc=2.5mm,boxrule=0.9pt,left=3mm]{\Large\bfseries\color{clit}b)}\quad\begin{minipage}[c]{5.3cm}\raggedright\Large $-15$\end{minipage}\end{tcolorbox}\par\vspace{2.5mm}
\begin{tcolorbox}[enhanced,nobeforeafter,width=7.1cm,height=2.1cm,valign=center,colback=clit!6,colframe=clit,arc=2.5mm,boxrule=0.9pt,left=3mm]{\Large\bfseries\color{clit}c)}\quad\begin{minipage}[c]{5.3cm}\raggedright\Large $0$\end{minipage}\end{tcolorbox}\hfill\begin{tcolorbox}[enhanced,nobeforeafter,width=7.1cm,height=2.1cm,valign=center,colback=clit!6,colframe=clit,arc=2.5mm,boxrule=0.9pt,left=3mm]{\Large\bfseries\color{clit}d)}\quad\begin{minipage}[c]{5.3cm}\raggedright\Large $-9$\end{minipage}\end{tcolorbox}
\end{tcolorbox}
\newpage
\begin{tcolorbox}[enhanced,width=16cm,height=14cm,colback=white,colframe=cfon,boxrule=2pt,arc=0mm,
left=6mm,right=6mm,top=0mm,bottom=4mm,sharp corners,
overlay={\fill[cfon] (frame.north west) rectangle ([yshift=-1.25cm]frame.north east);
\node[anchor=west,text=white,font=\sffamily\bfseries\LARGE] at ([xshift=5mm,yshift=-0.62cm]frame.north west) {Question 33};
\node[anchor=east,text=white,font=\sffamily\itshape\large] at ([xshift=-5mm,yshift=-0.62cm]frame.north east) {Fonctions};}]
\vspace*{1.75cm}
\begin{minipage}[t][6.9cm][c]{\linewidth}\LARGE\centering
\begin{minipage}{0.95\linewidth}\raggedright\linespread{1.25}\selectfont Soit $f(x)=3x-5$. Quelle est l'image de $4$ par $f$~?\end{minipage}
\end{minipage}

\vfill
\begin{tcolorbox}[enhanced,nobeforeafter,width=7.1cm,height=2.1cm,valign=center,colback=cfon!6,colframe=cfon,arc=2.5mm,boxrule=0.9pt,left=3mm]{\Large\bfseries\color{cfon}a)}\quad\begin{minipage}[c]{5.3cm}\raggedright\Large $7$\end{minipage}\end{tcolorbox}\hfill\begin{tcolorbox}[enhanced,nobeforeafter,width=7.1cm,height=2.1cm,valign=center,colback=cfon!6,colframe=cfon,arc=2.5mm,boxrule=0.9pt,left=3mm]{\Large\bfseries\color{cfon}b)}\quad\begin{minipage}[c]{5.3cm}\raggedright\Large $17$\end{minipage}\end{tcolorbox}\par\vspace{2.5mm}
\begin{tcolorbox}[enhanced,nobeforeafter,width=7.1cm,height=2.1cm,valign=center,colback=cfon!6,colframe=cfon,arc=2.5mm,boxrule=0.9pt,left=3mm]{\Large\bfseries\color{cfon}c)}\quad\begin{minipage}[c]{5.3cm}\raggedright\Large $3$\end{minipage}\end{tcolorbox}\hfill\begin{tcolorbox}[enhanced,nobeforeafter,width=7.1cm,height=2.1cm,valign=center,colback=cfon!6,colframe=cfon,arc=2.5mm,boxrule=0.9pt,left=3mm]{\Large\bfseries\color{cfon}d)}\quad\begin{minipage}[c]{5.3cm}\raggedright\Large $2$\end{minipage}\end{tcolorbox}
\end{tcolorbox}
\newpage
\begin{tcolorbox}[enhanced,width=16cm,height=14cm,colback=white,colframe=cfon,boxrule=2pt,arc=0mm,
left=6mm,right=6mm,top=0mm,bottom=4mm,sharp corners,
overlay={\fill[cfon] (frame.north west) rectangle ([yshift=-1.25cm]frame.north east);
\node[anchor=west,text=white,font=\sffamily\bfseries\LARGE] at ([xshift=5mm,yshift=-0.62cm]frame.north west) {Question 34};
\node[anchor=east,text=white,font=\sffamily\itshape\large] at ([xshift=-5mm,yshift=-0.62cm]frame.north east) {Fonctions};}]
\vspace*{1.75cm}
\begin{minipage}[t][6.9cm][c]{\linewidth}\LARGE\centering
\begin{minipage}{0.95\linewidth}\raggedright\linespread{1.25}\selectfont Soit $f(x)=-2x+7$. Quelle est l'image de $3$ par $f$~?\end{minipage}
\end{minipage}

\vfill
\begin{tcolorbox}[enhanced,nobeforeafter,width=7.1cm,height=2.1cm,valign=center,colback=cfon!6,colframe=cfon,arc=2.5mm,boxrule=0.9pt,left=3mm]{\Large\bfseries\color{cfon}a)}\quad\begin{minipage}[c]{5.3cm}\raggedright\Large $8$\end{minipage}\end{tcolorbox}\hfill\begin{tcolorbox}[enhanced,nobeforeafter,width=7.1cm,height=2.1cm,valign=center,colback=cfon!6,colframe=cfon,arc=2.5mm,boxrule=0.9pt,left=3mm]{\Large\bfseries\color{cfon}b)}\quad\begin{minipage}[c]{5.3cm}\raggedright\Large $1$\end{minipage}\end{tcolorbox}\par\vspace{2.5mm}
\begin{tcolorbox}[enhanced,nobeforeafter,width=7.1cm,height=2.1cm,valign=center,colback=cfon!6,colframe=cfon,arc=2.5mm,boxrule=0.9pt,left=3mm]{\Large\bfseries\color{cfon}c)}\quad\begin{minipage}[c]{5.3cm}\raggedright\Large $-13$\end{minipage}\end{tcolorbox}\hfill\begin{tcolorbox}[enhanced,nobeforeafter,width=7.1cm,height=2.1cm,valign=center,colback=cfon!6,colframe=cfon,arc=2.5mm,boxrule=0.9pt,left=3mm]{\Large\bfseries\color{cfon}d)}\quad\begin{minipage}[c]{5.3cm}\raggedright\Large $2$\end{minipage}\end{tcolorbox}
\end{tcolorbox}
\newpage
\begin{tcolorbox}[enhanced,width=16cm,height=14cm,colback=white,colframe=cfon,boxrule=2pt,arc=0mm,
left=6mm,right=6mm,top=0mm,bottom=4mm,sharp corners,
overlay={\fill[cfon] (frame.north west) rectangle ([yshift=-1.25cm]frame.north east);
\node[anchor=west,text=white,font=\sffamily\bfseries\LARGE] at ([xshift=5mm,yshift=-0.62cm]frame.north west) {Question 35};
\node[anchor=east,text=white,font=\sffamily\itshape\large] at ([xshift=-5mm,yshift=-0.62cm]frame.north east) {Fonctions};}]
\vspace*{1.75cm}
\begin{minipage}[t][6.9cm][c]{\linewidth}\LARGE\centering
\begin{minipage}{0.95\linewidth}\raggedright\linespread{1.25}\selectfont Soit $f(x)=2x+5$. Quel est l'antécédent de $13$ par $f$~?\end{minipage}
\end{minipage}

\vfill
\begin{tcolorbox}[enhanced,nobeforeafter,width=7.1cm,height=2.1cm,valign=center,colback=cfon!6,colframe=cfon,arc=2.5mm,boxrule=0.9pt,left=3mm]{\Large\bfseries\color{cfon}a)}\quad\begin{minipage}[c]{5.3cm}\raggedright\Large $9$\end{minipage}\end{tcolorbox}\hfill\begin{tcolorbox}[enhanced,nobeforeafter,width=7.1cm,height=2.1cm,valign=center,colback=cfon!6,colframe=cfon,arc=2.5mm,boxrule=0.9pt,left=3mm]{\Large\bfseries\color{cfon}b)}\quad\begin{minipage}[c]{5.3cm}\raggedright\Large $\dfrac{3}{2}$\end{minipage}\end{tcolorbox}\par\vspace{2.5mm}
\begin{tcolorbox}[enhanced,nobeforeafter,width=7.1cm,height=2.1cm,valign=center,colback=cfon!6,colframe=cfon,arc=2.5mm,boxrule=0.9pt,left=3mm]{\Large\bfseries\color{cfon}c)}\quad\begin{minipage}[c]{5.3cm}\raggedright\Large $31$\end{minipage}\end{tcolorbox}\hfill\begin{tcolorbox}[enhanced,nobeforeafter,width=7.1cm,height=2.1cm,valign=center,colback=cfon!6,colframe=cfon,arc=2.5mm,boxrule=0.9pt,left=3mm]{\Large\bfseries\color{cfon}d)}\quad\begin{minipage}[c]{5.3cm}\raggedright\Large $4$\end{minipage}\end{tcolorbox}
\end{tcolorbox}
\newpage
\begin{tcolorbox}[enhanced,width=16cm,height=14cm,colback=white,colframe=cfon,boxrule=2pt,arc=0mm,
left=6mm,right=6mm,top=0mm,bottom=4mm,sharp corners,
overlay={\fill[cfon] (frame.north west) rectangle ([yshift=-1.25cm]frame.north east);
\node[anchor=west,text=white,font=\sffamily\bfseries\LARGE] at ([xshift=5mm,yshift=-0.62cm]frame.north west) {Question 36};
\node[anchor=east,text=white,font=\sffamily\itshape\large] at ([xshift=-5mm,yshift=-0.62cm]frame.north east) {Fonctions};}]
\vspace*{1.75cm}
\begin{minipage}[t][6.9cm][c]{\linewidth}\LARGE\centering
\begin{minipage}{0.95\linewidth}\raggedright\linespread{1.25}\selectfont Soit $f(x)=-3x+4$. Quel est l'antécédent de $-8$ par $f$~?\end{minipage}
\end{minipage}

\vfill
\begin{tcolorbox}[enhanced,nobeforeafter,width=7.1cm,height=2.1cm,valign=center,colback=cfon!6,colframe=cfon,arc=2.5mm,boxrule=0.9pt,left=3mm]{\Large\bfseries\color{cfon}a)}\quad\begin{minipage}[c]{5.3cm}\raggedright\Large $-\dfrac{4}{3}$\end{minipage}\end{tcolorbox}\hfill\begin{tcolorbox}[enhanced,nobeforeafter,width=7.1cm,height=2.1cm,valign=center,colback=cfon!6,colframe=cfon,arc=2.5mm,boxrule=0.9pt,left=3mm]{\Large\bfseries\color{cfon}b)}\quad\begin{minipage}[c]{5.3cm}\raggedright\Large $28$\end{minipage}\end{tcolorbox}\par\vspace{2.5mm}
\begin{tcolorbox}[enhanced,nobeforeafter,width=7.1cm,height=2.1cm,valign=center,colback=cfon!6,colframe=cfon,arc=2.5mm,boxrule=0.9pt,left=3mm]{\Large\bfseries\color{cfon}c)}\quad\begin{minipage}[c]{5.3cm}\raggedright\Large $\dfrac{4}{3}$\end{minipage}\end{tcolorbox}\hfill\begin{tcolorbox}[enhanced,nobeforeafter,width=7.1cm,height=2.1cm,valign=center,colback=cfon!6,colframe=cfon,arc=2.5mm,boxrule=0.9pt,left=3mm]{\Large\bfseries\color{cfon}d)}\quad\begin{minipage}[c]{5.3cm}\raggedright\Large $4$\end{minipage}\end{tcolorbox}
\end{tcolorbox}
\newpage
\begin{tcolorbox}[enhanced,width=16cm,height=14cm,colback=white,colframe=cfon,boxrule=2pt,arc=0mm,
left=6mm,right=6mm,top=0mm,bottom=4mm,sharp corners,
overlay={\fill[cfon] (frame.north west) rectangle ([yshift=-1.25cm]frame.north east);
\node[anchor=west,text=white,font=\sffamily\bfseries\LARGE] at ([xshift=5mm,yshift=-0.62cm]frame.north west) {Question 37};
\node[anchor=east,text=white,font=\sffamily\itshape\large] at ([xshift=-5mm,yshift=-0.62cm]frame.north east) {Fonctions};}]
\vspace*{1.75cm}
\begin{minipage}[t][6.9cm][c]{\linewidth}\LARGE\centering
\begin{minipage}{0.95\linewidth}\raggedright\linespread{1.25}\selectfont $f$ est une fonction linéaire telle que $f(4)=12$. Quel est son coefficient~?\end{minipage}
\end{minipage}

\vfill
\begin{tcolorbox}[enhanced,nobeforeafter,width=7.1cm,height=2.1cm,valign=center,colback=cfon!6,colframe=cfon,arc=2.5mm,boxrule=0.9pt,left=3mm]{\Large\bfseries\color{cfon}a)}\quad\begin{minipage}[c]{5.3cm}\raggedright\Large $48$\end{minipage}\end{tcolorbox}\hfill\begin{tcolorbox}[enhanced,nobeforeafter,width=7.1cm,height=2.1cm,valign=center,colback=cfon!6,colframe=cfon,arc=2.5mm,boxrule=0.9pt,left=3mm]{\Large\bfseries\color{cfon}b)}\quad\begin{minipage}[c]{5.3cm}\raggedright\Large $8$\end{minipage}\end{tcolorbox}\par\vspace{2.5mm}
\begin{tcolorbox}[enhanced,nobeforeafter,width=7.1cm,height=2.1cm,valign=center,colback=cfon!6,colframe=cfon,arc=2.5mm,boxrule=0.9pt,left=3mm]{\Large\bfseries\color{cfon}c)}\quad\begin{minipage}[c]{5.3cm}\raggedright\Large $\dfrac{1}{3}$\end{minipage}\end{tcolorbox}\hfill\begin{tcolorbox}[enhanced,nobeforeafter,width=7.1cm,height=2.1cm,valign=center,colback=cfon!6,colframe=cfon,arc=2.5mm,boxrule=0.9pt,left=3mm]{\Large\bfseries\color{cfon}d)}\quad\begin{minipage}[c]{5.3cm}\raggedright\Large $3$\end{minipage}\end{tcolorbox}
\end{tcolorbox}
\newpage
\begin{tcolorbox}[enhanced,width=16cm,height=14cm,colback=white,colframe=cfon,boxrule=2pt,arc=0mm,
left=6mm,right=6mm,top=0mm,bottom=4mm,sharp corners,
overlay={\fill[cfon] (frame.north west) rectangle ([yshift=-1.25cm]frame.north east);
\node[anchor=west,text=white,font=\sffamily\bfseries\LARGE] at ([xshift=5mm,yshift=-0.62cm]frame.north west) {Question 38};
\node[anchor=east,text=white,font=\sffamily\itshape\large] at ([xshift=-5mm,yshift=-0.62cm]frame.north east) {Fonctions};}]
\vspace*{1.75cm}
\begin{minipage}[t][6.9cm][c]{\linewidth}\LARGE\centering
\begin{minipage}{0.95\linewidth}\raggedright\linespread{1.25}\selectfont $f$ est une fonction linéaire telle que $f(5)=-15$. Quel est son coefficient~?\end{minipage}
\end{minipage}

\vfill
\begin{tcolorbox}[enhanced,nobeforeafter,width=7.1cm,height=2.1cm,valign=center,colback=cfon!6,colframe=cfon,arc=2.5mm,boxrule=0.9pt,left=3mm]{\Large\bfseries\color{cfon}a)}\quad\begin{minipage}[c]{5.3cm}\raggedright\Large $-\dfrac{1}{3}$\end{minipage}\end{tcolorbox}\hfill\begin{tcolorbox}[enhanced,nobeforeafter,width=7.1cm,height=2.1cm,valign=center,colback=cfon!6,colframe=cfon,arc=2.5mm,boxrule=0.9pt,left=3mm]{\Large\bfseries\color{cfon}b)}\quad\begin{minipage}[c]{5.3cm}\raggedright\Large $-3$\end{minipage}\end{tcolorbox}\par\vspace{2.5mm}
\begin{tcolorbox}[enhanced,nobeforeafter,width=7.1cm,height=2.1cm,valign=center,colback=cfon!6,colframe=cfon,arc=2.5mm,boxrule=0.9pt,left=3mm]{\Large\bfseries\color{cfon}c)}\quad\begin{minipage}[c]{5.3cm}\raggedright\Large $-20$\end{minipage}\end{tcolorbox}\hfill\begin{tcolorbox}[enhanced,nobeforeafter,width=7.1cm,height=2.1cm,valign=center,colback=cfon!6,colframe=cfon,arc=2.5mm,boxrule=0.9pt,left=3mm]{\Large\bfseries\color{cfon}d)}\quad\begin{minipage}[c]{5.3cm}\raggedright\Large $-75$\end{minipage}\end{tcolorbox}
\end{tcolorbox}
\newpage
\begin{tcolorbox}[enhanced,width=16cm,height=14cm,colback=white,colframe=cfon,boxrule=2pt,arc=0mm,
left=6mm,right=6mm,top=0mm,bottom=4mm,sharp corners,
overlay={\fill[cfon] (frame.north west) rectangle ([yshift=-1.25cm]frame.north east);
\node[anchor=west,text=white,font=\sffamily\bfseries\LARGE] at ([xshift=5mm,yshift=-0.62cm]frame.north west) {Question 39};
\node[anchor=east,text=white,font=\sffamily\itshape\large] at ([xshift=-5mm,yshift=-0.62cm]frame.north east) {Fonctions};}]
\vspace*{1.75cm}
\begin{minipage}[t][6.9cm][c]{\linewidth}\LARGE\centering
\begin{minipage}{0.95\linewidth}\raggedright\linespread{1.25}\selectfont $f$ est une fonction affine telle que $f(1)=5$ et $f(3)=11$. Quel est le coefficient $a$ de $f(x)=ax+b$~?\end{minipage}
\end{minipage}

\vfill
\begin{tcolorbox}[enhanced,nobeforeafter,width=7.1cm,height=2.1cm,valign=center,colback=cfon!6,colframe=cfon,arc=2.5mm,boxrule=0.9pt,left=3mm]{\Large\bfseries\color{cfon}a)}\quad\begin{minipage}[c]{5.3cm}\raggedright\Large $6$\end{minipage}\end{tcolorbox}\hfill\begin{tcolorbox}[enhanced,nobeforeafter,width=7.1cm,height=2.1cm,valign=center,colback=cfon!6,colframe=cfon,arc=2.5mm,boxrule=0.9pt,left=3mm]{\Large\bfseries\color{cfon}b)}\quad\begin{minipage}[c]{5.3cm}\raggedright\Large $\dfrac{1}{3}$\end{minipage}\end{tcolorbox}\par\vspace{2.5mm}
\begin{tcolorbox}[enhanced,nobeforeafter,width=7.1cm,height=2.1cm,valign=center,colback=cfon!6,colframe=cfon,arc=2.5mm,boxrule=0.9pt,left=3mm]{\Large\bfseries\color{cfon}c)}\quad\begin{minipage}[c]{5.3cm}\raggedright\Large $-3$\end{minipage}\end{tcolorbox}\hfill\begin{tcolorbox}[enhanced,nobeforeafter,width=7.1cm,height=2.1cm,valign=center,colback=cfon!6,colframe=cfon,arc=2.5mm,boxrule=0.9pt,left=3mm]{\Large\bfseries\color{cfon}d)}\quad\begin{minipage}[c]{5.3cm}\raggedright\Large $3$\end{minipage}\end{tcolorbox}
\end{tcolorbox}
\newpage
\begin{tcolorbox}[enhanced,width=16cm,height=14cm,colback=white,colframe=cfon,boxrule=2pt,arc=0mm,
left=6mm,right=6mm,top=0mm,bottom=4mm,sharp corners,
overlay={\fill[cfon] (frame.north west) rectangle ([yshift=-1.25cm]frame.north east);
\node[anchor=west,text=white,font=\sffamily\bfseries\LARGE] at ([xshift=5mm,yshift=-0.62cm]frame.north west) {Question 40};
\node[anchor=east,text=white,font=\sffamily\itshape\large] at ([xshift=-5mm,yshift=-0.62cm]frame.north east) {Fonctions};}]
\vspace*{1.75cm}
\begin{minipage}[t][6.9cm][c]{\linewidth}\LARGE\centering
\begin{minipage}{0.95\linewidth}\raggedright\linespread{1.25}\selectfont $f$ est une fonction affine telle que $f(2)=1$ et $f(6)=-7$. Quel est le coefficient $a$ de $f(x)=ax+b$~?\end{minipage}
\end{minipage}

\vfill
\begin{tcolorbox}[enhanced,nobeforeafter,width=7.1cm,height=2.1cm,valign=center,colback=cfon!6,colframe=cfon,arc=2.5mm,boxrule=0.9pt,left=3mm]{\Large\bfseries\color{cfon}a)}\quad\begin{minipage}[c]{5.3cm}\raggedright\Large $-2$\end{minipage}\end{tcolorbox}\hfill\begin{tcolorbox}[enhanced,nobeforeafter,width=7.1cm,height=2.1cm,valign=center,colback=cfon!6,colframe=cfon,arc=2.5mm,boxrule=0.9pt,left=3mm]{\Large\bfseries\color{cfon}b)}\quad\begin{minipage}[c]{5.3cm}\raggedright\Large $-\dfrac{1}{2}$\end{minipage}\end{tcolorbox}\par\vspace{2.5mm}
\begin{tcolorbox}[enhanced,nobeforeafter,width=7.1cm,height=2.1cm,valign=center,colback=cfon!6,colframe=cfon,arc=2.5mm,boxrule=0.9pt,left=3mm]{\Large\bfseries\color{cfon}c)}\quad\begin{minipage}[c]{5.3cm}\raggedright\Large $2$\end{minipage}\end{tcolorbox}\hfill\begin{tcolorbox}[enhanced,nobeforeafter,width=7.1cm,height=2.1cm,valign=center,colback=cfon!6,colframe=cfon,arc=2.5mm,boxrule=0.9pt,left=3mm]{\Large\bfseries\color{cfon}d)}\quad\begin{minipage}[c]{5.3cm}\raggedright\Large $-8$\end{minipage}\end{tcolorbox}
\end{tcolorbox}
\newpage
\begin{tcolorbox}[enhanced,width=16cm,height=14cm,colback=white,colframe=cfon,boxrule=2pt,arc=0mm,
left=6mm,right=6mm,top=0mm,bottom=4mm,sharp corners,
overlay={\fill[cfon] (frame.north west) rectangle ([yshift=-1.25cm]frame.north east);
\node[anchor=west,text=white,font=\sffamily\bfseries\LARGE] at ([xshift=5mm,yshift=-0.62cm]frame.north west) {Question 41};
\node[anchor=east,text=white,font=\sffamily\itshape\large] at ([xshift=-5mm,yshift=-0.62cm]frame.north east) {Fonctions};}]
\vspace*{1.75cm}
\begin{minipage}[t][6.9cm][c]{\linewidth}\LARGE\centering
\begin{minipage}{0.95\linewidth}\raggedright\linespread{1.25}\selectfont $f$ est la fonction affine $f(x)=2x+b$ et $f(3)=11$. Que vaut $b$~?\end{minipage}
\end{minipage}

\vfill
\begin{tcolorbox}[enhanced,nobeforeafter,width=7.1cm,height=2.1cm,valign=center,colback=cfon!6,colframe=cfon,arc=2.5mm,boxrule=0.9pt,left=3mm]{\Large\bfseries\color{cfon}a)}\quad\begin{minipage}[c]{5.3cm}\raggedright\Large $\dfrac{11}{3}$\end{minipage}\end{tcolorbox}\hfill\begin{tcolorbox}[enhanced,nobeforeafter,width=7.1cm,height=2.1cm,valign=center,colback=cfon!6,colframe=cfon,arc=2.5mm,boxrule=0.9pt,left=3mm]{\Large\bfseries\color{cfon}b)}\quad\begin{minipage}[c]{5.3cm}\raggedright\Large $5$\end{minipage}\end{tcolorbox}\par\vspace{2.5mm}
\begin{tcolorbox}[enhanced,nobeforeafter,width=7.1cm,height=2.1cm,valign=center,colback=cfon!6,colframe=cfon,arc=2.5mm,boxrule=0.9pt,left=3mm]{\Large\bfseries\color{cfon}c)}\quad\begin{minipage}[c]{5.3cm}\raggedright\Large $17$\end{minipage}\end{tcolorbox}\hfill\begin{tcolorbox}[enhanced,nobeforeafter,width=7.1cm,height=2.1cm,valign=center,colback=cfon!6,colframe=cfon,arc=2.5mm,boxrule=0.9pt,left=3mm]{\Large\bfseries\color{cfon}d)}\quad\begin{minipage}[c]{5.3cm}\raggedright\Large $11$\end{minipage}\end{tcolorbox}
\end{tcolorbox}
\newpage
\begin{tcolorbox}[enhanced,width=16cm,height=14cm,colback=white,colframe=cfon,boxrule=2pt,arc=0mm,
left=6mm,right=6mm,top=0mm,bottom=4mm,sharp corners,
overlay={\fill[cfon] (frame.north west) rectangle ([yshift=-1.25cm]frame.north east);
\node[anchor=west,text=white,font=\sffamily\bfseries\LARGE] at ([xshift=5mm,yshift=-0.62cm]frame.north west) {Question 42};
\node[anchor=east,text=white,font=\sffamily\itshape\large] at ([xshift=-5mm,yshift=-0.62cm]frame.north east) {Fonctions};}]
\vspace*{1.75cm}
\begin{minipage}[t][6.9cm][c]{\linewidth}\LARGE\centering
\begin{minipage}{0.95\linewidth}\raggedright\linespread{1.25}\selectfont $f$ est la fonction affine $f(x)=-3x+b$ et $f(1)=4$. Que vaut $b$~?\end{minipage}
\end{minipage}

\vfill
\begin{tcolorbox}[enhanced,nobeforeafter,width=7.1cm,height=2.1cm,valign=center,colback=cfon!6,colframe=cfon,arc=2.5mm,boxrule=0.9pt,left=3mm]{\Large\bfseries\color{cfon}a)}\quad\begin{minipage}[c]{5.3cm}\raggedright\Large $1$\end{minipage}\end{tcolorbox}\hfill\begin{tcolorbox}[enhanced,nobeforeafter,width=7.1cm,height=2.1cm,valign=center,colback=cfon!6,colframe=cfon,arc=2.5mm,boxrule=0.9pt,left=3mm]{\Large\bfseries\color{cfon}b)}\quad\begin{minipage}[c]{5.3cm}\raggedright\Large $7$\end{minipage}\end{tcolorbox}\par\vspace{2.5mm}
\begin{tcolorbox}[enhanced,nobeforeafter,width=7.1cm,height=2.1cm,valign=center,colback=cfon!6,colframe=cfon,arc=2.5mm,boxrule=0.9pt,left=3mm]{\Large\bfseries\color{cfon}c)}\quad\begin{minipage}[c]{5.3cm}\raggedright\Large $-7$\end{minipage}\end{tcolorbox}\hfill\begin{tcolorbox}[enhanced,nobeforeafter,width=7.1cm,height=2.1cm,valign=center,colback=cfon!6,colframe=cfon,arc=2.5mm,boxrule=0.9pt,left=3mm]{\Large\bfseries\color{cfon}d)}\quad\begin{minipage}[c]{5.3cm}\raggedright\Large $4$\end{minipage}\end{tcolorbox}
\end{tcolorbox}
\newpage
\begin{tcolorbox}[enhanced,width=16cm,height=14cm,colback=white,colframe=cfon,boxrule=2pt,arc=0mm,
left=6mm,right=6mm,top=0mm,bottom=4mm,sharp corners,
overlay={\fill[cfon] (frame.north west) rectangle ([yshift=-1.25cm]frame.north east);
\node[anchor=west,text=white,font=\sffamily\bfseries\LARGE] at ([xshift=5mm,yshift=-0.62cm]frame.north west) {Question 43};
\node[anchor=east,text=white,font=\sffamily\itshape\large] at ([xshift=-5mm,yshift=-0.62cm]frame.north east) {Fonctions};}]
\vspace*{1.75cm}
\begin{minipage}[t][6.9cm][c]{\linewidth}\LARGE\centering
\begin{minipage}{0.95\linewidth}\raggedright\linespread{1.25}\selectfont Augmenter une quantité de $20\,\%$ revient à la multiplier par~:\end{minipage}
\end{minipage}

\vfill
\begin{tcolorbox}[enhanced,nobeforeafter,width=7.1cm,height=2.1cm,valign=center,colback=cfon!6,colframe=cfon,arc=2.5mm,boxrule=0.9pt,left=3mm]{\Large\bfseries\color{cfon}a)}\quad\begin{minipage}[c]{5.3cm}\raggedright\Large $0{,}2$\end{minipage}\end{tcolorbox}\hfill\begin{tcolorbox}[enhanced,nobeforeafter,width=7.1cm,height=2.1cm,valign=center,colback=cfon!6,colframe=cfon,arc=2.5mm,boxrule=0.9pt,left=3mm]{\Large\bfseries\color{cfon}b)}\quad\begin{minipage}[c]{5.3cm}\raggedright\Large $0{,}8$\end{minipage}\end{tcolorbox}\par\vspace{2.5mm}
\begin{tcolorbox}[enhanced,nobeforeafter,width=7.1cm,height=2.1cm,valign=center,colback=cfon!6,colframe=cfon,arc=2.5mm,boxrule=0.9pt,left=3mm]{\Large\bfseries\color{cfon}c)}\quad\begin{minipage}[c]{5.3cm}\raggedright\Large $1{,}2$\end{minipage}\end{tcolorbox}\hfill\begin{tcolorbox}[enhanced,nobeforeafter,width=7.1cm,height=2.1cm,valign=center,colback=cfon!6,colframe=cfon,arc=2.5mm,boxrule=0.9pt,left=3mm]{\Large\bfseries\color{cfon}d)}\quad\begin{minipage}[c]{5.3cm}\raggedright\Large $3$\end{minipage}\end{tcolorbox}
\end{tcolorbox}
\newpage
\begin{tcolorbox}[enhanced,width=16cm,height=14cm,colback=white,colframe=cfon,boxrule=2pt,arc=0mm,
left=6mm,right=6mm,top=0mm,bottom=4mm,sharp corners,
overlay={\fill[cfon] (frame.north west) rectangle ([yshift=-1.25cm]frame.north east);
\node[anchor=west,text=white,font=\sffamily\bfseries\LARGE] at ([xshift=5mm,yshift=-0.62cm]frame.north west) {Question 44};
\node[anchor=east,text=white,font=\sffamily\itshape\large] at ([xshift=-5mm,yshift=-0.62cm]frame.north east) {Fonctions};}]
\vspace*{1.75cm}
\begin{minipage}[t][6.9cm][c]{\linewidth}\LARGE\centering
\begin{minipage}{0.95\linewidth}\raggedright\linespread{1.25}\selectfont Diminuer une quantité de $15\,\%$ revient à la multiplier par~:\end{minipage}
\end{minipage}

\vfill
\begin{tcolorbox}[enhanced,nobeforeafter,width=7.1cm,height=2.1cm,valign=center,colback=cfon!6,colframe=cfon,arc=2.5mm,boxrule=0.9pt,left=3mm]{\Large\bfseries\color{cfon}a)}\quad\begin{minipage}[c]{5.3cm}\raggedright\Large $1{,}85$\end{minipage}\end{tcolorbox}\hfill\begin{tcolorbox}[enhanced,nobeforeafter,width=7.1cm,height=2.1cm,valign=center,colback=cfon!6,colframe=cfon,arc=2.5mm,boxrule=0.9pt,left=3mm]{\Large\bfseries\color{cfon}b)}\quad\begin{minipage}[c]{5.3cm}\raggedright\Large $1{,}15$\end{minipage}\end{tcolorbox}\par\vspace{2.5mm}
\begin{tcolorbox}[enhanced,nobeforeafter,width=7.1cm,height=2.1cm,valign=center,colback=cfon!6,colframe=cfon,arc=2.5mm,boxrule=0.9pt,left=3mm]{\Large\bfseries\color{cfon}c)}\quad\begin{minipage}[c]{5.3cm}\raggedright\Large $0{,}15$\end{minipage}\end{tcolorbox}\hfill\begin{tcolorbox}[enhanced,nobeforeafter,width=7.1cm,height=2.1cm,valign=center,colback=cfon!6,colframe=cfon,arc=2.5mm,boxrule=0.9pt,left=3mm]{\Large\bfseries\color{cfon}d)}\quad\begin{minipage}[c]{5.3cm}\raggedright\Large $0{,}85$\end{minipage}\end{tcolorbox}
\end{tcolorbox}
\newpage
\begin{tcolorbox}[enhanced,width=16cm,height=14cm,colback=white,colframe=cfon,boxrule=2pt,arc=0mm,
left=6mm,right=6mm,top=0mm,bottom=4mm,sharp corners,
overlay={\fill[cfon] (frame.north west) rectangle ([yshift=-1.25cm]frame.north east);
\node[anchor=west,text=white,font=\sffamily\bfseries\LARGE] at ([xshift=5mm,yshift=-0.62cm]frame.north west) {Question 45};
\node[anchor=east,text=white,font=\sffamily\itshape\large] at ([xshift=-5mm,yshift=-0.62cm]frame.north east) {Fonctions};}]
\vspace*{1.75cm}
\begin{minipage}[t][6.9cm][c]{\linewidth}\LARGE\centering
\begin{minipage}{0.95\linewidth}\raggedright\linespread{1.25}\selectfont Soit $g(x)=x^2-3x$. Que vaut $g(4)$~?\end{minipage}
\end{minipage}

\vfill
\begin{tcolorbox}[enhanced,nobeforeafter,width=7.1cm,height=2.1cm,valign=center,colback=cfon!6,colframe=cfon,arc=2.5mm,boxrule=0.9pt,left=3mm]{\Large\bfseries\color{cfon}a)}\quad\begin{minipage}[c]{5.3cm}\raggedright\Large $-28$\end{minipage}\end{tcolorbox}\hfill\begin{tcolorbox}[enhanced,nobeforeafter,width=7.1cm,height=2.1cm,valign=center,colback=cfon!6,colframe=cfon,arc=2.5mm,boxrule=0.9pt,left=3mm]{\Large\bfseries\color{cfon}b)}\quad\begin{minipage}[c]{5.3cm}\raggedright\Large $-4$\end{minipage}\end{tcolorbox}\par\vspace{2.5mm}
\begin{tcolorbox}[enhanced,nobeforeafter,width=7.1cm,height=2.1cm,valign=center,colback=cfon!6,colframe=cfon,arc=2.5mm,boxrule=0.9pt,left=3mm]{\Large\bfseries\color{cfon}c)}\quad\begin{minipage}[c]{5.3cm}\raggedright\Large $6$\end{minipage}\end{tcolorbox}\hfill\begin{tcolorbox}[enhanced,nobeforeafter,width=7.1cm,height=2.1cm,valign=center,colback=cfon!6,colframe=cfon,arc=2.5mm,boxrule=0.9pt,left=3mm]{\Large\bfseries\color{cfon}d)}\quad\begin{minipage}[c]{5.3cm}\raggedright\Large $4$\end{minipage}\end{tcolorbox}
\end{tcolorbox}
\newpage
\begin{tcolorbox}[enhanced,width=16cm,height=14cm,colback=white,colframe=cfon,boxrule=2pt,arc=0mm,
left=6mm,right=6mm,top=0mm,bottom=4mm,sharp corners,
overlay={\fill[cfon] (frame.north west) rectangle ([yshift=-1.25cm]frame.north east);
\node[anchor=west,text=white,font=\sffamily\bfseries\LARGE] at ([xshift=5mm,yshift=-0.62cm]frame.north west) {Question 46};
\node[anchor=east,text=white,font=\sffamily\itshape\large] at ([xshift=-5mm,yshift=-0.62cm]frame.north east) {Fonctions};}]
\vspace*{1.75cm}
\begin{minipage}[t][6.9cm][c]{\linewidth}\LARGE\centering
\begin{minipage}{0.95\linewidth}\raggedright\linespread{1.25}\selectfont Soit $g(x)=2x^2+x-5$. Que vaut $g(-2)$~?\end{minipage}
\end{minipage}

\vfill
\begin{tcolorbox}[enhanced,nobeforeafter,width=7.1cm,height=2.1cm,valign=center,colback=cfon!6,colframe=cfon,arc=2.5mm,boxrule=0.9pt,left=3mm]{\Large\bfseries\color{cfon}a)}\quad\begin{minipage}[c]{5.3cm}\raggedright\Large $3$\end{minipage}\end{tcolorbox}\hfill\begin{tcolorbox}[enhanced,nobeforeafter,width=7.1cm,height=2.1cm,valign=center,colback=cfon!6,colframe=cfon,arc=2.5mm,boxrule=0.9pt,left=3mm]{\Large\bfseries\color{cfon}b)}\quad\begin{minipage}[c]{5.3cm}\raggedright\Large $-15$\end{minipage}\end{tcolorbox}\par\vspace{2.5mm}
\begin{tcolorbox}[enhanced,nobeforeafter,width=7.1cm,height=2.1cm,valign=center,colback=cfon!6,colframe=cfon,arc=2.5mm,boxrule=0.9pt,left=3mm]{\Large\bfseries\color{cfon}c)}\quad\begin{minipage}[c]{5.3cm}\raggedright\Large $1$\end{minipage}\end{tcolorbox}\hfill\begin{tcolorbox}[enhanced,nobeforeafter,width=7.1cm,height=2.1cm,valign=center,colback=cfon!6,colframe=cfon,arc=2.5mm,boxrule=0.9pt,left=3mm]{\Large\bfseries\color{cfon}d)}\quad\begin{minipage}[c]{5.3cm}\raggedright\Large $-1$\end{minipage}\end{tcolorbox}
\end{tcolorbox}
\newpage
\begin{tcolorbox}[enhanced,width=16cm,height=14cm,colback=white,colframe=cfon,boxrule=2pt,arc=0mm,
left=6mm,right=6mm,top=0mm,bottom=4mm,sharp corners,
overlay={\fill[cfon] (frame.north west) rectangle ([yshift=-1.25cm]frame.north east);
\node[anchor=west,text=white,font=\sffamily\bfseries\LARGE] at ([xshift=5mm,yshift=-0.62cm]frame.north west) {Question 47};
\node[anchor=east,text=white,font=\sffamily\itshape\large] at ([xshift=-5mm,yshift=-0.62cm]frame.north east) {Fonctions};}]
\vspace*{1.75cm}
\begin{minipage}[t][6.9cm][c]{\linewidth}\LARGE\centering
\begin{minipage}{0.95\linewidth}\raggedright\linespread{1.25}\selectfont De quel type est la fonction $f:x\mapsto 4x$~?\end{minipage}
\end{minipage}

\vfill
\begin{tcolorbox}[enhanced,nobeforeafter,width=7.1cm,height=2.1cm,valign=center,colback=cfon!6,colframe=cfon,arc=2.5mm,boxrule=0.9pt,left=3mm]{\Large\bfseries\color{cfon}a)}\quad\begin{minipage}[c]{5.3cm}\raggedright\Large constante\end{minipage}\end{tcolorbox}\hfill\begin{tcolorbox}[enhanced,nobeforeafter,width=7.1cm,height=2.1cm,valign=center,colback=cfon!6,colframe=cfon,arc=2.5mm,boxrule=0.9pt,left=3mm]{\Large\bfseries\color{cfon}b)}\quad\begin{minipage}[c]{5.3cm}\raggedright\Large affine non linéaire\end{minipage}\end{tcolorbox}\par\vspace{2.5mm}
\begin{tcolorbox}[enhanced,nobeforeafter,width=7.1cm,height=2.1cm,valign=center,colback=cfon!6,colframe=cfon,arc=2.5mm,boxrule=0.9pt,left=3mm]{\Large\bfseries\color{cfon}c)}\quad\begin{minipage}[c]{5.3cm}\raggedright\Large ni affine ni linéaire\end{minipage}\end{tcolorbox}\hfill\begin{tcolorbox}[enhanced,nobeforeafter,width=7.1cm,height=2.1cm,valign=center,colback=cfon!6,colframe=cfon,arc=2.5mm,boxrule=0.9pt,left=3mm]{\Large\bfseries\color{cfon}d)}\quad\begin{minipage}[c]{5.3cm}\raggedright\Large linéaire\end{minipage}\end{tcolorbox}
\end{tcolorbox}
\newpage
\begin{tcolorbox}[enhanced,width=16cm,height=14cm,colback=white,colframe=cfon,boxrule=2pt,arc=0mm,
left=6mm,right=6mm,top=0mm,bottom=4mm,sharp corners,
overlay={\fill[cfon] (frame.north west) rectangle ([yshift=-1.25cm]frame.north east);
\node[anchor=west,text=white,font=\sffamily\bfseries\LARGE] at ([xshift=5mm,yshift=-0.62cm]frame.north west) {Question 48};
\node[anchor=east,text=white,font=\sffamily\itshape\large] at ([xshift=-5mm,yshift=-0.62cm]frame.north east) {Fonctions};}]
\vspace*{1.75cm}
\begin{minipage}[t][6.9cm][c]{\linewidth}\LARGE\centering
\begin{minipage}{0.95\linewidth}\raggedright\linespread{1.25}\selectfont De quel type est la fonction $f:x\mapsto 2x+3$~?\end{minipage}
\end{minipage}

\vfill
\begin{tcolorbox}[enhanced,nobeforeafter,width=7.1cm,height=2.1cm,valign=center,colback=cfon!6,colframe=cfon,arc=2.5mm,boxrule=0.9pt,left=3mm]{\Large\bfseries\color{cfon}a)}\quad\begin{minipage}[c]{5.3cm}\raggedright\Large affine non linéaire\end{minipage}\end{tcolorbox}\hfill\begin{tcolorbox}[enhanced,nobeforeafter,width=7.1cm,height=2.1cm,valign=center,colback=cfon!6,colframe=cfon,arc=2.5mm,boxrule=0.9pt,left=3mm]{\Large\bfseries\color{cfon}b)}\quad\begin{minipage}[c]{5.3cm}\raggedright\Large constante\end{minipage}\end{tcolorbox}\par\vspace{2.5mm}
\begin{tcolorbox}[enhanced,nobeforeafter,width=7.1cm,height=2.1cm,valign=center,colback=cfon!6,colframe=cfon,arc=2.5mm,boxrule=0.9pt,left=3mm]{\Large\bfseries\color{cfon}c)}\quad\begin{minipage}[c]{5.3cm}\raggedright\Large ni affine ni linéaire\end{minipage}\end{tcolorbox}\hfill\begin{tcolorbox}[enhanced,nobeforeafter,width=7.1cm,height=2.1cm,valign=center,colback=cfon!6,colframe=cfon,arc=2.5mm,boxrule=0.9pt,left=3mm]{\Large\bfseries\color{cfon}d)}\quad\begin{minipage}[c]{5.3cm}\raggedright\Large linéaire\end{minipage}\end{tcolorbox}
\end{tcolorbox}
\newpage
\begin{tcolorbox}[enhanced,width=16cm,height=14cm,colback=white,colframe=cgeo,boxrule=2pt,arc=0mm,
left=6mm,right=6mm,top=0mm,bottom=4mm,sharp corners,
overlay={\fill[cgeo] (frame.north west) rectangle ([yshift=-1.25cm]frame.north east);
\node[anchor=west,text=white,font=\sffamily\bfseries\LARGE] at ([xshift=5mm,yshift=-0.62cm]frame.north west) {Question 49};
\node[anchor=east,text=white,font=\sffamily\itshape\large] at ([xshift=-5mm,yshift=-0.62cm]frame.north east) {Géométrie};}]
\vspace*{1.75cm}
\begin{minipage}[t][6.9cm][c]{\linewidth}\Large\centering
\begin{minipage}{0.95\linewidth}\raggedright\linespread{1.25}\selectfont Les droites $(AC)$ et $(BD)$ se coupent en $O$, et les droites $(AB)$ et $(CD)$ sont parallèles. $OA=3$ cm, $OC=6$ cm et $AB=4$ cm. Combien mesure $CD$~?\end{minipage}
\end{minipage}

\vfill
\begin{tcolorbox}[enhanced,nobeforeafter,width=7.1cm,height=2.1cm,valign=center,colback=cgeo!6,colframe=cgeo,arc=2.5mm,boxrule=0.9pt,left=3mm]{\Large\bfseries\color{cgeo}a)}\quad\begin{minipage}[c]{5.3cm}\raggedright\Large $4{,}5~\text{cm}$\end{minipage}\end{tcolorbox}\hfill\begin{tcolorbox}[enhanced,nobeforeafter,width=7.1cm,height=2.1cm,valign=center,colback=cgeo!6,colframe=cgeo,arc=2.5mm,boxrule=0.9pt,left=3mm]{\Large\bfseries\color{cgeo}b)}\quad\begin{minipage}[c]{5.3cm}\raggedright\Large $8~\text{cm}$\end{minipage}\end{tcolorbox}\par\vspace{2.5mm}
\begin{tcolorbox}[enhanced,nobeforeafter,width=7.1cm,height=2.1cm,valign=center,colback=cgeo!6,colframe=cgeo,arc=2.5mm,boxrule=0.9pt,left=3mm]{\Large\bfseries\color{cgeo}c)}\quad\begin{minipage}[c]{5.3cm}\raggedright\Large $2~\text{cm}$\end{minipage}\end{tcolorbox}\hfill\begin{tcolorbox}[enhanced,nobeforeafter,width=7.1cm,height=2.1cm,valign=center,colback=cgeo!6,colframe=cgeo,arc=2.5mm,boxrule=0.9pt,left=3mm]{\Large\bfseries\color{cgeo}d)}\quad\begin{minipage}[c]{5.3cm}\raggedright\Large $7~\text{cm}$\end{minipage}\end{tcolorbox}
\end{tcolorbox}
\newpage
\begin{tcolorbox}[enhanced,width=16cm,height=14cm,colback=white,colframe=cgeo,boxrule=2pt,arc=0mm,
left=6mm,right=6mm,top=0mm,bottom=4mm,sharp corners,
overlay={\fill[cgeo] (frame.north west) rectangle ([yshift=-1.25cm]frame.north east);
\node[anchor=west,text=white,font=\sffamily\bfseries\LARGE] at ([xshift=5mm,yshift=-0.62cm]frame.north west) {Question 50};
\node[anchor=east,text=white,font=\sffamily\itshape\large] at ([xshift=-5mm,yshift=-0.62cm]frame.north east) {Géométrie};}]
\vspace*{1.75cm}
\begin{minipage}[t][6.9cm][c]{\linewidth}\Large\centering
\begin{minipage}{0.95\linewidth}\raggedright\linespread{1.25}\selectfont Les droites $(AC)$ et $(BD)$ se coupent en $O$, et les droites $(AB)$ et $(CD)$ sont parallèles. $OA=2$ cm, $OC=5$ cm et $AB=3$ cm. Combien mesure $CD$~?\end{minipage}
\end{minipage}

\vfill
\begin{tcolorbox}[enhanced,nobeforeafter,width=7.1cm,height=2.1cm,valign=center,colback=cgeo!6,colframe=cgeo,arc=2.5mm,boxrule=0.9pt,left=3mm]{\Large\bfseries\color{cgeo}a)}\quad\begin{minipage}[c]{5.3cm}\raggedright\Large $7{,}5~\text{cm}$\end{minipage}\end{tcolorbox}\hfill\begin{tcolorbox}[enhanced,nobeforeafter,width=7.1cm,height=2.1cm,valign=center,colback=cgeo!6,colframe=cgeo,arc=2.5mm,boxrule=0.9pt,left=3mm]{\Large\bfseries\color{cgeo}b)}\quad\begin{minipage}[c]{5.3cm}\raggedright\Large $\dfrac{10}{3}~\text{cm}$\end{minipage}\end{tcolorbox}\par\vspace{2.5mm}
\begin{tcolorbox}[enhanced,nobeforeafter,width=7.1cm,height=2.1cm,valign=center,colback=cgeo!6,colframe=cgeo,arc=2.5mm,boxrule=0.9pt,left=3mm]{\Large\bfseries\color{cgeo}c)}\quad\begin{minipage}[c]{5.3cm}\raggedright\Large $1{,}2~\text{cm}$\end{minipage}\end{tcolorbox}\hfill\begin{tcolorbox}[enhanced,nobeforeafter,width=7.1cm,height=2.1cm,valign=center,colback=cgeo!6,colframe=cgeo,arc=2.5mm,boxrule=0.9pt,left=3mm]{\Large\bfseries\color{cgeo}d)}\quad\begin{minipage}[c]{5.3cm}\raggedright\Large $6~\text{cm}$\end{minipage}\end{tcolorbox}
\end{tcolorbox}
\newpage
\begin{tcolorbox}[enhanced,width=16cm,height=14cm,colback=white,colframe=cgeo,boxrule=2pt,arc=0mm,
left=6mm,right=6mm,top=0mm,bottom=4mm,sharp corners,
overlay={\fill[cgeo] (frame.north west) rectangle ([yshift=-1.25cm]frame.north east);
\node[anchor=west,text=white,font=\sffamily\bfseries\LARGE] at ([xshift=5mm,yshift=-0.62cm]frame.north west) {Question 51};
\node[anchor=east,text=white,font=\sffamily\itshape\large] at ([xshift=-5mm,yshift=-0.62cm]frame.north east) {Géométrie};}]
\vspace*{1.75cm}
\begin{minipage}[t][6.9cm][c]{\linewidth}\Large\centering
\begin{minipage}{0.95\linewidth}\raggedright\linespread{1.25}\selectfont Les points $A$, $M$, $B$ sont alignés dans cet ordre, ainsi que $A$, $N$, $C$. $AM=3$, $AB=9$, $AN=4$ et $AC=12$. Les droites $(MN)$ et $(BC)$ sont-elles parallèles~?\end{minipage}
\end{minipage}

\vfill
\begin{tcolorbox}[enhanced,nobeforeafter,width=7.1cm,height=2.1cm,valign=center,colback=cgeo!6,colframe=cgeo,arc=2.5mm,boxrule=0.9pt,left=3mm]{\Large\bfseries\color{cgeo}a)}\quad\begin{minipage}[c]{5.3cm}\raggedright\Large non\end{minipage}\end{tcolorbox}\hfill\begin{tcolorbox}[enhanced,nobeforeafter,width=7.1cm,height=2.1cm,valign=center,colback=cgeo!6,colframe=cgeo,arc=2.5mm,boxrule=0.9pt,left=3mm]{\Large\bfseries\color{cgeo}b)}\quad\begin{minipage}[c]{5.3cm}\raggedright\Large oui\end{minipage}\end{tcolorbox}\par\vspace{2.5mm}
\begin{tcolorbox}[enhanced,nobeforeafter,width=7.1cm,height=2.1cm,valign=center,colback=cgeo!6,colframe=cgeo,arc=2.5mm,boxrule=0.9pt,left=3mm]{\Large\bfseries\color{cgeo}c)}\quad\begin{minipage}[c]{5.3cm}\raggedright\large seulement si le triangle est rectangle\end{minipage}\end{tcolorbox}\hfill\begin{tcolorbox}[enhanced,nobeforeafter,width=7.1cm,height=2.1cm,valign=center,colback=cgeo!6,colframe=cgeo,arc=2.5mm,boxrule=0.9pt,left=3mm]{\Large\bfseries\color{cgeo}d)}\quad\begin{minipage}[c]{5.3cm}\raggedright\Large on ne peut pas savoir\end{minipage}\end{tcolorbox}
\end{tcolorbox}
\newpage
\begin{tcolorbox}[enhanced,width=16cm,height=14cm,colback=white,colframe=cgeo,boxrule=2pt,arc=0mm,
left=6mm,right=6mm,top=0mm,bottom=4mm,sharp corners,
overlay={\fill[cgeo] (frame.north west) rectangle ([yshift=-1.25cm]frame.north east);
\node[anchor=west,text=white,font=\sffamily\bfseries\LARGE] at ([xshift=5mm,yshift=-0.62cm]frame.north west) {Question 52};
\node[anchor=east,text=white,font=\sffamily\itshape\large] at ([xshift=-5mm,yshift=-0.62cm]frame.north east) {Géométrie};}]
\vspace*{1.75cm}
\begin{minipage}[t][6.9cm][c]{\linewidth}\Large\centering
\begin{minipage}{0.95\linewidth}\raggedright\linespread{1.25}\selectfont Les points $A$, $M$, $B$ sont alignés dans cet ordre, ainsi que $A$, $N$, $C$. $AM=2$, $AB=5$, $AN=3$ et $AC=8$. Les droites $(MN)$ et $(BC)$ sont-elles parallèles~?\end{minipage}
\end{minipage}

\vfill
\begin{tcolorbox}[enhanced,nobeforeafter,width=7.1cm,height=2.1cm,valign=center,colback=cgeo!6,colframe=cgeo,arc=2.5mm,boxrule=0.9pt,left=3mm]{\Large\bfseries\color{cgeo}a)}\quad\begin{minipage}[c]{5.3cm}\raggedright\large seulement si le triangle est rectangle\end{minipage}\end{tcolorbox}\hfill\begin{tcolorbox}[enhanced,nobeforeafter,width=7.1cm,height=2.1cm,valign=center,colback=cgeo!6,colframe=cgeo,arc=2.5mm,boxrule=0.9pt,left=3mm]{\Large\bfseries\color{cgeo}b)}\quad\begin{minipage}[c]{5.3cm}\raggedright\Large on ne peut pas savoir\end{minipage}\end{tcolorbox}\par\vspace{2.5mm}
\begin{tcolorbox}[enhanced,nobeforeafter,width=7.1cm,height=2.1cm,valign=center,colback=cgeo!6,colframe=cgeo,arc=2.5mm,boxrule=0.9pt,left=3mm]{\Large\bfseries\color{cgeo}c)}\quad\begin{minipage}[c]{5.3cm}\raggedright\Large oui\end{minipage}\end{tcolorbox}\hfill\begin{tcolorbox}[enhanced,nobeforeafter,width=7.1cm,height=2.1cm,valign=center,colback=cgeo!6,colframe=cgeo,arc=2.5mm,boxrule=0.9pt,left=3mm]{\Large\bfseries\color{cgeo}d)}\quad\begin{minipage}[c]{5.3cm}\raggedright\Large non\end{minipage}\end{tcolorbox}
\end{tcolorbox}
\newpage
\begin{tcolorbox}[enhanced,width=16cm,height=14cm,colback=white,colframe=cgeo,boxrule=2pt,arc=0mm,
left=6mm,right=6mm,top=0mm,bottom=4mm,sharp corners,
overlay={\fill[cgeo] (frame.north west) rectangle ([yshift=-1.25cm]frame.north east);
\node[anchor=west,text=white,font=\sffamily\bfseries\LARGE] at ([xshift=5mm,yshift=-0.62cm]frame.north west) {Question 53};
\node[anchor=east,text=white,font=\sffamily\itshape\large] at ([xshift=-5mm,yshift=-0.62cm]frame.north east) {Géométrie};}]
\vspace*{1.75cm}
\begin{minipage}[t][6.9cm][c]{\linewidth}\LARGE\centering
\begin{minipage}{0.95\linewidth}\raggedright\linespread{1.25}\selectfont Le triangle $ABC$ est rectangle en $A$, avec $BC=10$ cm et $\sin\widehat{ABC}=0{,}6$. Combien mesure $AC$~?\end{minipage}
\end{minipage}

\vfill
\begin{tcolorbox}[enhanced,nobeforeafter,width=7.1cm,height=2.1cm,valign=center,colback=cgeo!6,colframe=cgeo,arc=2.5mm,boxrule=0.9pt,left=3mm]{\Large\bfseries\color{cgeo}a)}\quad\begin{minipage}[c]{5.3cm}\raggedright\Large $6~\text{cm}$\end{minipage}\end{tcolorbox}\hfill\begin{tcolorbox}[enhanced,nobeforeafter,width=7.1cm,height=2.1cm,valign=center,colback=cgeo!6,colframe=cgeo,arc=2.5mm,boxrule=0.9pt,left=3mm]{\Large\bfseries\color{cgeo}b)}\quad\begin{minipage}[c]{5.3cm}\raggedright\Large $12~\text{cm}$\end{minipage}\end{tcolorbox}\par\vspace{2.5mm}
\begin{tcolorbox}[enhanced,nobeforeafter,width=7.1cm,height=2.1cm,valign=center,colback=cgeo!6,colframe=cgeo,arc=2.5mm,boxrule=0.9pt,left=3mm]{\Large\bfseries\color{cgeo}c)}\quad\begin{minipage}[c]{5.3cm}\raggedright\Large $4~\text{cm}$\end{minipage}\end{tcolorbox}\hfill\begin{tcolorbox}[enhanced,nobeforeafter,width=7.1cm,height=2.1cm,valign=center,colback=cgeo!6,colframe=cgeo,arc=2.5mm,boxrule=0.9pt,left=3mm]{\Large\bfseries\color{cgeo}d)}\quad\begin{minipage}[c]{5.3cm}\raggedright\Large $3~\text{cm}$\end{minipage}\end{tcolorbox}
\end{tcolorbox}
\newpage
\begin{tcolorbox}[enhanced,width=16cm,height=14cm,colback=white,colframe=cgeo,boxrule=2pt,arc=0mm,
left=6mm,right=6mm,top=0mm,bottom=4mm,sharp corners,
overlay={\fill[cgeo] (frame.north west) rectangle ([yshift=-1.25cm]frame.north east);
\node[anchor=west,text=white,font=\sffamily\bfseries\LARGE] at ([xshift=5mm,yshift=-0.62cm]frame.north west) {Question 54};
\node[anchor=east,text=white,font=\sffamily\itshape\large] at ([xshift=-5mm,yshift=-0.62cm]frame.north east) {Géométrie};}]
\vspace*{1.75cm}
\begin{minipage}[t][6.9cm][c]{\linewidth}\LARGE\centering
\begin{minipage}{0.95\linewidth}\raggedright\linespread{1.25}\selectfont Le triangle $ABC$ est rectangle en $A$, avec $BC=8$ cm et $\sin\widehat{ABC}=0{,}5$. Combien mesure $AC$~?\end{minipage}
\end{minipage}

\vfill
\begin{tcolorbox}[enhanced,nobeforeafter,width=7.1cm,height=2.1cm,valign=center,colback=cgeo!6,colframe=cgeo,arc=2.5mm,boxrule=0.9pt,left=3mm]{\Large\bfseries\color{cgeo}a)}\quad\begin{minipage}[c]{5.3cm}\raggedright\Large $5~\text{cm}$\end{minipage}\end{tcolorbox}\hfill\begin{tcolorbox}[enhanced,nobeforeafter,width=7.1cm,height=2.1cm,valign=center,colback=cgeo!6,colframe=cgeo,arc=2.5mm,boxrule=0.9pt,left=3mm]{\Large\bfseries\color{cgeo}b)}\quad\begin{minipage}[c]{5.3cm}\raggedright\Large $2~\text{cm}$\end{minipage}\end{tcolorbox}\par\vspace{2.5mm}
\begin{tcolorbox}[enhanced,nobeforeafter,width=7.1cm,height=2.1cm,valign=center,colback=cgeo!6,colframe=cgeo,arc=2.5mm,boxrule=0.9pt,left=3mm]{\Large\bfseries\color{cgeo}c)}\quad\begin{minipage}[c]{5.3cm}\raggedright\Large $8~\text{cm}$\end{minipage}\end{tcolorbox}\hfill\begin{tcolorbox}[enhanced,nobeforeafter,width=7.1cm,height=2.1cm,valign=center,colback=cgeo!6,colframe=cgeo,arc=2.5mm,boxrule=0.9pt,left=3mm]{\Large\bfseries\color{cgeo}d)}\quad\begin{minipage}[c]{5.3cm}\raggedright\Large $4~\text{cm}$\end{minipage}\end{tcolorbox}
\end{tcolorbox}
\newpage
\begin{tcolorbox}[enhanced,width=16cm,height=14cm,colback=white,colframe=cgeo,boxrule=2pt,arc=0mm,
left=6mm,right=6mm,top=0mm,bottom=4mm,sharp corners,
overlay={\fill[cgeo] (frame.north west) rectangle ([yshift=-1.25cm]frame.north east);
\node[anchor=west,text=white,font=\sffamily\bfseries\LARGE] at ([xshift=5mm,yshift=-0.62cm]frame.north west) {Question 55};
\node[anchor=east,text=white,font=\sffamily\itshape\large] at ([xshift=-5mm,yshift=-0.62cm]frame.north east) {Géométrie};}]
\vspace*{1.75cm}
\begin{minipage}[t][6.9cm][c]{\linewidth}\LARGE\centering
\begin{minipage}{0.95\linewidth}\raggedright\linespread{1.25}\selectfont Le triangle $ABC$ est rectangle en $A$, avec $AB=4$ cm et $AC=3$ cm. Que vaut $\tan\widehat{ABC}$~?\end{minipage}
\end{minipage}

\vfill
\begin{tcolorbox}[enhanced,nobeforeafter,width=7.1cm,height=2.1cm,valign=center,colback=cgeo!6,colframe=cgeo,arc=2.5mm,boxrule=0.9pt,left=3mm]{\Large\bfseries\color{cgeo}a)}\quad\begin{minipage}[c]{5.3cm}\raggedright\Large $\dfrac{4}{3}$\end{minipage}\end{tcolorbox}\hfill\begin{tcolorbox}[enhanced,nobeforeafter,width=7.1cm,height=2.1cm,valign=center,colback=cgeo!6,colframe=cgeo,arc=2.5mm,boxrule=0.9pt,left=3mm]{\Large\bfseries\color{cgeo}b)}\quad\begin{minipage}[c]{5.3cm}\raggedright\Large $\dfrac{7}{4}$\end{minipage}\end{tcolorbox}\par\vspace{2.5mm}
\begin{tcolorbox}[enhanced,nobeforeafter,width=7.1cm,height=2.1cm,valign=center,colback=cgeo!6,colframe=cgeo,arc=2.5mm,boxrule=0.9pt,left=3mm]{\Large\bfseries\color{cgeo}c)}\quad\begin{minipage}[c]{5.3cm}\raggedright\Large $\dfrac{4}{7}$\end{minipage}\end{tcolorbox}\hfill\begin{tcolorbox}[enhanced,nobeforeafter,width=7.1cm,height=2.1cm,valign=center,colback=cgeo!6,colframe=cgeo,arc=2.5mm,boxrule=0.9pt,left=3mm]{\Large\bfseries\color{cgeo}d)}\quad\begin{minipage}[c]{5.3cm}\raggedright\Large $\dfrac{3}{4}$\end{minipage}\end{tcolorbox}
\end{tcolorbox}
\newpage
\begin{tcolorbox}[enhanced,width=16cm,height=14cm,colback=white,colframe=cgeo,boxrule=2pt,arc=0mm,
left=6mm,right=6mm,top=0mm,bottom=4mm,sharp corners,
overlay={\fill[cgeo] (frame.north west) rectangle ([yshift=-1.25cm]frame.north east);
\node[anchor=west,text=white,font=\sffamily\bfseries\LARGE] at ([xshift=5mm,yshift=-0.62cm]frame.north west) {Question 56};
\node[anchor=east,text=white,font=\sffamily\itshape\large] at ([xshift=-5mm,yshift=-0.62cm]frame.north east) {Géométrie};}]
\vspace*{1.75cm}
\begin{minipage}[t][6.9cm][c]{\linewidth}\LARGE\centering
\begin{minipage}{0.95\linewidth}\raggedright\linespread{1.25}\selectfont Le triangle $ABC$ est rectangle en $A$, avec $AB=5$ cm et $AC=12$ cm. Que vaut $\tan\widehat{ABC}$~?\end{minipage}
\end{minipage}

\vfill
\begin{tcolorbox}[enhanced,nobeforeafter,width=7.1cm,height=2.1cm,valign=center,colback=cgeo!6,colframe=cgeo,arc=2.5mm,boxrule=0.9pt,left=3mm]{\Large\bfseries\color{cgeo}a)}\quad\begin{minipage}[c]{5.3cm}\raggedright\Large $\dfrac{17}{5}$\end{minipage}\end{tcolorbox}\hfill\begin{tcolorbox}[enhanced,nobeforeafter,width=7.1cm,height=2.1cm,valign=center,colback=cgeo!6,colframe=cgeo,arc=2.5mm,boxrule=0.9pt,left=3mm]{\Large\bfseries\color{cgeo}b)}\quad\begin{minipage}[c]{5.3cm}\raggedright\Large $\dfrac{5}{12}$\end{minipage}\end{tcolorbox}\par\vspace{2.5mm}
\begin{tcolorbox}[enhanced,nobeforeafter,width=7.1cm,height=2.1cm,valign=center,colback=cgeo!6,colframe=cgeo,arc=2.5mm,boxrule=0.9pt,left=3mm]{\Large\bfseries\color{cgeo}c)}\quad\begin{minipage}[c]{5.3cm}\raggedright\Large $\dfrac{12}{5}$\end{minipage}\end{tcolorbox}\hfill\begin{tcolorbox}[enhanced,nobeforeafter,width=7.1cm,height=2.1cm,valign=center,colback=cgeo!6,colframe=cgeo,arc=2.5mm,boxrule=0.9pt,left=3mm]{\Large\bfseries\color{cgeo}d)}\quad\begin{minipage}[c]{5.3cm}\raggedright\Large $\dfrac{5}{17}$\end{minipage}\end{tcolorbox}
\end{tcolorbox}
\newpage
\begin{tcolorbox}[enhanced,width=16cm,height=14cm,colback=white,colframe=cgeo,boxrule=2pt,arc=0mm,
left=6mm,right=6mm,top=0mm,bottom=4mm,sharp corners,
overlay={\fill[cgeo] (frame.north west) rectangle ([yshift=-1.25cm]frame.north east);
\node[anchor=west,text=white,font=\sffamily\bfseries\LARGE] at ([xshift=5mm,yshift=-0.62cm]frame.north west) {Question 57};
\node[anchor=east,text=white,font=\sffamily\itshape\large] at ([xshift=-5mm,yshift=-0.62cm]frame.north east) {Géométrie};}]
\vspace*{1.75cm}
\begin{minipage}[t][6.9cm][c]{\linewidth}\LARGE\centering
\begin{minipage}{0.95\linewidth}\raggedright\linespread{1.25}\selectfont Une homothétie de rapport $-2$ transforme le segment $[AB]$, de longueur $3$ cm, en $[A'B']$. Combien mesure $A'B'$~?\end{minipage}
\end{minipage}

\vfill
\begin{tcolorbox}[enhanced,nobeforeafter,width=7.1cm,height=2.1cm,valign=center,colback=cgeo!6,colframe=cgeo,arc=2.5mm,boxrule=0.9pt,left=3mm]{\Large\bfseries\color{cgeo}a)}\quad\begin{minipage}[c]{5.3cm}\raggedright\Large $5~\text{cm}$\end{minipage}\end{tcolorbox}\hfill\begin{tcolorbox}[enhanced,nobeforeafter,width=7.1cm,height=2.1cm,valign=center,colback=cgeo!6,colframe=cgeo,arc=2.5mm,boxrule=0.9pt,left=3mm]{\Large\bfseries\color{cgeo}b)}\quad\begin{minipage}[c]{5.3cm}\raggedright\Large $6~\text{cm}$\end{minipage}\end{tcolorbox}\par\vspace{2.5mm}
\begin{tcolorbox}[enhanced,nobeforeafter,width=7.1cm,height=2.1cm,valign=center,colback=cgeo!6,colframe=cgeo,arc=2.5mm,boxrule=0.9pt,left=3mm]{\Large\bfseries\color{cgeo}c)}\quad\begin{minipage}[c]{5.3cm}\raggedright\Large $1{,}5~\text{cm}$\end{minipage}\end{tcolorbox}\hfill\begin{tcolorbox}[enhanced,nobeforeafter,width=7.1cm,height=2.1cm,valign=center,colback=cgeo!6,colframe=cgeo,arc=2.5mm,boxrule=0.9pt,left=3mm]{\Large\bfseries\color{cgeo}d)}\quad\begin{minipage}[c]{5.3cm}\raggedright\Large $7~\text{cm}$\end{minipage}\end{tcolorbox}
\end{tcolorbox}
\newpage
\begin{tcolorbox}[enhanced,width=16cm,height=14cm,colback=white,colframe=cgeo,boxrule=2pt,arc=0mm,
left=6mm,right=6mm,top=0mm,bottom=4mm,sharp corners,
overlay={\fill[cgeo] (frame.north west) rectangle ([yshift=-1.25cm]frame.north east);
\node[anchor=west,text=white,font=\sffamily\bfseries\LARGE] at ([xshift=5mm,yshift=-0.62cm]frame.north west) {Question 58};
\node[anchor=east,text=white,font=\sffamily\itshape\large] at ([xshift=-5mm,yshift=-0.62cm]frame.north east) {Géométrie};}]
\vspace*{1.75cm}
\begin{minipage}[t][6.9cm][c]{\linewidth}\LARGE\centering
\begin{minipage}{0.95\linewidth}\raggedright\linespread{1.25}\selectfont Une homothétie de rapport $3$ transforme le segment $[AB]$, de longueur $2{,}5$ cm, en $[A'B']$. Combien mesure $A'B'$~?\end{minipage}
\end{minipage}

\vfill
\begin{tcolorbox}[enhanced,nobeforeafter,width=7.1cm,height=2.1cm,valign=center,colback=cgeo!6,colframe=cgeo,arc=2.5mm,boxrule=0.9pt,left=3mm]{\Large\bfseries\color{cgeo}a)}\quad\begin{minipage}[c]{5.3cm}\raggedright\Large $5{,}5~\text{cm}$\end{minipage}\end{tcolorbox}\hfill\begin{tcolorbox}[enhanced,nobeforeafter,width=7.1cm,height=2.1cm,valign=center,colback=cgeo!6,colframe=cgeo,arc=2.5mm,boxrule=0.9pt,left=3mm]{\Large\bfseries\color{cgeo}b)}\quad\begin{minipage}[c]{5.3cm}\raggedright\Large $8{,}5~\text{cm}$\end{minipage}\end{tcolorbox}\par\vspace{2.5mm}
\begin{tcolorbox}[enhanced,nobeforeafter,width=7.1cm,height=2.1cm,valign=center,colback=cgeo!6,colframe=cgeo,arc=2.5mm,boxrule=0.9pt,left=3mm]{\Large\bfseries\color{cgeo}c)}\quad\begin{minipage}[c]{5.3cm}\raggedright\Large $\dfrac{5}{6}~\text{cm}$\end{minipage}\end{tcolorbox}\hfill\begin{tcolorbox}[enhanced,nobeforeafter,width=7.1cm,height=2.1cm,valign=center,colback=cgeo!6,colframe=cgeo,arc=2.5mm,boxrule=0.9pt,left=3mm]{\Large\bfseries\color{cgeo}d)}\quad\begin{minipage}[c]{5.3cm}\raggedright\Large $7{,}5~\text{cm}$\end{minipage}\end{tcolorbox}
\end{tcolorbox}
\newpage
\begin{tcolorbox}[enhanced,width=16cm,height=14cm,colback=white,colframe=cgeo,boxrule=2pt,arc=0mm,
left=6mm,right=6mm,top=0mm,bottom=4mm,sharp corners,
overlay={\fill[cgeo] (frame.north west) rectangle ([yshift=-1.25cm]frame.north east);
\node[anchor=west,text=white,font=\sffamily\bfseries\LARGE] at ([xshift=5mm,yshift=-0.62cm]frame.north west) {Question 59};
\node[anchor=east,text=white,font=\sffamily\itshape\large] at ([xshift=-5mm,yshift=-0.62cm]frame.north east) {Géométrie};}]
\vspace*{1.75cm}
\begin{minipage}[t][6.9cm][c]{\linewidth}\LARGE\centering
\begin{minipage}{0.95\linewidth}\raggedright\linespread{1.25}\selectfont Quel est le volume exact d'une boule de rayon $3$ cm~?\end{minipage}
\end{minipage}

\vfill
\begin{tcolorbox}[enhanced,nobeforeafter,width=7.1cm,height=2.1cm,valign=center,colback=cgeo!6,colframe=cgeo,arc=2.5mm,boxrule=0.9pt,left=3mm]{\Large\bfseries\color{cgeo}a)}\quad\begin{minipage}[c]{5.3cm}\raggedright\Large $27\pi~\text{cm}^3$\end{minipage}\end{tcolorbox}\hfill\begin{tcolorbox}[enhanced,nobeforeafter,width=7.1cm,height=2.1cm,valign=center,colback=cgeo!6,colframe=cgeo,arc=2.5mm,boxrule=0.9pt,left=3mm]{\Large\bfseries\color{cgeo}b)}\quad\begin{minipage}[c]{5.3cm}\raggedright\Large $72\pi~\text{cm}^3$\end{minipage}\end{tcolorbox}\par\vspace{2.5mm}
\begin{tcolorbox}[enhanced,nobeforeafter,width=7.1cm,height=2.1cm,valign=center,colback=cgeo!6,colframe=cgeo,arc=2.5mm,boxrule=0.9pt,left=3mm]{\Large\bfseries\color{cgeo}c)}\quad\begin{minipage}[c]{5.3cm}\raggedright\Large $12\pi~\text{cm}^3$\end{minipage}\end{tcolorbox}\hfill\begin{tcolorbox}[enhanced,nobeforeafter,width=7.1cm,height=2.1cm,valign=center,colback=cgeo!6,colframe=cgeo,arc=2.5mm,boxrule=0.9pt,left=3mm]{\Large\bfseries\color{cgeo}d)}\quad\begin{minipage}[c]{5.3cm}\raggedright\Large $36\pi~\text{cm}^3$\end{minipage}\end{tcolorbox}
\end{tcolorbox}
\newpage
\begin{tcolorbox}[enhanced,width=16cm,height=14cm,colback=white,colframe=cgeo,boxrule=2pt,arc=0mm,
left=6mm,right=6mm,top=0mm,bottom=4mm,sharp corners,
overlay={\fill[cgeo] (frame.north west) rectangle ([yshift=-1.25cm]frame.north east);
\node[anchor=west,text=white,font=\sffamily\bfseries\LARGE] at ([xshift=5mm,yshift=-0.62cm]frame.north west) {Question 60};
\node[anchor=east,text=white,font=\sffamily\itshape\large] at ([xshift=-5mm,yshift=-0.62cm]frame.north east) {Géométrie};}]
\vspace*{1.75cm}
\begin{minipage}[t][6.9cm][c]{\linewidth}\LARGE\centering
\begin{minipage}{0.95\linewidth}\raggedright\linespread{1.25}\selectfont Quel est le volume exact d'une boule de rayon $6$ cm~?\end{minipage}
\end{minipage}

\vfill
\begin{tcolorbox}[enhanced,nobeforeafter,width=7.1cm,height=2.1cm,valign=center,colback=cgeo!6,colframe=cgeo,arc=2.5mm,boxrule=0.9pt,left=3mm]{\Large\bfseries\color{cgeo}a)}\quad\begin{minipage}[c]{5.3cm}\raggedright\Large $48\pi~\text{cm}^3$\end{minipage}\end{tcolorbox}\hfill\begin{tcolorbox}[enhanced,nobeforeafter,width=7.1cm,height=2.1cm,valign=center,colback=cgeo!6,colframe=cgeo,arc=2.5mm,boxrule=0.9pt,left=3mm]{\Large\bfseries\color{cgeo}b)}\quad\begin{minipage}[c]{5.3cm}\raggedright\Large $216\pi~\text{cm}^3$\end{minipage}\end{tcolorbox}\par\vspace{2.5mm}
\begin{tcolorbox}[enhanced,nobeforeafter,width=7.1cm,height=2.1cm,valign=center,colback=cgeo!6,colframe=cgeo,arc=2.5mm,boxrule=0.9pt,left=3mm]{\Large\bfseries\color{cgeo}c)}\quad\begin{minipage}[c]{5.3cm}\raggedright\Large $144\pi~\text{cm}^3$\end{minipage}\end{tcolorbox}\hfill\begin{tcolorbox}[enhanced,nobeforeafter,width=7.1cm,height=2.1cm,valign=center,colback=cgeo!6,colframe=cgeo,arc=2.5mm,boxrule=0.9pt,left=3mm]{\Large\bfseries\color{cgeo}d)}\quad\begin{minipage}[c]{5.3cm}\raggedright\Large $288\pi~\text{cm}^3$\end{minipage}\end{tcolorbox}
\end{tcolorbox}
\newpage
\begin{tcolorbox}[enhanced,width=16cm,height=14cm,colback=white,colframe=cgeo,boxrule=2pt,arc=0mm,
left=6mm,right=6mm,top=0mm,bottom=4mm,sharp corners,
overlay={\fill[cgeo] (frame.north west) rectangle ([yshift=-1.25cm]frame.north east);
\node[anchor=west,text=white,font=\sffamily\bfseries\LARGE] at ([xshift=5mm,yshift=-0.62cm]frame.north west) {Question 61};
\node[anchor=east,text=white,font=\sffamily\itshape\large] at ([xshift=-5mm,yshift=-0.62cm]frame.north east) {Géométrie};}]
\vspace*{1.75cm}
\begin{minipage}[t][6.9cm][c]{\linewidth}\LARGE\centering
\begin{minipage}{0.95\linewidth}\raggedright\linespread{1.25}\selectfont Quelle est l'aire exacte d'une sphère de rayon $5$ cm~?\end{minipage}
\end{minipage}

\vfill
\begin{tcolorbox}[enhanced,nobeforeafter,width=7.1cm,height=2.1cm,valign=center,colback=cgeo!6,colframe=cgeo,arc=2.5mm,boxrule=0.9pt,left=3mm]{\Large\bfseries\color{cgeo}a)}\quad\begin{minipage}[c]{5.3cm}\raggedright\Large $50\pi~\text{cm}^2$\end{minipage}\end{tcolorbox}\hfill\begin{tcolorbox}[enhanced,nobeforeafter,width=7.1cm,height=2.1cm,valign=center,colback=cgeo!6,colframe=cgeo,arc=2.5mm,boxrule=0.9pt,left=3mm]{\Large\bfseries\color{cgeo}b)}\quad\begin{minipage}[c]{5.3cm}\raggedright\Large $25\pi~\text{cm}^2$\end{minipage}\end{tcolorbox}\par\vspace{2.5mm}
\begin{tcolorbox}[enhanced,nobeforeafter,width=7.1cm,height=2.1cm,valign=center,colback=cgeo!6,colframe=cgeo,arc=2.5mm,boxrule=0.9pt,left=3mm]{\Large\bfseries\color{cgeo}c)}\quad\begin{minipage}[c]{5.3cm}\raggedright\Large $100\pi~\text{cm}^2$\end{minipage}\end{tcolorbox}\hfill\begin{tcolorbox}[enhanced,nobeforeafter,width=7.1cm,height=2.1cm,valign=center,colback=cgeo!6,colframe=cgeo,arc=2.5mm,boxrule=0.9pt,left=3mm]{\Large\bfseries\color{cgeo}d)}\quad\begin{minipage}[c]{5.3cm}\raggedright\large $\dfrac{500}{3}\pi~\text{cm}^2$\end{minipage}\end{tcolorbox}
\end{tcolorbox}
\newpage
\begin{tcolorbox}[enhanced,width=16cm,height=14cm,colback=white,colframe=cgeo,boxrule=2pt,arc=0mm,
left=6mm,right=6mm,top=0mm,bottom=4mm,sharp corners,
overlay={\fill[cgeo] (frame.north west) rectangle ([yshift=-1.25cm]frame.north east);
\node[anchor=west,text=white,font=\sffamily\bfseries\LARGE] at ([xshift=5mm,yshift=-0.62cm]frame.north west) {Question 62};
\node[anchor=east,text=white,font=\sffamily\itshape\large] at ([xshift=-5mm,yshift=-0.62cm]frame.north east) {Géométrie};}]
\vspace*{1.75cm}
\begin{minipage}[t][6.9cm][c]{\linewidth}\LARGE\centering
\begin{minipage}{0.95\linewidth}\raggedright\linespread{1.25}\selectfont Quelle est l'aire exacte d'une sphère de rayon $3$ cm~?\end{minipage}
\end{minipage}

\vfill
\begin{tcolorbox}[enhanced,nobeforeafter,width=7.1cm,height=2.1cm,valign=center,colback=cgeo!6,colframe=cgeo,arc=2.5mm,boxrule=0.9pt,left=3mm]{\Large\bfseries\color{cgeo}a)}\quad\begin{minipage}[c]{5.3cm}\raggedright\Large $9\pi~\text{cm}^2$\end{minipage}\end{tcolorbox}\hfill\begin{tcolorbox}[enhanced,nobeforeafter,width=7.1cm,height=2.1cm,valign=center,colback=cgeo!6,colframe=cgeo,arc=2.5mm,boxrule=0.9pt,left=3mm]{\Large\bfseries\color{cgeo}b)}\quad\begin{minipage}[c]{5.3cm}\raggedright\Large $36\pi~\text{cm}^2$\end{minipage}\end{tcolorbox}\par\vspace{2.5mm}
\begin{tcolorbox}[enhanced,nobeforeafter,width=7.1cm,height=2.1cm,valign=center,colback=cgeo!6,colframe=cgeo,arc=2.5mm,boxrule=0.9pt,left=3mm]{\Large\bfseries\color{cgeo}c)}\quad\begin{minipage}[c]{5.3cm}\raggedright\Large $12\pi~\text{cm}^2$\end{minipage}\end{tcolorbox}\hfill\begin{tcolorbox}[enhanced,nobeforeafter,width=7.1cm,height=2.1cm,valign=center,colback=cgeo!6,colframe=cgeo,arc=2.5mm,boxrule=0.9pt,left=3mm]{\Large\bfseries\color{cgeo}d)}\quad\begin{minipage}[c]{5.3cm}\raggedright\Large $18\pi~\text{cm}^2$\end{minipage}\end{tcolorbox}
\end{tcolorbox}
\newpage
\begin{tcolorbox}[enhanced,width=16cm,height=14cm,colback=white,colframe=cgeo,boxrule=2pt,arc=0mm,
left=6mm,right=6mm,top=0mm,bottom=4mm,sharp corners,
overlay={\fill[cgeo] (frame.north west) rectangle ([yshift=-1.25cm]frame.north east);
\node[anchor=west,text=white,font=\sffamily\bfseries\LARGE] at ([xshift=5mm,yshift=-0.62cm]frame.north west) {Question 63};
\node[anchor=east,text=white,font=\sffamily\itshape\large] at ([xshift=-5mm,yshift=-0.62cm]frame.north east) {Géométrie};}]
\vspace*{1.75cm}
\begin{minipage}[t][6.9cm][c]{\linewidth}\LARGE\centering
\begin{minipage}{0.95\linewidth}\raggedright\linespread{1.25}\selectfont Un solide est agrandi à l'échelle $2$. Par combien son volume est-il multiplié~?\end{minipage}
\end{minipage}

\vfill
\begin{tcolorbox}[enhanced,nobeforeafter,width=7.1cm,height=2.1cm,valign=center,colback=cgeo!6,colframe=cgeo,arc=2.5mm,boxrule=0.9pt,left=3mm]{\Large\bfseries\color{cgeo}a)}\quad\begin{minipage}[c]{5.3cm}\raggedright\Large $2$\end{minipage}\end{tcolorbox}\hfill\begin{tcolorbox}[enhanced,nobeforeafter,width=7.1cm,height=2.1cm,valign=center,colback=cgeo!6,colframe=cgeo,arc=2.5mm,boxrule=0.9pt,left=3mm]{\Large\bfseries\color{cgeo}b)}\quad\begin{minipage}[c]{5.3cm}\raggedright\Large $4$\end{minipage}\end{tcolorbox}\par\vspace{2.5mm}
\begin{tcolorbox}[enhanced,nobeforeafter,width=7.1cm,height=2.1cm,valign=center,colback=cgeo!6,colframe=cgeo,arc=2.5mm,boxrule=0.9pt,left=3mm]{\Large\bfseries\color{cgeo}c)}\quad\begin{minipage}[c]{5.3cm}\raggedright\Large $8$\end{minipage}\end{tcolorbox}\hfill\begin{tcolorbox}[enhanced,nobeforeafter,width=7.1cm,height=2.1cm,valign=center,colback=cgeo!6,colframe=cgeo,arc=2.5mm,boxrule=0.9pt,left=3mm]{\Large\bfseries\color{cgeo}d)}\quad\begin{minipage}[c]{5.3cm}\raggedright\Large $6$\end{minipage}\end{tcolorbox}
\end{tcolorbox}
\newpage
\begin{tcolorbox}[enhanced,width=16cm,height=14cm,colback=white,colframe=cgeo,boxrule=2pt,arc=0mm,
left=6mm,right=6mm,top=0mm,bottom=4mm,sharp corners,
overlay={\fill[cgeo] (frame.north west) rectangle ([yshift=-1.25cm]frame.north east);
\node[anchor=west,text=white,font=\sffamily\bfseries\LARGE] at ([xshift=5mm,yshift=-0.62cm]frame.north west) {Question 64};
\node[anchor=east,text=white,font=\sffamily\itshape\large] at ([xshift=-5mm,yshift=-0.62cm]frame.north east) {Géométrie};}]
\vspace*{1.75cm}
\begin{minipage}[t][6.9cm][c]{\linewidth}\LARGE\centering
\begin{minipage}{0.95\linewidth}\raggedright\linespread{1.25}\selectfont Un solide est agrandi à l'échelle $3$. Par combien son volume est-il multiplié~?\end{minipage}
\end{minipage}

\vfill
\begin{tcolorbox}[enhanced,nobeforeafter,width=7.1cm,height=2.1cm,valign=center,colback=cgeo!6,colframe=cgeo,arc=2.5mm,boxrule=0.9pt,left=3mm]{\Large\bfseries\color{cgeo}a)}\quad\begin{minipage}[c]{5.3cm}\raggedright\Large $9$\end{minipage}\end{tcolorbox}\hfill\begin{tcolorbox}[enhanced,nobeforeafter,width=7.1cm,height=2.1cm,valign=center,colback=cgeo!6,colframe=cgeo,arc=2.5mm,boxrule=0.9pt,left=3mm]{\Large\bfseries\color{cgeo}b)}\quad\begin{minipage}[c]{5.3cm}\raggedright\Large $27$\end{minipage}\end{tcolorbox}\par\vspace{2.5mm}
\begin{tcolorbox}[enhanced,nobeforeafter,width=7.1cm,height=2.1cm,valign=center,colback=cgeo!6,colframe=cgeo,arc=2.5mm,boxrule=0.9pt,left=3mm]{\Large\bfseries\color{cgeo}c)}\quad\begin{minipage}[c]{5.3cm}\raggedright\Large $3$\end{minipage}\end{tcolorbox}\hfill\begin{tcolorbox}[enhanced,nobeforeafter,width=7.1cm,height=2.1cm,valign=center,colback=cgeo!6,colframe=cgeo,arc=2.5mm,boxrule=0.9pt,left=3mm]{\Large\bfseries\color{cgeo}d)}\quad\begin{minipage}[c]{5.3cm}\raggedright\Large $54$\end{minipage}\end{tcolorbox}
\end{tcolorbox}
\newpage
\begin{tcolorbox}[enhanced,width=16cm,height=14cm,colback=white,colframe=csta,boxrule=2pt,arc=0mm,
left=6mm,right=6mm,top=0mm,bottom=4mm,sharp corners,
overlay={\fill[csta] (frame.north west) rectangle ([yshift=-1.25cm]frame.north east);
\node[anchor=west,text=white,font=\sffamily\bfseries\LARGE] at ([xshift=5mm,yshift=-0.62cm]frame.north west) {Question 65};
\node[anchor=east,text=white,font=\sffamily\itshape\large] at ([xshift=-5mm,yshift=-0.62cm]frame.north east) {Statistiques et probabilités};}]
\vspace*{1.75cm}
\begin{minipage}[t][6.9cm][c]{\linewidth}\LARGE\centering
\begin{minipage}{0.95\linewidth}\raggedright\linespread{1.25}\selectfont Quelle est la médiane de la série $12\ ;\ 7\ ;\ 15\ ;\ 9\ ;\ 11$~?\end{minipage}
\end{minipage}

\vfill
\begin{tcolorbox}[enhanced,nobeforeafter,width=7.1cm,height=2.1cm,valign=center,colback=csta!6,colframe=csta,arc=2.5mm,boxrule=0.9pt,left=3mm]{\Large\bfseries\color{csta}a)}\quad\begin{minipage}[c]{5.3cm}\raggedright\Large $10{,}8$\end{minipage}\end{tcolorbox}\hfill\begin{tcolorbox}[enhanced,nobeforeafter,width=7.1cm,height=2.1cm,valign=center,colback=csta!6,colframe=csta,arc=2.5mm,boxrule=0.9pt,left=3mm]{\Large\bfseries\color{csta}b)}\quad\begin{minipage}[c]{5.3cm}\raggedright\Large $15$\end{minipage}\end{tcolorbox}\par\vspace{2.5mm}
\begin{tcolorbox}[enhanced,nobeforeafter,width=7.1cm,height=2.1cm,valign=center,colback=csta!6,colframe=csta,arc=2.5mm,boxrule=0.9pt,left=3mm]{\Large\bfseries\color{csta}c)}\quad\begin{minipage}[c]{5.3cm}\raggedright\Large $11$\end{minipage}\end{tcolorbox}\hfill\begin{tcolorbox}[enhanced,nobeforeafter,width=7.1cm,height=2.1cm,valign=center,colback=csta!6,colframe=csta,arc=2.5mm,boxrule=0.9pt,left=3mm]{\Large\bfseries\color{csta}d)}\quad\begin{minipage}[c]{5.3cm}\raggedright\Large $8$\end{minipage}\end{tcolorbox}
\end{tcolorbox}
\newpage
\begin{tcolorbox}[enhanced,width=16cm,height=14cm,colback=white,colframe=csta,boxrule=2pt,arc=0mm,
left=6mm,right=6mm,top=0mm,bottom=4mm,sharp corners,
overlay={\fill[csta] (frame.north west) rectangle ([yshift=-1.25cm]frame.north east);
\node[anchor=west,text=white,font=\sffamily\bfseries\LARGE] at ([xshift=5mm,yshift=-0.62cm]frame.north west) {Question 66};
\node[anchor=east,text=white,font=\sffamily\itshape\large] at ([xshift=-5mm,yshift=-0.62cm]frame.north east) {Statistiques et probabilités};}]
\vspace*{1.75cm}
\begin{minipage}[t][6.9cm][c]{\linewidth}\LARGE\centering
\begin{minipage}{0.95\linewidth}\raggedright\linespread{1.25}\selectfont Quelle est la médiane de la série $3\ ;\ 8\ ;\ 14\ ;\ 6\ ;\ 10\ ;\ 2$~?\end{minipage}
\end{minipage}

\vfill
\begin{tcolorbox}[enhanced,nobeforeafter,width=7.1cm,height=2.1cm,valign=center,colback=csta!6,colframe=csta,arc=2.5mm,boxrule=0.9pt,left=3mm]{\Large\bfseries\color{csta}a)}\quad\begin{minipage}[c]{5.3cm}\raggedright\Large $10$\end{minipage}\end{tcolorbox}\hfill\begin{tcolorbox}[enhanced,nobeforeafter,width=7.1cm,height=2.1cm,valign=center,colback=csta!6,colframe=csta,arc=2.5mm,boxrule=0.9pt,left=3mm]{\Large\bfseries\color{csta}b)}\quad\begin{minipage}[c]{5.3cm}\raggedright\Large $12$\end{minipage}\end{tcolorbox}\par\vspace{2.5mm}
\begin{tcolorbox}[enhanced,nobeforeafter,width=7.1cm,height=2.1cm,valign=center,colback=csta!6,colframe=csta,arc=2.5mm,boxrule=0.9pt,left=3mm]{\Large\bfseries\color{csta}c)}\quad\begin{minipage}[c]{5.3cm}\raggedright\Large $7$\end{minipage}\end{tcolorbox}\hfill\begin{tcolorbox}[enhanced,nobeforeafter,width=7.1cm,height=2.1cm,valign=center,colback=csta!6,colframe=csta,arc=2.5mm,boxrule=0.9pt,left=3mm]{\Large\bfseries\color{csta}d)}\quad\begin{minipage}[c]{5.3cm}\raggedright\Large $\dfrac{43}{6}$\end{minipage}\end{tcolorbox}
\end{tcolorbox}
\newpage
\begin{tcolorbox}[enhanced,width=16cm,height=14cm,colback=white,colframe=csta,boxrule=2pt,arc=0mm,
left=6mm,right=6mm,top=0mm,bottom=4mm,sharp corners,
overlay={\fill[csta] (frame.north west) rectangle ([yshift=-1.25cm]frame.north east);
\node[anchor=west,text=white,font=\sffamily\bfseries\LARGE] at ([xshift=5mm,yshift=-0.62cm]frame.north west) {Question 67};
\node[anchor=east,text=white,font=\sffamily\itshape\large] at ([xshift=-5mm,yshift=-0.62cm]frame.north east) {Statistiques et probabilités};}]
\vspace*{1.75cm}
\begin{minipage}[t][6.9cm][c]{\linewidth}\LARGE\centering
\begin{minipage}{0.95\linewidth}\raggedright\linespread{1.25}\selectfont Quelle est la moyenne de la série $14\ ;\ 9\ ;\ 16\ ;\ 11$~?\end{minipage}
\end{minipage}

\vfill
\begin{tcolorbox}[enhanced,nobeforeafter,width=7.1cm,height=2.1cm,valign=center,colback=csta!6,colframe=csta,arc=2.5mm,boxrule=0.9pt,left=3mm]{\Large\bfseries\color{csta}a)}\quad\begin{minipage}[c]{5.3cm}\raggedright\Large $12{,}5$\end{minipage}\end{tcolorbox}\hfill\begin{tcolorbox}[enhanced,nobeforeafter,width=7.1cm,height=2.1cm,valign=center,colback=csta!6,colframe=csta,arc=2.5mm,boxrule=0.9pt,left=3mm]{\Large\bfseries\color{csta}b)}\quad\begin{minipage}[c]{5.3cm}\raggedright\Large $14{,}5$\end{minipage}\end{tcolorbox}\par\vspace{2.5mm}
\begin{tcolorbox}[enhanced,nobeforeafter,width=7.1cm,height=2.1cm,valign=center,colback=csta!6,colframe=csta,arc=2.5mm,boxrule=0.9pt,left=3mm]{\Large\bfseries\color{csta}c)}\quad\begin{minipage}[c]{5.3cm}\raggedright\Large $13{,}5$\end{minipage}\end{tcolorbox}\hfill\begin{tcolorbox}[enhanced,nobeforeafter,width=7.1cm,height=2.1cm,valign=center,colback=csta!6,colframe=csta,arc=2.5mm,boxrule=0.9pt,left=3mm]{\Large\bfseries\color{csta}d)}\quad\begin{minipage}[c]{5.3cm}\raggedright\Large $50$\end{minipage}\end{tcolorbox}
\end{tcolorbox}
\newpage
\begin{tcolorbox}[enhanced,width=16cm,height=14cm,colback=white,colframe=csta,boxrule=2pt,arc=0mm,
left=6mm,right=6mm,top=0mm,bottom=4mm,sharp corners,
overlay={\fill[csta] (frame.north west) rectangle ([yshift=-1.25cm]frame.north east);
\node[anchor=west,text=white,font=\sffamily\bfseries\LARGE] at ([xshift=5mm,yshift=-0.62cm]frame.north west) {Question 68};
\node[anchor=east,text=white,font=\sffamily\itshape\large] at ([xshift=-5mm,yshift=-0.62cm]frame.north east) {Statistiques et probabilités};}]
\vspace*{1.75cm}
\begin{minipage}[t][6.9cm][c]{\linewidth}\LARGE\centering
\begin{minipage}{0.95\linewidth}\raggedright\linespread{1.25}\selectfont Quelle est la moyenne de la série $8\ ;\ 12\ ;\ 10\ ;\ 15\ ;\ 5$~?\end{minipage}
\end{minipage}

\vfill
\begin{tcolorbox}[enhanced,nobeforeafter,width=7.1cm,height=2.1cm,valign=center,colback=csta!6,colframe=csta,arc=2.5mm,boxrule=0.9pt,left=3mm]{\Large\bfseries\color{csta}a)}\quad\begin{minipage}[c]{5.3cm}\raggedright\Large $12$\end{minipage}\end{tcolorbox}\hfill\begin{tcolorbox}[enhanced,nobeforeafter,width=7.1cm,height=2.1cm,valign=center,colback=csta!6,colframe=csta,arc=2.5mm,boxrule=0.9pt,left=3mm]{\Large\bfseries\color{csta}b)}\quad\begin{minipage}[c]{5.3cm}\raggedright\Large $50$\end{minipage}\end{tcolorbox}\par\vspace{2.5mm}
\begin{tcolorbox}[enhanced,nobeforeafter,width=7.1cm,height=2.1cm,valign=center,colback=csta!6,colframe=csta,arc=2.5mm,boxrule=0.9pt,left=3mm]{\Large\bfseries\color{csta}c)}\quad\begin{minipage}[c]{5.3cm}\raggedright\Large $10$\end{minipage}\end{tcolorbox}\hfill\begin{tcolorbox}[enhanced,nobeforeafter,width=7.1cm,height=2.1cm,valign=center,colback=csta!6,colframe=csta,arc=2.5mm,boxrule=0.9pt,left=3mm]{\Large\bfseries\color{csta}d)}\quad\begin{minipage}[c]{5.3cm}\raggedright\Large $11$\end{minipage}\end{tcolorbox}
\end{tcolorbox}
\newpage
\begin{tcolorbox}[enhanced,width=16cm,height=14cm,colback=white,colframe=csta,boxrule=2pt,arc=0mm,
left=6mm,right=6mm,top=0mm,bottom=4mm,sharp corners,
overlay={\fill[csta] (frame.north west) rectangle ([yshift=-1.25cm]frame.north east);
\node[anchor=west,text=white,font=\sffamily\bfseries\LARGE] at ([xshift=5mm,yshift=-0.62cm]frame.north west) {Question 69};
\node[anchor=east,text=white,font=\sffamily\itshape\large] at ([xshift=-5mm,yshift=-0.62cm]frame.north east) {Statistiques et probabilités};}]
\vspace*{1.75cm}
\begin{minipage}[t][6.9cm][c]{\linewidth}\LARGE\centering
\begin{minipage}{0.95\linewidth}\raggedright\linespread{1.25}\selectfont Quelle est l'étendue de la série $14\ ;\ 9\ ;\ 16\ ;\ 11$~?\end{minipage}
\end{minipage}

\vfill
\begin{tcolorbox}[enhanced,nobeforeafter,width=7.1cm,height=2.1cm,valign=center,colback=csta!6,colframe=csta,arc=2.5mm,boxrule=0.9pt,left=3mm]{\Large\bfseries\color{csta}a)}\quad\begin{minipage}[c]{5.3cm}\raggedright\Large $7$\end{minipage}\end{tcolorbox}\hfill\begin{tcolorbox}[enhanced,nobeforeafter,width=7.1cm,height=2.1cm,valign=center,colback=csta!6,colframe=csta,arc=2.5mm,boxrule=0.9pt,left=3mm]{\Large\bfseries\color{csta}b)}\quad\begin{minipage}[c]{5.3cm}\raggedright\Large $25$\end{minipage}\end{tcolorbox}\par\vspace{2.5mm}
\begin{tcolorbox}[enhanced,nobeforeafter,width=7.1cm,height=2.1cm,valign=center,colback=csta!6,colframe=csta,arc=2.5mm,boxrule=0.9pt,left=3mm]{\Large\bfseries\color{csta}c)}\quad\begin{minipage}[c]{5.3cm}\raggedright\Large $16$\end{minipage}\end{tcolorbox}\hfill\begin{tcolorbox}[enhanced,nobeforeafter,width=7.1cm,height=2.1cm,valign=center,colback=csta!6,colframe=csta,arc=2.5mm,boxrule=0.9pt,left=3mm]{\Large\bfseries\color{csta}d)}\quad\begin{minipage}[c]{5.3cm}\raggedright\Large $8$\end{minipage}\end{tcolorbox}
\end{tcolorbox}
\newpage
\begin{tcolorbox}[enhanced,width=16cm,height=14cm,colback=white,colframe=csta,boxrule=2pt,arc=0mm,
left=6mm,right=6mm,top=0mm,bottom=4mm,sharp corners,
overlay={\fill[csta] (frame.north west) rectangle ([yshift=-1.25cm]frame.north east);
\node[anchor=west,text=white,font=\sffamily\bfseries\LARGE] at ([xshift=5mm,yshift=-0.62cm]frame.north west) {Question 70};
\node[anchor=east,text=white,font=\sffamily\itshape\large] at ([xshift=-5mm,yshift=-0.62cm]frame.north east) {Statistiques et probabilités};}]
\vspace*{1.75cm}
\begin{minipage}[t][6.9cm][c]{\linewidth}\LARGE\centering
\begin{minipage}{0.95\linewidth}\raggedright\linespread{1.25}\selectfont Quelle est l'étendue de la série $8\ ;\ 12\ ;\ 10\ ;\ 15\ ;\ 5$~?\end{minipage}
\end{minipage}

\vfill
\begin{tcolorbox}[enhanced,nobeforeafter,width=7.1cm,height=2.1cm,valign=center,colback=csta!6,colframe=csta,arc=2.5mm,boxrule=0.9pt,left=3mm]{\Large\bfseries\color{csta}a)}\quad\begin{minipage}[c]{5.3cm}\raggedright\Large $15$\end{minipage}\end{tcolorbox}\hfill\begin{tcolorbox}[enhanced,nobeforeafter,width=7.1cm,height=2.1cm,valign=center,colback=csta!6,colframe=csta,arc=2.5mm,boxrule=0.9pt,left=3mm]{\Large\bfseries\color{csta}b)}\quad\begin{minipage}[c]{5.3cm}\raggedright\Large $10$\end{minipage}\end{tcolorbox}\par\vspace{2.5mm}
\begin{tcolorbox}[enhanced,nobeforeafter,width=7.1cm,height=2.1cm,valign=center,colback=csta!6,colframe=csta,arc=2.5mm,boxrule=0.9pt,left=3mm]{\Large\bfseries\color{csta}c)}\quad\begin{minipage}[c]{5.3cm}\raggedright\Large $11$\end{minipage}\end{tcolorbox}\hfill\begin{tcolorbox}[enhanced,nobeforeafter,width=7.1cm,height=2.1cm,valign=center,colback=csta!6,colframe=csta,arc=2.5mm,boxrule=0.9pt,left=3mm]{\Large\bfseries\color{csta}d)}\quad\begin{minipage}[c]{5.3cm}\raggedright\Large $20$\end{minipage}\end{tcolorbox}
\end{tcolorbox}
\newpage
\begin{tcolorbox}[enhanced,width=16cm,height=14cm,colback=white,colframe=csta,boxrule=2pt,arc=0mm,
left=6mm,right=6mm,top=0mm,bottom=4mm,sharp corners,
overlay={\fill[csta] (frame.north west) rectangle ([yshift=-1.25cm]frame.north east);
\node[anchor=west,text=white,font=\sffamily\bfseries\LARGE] at ([xshift=5mm,yshift=-0.62cm]frame.north west) {Question 71};
\node[anchor=east,text=white,font=\sffamily\itshape\large] at ([xshift=-5mm,yshift=-0.62cm]frame.north east) {Statistiques et probabilités};}]
\vspace*{1.75cm}
\begin{minipage}[t][6.9cm][c]{\linewidth}\Large\centering
\begin{minipage}{0.95\linewidth}\raggedright\linespread{1.25}\selectfont Une urne contient $5$ boules dont $3$ rouges. On tire successivement deux boules, sans remise. Quelle est la probabilité de tirer deux boules rouges~?\end{minipage}
\end{minipage}

\vfill
\begin{tcolorbox}[enhanced,nobeforeafter,width=7.1cm,height=2.1cm,valign=center,colback=csta!6,colframe=csta,arc=2.5mm,boxrule=0.9pt,left=3mm]{\Large\bfseries\color{csta}a)}\quad\begin{minipage}[c]{5.3cm}\raggedright\Large $\dfrac{3}{10}$\end{minipage}\end{tcolorbox}\hfill\begin{tcolorbox}[enhanced,nobeforeafter,width=7.1cm,height=2.1cm,valign=center,colback=csta!6,colframe=csta,arc=2.5mm,boxrule=0.9pt,left=3mm]{\Large\bfseries\color{csta}b)}\quad\begin{minipage}[c]{5.3cm}\raggedright\Large $\dfrac{9}{25}$\end{minipage}\end{tcolorbox}\par\vspace{2.5mm}
\begin{tcolorbox}[enhanced,nobeforeafter,width=7.1cm,height=2.1cm,valign=center,colback=csta!6,colframe=csta,arc=2.5mm,boxrule=0.9pt,left=3mm]{\Large\bfseries\color{csta}c)}\quad\begin{minipage}[c]{5.3cm}\raggedright\Large $\dfrac{6}{25}$\end{minipage}\end{tcolorbox}\hfill\begin{tcolorbox}[enhanced,nobeforeafter,width=7.1cm,height=2.1cm,valign=center,colback=csta!6,colframe=csta,arc=2.5mm,boxrule=0.9pt,left=3mm]{\Large\bfseries\color{csta}d)}\quad\begin{minipage}[c]{5.3cm}\raggedright\Large $\dfrac{3}{5}$\end{minipage}\end{tcolorbox}
\end{tcolorbox}
\newpage
\begin{tcolorbox}[enhanced,width=16cm,height=14cm,colback=white,colframe=csta,boxrule=2pt,arc=0mm,
left=6mm,right=6mm,top=0mm,bottom=4mm,sharp corners,
overlay={\fill[csta] (frame.north west) rectangle ([yshift=-1.25cm]frame.north east);
\node[anchor=west,text=white,font=\sffamily\bfseries\LARGE] at ([xshift=5mm,yshift=-0.62cm]frame.north west) {Question 72};
\node[anchor=east,text=white,font=\sffamily\itshape\large] at ([xshift=-5mm,yshift=-0.62cm]frame.north east) {Statistiques et probabilités};}]
\vspace*{1.75cm}
\begin{minipage}[t][6.9cm][c]{\linewidth}\Large\centering
\begin{minipage}{0.95\linewidth}\raggedright\linespread{1.25}\selectfont Une urne contient $4$ boules dont $2$ rouges. On tire successivement deux boules, sans remise. Quelle est la probabilité de tirer deux boules rouges~?\end{minipage}
\end{minipage}

\vfill
\begin{tcolorbox}[enhanced,nobeforeafter,width=7.1cm,height=2.1cm,valign=center,colback=csta!6,colframe=csta,arc=2.5mm,boxrule=0.9pt,left=3mm]{\Large\bfseries\color{csta}a)}\quad\begin{minipage}[c]{5.3cm}\raggedright\Large $1$\end{minipage}\end{tcolorbox}\hfill\begin{tcolorbox}[enhanced,nobeforeafter,width=7.1cm,height=2.1cm,valign=center,colback=csta!6,colframe=csta,arc=2.5mm,boxrule=0.9pt,left=3mm]{\Large\bfseries\color{csta}b)}\quad\begin{minipage}[c]{5.3cm}\raggedright\Large $\dfrac{1}{8}$\end{minipage}\end{tcolorbox}\par\vspace{2.5mm}
\begin{tcolorbox}[enhanced,nobeforeafter,width=7.1cm,height=2.1cm,valign=center,colback=csta!6,colframe=csta,arc=2.5mm,boxrule=0.9pt,left=3mm]{\Large\bfseries\color{csta}c)}\quad\begin{minipage}[c]{5.3cm}\raggedright\Large $\dfrac{1}{4}$\end{minipage}\end{tcolorbox}\hfill\begin{tcolorbox}[enhanced,nobeforeafter,width=7.1cm,height=2.1cm,valign=center,colback=csta!6,colframe=csta,arc=2.5mm,boxrule=0.9pt,left=3mm]{\Large\bfseries\color{csta}d)}\quad\begin{minipage}[c]{5.3cm}\raggedright\Large $\dfrac{1}{6}$\end{minipage}\end{tcolorbox}
\end{tcolorbox}
\newpage
\begin{tcolorbox}[enhanced,width=16cm,height=14cm,colback=white,colframe=csta,boxrule=2pt,arc=0mm,
left=6mm,right=6mm,top=0mm,bottom=4mm,sharp corners,
overlay={\fill[csta] (frame.north west) rectangle ([yshift=-1.25cm]frame.north east);
\node[anchor=west,text=white,font=\sffamily\bfseries\LARGE] at ([xshift=5mm,yshift=-0.62cm]frame.north west) {Question 73};
\node[anchor=east,text=white,font=\sffamily\itshape\large] at ([xshift=-5mm,yshift=-0.62cm]frame.north east) {Statistiques et probabilités};}]
\vspace*{1.75cm}
\begin{minipage}[t][6.9cm][c]{\linewidth}\Large\centering
\begin{minipage}{0.95\linewidth}\raggedright\linespread{1.25}\selectfont On lance un dé équilibré à six faces et une pièce équilibrée (ou un deuxième dé). Quelle est la probabilité d'obtenir un 6 avec le dé et PILE avec la pièce~?\end{minipage}
\end{minipage}

\vfill
\begin{tcolorbox}[enhanced,nobeforeafter,width=7.1cm,height=2.1cm,valign=center,colback=csta!6,colframe=csta,arc=2.5mm,boxrule=0.9pt,left=3mm]{\Large\bfseries\color{csta}a)}\quad\begin{minipage}[c]{5.3cm}\raggedright\Large $\dfrac{1}{2}$\end{minipage}\end{tcolorbox}\hfill\begin{tcolorbox}[enhanced,nobeforeafter,width=7.1cm,height=2.1cm,valign=center,colback=csta!6,colframe=csta,arc=2.5mm,boxrule=0.9pt,left=3mm]{\Large\bfseries\color{csta}b)}\quad\begin{minipage}[c]{5.3cm}\raggedright\Large $\dfrac{1}{6}$\end{minipage}\end{tcolorbox}\par\vspace{2.5mm}
\begin{tcolorbox}[enhanced,nobeforeafter,width=7.1cm,height=2.1cm,valign=center,colback=csta!6,colframe=csta,arc=2.5mm,boxrule=0.9pt,left=3mm]{\Large\bfseries\color{csta}c)}\quad\begin{minipage}[c]{5.3cm}\raggedright\Large $\dfrac{1}{12}$\end{minipage}\end{tcolorbox}\hfill\begin{tcolorbox}[enhanced,nobeforeafter,width=7.1cm,height=2.1cm,valign=center,colback=csta!6,colframe=csta,arc=2.5mm,boxrule=0.9pt,left=3mm]{\Large\bfseries\color{csta}d)}\quad\begin{minipage}[c]{5.3cm}\raggedright\Large $\dfrac{2}{3}$\end{minipage}\end{tcolorbox}
\end{tcolorbox}
\newpage
\begin{tcolorbox}[enhanced,width=16cm,height=14cm,colback=white,colframe=csta,boxrule=2pt,arc=0mm,
left=6mm,right=6mm,top=0mm,bottom=4mm,sharp corners,
overlay={\fill[csta] (frame.north west) rectangle ([yshift=-1.25cm]frame.north east);
\node[anchor=west,text=white,font=\sffamily\bfseries\LARGE] at ([xshift=5mm,yshift=-0.62cm]frame.north west) {Question 74};
\node[anchor=east,text=white,font=\sffamily\itshape\large] at ([xshift=-5mm,yshift=-0.62cm]frame.north east) {Statistiques et probabilités};}]
\vspace*{1.75cm}
\begin{minipage}[t][6.9cm][c]{\linewidth}\Large\centering
\begin{minipage}{0.95\linewidth}\raggedright\linespread{1.25}\selectfont On lance un dé équilibré à six faces et une pièce équilibrée (ou un deuxième dé). Quelle est la probabilité d'obtenir un nombre pair avec le dé et PILE avec la pièce~?\end{minipage}
\end{minipage}

\vfill
\begin{tcolorbox}[enhanced,nobeforeafter,width=7.1cm,height=2.1cm,valign=center,colback=csta!6,colframe=csta,arc=2.5mm,boxrule=0.9pt,left=3mm]{\Large\bfseries\color{csta}a)}\quad\begin{minipage}[c]{5.3cm}\raggedright\Large $\dfrac{1}{2}$\end{minipage}\end{tcolorbox}\hfill\begin{tcolorbox}[enhanced,nobeforeafter,width=7.1cm,height=2.1cm,valign=center,colback=csta!6,colframe=csta,arc=2.5mm,boxrule=0.9pt,left=3mm]{\Large\bfseries\color{csta}b)}\quad\begin{minipage}[c]{5.3cm}\raggedright\Large $\dfrac{1}{4}$\end{minipage}\end{tcolorbox}\par\vspace{2.5mm}
\begin{tcolorbox}[enhanced,nobeforeafter,width=7.1cm,height=2.1cm,valign=center,colback=csta!6,colframe=csta,arc=2.5mm,boxrule=0.9pt,left=3mm]{\Large\bfseries\color{csta}c)}\quad\begin{minipage}[c]{5.3cm}\raggedright\Large $\dfrac{1}{8}$\end{minipage}\end{tcolorbox}\hfill\begin{tcolorbox}[enhanced,nobeforeafter,width=7.1cm,height=2.1cm,valign=center,colback=csta!6,colframe=csta,arc=2.5mm,boxrule=0.9pt,left=3mm]{\Large\bfseries\color{csta}d)}\quad\begin{minipage}[c]{5.3cm}\raggedright\Large $1$\end{minipage}\end{tcolorbox}
\end{tcolorbox}
\newpage
\begin{tcolorbox}[enhanced,width=16cm,height=14cm,colback=white,colframe=csta,boxrule=2pt,arc=0mm,
left=6mm,right=6mm,top=0mm,bottom=4mm,sharp corners,
overlay={\fill[csta] (frame.north west) rectangle ([yshift=-1.25cm]frame.north east);
\node[anchor=west,text=white,font=\sffamily\bfseries\LARGE] at ([xshift=5mm,yshift=-0.62cm]frame.north west) {Question 75};
\node[anchor=east,text=white,font=\sffamily\itshape\large] at ([xshift=-5mm,yshift=-0.62cm]frame.north east) {Statistiques et probabilités};}]
\vspace*{1.75cm}
\begin{minipage}[t][6.9cm][c]{\linewidth}\LARGE\centering
\begin{minipage}{0.95\linewidth}\raggedright\linespread{1.25}\selectfont Dans une classe de $25$ élèves, $8$ pratiquent un sport en club. Quelle est la fréquence, en pourcentage, de ces élèves~?\end{minipage}
\end{minipage}

\vfill
\begin{tcolorbox}[enhanced,nobeforeafter,width=7.1cm,height=2.1cm,valign=center,colback=csta!6,colframe=csta,arc=2.5mm,boxrule=0.9pt,left=3mm]{\Large\bfseries\color{csta}a)}\quad\begin{minipage}[c]{5.3cm}\raggedright\Large $32\,\%$\end{minipage}\end{tcolorbox}\hfill\begin{tcolorbox}[enhanced,nobeforeafter,width=7.1cm,height=2.1cm,valign=center,colback=csta!6,colframe=csta,arc=2.5mm,boxrule=0.9pt,left=3mm]{\Large\bfseries\color{csta}b)}\quad\begin{minipage}[c]{5.3cm}\raggedright\Large $8\,\%$\end{minipage}\end{tcolorbox}\par\vspace{2.5mm}
\begin{tcolorbox}[enhanced,nobeforeafter,width=7.1cm,height=2.1cm,valign=center,colback=csta!6,colframe=csta,arc=2.5mm,boxrule=0.9pt,left=3mm]{\Large\bfseries\color{csta}c)}\quad\begin{minipage}[c]{5.3cm}\raggedright\Large $68\,\%$\end{minipage}\end{tcolorbox}\hfill\begin{tcolorbox}[enhanced,nobeforeafter,width=7.1cm,height=2.1cm,valign=center,colback=csta!6,colframe=csta,arc=2.5mm,boxrule=0.9pt,left=3mm]{\Large\bfseries\color{csta}d)}\quad\begin{minipage}[c]{5.3cm}\raggedright\Large $3{,}2\,\%$\end{minipage}\end{tcolorbox}
\end{tcolorbox}
\newpage
\begin{tcolorbox}[enhanced,width=16cm,height=14cm,colback=white,colframe=csta,boxrule=2pt,arc=0mm,
left=6mm,right=6mm,top=0mm,bottom=4mm,sharp corners,
overlay={\fill[csta] (frame.north west) rectangle ([yshift=-1.25cm]frame.north east);
\node[anchor=west,text=white,font=\sffamily\bfseries\LARGE] at ([xshift=5mm,yshift=-0.62cm]frame.north west) {Question 76};
\node[anchor=east,text=white,font=\sffamily\itshape\large] at ([xshift=-5mm,yshift=-0.62cm]frame.north east) {Statistiques et probabilités};}]
\vspace*{1.75cm}
\begin{minipage}[t][6.9cm][c]{\linewidth}\LARGE\centering
\begin{minipage}{0.95\linewidth}\raggedright\linespread{1.25}\selectfont Dans une classe de $40$ élèves, $12$ pratiquent un sport en club. Quelle est la fréquence, en pourcentage, de ces élèves~?\end{minipage}
\end{minipage}

\vfill
\begin{tcolorbox}[enhanced,nobeforeafter,width=7.1cm,height=2.1cm,valign=center,colback=csta!6,colframe=csta,arc=2.5mm,boxrule=0.9pt,left=3mm]{\Large\bfseries\color{csta}a)}\quad\begin{minipage}[c]{5.3cm}\raggedright\Large $70\,\%$\end{minipage}\end{tcolorbox}\hfill\begin{tcolorbox}[enhanced,nobeforeafter,width=7.1cm,height=2.1cm,valign=center,colback=csta!6,colframe=csta,arc=2.5mm,boxrule=0.9pt,left=3mm]{\Large\bfseries\color{csta}b)}\quad\begin{minipage}[c]{5.3cm}\raggedright\Large $30\,\%$\end{minipage}\end{tcolorbox}\par\vspace{2.5mm}
\begin{tcolorbox}[enhanced,nobeforeafter,width=7.1cm,height=2.1cm,valign=center,colback=csta!6,colframe=csta,arc=2.5mm,boxrule=0.9pt,left=3mm]{\Large\bfseries\color{csta}c)}\quad\begin{minipage}[c]{5.3cm}\raggedright\Large $12\,\%$\end{minipage}\end{tcolorbox}\hfill\begin{tcolorbox}[enhanced,nobeforeafter,width=7.1cm,height=2.1cm,valign=center,colback=csta!6,colframe=csta,arc=2.5mm,boxrule=0.9pt,left=3mm]{\Large\bfseries\color{csta}d)}\quad\begin{minipage}[c]{5.3cm}\raggedright\Large $3\,\%$\end{minipage}\end{tcolorbox}
\end{tcolorbox}
\newpage
\begin{tcolorbox}[enhanced,width=16cm,height=14cm,colback=white,colframe=csta,boxrule=2pt,arc=0mm,
left=6mm,right=6mm,top=0mm,bottom=4mm,sharp corners,
overlay={\fill[csta] (frame.north west) rectangle ([yshift=-1.25cm]frame.north east);
\node[anchor=west,text=white,font=\sffamily\bfseries\LARGE] at ([xshift=5mm,yshift=-0.62cm]frame.north west) {Question 77};
\node[anchor=east,text=white,font=\sffamily\itshape\large] at ([xshift=-5mm,yshift=-0.62cm]frame.north east) {Statistiques et probabilités};}]
\vspace*{1.75cm}
\begin{minipage}[t][6.9cm][c]{\linewidth}\Large\centering
\begin{minipage}{0.95\linewidth}\raggedright\linespread{1.25}\selectfont Un sac contient $4$ boules rouges, $6$ boules bleues et $2$ boules vertes. On tire une boule au hasard. Quelle est la probabilité de ne pas tirer une boule rouge~?\end{minipage}
\end{minipage}

\vfill
\begin{tcolorbox}[enhanced,nobeforeafter,width=7.1cm,height=2.1cm,valign=center,colback=csta!6,colframe=csta,arc=2.5mm,boxrule=0.9pt,left=3mm]{\Large\bfseries\color{csta}a)}\quad\begin{minipage}[c]{5.3cm}\raggedright\Large $2$\end{minipage}\end{tcolorbox}\hfill\begin{tcolorbox}[enhanced,nobeforeafter,width=7.1cm,height=2.1cm,valign=center,colback=csta!6,colframe=csta,arc=2.5mm,boxrule=0.9pt,left=3mm]{\Large\bfseries\color{csta}b)}\quad\begin{minipage}[c]{5.3cm}\raggedright\Large $\dfrac{1}{3}$\end{minipage}\end{tcolorbox}\par\vspace{2.5mm}
\begin{tcolorbox}[enhanced,nobeforeafter,width=7.1cm,height=2.1cm,valign=center,colback=csta!6,colframe=csta,arc=2.5mm,boxrule=0.9pt,left=3mm]{\Large\bfseries\color{csta}c)}\quad\begin{minipage}[c]{5.3cm}\raggedright\Large $\dfrac{2}{3}$\end{minipage}\end{tcolorbox}\hfill\begin{tcolorbox}[enhanced,nobeforeafter,width=7.1cm,height=2.1cm,valign=center,colback=csta!6,colframe=csta,arc=2.5mm,boxrule=0.9pt,left=3mm]{\Large\bfseries\color{csta}d)}\quad\begin{minipage}[c]{5.3cm}\raggedright\Large $\dfrac{1}{2}$\end{minipage}\end{tcolorbox}
\end{tcolorbox}
\newpage
\begin{tcolorbox}[enhanced,width=16cm,height=14cm,colback=white,colframe=csta,boxrule=2pt,arc=0mm,
left=6mm,right=6mm,top=0mm,bottom=4mm,sharp corners,
overlay={\fill[csta] (frame.north west) rectangle ([yshift=-1.25cm]frame.north east);
\node[anchor=west,text=white,font=\sffamily\bfseries\LARGE] at ([xshift=5mm,yshift=-0.62cm]frame.north west) {Question 78};
\node[anchor=east,text=white,font=\sffamily\itshape\large] at ([xshift=-5mm,yshift=-0.62cm]frame.north east) {Statistiques et probabilités};}]
\vspace*{1.75cm}
\begin{minipage}[t][6.9cm][c]{\linewidth}\Large\centering
\begin{minipage}{0.95\linewidth}\raggedright\linespread{1.25}\selectfont Un sac contient $5$ boules rouges, $3$ boules bleues et $4$ boules vertes. On tire une boule au hasard. Quelle est la probabilité de ne pas tirer une boule rouge~?\end{minipage}
\end{minipage}

\vfill
\begin{tcolorbox}[enhanced,nobeforeafter,width=7.1cm,height=2.1cm,valign=center,colback=csta!6,colframe=csta,arc=2.5mm,boxrule=0.9pt,left=3mm]{\Large\bfseries\color{csta}a)}\quad\begin{minipage}[c]{5.3cm}\raggedright\Large $\dfrac{5}{12}$\end{minipage}\end{tcolorbox}\hfill\begin{tcolorbox}[enhanced,nobeforeafter,width=7.1cm,height=2.1cm,valign=center,colback=csta!6,colframe=csta,arc=2.5mm,boxrule=0.9pt,left=3mm]{\Large\bfseries\color{csta}b)}\quad\begin{minipage}[c]{5.3cm}\raggedright\Large $\dfrac{1}{4}$\end{minipage}\end{tcolorbox}\par\vspace{2.5mm}
\begin{tcolorbox}[enhanced,nobeforeafter,width=7.1cm,height=2.1cm,valign=center,colback=csta!6,colframe=csta,arc=2.5mm,boxrule=0.9pt,left=3mm]{\Large\bfseries\color{csta}c)}\quad\begin{minipage}[c]{5.3cm}\raggedright\Large $\dfrac{7}{12}$\end{minipage}\end{tcolorbox}\hfill\begin{tcolorbox}[enhanced,nobeforeafter,width=7.1cm,height=2.1cm,valign=center,colback=csta!6,colframe=csta,arc=2.5mm,boxrule=0.9pt,left=3mm]{\Large\bfseries\color{csta}d)}\quad\begin{minipage}[c]{5.3cm}\raggedright\Large $\dfrac{7}{5}$\end{minipage}\end{tcolorbox}
\end{tcolorbox}
\newpage
\begin{tcolorbox}[enhanced,width=16cm,height=14cm,colback=white,colframe=csta,boxrule=2pt,arc=0mm,
left=6mm,right=6mm,top=0mm,bottom=4mm,sharp corners,
overlay={\fill[csta] (frame.north west) rectangle ([yshift=-1.25cm]frame.north east);
\node[anchor=west,text=white,font=\sffamily\bfseries\LARGE] at ([xshift=5mm,yshift=-0.62cm]frame.north west) {Question 79};
\node[anchor=east,text=white,font=\sffamily\itshape\large] at ([xshift=-5mm,yshift=-0.62cm]frame.north east) {Statistiques et probabilités};}]
\vspace*{1.75cm}
\begin{minipage}[t][6.9cm][c]{\linewidth}\LARGE\centering
\begin{minipage}{0.95\linewidth}\raggedright\linespread{1.25}\selectfont La probabilité d'un événement $A$ est $\dfrac{3}{7}$. Quelle est la probabilité de l'événement contraire de $A$~?\end{minipage}
\end{minipage}

\vfill
\begin{tcolorbox}[enhanced,nobeforeafter,width=7.1cm,height=2.1cm,valign=center,colback=csta!6,colframe=csta,arc=2.5mm,boxrule=0.9pt,left=3mm]{\Large\bfseries\color{csta}a)}\quad\begin{minipage}[c]{5.3cm}\raggedright\Large $\dfrac{3}{7}$\end{minipage}\end{tcolorbox}\hfill\begin{tcolorbox}[enhanced,nobeforeafter,width=7.1cm,height=2.1cm,valign=center,colback=csta!6,colframe=csta,arc=2.5mm,boxrule=0.9pt,left=3mm]{\Large\bfseries\color{csta}b)}\quad\begin{minipage}[c]{5.3cm}\raggedright\Large $\dfrac{4}{7}$\end{minipage}\end{tcolorbox}\par\vspace{2.5mm}
\begin{tcolorbox}[enhanced,nobeforeafter,width=7.1cm,height=2.1cm,valign=center,colback=csta!6,colframe=csta,arc=2.5mm,boxrule=0.9pt,left=3mm]{\Large\bfseries\color{csta}c)}\quad\begin{minipage}[c]{5.3cm}\raggedright\Large $\dfrac{1}{14}$\end{minipage}\end{tcolorbox}\hfill\begin{tcolorbox}[enhanced,nobeforeafter,width=7.1cm,height=2.1cm,valign=center,colback=csta!6,colframe=csta,arc=2.5mm,boxrule=0.9pt,left=3mm]{\Large\bfseries\color{csta}d)}\quad\begin{minipage}[c]{5.3cm}\raggedright\Large $\dfrac{2}{7}$\end{minipage}\end{tcolorbox}
\end{tcolorbox}
\newpage
\begin{tcolorbox}[enhanced,width=16cm,height=14cm,colback=white,colframe=csta,boxrule=2pt,arc=0mm,
left=6mm,right=6mm,top=0mm,bottom=4mm,sharp corners,
overlay={\fill[csta] (frame.north west) rectangle ([yshift=-1.25cm]frame.north east);
\node[anchor=west,text=white,font=\sffamily\bfseries\LARGE] at ([xshift=5mm,yshift=-0.62cm]frame.north west) {Question 80};
\node[anchor=east,text=white,font=\sffamily\itshape\large] at ([xshift=-5mm,yshift=-0.62cm]frame.north east) {Statistiques et probabilités};}]
\vspace*{1.75cm}
\begin{minipage}[t][6.9cm][c]{\linewidth}\LARGE\centering
\begin{minipage}{0.95\linewidth}\raggedright\linespread{1.25}\selectfont La probabilité d'un événement $A$ est $0{,}15$. Quelle est la probabilité de l'événement contraire de $A$~?\end{minipage}
\end{minipage}

\vfill
\begin{tcolorbox}[enhanced,nobeforeafter,width=7.1cm,height=2.1cm,valign=center,colback=csta!6,colframe=csta,arc=2.5mm,boxrule=0.9pt,left=3mm]{\Large\bfseries\color{csta}a)}\quad\begin{minipage}[c]{5.3cm}\raggedright\Large $0{,}35$\end{minipage}\end{tcolorbox}\hfill\begin{tcolorbox}[enhanced,nobeforeafter,width=7.1cm,height=2.1cm,valign=center,colback=csta!6,colframe=csta,arc=2.5mm,boxrule=0.9pt,left=3mm]{\Large\bfseries\color{csta}b)}\quad\begin{minipage}[c]{5.3cm}\raggedright\Large $0{,}425$\end{minipage}\end{tcolorbox}\par\vspace{2.5mm}
\begin{tcolorbox}[enhanced,nobeforeafter,width=7.1cm,height=2.1cm,valign=center,colback=csta!6,colframe=csta,arc=2.5mm,boxrule=0.9pt,left=3mm]{\Large\bfseries\color{csta}c)}\quad\begin{minipage}[c]{5.3cm}\raggedright\Large $0{,}85$\end{minipage}\end{tcolorbox}\hfill\begin{tcolorbox}[enhanced,nobeforeafter,width=7.1cm,height=2.1cm,valign=center,colback=csta!6,colframe=csta,arc=2.5mm,boxrule=0.9pt,left=3mm]{\Large\bfseries\color{csta}d)}\quad\begin{minipage}[c]{5.3cm}\raggedright\Large $0{,}15$\end{minipage}\end{tcolorbox}
\end{tcolorbox}
\newpage
\begin{tcolorbox}[enhanced,width=16cm,height=14cm,colback=white,colframe=cgra,boxrule=2pt,arc=0mm,
left=6mm,right=6mm,top=0mm,bottom=4mm,sharp corners,
overlay={\fill[cgra] (frame.north west) rectangle ([yshift=-1.25cm]frame.north east);
\node[anchor=west,text=white,font=\sffamily\bfseries\LARGE] at ([xshift=5mm,yshift=-0.62cm]frame.north west) {Question 81};
\node[anchor=east,text=white,font=\sffamily\itshape\large] at ([xshift=-5mm,yshift=-0.62cm]frame.north east) {Grandeurs et proportionnalité};}]
\vspace*{1.75cm}
\begin{minipage}[t][6.9cm][c]{\linewidth}\LARGE\centering
\begin{minipage}{0.95\linewidth}\raggedright\linespread{1.25}\selectfont Une vitesse de $10$ m/s correspond à~:\end{minipage}
\end{minipage}

\vfill
\begin{tcolorbox}[enhanced,nobeforeafter,width=7.1cm,height=2.1cm,valign=center,colback=cgra!6,colframe=cgra,arc=2.5mm,boxrule=0.9pt,left=3mm]{\Large\bfseries\color{cgra}a)}\quad\begin{minipage}[c]{5.3cm}\raggedright\Large $10\,000~\text{km/h}$\end{minipage}\end{tcolorbox}\hfill\begin{tcolorbox}[enhanced,nobeforeafter,width=7.1cm,height=2.1cm,valign=center,colback=cgra!6,colframe=cgra,arc=2.5mm,boxrule=0.9pt,left=3mm]{\Large\bfseries\color{cgra}b)}\quad\begin{minipage}[c]{5.3cm}\raggedright\Large $100~\text{km/h}$\end{minipage}\end{tcolorbox}\par\vspace{2.5mm}
\begin{tcolorbox}[enhanced,nobeforeafter,width=7.1cm,height=2.1cm,valign=center,colback=cgra!6,colframe=cgra,arc=2.5mm,boxrule=0.9pt,left=3mm]{\Large\bfseries\color{cgra}c)}\quad\begin{minipage}[c]{5.3cm}\raggedright\Large $600~\text{km/h}$\end{minipage}\end{tcolorbox}\hfill\begin{tcolorbox}[enhanced,nobeforeafter,width=7.1cm,height=2.1cm,valign=center,colback=cgra!6,colframe=cgra,arc=2.5mm,boxrule=0.9pt,left=3mm]{\Large\bfseries\color{cgra}d)}\quad\begin{minipage}[c]{5.3cm}\raggedright\Large $36~\text{km/h}$\end{minipage}\end{tcolorbox}
\end{tcolorbox}
\newpage
\begin{tcolorbox}[enhanced,width=16cm,height=14cm,colback=white,colframe=cgra,boxrule=2pt,arc=0mm,
left=6mm,right=6mm,top=0mm,bottom=4mm,sharp corners,
overlay={\fill[cgra] (frame.north west) rectangle ([yshift=-1.25cm]frame.north east);
\node[anchor=west,text=white,font=\sffamily\bfseries\LARGE] at ([xshift=5mm,yshift=-0.62cm]frame.north west) {Question 82};
\node[anchor=east,text=white,font=\sffamily\itshape\large] at ([xshift=-5mm,yshift=-0.62cm]frame.north east) {Grandeurs et proportionnalité};}]
\vspace*{1.75cm}
\begin{minipage}[t][6.9cm][c]{\linewidth}\LARGE\centering
\begin{minipage}{0.95\linewidth}\raggedright\linespread{1.25}\selectfont Une vitesse de $25$ m/s correspond à~:\end{minipage}
\end{minipage}

\vfill
\begin{tcolorbox}[enhanced,nobeforeafter,width=7.1cm,height=2.1cm,valign=center,colback=cgra!6,colframe=cgra,arc=2.5mm,boxrule=0.9pt,left=3mm]{\Large\bfseries\color{cgra}a)}\quad\begin{minipage}[c]{5.3cm}\raggedright\Large $90~\text{km/h}$\end{minipage}\end{tcolorbox}\hfill\begin{tcolorbox}[enhanced,nobeforeafter,width=7.1cm,height=2.1cm,valign=center,colback=cgra!6,colframe=cgra,arc=2.5mm,boxrule=0.9pt,left=3mm]{\Large\bfseries\color{cgra}b)}\quad\begin{minipage}[c]{5.3cm}\raggedright\Large $25\,000~\text{km/h}$\end{minipage}\end{tcolorbox}\par\vspace{2.5mm}
\begin{tcolorbox}[enhanced,nobeforeafter,width=7.1cm,height=2.1cm,valign=center,colback=cgra!6,colframe=cgra,arc=2.5mm,boxrule=0.9pt,left=3mm]{\Large\bfseries\color{cgra}c)}\quad\begin{minipage}[c]{5.3cm}\raggedright\Large $250~\text{km/h}$\end{minipage}\end{tcolorbox}\hfill\begin{tcolorbox}[enhanced,nobeforeafter,width=7.1cm,height=2.1cm,valign=center,colback=cgra!6,colframe=cgra,arc=2.5mm,boxrule=0.9pt,left=3mm]{\Large\bfseries\color{cgra}d)}\quad\begin{minipage}[c]{5.3cm}\raggedright\Large $1\,500~\text{km/h}$\end{minipage}\end{tcolorbox}
\end{tcolorbox}
\newpage
\begin{tcolorbox}[enhanced,width=16cm,height=14cm,colback=white,colframe=cgra,boxrule=2pt,arc=0mm,
left=6mm,right=6mm,top=0mm,bottom=4mm,sharp corners,
overlay={\fill[cgra] (frame.north west) rectangle ([yshift=-1.25cm]frame.north east);
\node[anchor=west,text=white,font=\sffamily\bfseries\LARGE] at ([xshift=5mm,yshift=-0.62cm]frame.north west) {Question 83};
\node[anchor=east,text=white,font=\sffamily\itshape\large] at ([xshift=-5mm,yshift=-0.62cm]frame.north east) {Grandeurs et proportionnalité};}]
\vspace*{1.75cm}
\begin{minipage}[t][6.9cm][c]{\linewidth}\LARGE\centering
\begin{minipage}{0.95\linewidth}\raggedright\linespread{1.25}\selectfont Un véhicule parcourt $180$ km en 2 h. Quelle est sa vitesse moyenne~?\end{minipage}
\end{minipage}

\vfill
\begin{tcolorbox}[enhanced,nobeforeafter,width=7.1cm,height=2.1cm,valign=center,colback=cgra!6,colframe=cgra,arc=2.5mm,boxrule=0.9pt,left=3mm]{\Large\bfseries\color{cgra}a)}\quad\begin{minipage}[c]{5.3cm}\raggedright\Large $90~\text{km/h}$\end{minipage}\end{tcolorbox}\hfill\begin{tcolorbox}[enhanced,nobeforeafter,width=7.1cm,height=2.1cm,valign=center,colback=cgra!6,colframe=cgra,arc=2.5mm,boxrule=0.9pt,left=3mm]{\Large\bfseries\color{cgra}b)}\quad\begin{minipage}[c]{5.3cm}\raggedright\Large $360~\text{km/h}$\end{minipage}\end{tcolorbox}\par\vspace{2.5mm}
\begin{tcolorbox}[enhanced,nobeforeafter,width=7.1cm,height=2.1cm,valign=center,colback=cgra!6,colframe=cgra,arc=2.5mm,boxrule=0.9pt,left=3mm]{\Large\bfseries\color{cgra}c)}\quad\begin{minipage}[c]{5.3cm}\raggedright\Large $100~\text{km/h}$\end{minipage}\end{tcolorbox}\hfill\begin{tcolorbox}[enhanced,nobeforeafter,width=7.1cm,height=2.1cm,valign=center,colback=cgra!6,colframe=cgra,arc=2.5mm,boxrule=0.9pt,left=3mm]{\Large\bfseries\color{cgra}d)}\quad\begin{minipage}[c]{5.3cm}\raggedright\Large $45~\text{km/h}$\end{minipage}\end{tcolorbox}
\end{tcolorbox}
\newpage
\begin{tcolorbox}[enhanced,width=16cm,height=14cm,colback=white,colframe=cgra,boxrule=2pt,arc=0mm,
left=6mm,right=6mm,top=0mm,bottom=4mm,sharp corners,
overlay={\fill[cgra] (frame.north west) rectangle ([yshift=-1.25cm]frame.north east);
\node[anchor=west,text=white,font=\sffamily\bfseries\LARGE] at ([xshift=5mm,yshift=-0.62cm]frame.north west) {Question 84};
\node[anchor=east,text=white,font=\sffamily\itshape\large] at ([xshift=-5mm,yshift=-0.62cm]frame.north east) {Grandeurs et proportionnalité};}]
\vspace*{1.75cm}
\begin{minipage}[t][6.9cm][c]{\linewidth}\LARGE\centering
\begin{minipage}{0.95\linewidth}\raggedright\linespread{1.25}\selectfont Un véhicule parcourt $135$ km en 1 h 30 min. Quelle est sa vitesse moyenne~?\end{minipage}
\end{minipage}

\vfill
\begin{tcolorbox}[enhanced,nobeforeafter,width=7.1cm,height=2.1cm,valign=center,colback=cgra!6,colframe=cgra,arc=2.5mm,boxrule=0.9pt,left=3mm]{\Large\bfseries\color{cgra}a)}\quad\begin{minipage}[c]{5.3cm}\raggedright\Large $100~\text{km/h}$\end{minipage}\end{tcolorbox}\hfill\begin{tcolorbox}[enhanced,nobeforeafter,width=7.1cm,height=2.1cm,valign=center,colback=cgra!6,colframe=cgra,arc=2.5mm,boxrule=0.9pt,left=3mm]{\Large\bfseries\color{cgra}b)}\quad\begin{minipage}[c]{5.3cm}\raggedright\Large $202{,}5~\text{km/h}$\end{minipage}\end{tcolorbox}\par\vspace{2.5mm}
\begin{tcolorbox}[enhanced,nobeforeafter,width=7.1cm,height=2.1cm,valign=center,colback=cgra!6,colframe=cgra,arc=2.5mm,boxrule=0.9pt,left=3mm]{\Large\bfseries\color{cgra}c)}\quad\begin{minipage}[c]{5.3cm}\raggedright\Large $90~\text{km/h}$\end{minipage}\end{tcolorbox}\hfill\begin{tcolorbox}[enhanced,nobeforeafter,width=7.1cm,height=2.1cm,valign=center,colback=cgra!6,colframe=cgra,arc=2.5mm,boxrule=0.9pt,left=3mm]{\Large\bfseries\color{cgra}d)}\quad\begin{minipage}[c]{5.3cm}\raggedright\Large $45~\text{km/h}$\end{minipage}\end{tcolorbox}
\end{tcolorbox}
\newpage
\begin{tcolorbox}[enhanced,width=16cm,height=14cm,colback=white,colframe=cgra,boxrule=2pt,arc=0mm,
left=6mm,right=6mm,top=0mm,bottom=4mm,sharp corners,
overlay={\fill[cgra] (frame.north west) rectangle ([yshift=-1.25cm]frame.north east);
\node[anchor=west,text=white,font=\sffamily\bfseries\LARGE] at ([xshift=5mm,yshift=-0.62cm]frame.north west) {Question 85};
\node[anchor=east,text=white,font=\sffamily\itshape\large] at ([xshift=-5mm,yshift=-0.62cm]frame.north east) {Grandeurs et proportionnalité};}]
\vspace*{1.75cm}
\begin{minipage}[t][6.9cm][c]{\linewidth}\LARGE\centering
\begin{minipage}{0.95\linewidth}\raggedright\linespread{1.25}\selectfont Un article coûte $60$ F. Son prix augmente de $15\,\%$. Quel est son nouveau prix~?\end{minipage}
\end{minipage}

\vfill
\begin{tcolorbox}[enhanced,nobeforeafter,width=7.1cm,height=2.1cm,valign=center,colback=cgra!6,colframe=cgra,arc=2.5mm,boxrule=0.9pt,left=3mm]{\Large\bfseries\color{cgra}a)}\quad\begin{minipage}[c]{5.3cm}\raggedright\Large $75~\text{F}$\end{minipage}\end{tcolorbox}\hfill\begin{tcolorbox}[enhanced,nobeforeafter,width=7.1cm,height=2.1cm,valign=center,colback=cgra!6,colframe=cgra,arc=2.5mm,boxrule=0.9pt,left=3mm]{\Large\bfseries\color{cgra}b)}\quad\begin{minipage}[c]{5.3cm}\raggedright\Large $9~\text{F}$\end{minipage}\end{tcolorbox}\par\vspace{2.5mm}
\begin{tcolorbox}[enhanced,nobeforeafter,width=7.1cm,height=2.1cm,valign=center,colback=cgra!6,colframe=cgra,arc=2.5mm,boxrule=0.9pt,left=3mm]{\Large\bfseries\color{cgra}c)}\quad\begin{minipage}[c]{5.3cm}\raggedright\Large $51~\text{F}$\end{minipage}\end{tcolorbox}\hfill\begin{tcolorbox}[enhanced,nobeforeafter,width=7.1cm,height=2.1cm,valign=center,colback=cgra!6,colframe=cgra,arc=2.5mm,boxrule=0.9pt,left=3mm]{\Large\bfseries\color{cgra}d)}\quad\begin{minipage}[c]{5.3cm}\raggedright\Large $69~\text{F}$\end{minipage}\end{tcolorbox}
\end{tcolorbox}
\newpage
\begin{tcolorbox}[enhanced,width=16cm,height=14cm,colback=white,colframe=cgra,boxrule=2pt,arc=0mm,
left=6mm,right=6mm,top=0mm,bottom=4mm,sharp corners,
overlay={\fill[cgra] (frame.north west) rectangle ([yshift=-1.25cm]frame.north east);
\node[anchor=west,text=white,font=\sffamily\bfseries\LARGE] at ([xshift=5mm,yshift=-0.62cm]frame.north west) {Question 86};
\node[anchor=east,text=white,font=\sffamily\itshape\large] at ([xshift=-5mm,yshift=-0.62cm]frame.north east) {Grandeurs et proportionnalité};}]
\vspace*{1.75cm}
\begin{minipage}[t][6.9cm][c]{\linewidth}\LARGE\centering
\begin{minipage}{0.95\linewidth}\raggedright\linespread{1.25}\selectfont Un article coûte $240$ F. Son prix baisse de $25\,\%$. Quel est son nouveau prix~?\end{minipage}
\end{minipage}

\vfill
\begin{tcolorbox}[enhanced,nobeforeafter,width=7.1cm,height=2.1cm,valign=center,colback=cgra!6,colframe=cgra,arc=2.5mm,boxrule=0.9pt,left=3mm]{\Large\bfseries\color{cgra}a)}\quad\begin{minipage}[c]{5.3cm}\raggedright\Large $60~\text{F}$\end{minipage}\end{tcolorbox}\hfill\begin{tcolorbox}[enhanced,nobeforeafter,width=7.1cm,height=2.1cm,valign=center,colback=cgra!6,colframe=cgra,arc=2.5mm,boxrule=0.9pt,left=3mm]{\Large\bfseries\color{cgra}b)}\quad\begin{minipage}[c]{5.3cm}\raggedright\Large $180~\text{F}$\end{minipage}\end{tcolorbox}\par\vspace{2.5mm}
\begin{tcolorbox}[enhanced,nobeforeafter,width=7.1cm,height=2.1cm,valign=center,colback=cgra!6,colframe=cgra,arc=2.5mm,boxrule=0.9pt,left=3mm]{\Large\bfseries\color{cgra}c)}\quad\begin{minipage}[c]{5.3cm}\raggedright\Large $215~\text{F}$\end{minipage}\end{tcolorbox}\hfill\begin{tcolorbox}[enhanced,nobeforeafter,width=7.1cm,height=2.1cm,valign=center,colback=cgra!6,colframe=cgra,arc=2.5mm,boxrule=0.9pt,left=3mm]{\Large\bfseries\color{cgra}d)}\quad\begin{minipage}[c]{5.3cm}\raggedright\Large $300~\text{F}$\end{minipage}\end{tcolorbox}
\end{tcolorbox}
\newpage
\begin{tcolorbox}[enhanced,width=16cm,height=14cm,colback=white,colframe=cgra,boxrule=2pt,arc=0mm,
left=6mm,right=6mm,top=0mm,bottom=4mm,sharp corners,
overlay={\fill[cgra] (frame.north west) rectangle ([yshift=-1.25cm]frame.north east);
\node[anchor=west,text=white,font=\sffamily\bfseries\LARGE] at ([xshift=5mm,yshift=-0.62cm]frame.north west) {Question 87};
\node[anchor=east,text=white,font=\sffamily\itshape\large] at ([xshift=-5mm,yshift=-0.62cm]frame.north east) {Grandeurs et proportionnalité};}]
\vspace*{1.75cm}
\begin{minipage}[t][6.9cm][c]{\linewidth}\LARGE\centering
\begin{minipage}{0.95\linewidth}\raggedright\linespread{1.25}\selectfont Un prix passe de $80$ F à $100$ F. Quel est le pourcentage d'évolution~?\end{minipage}
\end{minipage}

\vfill
\begin{tcolorbox}[enhanced,nobeforeafter,width=7.1cm,height=2.1cm,valign=center,colback=cgra!6,colframe=cgra,arc=2.5mm,boxrule=0.9pt,left=3mm]{\Large\bfseries\color{cgra}a)}\quad\begin{minipage}[c]{5.3cm}\raggedright\Large $+25\,\%$\end{minipage}\end{tcolorbox}\hfill\begin{tcolorbox}[enhanced,nobeforeafter,width=7.1cm,height=2.1cm,valign=center,colback=cgra!6,colframe=cgra,arc=2.5mm,boxrule=0.9pt,left=3mm]{\Large\bfseries\color{cgra}b)}\quad\begin{minipage}[c]{5.3cm}\raggedright\Large $+12{,}5\,\%$\end{minipage}\end{tcolorbox}\par\vspace{2.5mm}
\begin{tcolorbox}[enhanced,nobeforeafter,width=7.1cm,height=2.1cm,valign=center,colback=cgra!6,colframe=cgra,arc=2.5mm,boxrule=0.9pt,left=3mm]{\Large\bfseries\color{cgra}c)}\quad\begin{minipage}[c]{5.3cm}\raggedright\Large $-25\,\%$\end{minipage}\end{tcolorbox}\hfill\begin{tcolorbox}[enhanced,nobeforeafter,width=7.1cm,height=2.1cm,valign=center,colback=cgra!6,colframe=cgra,arc=2.5mm,boxrule=0.9pt,left=3mm]{\Large\bfseries\color{cgra}d)}\quad\begin{minipage}[c]{5.3cm}\raggedright\Large $+20\,\%$\end{minipage}\end{tcolorbox}
\end{tcolorbox}
\newpage
\begin{tcolorbox}[enhanced,width=16cm,height=14cm,colback=white,colframe=cgra,boxrule=2pt,arc=0mm,
left=6mm,right=6mm,top=0mm,bottom=4mm,sharp corners,
overlay={\fill[cgra] (frame.north west) rectangle ([yshift=-1.25cm]frame.north east);
\node[anchor=west,text=white,font=\sffamily\bfseries\LARGE] at ([xshift=5mm,yshift=-0.62cm]frame.north west) {Question 88};
\node[anchor=east,text=white,font=\sffamily\itshape\large] at ([xshift=-5mm,yshift=-0.62cm]frame.north east) {Grandeurs et proportionnalité};}]
\vspace*{1.75cm}
\begin{minipage}[t][6.9cm][c]{\linewidth}\LARGE\centering
\begin{minipage}{0.95\linewidth}\raggedright\linespread{1.25}\selectfont Un prix passe de $250$ F à $200$ F. Quel est le pourcentage d'évolution~?\end{minipage}
\end{minipage}

\vfill
\begin{tcolorbox}[enhanced,nobeforeafter,width=7.1cm,height=2.1cm,valign=center,colback=cgra!6,colframe=cgra,arc=2.5mm,boxrule=0.9pt,left=3mm]{\Large\bfseries\color{cgra}a)}\quad\begin{minipage}[c]{5.3cm}\raggedright\Large $-10\,\%$\end{minipage}\end{tcolorbox}\hfill\begin{tcolorbox}[enhanced,nobeforeafter,width=7.1cm,height=2.1cm,valign=center,colback=cgra!6,colframe=cgra,arc=2.5mm,boxrule=0.9pt,left=3mm]{\Large\bfseries\color{cgra}b)}\quad\begin{minipage}[c]{5.3cm}\raggedright\Large $-50\,\%$\end{minipage}\end{tcolorbox}\par\vspace{2.5mm}
\begin{tcolorbox}[enhanced,nobeforeafter,width=7.1cm,height=2.1cm,valign=center,colback=cgra!6,colframe=cgra,arc=2.5mm,boxrule=0.9pt,left=3mm]{\Large\bfseries\color{cgra}c)}\quad\begin{minipage}[c]{5.3cm}\raggedright\Large $-20\,\%$\end{minipage}\end{tcolorbox}\hfill\begin{tcolorbox}[enhanced,nobeforeafter,width=7.1cm,height=2.1cm,valign=center,colback=cgra!6,colframe=cgra,arc=2.5mm,boxrule=0.9pt,left=3mm]{\Large\bfseries\color{cgra}d)}\quad\begin{minipage}[c]{5.3cm}\raggedright\Large $+20\,\%$\end{minipage}\end{tcolorbox}
\end{tcolorbox}
\newpage
\begin{tcolorbox}[enhanced,width=16cm,height=14cm,colback=white,colframe=cgra,boxrule=2pt,arc=0mm,
left=6mm,right=6mm,top=0mm,bottom=4mm,sharp corners,
overlay={\fill[cgra] (frame.north west) rectangle ([yshift=-1.25cm]frame.north east);
\node[anchor=west,text=white,font=\sffamily\bfseries\LARGE] at ([xshift=5mm,yshift=-0.62cm]frame.north west) {Question 89};
\node[anchor=east,text=white,font=\sffamily\itshape\large] at ([xshift=-5mm,yshift=-0.62cm]frame.north east) {Grandeurs et proportionnalité};}]
\vspace*{1.75cm}
\begin{minipage}[t][6.9cm][c]{\linewidth}\LARGE\centering
\begin{minipage}{0.95\linewidth}\raggedright\linespread{1.25}\selectfont Un robinet remplit une cuve de $120$ L en $5$ min. Quel est son débit~?\end{minipage}
\end{minipage}

\vfill
\begin{tcolorbox}[enhanced,nobeforeafter,width=7.1cm,height=2.1cm,valign=center,colback=cgra!6,colframe=cgra,arc=2.5mm,boxrule=0.9pt,left=3mm]{\Large\bfseries\color{cgra}a)}\quad\begin{minipage}[c]{5.3cm}\raggedright\large $\dfrac{25}{6}~\text{L/min}$\end{minipage}\end{tcolorbox}\hfill\begin{tcolorbox}[enhanced,nobeforeafter,width=7.1cm,height=2.1cm,valign=center,colback=cgra!6,colframe=cgra,arc=2.5mm,boxrule=0.9pt,left=3mm]{\Large\bfseries\color{cgra}b)}\quad\begin{minipage}[c]{5.3cm}\raggedright\Large $24~\text{L/min}$\end{minipage}\end{tcolorbox}\par\vspace{2.5mm}
\begin{tcolorbox}[enhanced,nobeforeafter,width=7.1cm,height=2.1cm,valign=center,colback=cgra!6,colframe=cgra,arc=2.5mm,boxrule=0.9pt,left=3mm]{\Large\bfseries\color{cgra}c)}\quad\begin{minipage}[c]{5.3cm}\raggedright\Large $600~\text{L/min}$\end{minipage}\end{tcolorbox}\hfill\begin{tcolorbox}[enhanced,nobeforeafter,width=7.1cm,height=2.1cm,valign=center,colback=cgra!6,colframe=cgra,arc=2.5mm,boxrule=0.9pt,left=3mm]{\Large\bfseries\color{cgra}d)}\quad\begin{minipage}[c]{5.3cm}\raggedright\Large $26~\text{L/min}$\end{minipage}\end{tcolorbox}
\end{tcolorbox}
\newpage
\begin{tcolorbox}[enhanced,width=16cm,height=14cm,colback=white,colframe=cgra,boxrule=2pt,arc=0mm,
left=6mm,right=6mm,top=0mm,bottom=4mm,sharp corners,
overlay={\fill[cgra] (frame.north west) rectangle ([yshift=-1.25cm]frame.north east);
\node[anchor=west,text=white,font=\sffamily\bfseries\LARGE] at ([xshift=5mm,yshift=-0.62cm]frame.north west) {Question 90};
\node[anchor=east,text=white,font=\sffamily\itshape\large] at ([xshift=-5mm,yshift=-0.62cm]frame.north east) {Grandeurs et proportionnalité};}]
\vspace*{1.75cm}
\begin{minipage}[t][6.9cm][c]{\linewidth}\LARGE\centering
\begin{minipage}{0.95\linewidth}\raggedright\linespread{1.25}\selectfont Un robinet remplit une cuve de $90$ L en $6$ min. Quel est son débit~?\end{minipage}
\end{minipage}

\vfill
\begin{tcolorbox}[enhanced,nobeforeafter,width=7.1cm,height=2.1cm,valign=center,colback=cgra!6,colframe=cgra,arc=2.5mm,boxrule=0.9pt,left=3mm]{\Large\bfseries\color{cgra}a)}\quad\begin{minipage}[c]{5.3cm}\raggedright\Large $540~\text{L/min}$\end{minipage}\end{tcolorbox}\hfill\begin{tcolorbox}[enhanced,nobeforeafter,width=7.1cm,height=2.1cm,valign=center,colback=cgra!6,colframe=cgra,arc=2.5mm,boxrule=0.9pt,left=3mm]{\Large\bfseries\color{cgra}b)}\quad\begin{minipage}[c]{5.3cm}\raggedright\Large $15~\text{L/min}$\end{minipage}\end{tcolorbox}\par\vspace{2.5mm}
\begin{tcolorbox}[enhanced,nobeforeafter,width=7.1cm,height=2.1cm,valign=center,colback=cgra!6,colframe=cgra,arc=2.5mm,boxrule=0.9pt,left=3mm]{\Large\bfseries\color{cgra}c)}\quad\begin{minipage}[c]{5.3cm}\raggedright\Large $17~\text{L/min}$\end{minipage}\end{tcolorbox}\hfill\begin{tcolorbox}[enhanced,nobeforeafter,width=7.1cm,height=2.1cm,valign=center,colback=cgra!6,colframe=cgra,arc=2.5mm,boxrule=0.9pt,left=3mm]{\Large\bfseries\color{cgra}d)}\quad\begin{minipage}[c]{5.3cm}\raggedright\large $\dfrac{20}{3}~\text{L/min}$\end{minipage}\end{tcolorbox}
\end{tcolorbox}
\newpage
\begin{tcolorbox}[enhanced,width=16cm,height=14cm,colback=white,colframe=cgra,boxrule=2pt,arc=0mm,
left=6mm,right=6mm,top=0mm,bottom=4mm,sharp corners,
overlay={\fill[cgra] (frame.north west) rectangle ([yshift=-1.25cm]frame.north east);
\node[anchor=west,text=white,font=\sffamily\bfseries\LARGE] at ([xshift=5mm,yshift=-0.62cm]frame.north west) {Question 91};
\node[anchor=east,text=white,font=\sffamily\itshape\large] at ([xshift=-5mm,yshift=-0.62cm]frame.north east) {Grandeurs et proportionnalité};}]
\vspace*{1.75cm}
\begin{minipage}[t][6.9cm][c]{\linewidth}\LARGE\centering
\begin{minipage}{0.95\linewidth}\raggedright\linespread{1.25}\selectfont Sur une carte à l'échelle $\dfrac{1}{1\,000}$, deux points sont distants de $4$ cm. Quelle est la distance réelle~?\end{minipage}
\end{minipage}

\vfill
\begin{tcolorbox}[enhanced,nobeforeafter,width=7.1cm,height=2.1cm,valign=center,colback=cgra!6,colframe=cgra,arc=2.5mm,boxrule=0.9pt,left=3mm]{\Large\bfseries\color{cgra}a)}\quad\begin{minipage}[c]{5.3cm}\raggedright\Large $40~\text{m}$\end{minipage}\end{tcolorbox}\hfill\begin{tcolorbox}[enhanced,nobeforeafter,width=7.1cm,height=2.1cm,valign=center,colback=cgra!6,colframe=cgra,arc=2.5mm,boxrule=0.9pt,left=3mm]{\Large\bfseries\color{cgra}b)}\quad\begin{minipage}[c]{5.3cm}\raggedright\Large $400~\text{m}$\end{minipage}\end{tcolorbox}\par\vspace{2.5mm}
\begin{tcolorbox}[enhanced,nobeforeafter,width=7.1cm,height=2.1cm,valign=center,colback=cgra!6,colframe=cgra,arc=2.5mm,boxrule=0.9pt,left=3mm]{\Large\bfseries\color{cgra}c)}\quad\begin{minipage}[c]{5.3cm}\raggedright\Large $4~\text{m}$\end{minipage}\end{tcolorbox}\hfill\begin{tcolorbox}[enhanced,nobeforeafter,width=7.1cm,height=2.1cm,valign=center,colback=cgra!6,colframe=cgra,arc=2.5mm,boxrule=0.9pt,left=3mm]{\Large\bfseries\color{cgra}d)}\quad\begin{minipage}[c]{5.3cm}\raggedright\Large $4\,000~\text{m}$\end{minipage}\end{tcolorbox}
\end{tcolorbox}
\newpage
\begin{tcolorbox}[enhanced,width=16cm,height=14cm,colback=white,colframe=cgra,boxrule=2pt,arc=0mm,
left=6mm,right=6mm,top=0mm,bottom=4mm,sharp corners,
overlay={\fill[cgra] (frame.north west) rectangle ([yshift=-1.25cm]frame.north east);
\node[anchor=west,text=white,font=\sffamily\bfseries\LARGE] at ([xshift=5mm,yshift=-0.62cm]frame.north west) {Question 92};
\node[anchor=east,text=white,font=\sffamily\itshape\large] at ([xshift=-5mm,yshift=-0.62cm]frame.north east) {Grandeurs et proportionnalité};}]
\vspace*{1.75cm}
\begin{minipage}[t][6.9cm][c]{\linewidth}\LARGE\centering
\begin{minipage}{0.95\linewidth}\raggedright\linespread{1.25}\selectfont Sur une carte à l'échelle $\dfrac{1}{500}$, deux points sont distants de $6$ cm. Quelle est la distance réelle~?\end{minipage}
\end{minipage}

\vfill
\begin{tcolorbox}[enhanced,nobeforeafter,width=7.1cm,height=2.1cm,valign=center,colback=cgra!6,colframe=cgra,arc=2.5mm,boxrule=0.9pt,left=3mm]{\Large\bfseries\color{cgra}a)}\quad\begin{minipage}[c]{5.3cm}\raggedright\Large $3~\text{m}$\end{minipage}\end{tcolorbox}\hfill\begin{tcolorbox}[enhanced,nobeforeafter,width=7.1cm,height=2.1cm,valign=center,colback=cgra!6,colframe=cgra,arc=2.5mm,boxrule=0.9pt,left=3mm]{\Large\bfseries\color{cgra}b)}\quad\begin{minipage}[c]{5.3cm}\raggedright\Large $30~\text{m}$\end{minipage}\end{tcolorbox}\par\vspace{2.5mm}
\begin{tcolorbox}[enhanced,nobeforeafter,width=7.1cm,height=2.1cm,valign=center,colback=cgra!6,colframe=cgra,arc=2.5mm,boxrule=0.9pt,left=3mm]{\Large\bfseries\color{cgra}c)}\quad\begin{minipage}[c]{5.3cm}\raggedright\Large $3\,000~\text{m}$\end{minipage}\end{tcolorbox}\hfill\begin{tcolorbox}[enhanced,nobeforeafter,width=7.1cm,height=2.1cm,valign=center,colback=cgra!6,colframe=cgra,arc=2.5mm,boxrule=0.9pt,left=3mm]{\Large\bfseries\color{cgra}d)}\quad\begin{minipage}[c]{5.3cm}\raggedright\Large $300~\text{m}$\end{minipage}\end{tcolorbox}
\end{tcolorbox}
\newpage
\begin{tcolorbox}[enhanced,width=16cm,height=14cm,colback=white,colframe=cgra,boxrule=2pt,arc=0mm,
left=6mm,right=6mm,top=0mm,bottom=4mm,sharp corners,
overlay={\fill[cgra] (frame.north west) rectangle ([yshift=-1.25cm]frame.north east);
\node[anchor=west,text=white,font=\sffamily\bfseries\LARGE] at ([xshift=5mm,yshift=-0.62cm]frame.north west) {Question 93};
\node[anchor=east,text=white,font=\sffamily\itshape\large] at ([xshift=-5mm,yshift=-0.62cm]frame.north east) {Grandeurs et proportionnalité};}]
\vspace*{1.75cm}
\begin{minipage}[t][6.9cm][c]{\linewidth}\LARGE\centering
\begin{minipage}{0.95\linewidth}\raggedright\linespread{1.25}\selectfont On partage $60$ bonbons entre Léa et Noa dans le ratio $2:3$. Combien chacun en reçoit-il (Léa, puis Noa)~?\end{minipage}
\end{minipage}

\vfill
\begin{tcolorbox}[enhanced,nobeforeafter,width=7.1cm,height=2.1cm,valign=center,colback=cgra!6,colframe=cgra,arc=2.5mm,boxrule=0.9pt,left=3mm]{\Large\bfseries\color{cgra}a)}\quad\begin{minipage}[c]{5.3cm}\raggedright\Large $36\text{ et }24$\end{minipage}\end{tcolorbox}\hfill\begin{tcolorbox}[enhanced,nobeforeafter,width=7.1cm,height=2.1cm,valign=center,colback=cgra!6,colframe=cgra,arc=2.5mm,boxrule=0.9pt,left=3mm]{\Large\bfseries\color{cgra}b)}\quad\begin{minipage}[c]{5.3cm}\raggedright\Large $24\text{ et }36$\end{minipage}\end{tcolorbox}\par\vspace{2.5mm}
\begin{tcolorbox}[enhanced,nobeforeafter,width=7.1cm,height=2.1cm,valign=center,colback=cgra!6,colframe=cgra,arc=2.5mm,boxrule=0.9pt,left=3mm]{\Large\bfseries\color{cgra}c)}\quad\begin{minipage}[c]{5.3cm}\raggedright\Large $20\text{ et }40$\end{minipage}\end{tcolorbox}\hfill\begin{tcolorbox}[enhanced,nobeforeafter,width=7.1cm,height=2.1cm,valign=center,colback=cgra!6,colframe=cgra,arc=2.5mm,boxrule=0.9pt,left=3mm]{\Large\bfseries\color{cgra}d)}\quad\begin{minipage}[c]{5.3cm}\raggedright\Large $30\text{ et }30$\end{minipage}\end{tcolorbox}
\end{tcolorbox}
\newpage
\begin{tcolorbox}[enhanced,width=16cm,height=14cm,colback=white,colframe=cgra,boxrule=2pt,arc=0mm,
left=6mm,right=6mm,top=0mm,bottom=4mm,sharp corners,
overlay={\fill[cgra] (frame.north west) rectangle ([yshift=-1.25cm]frame.north east);
\node[anchor=west,text=white,font=\sffamily\bfseries\LARGE] at ([xshift=5mm,yshift=-0.62cm]frame.north west) {Question 94};
\node[anchor=east,text=white,font=\sffamily\itshape\large] at ([xshift=-5mm,yshift=-0.62cm]frame.north east) {Grandeurs et proportionnalité};}]
\vspace*{1.75cm}
\begin{minipage}[t][6.9cm][c]{\linewidth}\LARGE\centering
\begin{minipage}{0.95\linewidth}\raggedright\linespread{1.25}\selectfont On partage $84$ bonbons entre Léa et Noa dans le ratio $3:4$. Combien chacun en reçoit-il (Léa, puis Noa)~?\end{minipage}
\end{minipage}

\vfill
\begin{tcolorbox}[enhanced,nobeforeafter,width=7.1cm,height=2.1cm,valign=center,colback=cgra!6,colframe=cgra,arc=2.5mm,boxrule=0.9pt,left=3mm]{\Large\bfseries\color{cgra}a)}\quad\begin{minipage}[c]{5.3cm}\raggedright\Large $36\text{ et }48$\end{minipage}\end{tcolorbox}\hfill\begin{tcolorbox}[enhanced,nobeforeafter,width=7.1cm,height=2.1cm,valign=center,colback=cgra!6,colframe=cgra,arc=2.5mm,boxrule=0.9pt,left=3mm]{\Large\bfseries\color{cgra}b)}\quad\begin{minipage}[c]{5.3cm}\raggedright\Large $35\text{ et }49$\end{minipage}\end{tcolorbox}\par\vspace{2.5mm}
\begin{tcolorbox}[enhanced,nobeforeafter,width=7.1cm,height=2.1cm,valign=center,colback=cgra!6,colframe=cgra,arc=2.5mm,boxrule=0.9pt,left=3mm]{\Large\bfseries\color{cgra}c)}\quad\begin{minipage}[c]{5.3cm}\raggedright\Large $48\text{ et }36$\end{minipage}\end{tcolorbox}\hfill\begin{tcolorbox}[enhanced,nobeforeafter,width=7.1cm,height=2.1cm,valign=center,colback=cgra!6,colframe=cgra,arc=2.5mm,boxrule=0.9pt,left=3mm]{\Large\bfseries\color{cgra}d)}\quad\begin{minipage}[c]{5.3cm}\raggedright\Large $42\text{ et }42$\end{minipage}\end{tcolorbox}
\end{tcolorbox}
\newpage
\begin{tcolorbox}[enhanced,width=16cm,height=14cm,colback=white,colframe=cgra,boxrule=2pt,arc=0mm,
left=6mm,right=6mm,top=0mm,bottom=4mm,sharp corners,
overlay={\fill[cgra] (frame.north west) rectangle ([yshift=-1.25cm]frame.north east);
\node[anchor=west,text=white,font=\sffamily\bfseries\LARGE] at ([xshift=5mm,yshift=-0.62cm]frame.north west) {Question 95};
\node[anchor=east,text=white,font=\sffamily\itshape\large] at ([xshift=-5mm,yshift=-0.62cm]frame.north east) {Grandeurs et proportionnalité};}]
\vspace*{1.75cm}
\begin{minipage}[t][6.9cm][c]{\linewidth}\LARGE\centering
\begin{minipage}{0.95\linewidth}\raggedright\linespread{1.25}\selectfont Compléter~: $2{,}5~\text{m}^3=\;?\;\text{L}$\end{minipage}
\end{minipage}

\vfill
\begin{tcolorbox}[enhanced,nobeforeafter,width=7.1cm,height=2.1cm,valign=center,colback=cgra!6,colframe=cgra,arc=2.5mm,boxrule=0.9pt,left=3mm]{\Large\bfseries\color{cgra}a)}\quad\begin{minipage}[c]{5.3cm}\raggedright\Large $250~\text{L}$\end{minipage}\end{tcolorbox}\hfill\begin{tcolorbox}[enhanced,nobeforeafter,width=7.1cm,height=2.1cm,valign=center,colback=cgra!6,colframe=cgra,arc=2.5mm,boxrule=0.9pt,left=3mm]{\Large\bfseries\color{cgra}b)}\quad\begin{minipage}[c]{5.3cm}\raggedright\Large $25\,000~\text{L}$\end{minipage}\end{tcolorbox}\par\vspace{2.5mm}
\begin{tcolorbox}[enhanced,nobeforeafter,width=7.1cm,height=2.1cm,valign=center,colback=cgra!6,colframe=cgra,arc=2.5mm,boxrule=0.9pt,left=3mm]{\Large\bfseries\color{cgra}c)}\quad\begin{minipage}[c]{5.3cm}\raggedright\Large $250\,000~\text{L}$\end{minipage}\end{tcolorbox}\hfill\begin{tcolorbox}[enhanced,nobeforeafter,width=7.1cm,height=2.1cm,valign=center,colback=cgra!6,colframe=cgra,arc=2.5mm,boxrule=0.9pt,left=3mm]{\Large\bfseries\color{cgra}d)}\quad\begin{minipage}[c]{5.3cm}\raggedright\Large $2\,500~\text{L}$\end{minipage}\end{tcolorbox}
\end{tcolorbox}
\newpage
\begin{tcolorbox}[enhanced,width=16cm,height=14cm,colback=white,colframe=cgra,boxrule=2pt,arc=0mm,
left=6mm,right=6mm,top=0mm,bottom=4mm,sharp corners,
overlay={\fill[cgra] (frame.north west) rectangle ([yshift=-1.25cm]frame.north east);
\node[anchor=west,text=white,font=\sffamily\bfseries\LARGE] at ([xshift=5mm,yshift=-0.62cm]frame.north west) {Question 96};
\node[anchor=east,text=white,font=\sffamily\itshape\large] at ([xshift=-5mm,yshift=-0.62cm]frame.north east) {Grandeurs et proportionnalité};}]
\vspace*{1.75cm}
\begin{minipage}[t][6.9cm][c]{\linewidth}\LARGE\centering
\begin{minipage}{0.95\linewidth}\raggedright\linespread{1.25}\selectfont Compléter~: $750~\text{mL}=\;?\;\text{cm}^3$\end{minipage}
\end{minipage}

\vfill
\begin{tcolorbox}[enhanced,nobeforeafter,width=7.1cm,height=2.1cm,valign=center,colback=cgra!6,colframe=cgra,arc=2.5mm,boxrule=0.9pt,left=3mm]{\Large\bfseries\color{cgra}a)}\quad\begin{minipage}[c]{5.3cm}\raggedright\Large $0{,}75~\text{cm}^3$\end{minipage}\end{tcolorbox}\hfill\begin{tcolorbox}[enhanced,nobeforeafter,width=7.1cm,height=2.1cm,valign=center,colback=cgra!6,colframe=cgra,arc=2.5mm,boxrule=0.9pt,left=3mm]{\Large\bfseries\color{cgra}b)}\quad\begin{minipage}[c]{5.3cm}\raggedright\Large $750~\text{cm}^3$\end{minipage}\end{tcolorbox}\par\vspace{2.5mm}
\begin{tcolorbox}[enhanced,nobeforeafter,width=7.1cm,height=2.1cm,valign=center,colback=cgra!6,colframe=cgra,arc=2.5mm,boxrule=0.9pt,left=3mm]{\Large\bfseries\color{cgra}c)}\quad\begin{minipage}[c]{5.3cm}\raggedright\Large $7\,500~\text{cm}^3$\end{minipage}\end{tcolorbox}\hfill\begin{tcolorbox}[enhanced,nobeforeafter,width=7.1cm,height=2.1cm,valign=center,colback=cgra!6,colframe=cgra,arc=2.5mm,boxrule=0.9pt,left=3mm]{\Large\bfseries\color{cgra}d)}\quad\begin{minipage}[c]{5.3cm}\raggedright\Large $750\,000~\text{cm}^3$\end{minipage}\end{tcolorbox}
\end{tcolorbox}
\newpage
\begin{tcolorbox}[enhanced,width=16cm,height=14cm,colback=white,colframe=cari,boxrule=2pt,arc=0mm,
left=6mm,right=6mm,top=0mm,bottom=4mm,sharp corners,
overlay={\fill[cari] (frame.north west) rectangle ([yshift=-1.25cm]frame.north east);
\node[anchor=west,text=white,font=\sffamily\bfseries\LARGE] at ([xshift=5mm,yshift=-0.62cm]frame.north west) {Question 1 B};
\node[anchor=east,text=white,font=\sffamily\itshape\large] at ([xshift=-5mm,yshift=-0.62cm]frame.north east) {Arithmétique et puissances};}]
\vspace*{1.75cm}
\begin{minipage}[t][6.9cm][c]{\linewidth}\LARGE\centering
\begin{minipage}{0.95\linewidth}\raggedright\linespread{1.25}\selectfont Lequel de ces nombres est premier~?\end{minipage}
\end{minipage}

\vfill
\begin{tcolorbox}[enhanced,nobeforeafter,width=7.1cm,height=2.1cm,valign=center,colback=cari!6,colframe=cari,arc=2.5mm,boxrule=0.9pt,left=3mm]{\Large\bfseries\color{cari}a)}\quad\begin{minipage}[c]{5.3cm}\raggedright\Large $87$\end{minipage}\end{tcolorbox}\hfill\begin{tcolorbox}[enhanced,nobeforeafter,width=7.1cm,height=2.1cm,valign=center,colback=cari!6,colframe=cari,arc=2.5mm,boxrule=0.9pt,left=3mm]{\Large\bfseries\color{cari}b)}\quad\begin{minipage}[c]{5.3cm}\raggedright\Large $91$\end{minipage}\end{tcolorbox}\par\vspace{2.5mm}
\begin{tcolorbox}[enhanced,nobeforeafter,width=7.1cm,height=2.1cm,valign=center,colback=cari!6,colframe=cari,arc=2.5mm,boxrule=0.9pt,left=3mm]{\Large\bfseries\color{cari}c)}\quad\begin{minipage}[c]{5.3cm}\raggedright\Large $93$\end{minipage}\end{tcolorbox}\hfill\begin{tcolorbox}[enhanced,nobeforeafter,width=7.1cm,height=2.1cm,valign=center,colback=cari!6,colframe=cari,arc=2.5mm,boxrule=0.9pt,left=3mm]{\Large\bfseries\color{cari}d)}\quad\begin{minipage}[c]{5.3cm}\raggedright\Large $83$\end{minipage}\end{tcolorbox}
\end{tcolorbox}
\newpage
\begin{tcolorbox}[enhanced,width=16cm,height=14cm,colback=white,colframe=cari,boxrule=2pt,arc=0mm,
left=6mm,right=6mm,top=0mm,bottom=4mm,sharp corners,
overlay={\fill[cari] (frame.north west) rectangle ([yshift=-1.25cm]frame.north east);
\node[anchor=west,text=white,font=\sffamily\bfseries\LARGE] at ([xshift=5mm,yshift=-0.62cm]frame.north west) {Question 2 B};
\node[anchor=east,text=white,font=\sffamily\itshape\large] at ([xshift=-5mm,yshift=-0.62cm]frame.north east) {Arithmétique et puissances};}]
\vspace*{1.75cm}
\begin{minipage}[t][6.9cm][c]{\linewidth}\LARGE\centering
\begin{minipage}{0.95\linewidth}\raggedright\linespread{1.25}\selectfont Lequel de ces nombres est premier~?\end{minipage}
\end{minipage}

\vfill
\begin{tcolorbox}[enhanced,nobeforeafter,width=7.1cm,height=2.1cm,valign=center,colback=cari!6,colframe=cari,arc=2.5mm,boxrule=0.9pt,left=3mm]{\Large\bfseries\color{cari}a)}\quad\begin{minipage}[c]{5.3cm}\raggedright\Large $71$\end{minipage}\end{tcolorbox}\hfill\begin{tcolorbox}[enhanced,nobeforeafter,width=7.1cm,height=2.1cm,valign=center,colback=cari!6,colframe=cari,arc=2.5mm,boxrule=0.9pt,left=3mm]{\Large\bfseries\color{cari}b)}\quad\begin{minipage}[c]{5.3cm}\raggedright\Large $75$\end{minipage}\end{tcolorbox}\par\vspace{2.5mm}
\begin{tcolorbox}[enhanced,nobeforeafter,width=7.1cm,height=2.1cm,valign=center,colback=cari!6,colframe=cari,arc=2.5mm,boxrule=0.9pt,left=3mm]{\Large\bfseries\color{cari}c)}\quad\begin{minipage}[c]{5.3cm}\raggedright\Large $69$\end{minipage}\end{tcolorbox}\hfill\begin{tcolorbox}[enhanced,nobeforeafter,width=7.1cm,height=2.1cm,valign=center,colback=cari!6,colframe=cari,arc=2.5mm,boxrule=0.9pt,left=3mm]{\Large\bfseries\color{cari}d)}\quad\begin{minipage}[c]{5.3cm}\raggedright\Large $77$\end{minipage}\end{tcolorbox}
\end{tcolorbox}
\newpage
\begin{tcolorbox}[enhanced,width=16cm,height=14cm,colback=white,colframe=cari,boxrule=2pt,arc=0mm,
left=6mm,right=6mm,top=0mm,bottom=4mm,sharp corners,
overlay={\fill[cari] (frame.north west) rectangle ([yshift=-1.25cm]frame.north east);
\node[anchor=west,text=white,font=\sffamily\bfseries\LARGE] at ([xshift=5mm,yshift=-0.62cm]frame.north west) {Question 3 B};
\node[anchor=east,text=white,font=\sffamily\itshape\large] at ([xshift=-5mm,yshift=-0.62cm]frame.north east) {Arithmétique et puissances};}]
\vspace*{1.75cm}
\begin{minipage}[t][6.9cm][c]{\linewidth}\LARGE\centering
\begin{minipage}{0.95\linewidth}\raggedright\linespread{1.25}\selectfont Quelle est la décomposition de $126$ en produit de facteurs premiers~?\end{minipage}
\end{minipage}

\vfill
\begin{tcolorbox}[enhanced,nobeforeafter,width=7.1cm,height=2.1cm,valign=center,colback=cari!6,colframe=cari,arc=2.5mm,boxrule=0.9pt,left=3mm]{\Large\bfseries\color{cari}a)}\quad\begin{minipage}[c]{5.3cm}\raggedright\Large $2^{2}\times3^{2}\times7$\end{minipage}\end{tcolorbox}\hfill\begin{tcolorbox}[enhanced,nobeforeafter,width=7.1cm,height=2.1cm,valign=center,colback=cari!6,colframe=cari,arc=2.5mm,boxrule=0.9pt,left=3mm]{\Large\bfseries\color{cari}b)}\quad\begin{minipage}[c]{5.3cm}\raggedright\Large $2\times3^{2}\times7^{2}$\end{minipage}\end{tcolorbox}\par\vspace{2.5mm}
\begin{tcolorbox}[enhanced,nobeforeafter,width=7.1cm,height=2.1cm,valign=center,colback=cari!6,colframe=cari,arc=2.5mm,boxrule=0.9pt,left=3mm]{\Large\bfseries\color{cari}c)}\quad\begin{minipage}[c]{5.3cm}\raggedright\Large $2\times3^{2}\times7$\end{minipage}\end{tcolorbox}\hfill\begin{tcolorbox}[enhanced,nobeforeafter,width=7.1cm,height=2.1cm,valign=center,colback=cari!6,colframe=cari,arc=2.5mm,boxrule=0.9pt,left=3mm]{\Large\bfseries\color{cari}d)}\quad\begin{minipage}[c]{5.3cm}\raggedright\Large $6\times21$\end{minipage}\end{tcolorbox}
\end{tcolorbox}
\newpage
\begin{tcolorbox}[enhanced,width=16cm,height=14cm,colback=white,colframe=cari,boxrule=2pt,arc=0mm,
left=6mm,right=6mm,top=0mm,bottom=4mm,sharp corners,
overlay={\fill[cari] (frame.north west) rectangle ([yshift=-1.25cm]frame.north east);
\node[anchor=west,text=white,font=\sffamily\bfseries\LARGE] at ([xshift=5mm,yshift=-0.62cm]frame.north west) {Question 4 B};
\node[anchor=east,text=white,font=\sffamily\itshape\large] at ([xshift=-5mm,yshift=-0.62cm]frame.north east) {Arithmétique et puissances};}]
\vspace*{1.75cm}
\begin{minipage}[t][6.9cm][c]{\linewidth}\LARGE\centering
\begin{minipage}{0.95\linewidth}\raggedright\linespread{1.25}\selectfont Quelle est la décomposition de $600$ en produit de facteurs premiers~?\end{minipage}
\end{minipage}

\vfill
\begin{tcolorbox}[enhanced,nobeforeafter,width=7.1cm,height=2.1cm,valign=center,colback=cari!6,colframe=cari,arc=2.5mm,boxrule=0.9pt,left=3mm]{\Large\bfseries\color{cari}a)}\quad\begin{minipage}[c]{5.3cm}\raggedright\Large $24\times25$\end{minipage}\end{tcolorbox}\hfill\begin{tcolorbox}[enhanced,nobeforeafter,width=7.1cm,height=2.1cm,valign=center,colback=cari!6,colframe=cari,arc=2.5mm,boxrule=0.9pt,left=3mm]{\Large\bfseries\color{cari}b)}\quad\begin{minipage}[c]{5.3cm}\raggedright\Large $2^{3}\times3\times5^{2}$\end{minipage}\end{tcolorbox}\par\vspace{2.5mm}
\begin{tcolorbox}[enhanced,nobeforeafter,width=7.1cm,height=2.1cm,valign=center,colback=cari!6,colframe=cari,arc=2.5mm,boxrule=0.9pt,left=3mm]{\Large\bfseries\color{cari}c)}\quad\begin{minipage}[c]{5.3cm}\raggedright\Large $2^{3}\times3\times5^{3}$\end{minipage}\end{tcolorbox}\hfill\begin{tcolorbox}[enhanced,nobeforeafter,width=7.1cm,height=2.1cm,valign=center,colback=cari!6,colframe=cari,arc=2.5mm,boxrule=0.9pt,left=3mm]{\Large\bfseries\color{cari}d)}\quad\begin{minipage}[c]{5.3cm}\raggedright\Large $2^{4}\times3\times5^{2}$\end{minipage}\end{tcolorbox}
\end{tcolorbox}
\newpage
\begin{tcolorbox}[enhanced,width=16cm,height=14cm,colback=white,colframe=cari,boxrule=2pt,arc=0mm,
left=6mm,right=6mm,top=0mm,bottom=4mm,sharp corners,
overlay={\fill[cari] (frame.north west) rectangle ([yshift=-1.25cm]frame.north east);
\node[anchor=west,text=white,font=\sffamily\bfseries\LARGE] at ([xshift=5mm,yshift=-0.62cm]frame.north west) {Question 5 B};
\node[anchor=east,text=white,font=\sffamily\itshape\large] at ([xshift=-5mm,yshift=-0.62cm]frame.north east) {Arithmétique et puissances};}]
\vspace*{1.75cm}
\begin{minipage}[t][6.9cm][c]{\linewidth}\LARGE\centering
\begin{minipage}{0.95\linewidth}\raggedright\linespread{1.25}\selectfont Quelle est la forme irréductible de $\dfrac{150}{375}$~?\end{minipage}
\end{minipage}

\vfill
\begin{tcolorbox}[enhanced,nobeforeafter,width=7.1cm,height=2.1cm,valign=center,colback=cari!6,colframe=cari,arc=2.5mm,boxrule=0.9pt,left=3mm]{\Large\bfseries\color{cari}a)}\quad\begin{minipage}[c]{5.3cm}\raggedright\Large $\dfrac{2}{5}$\end{minipage}\end{tcolorbox}\hfill\begin{tcolorbox}[enhanced,nobeforeafter,width=7.1cm,height=2.1cm,valign=center,colback=cari!6,colframe=cari,arc=2.5mm,boxrule=0.9pt,left=3mm]{\Large\bfseries\color{cari}b)}\quad\begin{minipage}[c]{5.3cm}\raggedright\Large $\dfrac{30}{75}$\end{minipage}\end{tcolorbox}\par\vspace{2.5mm}
\begin{tcolorbox}[enhanced,nobeforeafter,width=7.1cm,height=2.1cm,valign=center,colback=cari!6,colframe=cari,arc=2.5mm,boxrule=0.9pt,left=3mm]{\Large\bfseries\color{cari}c)}\quad\begin{minipage}[c]{5.3cm}\raggedright\Large $\dfrac{5}{2}$\end{minipage}\end{tcolorbox}\hfill\begin{tcolorbox}[enhanced,nobeforeafter,width=7.1cm,height=2.1cm,valign=center,colback=cari!6,colframe=cari,arc=2.5mm,boxrule=0.9pt,left=3mm]{\Large\bfseries\color{cari}d)}\quad\begin{minipage}[c]{5.3cm}\raggedright\Large $\dfrac{50}{125}$\end{minipage}\end{tcolorbox}
\end{tcolorbox}
\newpage
\begin{tcolorbox}[enhanced,width=16cm,height=14cm,colback=white,colframe=cari,boxrule=2pt,arc=0mm,
left=6mm,right=6mm,top=0mm,bottom=4mm,sharp corners,
overlay={\fill[cari] (frame.north west) rectangle ([yshift=-1.25cm]frame.north east);
\node[anchor=west,text=white,font=\sffamily\bfseries\LARGE] at ([xshift=5mm,yshift=-0.62cm]frame.north west) {Question 6 B};
\node[anchor=east,text=white,font=\sffamily\itshape\large] at ([xshift=-5mm,yshift=-0.62cm]frame.north east) {Arithmétique et puissances};}]
\vspace*{1.75cm}
\begin{minipage}[t][6.9cm][c]{\linewidth}\LARGE\centering
\begin{minipage}{0.95\linewidth}\raggedright\linespread{1.25}\selectfont Quelle est la forme irréductible de $\dfrac{84}{196}$~?\end{minipage}
\end{minipage}

\vfill
\begin{tcolorbox}[enhanced,nobeforeafter,width=7.1cm,height=2.1cm,valign=center,colback=cari!6,colframe=cari,arc=2.5mm,boxrule=0.9pt,left=3mm]{\Large\bfseries\color{cari}a)}\quad\begin{minipage}[c]{5.3cm}\raggedright\Large $\dfrac{12}{28}$\end{minipage}\end{tcolorbox}\hfill\begin{tcolorbox}[enhanced,nobeforeafter,width=7.1cm,height=2.1cm,valign=center,colback=cari!6,colframe=cari,arc=2.5mm,boxrule=0.9pt,left=3mm]{\Large\bfseries\color{cari}b)}\quad\begin{minipage}[c]{5.3cm}\raggedright\Large $\dfrac{7}{3}$\end{minipage}\end{tcolorbox}\par\vspace{2.5mm}
\begin{tcolorbox}[enhanced,nobeforeafter,width=7.1cm,height=2.1cm,valign=center,colback=cari!6,colframe=cari,arc=2.5mm,boxrule=0.9pt,left=3mm]{\Large\bfseries\color{cari}c)}\quad\begin{minipage}[c]{5.3cm}\raggedright\Large $\dfrac{42}{98}$\end{minipage}\end{tcolorbox}\hfill\begin{tcolorbox}[enhanced,nobeforeafter,width=7.1cm,height=2.1cm,valign=center,colback=cari!6,colframe=cari,arc=2.5mm,boxrule=0.9pt,left=3mm]{\Large\bfseries\color{cari}d)}\quad\begin{minipage}[c]{5.3cm}\raggedright\Large $\dfrac{3}{7}$\end{minipage}\end{tcolorbox}
\end{tcolorbox}
\newpage
\begin{tcolorbox}[enhanced,width=16cm,height=14cm,colback=white,colframe=cari,boxrule=2pt,arc=0mm,
left=6mm,right=6mm,top=0mm,bottom=4mm,sharp corners,
overlay={\fill[cari] (frame.north west) rectangle ([yshift=-1.25cm]frame.north east);
\node[anchor=west,text=white,font=\sffamily\bfseries\LARGE] at ([xshift=5mm,yshift=-0.62cm]frame.north west) {Question 7 B};
\node[anchor=east,text=white,font=\sffamily\itshape\large] at ([xshift=-5mm,yshift=-0.62cm]frame.north east) {Arithmétique et puissances};}]
\vspace*{1.75cm}
\begin{minipage}[t][6.9cm][c]{\linewidth}\LARGE\centering
\begin{minipage}{0.95\linewidth}\raggedright\linespread{1.25}\selectfont Quel est le plus grand diviseur commun à $72$ et $120$~?\end{minipage}
\end{minipage}

\vfill
\begin{tcolorbox}[enhanced,nobeforeafter,width=7.1cm,height=2.1cm,valign=center,colback=cari!6,colframe=cari,arc=2.5mm,boxrule=0.9pt,left=3mm]{\Large\bfseries\color{cari}a)}\quad\begin{minipage}[c]{5.3cm}\raggedright\Large $24$\end{minipage}\end{tcolorbox}\hfill\begin{tcolorbox}[enhanced,nobeforeafter,width=7.1cm,height=2.1cm,valign=center,colback=cari!6,colframe=cari,arc=2.5mm,boxrule=0.9pt,left=3mm]{\Large\bfseries\color{cari}b)}\quad\begin{minipage}[c]{5.3cm}\raggedright\Large $360$\end{minipage}\end{tcolorbox}\par\vspace{2.5mm}
\begin{tcolorbox}[enhanced,nobeforeafter,width=7.1cm,height=2.1cm,valign=center,colback=cari!6,colframe=cari,arc=2.5mm,boxrule=0.9pt,left=3mm]{\Large\bfseries\color{cari}c)}\quad\begin{minipage}[c]{5.3cm}\raggedright\Large $12$\end{minipage}\end{tcolorbox}\hfill\begin{tcolorbox}[enhanced,nobeforeafter,width=7.1cm,height=2.1cm,valign=center,colback=cari!6,colframe=cari,arc=2.5mm,boxrule=0.9pt,left=3mm]{\Large\bfseries\color{cari}d)}\quad\begin{minipage}[c]{5.3cm}\raggedright\Large $48$\end{minipage}\end{tcolorbox}
\end{tcolorbox}
\newpage
\begin{tcolorbox}[enhanced,width=16cm,height=14cm,colback=white,colframe=cari,boxrule=2pt,arc=0mm,
left=6mm,right=6mm,top=0mm,bottom=4mm,sharp corners,
overlay={\fill[cari] (frame.north west) rectangle ([yshift=-1.25cm]frame.north east);
\node[anchor=west,text=white,font=\sffamily\bfseries\LARGE] at ([xshift=5mm,yshift=-0.62cm]frame.north west) {Question 8 B};
\node[anchor=east,text=white,font=\sffamily\itshape\large] at ([xshift=-5mm,yshift=-0.62cm]frame.north east) {Arithmétique et puissances};}]
\vspace*{1.75cm}
\begin{minipage}[t][6.9cm][c]{\linewidth}\LARGE\centering
\begin{minipage}{0.95\linewidth}\raggedright\linespread{1.25}\selectfont Quel est le plus grand diviseur commun à $105$ et $75$~?\end{minipage}
\end{minipage}

\vfill
\begin{tcolorbox}[enhanced,nobeforeafter,width=7.1cm,height=2.1cm,valign=center,colback=cari!6,colframe=cari,arc=2.5mm,boxrule=0.9pt,left=3mm]{\Large\bfseries\color{cari}a)}\quad\begin{minipage}[c]{5.3cm}\raggedright\Large $5$\end{minipage}\end{tcolorbox}\hfill\begin{tcolorbox}[enhanced,nobeforeafter,width=7.1cm,height=2.1cm,valign=center,colback=cari!6,colframe=cari,arc=2.5mm,boxrule=0.9pt,left=3mm]{\Large\bfseries\color{cari}b)}\quad\begin{minipage}[c]{5.3cm}\raggedright\Large $15$\end{minipage}\end{tcolorbox}\par\vspace{2.5mm}
\begin{tcolorbox}[enhanced,nobeforeafter,width=7.1cm,height=2.1cm,valign=center,colback=cari!6,colframe=cari,arc=2.5mm,boxrule=0.9pt,left=3mm]{\Large\bfseries\color{cari}c)}\quad\begin{minipage}[c]{5.3cm}\raggedright\Large $525$\end{minipage}\end{tcolorbox}\hfill\begin{tcolorbox}[enhanced,nobeforeafter,width=7.1cm,height=2.1cm,valign=center,colback=cari!6,colframe=cari,arc=2.5mm,boxrule=0.9pt,left=3mm]{\Large\bfseries\color{cari}d)}\quad\begin{minipage}[c]{5.3cm}\raggedright\Large $30$\end{minipage}\end{tcolorbox}
\end{tcolorbox}
\newpage
\begin{tcolorbox}[enhanced,width=16cm,height=14cm,colback=white,colframe=cari,boxrule=2pt,arc=0mm,
left=6mm,right=6mm,top=0mm,bottom=4mm,sharp corners,
overlay={\fill[cari] (frame.north west) rectangle ([yshift=-1.25cm]frame.north east);
\node[anchor=west,text=white,font=\sffamily\bfseries\LARGE] at ([xshift=5mm,yshift=-0.62cm]frame.north west) {Question 9 B};
\node[anchor=east,text=white,font=\sffamily\itshape\large] at ([xshift=-5mm,yshift=-0.62cm]frame.north east) {Arithmétique et puissances};}]
\vspace*{1.75cm}
\begin{minipage}[t][6.9cm][c]{\linewidth}\Large\centering
\begin{minipage}{0.95\linewidth}\raggedright\linespread{1.25}\selectfont Un fleuriste a $90$ roses et $126$ tulipes. Il veut faire le plus grand nombre de bouquets identiques en utilisant toutes les fleurs. Combien de bouquets peut-il faire~?\end{minipage}
\end{minipage}

\vfill
\begin{tcolorbox}[enhanced,nobeforeafter,width=7.1cm,height=2.1cm,valign=center,colback=cari!6,colframe=cari,arc=2.5mm,boxrule=0.9pt,left=3mm]{\Large\bfseries\color{cari}a)}\quad\begin{minipage}[c]{5.3cm}\raggedright\Large $36$\end{minipage}\end{tcolorbox}\hfill\begin{tcolorbox}[enhanced,nobeforeafter,width=7.1cm,height=2.1cm,valign=center,colback=cari!6,colframe=cari,arc=2.5mm,boxrule=0.9pt,left=3mm]{\Large\bfseries\color{cari}b)}\quad\begin{minipage}[c]{5.3cm}\raggedright\Large $9$\end{minipage}\end{tcolorbox}\par\vspace{2.5mm}
\begin{tcolorbox}[enhanced,nobeforeafter,width=7.1cm,height=2.1cm,valign=center,colback=cari!6,colframe=cari,arc=2.5mm,boxrule=0.9pt,left=3mm]{\Large\bfseries\color{cari}c)}\quad\begin{minipage}[c]{5.3cm}\raggedright\Large $216$\end{minipage}\end{tcolorbox}\hfill\begin{tcolorbox}[enhanced,nobeforeafter,width=7.1cm,height=2.1cm,valign=center,colback=cari!6,colframe=cari,arc=2.5mm,boxrule=0.9pt,left=3mm]{\Large\bfseries\color{cari}d)}\quad\begin{minipage}[c]{5.3cm}\raggedright\Large $18$\end{minipage}\end{tcolorbox}
\end{tcolorbox}
\newpage
\begin{tcolorbox}[enhanced,width=16cm,height=14cm,colback=white,colframe=cari,boxrule=2pt,arc=0mm,
left=6mm,right=6mm,top=0mm,bottom=4mm,sharp corners,
overlay={\fill[cari] (frame.north west) rectangle ([yshift=-1.25cm]frame.north east);
\node[anchor=west,text=white,font=\sffamily\bfseries\LARGE] at ([xshift=5mm,yshift=-0.62cm]frame.north west) {Question 10 B};
\node[anchor=east,text=white,font=\sffamily\itshape\large] at ([xshift=-5mm,yshift=-0.62cm]frame.north east) {Arithmétique et puissances};}]
\vspace*{1.75cm}
\begin{minipage}[t][6.9cm][c]{\linewidth}\Large\centering
\begin{minipage}{0.95\linewidth}\raggedright\linespread{1.25}\selectfont Un fleuriste a $75$ roses et $105$ tulipes. Il veut faire le plus grand nombre de bouquets identiques en utilisant toutes les fleurs. Combien de bouquets peut-il faire~?\end{minipage}
\end{minipage}

\vfill
\begin{tcolorbox}[enhanced,nobeforeafter,width=7.1cm,height=2.1cm,valign=center,colback=cari!6,colframe=cari,arc=2.5mm,boxrule=0.9pt,left=3mm]{\Large\bfseries\color{cari}a)}\quad\begin{minipage}[c]{5.3cm}\raggedright\Large $30$\end{minipage}\end{tcolorbox}\hfill\begin{tcolorbox}[enhanced,nobeforeafter,width=7.1cm,height=2.1cm,valign=center,colback=cari!6,colframe=cari,arc=2.5mm,boxrule=0.9pt,left=3mm]{\Large\bfseries\color{cari}b)}\quad\begin{minipage}[c]{5.3cm}\raggedright\Large $45$\end{minipage}\end{tcolorbox}\par\vspace{2.5mm}
\begin{tcolorbox}[enhanced,nobeforeafter,width=7.1cm,height=2.1cm,valign=center,colback=cari!6,colframe=cari,arc=2.5mm,boxrule=0.9pt,left=3mm]{\Large\bfseries\color{cari}c)}\quad\begin{minipage}[c]{5.3cm}\raggedright\Large $180$\end{minipage}\end{tcolorbox}\hfill\begin{tcolorbox}[enhanced,nobeforeafter,width=7.1cm,height=2.1cm,valign=center,colback=cari!6,colframe=cari,arc=2.5mm,boxrule=0.9pt,left=3mm]{\Large\bfseries\color{cari}d)}\quad\begin{minipage}[c]{5.3cm}\raggedright\Large $15$\end{minipage}\end{tcolorbox}
\end{tcolorbox}
\newpage
\begin{tcolorbox}[enhanced,width=16cm,height=14cm,colback=white,colframe=cari,boxrule=2pt,arc=0mm,
left=6mm,right=6mm,top=0mm,bottom=4mm,sharp corners,
overlay={\fill[cari] (frame.north west) rectangle ([yshift=-1.25cm]frame.north east);
\node[anchor=west,text=white,font=\sffamily\bfseries\LARGE] at ([xshift=5mm,yshift=-0.62cm]frame.north west) {Question 11 B};
\node[anchor=east,text=white,font=\sffamily\itshape\large] at ([xshift=-5mm,yshift=-0.62cm]frame.north east) {Arithmétique et puissances};}]
\vspace*{1.75cm}
\begin{minipage}[t][6.9cm][c]{\linewidth}\Large\centering
\begin{minipage}{0.95\linewidth}\raggedright\linespread{1.25}\selectfont Un phare clignote toutes les $9$ secondes et un autre toutes les $12$ secondes. Ils viennent de clignoter ensemble. Dans combien de secondes clignoteront-ils de nouveau ensemble~?\end{minipage}
\end{minipage}

\vfill
\begin{tcolorbox}[enhanced,nobeforeafter,width=7.1cm,height=2.1cm,valign=center,colback=cari!6,colframe=cari,arc=2.5mm,boxrule=0.9pt,left=3mm]{\Large\bfseries\color{cari}a)}\quad\begin{minipage}[c]{5.3cm}\raggedright\Large $3~\text{s}$\end{minipage}\end{tcolorbox}\hfill\begin{tcolorbox}[enhanced,nobeforeafter,width=7.1cm,height=2.1cm,valign=center,colback=cari!6,colframe=cari,arc=2.5mm,boxrule=0.9pt,left=3mm]{\Large\bfseries\color{cari}b)}\quad\begin{minipage}[c]{5.3cm}\raggedright\Large $36~\text{s}$\end{minipage}\end{tcolorbox}\par\vspace{2.5mm}
\begin{tcolorbox}[enhanced,nobeforeafter,width=7.1cm,height=2.1cm,valign=center,colback=cari!6,colframe=cari,arc=2.5mm,boxrule=0.9pt,left=3mm]{\Large\bfseries\color{cari}c)}\quad\begin{minipage}[c]{5.3cm}\raggedright\Large $108~\text{s}$\end{minipage}\end{tcolorbox}\hfill\begin{tcolorbox}[enhanced,nobeforeafter,width=7.1cm,height=2.1cm,valign=center,colback=cari!6,colframe=cari,arc=2.5mm,boxrule=0.9pt,left=3mm]{\Large\bfseries\color{cari}d)}\quad\begin{minipage}[c]{5.3cm}\raggedright\Large $21~\text{s}$\end{minipage}\end{tcolorbox}
\end{tcolorbox}
\newpage
\begin{tcolorbox}[enhanced,width=16cm,height=14cm,colback=white,colframe=cari,boxrule=2pt,arc=0mm,
left=6mm,right=6mm,top=0mm,bottom=4mm,sharp corners,
overlay={\fill[cari] (frame.north west) rectangle ([yshift=-1.25cm]frame.north east);
\node[anchor=west,text=white,font=\sffamily\bfseries\LARGE] at ([xshift=5mm,yshift=-0.62cm]frame.north west) {Question 12 B};
\node[anchor=east,text=white,font=\sffamily\itshape\large] at ([xshift=-5mm,yshift=-0.62cm]frame.north east) {Arithmétique et puissances};}]
\vspace*{1.75cm}
\begin{minipage}[t][6.9cm][c]{\linewidth}\Large\centering
\begin{minipage}{0.95\linewidth}\raggedright\linespread{1.25}\selectfont Un phare clignote toutes les $15$ secondes et un autre toutes les $20$ secondes. Ils viennent de clignoter ensemble. Dans combien de secondes clignoteront-ils de nouveau ensemble~?\end{minipage}
\end{minipage}

\vfill
\begin{tcolorbox}[enhanced,nobeforeafter,width=7.1cm,height=2.1cm,valign=center,colback=cari!6,colframe=cari,arc=2.5mm,boxrule=0.9pt,left=3mm]{\Large\bfseries\color{cari}a)}\quad\begin{minipage}[c]{5.3cm}\raggedright\Large $35~\text{s}$\end{minipage}\end{tcolorbox}\hfill\begin{tcolorbox}[enhanced,nobeforeafter,width=7.1cm,height=2.1cm,valign=center,colback=cari!6,colframe=cari,arc=2.5mm,boxrule=0.9pt,left=3mm]{\Large\bfseries\color{cari}b)}\quad\begin{minipage}[c]{5.3cm}\raggedright\Large $300~\text{s}$\end{minipage}\end{tcolorbox}\par\vspace{2.5mm}
\begin{tcolorbox}[enhanced,nobeforeafter,width=7.1cm,height=2.1cm,valign=center,colback=cari!6,colframe=cari,arc=2.5mm,boxrule=0.9pt,left=3mm]{\Large\bfseries\color{cari}c)}\quad\begin{minipage}[c]{5.3cm}\raggedright\Large $60~\text{s}$\end{minipage}\end{tcolorbox}\hfill\begin{tcolorbox}[enhanced,nobeforeafter,width=7.1cm,height=2.1cm,valign=center,colback=cari!6,colframe=cari,arc=2.5mm,boxrule=0.9pt,left=3mm]{\Large\bfseries\color{cari}d)}\quad\begin{minipage}[c]{5.3cm}\raggedright\Large $5~\text{s}$\end{minipage}\end{tcolorbox}
\end{tcolorbox}
\newpage
\begin{tcolorbox}[enhanced,width=16cm,height=14cm,colback=white,colframe=cari,boxrule=2pt,arc=0mm,
left=6mm,right=6mm,top=0mm,bottom=4mm,sharp corners,
overlay={\fill[cari] (frame.north west) rectangle ([yshift=-1.25cm]frame.north east);
\node[anchor=west,text=white,font=\sffamily\bfseries\LARGE] at ([xshift=5mm,yshift=-0.62cm]frame.north west) {Question 13 B};
\node[anchor=east,text=white,font=\sffamily\itshape\large] at ([xshift=-5mm,yshift=-0.62cm]frame.north east) {Arithmétique et puissances};}]
\vspace*{1.75cm}
\begin{minipage}[t][6.9cm][c]{\linewidth}\LARGE\centering
\begin{minipage}{0.95\linewidth}\raggedright\linespread{1.25}\selectfont Que vaut $(5\times10^{3})\times(3\times10^{4})$, en notation scientifique~?\end{minipage}
\end{minipage}

\vfill
\begin{tcolorbox}[enhanced,nobeforeafter,width=7.1cm,height=2.1cm,valign=center,colback=cari!6,colframe=cari,arc=2.5mm,boxrule=0.9pt,left=3mm]{\Large\bfseries\color{cari}a)}\quad\begin{minipage}[c]{5.3cm}\raggedright\Large $1{,}5\times10^{8}$\end{minipage}\end{tcolorbox}\hfill\begin{tcolorbox}[enhanced,nobeforeafter,width=7.1cm,height=2.1cm,valign=center,colback=cari!6,colframe=cari,arc=2.5mm,boxrule=0.9pt,left=3mm]{\Large\bfseries\color{cari}b)}\quad\begin{minipage}[c]{5.3cm}\raggedright\Large $15\times10^{7}$\end{minipage}\end{tcolorbox}\par\vspace{2.5mm}
\begin{tcolorbox}[enhanced,nobeforeafter,width=7.1cm,height=2.1cm,valign=center,colback=cari!6,colframe=cari,arc=2.5mm,boxrule=0.9pt,left=3mm]{\Large\bfseries\color{cari}c)}\quad\begin{minipage}[c]{5.3cm}\raggedright\Large $1{,}5\times10^{7}$\end{minipage}\end{tcolorbox}\hfill\begin{tcolorbox}[enhanced,nobeforeafter,width=7.1cm,height=2.1cm,valign=center,colback=cari!6,colframe=cari,arc=2.5mm,boxrule=0.9pt,left=3mm]{\Large\bfseries\color{cari}d)}\quad\begin{minipage}[c]{5.3cm}\raggedright\Large $1{,}5\times10^{12}$\end{minipage}\end{tcolorbox}
\end{tcolorbox}
\newpage
\begin{tcolorbox}[enhanced,width=16cm,height=14cm,colback=white,colframe=cari,boxrule=2pt,arc=0mm,
left=6mm,right=6mm,top=0mm,bottom=4mm,sharp corners,
overlay={\fill[cari] (frame.north west) rectangle ([yshift=-1.25cm]frame.north east);
\node[anchor=west,text=white,font=\sffamily\bfseries\LARGE] at ([xshift=5mm,yshift=-0.62cm]frame.north west) {Question 14 B};
\node[anchor=east,text=white,font=\sffamily\itshape\large] at ([xshift=-5mm,yshift=-0.62cm]frame.north east) {Arithmétique et puissances};}]
\vspace*{1.75cm}
\begin{minipage}[t][6.9cm][c]{\linewidth}\LARGE\centering
\begin{minipage}{0.95\linewidth}\raggedright\linespread{1.25}\selectfont Que vaut $(1{,}5\times10^{7})\times(4\times10^{-2})$, en notation scientifique~?\end{minipage}
\end{minipage}

\vfill
\begin{tcolorbox}[enhanced,nobeforeafter,width=7.1cm,height=2.1cm,valign=center,colback=cari!6,colframe=cari,arc=2.5mm,boxrule=0.9pt,left=3mm]{\Large\bfseries\color{cari}a)}\quad\begin{minipage}[c]{5.3cm}\raggedright\Large $60\times10^{4}$\end{minipage}\end{tcolorbox}\hfill\begin{tcolorbox}[enhanced,nobeforeafter,width=7.1cm,height=2.1cm,valign=center,colback=cari!6,colframe=cari,arc=2.5mm,boxrule=0.9pt,left=3mm]{\Large\bfseries\color{cari}b)}\quad\begin{minipage}[c]{5.3cm}\raggedright\Large $6\times10^{-14}$\end{minipage}\end{tcolorbox}\par\vspace{2.5mm}
\begin{tcolorbox}[enhanced,nobeforeafter,width=7.1cm,height=2.1cm,valign=center,colback=cari!6,colframe=cari,arc=2.5mm,boxrule=0.9pt,left=3mm]{\Large\bfseries\color{cari}c)}\quad\begin{minipage}[c]{5.3cm}\raggedright\Large $6\times10^{5}$\end{minipage}\end{tcolorbox}\hfill\begin{tcolorbox}[enhanced,nobeforeafter,width=7.1cm,height=2.1cm,valign=center,colback=cari!6,colframe=cari,arc=2.5mm,boxrule=0.9pt,left=3mm]{\Large\bfseries\color{cari}d)}\quad\begin{minipage}[c]{5.3cm}\raggedright\Large $6\times10^{4}$\end{minipage}\end{tcolorbox}
\end{tcolorbox}
\newpage
\begin{tcolorbox}[enhanced,width=16cm,height=14cm,colback=white,colframe=cari,boxrule=2pt,arc=0mm,
left=6mm,right=6mm,top=0mm,bottom=4mm,sharp corners,
overlay={\fill[cari] (frame.north west) rectangle ([yshift=-1.25cm]frame.north east);
\node[anchor=west,text=white,font=\sffamily\bfseries\LARGE] at ([xshift=5mm,yshift=-0.62cm]frame.north west) {Question 15 B};
\node[anchor=east,text=white,font=\sffamily\itshape\large] at ([xshift=-5mm,yshift=-0.62cm]frame.north east) {Arithmétique et puissances};}]
\vspace*{1.75cm}
\begin{minipage}[t][6.9cm][c]{\linewidth}\LARGE\centering
\begin{minipage}{0.95\linewidth}\raggedright\linespread{1.25}\selectfont Quelle est l'écriture décimale de $10^{-3}$~?\end{minipage}
\end{minipage}

\vfill
\begin{tcolorbox}[enhanced,nobeforeafter,width=7.1cm,height=2.1cm,valign=center,colback=cari!6,colframe=cari,arc=2.5mm,boxrule=0.9pt,left=3mm]{\Large\bfseries\color{cari}a)}\quad\begin{minipage}[c]{5.3cm}\raggedright\Large $-1\,000$\end{minipage}\end{tcolorbox}\hfill\begin{tcolorbox}[enhanced,nobeforeafter,width=7.1cm,height=2.1cm,valign=center,colback=cari!6,colframe=cari,arc=2.5mm,boxrule=0.9pt,left=3mm]{\Large\bfseries\color{cari}b)}\quad\begin{minipage}[c]{5.3cm}\raggedright\Large $-30$\end{minipage}\end{tcolorbox}\par\vspace{2.5mm}
\begin{tcolorbox}[enhanced,nobeforeafter,width=7.1cm,height=2.1cm,valign=center,colback=cari!6,colframe=cari,arc=2.5mm,boxrule=0.9pt,left=3mm]{\Large\bfseries\color{cari}c)}\quad\begin{minipage}[c]{5.3cm}\raggedright\Large $\dfrac{1}{30}$\end{minipage}\end{tcolorbox}\hfill\begin{tcolorbox}[enhanced,nobeforeafter,width=7.1cm,height=2.1cm,valign=center,colback=cari!6,colframe=cari,arc=2.5mm,boxrule=0.9pt,left=3mm]{\Large\bfseries\color{cari}d)}\quad\begin{minipage}[c]{5.3cm}\raggedright\Large $0{,}001$\end{minipage}\end{tcolorbox}
\end{tcolorbox}
\newpage
\begin{tcolorbox}[enhanced,width=16cm,height=14cm,colback=white,colframe=cari,boxrule=2pt,arc=0mm,
left=6mm,right=6mm,top=0mm,bottom=4mm,sharp corners,
overlay={\fill[cari] (frame.north west) rectangle ([yshift=-1.25cm]frame.north east);
\node[anchor=west,text=white,font=\sffamily\bfseries\LARGE] at ([xshift=5mm,yshift=-0.62cm]frame.north west) {Question 16 B};
\node[anchor=east,text=white,font=\sffamily\itshape\large] at ([xshift=-5mm,yshift=-0.62cm]frame.north east) {Arithmétique et puissances};}]
\vspace*{1.75cm}
\begin{minipage}[t][6.9cm][c]{\linewidth}\LARGE\centering
\begin{minipage}{0.95\linewidth}\raggedright\linespread{1.25}\selectfont Quelle est l'écriture décimale de $4^{-2}$~?\end{minipage}
\end{minipage}

\vfill
\begin{tcolorbox}[enhanced,nobeforeafter,width=7.1cm,height=2.1cm,valign=center,colback=cari!6,colframe=cari,arc=2.5mm,boxrule=0.9pt,left=3mm]{\Large\bfseries\color{cari}a)}\quad\begin{minipage}[c]{5.3cm}\raggedright\Large $0{,}0625$\end{minipage}\end{tcolorbox}\hfill\begin{tcolorbox}[enhanced,nobeforeafter,width=7.1cm,height=2.1cm,valign=center,colback=cari!6,colframe=cari,arc=2.5mm,boxrule=0.9pt,left=3mm]{\Large\bfseries\color{cari}b)}\quad\begin{minipage}[c]{5.3cm}\raggedright\Large $0{,}125$\end{minipage}\end{tcolorbox}\par\vspace{2.5mm}
\begin{tcolorbox}[enhanced,nobeforeafter,width=7.1cm,height=2.1cm,valign=center,colback=cari!6,colframe=cari,arc=2.5mm,boxrule=0.9pt,left=3mm]{\Large\bfseries\color{cari}c)}\quad\begin{minipage}[c]{5.3cm}\raggedright\Large $-8$\end{minipage}\end{tcolorbox}\hfill\begin{tcolorbox}[enhanced,nobeforeafter,width=7.1cm,height=2.1cm,valign=center,colback=cari!6,colframe=cari,arc=2.5mm,boxrule=0.9pt,left=3mm]{\Large\bfseries\color{cari}d)}\quad\begin{minipage}[c]{5.3cm}\raggedright\Large $-16$\end{minipage}\end{tcolorbox}
\end{tcolorbox}
\newpage
\begin{tcolorbox}[enhanced,width=16cm,height=14cm,colback=white,colframe=clit,boxrule=2pt,arc=0mm,
left=6mm,right=6mm,top=0mm,bottom=4mm,sharp corners,
overlay={\fill[clit] (frame.north west) rectangle ([yshift=-1.25cm]frame.north east);
\node[anchor=west,text=white,font=\sffamily\bfseries\LARGE] at ([xshift=5mm,yshift=-0.62cm]frame.north west) {Question 17 B};
\node[anchor=east,text=white,font=\sffamily\itshape\large] at ([xshift=-5mm,yshift=-0.62cm]frame.north east) {Calcul littéral et équations};}]
\vspace*{1.75cm}
\begin{minipage}[t][6.9cm][c]{\linewidth}\LARGE\centering
\begin{minipage}{0.95\linewidth}\raggedright\linespread{1.25}\selectfont Développer et réduire $(x+6)(4x-2)$.\end{minipage}
\end{minipage}

\vfill
\begin{tcolorbox}[enhanced,nobeforeafter,width=7.1cm,height=2.1cm,valign=center,colback=clit!6,colframe=clit,arc=2.5mm,boxrule=0.9pt,left=3mm]{\Large\bfseries\color{clit}a)}\quad\begin{minipage}[c]{5.3cm}\raggedright\Large $5x^2+22x-12$\end{minipage}\end{tcolorbox}\hfill\begin{tcolorbox}[enhanced,nobeforeafter,width=7.1cm,height=2.1cm,valign=center,colback=clit!6,colframe=clit,arc=2.5mm,boxrule=0.9pt,left=3mm]{\Large\bfseries\color{clit}b)}\quad\begin{minipage}[c]{5.3cm}\raggedright\Large $4x^2-26x-12$\end{minipage}\end{tcolorbox}\par\vspace{2.5mm}
\begin{tcolorbox}[enhanced,nobeforeafter,width=7.1cm,height=2.1cm,valign=center,colback=clit!6,colframe=clit,arc=2.5mm,boxrule=0.9pt,left=3mm]{\Large\bfseries\color{clit}c)}\quad\begin{minipage}[c]{5.3cm}\raggedright\Large $4x^2+22x-12$\end{minipage}\end{tcolorbox}\hfill\begin{tcolorbox}[enhanced,nobeforeafter,width=7.1cm,height=2.1cm,valign=center,colback=clit!6,colframe=clit,arc=2.5mm,boxrule=0.9pt,left=3mm]{\Large\bfseries\color{clit}d)}\quad\begin{minipage}[c]{5.3cm}\raggedright\Large $4x^2-12$\end{minipage}\end{tcolorbox}
\end{tcolorbox}
\newpage
\begin{tcolorbox}[enhanced,width=16cm,height=14cm,colback=white,colframe=clit,boxrule=2pt,arc=0mm,
left=6mm,right=6mm,top=0mm,bottom=4mm,sharp corners,
overlay={\fill[clit] (frame.north west) rectangle ([yshift=-1.25cm]frame.north east);
\node[anchor=west,text=white,font=\sffamily\bfseries\LARGE] at ([xshift=5mm,yshift=-0.62cm]frame.north west) {Question 18 B};
\node[anchor=east,text=white,font=\sffamily\itshape\large] at ([xshift=-5mm,yshift=-0.62cm]frame.north east) {Calcul littéral et équations};}]
\vspace*{1.75cm}
\begin{minipage}[t][6.9cm][c]{\linewidth}\LARGE\centering
\begin{minipage}{0.95\linewidth}\raggedright\linespread{1.25}\selectfont Développer et réduire $(5x-2)(x-3)$.\end{minipage}
\end{minipage}

\vfill
\begin{tcolorbox}[enhanced,nobeforeafter,width=7.1cm,height=2.1cm,valign=center,colback=clit!6,colframe=clit,arc=2.5mm,boxrule=0.9pt,left=3mm]{\Large\bfseries\color{clit}a)}\quad\begin{minipage}[c]{5.3cm}\raggedright\Large $6x^2-17x+6$\end{minipage}\end{tcolorbox}\hfill\begin{tcolorbox}[enhanced,nobeforeafter,width=7.1cm,height=2.1cm,valign=center,colback=clit!6,colframe=clit,arc=2.5mm,boxrule=0.9pt,left=3mm]{\Large\bfseries\color{clit}b)}\quad\begin{minipage}[c]{5.3cm}\raggedright\Large $5x^2+6$\end{minipage}\end{tcolorbox}\par\vspace{2.5mm}
\begin{tcolorbox}[enhanced,nobeforeafter,width=7.1cm,height=2.1cm,valign=center,colback=clit!6,colframe=clit,arc=2.5mm,boxrule=0.9pt,left=3mm]{\Large\bfseries\color{clit}c)}\quad\begin{minipage}[c]{5.3cm}\raggedright\Large $5x^2-17x+6$\end{minipage}\end{tcolorbox}\hfill\begin{tcolorbox}[enhanced,nobeforeafter,width=7.1cm,height=2.1cm,valign=center,colback=clit!6,colframe=clit,arc=2.5mm,boxrule=0.9pt,left=3mm]{\Large\bfseries\color{clit}d)}\quad\begin{minipage}[c]{5.3cm}\raggedright\Large $5x^2-13x+6$\end{minipage}\end{tcolorbox}
\end{tcolorbox}
\newpage
\begin{tcolorbox}[enhanced,width=16cm,height=14cm,colback=white,colframe=clit,boxrule=2pt,arc=0mm,
left=6mm,right=6mm,top=0mm,bottom=4mm,sharp corners,
overlay={\fill[clit] (frame.north west) rectangle ([yshift=-1.25cm]frame.north east);
\node[anchor=west,text=white,font=\sffamily\bfseries\LARGE] at ([xshift=5mm,yshift=-0.62cm]frame.north west) {Question 19 B};
\node[anchor=east,text=white,font=\sffamily\itshape\large] at ([xshift=-5mm,yshift=-0.62cm]frame.north east) {Calcul littéral et équations};}]
\vspace*{1.75cm}
\begin{minipage}[t][6.9cm][c]{\linewidth}\LARGE\centering
\begin{minipage}{0.95\linewidth}\raggedright\linespread{1.25}\selectfont Factoriser $x^2-100$.\end{minipage}
\end{minipage}

\vfill
\begin{tcolorbox}[enhanced,nobeforeafter,width=7.1cm,height=2.1cm,valign=center,colback=clit!6,colframe=clit,arc=2.5mm,boxrule=0.9pt,left=3mm]{\Large\bfseries\color{clit}a)}\quad\begin{minipage}[c]{5.3cm}\raggedright\Large $(x-10)^2$\end{minipage}\end{tcolorbox}\hfill\begin{tcolorbox}[enhanced,nobeforeafter,width=7.1cm,height=2.1cm,valign=center,colback=clit!6,colframe=clit,arc=2.5mm,boxrule=0.9pt,left=3mm]{\Large\bfseries\color{clit}b)}\quad\begin{minipage}[c]{5.3cm}\raggedright\Large $(x-100)(x+100)$\end{minipage}\end{tcolorbox}\par\vspace{2.5mm}
\begin{tcolorbox}[enhanced,nobeforeafter,width=7.1cm,height=2.1cm,valign=center,colback=clit!6,colframe=clit,arc=2.5mm,boxrule=0.9pt,left=3mm]{\Large\bfseries\color{clit}c)}\quad\begin{minipage}[c]{5.3cm}\raggedright\Large $(x+10)^2$\end{minipage}\end{tcolorbox}\hfill\begin{tcolorbox}[enhanced,nobeforeafter,width=7.1cm,height=2.1cm,valign=center,colback=clit!6,colframe=clit,arc=2.5mm,boxrule=0.9pt,left=3mm]{\Large\bfseries\color{clit}d)}\quad\begin{minipage}[c]{5.3cm}\raggedright\Large $(x-10)(x+10)$\end{minipage}\end{tcolorbox}
\end{tcolorbox}
\newpage
\begin{tcolorbox}[enhanced,width=16cm,height=14cm,colback=white,colframe=clit,boxrule=2pt,arc=0mm,
left=6mm,right=6mm,top=0mm,bottom=4mm,sharp corners,
overlay={\fill[clit] (frame.north west) rectangle ([yshift=-1.25cm]frame.north east);
\node[anchor=west,text=white,font=\sffamily\bfseries\LARGE] at ([xshift=5mm,yshift=-0.62cm]frame.north west) {Question 20 B};
\node[anchor=east,text=white,font=\sffamily\itshape\large] at ([xshift=-5mm,yshift=-0.62cm]frame.north east) {Calcul littéral et équations};}]
\vspace*{1.75cm}
\begin{minipage}[t][6.9cm][c]{\linewidth}\LARGE\centering
\begin{minipage}{0.95\linewidth}\raggedright\linespread{1.25}\selectfont Factoriser $4x^2-49$.\end{minipage}
\end{minipage}

\vfill
\begin{tcolorbox}[enhanced,nobeforeafter,width=7.1cm,height=2.1cm,valign=center,colback=clit!6,colframe=clit,arc=2.5mm,boxrule=0.9pt,left=3mm]{\Large\bfseries\color{clit}a)}\quad\begin{minipage}[c]{5.3cm}\raggedright\Large $(2x-7)^2$\end{minipage}\end{tcolorbox}\hfill\begin{tcolorbox}[enhanced,nobeforeafter,width=7.1cm,height=2.1cm,valign=center,colback=clit!6,colframe=clit,arc=2.5mm,boxrule=0.9pt,left=3mm]{\Large\bfseries\color{clit}b)}\quad\begin{minipage}[c]{5.3cm}\raggedright\Large $(2x-49)(2x+49)$\end{minipage}\end{tcolorbox}\par\vspace{2.5mm}
\begin{tcolorbox}[enhanced,nobeforeafter,width=7.1cm,height=2.1cm,valign=center,colback=clit!6,colframe=clit,arc=2.5mm,boxrule=0.9pt,left=3mm]{\Large\bfseries\color{clit}c)}\quad\begin{minipage}[c]{5.3cm}\raggedright\Large $(2x+7)^2$\end{minipage}\end{tcolorbox}\hfill\begin{tcolorbox}[enhanced,nobeforeafter,width=7.1cm,height=2.1cm,valign=center,colback=clit!6,colframe=clit,arc=2.5mm,boxrule=0.9pt,left=3mm]{\Large\bfseries\color{clit}d)}\quad\begin{minipage}[c]{5.3cm}\raggedright\Large $(2x-7)(2x+7)$\end{minipage}\end{tcolorbox}
\end{tcolorbox}
\newpage
\begin{tcolorbox}[enhanced,width=16cm,height=14cm,colback=white,colframe=clit,boxrule=2pt,arc=0mm,
left=6mm,right=6mm,top=0mm,bottom=4mm,sharp corners,
overlay={\fill[clit] (frame.north west) rectangle ([yshift=-1.25cm]frame.north east);
\node[anchor=west,text=white,font=\sffamily\bfseries\LARGE] at ([xshift=5mm,yshift=-0.62cm]frame.north west) {Question 21 B};
\node[anchor=east,text=white,font=\sffamily\itshape\large] at ([xshift=-5mm,yshift=-0.62cm]frame.north east) {Calcul littéral et équations};}]
\vspace*{1.75cm}
\begin{minipage}[t][6.9cm][c]{\linewidth}\LARGE\centering
\begin{minipage}{0.95\linewidth}\raggedright\linespread{1.25}\selectfont Quelles sont les solutions de l'équation $(3x-6)(x+2)=0$~?\end{minipage}
\end{minipage}

\vfill
\begin{tcolorbox}[enhanced,nobeforeafter,width=7.1cm,height=2.1cm,valign=center,colback=clit!6,colframe=clit,arc=2.5mm,boxrule=0.9pt,left=3mm]{\Large\bfseries\color{clit}a)}\quad\begin{minipage}[c]{5.3cm}\raggedright\Large $\{-2\,;\,2\}$\end{minipage}\end{tcolorbox}\hfill\begin{tcolorbox}[enhanced,nobeforeafter,width=7.1cm,height=2.1cm,valign=center,colback=clit!6,colframe=clit,arc=2.5mm,boxrule=0.9pt,left=3mm]{\Large\bfseries\color{clit}b)}\quad\begin{minipage}[c]{5.3cm}\raggedright\Large $\{-1\,;\,2\}$\end{minipage}\end{tcolorbox}\par\vspace{2.5mm}
\begin{tcolorbox}[enhanced,nobeforeafter,width=7.1cm,height=2.1cm,valign=center,colback=clit!6,colframe=clit,arc=2.5mm,boxrule=0.9pt,left=3mm]{\Large\bfseries\color{clit}c)}\quad\begin{minipage}[c]{5.3cm}\raggedright\Large $\{-6\,;\,2\}$\end{minipage}\end{tcolorbox}\hfill\begin{tcolorbox}[enhanced,nobeforeafter,width=7.1cm,height=2.1cm,valign=center,colback=clit!6,colframe=clit,arc=2.5mm,boxrule=0.9pt,left=3mm]{\Large\bfseries\color{clit}d)}\quad\begin{minipage}[c]{5.3cm}\raggedright\Large $\{2\}$\end{minipage}\end{tcolorbox}
\end{tcolorbox}
\newpage
\begin{tcolorbox}[enhanced,width=16cm,height=14cm,colback=white,colframe=clit,boxrule=2pt,arc=0mm,
left=6mm,right=6mm,top=0mm,bottom=4mm,sharp corners,
overlay={\fill[clit] (frame.north west) rectangle ([yshift=-1.25cm]frame.north east);
\node[anchor=west,text=white,font=\sffamily\bfseries\LARGE] at ([xshift=5mm,yshift=-0.62cm]frame.north west) {Question 22 B};
\node[anchor=east,text=white,font=\sffamily\itshape\large] at ([xshift=-5mm,yshift=-0.62cm]frame.north east) {Calcul littéral et équations};}]
\vspace*{1.75cm}
\begin{minipage}[t][6.9cm][c]{\linewidth}\LARGE\centering
\begin{minipage}{0.95\linewidth}\raggedright\linespread{1.25}\selectfont Quelles sont les solutions de l'équation $(2x+1)(x-5)=0$~?\end{minipage}
\end{minipage}

\vfill
\begin{tcolorbox}[enhanced,nobeforeafter,width=7.1cm,height=2.1cm,valign=center,colback=clit!6,colframe=clit,arc=2.5mm,boxrule=0.9pt,left=3mm]{\Large\bfseries\color{clit}a)}\quad\begin{minipage}[c]{5.3cm}\raggedright\Large $\{-\dfrac{1}{2}\,;\,5\}$\end{minipage}\end{tcolorbox}\hfill\begin{tcolorbox}[enhanced,nobeforeafter,width=7.1cm,height=2.1cm,valign=center,colback=clit!6,colframe=clit,arc=2.5mm,boxrule=0.9pt,left=3mm]{\Large\bfseries\color{clit}b)}\quad\begin{minipage}[c]{5.3cm}\raggedright\Large $\{-5\,;\,1\}$\end{minipage}\end{tcolorbox}\par\vspace{2.5mm}
\begin{tcolorbox}[enhanced,nobeforeafter,width=7.1cm,height=2.1cm,valign=center,colback=clit!6,colframe=clit,arc=2.5mm,boxrule=0.9pt,left=3mm]{\Large\bfseries\color{clit}c)}\quad\begin{minipage}[c]{5.3cm}\raggedright\Large $\{-\dfrac{1}{2}\}$\end{minipage}\end{tcolorbox}\hfill\begin{tcolorbox}[enhanced,nobeforeafter,width=7.1cm,height=2.1cm,valign=center,colback=clit!6,colframe=clit,arc=2.5mm,boxrule=0.9pt,left=3mm]{\Large\bfseries\color{clit}d)}\quad\begin{minipage}[c]{5.3cm}\raggedright\Large $\{-5\,;\,\dfrac{1}{2}\}$\end{minipage}\end{tcolorbox}
\end{tcolorbox}
\newpage
\begin{tcolorbox}[enhanced,width=16cm,height=14cm,colback=white,colframe=clit,boxrule=2pt,arc=0mm,
left=6mm,right=6mm,top=0mm,bottom=4mm,sharp corners,
overlay={\fill[clit] (frame.north west) rectangle ([yshift=-1.25cm]frame.north east);
\node[anchor=west,text=white,font=\sffamily\bfseries\LARGE] at ([xshift=5mm,yshift=-0.62cm]frame.north west) {Question 23 B};
\node[anchor=east,text=white,font=\sffamily\itshape\large] at ([xshift=-5mm,yshift=-0.62cm]frame.north east) {Calcul littéral et équations};}]
\vspace*{1.75cm}
\begin{minipage}[t][6.9cm][c]{\linewidth}\LARGE\centering
\begin{minipage}{0.95\linewidth}\raggedright\linespread{1.25}\selectfont Factoriser $(2x+1)(x+6)+(2x+1)(x-2)$.\end{minipage}
\end{minipage}

\vfill
\begin{tcolorbox}[enhanced,nobeforeafter,width=7.1cm,height=2.1cm,valign=center,colback=clit!6,colframe=clit,arc=2.5mm,boxrule=0.9pt,left=3mm]{\Large\bfseries\color{clit}a)}\quad\begin{minipage}[c]{5.3cm}\raggedright\Large $(2x+1)^2(2x+4)$\end{minipage}\end{tcolorbox}\hfill\begin{tcolorbox}[enhanced,nobeforeafter,width=7.1cm,height=2.1cm,valign=center,colback=clit!6,colframe=clit,arc=2.5mm,boxrule=0.9pt,left=3mm]{\Large\bfseries\color{clit}b)}\quad\begin{minipage}[c]{5.3cm}\raggedright\Large $(2x+1)(8)$\end{minipage}\end{tcolorbox}\par\vspace{2.5mm}
\begin{tcolorbox}[enhanced,nobeforeafter,width=7.1cm,height=2.1cm,valign=center,colback=clit!6,colframe=clit,arc=2.5mm,boxrule=0.9pt,left=3mm]{\Large\bfseries\color{clit}c)}\quad\begin{minipage}[c]{5.3cm}\raggedright\Large $(2x+1)(x-12)$\end{minipage}\end{tcolorbox}\hfill\begin{tcolorbox}[enhanced,nobeforeafter,width=7.1cm,height=2.1cm,valign=center,colback=clit!6,colframe=clit,arc=2.5mm,boxrule=0.9pt,left=3mm]{\Large\bfseries\color{clit}d)}\quad\begin{minipage}[c]{5.3cm}\raggedright\Large $(2x+1)(2x+4)$\end{minipage}\end{tcolorbox}
\end{tcolorbox}
\newpage
\begin{tcolorbox}[enhanced,width=16cm,height=14cm,colback=white,colframe=clit,boxrule=2pt,arc=0mm,
left=6mm,right=6mm,top=0mm,bottom=4mm,sharp corners,
overlay={\fill[clit] (frame.north west) rectangle ([yshift=-1.25cm]frame.north east);
\node[anchor=west,text=white,font=\sffamily\bfseries\LARGE] at ([xshift=5mm,yshift=-0.62cm]frame.north west) {Question 24 B};
\node[anchor=east,text=white,font=\sffamily\itshape\large] at ([xshift=-5mm,yshift=-0.62cm]frame.north east) {Calcul littéral et équations};}]
\vspace*{1.75cm}
\begin{minipage}[t][6.9cm][c]{\linewidth}\LARGE\centering
\begin{minipage}{0.95\linewidth}\raggedright\linespread{1.25}\selectfont Factoriser $(x+3)(4x-1)+(x+3)(2x+2)$.\end{minipage}
\end{minipage}

\vfill
\begin{tcolorbox}[enhanced,nobeforeafter,width=7.1cm,height=2.1cm,valign=center,colback=clit!6,colframe=clit,arc=2.5mm,boxrule=0.9pt,left=3mm]{\Large\bfseries\color{clit}a)}\quad\begin{minipage}[c]{5.3cm}\raggedright\Large $(x+3)(2x-3)$\end{minipage}\end{tcolorbox}\hfill\begin{tcolorbox}[enhanced,nobeforeafter,width=7.1cm,height=2.1cm,valign=center,colback=clit!6,colframe=clit,arc=2.5mm,boxrule=0.9pt,left=3mm]{\Large\bfseries\color{clit}b)}\quad\begin{minipage}[c]{5.3cm}\raggedright\Large $(x+3)(8x-2)$\end{minipage}\end{tcolorbox}\par\vspace{2.5mm}
\begin{tcolorbox}[enhanced,nobeforeafter,width=7.1cm,height=2.1cm,valign=center,colback=clit!6,colframe=clit,arc=2.5mm,boxrule=0.9pt,left=3mm]{\Large\bfseries\color{clit}c)}\quad\begin{minipage}[c]{5.3cm}\raggedright\Large $(x+3)^2(6x+1)$\end{minipage}\end{tcolorbox}\hfill\begin{tcolorbox}[enhanced,nobeforeafter,width=7.1cm,height=2.1cm,valign=center,colback=clit!6,colframe=clit,arc=2.5mm,boxrule=0.9pt,left=3mm]{\Large\bfseries\color{clit}d)}\quad\begin{minipage}[c]{5.3cm}\raggedright\Large $(x+3)(6x+1)$\end{minipage}\end{tcolorbox}
\end{tcolorbox}
\newpage
\begin{tcolorbox}[enhanced,width=16cm,height=14cm,colback=white,colframe=clit,boxrule=2pt,arc=0mm,
left=6mm,right=6mm,top=0mm,bottom=4mm,sharp corners,
overlay={\fill[clit] (frame.north west) rectangle ([yshift=-1.25cm]frame.north east);
\node[anchor=west,text=white,font=\sffamily\bfseries\LARGE] at ([xshift=5mm,yshift=-0.62cm]frame.north west) {Question 25 B};
\node[anchor=east,text=white,font=\sffamily\itshape\large] at ([xshift=-5mm,yshift=-0.62cm]frame.north east) {Calcul littéral et équations};}]
\vspace*{1.75cm}
\begin{minipage}[t][6.9cm][c]{\linewidth}\LARGE\centering
\begin{minipage}{0.95\linewidth}\raggedright\linespread{1.25}\selectfont Quelles sont les solutions de l'inéquation $5x+1>-9$~?\end{minipage}
\end{minipage}

\vfill
\begin{tcolorbox}[enhanced,nobeforeafter,width=7.1cm,height=2.1cm,valign=center,colback=clit!6,colframe=clit,arc=2.5mm,boxrule=0.9pt,left=3mm]{\Large\bfseries\color{clit}a)}\quad\begin{minipage}[c]{5.3cm}\raggedright\Large $x>-2$\end{minipage}\end{tcolorbox}\hfill\begin{tcolorbox}[enhanced,nobeforeafter,width=7.1cm,height=2.1cm,valign=center,colback=clit!6,colframe=clit,arc=2.5mm,boxrule=0.9pt,left=3mm]{\Large\bfseries\color{clit}b)}\quad\begin{minipage}[c]{5.3cm}\raggedright\Large $x<-2$\end{minipage}\end{tcolorbox}\par\vspace{2.5mm}
\begin{tcolorbox}[enhanced,nobeforeafter,width=7.1cm,height=2.1cm,valign=center,colback=clit!6,colframe=clit,arc=2.5mm,boxrule=0.9pt,left=3mm]{\Large\bfseries\color{clit}c)}\quad\begin{minipage}[c]{5.3cm}\raggedright\Large $x>2$\end{minipage}\end{tcolorbox}\hfill\begin{tcolorbox}[enhanced,nobeforeafter,width=7.1cm,height=2.1cm,valign=center,colback=clit!6,colframe=clit,arc=2.5mm,boxrule=0.9pt,left=3mm]{\Large\bfseries\color{clit}d)}\quad\begin{minipage}[c]{5.3cm}\raggedright\Large $x>-\dfrac{8}{5}$\end{minipage}\end{tcolorbox}
\end{tcolorbox}
\newpage
\begin{tcolorbox}[enhanced,width=16cm,height=14cm,colback=white,colframe=clit,boxrule=2pt,arc=0mm,
left=6mm,right=6mm,top=0mm,bottom=4mm,sharp corners,
overlay={\fill[clit] (frame.north west) rectangle ([yshift=-1.25cm]frame.north east);
\node[anchor=west,text=white,font=\sffamily\bfseries\LARGE] at ([xshift=5mm,yshift=-0.62cm]frame.north west) {Question 26 B};
\node[anchor=east,text=white,font=\sffamily\itshape\large] at ([xshift=-5mm,yshift=-0.62cm]frame.north east) {Calcul littéral et équations};}]
\vspace*{1.75cm}
\begin{minipage}[t][6.9cm][c]{\linewidth}\LARGE\centering
\begin{minipage}{0.95\linewidth}\raggedright\linespread{1.25}\selectfont Quelles sont les solutions de l'inéquation $-4x-3>9$~?\end{minipage}
\end{minipage}

\vfill
\begin{tcolorbox}[enhanced,nobeforeafter,width=7.1cm,height=2.1cm,valign=center,colback=clit!6,colframe=clit,arc=2.5mm,boxrule=0.9pt,left=3mm]{\Large\bfseries\color{clit}a)}\quad\begin{minipage}[c]{5.3cm}\raggedright\Large $x>-3$\end{minipage}\end{tcolorbox}\hfill\begin{tcolorbox}[enhanced,nobeforeafter,width=7.1cm,height=2.1cm,valign=center,colback=clit!6,colframe=clit,arc=2.5mm,boxrule=0.9pt,left=3mm]{\Large\bfseries\color{clit}b)}\quad\begin{minipage}[c]{5.3cm}\raggedright\Large $x<-\dfrac{3}{2}$\end{minipage}\end{tcolorbox}\par\vspace{2.5mm}
\begin{tcolorbox}[enhanced,nobeforeafter,width=7.1cm,height=2.1cm,valign=center,colback=clit!6,colframe=clit,arc=2.5mm,boxrule=0.9pt,left=3mm]{\Large\bfseries\color{clit}c)}\quad\begin{minipage}[c]{5.3cm}\raggedright\Large $x<3$\end{minipage}\end{tcolorbox}\hfill\begin{tcolorbox}[enhanced,nobeforeafter,width=7.1cm,height=2.1cm,valign=center,colback=clit!6,colframe=clit,arc=2.5mm,boxrule=0.9pt,left=3mm]{\Large\bfseries\color{clit}d)}\quad\begin{minipage}[c]{5.3cm}\raggedright\Large $x<-3$\end{minipage}\end{tcolorbox}
\end{tcolorbox}
\newpage
\begin{tcolorbox}[enhanced,width=16cm,height=14cm,colback=white,colframe=clit,boxrule=2pt,arc=0mm,
left=6mm,right=6mm,top=0mm,bottom=4mm,sharp corners,
overlay={\fill[clit] (frame.north west) rectangle ([yshift=-1.25cm]frame.north east);
\node[anchor=west,text=white,font=\sffamily\bfseries\LARGE] at ([xshift=5mm,yshift=-0.62cm]frame.north west) {Question 27 B};
\node[anchor=east,text=white,font=\sffamily\itshape\large] at ([xshift=-5mm,yshift=-0.62cm]frame.north east) {Calcul littéral et équations};}]
\vspace*{1.75cm}
\begin{minipage}[t][6.9cm][c]{\linewidth}\LARGE\centering
\begin{minipage}{0.95\linewidth}\raggedright\linespread{1.25}\selectfont Quelles sont les solutions de l'équation $x^2=-9$~?\end{minipage}
\end{minipage}

\vfill
\begin{tcolorbox}[enhanced,nobeforeafter,width=7.1cm,height=2.1cm,valign=center,colback=clit!6,colframe=clit,arc=2.5mm,boxrule=0.9pt,left=3mm]{\Large\bfseries\color{clit}a)}\quad\begin{minipage}[c]{5.3cm}\raggedright\Large $\{-9\}$\end{minipage}\end{tcolorbox}\hfill\begin{tcolorbox}[enhanced,nobeforeafter,width=7.1cm,height=2.1cm,valign=center,colback=clit!6,colframe=clit,arc=2.5mm,boxrule=0.9pt,left=3mm]{\Large\bfseries\color{clit}b)}\quad\begin{minipage}[c]{5.3cm}\raggedright\Large aucune solution\end{minipage}\end{tcolorbox}\par\vspace{2.5mm}
\begin{tcolorbox}[enhanced,nobeforeafter,width=7.1cm,height=2.1cm,valign=center,colback=clit!6,colframe=clit,arc=2.5mm,boxrule=0.9pt,left=3mm]{\Large\bfseries\color{clit}c)}\quad\begin{minipage}[c]{5.3cm}\raggedright\Large $\{-3\,;\,3\}$\end{minipage}\end{tcolorbox}\hfill\begin{tcolorbox}[enhanced,nobeforeafter,width=7.1cm,height=2.1cm,valign=center,colback=clit!6,colframe=clit,arc=2.5mm,boxrule=0.9pt,left=3mm]{\Large\bfseries\color{clit}d)}\quad\begin{minipage}[c]{5.3cm}\raggedright\Large $\{-3\}$\end{minipage}\end{tcolorbox}
\end{tcolorbox}
\newpage
\begin{tcolorbox}[enhanced,width=16cm,height=14cm,colback=white,colframe=clit,boxrule=2pt,arc=0mm,
left=6mm,right=6mm,top=0mm,bottom=4mm,sharp corners,
overlay={\fill[clit] (frame.north west) rectangle ([yshift=-1.25cm]frame.north east);
\node[anchor=west,text=white,font=\sffamily\bfseries\LARGE] at ([xshift=5mm,yshift=-0.62cm]frame.north west) {Question 28 B};
\node[anchor=east,text=white,font=\sffamily\itshape\large] at ([xshift=-5mm,yshift=-0.62cm]frame.north east) {Calcul littéral et équations};}]
\vspace*{1.75cm}
\begin{minipage}[t][6.9cm][c]{\linewidth}\LARGE\centering
\begin{minipage}{0.95\linewidth}\raggedright\linespread{1.25}\selectfont Quelles sont les solutions de l'équation $x^2=81$~?\end{minipage}
\end{minipage}

\vfill
\begin{tcolorbox}[enhanced,nobeforeafter,width=7.1cm,height=2.1cm,valign=center,colback=clit!6,colframe=clit,arc=2.5mm,boxrule=0.9pt,left=3mm]{\Large\bfseries\color{clit}a)}\quad\begin{minipage}[c]{5.3cm}\raggedright\Large aucune solution\end{minipage}\end{tcolorbox}\hfill\begin{tcolorbox}[enhanced,nobeforeafter,width=7.1cm,height=2.1cm,valign=center,colback=clit!6,colframe=clit,arc=2.5mm,boxrule=0.9pt,left=3mm]{\Large\bfseries\color{clit}b)}\quad\begin{minipage}[c]{5.3cm}\raggedright\Large $\{-40{,}5\,;\,40{,}5\}$\end{minipage}\end{tcolorbox}\par\vspace{2.5mm}
\begin{tcolorbox}[enhanced,nobeforeafter,width=7.1cm,height=2.1cm,valign=center,colback=clit!6,colframe=clit,arc=2.5mm,boxrule=0.9pt,left=3mm]{\Large\bfseries\color{clit}c)}\quad\begin{minipage}[c]{5.3cm}\raggedright\Large $\{9\}$\end{minipage}\end{tcolorbox}\hfill\begin{tcolorbox}[enhanced,nobeforeafter,width=7.1cm,height=2.1cm,valign=center,colback=clit!6,colframe=clit,arc=2.5mm,boxrule=0.9pt,left=3mm]{\Large\bfseries\color{clit}d)}\quad\begin{minipage}[c]{5.3cm}\raggedright\Large $\{-9\,;\,9\}$\end{minipage}\end{tcolorbox}
\end{tcolorbox}
\newpage
\begin{tcolorbox}[enhanced,width=16cm,height=14cm,colback=white,colframe=clit,boxrule=2pt,arc=0mm,
left=6mm,right=6mm,top=0mm,bottom=4mm,sharp corners,
overlay={\fill[clit] (frame.north west) rectangle ([yshift=-1.25cm]frame.north east);
\node[anchor=west,text=white,font=\sffamily\bfseries\LARGE] at ([xshift=5mm,yshift=-0.62cm]frame.north west) {Question 29 B};
\node[anchor=east,text=white,font=\sffamily\itshape\large] at ([xshift=-5mm,yshift=-0.62cm]frame.north east) {Calcul littéral et équations};}]
\vspace*{1.75cm}
\begin{minipage}[t][6.9cm][c]{\linewidth}\LARGE\centering
\begin{minipage}{0.95\linewidth}\raggedright\linespread{1.25}\selectfont Le nombre $1$ est-il solution de l'équation $2x^2+x-4=0$~?\end{minipage}
\end{minipage}

\vfill
\begin{tcolorbox}[enhanced,nobeforeafter,width=7.1cm,height=2.1cm,valign=center,colback=clit!6,colframe=clit,arc=2.5mm,boxrule=0.9pt,left=3mm]{\Large\bfseries\color{clit}a)}\quad\begin{minipage}[c]{5.3cm}\raggedright\Large non\end{minipage}\end{tcolorbox}\hfill\begin{tcolorbox}[enhanced,nobeforeafter,width=7.1cm,height=2.1cm,valign=center,colback=clit!6,colframe=clit,arc=2.5mm,boxrule=0.9pt,left=3mm]{\Large\bfseries\color{clit}b)}\quad\begin{minipage}[c]{5.3cm}\raggedright\Large on ne peut pas savoir\end{minipage}\end{tcolorbox}\par\vspace{2.5mm}
\begin{tcolorbox}[enhanced,nobeforeafter,width=7.1cm,height=2.1cm,valign=center,colback=clit!6,colframe=clit,arc=2.5mm,boxrule=0.9pt,left=3mm]{\Large\bfseries\color{clit}c)}\quad\begin{minipage}[c]{5.3cm}\raggedright\large seulement si $x$ est positif\end{minipage}\end{tcolorbox}\hfill\begin{tcolorbox}[enhanced,nobeforeafter,width=7.1cm,height=2.1cm,valign=center,colback=clit!6,colframe=clit,arc=2.5mm,boxrule=0.9pt,left=3mm]{\Large\bfseries\color{clit}d)}\quad\begin{minipage}[c]{5.3cm}\raggedright\Large oui\end{minipage}\end{tcolorbox}
\end{tcolorbox}
\newpage
\begin{tcolorbox}[enhanced,width=16cm,height=14cm,colback=white,colframe=clit,boxrule=2pt,arc=0mm,
left=6mm,right=6mm,top=0mm,bottom=4mm,sharp corners,
overlay={\fill[clit] (frame.north west) rectangle ([yshift=-1.25cm]frame.north east);
\node[anchor=west,text=white,font=\sffamily\bfseries\LARGE] at ([xshift=5mm,yshift=-0.62cm]frame.north west) {Question 30 B};
\node[anchor=east,text=white,font=\sffamily\itshape\large] at ([xshift=-5mm,yshift=-0.62cm]frame.north east) {Calcul littéral et équations};}]
\vspace*{1.75cm}
\begin{minipage}[t][6.9cm][c]{\linewidth}\LARGE\centering
\begin{minipage}{0.95\linewidth}\raggedright\linespread{1.25}\selectfont Le nombre $-1$ est-il solution de l'équation $x^2-4x-5=0$~?\end{minipage}
\end{minipage}

\vfill
\begin{tcolorbox}[enhanced,nobeforeafter,width=7.1cm,height=2.1cm,valign=center,colback=clit!6,colframe=clit,arc=2.5mm,boxrule=0.9pt,left=3mm]{\Large\bfseries\color{clit}a)}\quad\begin{minipage}[c]{5.3cm}\raggedright\Large oui\end{minipage}\end{tcolorbox}\hfill\begin{tcolorbox}[enhanced,nobeforeafter,width=7.1cm,height=2.1cm,valign=center,colback=clit!6,colframe=clit,arc=2.5mm,boxrule=0.9pt,left=3mm]{\Large\bfseries\color{clit}b)}\quad\begin{minipage}[c]{5.3cm}\raggedright\Large on ne peut pas savoir\end{minipage}\end{tcolorbox}\par\vspace{2.5mm}
\begin{tcolorbox}[enhanced,nobeforeafter,width=7.1cm,height=2.1cm,valign=center,colback=clit!6,colframe=clit,arc=2.5mm,boxrule=0.9pt,left=3mm]{\Large\bfseries\color{clit}c)}\quad\begin{minipage}[c]{5.3cm}\raggedright\large seulement si $x$ est positif\end{minipage}\end{tcolorbox}\hfill\begin{tcolorbox}[enhanced,nobeforeafter,width=7.1cm,height=2.1cm,valign=center,colback=clit!6,colframe=clit,arc=2.5mm,boxrule=0.9pt,left=3mm]{\Large\bfseries\color{clit}d)}\quad\begin{minipage}[c]{5.3cm}\raggedright\Large non\end{minipage}\end{tcolorbox}
\end{tcolorbox}
\newpage
\begin{tcolorbox}[enhanced,width=16cm,height=14cm,colback=white,colframe=clit,boxrule=2pt,arc=0mm,
left=6mm,right=6mm,top=0mm,bottom=4mm,sharp corners,
overlay={\fill[clit] (frame.north west) rectangle ([yshift=-1.25cm]frame.north east);
\node[anchor=west,text=white,font=\sffamily\bfseries\LARGE] at ([xshift=5mm,yshift=-0.62cm]frame.north west) {Question 31 B};
\node[anchor=east,text=white,font=\sffamily\itshape\large] at ([xshift=-5mm,yshift=-0.62cm]frame.north east) {Calcul littéral et équations};}]
\vspace*{1.75cm}
\begin{minipage}[t][6.9cm][c]{\linewidth}\Large\centering
\begin{minipage}{0.95\linewidth}\raggedright\linespread{1.25}\selectfont Programme de calcul~: choisir un nombre ; soustraire $1$ ; multiplier le résultat par le nombre choisi ; soustraire $4$. Quel résultat obtient-on en choisissant $5$~?\end{minipage}
\end{minipage}

\vfill
\begin{tcolorbox}[enhanced,nobeforeafter,width=7.1cm,height=2.1cm,valign=center,colback=clit!6,colframe=clit,arc=2.5mm,boxrule=0.9pt,left=3mm]{\Large\bfseries\color{clit}a)}\quad\begin{minipage}[c]{5.3cm}\raggedright\Large $16$\end{minipage}\end{tcolorbox}\hfill\begin{tcolorbox}[enhanced,nobeforeafter,width=7.1cm,height=2.1cm,valign=center,colback=clit!6,colframe=clit,arc=2.5mm,boxrule=0.9pt,left=3mm]{\Large\bfseries\color{clit}b)}\quad\begin{minipage}[c]{5.3cm}\raggedright\Large $-4$\end{minipage}\end{tcolorbox}\par\vspace{2.5mm}
\begin{tcolorbox}[enhanced,nobeforeafter,width=7.1cm,height=2.1cm,valign=center,colback=clit!6,colframe=clit,arc=2.5mm,boxrule=0.9pt,left=3mm]{\Large\bfseries\color{clit}c)}\quad\begin{minipage}[c]{5.3cm}\raggedright\Large $24$\end{minipage}\end{tcolorbox}\hfill\begin{tcolorbox}[enhanced,nobeforeafter,width=7.1cm,height=2.1cm,valign=center,colback=clit!6,colframe=clit,arc=2.5mm,boxrule=0.9pt,left=3mm]{\Large\bfseries\color{clit}d)}\quad\begin{minipage}[c]{5.3cm}\raggedright\Large $0$\end{minipage}\end{tcolorbox}
\end{tcolorbox}
\newpage
\begin{tcolorbox}[enhanced,width=16cm,height=14cm,colback=white,colframe=clit,boxrule=2pt,arc=0mm,
left=6mm,right=6mm,top=0mm,bottom=4mm,sharp corners,
overlay={\fill[clit] (frame.north west) rectangle ([yshift=-1.25cm]frame.north east);
\node[anchor=west,text=white,font=\sffamily\bfseries\LARGE] at ([xshift=5mm,yshift=-0.62cm]frame.north west) {Question 32 B};
\node[anchor=east,text=white,font=\sffamily\itshape\large] at ([xshift=-5mm,yshift=-0.62cm]frame.north east) {Calcul littéral et équations};}]
\vspace*{1.75cm}
\begin{minipage}[t][6.9cm][c]{\linewidth}\Large\centering
\begin{minipage}{0.95\linewidth}\raggedright\linespread{1.25}\selectfont Programme de calcul~: choisir un nombre ; ajouter $6$ ; multiplier le résultat par le nombre choisi ; ajouter $1$. Quel résultat obtient-on en choisissant $4$~?\end{minipage}
\end{minipage}

\vfill
\begin{tcolorbox}[enhanced,nobeforeafter,width=7.1cm,height=2.1cm,valign=center,colback=clit!6,colframe=clit,arc=2.5mm,boxrule=0.9pt,left=3mm]{\Large\bfseries\color{clit}a)}\quad\begin{minipage}[c]{5.3cm}\raggedright\Large $11$\end{minipage}\end{tcolorbox}\hfill\begin{tcolorbox}[enhanced,nobeforeafter,width=7.1cm,height=2.1cm,valign=center,colback=clit!6,colframe=clit,arc=2.5mm,boxrule=0.9pt,left=3mm]{\Large\bfseries\color{clit}b)}\quad\begin{minipage}[c]{5.3cm}\raggedright\Large $41$\end{minipage}\end{tcolorbox}\par\vspace{2.5mm}
\begin{tcolorbox}[enhanced,nobeforeafter,width=7.1cm,height=2.1cm,valign=center,colback=clit!6,colframe=clit,arc=2.5mm,boxrule=0.9pt,left=3mm]{\Large\bfseries\color{clit}c)}\quad\begin{minipage}[c]{5.3cm}\raggedright\Large $29$\end{minipage}\end{tcolorbox}\hfill\begin{tcolorbox}[enhanced,nobeforeafter,width=7.1cm,height=2.1cm,valign=center,colback=clit!6,colframe=clit,arc=2.5mm,boxrule=0.9pt,left=3mm]{\Large\bfseries\color{clit}d)}\quad\begin{minipage}[c]{5.3cm}\raggedright\Large $39$\end{minipage}\end{tcolorbox}
\end{tcolorbox}
\newpage
\begin{tcolorbox}[enhanced,width=16cm,height=14cm,colback=white,colframe=cfon,boxrule=2pt,arc=0mm,
left=6mm,right=6mm,top=0mm,bottom=4mm,sharp corners,
overlay={\fill[cfon] (frame.north west) rectangle ([yshift=-1.25cm]frame.north east);
\node[anchor=west,text=white,font=\sffamily\bfseries\LARGE] at ([xshift=5mm,yshift=-0.62cm]frame.north west) {Question 33 B};
\node[anchor=east,text=white,font=\sffamily\itshape\large] at ([xshift=-5mm,yshift=-0.62cm]frame.north east) {Fonctions};}]
\vspace*{1.75cm}
\begin{minipage}[t][6.9cm][c]{\linewidth}\LARGE\centering
\begin{minipage}{0.95\linewidth}\raggedright\linespread{1.25}\selectfont Soit $f(x)=5x+1$. Quelle est l'image de $-2$ par $f$~?\end{minipage}
\end{minipage}

\vfill
\begin{tcolorbox}[enhanced,nobeforeafter,width=7.1cm,height=2.1cm,valign=center,colback=cfon!6,colframe=cfon,arc=2.5mm,boxrule=0.9pt,left=3mm]{\Large\bfseries\color{cfon}a)}\quad\begin{minipage}[c]{5.3cm}\raggedright\Large $-11$\end{minipage}\end{tcolorbox}\hfill\begin{tcolorbox}[enhanced,nobeforeafter,width=7.1cm,height=2.1cm,valign=center,colback=cfon!6,colframe=cfon,arc=2.5mm,boxrule=0.9pt,left=3mm]{\Large\bfseries\color{cfon}b)}\quad\begin{minipage}[c]{5.3cm}\raggedright\Large $-\dfrac{3}{5}$\end{minipage}\end{tcolorbox}\par\vspace{2.5mm}
\begin{tcolorbox}[enhanced,nobeforeafter,width=7.1cm,height=2.1cm,valign=center,colback=cfon!6,colframe=cfon,arc=2.5mm,boxrule=0.9pt,left=3mm]{\Large\bfseries\color{cfon}c)}\quad\begin{minipage}[c]{5.3cm}\raggedright\Large $4$\end{minipage}\end{tcolorbox}\hfill\begin{tcolorbox}[enhanced,nobeforeafter,width=7.1cm,height=2.1cm,valign=center,colback=cfon!6,colframe=cfon,arc=2.5mm,boxrule=0.9pt,left=3mm]{\Large\bfseries\color{cfon}d)}\quad\begin{minipage}[c]{5.3cm}\raggedright\Large $-9$\end{minipage}\end{tcolorbox}
\end{tcolorbox}
\newpage
\begin{tcolorbox}[enhanced,width=16cm,height=14cm,colback=white,colframe=cfon,boxrule=2pt,arc=0mm,
left=6mm,right=6mm,top=0mm,bottom=4mm,sharp corners,
overlay={\fill[cfon] (frame.north west) rectangle ([yshift=-1.25cm]frame.north east);
\node[anchor=west,text=white,font=\sffamily\bfseries\LARGE] at ([xshift=5mm,yshift=-0.62cm]frame.north west) {Question 34 B};
\node[anchor=east,text=white,font=\sffamily\itshape\large] at ([xshift=-5mm,yshift=-0.62cm]frame.north east) {Fonctions};}]
\vspace*{1.75cm}
\begin{minipage}[t][6.9cm][c]{\linewidth}\LARGE\centering
\begin{minipage}{0.95\linewidth}\raggedright\linespread{1.25}\selectfont Soit $f(x)=-4x-3$. Quelle est l'image de $-1$ par $f$~?\end{minipage}
\end{minipage}

\vfill
\begin{tcolorbox}[enhanced,nobeforeafter,width=7.1cm,height=2.1cm,valign=center,colback=cfon!6,colframe=cfon,arc=2.5mm,boxrule=0.9pt,left=3mm]{\Large\bfseries\color{cfon}a)}\quad\begin{minipage}[c]{5.3cm}\raggedright\Large $1$\end{minipage}\end{tcolorbox}\hfill\begin{tcolorbox}[enhanced,nobeforeafter,width=7.1cm,height=2.1cm,valign=center,colback=cfon!6,colframe=cfon,arc=2.5mm,boxrule=0.9pt,left=3mm]{\Large\bfseries\color{cfon}b)}\quad\begin{minipage}[c]{5.3cm}\raggedright\Large $-\dfrac{1}{2}$\end{minipage}\end{tcolorbox}\par\vspace{2.5mm}
\begin{tcolorbox}[enhanced,nobeforeafter,width=7.1cm,height=2.1cm,valign=center,colback=cfon!6,colframe=cfon,arc=2.5mm,boxrule=0.9pt,left=3mm]{\Large\bfseries\color{cfon}c)}\quad\begin{minipage}[c]{5.3cm}\raggedright\Large $-8$\end{minipage}\end{tcolorbox}\hfill\begin{tcolorbox}[enhanced,nobeforeafter,width=7.1cm,height=2.1cm,valign=center,colback=cfon!6,colframe=cfon,arc=2.5mm,boxrule=0.9pt,left=3mm]{\Large\bfseries\color{cfon}d)}\quad\begin{minipage}[c]{5.3cm}\raggedright\Large $7$\end{minipage}\end{tcolorbox}
\end{tcolorbox}
\newpage
\begin{tcolorbox}[enhanced,width=16cm,height=14cm,colback=white,colframe=cfon,boxrule=2pt,arc=0mm,
left=6mm,right=6mm,top=0mm,bottom=4mm,sharp corners,
overlay={\fill[cfon] (frame.north west) rectangle ([yshift=-1.25cm]frame.north east);
\node[anchor=west,text=white,font=\sffamily\bfseries\LARGE] at ([xshift=5mm,yshift=-0.62cm]frame.north west) {Question 35 B};
\node[anchor=east,text=white,font=\sffamily\itshape\large] at ([xshift=-5mm,yshift=-0.62cm]frame.north east) {Fonctions};}]
\vspace*{1.75cm}
\begin{minipage}[t][6.9cm][c]{\linewidth}\LARGE\centering
\begin{minipage}{0.95\linewidth}\raggedright\linespread{1.25}\selectfont Soit $f(x)=5x-2$. Quel est l'antécédent de $23$ par $f$~?\end{minipage}
\end{minipage}

\vfill
\begin{tcolorbox}[enhanced,nobeforeafter,width=7.1cm,height=2.1cm,valign=center,colback=cfon!6,colframe=cfon,arc=2.5mm,boxrule=0.9pt,left=3mm]{\Large\bfseries\color{cfon}a)}\quad\begin{minipage}[c]{5.3cm}\raggedright\Large $113$\end{minipage}\end{tcolorbox}\hfill\begin{tcolorbox}[enhanced,nobeforeafter,width=7.1cm,height=2.1cm,valign=center,colback=cfon!6,colframe=cfon,arc=2.5mm,boxrule=0.9pt,left=3mm]{\Large\bfseries\color{cfon}b)}\quad\begin{minipage}[c]{5.3cm}\raggedright\Large $\dfrac{21}{5}$\end{minipage}\end{tcolorbox}\par\vspace{2.5mm}
\begin{tcolorbox}[enhanced,nobeforeafter,width=7.1cm,height=2.1cm,valign=center,colback=cfon!6,colframe=cfon,arc=2.5mm,boxrule=0.9pt,left=3mm]{\Large\bfseries\color{cfon}c)}\quad\begin{minipage}[c]{5.3cm}\raggedright\Large $\dfrac{33}{5}$\end{minipage}\end{tcolorbox}\hfill\begin{tcolorbox}[enhanced,nobeforeafter,width=7.1cm,height=2.1cm,valign=center,colback=cfon!6,colframe=cfon,arc=2.5mm,boxrule=0.9pt,left=3mm]{\Large\bfseries\color{cfon}d)}\quad\begin{minipage}[c]{5.3cm}\raggedright\Large $5$\end{minipage}\end{tcolorbox}
\end{tcolorbox}
\newpage
\begin{tcolorbox}[enhanced,width=16cm,height=14cm,colback=white,colframe=cfon,boxrule=2pt,arc=0mm,
left=6mm,right=6mm,top=0mm,bottom=4mm,sharp corners,
overlay={\fill[cfon] (frame.north west) rectangle ([yshift=-1.25cm]frame.north east);
\node[anchor=west,text=white,font=\sffamily\bfseries\LARGE] at ([xshift=5mm,yshift=-0.62cm]frame.north west) {Question 36 B};
\node[anchor=east,text=white,font=\sffamily\itshape\large] at ([xshift=-5mm,yshift=-0.62cm]frame.north east) {Fonctions};}]
\vspace*{1.75cm}
\begin{minipage}[t][6.9cm][c]{\linewidth}\LARGE\centering
\begin{minipage}{0.95\linewidth}\raggedright\linespread{1.25}\selectfont Soit $f(x)=4x+1$. Quel est l'antécédent de $-7$ par $f$~?\end{minipage}
\end{minipage}

\vfill
\begin{tcolorbox}[enhanced,nobeforeafter,width=7.1cm,height=2.1cm,valign=center,colback=cfon!6,colframe=cfon,arc=2.5mm,boxrule=0.9pt,left=3mm]{\Large\bfseries\color{cfon}a)}\quad\begin{minipage}[c]{5.3cm}\raggedright\Large $-27$\end{minipage}\end{tcolorbox}\hfill\begin{tcolorbox}[enhanced,nobeforeafter,width=7.1cm,height=2.1cm,valign=center,colback=cfon!6,colframe=cfon,arc=2.5mm,boxrule=0.9pt,left=3mm]{\Large\bfseries\color{cfon}b)}\quad\begin{minipage}[c]{5.3cm}\raggedright\Large $-\dfrac{11}{4}$\end{minipage}\end{tcolorbox}\par\vspace{2.5mm}
\begin{tcolorbox}[enhanced,nobeforeafter,width=7.1cm,height=2.1cm,valign=center,colback=cfon!6,colframe=cfon,arc=2.5mm,boxrule=0.9pt,left=3mm]{\Large\bfseries\color{cfon}c)}\quad\begin{minipage}[c]{5.3cm}\raggedright\Large $-2$\end{minipage}\end{tcolorbox}\hfill\begin{tcolorbox}[enhanced,nobeforeafter,width=7.1cm,height=2.1cm,valign=center,colback=cfon!6,colframe=cfon,arc=2.5mm,boxrule=0.9pt,left=3mm]{\Large\bfseries\color{cfon}d)}\quad\begin{minipage}[c]{5.3cm}\raggedright\Large $-\dfrac{3}{2}$\end{minipage}\end{tcolorbox}
\end{tcolorbox}
\newpage
\begin{tcolorbox}[enhanced,width=16cm,height=14cm,colback=white,colframe=cfon,boxrule=2pt,arc=0mm,
left=6mm,right=6mm,top=0mm,bottom=4mm,sharp corners,
overlay={\fill[cfon] (frame.north west) rectangle ([yshift=-1.25cm]frame.north east);
\node[anchor=west,text=white,font=\sffamily\bfseries\LARGE] at ([xshift=5mm,yshift=-0.62cm]frame.north west) {Question 37 B};
\node[anchor=east,text=white,font=\sffamily\itshape\large] at ([xshift=-5mm,yshift=-0.62cm]frame.north east) {Fonctions};}]
\vspace*{1.75cm}
\begin{minipage}[t][6.9cm][c]{\linewidth}\LARGE\centering
\begin{minipage}{0.95\linewidth}\raggedright\linespread{1.25}\selectfont $f$ est une fonction linéaire telle que $f(-2)=7$. Quel est son coefficient~?\end{minipage}
\end{minipage}

\vfill
\begin{tcolorbox}[enhanced,nobeforeafter,width=7.1cm,height=2.1cm,valign=center,colback=cfon!6,colframe=cfon,arc=2.5mm,boxrule=0.9pt,left=3mm]{\Large\bfseries\color{cfon}a)}\quad\begin{minipage}[c]{5.3cm}\raggedright\Large $9$\end{minipage}\end{tcolorbox}\hfill\begin{tcolorbox}[enhanced,nobeforeafter,width=7.1cm,height=2.1cm,valign=center,colback=cfon!6,colframe=cfon,arc=2.5mm,boxrule=0.9pt,left=3mm]{\Large\bfseries\color{cfon}b)}\quad\begin{minipage}[c]{5.3cm}\raggedright\Large $-14$\end{minipage}\end{tcolorbox}\par\vspace{2.5mm}
\begin{tcolorbox}[enhanced,nobeforeafter,width=7.1cm,height=2.1cm,valign=center,colback=cfon!6,colframe=cfon,arc=2.5mm,boxrule=0.9pt,left=3mm]{\Large\bfseries\color{cfon}c)}\quad\begin{minipage}[c]{5.3cm}\raggedright\Large $-\dfrac{2}{7}$\end{minipage}\end{tcolorbox}\hfill\begin{tcolorbox}[enhanced,nobeforeafter,width=7.1cm,height=2.1cm,valign=center,colback=cfon!6,colframe=cfon,arc=2.5mm,boxrule=0.9pt,left=3mm]{\Large\bfseries\color{cfon}d)}\quad\begin{minipage}[c]{5.3cm}\raggedright\Large $-\dfrac{7}{2}$\end{minipage}\end{tcolorbox}
\end{tcolorbox}
\newpage
\begin{tcolorbox}[enhanced,width=16cm,height=14cm,colback=white,colframe=cfon,boxrule=2pt,arc=0mm,
left=6mm,right=6mm,top=0mm,bottom=4mm,sharp corners,
overlay={\fill[cfon] (frame.north west) rectangle ([yshift=-1.25cm]frame.north east);
\node[anchor=west,text=white,font=\sffamily\bfseries\LARGE] at ([xshift=5mm,yshift=-0.62cm]frame.north west) {Question 38 B};
\node[anchor=east,text=white,font=\sffamily\itshape\large] at ([xshift=-5mm,yshift=-0.62cm]frame.north east) {Fonctions};}]
\vspace*{1.75cm}
\begin{minipage}[t][6.9cm][c]{\linewidth}\LARGE\centering
\begin{minipage}{0.95\linewidth}\raggedright\linespread{1.25}\selectfont $f$ est une fonction linéaire telle que $f(3)=1$. Quel est son coefficient~?\end{minipage}
\end{minipage}

\vfill
\begin{tcolorbox}[enhanced,nobeforeafter,width=7.1cm,height=2.1cm,valign=center,colback=cfon!6,colframe=cfon,arc=2.5mm,boxrule=0.9pt,left=3mm]{\Large\bfseries\color{cfon}a)}\quad\begin{minipage}[c]{5.3cm}\raggedright\Large $-2$\end{minipage}\end{tcolorbox}\hfill\begin{tcolorbox}[enhanced,nobeforeafter,width=7.1cm,height=2.1cm,valign=center,colback=cfon!6,colframe=cfon,arc=2.5mm,boxrule=0.9pt,left=3mm]{\Large\bfseries\color{cfon}b)}\quad\begin{minipage}[c]{5.3cm}\raggedright\Large $-\dfrac{1}{3}$\end{minipage}\end{tcolorbox}\par\vspace{2.5mm}
\begin{tcolorbox}[enhanced,nobeforeafter,width=7.1cm,height=2.1cm,valign=center,colback=cfon!6,colframe=cfon,arc=2.5mm,boxrule=0.9pt,left=3mm]{\Large\bfseries\color{cfon}c)}\quad\begin{minipage}[c]{5.3cm}\raggedright\Large $\dfrac{1}{3}$\end{minipage}\end{tcolorbox}\hfill\begin{tcolorbox}[enhanced,nobeforeafter,width=7.1cm,height=2.1cm,valign=center,colback=cfon!6,colframe=cfon,arc=2.5mm,boxrule=0.9pt,left=3mm]{\Large\bfseries\color{cfon}d)}\quad\begin{minipage}[c]{5.3cm}\raggedright\Large $3$\end{minipage}\end{tcolorbox}
\end{tcolorbox}
\newpage
\begin{tcolorbox}[enhanced,width=16cm,height=14cm,colback=white,colframe=cfon,boxrule=2pt,arc=0mm,
left=6mm,right=6mm,top=0mm,bottom=4mm,sharp corners,
overlay={\fill[cfon] (frame.north west) rectangle ([yshift=-1.25cm]frame.north east);
\node[anchor=west,text=white,font=\sffamily\bfseries\LARGE] at ([xshift=5mm,yshift=-0.62cm]frame.north west) {Question 39 B};
\node[anchor=east,text=white,font=\sffamily\itshape\large] at ([xshift=-5mm,yshift=-0.62cm]frame.north east) {Fonctions};}]
\vspace*{1.75cm}
\begin{minipage}[t][6.9cm][c]{\linewidth}\LARGE\centering
\begin{minipage}{0.95\linewidth}\raggedright\linespread{1.25}\selectfont $f$ est une fonction affine telle que $f(0)=4$ et $f(5)=14$. Quel est le coefficient $a$ de $f(x)=ax+b$~?\end{minipage}
\end{minipage}

\vfill
\begin{tcolorbox}[enhanced,nobeforeafter,width=7.1cm,height=2.1cm,valign=center,colback=cfon!6,colframe=cfon,arc=2.5mm,boxrule=0.9pt,left=3mm]{\Large\bfseries\color{cfon}a)}\quad\begin{minipage}[c]{5.3cm}\raggedright\Large $-2$\end{minipage}\end{tcolorbox}\hfill\begin{tcolorbox}[enhanced,nobeforeafter,width=7.1cm,height=2.1cm,valign=center,colback=cfon!6,colframe=cfon,arc=2.5mm,boxrule=0.9pt,left=3mm]{\Large\bfseries\color{cfon}b)}\quad\begin{minipage}[c]{5.3cm}\raggedright\Large $2$\end{minipage}\end{tcolorbox}\par\vspace{2.5mm}
\begin{tcolorbox}[enhanced,nobeforeafter,width=7.1cm,height=2.1cm,valign=center,colback=cfon!6,colframe=cfon,arc=2.5mm,boxrule=0.9pt,left=3mm]{\Large\bfseries\color{cfon}c)}\quad\begin{minipage}[c]{5.3cm}\raggedright\Large $10$\end{minipage}\end{tcolorbox}\hfill\begin{tcolorbox}[enhanced,nobeforeafter,width=7.1cm,height=2.1cm,valign=center,colback=cfon!6,colframe=cfon,arc=2.5mm,boxrule=0.9pt,left=3mm]{\Large\bfseries\color{cfon}d)}\quad\begin{minipage}[c]{5.3cm}\raggedright\Large $\dfrac{1}{2}$\end{minipage}\end{tcolorbox}
\end{tcolorbox}
\newpage
\begin{tcolorbox}[enhanced,width=16cm,height=14cm,colback=white,colframe=cfon,boxrule=2pt,arc=0mm,
left=6mm,right=6mm,top=0mm,bottom=4mm,sharp corners,
overlay={\fill[cfon] (frame.north west) rectangle ([yshift=-1.25cm]frame.north east);
\node[anchor=west,text=white,font=\sffamily\bfseries\LARGE] at ([xshift=5mm,yshift=-0.62cm]frame.north west) {Question 40 B};
\node[anchor=east,text=white,font=\sffamily\itshape\large] at ([xshift=-5mm,yshift=-0.62cm]frame.north east) {Fonctions};}]
\vspace*{1.75cm}
\begin{minipage}[t][6.9cm][c]{\linewidth}\LARGE\centering
\begin{minipage}{0.95\linewidth}\raggedright\linespread{1.25}\selectfont $f$ est une fonction affine telle que $f(-1)=3$ et $f(2)=-3$. Quel est le coefficient $a$ de $f(x)=ax+b$~?\end{minipage}
\end{minipage}

\vfill
\begin{tcolorbox}[enhanced,nobeforeafter,width=7.1cm,height=2.1cm,valign=center,colback=cfon!6,colframe=cfon,arc=2.5mm,boxrule=0.9pt,left=3mm]{\Large\bfseries\color{cfon}a)}\quad\begin{minipage}[c]{5.3cm}\raggedright\Large $-2$\end{minipage}\end{tcolorbox}\hfill\begin{tcolorbox}[enhanced,nobeforeafter,width=7.1cm,height=2.1cm,valign=center,colback=cfon!6,colframe=cfon,arc=2.5mm,boxrule=0.9pt,left=3mm]{\Large\bfseries\color{cfon}b)}\quad\begin{minipage}[c]{5.3cm}\raggedright\Large $-\dfrac{1}{2}$\end{minipage}\end{tcolorbox}\par\vspace{2.5mm}
\begin{tcolorbox}[enhanced,nobeforeafter,width=7.1cm,height=2.1cm,valign=center,colback=cfon!6,colframe=cfon,arc=2.5mm,boxrule=0.9pt,left=3mm]{\Large\bfseries\color{cfon}c)}\quad\begin{minipage}[c]{5.3cm}\raggedright\Large $2$\end{minipage}\end{tcolorbox}\hfill\begin{tcolorbox}[enhanced,nobeforeafter,width=7.1cm,height=2.1cm,valign=center,colback=cfon!6,colframe=cfon,arc=2.5mm,boxrule=0.9pt,left=3mm]{\Large\bfseries\color{cfon}d)}\quad\begin{minipage}[c]{5.3cm}\raggedright\Large $-6$\end{minipage}\end{tcolorbox}
\end{tcolorbox}
\newpage
\begin{tcolorbox}[enhanced,width=16cm,height=14cm,colback=white,colframe=cfon,boxrule=2pt,arc=0mm,
left=6mm,right=6mm,top=0mm,bottom=4mm,sharp corners,
overlay={\fill[cfon] (frame.north west) rectangle ([yshift=-1.25cm]frame.north east);
\node[anchor=west,text=white,font=\sffamily\bfseries\LARGE] at ([xshift=5mm,yshift=-0.62cm]frame.north west) {Question 41 B};
\node[anchor=east,text=white,font=\sffamily\itshape\large] at ([xshift=-5mm,yshift=-0.62cm]frame.north east) {Fonctions};}]
\vspace*{1.75cm}
\begin{minipage}[t][6.9cm][c]{\linewidth}\LARGE\centering
\begin{minipage}{0.95\linewidth}\raggedright\linespread{1.25}\selectfont $f$ est la fonction affine $f(x)=4x+b$ et $f(-2)=-5$. Que vaut $b$~?\end{minipage}
\end{minipage}

\vfill
\begin{tcolorbox}[enhanced,nobeforeafter,width=7.1cm,height=2.1cm,valign=center,colback=cfon!6,colframe=cfon,arc=2.5mm,boxrule=0.9pt,left=3mm]{\Large\bfseries\color{cfon}a)}\quad\begin{minipage}[c]{5.3cm}\raggedright\Large $-5$\end{minipage}\end{tcolorbox}\hfill\begin{tcolorbox}[enhanced,nobeforeafter,width=7.1cm,height=2.1cm,valign=center,colback=cfon!6,colframe=cfon,arc=2.5mm,boxrule=0.9pt,left=3mm]{\Large\bfseries\color{cfon}b)}\quad\begin{minipage}[c]{5.3cm}\raggedright\Large $\dfrac{5}{2}$\end{minipage}\end{tcolorbox}\par\vspace{2.5mm}
\begin{tcolorbox}[enhanced,nobeforeafter,width=7.1cm,height=2.1cm,valign=center,colback=cfon!6,colframe=cfon,arc=2.5mm,boxrule=0.9pt,left=3mm]{\Large\bfseries\color{cfon}c)}\quad\begin{minipage}[c]{5.3cm}\raggedright\Large $-13$\end{minipage}\end{tcolorbox}\hfill\begin{tcolorbox}[enhanced,nobeforeafter,width=7.1cm,height=2.1cm,valign=center,colback=cfon!6,colframe=cfon,arc=2.5mm,boxrule=0.9pt,left=3mm]{\Large\bfseries\color{cfon}d)}\quad\begin{minipage}[c]{5.3cm}\raggedright\Large $3$\end{minipage}\end{tcolorbox}
\end{tcolorbox}
\newpage
\begin{tcolorbox}[enhanced,width=16cm,height=14cm,colback=white,colframe=cfon,boxrule=2pt,arc=0mm,
left=6mm,right=6mm,top=0mm,bottom=4mm,sharp corners,
overlay={\fill[cfon] (frame.north west) rectangle ([yshift=-1.25cm]frame.north east);
\node[anchor=west,text=white,font=\sffamily\bfseries\LARGE] at ([xshift=5mm,yshift=-0.62cm]frame.north west) {Question 42 B};
\node[anchor=east,text=white,font=\sffamily\itshape\large] at ([xshift=-5mm,yshift=-0.62cm]frame.north east) {Fonctions};}]
\vspace*{1.75cm}
\begin{minipage}[t][6.9cm][c]{\linewidth}\LARGE\centering
\begin{minipage}{0.95\linewidth}\raggedright\linespread{1.25}\selectfont $f$ est la fonction affine $f(x)=5x+b$ et $f(2)=7$. Que vaut $b$~?\end{minipage}
\end{minipage}

\vfill
\begin{tcolorbox}[enhanced,nobeforeafter,width=7.1cm,height=2.1cm,valign=center,colback=cfon!6,colframe=cfon,arc=2.5mm,boxrule=0.9pt,left=3mm]{\Large\bfseries\color{cfon}a)}\quad\begin{minipage}[c]{5.3cm}\raggedright\Large $7$\end{minipage}\end{tcolorbox}\hfill\begin{tcolorbox}[enhanced,nobeforeafter,width=7.1cm,height=2.1cm,valign=center,colback=cfon!6,colframe=cfon,arc=2.5mm,boxrule=0.9pt,left=3mm]{\Large\bfseries\color{cfon}b)}\quad\begin{minipage}[c]{5.3cm}\raggedright\Large $-3$\end{minipage}\end{tcolorbox}\par\vspace{2.5mm}
\begin{tcolorbox}[enhanced,nobeforeafter,width=7.1cm,height=2.1cm,valign=center,colback=cfon!6,colframe=cfon,arc=2.5mm,boxrule=0.9pt,left=3mm]{\Large\bfseries\color{cfon}c)}\quad\begin{minipage}[c]{5.3cm}\raggedright\Large $\dfrac{7}{2}$\end{minipage}\end{tcolorbox}\hfill\begin{tcolorbox}[enhanced,nobeforeafter,width=7.1cm,height=2.1cm,valign=center,colback=cfon!6,colframe=cfon,arc=2.5mm,boxrule=0.9pt,left=3mm]{\Large\bfseries\color{cfon}d)}\quad\begin{minipage}[c]{5.3cm}\raggedright\Large $17$\end{minipage}\end{tcolorbox}
\end{tcolorbox}
\newpage
\begin{tcolorbox}[enhanced,width=16cm,height=14cm,colback=white,colframe=cfon,boxrule=2pt,arc=0mm,
left=6mm,right=6mm,top=0mm,bottom=4mm,sharp corners,
overlay={\fill[cfon] (frame.north west) rectangle ([yshift=-1.25cm]frame.north east);
\node[anchor=west,text=white,font=\sffamily\bfseries\LARGE] at ([xshift=5mm,yshift=-0.62cm]frame.north west) {Question 43 B};
\node[anchor=east,text=white,font=\sffamily\itshape\large] at ([xshift=-5mm,yshift=-0.62cm]frame.north east) {Fonctions};}]
\vspace*{1.75cm}
\begin{minipage}[t][6.9cm][c]{\linewidth}\LARGE\centering
\begin{minipage}{0.95\linewidth}\raggedright\linespread{1.25}\selectfont Augmenter une quantité de $8\,\%$ revient à la multiplier par~:\end{minipage}
\end{minipage}

\vfill
\begin{tcolorbox}[enhanced,nobeforeafter,width=7.1cm,height=2.1cm,valign=center,colback=cfon!6,colframe=cfon,arc=2.5mm,boxrule=0.9pt,left=3mm]{\Large\bfseries\color{cfon}a)}\quad\begin{minipage}[c]{5.3cm}\raggedright\Large $1{,}08$\end{minipage}\end{tcolorbox}\hfill\begin{tcolorbox}[enhanced,nobeforeafter,width=7.1cm,height=2.1cm,valign=center,colback=cfon!6,colframe=cfon,arc=2.5mm,boxrule=0.9pt,left=3mm]{\Large\bfseries\color{cfon}b)}\quad\begin{minipage}[c]{5.3cm}\raggedright\Large $0{,}08$\end{minipage}\end{tcolorbox}\par\vspace{2.5mm}
\begin{tcolorbox}[enhanced,nobeforeafter,width=7.1cm,height=2.1cm,valign=center,colback=cfon!6,colframe=cfon,arc=2.5mm,boxrule=0.9pt,left=3mm]{\Large\bfseries\color{cfon}c)}\quad\begin{minipage}[c]{5.3cm}\raggedright\Large $0{,}92$\end{minipage}\end{tcolorbox}\hfill\begin{tcolorbox}[enhanced,nobeforeafter,width=7.1cm,height=2.1cm,valign=center,colback=cfon!6,colframe=cfon,arc=2.5mm,boxrule=0.9pt,left=3mm]{\Large\bfseries\color{cfon}d)}\quad\begin{minipage}[c]{5.3cm}\raggedright\Large $1{,}8$\end{minipage}\end{tcolorbox}
\end{tcolorbox}
\newpage
\begin{tcolorbox}[enhanced,width=16cm,height=14cm,colback=white,colframe=cfon,boxrule=2pt,arc=0mm,
left=6mm,right=6mm,top=0mm,bottom=4mm,sharp corners,
overlay={\fill[cfon] (frame.north west) rectangle ([yshift=-1.25cm]frame.north east);
\node[anchor=west,text=white,font=\sffamily\bfseries\LARGE] at ([xshift=5mm,yshift=-0.62cm]frame.north west) {Question 44 B};
\node[anchor=east,text=white,font=\sffamily\itshape\large] at ([xshift=-5mm,yshift=-0.62cm]frame.north east) {Fonctions};}]
\vspace*{1.75cm}
\begin{minipage}[t][6.9cm][c]{\linewidth}\LARGE\centering
\begin{minipage}{0.95\linewidth}\raggedright\linespread{1.25}\selectfont Diminuer une quantité de $35\,\%$ revient à la multiplier par~:\end{minipage}
\end{minipage}

\vfill
\begin{tcolorbox}[enhanced,nobeforeafter,width=7.1cm,height=2.1cm,valign=center,colback=cfon!6,colframe=cfon,arc=2.5mm,boxrule=0.9pt,left=3mm]{\Large\bfseries\color{cfon}a)}\quad\begin{minipage}[c]{5.3cm}\raggedright\Large $0{,}35$\end{minipage}\end{tcolorbox}\hfill\begin{tcolorbox}[enhanced,nobeforeafter,width=7.1cm,height=2.1cm,valign=center,colback=cfon!6,colframe=cfon,arc=2.5mm,boxrule=0.9pt,left=3mm]{\Large\bfseries\color{cfon}b)}\quad\begin{minipage}[c]{5.3cm}\raggedright\Large $1{,}35$\end{minipage}\end{tcolorbox}\par\vspace{2.5mm}
\begin{tcolorbox}[enhanced,nobeforeafter,width=7.1cm,height=2.1cm,valign=center,colback=cfon!6,colframe=cfon,arc=2.5mm,boxrule=0.9pt,left=3mm]{\Large\bfseries\color{cfon}c)}\quad\begin{minipage}[c]{5.3cm}\raggedright\Large $0{,}65$\end{minipage}\end{tcolorbox}\hfill\begin{tcolorbox}[enhanced,nobeforeafter,width=7.1cm,height=2.1cm,valign=center,colback=cfon!6,colframe=cfon,arc=2.5mm,boxrule=0.9pt,left=3mm]{\Large\bfseries\color{cfon}d)}\quad\begin{minipage}[c]{5.3cm}\raggedright\Large $1{,}65$\end{minipage}\end{tcolorbox}
\end{tcolorbox}
\newpage
\begin{tcolorbox}[enhanced,width=16cm,height=14cm,colback=white,colframe=cfon,boxrule=2pt,arc=0mm,
left=6mm,right=6mm,top=0mm,bottom=4mm,sharp corners,
overlay={\fill[cfon] (frame.north west) rectangle ([yshift=-1.25cm]frame.north east);
\node[anchor=west,text=white,font=\sffamily\bfseries\LARGE] at ([xshift=5mm,yshift=-0.62cm]frame.north west) {Question 45 B};
\node[anchor=east,text=white,font=\sffamily\itshape\large] at ([xshift=-5mm,yshift=-0.62cm]frame.north east) {Fonctions};}]
\vspace*{1.75cm}
\begin{minipage}[t][6.9cm][c]{\linewidth}\LARGE\centering
\begin{minipage}{0.95\linewidth}\raggedright\linespread{1.25}\selectfont Soit $g(x)=-x^2+4x+3$. Que vaut $g(3)$~?\end{minipage}
\end{minipage}

\vfill
\begin{tcolorbox}[enhanced,nobeforeafter,width=7.1cm,height=2.1cm,valign=center,colback=cfon!6,colframe=cfon,arc=2.5mm,boxrule=0.9pt,left=3mm]{\Large\bfseries\color{cfon}a)}\quad\begin{minipage}[c]{5.3cm}\raggedright\Large $24$\end{minipage}\end{tcolorbox}\hfill\begin{tcolorbox}[enhanced,nobeforeafter,width=7.1cm,height=2.1cm,valign=center,colback=cfon!6,colframe=cfon,arc=2.5mm,boxrule=0.9pt,left=3mm]{\Large\bfseries\color{cfon}b)}\quad\begin{minipage}[c]{5.3cm}\raggedright\Large $-6$\end{minipage}\end{tcolorbox}\par\vspace{2.5mm}
\begin{tcolorbox}[enhanced,nobeforeafter,width=7.1cm,height=2.1cm,valign=center,colback=cfon!6,colframe=cfon,arc=2.5mm,boxrule=0.9pt,left=3mm]{\Large\bfseries\color{cfon}c)}\quad\begin{minipage}[c]{5.3cm}\raggedright\Large $9$\end{minipage}\end{tcolorbox}\hfill\begin{tcolorbox}[enhanced,nobeforeafter,width=7.1cm,height=2.1cm,valign=center,colback=cfon!6,colframe=cfon,arc=2.5mm,boxrule=0.9pt,left=3mm]{\Large\bfseries\color{cfon}d)}\quad\begin{minipage}[c]{5.3cm}\raggedright\Large $6$\end{minipage}\end{tcolorbox}
\end{tcolorbox}
\newpage
\begin{tcolorbox}[enhanced,width=16cm,height=14cm,colback=white,colframe=cfon,boxrule=2pt,arc=0mm,
left=6mm,right=6mm,top=0mm,bottom=4mm,sharp corners,
overlay={\fill[cfon] (frame.north west) rectangle ([yshift=-1.25cm]frame.north east);
\node[anchor=west,text=white,font=\sffamily\bfseries\LARGE] at ([xshift=5mm,yshift=-0.62cm]frame.north west) {Question 46 B};
\node[anchor=east,text=white,font=\sffamily\itshape\large] at ([xshift=-5mm,yshift=-0.62cm]frame.north east) {Fonctions};}]
\vspace*{1.75cm}
\begin{minipage}[t][6.9cm][c]{\linewidth}\LARGE\centering
\begin{minipage}{0.95\linewidth}\raggedright\linespread{1.25}\selectfont Soit $g(x)=x^2+5x-2$. Que vaut $g(-1)$~?\end{minipage}
\end{minipage}

\vfill
\begin{tcolorbox}[enhanced,nobeforeafter,width=7.1cm,height=2.1cm,valign=center,colback=cfon!6,colframe=cfon,arc=2.5mm,boxrule=0.9pt,left=3mm]{\Large\bfseries\color{cfon}a)}\quad\begin{minipage}[c]{5.3cm}\raggedright\Large $-6$\end{minipage}\end{tcolorbox}\hfill\begin{tcolorbox}[enhanced,nobeforeafter,width=7.1cm,height=2.1cm,valign=center,colback=cfon!6,colframe=cfon,arc=2.5mm,boxrule=0.9pt,left=3mm]{\Large\bfseries\color{cfon}b)}\quad\begin{minipage}[c]{5.3cm}\raggedright\Large $-9$\end{minipage}\end{tcolorbox}\par\vspace{2.5mm}
\begin{tcolorbox}[enhanced,nobeforeafter,width=7.1cm,height=2.1cm,valign=center,colback=cfon!6,colframe=cfon,arc=2.5mm,boxrule=0.9pt,left=3mm]{\Large\bfseries\color{cfon}c)}\quad\begin{minipage}[c]{5.3cm}\raggedright\Large $6$\end{minipage}\end{tcolorbox}\hfill\begin{tcolorbox}[enhanced,nobeforeafter,width=7.1cm,height=2.1cm,valign=center,colback=cfon!6,colframe=cfon,arc=2.5mm,boxrule=0.9pt,left=3mm]{\Large\bfseries\color{cfon}d)}\quad\begin{minipage}[c]{5.3cm}\raggedright\Large $-8$\end{minipage}\end{tcolorbox}
\end{tcolorbox}
\newpage
\begin{tcolorbox}[enhanced,width=16cm,height=14cm,colback=white,colframe=cfon,boxrule=2pt,arc=0mm,
left=6mm,right=6mm,top=0mm,bottom=4mm,sharp corners,
overlay={\fill[cfon] (frame.north west) rectangle ([yshift=-1.25cm]frame.north east);
\node[anchor=west,text=white,font=\sffamily\bfseries\LARGE] at ([xshift=5mm,yshift=-0.62cm]frame.north west) {Question 47 B};
\node[anchor=east,text=white,font=\sffamily\itshape\large] at ([xshift=-5mm,yshift=-0.62cm]frame.north east) {Fonctions};}]
\vspace*{1.75cm}
\begin{minipage}[t][6.9cm][c]{\linewidth}\LARGE\centering
\begin{minipage}{0.95\linewidth}\raggedright\linespread{1.25}\selectfont De quel type est la fonction $f:x\mapsto x^2$~?\end{minipage}
\end{minipage}

\vfill
\begin{tcolorbox}[enhanced,nobeforeafter,width=7.1cm,height=2.1cm,valign=center,colback=cfon!6,colframe=cfon,arc=2.5mm,boxrule=0.9pt,left=3mm]{\Large\bfseries\color{cfon}a)}\quad\begin{minipage}[c]{5.3cm}\raggedright\Large ni affine ni linéaire\end{minipage}\end{tcolorbox}\hfill\begin{tcolorbox}[enhanced,nobeforeafter,width=7.1cm,height=2.1cm,valign=center,colback=cfon!6,colframe=cfon,arc=2.5mm,boxrule=0.9pt,left=3mm]{\Large\bfseries\color{cfon}b)}\quad\begin{minipage}[c]{5.3cm}\raggedright\Large constante\end{minipage}\end{tcolorbox}\par\vspace{2.5mm}
\begin{tcolorbox}[enhanced,nobeforeafter,width=7.1cm,height=2.1cm,valign=center,colback=cfon!6,colframe=cfon,arc=2.5mm,boxrule=0.9pt,left=3mm]{\Large\bfseries\color{cfon}c)}\quad\begin{minipage}[c]{5.3cm}\raggedright\Large affine non linéaire\end{minipage}\end{tcolorbox}\hfill\begin{tcolorbox}[enhanced,nobeforeafter,width=7.1cm,height=2.1cm,valign=center,colback=cfon!6,colframe=cfon,arc=2.5mm,boxrule=0.9pt,left=3mm]{\Large\bfseries\color{cfon}d)}\quad\begin{minipage}[c]{5.3cm}\raggedright\Large linéaire\end{minipage}\end{tcolorbox}
\end{tcolorbox}
\newpage
\begin{tcolorbox}[enhanced,width=16cm,height=14cm,colback=white,colframe=cfon,boxrule=2pt,arc=0mm,
left=6mm,right=6mm,top=0mm,bottom=4mm,sharp corners,
overlay={\fill[cfon] (frame.north west) rectangle ([yshift=-1.25cm]frame.north east);
\node[anchor=west,text=white,font=\sffamily\bfseries\LARGE] at ([xshift=5mm,yshift=-0.62cm]frame.north west) {Question 48 B};
\node[anchor=east,text=white,font=\sffamily\itshape\large] at ([xshift=-5mm,yshift=-0.62cm]frame.north east) {Fonctions};}]
\vspace*{1.75cm}
\begin{minipage}[t][6.9cm][c]{\linewidth}\LARGE\centering
\begin{minipage}{0.95\linewidth}\raggedright\linespread{1.25}\selectfont De quel type est la fonction $f:x\mapsto 5$~?\end{minipage}
\end{minipage}

\vfill
\begin{tcolorbox}[enhanced,nobeforeafter,width=7.1cm,height=2.1cm,valign=center,colback=cfon!6,colframe=cfon,arc=2.5mm,boxrule=0.9pt,left=3mm]{\Large\bfseries\color{cfon}a)}\quad\begin{minipage}[c]{5.3cm}\raggedright\Large affine non linéaire\end{minipage}\end{tcolorbox}\hfill\begin{tcolorbox}[enhanced,nobeforeafter,width=7.1cm,height=2.1cm,valign=center,colback=cfon!6,colframe=cfon,arc=2.5mm,boxrule=0.9pt,left=3mm]{\Large\bfseries\color{cfon}b)}\quad\begin{minipage}[c]{5.3cm}\raggedright\Large ni affine ni linéaire\end{minipage}\end{tcolorbox}\par\vspace{2.5mm}
\begin{tcolorbox}[enhanced,nobeforeafter,width=7.1cm,height=2.1cm,valign=center,colback=cfon!6,colframe=cfon,arc=2.5mm,boxrule=0.9pt,left=3mm]{\Large\bfseries\color{cfon}c)}\quad\begin{minipage}[c]{5.3cm}\raggedright\Large constante\end{minipage}\end{tcolorbox}\hfill\begin{tcolorbox}[enhanced,nobeforeafter,width=7.1cm,height=2.1cm,valign=center,colback=cfon!6,colframe=cfon,arc=2.5mm,boxrule=0.9pt,left=3mm]{\Large\bfseries\color{cfon}d)}\quad\begin{minipage}[c]{5.3cm}\raggedright\Large linéaire\end{minipage}\end{tcolorbox}
\end{tcolorbox}
\newpage
\begin{tcolorbox}[enhanced,width=16cm,height=14cm,colback=white,colframe=cgeo,boxrule=2pt,arc=0mm,
left=6mm,right=6mm,top=0mm,bottom=4mm,sharp corners,
overlay={\fill[cgeo] (frame.north west) rectangle ([yshift=-1.25cm]frame.north east);
\node[anchor=west,text=white,font=\sffamily\bfseries\LARGE] at ([xshift=5mm,yshift=-0.62cm]frame.north west) {Question 49 B};
\node[anchor=east,text=white,font=\sffamily\itshape\large] at ([xshift=-5mm,yshift=-0.62cm]frame.north east) {Géométrie};}]
\vspace*{1.75cm}
\begin{minipage}[t][6.9cm][c]{\linewidth}\Large\centering
\begin{minipage}{0.95\linewidth}\raggedright\linespread{1.25}\selectfont Les droites $(AC)$ et $(BD)$ se coupent en $O$, et les droites $(AB)$ et $(CD)$ sont parallèles. $OA=4$ cm, $OC=10$ cm et $AB=7$ cm. Combien mesure $CD$~?\end{minipage}
\end{minipage}

\vfill
\begin{tcolorbox}[enhanced,nobeforeafter,width=7.1cm,height=2.1cm,valign=center,colback=cgeo!6,colframe=cgeo,arc=2.5mm,boxrule=0.9pt,left=3mm]{\Large\bfseries\color{cgeo}a)}\quad\begin{minipage}[c]{5.3cm}\raggedright\Large $2{,}8~\text{cm}$\end{minipage}\end{tcolorbox}\hfill\begin{tcolorbox}[enhanced,nobeforeafter,width=7.1cm,height=2.1cm,valign=center,colback=cgeo!6,colframe=cgeo,arc=2.5mm,boxrule=0.9pt,left=3mm]{\Large\bfseries\color{cgeo}b)}\quad\begin{minipage}[c]{5.3cm}\raggedright\Large $\dfrac{40}{7}~\text{cm}$\end{minipage}\end{tcolorbox}\par\vspace{2.5mm}
\begin{tcolorbox}[enhanced,nobeforeafter,width=7.1cm,height=2.1cm,valign=center,colback=cgeo!6,colframe=cgeo,arc=2.5mm,boxrule=0.9pt,left=3mm]{\Large\bfseries\color{cgeo}c)}\quad\begin{minipage}[c]{5.3cm}\raggedright\Large $17{,}5~\text{cm}$\end{minipage}\end{tcolorbox}\hfill\begin{tcolorbox}[enhanced,nobeforeafter,width=7.1cm,height=2.1cm,valign=center,colback=cgeo!6,colframe=cgeo,arc=2.5mm,boxrule=0.9pt,left=3mm]{\Large\bfseries\color{cgeo}d)}\quad\begin{minipage}[c]{5.3cm}\raggedright\Large $13~\text{cm}$\end{minipage}\end{tcolorbox}
\end{tcolorbox}
\newpage
\begin{tcolorbox}[enhanced,width=16cm,height=14cm,colback=white,colframe=cgeo,boxrule=2pt,arc=0mm,
left=6mm,right=6mm,top=0mm,bottom=4mm,sharp corners,
overlay={\fill[cgeo] (frame.north west) rectangle ([yshift=-1.25cm]frame.north east);
\node[anchor=west,text=white,font=\sffamily\bfseries\LARGE] at ([xshift=5mm,yshift=-0.62cm]frame.north west) {Question 50 B};
\node[anchor=east,text=white,font=\sffamily\itshape\large] at ([xshift=-5mm,yshift=-0.62cm]frame.north east) {Géométrie};}]
\vspace*{1.75cm}
\begin{minipage}[t][6.9cm][c]{\linewidth}\Large\centering
\begin{minipage}{0.95\linewidth}\raggedright\linespread{1.25}\selectfont Les droites $(AC)$ et $(BD)$ se coupent en $O$, et les droites $(AB)$ et $(CD)$ sont parallèles. $OA=5$ cm, $OC=3$ cm et $AB=8$ cm. Combien mesure $CD$~?\end{minipage}
\end{minipage}

\vfill
\begin{tcolorbox}[enhanced,nobeforeafter,width=7.1cm,height=2.1cm,valign=center,colback=cgeo!6,colframe=cgeo,arc=2.5mm,boxrule=0.9pt,left=3mm]{\Large\bfseries\color{cgeo}a)}\quad\begin{minipage}[c]{5.3cm}\raggedright\Large $6~\text{cm}$\end{minipage}\end{tcolorbox}\hfill\begin{tcolorbox}[enhanced,nobeforeafter,width=7.1cm,height=2.1cm,valign=center,colback=cgeo!6,colframe=cgeo,arc=2.5mm,boxrule=0.9pt,left=3mm]{\Large\bfseries\color{cgeo}b)}\quad\begin{minipage}[c]{5.3cm}\raggedright\Large $4{,}8~\text{cm}$\end{minipage}\end{tcolorbox}\par\vspace{2.5mm}
\begin{tcolorbox}[enhanced,nobeforeafter,width=7.1cm,height=2.1cm,valign=center,colback=cgeo!6,colframe=cgeo,arc=2.5mm,boxrule=0.9pt,left=3mm]{\Large\bfseries\color{cgeo}c)}\quad\begin{minipage}[c]{5.3cm}\raggedright\Large $1{,}875~\text{cm}$\end{minipage}\end{tcolorbox}\hfill\begin{tcolorbox}[enhanced,nobeforeafter,width=7.1cm,height=2.1cm,valign=center,colback=cgeo!6,colframe=cgeo,arc=2.5mm,boxrule=0.9pt,left=3mm]{\Large\bfseries\color{cgeo}d)}\quad\begin{minipage}[c]{5.3cm}\raggedright\Large $\dfrac{40}{3}~\text{cm}$\end{minipage}\end{tcolorbox}
\end{tcolorbox}
\newpage
\begin{tcolorbox}[enhanced,width=16cm,height=14cm,colback=white,colframe=cgeo,boxrule=2pt,arc=0mm,
left=6mm,right=6mm,top=0mm,bottom=4mm,sharp corners,
overlay={\fill[cgeo] (frame.north west) rectangle ([yshift=-1.25cm]frame.north east);
\node[anchor=west,text=white,font=\sffamily\bfseries\LARGE] at ([xshift=5mm,yshift=-0.62cm]frame.north west) {Question 51 B};
\node[anchor=east,text=white,font=\sffamily\itshape\large] at ([xshift=-5mm,yshift=-0.62cm]frame.north east) {Géométrie};}]
\vspace*{1.75cm}
\begin{minipage}[t][6.9cm][c]{\linewidth}\Large\centering
\begin{minipage}{0.95\linewidth}\raggedright\linespread{1.25}\selectfont Les points $A$, $M$, $B$ sont alignés dans cet ordre, ainsi que $A$, $N$, $C$. $AM=5$, $AB=15$, $AN=6$ et $AC=18$. Les droites $(MN)$ et $(BC)$ sont-elles parallèles~?\end{minipage}
\end{minipage}

\vfill
\begin{tcolorbox}[enhanced,nobeforeafter,width=7.1cm,height=2.1cm,valign=center,colback=cgeo!6,colframe=cgeo,arc=2.5mm,boxrule=0.9pt,left=3mm]{\Large\bfseries\color{cgeo}a)}\quad\begin{minipage}[c]{5.3cm}\raggedright\Large oui\end{minipage}\end{tcolorbox}\hfill\begin{tcolorbox}[enhanced,nobeforeafter,width=7.1cm,height=2.1cm,valign=center,colback=cgeo!6,colframe=cgeo,arc=2.5mm,boxrule=0.9pt,left=3mm]{\Large\bfseries\color{cgeo}b)}\quad\begin{minipage}[c]{5.3cm}\raggedright\Large non\end{minipage}\end{tcolorbox}\par\vspace{2.5mm}
\begin{tcolorbox}[enhanced,nobeforeafter,width=7.1cm,height=2.1cm,valign=center,colback=cgeo!6,colframe=cgeo,arc=2.5mm,boxrule=0.9pt,left=3mm]{\Large\bfseries\color{cgeo}c)}\quad\begin{minipage}[c]{5.3cm}\raggedright\large seulement si le triangle est rectangle\end{minipage}\end{tcolorbox}\hfill\begin{tcolorbox}[enhanced,nobeforeafter,width=7.1cm,height=2.1cm,valign=center,colback=cgeo!6,colframe=cgeo,arc=2.5mm,boxrule=0.9pt,left=3mm]{\Large\bfseries\color{cgeo}d)}\quad\begin{minipage}[c]{5.3cm}\raggedright\Large on ne peut pas savoir\end{minipage}\end{tcolorbox}
\end{tcolorbox}
\newpage
\begin{tcolorbox}[enhanced,width=16cm,height=14cm,colback=white,colframe=cgeo,boxrule=2pt,arc=0mm,
left=6mm,right=6mm,top=0mm,bottom=4mm,sharp corners,
overlay={\fill[cgeo] (frame.north west) rectangle ([yshift=-1.25cm]frame.north east);
\node[anchor=west,text=white,font=\sffamily\bfseries\LARGE] at ([xshift=5mm,yshift=-0.62cm]frame.north west) {Question 52 B};
\node[anchor=east,text=white,font=\sffamily\itshape\large] at ([xshift=-5mm,yshift=-0.62cm]frame.north east) {Géométrie};}]
\vspace*{1.75cm}
\begin{minipage}[t][6.9cm][c]{\linewidth}\Large\centering
\begin{minipage}{0.95\linewidth}\raggedright\linespread{1.25}\selectfont Les points $A$, $M$, $B$ sont alignés dans cet ordre, ainsi que $A$, $N$, $C$. $AM=4$, $AB=10$, $AN=6$ et $AC=14$. Les droites $(MN)$ et $(BC)$ sont-elles parallèles~?\end{minipage}
\end{minipage}

\vfill
\begin{tcolorbox}[enhanced,nobeforeafter,width=7.1cm,height=2.1cm,valign=center,colback=cgeo!6,colframe=cgeo,arc=2.5mm,boxrule=0.9pt,left=3mm]{\Large\bfseries\color{cgeo}a)}\quad\begin{minipage}[c]{5.3cm}\raggedright\Large oui\end{minipage}\end{tcolorbox}\hfill\begin{tcolorbox}[enhanced,nobeforeafter,width=7.1cm,height=2.1cm,valign=center,colback=cgeo!6,colframe=cgeo,arc=2.5mm,boxrule=0.9pt,left=3mm]{\Large\bfseries\color{cgeo}b)}\quad\begin{minipage}[c]{5.3cm}\raggedright\Large on ne peut pas savoir\end{minipage}\end{tcolorbox}\par\vspace{2.5mm}
\begin{tcolorbox}[enhanced,nobeforeafter,width=7.1cm,height=2.1cm,valign=center,colback=cgeo!6,colframe=cgeo,arc=2.5mm,boxrule=0.9pt,left=3mm]{\Large\bfseries\color{cgeo}c)}\quad\begin{minipage}[c]{5.3cm}\raggedright\Large non\end{minipage}\end{tcolorbox}\hfill\begin{tcolorbox}[enhanced,nobeforeafter,width=7.1cm,height=2.1cm,valign=center,colback=cgeo!6,colframe=cgeo,arc=2.5mm,boxrule=0.9pt,left=3mm]{\Large\bfseries\color{cgeo}d)}\quad\begin{minipage}[c]{5.3cm}\raggedright\large seulement si le triangle est rectangle\end{minipage}\end{tcolorbox}
\end{tcolorbox}
\newpage
\begin{tcolorbox}[enhanced,width=16cm,height=14cm,colback=white,colframe=cgeo,boxrule=2pt,arc=0mm,
left=6mm,right=6mm,top=0mm,bottom=4mm,sharp corners,
overlay={\fill[cgeo] (frame.north west) rectangle ([yshift=-1.25cm]frame.north east);
\node[anchor=west,text=white,font=\sffamily\bfseries\LARGE] at ([xshift=5mm,yshift=-0.62cm]frame.north west) {Question 53 B};
\node[anchor=east,text=white,font=\sffamily\itshape\large] at ([xshift=-5mm,yshift=-0.62cm]frame.north east) {Géométrie};}]
\vspace*{1.75cm}
\begin{minipage}[t][6.9cm][c]{\linewidth}\LARGE\centering
\begin{minipage}{0.95\linewidth}\raggedright\linespread{1.25}\selectfont Le triangle $ABC$ est rectangle en $A$, avec $BC=20$ cm et $\sin\widehat{ABC}=0{,}45$. Combien mesure $AC$~?\end{minipage}
\end{minipage}

\vfill
\begin{tcolorbox}[enhanced,nobeforeafter,width=7.1cm,height=2.1cm,valign=center,colback=cgeo!6,colframe=cgeo,arc=2.5mm,boxrule=0.9pt,left=3mm]{\Large\bfseries\color{cgeo}a)}\quad\begin{minipage}[c]{5.3cm}\raggedright\Large $11~\text{cm}$\end{minipage}\end{tcolorbox}\hfill\begin{tcolorbox}[enhanced,nobeforeafter,width=7.1cm,height=2.1cm,valign=center,colback=cgeo!6,colframe=cgeo,arc=2.5mm,boxrule=0.9pt,left=3mm]{\Large\bfseries\color{cgeo}b)}\quad\begin{minipage}[c]{5.3cm}\raggedright\Large $18~\text{cm}$\end{minipage}\end{tcolorbox}\par\vspace{2.5mm}
\begin{tcolorbox}[enhanced,nobeforeafter,width=7.1cm,height=2.1cm,valign=center,colback=cgeo!6,colframe=cgeo,arc=2.5mm,boxrule=0.9pt,left=3mm]{\Large\bfseries\color{cgeo}c)}\quad\begin{minipage}[c]{5.3cm}\raggedright\Large $4{,}5~\text{cm}$\end{minipage}\end{tcolorbox}\hfill\begin{tcolorbox}[enhanced,nobeforeafter,width=7.1cm,height=2.1cm,valign=center,colback=cgeo!6,colframe=cgeo,arc=2.5mm,boxrule=0.9pt,left=3mm]{\Large\bfseries\color{cgeo}d)}\quad\begin{minipage}[c]{5.3cm}\raggedright\Large $9~\text{cm}$\end{minipage}\end{tcolorbox}
\end{tcolorbox}
\newpage
\begin{tcolorbox}[enhanced,width=16cm,height=14cm,colback=white,colframe=cgeo,boxrule=2pt,arc=0mm,
left=6mm,right=6mm,top=0mm,bottom=4mm,sharp corners,
overlay={\fill[cgeo] (frame.north west) rectangle ([yshift=-1.25cm]frame.north east);
\node[anchor=west,text=white,font=\sffamily\bfseries\LARGE] at ([xshift=5mm,yshift=-0.62cm]frame.north west) {Question 54 B};
\node[anchor=east,text=white,font=\sffamily\itshape\large] at ([xshift=-5mm,yshift=-0.62cm]frame.north east) {Géométrie};}]
\vspace*{1.75cm}
\begin{minipage}[t][6.9cm][c]{\linewidth}\LARGE\centering
\begin{minipage}{0.95\linewidth}\raggedright\linespread{1.25}\selectfont Le triangle $ABC$ est rectangle en $A$, avec $BC=12$ cm et $\sin\widehat{ABC}=0{,}75$. Combien mesure $AC$~?\end{minipage}
\end{minipage}

\vfill
\begin{tcolorbox}[enhanced,nobeforeafter,width=7.1cm,height=2.1cm,valign=center,colback=cgeo!6,colframe=cgeo,arc=2.5mm,boxrule=0.9pt,left=3mm]{\Large\bfseries\color{cgeo}a)}\quad\begin{minipage}[c]{5.3cm}\raggedright\Large $18~\text{cm}$\end{minipage}\end{tcolorbox}\hfill\begin{tcolorbox}[enhanced,nobeforeafter,width=7.1cm,height=2.1cm,valign=center,colback=cgeo!6,colframe=cgeo,arc=2.5mm,boxrule=0.9pt,left=3mm]{\Large\bfseries\color{cgeo}b)}\quad\begin{minipage}[c]{5.3cm}\raggedright\Large $3~\text{cm}$\end{minipage}\end{tcolorbox}\par\vspace{2.5mm}
\begin{tcolorbox}[enhanced,nobeforeafter,width=7.1cm,height=2.1cm,valign=center,colback=cgeo!6,colframe=cgeo,arc=2.5mm,boxrule=0.9pt,left=3mm]{\Large\bfseries\color{cgeo}c)}\quad\begin{minipage}[c]{5.3cm}\raggedright\Large $9~\text{cm}$\end{minipage}\end{tcolorbox}\hfill\begin{tcolorbox}[enhanced,nobeforeafter,width=7.1cm,height=2.1cm,valign=center,colback=cgeo!6,colframe=cgeo,arc=2.5mm,boxrule=0.9pt,left=3mm]{\Large\bfseries\color{cgeo}d)}\quad\begin{minipage}[c]{5.3cm}\raggedright\Large $4{,}5~\text{cm}$\end{minipage}\end{tcolorbox}
\end{tcolorbox}
\newpage
\begin{tcolorbox}[enhanced,width=16cm,height=14cm,colback=white,colframe=cgeo,boxrule=2pt,arc=0mm,
left=6mm,right=6mm,top=0mm,bottom=4mm,sharp corners,
overlay={\fill[cgeo] (frame.north west) rectangle ([yshift=-1.25cm]frame.north east);
\node[anchor=west,text=white,font=\sffamily\bfseries\LARGE] at ([xshift=5mm,yshift=-0.62cm]frame.north west) {Question 55 B};
\node[anchor=east,text=white,font=\sffamily\itshape\large] at ([xshift=-5mm,yshift=-0.62cm]frame.north east) {Géométrie};}]
\vspace*{1.75cm}
\begin{minipage}[t][6.9cm][c]{\linewidth}\LARGE\centering
\begin{minipage}{0.95\linewidth}\raggedright\linespread{1.25}\selectfont Le triangle $ABC$ est rectangle en $A$, avec $AB=8$ cm et $AC=6$ cm. Que vaut $\tan\widehat{ABC}$~?\end{minipage}
\end{minipage}

\vfill
\begin{tcolorbox}[enhanced,nobeforeafter,width=7.1cm,height=2.1cm,valign=center,colback=cgeo!6,colframe=cgeo,arc=2.5mm,boxrule=0.9pt,left=3mm]{\Large\bfseries\color{cgeo}a)}\quad\begin{minipage}[c]{5.3cm}\raggedright\Large $\dfrac{7}{4}$\end{minipage}\end{tcolorbox}\hfill\begin{tcolorbox}[enhanced,nobeforeafter,width=7.1cm,height=2.1cm,valign=center,colback=cgeo!6,colframe=cgeo,arc=2.5mm,boxrule=0.9pt,left=3mm]{\Large\bfseries\color{cgeo}b)}\quad\begin{minipage}[c]{5.3cm}\raggedright\Large $\dfrac{3}{4}$\end{minipage}\end{tcolorbox}\par\vspace{2.5mm}
\begin{tcolorbox}[enhanced,nobeforeafter,width=7.1cm,height=2.1cm,valign=center,colback=cgeo!6,colframe=cgeo,arc=2.5mm,boxrule=0.9pt,left=3mm]{\Large\bfseries\color{cgeo}c)}\quad\begin{minipage}[c]{5.3cm}\raggedright\Large $\dfrac{4}{7}$\end{minipage}\end{tcolorbox}\hfill\begin{tcolorbox}[enhanced,nobeforeafter,width=7.1cm,height=2.1cm,valign=center,colback=cgeo!6,colframe=cgeo,arc=2.5mm,boxrule=0.9pt,left=3mm]{\Large\bfseries\color{cgeo}d)}\quad\begin{minipage}[c]{5.3cm}\raggedright\Large $\dfrac{4}{3}$\end{minipage}\end{tcolorbox}
\end{tcolorbox}
\newpage
\begin{tcolorbox}[enhanced,width=16cm,height=14cm,colback=white,colframe=cgeo,boxrule=2pt,arc=0mm,
left=6mm,right=6mm,top=0mm,bottom=4mm,sharp corners,
overlay={\fill[cgeo] (frame.north west) rectangle ([yshift=-1.25cm]frame.north east);
\node[anchor=west,text=white,font=\sffamily\bfseries\LARGE] at ([xshift=5mm,yshift=-0.62cm]frame.north west) {Question 56 B};
\node[anchor=east,text=white,font=\sffamily\itshape\large] at ([xshift=-5mm,yshift=-0.62cm]frame.north east) {Géométrie};}]
\vspace*{1.75cm}
\begin{minipage}[t][6.9cm][c]{\linewidth}\LARGE\centering
\begin{minipage}{0.95\linewidth}\raggedright\linespread{1.25}\selectfont Le triangle $ABC$ est rectangle en $A$, avec $AB=7$ cm et $AC=7$ cm. Que vaut $\tan\widehat{ABC}$~?\end{minipage}
\end{minipage}

\vfill
\begin{tcolorbox}[enhanced,nobeforeafter,width=7.1cm,height=2.1cm,valign=center,colback=cgeo!6,colframe=cgeo,arc=2.5mm,boxrule=0.9pt,left=3mm]{\Large\bfseries\color{cgeo}a)}\quad\begin{minipage}[c]{5.3cm}\raggedright\Large $0$\end{minipage}\end{tcolorbox}\hfill\begin{tcolorbox}[enhanced,nobeforeafter,width=7.1cm,height=2.1cm,valign=center,colback=cgeo!6,colframe=cgeo,arc=2.5mm,boxrule=0.9pt,left=3mm]{\Large\bfseries\color{cgeo}b)}\quad\begin{minipage}[c]{5.3cm}\raggedright\Large $2$\end{minipage}\end{tcolorbox}\par\vspace{2.5mm}
\begin{tcolorbox}[enhanced,nobeforeafter,width=7.1cm,height=2.1cm,valign=center,colback=cgeo!6,colframe=cgeo,arc=2.5mm,boxrule=0.9pt,left=3mm]{\Large\bfseries\color{cgeo}c)}\quad\begin{minipage}[c]{5.3cm}\raggedright\Large $1$\end{minipage}\end{tcolorbox}\hfill\begin{tcolorbox}[enhanced,nobeforeafter,width=7.1cm,height=2.1cm,valign=center,colback=cgeo!6,colframe=cgeo,arc=2.5mm,boxrule=0.9pt,left=3mm]{\Large\bfseries\color{cgeo}d)}\quad\begin{minipage}[c]{5.3cm}\raggedright\Large $\dfrac{1}{2}$\end{minipage}\end{tcolorbox}
\end{tcolorbox}
\newpage
\begin{tcolorbox}[enhanced,width=16cm,height=14cm,colback=white,colframe=cgeo,boxrule=2pt,arc=0mm,
left=6mm,right=6mm,top=0mm,bottom=4mm,sharp corners,
overlay={\fill[cgeo] (frame.north west) rectangle ([yshift=-1.25cm]frame.north east);
\node[anchor=west,text=white,font=\sffamily\bfseries\LARGE] at ([xshift=5mm,yshift=-0.62cm]frame.north west) {Question 57 B};
\node[anchor=east,text=white,font=\sffamily\itshape\large] at ([xshift=-5mm,yshift=-0.62cm]frame.north east) {Géométrie};}]
\vspace*{1.75cm}
\begin{minipage}[t][6.9cm][c]{\linewidth}\LARGE\centering
\begin{minipage}{0.95\linewidth}\raggedright\linespread{1.25}\selectfont Une homothétie de rapport $\dfrac{1}{2}$ transforme le segment $[AB]$, de longueur $8$ cm, en $[A'B']$. Combien mesure $A'B'$~?\end{minipage}
\end{minipage}

\vfill
\begin{tcolorbox}[enhanced,nobeforeafter,width=7.1cm,height=2.1cm,valign=center,colback=cgeo!6,colframe=cgeo,arc=2.5mm,boxrule=0.9pt,left=3mm]{\Large\bfseries\color{cgeo}a)}\quad\begin{minipage}[c]{5.3cm}\raggedright\Large $8{,}5~\text{cm}$\end{minipage}\end{tcolorbox}\hfill\begin{tcolorbox}[enhanced,nobeforeafter,width=7.1cm,height=2.1cm,valign=center,colback=cgeo!6,colframe=cgeo,arc=2.5mm,boxrule=0.9pt,left=3mm]{\Large\bfseries\color{cgeo}b)}\quad\begin{minipage}[c]{5.3cm}\raggedright\Large $16~\text{cm}$\end{minipage}\end{tcolorbox}\par\vspace{2.5mm}
\begin{tcolorbox}[enhanced,nobeforeafter,width=7.1cm,height=2.1cm,valign=center,colback=cgeo!6,colframe=cgeo,arc=2.5mm,boxrule=0.9pt,left=3mm]{\Large\bfseries\color{cgeo}c)}\quad\begin{minipage}[c]{5.3cm}\raggedright\Large $5~\text{cm}$\end{minipage}\end{tcolorbox}\hfill\begin{tcolorbox}[enhanced,nobeforeafter,width=7.1cm,height=2.1cm,valign=center,colback=cgeo!6,colframe=cgeo,arc=2.5mm,boxrule=0.9pt,left=3mm]{\Large\bfseries\color{cgeo}d)}\quad\begin{minipage}[c]{5.3cm}\raggedright\Large $4~\text{cm}$\end{minipage}\end{tcolorbox}
\end{tcolorbox}
\newpage
\begin{tcolorbox}[enhanced,width=16cm,height=14cm,colback=white,colframe=cgeo,boxrule=2pt,arc=0mm,
left=6mm,right=6mm,top=0mm,bottom=4mm,sharp corners,
overlay={\fill[cgeo] (frame.north west) rectangle ([yshift=-1.25cm]frame.north east);
\node[anchor=west,text=white,font=\sffamily\bfseries\LARGE] at ([xshift=5mm,yshift=-0.62cm]frame.north west) {Question 58 B};
\node[anchor=east,text=white,font=\sffamily\itshape\large] at ([xshift=-5mm,yshift=-0.62cm]frame.north east) {Géométrie};}]
\vspace*{1.75cm}
\begin{minipage}[t][6.9cm][c]{\linewidth}\LARGE\centering
\begin{minipage}{0.95\linewidth}\raggedright\linespread{1.25}\selectfont Une homothétie de rapport $-3$ transforme le segment $[AB]$, de longueur $4$ cm, en $[A'B']$. Combien mesure $A'B'$~?\end{minipage}
\end{minipage}

\vfill
\begin{tcolorbox}[enhanced,nobeforeafter,width=7.1cm,height=2.1cm,valign=center,colback=cgeo!6,colframe=cgeo,arc=2.5mm,boxrule=0.9pt,left=3mm]{\Large\bfseries\color{cgeo}a)}\quad\begin{minipage}[c]{5.3cm}\raggedright\Large $7~\text{cm}$\end{minipage}\end{tcolorbox}\hfill\begin{tcolorbox}[enhanced,nobeforeafter,width=7.1cm,height=2.1cm,valign=center,colback=cgeo!6,colframe=cgeo,arc=2.5mm,boxrule=0.9pt,left=3mm]{\Large\bfseries\color{cgeo}b)}\quad\begin{minipage}[c]{5.3cm}\raggedright\Large $12~\text{cm}$\end{minipage}\end{tcolorbox}\par\vspace{2.5mm}
\begin{tcolorbox}[enhanced,nobeforeafter,width=7.1cm,height=2.1cm,valign=center,colback=cgeo!6,colframe=cgeo,arc=2.5mm,boxrule=0.9pt,left=3mm]{\Large\bfseries\color{cgeo}c)}\quad\begin{minipage}[c]{5.3cm}\raggedright\Large $13~\text{cm}$\end{minipage}\end{tcolorbox}\hfill\begin{tcolorbox}[enhanced,nobeforeafter,width=7.1cm,height=2.1cm,valign=center,colback=cgeo!6,colframe=cgeo,arc=2.5mm,boxrule=0.9pt,left=3mm]{\Large\bfseries\color{cgeo}d)}\quad\begin{minipage}[c]{5.3cm}\raggedright\Large $\dfrac{4}{3}~\text{cm}$\end{minipage}\end{tcolorbox}
\end{tcolorbox}
\newpage
\begin{tcolorbox}[enhanced,width=16cm,height=14cm,colback=white,colframe=cgeo,boxrule=2pt,arc=0mm,
left=6mm,right=6mm,top=0mm,bottom=4mm,sharp corners,
overlay={\fill[cgeo] (frame.north west) rectangle ([yshift=-1.25cm]frame.north east);
\node[anchor=west,text=white,font=\sffamily\bfseries\LARGE] at ([xshift=5mm,yshift=-0.62cm]frame.north west) {Question 59 B};
\node[anchor=east,text=white,font=\sffamily\itshape\large] at ([xshift=-5mm,yshift=-0.62cm]frame.north east) {Géométrie};}]
\vspace*{1.75cm}
\begin{minipage}[t][6.9cm][c]{\linewidth}\LARGE\centering
\begin{minipage}{0.95\linewidth}\raggedright\linespread{1.25}\selectfont Quel est le volume exact d'une boule de rayon $1$ cm~?\end{minipage}
\end{minipage}

\vfill
\begin{tcolorbox}[enhanced,nobeforeafter,width=7.1cm,height=2.1cm,valign=center,colback=cgeo!6,colframe=cgeo,arc=2.5mm,boxrule=0.9pt,left=3mm]{\Large\bfseries\color{cgeo}a)}\quad\begin{minipage}[c]{5.3cm}\raggedright\large $\dfrac{8}{3}\pi~\text{cm}^3$\end{minipage}\end{tcolorbox}\hfill\begin{tcolorbox}[enhanced,nobeforeafter,width=7.1cm,height=2.1cm,valign=center,colback=cgeo!6,colframe=cgeo,arc=2.5mm,boxrule=0.9pt,left=3mm]{\Large\bfseries\color{cgeo}b)}\quad\begin{minipage}[c]{5.3cm}\raggedright\large $\dfrac{4}{3}\pi~\text{cm}^3$\end{minipage}\end{tcolorbox}\par\vspace{2.5mm}
\begin{tcolorbox}[enhanced,nobeforeafter,width=7.1cm,height=2.1cm,valign=center,colback=cgeo!6,colframe=cgeo,arc=2.5mm,boxrule=0.9pt,left=3mm]{\Large\bfseries\color{cgeo}c)}\quad\begin{minipage}[c]{5.3cm}\raggedright\Large $4\pi~\text{cm}^3$\end{minipage}\end{tcolorbox}\hfill\begin{tcolorbox}[enhanced,nobeforeafter,width=7.1cm,height=2.1cm,valign=center,colback=cgeo!6,colframe=cgeo,arc=2.5mm,boxrule=0.9pt,left=3mm]{\Large\bfseries\color{cgeo}d)}\quad\begin{minipage}[c]{5.3cm}\raggedright\Large $\pi~\text{cm}^3$\end{minipage}\end{tcolorbox}
\end{tcolorbox}
\newpage
\begin{tcolorbox}[enhanced,width=16cm,height=14cm,colback=white,colframe=cgeo,boxrule=2pt,arc=0mm,
left=6mm,right=6mm,top=0mm,bottom=4mm,sharp corners,
overlay={\fill[cgeo] (frame.north west) rectangle ([yshift=-1.25cm]frame.north east);
\node[anchor=west,text=white,font=\sffamily\bfseries\LARGE] at ([xshift=5mm,yshift=-0.62cm]frame.north west) {Question 60 B};
\node[anchor=east,text=white,font=\sffamily\itshape\large] at ([xshift=-5mm,yshift=-0.62cm]frame.north east) {Géométrie};}]
\vspace*{1.75cm}
\begin{minipage}[t][6.9cm][c]{\linewidth}\LARGE\centering
\begin{minipage}{0.95\linewidth}\raggedright\linespread{1.25}\selectfont Quel est le volume exact d'une boule de rayon $9$ cm~?\end{minipage}
\end{minipage}

\vfill
\begin{tcolorbox}[enhanced,nobeforeafter,width=7.1cm,height=2.1cm,valign=center,colback=cgeo!6,colframe=cgeo,arc=2.5mm,boxrule=0.9pt,left=3mm]{\Large\bfseries\color{cgeo}a)}\quad\begin{minipage}[c]{5.3cm}\raggedright\Large $729\pi~\text{cm}^3$\end{minipage}\end{tcolorbox}\hfill\begin{tcolorbox}[enhanced,nobeforeafter,width=7.1cm,height=2.1cm,valign=center,colback=cgeo!6,colframe=cgeo,arc=2.5mm,boxrule=0.9pt,left=3mm]{\Large\bfseries\color{cgeo}b)}\quad\begin{minipage}[c]{5.3cm}\raggedright\Large $972\pi~\text{cm}^3$\end{minipage}\end{tcolorbox}\par\vspace{2.5mm}
\begin{tcolorbox}[enhanced,nobeforeafter,width=7.1cm,height=2.1cm,valign=center,colback=cgeo!6,colframe=cgeo,arc=2.5mm,boxrule=0.9pt,left=3mm]{\Large\bfseries\color{cgeo}c)}\quad\begin{minipage}[c]{5.3cm}\raggedright\Large $108\pi~\text{cm}^3$\end{minipage}\end{tcolorbox}\hfill\begin{tcolorbox}[enhanced,nobeforeafter,width=7.1cm,height=2.1cm,valign=center,colback=cgeo!6,colframe=cgeo,arc=2.5mm,boxrule=0.9pt,left=3mm]{\Large\bfseries\color{cgeo}d)}\quad\begin{minipage}[c]{5.3cm}\raggedright\Large $324\pi~\text{cm}^3$\end{minipage}\end{tcolorbox}
\end{tcolorbox}
\newpage
\begin{tcolorbox}[enhanced,width=16cm,height=14cm,colback=white,colframe=cgeo,boxrule=2pt,arc=0mm,
left=6mm,right=6mm,top=0mm,bottom=4mm,sharp corners,
overlay={\fill[cgeo] (frame.north west) rectangle ([yshift=-1.25cm]frame.north east);
\node[anchor=west,text=white,font=\sffamily\bfseries\LARGE] at ([xshift=5mm,yshift=-0.62cm]frame.north west) {Question 61 B};
\node[anchor=east,text=white,font=\sffamily\itshape\large] at ([xshift=-5mm,yshift=-0.62cm]frame.north east) {Géométrie};}]
\vspace*{1.75cm}
\begin{minipage}[t][6.9cm][c]{\linewidth}\LARGE\centering
\begin{minipage}{0.95\linewidth}\raggedright\linespread{1.25}\selectfont Quelle est l'aire exacte d'une sphère de rayon $10$ cm~?\end{minipage}
\end{minipage}

\vfill
\begin{tcolorbox}[enhanced,nobeforeafter,width=7.1cm,height=2.1cm,valign=center,colback=cgeo!6,colframe=cgeo,arc=2.5mm,boxrule=0.9pt,left=3mm]{\Large\bfseries\color{cgeo}a)}\quad\begin{minipage}[c]{5.3cm}\raggedright\Large $400\pi~\text{cm}^2$\end{minipage}\end{tcolorbox}\hfill\begin{tcolorbox}[enhanced,nobeforeafter,width=7.1cm,height=2.1cm,valign=center,colback=cgeo!6,colframe=cgeo,arc=2.5mm,boxrule=0.9pt,left=3mm]{\Large\bfseries\color{cgeo}b)}\quad\begin{minipage}[c]{5.3cm}\raggedright\large $\dfrac{4000}{3}\pi~\text{cm}^2$\end{minipage}\end{tcolorbox}\par\vspace{2.5mm}
\begin{tcolorbox}[enhanced,nobeforeafter,width=7.1cm,height=2.1cm,valign=center,colback=cgeo!6,colframe=cgeo,arc=2.5mm,boxrule=0.9pt,left=3mm]{\Large\bfseries\color{cgeo}c)}\quad\begin{minipage}[c]{5.3cm}\raggedright\Large $200\pi~\text{cm}^2$\end{minipage}\end{tcolorbox}\hfill\begin{tcolorbox}[enhanced,nobeforeafter,width=7.1cm,height=2.1cm,valign=center,colback=cgeo!6,colframe=cgeo,arc=2.5mm,boxrule=0.9pt,left=3mm]{\Large\bfseries\color{cgeo}d)}\quad\begin{minipage}[c]{5.3cm}\raggedright\Large $100\pi~\text{cm}^2$\end{minipage}\end{tcolorbox}
\end{tcolorbox}
\newpage
\begin{tcolorbox}[enhanced,width=16cm,height=14cm,colback=white,colframe=cgeo,boxrule=2pt,arc=0mm,
left=6mm,right=6mm,top=0mm,bottom=4mm,sharp corners,
overlay={\fill[cgeo] (frame.north west) rectangle ([yshift=-1.25cm]frame.north east);
\node[anchor=west,text=white,font=\sffamily\bfseries\LARGE] at ([xshift=5mm,yshift=-0.62cm]frame.north west) {Question 62 B};
\node[anchor=east,text=white,font=\sffamily\itshape\large] at ([xshift=-5mm,yshift=-0.62cm]frame.north east) {Géométrie};}]
\vspace*{1.75cm}
\begin{minipage}[t][6.9cm][c]{\linewidth}\LARGE\centering
\begin{minipage}{0.95\linewidth}\raggedright\linespread{1.25}\selectfont Quelle est l'aire exacte d'une sphère de rayon $1$ cm~?\end{minipage}
\end{minipage}

\vfill
\begin{tcolorbox}[enhanced,nobeforeafter,width=7.1cm,height=2.1cm,valign=center,colback=cgeo!6,colframe=cgeo,arc=2.5mm,boxrule=0.9pt,left=3mm]{\Large\bfseries\color{cgeo}a)}\quad\begin{minipage}[c]{5.3cm}\raggedright\Large $\pi~\text{cm}^2$\end{minipage}\end{tcolorbox}\hfill\begin{tcolorbox}[enhanced,nobeforeafter,width=7.1cm,height=2.1cm,valign=center,colback=cgeo!6,colframe=cgeo,arc=2.5mm,boxrule=0.9pt,left=3mm]{\Large\bfseries\color{cgeo}b)}\quad\begin{minipage}[c]{5.3cm}\raggedright\Large $4\pi~\text{cm}^2$\end{minipage}\end{tcolorbox}\par\vspace{2.5mm}
\begin{tcolorbox}[enhanced,nobeforeafter,width=7.1cm,height=2.1cm,valign=center,colback=cgeo!6,colframe=cgeo,arc=2.5mm,boxrule=0.9pt,left=3mm]{\Large\bfseries\color{cgeo}c)}\quad\begin{minipage}[c]{5.3cm}\raggedright\large $\dfrac{4}{3}\pi~\text{cm}^2$\end{minipage}\end{tcolorbox}\hfill\begin{tcolorbox}[enhanced,nobeforeafter,width=7.1cm,height=2.1cm,valign=center,colback=cgeo!6,colframe=cgeo,arc=2.5mm,boxrule=0.9pt,left=3mm]{\Large\bfseries\color{cgeo}d)}\quad\begin{minipage}[c]{5.3cm}\raggedright\Large $2\pi~\text{cm}^2$\end{minipage}\end{tcolorbox}
\end{tcolorbox}
\newpage
\begin{tcolorbox}[enhanced,width=16cm,height=14cm,colback=white,colframe=cgeo,boxrule=2pt,arc=0mm,
left=6mm,right=6mm,top=0mm,bottom=4mm,sharp corners,
overlay={\fill[cgeo] (frame.north west) rectangle ([yshift=-1.25cm]frame.north east);
\node[anchor=west,text=white,font=\sffamily\bfseries\LARGE] at ([xshift=5mm,yshift=-0.62cm]frame.north west) {Question 63 B};
\node[anchor=east,text=white,font=\sffamily\itshape\large] at ([xshift=-5mm,yshift=-0.62cm]frame.north east) {Géométrie};}]
\vspace*{1.75cm}
\begin{minipage}[t][6.9cm][c]{\linewidth}\LARGE\centering
\begin{minipage}{0.95\linewidth}\raggedright\linespread{1.25}\selectfont Un solide est réduit à l'échelle $\dfrac{1}{2}$. Par combien son volume est-il multiplié~?\end{minipage}
\end{minipage}

\vfill
\begin{tcolorbox}[enhanced,nobeforeafter,width=7.1cm,height=2.1cm,valign=center,colback=cgeo!6,colframe=cgeo,arc=2.5mm,boxrule=0.9pt,left=3mm]{\Large\bfseries\color{cgeo}a)}\quad\begin{minipage}[c]{5.3cm}\raggedright\Large $\dfrac{1}{2}$\end{minipage}\end{tcolorbox}\hfill\begin{tcolorbox}[enhanced,nobeforeafter,width=7.1cm,height=2.1cm,valign=center,colback=cgeo!6,colframe=cgeo,arc=2.5mm,boxrule=0.9pt,left=3mm]{\Large\bfseries\color{cgeo}b)}\quad\begin{minipage}[c]{5.3cm}\raggedright\Large $\dfrac{3}{2}$\end{minipage}\end{tcolorbox}\par\vspace{2.5mm}
\begin{tcolorbox}[enhanced,nobeforeafter,width=7.1cm,height=2.1cm,valign=center,colback=cgeo!6,colframe=cgeo,arc=2.5mm,boxrule=0.9pt,left=3mm]{\Large\bfseries\color{cgeo}c)}\quad\begin{minipage}[c]{5.3cm}\raggedright\Large $\dfrac{1}{8}$\end{minipage}\end{tcolorbox}\hfill\begin{tcolorbox}[enhanced,nobeforeafter,width=7.1cm,height=2.1cm,valign=center,colback=cgeo!6,colframe=cgeo,arc=2.5mm,boxrule=0.9pt,left=3mm]{\Large\bfseries\color{cgeo}d)}\quad\begin{minipage}[c]{5.3cm}\raggedright\Large $\dfrac{1}{4}$\end{minipage}\end{tcolorbox}
\end{tcolorbox}
\newpage
\begin{tcolorbox}[enhanced,width=16cm,height=14cm,colback=white,colframe=cgeo,boxrule=2pt,arc=0mm,
left=6mm,right=6mm,top=0mm,bottom=4mm,sharp corners,
overlay={\fill[cgeo] (frame.north west) rectangle ([yshift=-1.25cm]frame.north east);
\node[anchor=west,text=white,font=\sffamily\bfseries\LARGE] at ([xshift=5mm,yshift=-0.62cm]frame.north west) {Question 64 B};
\node[anchor=east,text=white,font=\sffamily\itshape\large] at ([xshift=-5mm,yshift=-0.62cm]frame.north east) {Géométrie};}]
\vspace*{1.75cm}
\begin{minipage}[t][6.9cm][c]{\linewidth}\LARGE\centering
\begin{minipage}{0.95\linewidth}\raggedright\linespread{1.25}\selectfont Un solide est agrandi à l'échelle $10$. Par combien son volume est-il multiplié~?\end{minipage}
\end{minipage}

\vfill
\begin{tcolorbox}[enhanced,nobeforeafter,width=7.1cm,height=2.1cm,valign=center,colback=cgeo!6,colframe=cgeo,arc=2.5mm,boxrule=0.9pt,left=3mm]{\Large\bfseries\color{cgeo}a)}\quad\begin{minipage}[c]{5.3cm}\raggedright\Large $1\,000$\end{minipage}\end{tcolorbox}\hfill\begin{tcolorbox}[enhanced,nobeforeafter,width=7.1cm,height=2.1cm,valign=center,colback=cgeo!6,colframe=cgeo,arc=2.5mm,boxrule=0.9pt,left=3mm]{\Large\bfseries\color{cgeo}b)}\quad\begin{minipage}[c]{5.3cm}\raggedright\Large $10$\end{minipage}\end{tcolorbox}\par\vspace{2.5mm}
\begin{tcolorbox}[enhanced,nobeforeafter,width=7.1cm,height=2.1cm,valign=center,colback=cgeo!6,colframe=cgeo,arc=2.5mm,boxrule=0.9pt,left=3mm]{\Large\bfseries\color{cgeo}c)}\quad\begin{minipage}[c]{5.3cm}\raggedright\Large $100$\end{minipage}\end{tcolorbox}\hfill\begin{tcolorbox}[enhanced,nobeforeafter,width=7.1cm,height=2.1cm,valign=center,colback=cgeo!6,colframe=cgeo,arc=2.5mm,boxrule=0.9pt,left=3mm]{\Large\bfseries\color{cgeo}d)}\quad\begin{minipage}[c]{5.3cm}\raggedright\Large $30$\end{minipage}\end{tcolorbox}
\end{tcolorbox}
\newpage
\begin{tcolorbox}[enhanced,width=16cm,height=14cm,colback=white,colframe=csta,boxrule=2pt,arc=0mm,
left=6mm,right=6mm,top=0mm,bottom=4mm,sharp corners,
overlay={\fill[csta] (frame.north west) rectangle ([yshift=-1.25cm]frame.north east);
\node[anchor=west,text=white,font=\sffamily\bfseries\LARGE] at ([xshift=5mm,yshift=-0.62cm]frame.north west) {Question 65 B};
\node[anchor=east,text=white,font=\sffamily\itshape\large] at ([xshift=-5mm,yshift=-0.62cm]frame.north east) {Statistiques et probabilités};}]
\vspace*{1.75cm}
\begin{minipage}[t][6.9cm][c]{\linewidth}\LARGE\centering
\begin{minipage}{0.95\linewidth}\raggedright\linespread{1.25}\selectfont Quelle est la médiane de la série $18\ ;\ 25\ ;\ 13\ ;\ 20\ ;\ 16$~?\end{minipage}
\end{minipage}

\vfill
\begin{tcolorbox}[enhanced,nobeforeafter,width=7.1cm,height=2.1cm,valign=center,colback=csta!6,colframe=csta,arc=2.5mm,boxrule=0.9pt,left=3mm]{\Large\bfseries\color{csta}a)}\quad\begin{minipage}[c]{5.3cm}\raggedright\Large $13$\end{minipage}\end{tcolorbox}\hfill\begin{tcolorbox}[enhanced,nobeforeafter,width=7.1cm,height=2.1cm,valign=center,colback=csta!6,colframe=csta,arc=2.5mm,boxrule=0.9pt,left=3mm]{\Large\bfseries\color{csta}b)}\quad\begin{minipage}[c]{5.3cm}\raggedright\Large $18{,}4$\end{minipage}\end{tcolorbox}\par\vspace{2.5mm}
\begin{tcolorbox}[enhanced,nobeforeafter,width=7.1cm,height=2.1cm,valign=center,colback=csta!6,colframe=csta,arc=2.5mm,boxrule=0.9pt,left=3mm]{\Large\bfseries\color{csta}c)}\quad\begin{minipage}[c]{5.3cm}\raggedright\Large $18$\end{minipage}\end{tcolorbox}\hfill\begin{tcolorbox}[enhanced,nobeforeafter,width=7.1cm,height=2.1cm,valign=center,colback=csta!6,colframe=csta,arc=2.5mm,boxrule=0.9pt,left=3mm]{\Large\bfseries\color{csta}d)}\quad\begin{minipage}[c]{5.3cm}\raggedright\Large $12$\end{minipage}\end{tcolorbox}
\end{tcolorbox}
\newpage
\begin{tcolorbox}[enhanced,width=16cm,height=14cm,colback=white,colframe=csta,boxrule=2pt,arc=0mm,
left=6mm,right=6mm,top=0mm,bottom=4mm,sharp corners,
overlay={\fill[csta] (frame.north west) rectangle ([yshift=-1.25cm]frame.north east);
\node[anchor=west,text=white,font=\sffamily\bfseries\LARGE] at ([xshift=5mm,yshift=-0.62cm]frame.north west) {Question 66 B};
\node[anchor=east,text=white,font=\sffamily\itshape\large] at ([xshift=-5mm,yshift=-0.62cm]frame.north east) {Statistiques et probabilités};}]
\vspace*{1.75cm}
\begin{minipage}[t][6.9cm][c]{\linewidth}\LARGE\centering
\begin{minipage}{0.95\linewidth}\raggedright\linespread{1.25}\selectfont Quelle est la médiane de la série $5\ ;\ 9\ ;\ 1\ ;\ 12$~?\end{minipage}
\end{minipage}

\vfill
\begin{tcolorbox}[enhanced,nobeforeafter,width=7.1cm,height=2.1cm,valign=center,colback=csta!6,colframe=csta,arc=2.5mm,boxrule=0.9pt,left=3mm]{\Large\bfseries\color{csta}a)}\quad\begin{minipage}[c]{5.3cm}\raggedright\Large $7$\end{minipage}\end{tcolorbox}\hfill\begin{tcolorbox}[enhanced,nobeforeafter,width=7.1cm,height=2.1cm,valign=center,colback=csta!6,colframe=csta,arc=2.5mm,boxrule=0.9pt,left=3mm]{\Large\bfseries\color{csta}b)}\quad\begin{minipage}[c]{5.3cm}\raggedright\Large $5$\end{minipage}\end{tcolorbox}\par\vspace{2.5mm}
\begin{tcolorbox}[enhanced,nobeforeafter,width=7.1cm,height=2.1cm,valign=center,colback=csta!6,colframe=csta,arc=2.5mm,boxrule=0.9pt,left=3mm]{\Large\bfseries\color{csta}c)}\quad\begin{minipage}[c]{5.3cm}\raggedright\Large $11$\end{minipage}\end{tcolorbox}\hfill\begin{tcolorbox}[enhanced,nobeforeafter,width=7.1cm,height=2.1cm,valign=center,colback=csta!6,colframe=csta,arc=2.5mm,boxrule=0.9pt,left=3mm]{\Large\bfseries\color{csta}d)}\quad\begin{minipage}[c]{5.3cm}\raggedright\Large $6{,}75$\end{minipage}\end{tcolorbox}
\end{tcolorbox}
\newpage
\begin{tcolorbox}[enhanced,width=16cm,height=14cm,colback=white,colframe=csta,boxrule=2pt,arc=0mm,
left=6mm,right=6mm,top=0mm,bottom=4mm,sharp corners,
overlay={\fill[csta] (frame.north west) rectangle ([yshift=-1.25cm]frame.north east);
\node[anchor=west,text=white,font=\sffamily\bfseries\LARGE] at ([xshift=5mm,yshift=-0.62cm]frame.north west) {Question 67 B};
\node[anchor=east,text=white,font=\sffamily\itshape\large] at ([xshift=-5mm,yshift=-0.62cm]frame.north east) {Statistiques et probabilités};}]
\vspace*{1.75cm}
\begin{minipage}[t][6.9cm][c]{\linewidth}\LARGE\centering
\begin{minipage}{0.95\linewidth}\raggedright\linespread{1.25}\selectfont Quelle est la moyenne de la série $20\ ;\ 17\ ;\ 11$~?\end{minipage}
\end{minipage}

\vfill
\begin{tcolorbox}[enhanced,nobeforeafter,width=7.1cm,height=2.1cm,valign=center,colback=csta!6,colframe=csta,arc=2.5mm,boxrule=0.9pt,left=3mm]{\Large\bfseries\color{csta}a)}\quad\begin{minipage}[c]{5.3cm}\raggedright\Large $48$\end{minipage}\end{tcolorbox}\hfill\begin{tcolorbox}[enhanced,nobeforeafter,width=7.1cm,height=2.1cm,valign=center,colback=csta!6,colframe=csta,arc=2.5mm,boxrule=0.9pt,left=3mm]{\Large\bfseries\color{csta}b)}\quad\begin{minipage}[c]{5.3cm}\raggedright\Large $16$\end{minipage}\end{tcolorbox}\par\vspace{2.5mm}
\begin{tcolorbox}[enhanced,nobeforeafter,width=7.1cm,height=2.1cm,valign=center,colback=csta!6,colframe=csta,arc=2.5mm,boxrule=0.9pt,left=3mm]{\Large\bfseries\color{csta}c)}\quad\begin{minipage}[c]{5.3cm}\raggedright\Large $17$\end{minipage}\end{tcolorbox}\hfill\begin{tcolorbox}[enhanced,nobeforeafter,width=7.1cm,height=2.1cm,valign=center,colback=csta!6,colframe=csta,arc=2.5mm,boxrule=0.9pt,left=3mm]{\Large\bfseries\color{csta}d)}\quad\begin{minipage}[c]{5.3cm}\raggedright\Large $20$\end{minipage}\end{tcolorbox}
\end{tcolorbox}
\newpage
\begin{tcolorbox}[enhanced,width=16cm,height=14cm,colback=white,colframe=csta,boxrule=2pt,arc=0mm,
left=6mm,right=6mm,top=0mm,bottom=4mm,sharp corners,
overlay={\fill[csta] (frame.north west) rectangle ([yshift=-1.25cm]frame.north east);
\node[anchor=west,text=white,font=\sffamily\bfseries\LARGE] at ([xshift=5mm,yshift=-0.62cm]frame.north west) {Question 68 B};
\node[anchor=east,text=white,font=\sffamily\itshape\large] at ([xshift=-5mm,yshift=-0.62cm]frame.north east) {Statistiques et probabilités};}]
\vspace*{1.75cm}
\begin{minipage}[t][6.9cm][c]{\linewidth}\LARGE\centering
\begin{minipage}{0.95\linewidth}\raggedright\linespread{1.25}\selectfont Quelle est la moyenne de la série $6\ ;\ 9\ ;\ 13\ ;\ 4$~?\end{minipage}
\end{minipage}

\vfill
\begin{tcolorbox}[enhanced,nobeforeafter,width=7.1cm,height=2.1cm,valign=center,colback=csta!6,colframe=csta,arc=2.5mm,boxrule=0.9pt,left=3mm]{\Large\bfseries\color{csta}a)}\quad\begin{minipage}[c]{5.3cm}\raggedright\Large $8$\end{minipage}\end{tcolorbox}\hfill\begin{tcolorbox}[enhanced,nobeforeafter,width=7.1cm,height=2.1cm,valign=center,colback=csta!6,colframe=csta,arc=2.5mm,boxrule=0.9pt,left=3mm]{\Large\bfseries\color{csta}b)}\quad\begin{minipage}[c]{5.3cm}\raggedright\Large $7{,}5$\end{minipage}\end{tcolorbox}\par\vspace{2.5mm}
\begin{tcolorbox}[enhanced,nobeforeafter,width=7.1cm,height=2.1cm,valign=center,colback=csta!6,colframe=csta,arc=2.5mm,boxrule=0.9pt,left=3mm]{\Large\bfseries\color{csta}c)}\quad\begin{minipage}[c]{5.3cm}\raggedright\Large $9$\end{minipage}\end{tcolorbox}\hfill\begin{tcolorbox}[enhanced,nobeforeafter,width=7.1cm,height=2.1cm,valign=center,colback=csta!6,colframe=csta,arc=2.5mm,boxrule=0.9pt,left=3mm]{\Large\bfseries\color{csta}d)}\quad\begin{minipage}[c]{5.3cm}\raggedright\Large $32$\end{minipage}\end{tcolorbox}
\end{tcolorbox}
\newpage
\begin{tcolorbox}[enhanced,width=16cm,height=14cm,colback=white,colframe=csta,boxrule=2pt,arc=0mm,
left=6mm,right=6mm,top=0mm,bottom=4mm,sharp corners,
overlay={\fill[csta] (frame.north west) rectangle ([yshift=-1.25cm]frame.north east);
\node[anchor=west,text=white,font=\sffamily\bfseries\LARGE] at ([xshift=5mm,yshift=-0.62cm]frame.north west) {Question 69 B};
\node[anchor=east,text=white,font=\sffamily\itshape\large] at ([xshift=-5mm,yshift=-0.62cm]frame.north east) {Statistiques et probabilités};}]
\vspace*{1.75cm}
\begin{minipage}[t][6.9cm][c]{\linewidth}\LARGE\centering
\begin{minipage}{0.95\linewidth}\raggedright\linespread{1.25}\selectfont Quelle est l'étendue de la série $20\ ;\ 17\ ;\ 11\ ;\ 26$~?\end{minipage}
\end{minipage}

\vfill
\begin{tcolorbox}[enhanced,nobeforeafter,width=7.1cm,height=2.1cm,valign=center,colback=csta!6,colframe=csta,arc=2.5mm,boxrule=0.9pt,left=3mm]{\Large\bfseries\color{csta}a)}\quad\begin{minipage}[c]{5.3cm}\raggedright\Large $16$\end{minipage}\end{tcolorbox}\hfill\begin{tcolorbox}[enhanced,nobeforeafter,width=7.1cm,height=2.1cm,valign=center,colback=csta!6,colframe=csta,arc=2.5mm,boxrule=0.9pt,left=3mm]{\Large\bfseries\color{csta}b)}\quad\begin{minipage}[c]{5.3cm}\raggedright\Large $26$\end{minipage}\end{tcolorbox}\par\vspace{2.5mm}
\begin{tcolorbox}[enhanced,nobeforeafter,width=7.1cm,height=2.1cm,valign=center,colback=csta!6,colframe=csta,arc=2.5mm,boxrule=0.9pt,left=3mm]{\Large\bfseries\color{csta}c)}\quad\begin{minipage}[c]{5.3cm}\raggedright\Large $37$\end{minipage}\end{tcolorbox}\hfill\begin{tcolorbox}[enhanced,nobeforeafter,width=7.1cm,height=2.1cm,valign=center,colback=csta!6,colframe=csta,arc=2.5mm,boxrule=0.9pt,left=3mm]{\Large\bfseries\color{csta}d)}\quad\begin{minipage}[c]{5.3cm}\raggedright\Large $15$\end{minipage}\end{tcolorbox}
\end{tcolorbox}
\newpage
\begin{tcolorbox}[enhanced,width=16cm,height=14cm,colback=white,colframe=csta,boxrule=2pt,arc=0mm,
left=6mm,right=6mm,top=0mm,bottom=4mm,sharp corners,
overlay={\fill[csta] (frame.north west) rectangle ([yshift=-1.25cm]frame.north east);
\node[anchor=west,text=white,font=\sffamily\bfseries\LARGE] at ([xshift=5mm,yshift=-0.62cm]frame.north west) {Question 70 B};
\node[anchor=east,text=white,font=\sffamily\itshape\large] at ([xshift=-5mm,yshift=-0.62cm]frame.north east) {Statistiques et probabilités};}]
\vspace*{1.75cm}
\begin{minipage}[t][6.9cm][c]{\linewidth}\LARGE\centering
\begin{minipage}{0.95\linewidth}\raggedright\linespread{1.25}\selectfont Quelle est l'étendue de la série $6\ ;\ 9\ ;\ 13\ ;\ 4$~?\end{minipage}
\end{minipage}

\vfill
\begin{tcolorbox}[enhanced,nobeforeafter,width=7.1cm,height=2.1cm,valign=center,colback=csta!6,colframe=csta,arc=2.5mm,boxrule=0.9pt,left=3mm]{\Large\bfseries\color{csta}a)}\quad\begin{minipage}[c]{5.3cm}\raggedright\Large $9$\end{minipage}\end{tcolorbox}\hfill\begin{tcolorbox}[enhanced,nobeforeafter,width=7.1cm,height=2.1cm,valign=center,colback=csta!6,colframe=csta,arc=2.5mm,boxrule=0.9pt,left=3mm]{\Large\bfseries\color{csta}b)}\quad\begin{minipage}[c]{5.3cm}\raggedright\Large $10$\end{minipage}\end{tcolorbox}\par\vspace{2.5mm}
\begin{tcolorbox}[enhanced,nobeforeafter,width=7.1cm,height=2.1cm,valign=center,colback=csta!6,colframe=csta,arc=2.5mm,boxrule=0.9pt,left=3mm]{\Large\bfseries\color{csta}c)}\quad\begin{minipage}[c]{5.3cm}\raggedright\Large $17$\end{minipage}\end{tcolorbox}\hfill\begin{tcolorbox}[enhanced,nobeforeafter,width=7.1cm,height=2.1cm,valign=center,colback=csta!6,colframe=csta,arc=2.5mm,boxrule=0.9pt,left=3mm]{\Large\bfseries\color{csta}d)}\quad\begin{minipage}[c]{5.3cm}\raggedright\Large $13$\end{minipage}\end{tcolorbox}
\end{tcolorbox}
\newpage
\begin{tcolorbox}[enhanced,width=16cm,height=14cm,colback=white,colframe=csta,boxrule=2pt,arc=0mm,
left=6mm,right=6mm,top=0mm,bottom=4mm,sharp corners,
overlay={\fill[csta] (frame.north west) rectangle ([yshift=-1.25cm]frame.north east);
\node[anchor=west,text=white,font=\sffamily\bfseries\LARGE] at ([xshift=5mm,yshift=-0.62cm]frame.north west) {Question 71 B};
\node[anchor=east,text=white,font=\sffamily\itshape\large] at ([xshift=-5mm,yshift=-0.62cm]frame.north east) {Statistiques et probabilités};}]
\vspace*{1.75cm}
\begin{minipage}[t][6.9cm][c]{\linewidth}\Large\centering
\begin{minipage}{0.95\linewidth}\raggedright\linespread{1.25}\selectfont Une urne contient $6$ boules dont $4$ rouges. On tire successivement deux boules, sans remise. Quelle est la probabilité de tirer deux boules rouges~?\end{minipage}
\end{minipage}

\vfill
\begin{tcolorbox}[enhanced,nobeforeafter,width=7.1cm,height=2.1cm,valign=center,colback=csta!6,colframe=csta,arc=2.5mm,boxrule=0.9pt,left=3mm]{\Large\bfseries\color{csta}a)}\quad\begin{minipage}[c]{5.3cm}\raggedright\Large $\dfrac{4}{9}$\end{minipage}\end{tcolorbox}\hfill\begin{tcolorbox}[enhanced,nobeforeafter,width=7.1cm,height=2.1cm,valign=center,colback=csta!6,colframe=csta,arc=2.5mm,boxrule=0.9pt,left=3mm]{\Large\bfseries\color{csta}b)}\quad\begin{minipage}[c]{5.3cm}\raggedright\Large $\dfrac{2}{3}$\end{minipage}\end{tcolorbox}\par\vspace{2.5mm}
\begin{tcolorbox}[enhanced,nobeforeafter,width=7.1cm,height=2.1cm,valign=center,colback=csta!6,colframe=csta,arc=2.5mm,boxrule=0.9pt,left=3mm]{\Large\bfseries\color{csta}c)}\quad\begin{minipage}[c]{5.3cm}\raggedright\Large $\dfrac{1}{3}$\end{minipage}\end{tcolorbox}\hfill\begin{tcolorbox}[enhanced,nobeforeafter,width=7.1cm,height=2.1cm,valign=center,colback=csta!6,colframe=csta,arc=2.5mm,boxrule=0.9pt,left=3mm]{\Large\bfseries\color{csta}d)}\quad\begin{minipage}[c]{5.3cm}\raggedright\Large $\dfrac{2}{5}$\end{minipage}\end{tcolorbox}
\end{tcolorbox}
\newpage
\begin{tcolorbox}[enhanced,width=16cm,height=14cm,colback=white,colframe=csta,boxrule=2pt,arc=0mm,
left=6mm,right=6mm,top=0mm,bottom=4mm,sharp corners,
overlay={\fill[csta] (frame.north west) rectangle ([yshift=-1.25cm]frame.north east);
\node[anchor=west,text=white,font=\sffamily\bfseries\LARGE] at ([xshift=5mm,yshift=-0.62cm]frame.north west) {Question 72 B};
\node[anchor=east,text=white,font=\sffamily\itshape\large] at ([xshift=-5mm,yshift=-0.62cm]frame.north east) {Statistiques et probabilités};}]
\vspace*{1.75cm}
\begin{minipage}[t][6.9cm][c]{\linewidth}\Large\centering
\begin{minipage}{0.95\linewidth}\raggedright\linespread{1.25}\selectfont Une urne contient $5$ boules dont $2$ rouges. On tire successivement deux boules, sans remise. Quelle est la probabilité de tirer deux boules rouges~?\end{minipage}
\end{minipage}

\vfill
\begin{tcolorbox}[enhanced,nobeforeafter,width=7.1cm,height=2.1cm,valign=center,colback=csta!6,colframe=csta,arc=2.5mm,boxrule=0.9pt,left=3mm]{\Large\bfseries\color{csta}a)}\quad\begin{minipage}[c]{5.3cm}\raggedright\Large $\dfrac{4}{25}$\end{minipage}\end{tcolorbox}\hfill\begin{tcolorbox}[enhanced,nobeforeafter,width=7.1cm,height=2.1cm,valign=center,colback=csta!6,colframe=csta,arc=2.5mm,boxrule=0.9pt,left=3mm]{\Large\bfseries\color{csta}b)}\quad\begin{minipage}[c]{5.3cm}\raggedright\Large $\dfrac{4}{5}$\end{minipage}\end{tcolorbox}\par\vspace{2.5mm}
\begin{tcolorbox}[enhanced,nobeforeafter,width=7.1cm,height=2.1cm,valign=center,colback=csta!6,colframe=csta,arc=2.5mm,boxrule=0.9pt,left=3mm]{\Large\bfseries\color{csta}c)}\quad\begin{minipage}[c]{5.3cm}\raggedright\Large $\dfrac{2}{25}$\end{minipage}\end{tcolorbox}\hfill\begin{tcolorbox}[enhanced,nobeforeafter,width=7.1cm,height=2.1cm,valign=center,colback=csta!6,colframe=csta,arc=2.5mm,boxrule=0.9pt,left=3mm]{\Large\bfseries\color{csta}d)}\quad\begin{minipage}[c]{5.3cm}\raggedright\Large $\dfrac{1}{10}$\end{minipage}\end{tcolorbox}
\end{tcolorbox}
\newpage
\begin{tcolorbox}[enhanced,width=16cm,height=14cm,colback=white,colframe=csta,boxrule=2pt,arc=0mm,
left=6mm,right=6mm,top=0mm,bottom=4mm,sharp corners,
overlay={\fill[csta] (frame.north west) rectangle ([yshift=-1.25cm]frame.north east);
\node[anchor=west,text=white,font=\sffamily\bfseries\LARGE] at ([xshift=5mm,yshift=-0.62cm]frame.north west) {Question 73 B};
\node[anchor=east,text=white,font=\sffamily\itshape\large] at ([xshift=-5mm,yshift=-0.62cm]frame.north east) {Statistiques et probabilités};}]
\vspace*{1.75cm}
\begin{minipage}[t][6.9cm][c]{\linewidth}\Large\centering
\begin{minipage}{0.95\linewidth}\raggedright\linespread{1.25}\selectfont On lance un dé équilibré à six faces et une pièce équilibrée (ou un deuxième dé). Quelle est la probabilité d'obtenir un 1 ou un 2 avec le dé et FACE avec la pièce~?\end{minipage}
\end{minipage}

\vfill
\begin{tcolorbox}[enhanced,nobeforeafter,width=7.1cm,height=2.1cm,valign=center,colback=csta!6,colframe=csta,arc=2.5mm,boxrule=0.9pt,left=3mm]{\Large\bfseries\color{csta}a)}\quad\begin{minipage}[c]{5.3cm}\raggedright\Large $\dfrac{5}{6}$\end{minipage}\end{tcolorbox}\hfill\begin{tcolorbox}[enhanced,nobeforeafter,width=7.1cm,height=2.1cm,valign=center,colback=csta!6,colframe=csta,arc=2.5mm,boxrule=0.9pt,left=3mm]{\Large\bfseries\color{csta}b)}\quad\begin{minipage}[c]{5.3cm}\raggedright\Large $\dfrac{1}{6}$\end{minipage}\end{tcolorbox}\par\vspace{2.5mm}
\begin{tcolorbox}[enhanced,nobeforeafter,width=7.1cm,height=2.1cm,valign=center,colback=csta!6,colframe=csta,arc=2.5mm,boxrule=0.9pt,left=3mm]{\Large\bfseries\color{csta}c)}\quad\begin{minipage}[c]{5.3cm}\raggedright\Large $\dfrac{1}{3}$\end{minipage}\end{tcolorbox}\hfill\begin{tcolorbox}[enhanced,nobeforeafter,width=7.1cm,height=2.1cm,valign=center,colback=csta!6,colframe=csta,arc=2.5mm,boxrule=0.9pt,left=3mm]{\Large\bfseries\color{csta}d)}\quad\begin{minipage}[c]{5.3cm}\raggedright\Large $\dfrac{1}{2}$\end{minipage}\end{tcolorbox}
\end{tcolorbox}
\newpage
\begin{tcolorbox}[enhanced,width=16cm,height=14cm,colback=white,colframe=csta,boxrule=2pt,arc=0mm,
left=6mm,right=6mm,top=0mm,bottom=4mm,sharp corners,
overlay={\fill[csta] (frame.north west) rectangle ([yshift=-1.25cm]frame.north east);
\node[anchor=west,text=white,font=\sffamily\bfseries\LARGE] at ([xshift=5mm,yshift=-0.62cm]frame.north west) {Question 74 B};
\node[anchor=east,text=white,font=\sffamily\itshape\large] at ([xshift=-5mm,yshift=-0.62cm]frame.north east) {Statistiques et probabilités};}]
\vspace*{1.75cm}
\begin{minipage}[t][6.9cm][c]{\linewidth}\Large\centering
\begin{minipage}{0.95\linewidth}\raggedright\linespread{1.25}\selectfont On lance un dé équilibré à six faces et une pièce équilibrée (ou un deuxième dé). Quelle est la probabilité d'obtenir un 5 avec le dé et un 5 avec un deuxième dé~?\end{minipage}
\end{minipage}

\vfill
\begin{tcolorbox}[enhanced,nobeforeafter,width=7.1cm,height=2.1cm,valign=center,colback=csta!6,colframe=csta,arc=2.5mm,boxrule=0.9pt,left=3mm]{\Large\bfseries\color{csta}a)}\quad\begin{minipage}[c]{5.3cm}\raggedright\Large $\dfrac{1}{36}$\end{minipage}\end{tcolorbox}\hfill\begin{tcolorbox}[enhanced,nobeforeafter,width=7.1cm,height=2.1cm,valign=center,colback=csta!6,colframe=csta,arc=2.5mm,boxrule=0.9pt,left=3mm]{\Large\bfseries\color{csta}b)}\quad\begin{minipage}[c]{5.3cm}\raggedright\Large $\dfrac{1}{18}$\end{minipage}\end{tcolorbox}\par\vspace{2.5mm}
\begin{tcolorbox}[enhanced,nobeforeafter,width=7.1cm,height=2.1cm,valign=center,colback=csta!6,colframe=csta,arc=2.5mm,boxrule=0.9pt,left=3mm]{\Large\bfseries\color{csta}c)}\quad\begin{minipage}[c]{5.3cm}\raggedright\Large $\dfrac{1}{6}$\end{minipage}\end{tcolorbox}\hfill\begin{tcolorbox}[enhanced,nobeforeafter,width=7.1cm,height=2.1cm,valign=center,colback=csta!6,colframe=csta,arc=2.5mm,boxrule=0.9pt,left=3mm]{\Large\bfseries\color{csta}d)}\quad\begin{minipage}[c]{5.3cm}\raggedright\Large $\dfrac{1}{3}$\end{minipage}\end{tcolorbox}
\end{tcolorbox}
\newpage
\begin{tcolorbox}[enhanced,width=16cm,height=14cm,colback=white,colframe=csta,boxrule=2pt,arc=0mm,
left=6mm,right=6mm,top=0mm,bottom=4mm,sharp corners,
overlay={\fill[csta] (frame.north west) rectangle ([yshift=-1.25cm]frame.north east);
\node[anchor=west,text=white,font=\sffamily\bfseries\LARGE] at ([xshift=5mm,yshift=-0.62cm]frame.north west) {Question 75 B};
\node[anchor=east,text=white,font=\sffamily\itshape\large] at ([xshift=-5mm,yshift=-0.62cm]frame.north east) {Statistiques et probabilités};}]
\vspace*{1.75cm}
\begin{minipage}[t][6.9cm][c]{\linewidth}\LARGE\centering
\begin{minipage}{0.95\linewidth}\raggedright\linespread{1.25}\selectfont Dans une classe de $20$ élèves, $9$ pratiquent un sport en club. Quelle est la fréquence, en pourcentage, de ces élèves~?\end{minipage}
\end{minipage}

\vfill
\begin{tcolorbox}[enhanced,nobeforeafter,width=7.1cm,height=2.1cm,valign=center,colback=csta!6,colframe=csta,arc=2.5mm,boxrule=0.9pt,left=3mm]{\Large\bfseries\color{csta}a)}\quad\begin{minipage}[c]{5.3cm}\raggedright\Large $45\,\%$\end{minipage}\end{tcolorbox}\hfill\begin{tcolorbox}[enhanced,nobeforeafter,width=7.1cm,height=2.1cm,valign=center,colback=csta!6,colframe=csta,arc=2.5mm,boxrule=0.9pt,left=3mm]{\Large\bfseries\color{csta}b)}\quad\begin{minipage}[c]{5.3cm}\raggedright\Large $9\,\%$\end{minipage}\end{tcolorbox}\par\vspace{2.5mm}
\begin{tcolorbox}[enhanced,nobeforeafter,width=7.1cm,height=2.1cm,valign=center,colback=csta!6,colframe=csta,arc=2.5mm,boxrule=0.9pt,left=3mm]{\Large\bfseries\color{csta}c)}\quad\begin{minipage}[c]{5.3cm}\raggedright\Large $55\,\%$\end{minipage}\end{tcolorbox}\hfill\begin{tcolorbox}[enhanced,nobeforeafter,width=7.1cm,height=2.1cm,valign=center,colback=csta!6,colframe=csta,arc=2.5mm,boxrule=0.9pt,left=3mm]{\Large\bfseries\color{csta}d)}\quad\begin{minipage}[c]{5.3cm}\raggedright\Large $4{,}5\,\%$\end{minipage}\end{tcolorbox}
\end{tcolorbox}
\newpage
\begin{tcolorbox}[enhanced,width=16cm,height=14cm,colback=white,colframe=csta,boxrule=2pt,arc=0mm,
left=6mm,right=6mm,top=0mm,bottom=4mm,sharp corners,
overlay={\fill[csta] (frame.north west) rectangle ([yshift=-1.25cm]frame.north east);
\node[anchor=west,text=white,font=\sffamily\bfseries\LARGE] at ([xshift=5mm,yshift=-0.62cm]frame.north west) {Question 76 B};
\node[anchor=east,text=white,font=\sffamily\itshape\large] at ([xshift=-5mm,yshift=-0.62cm]frame.north east) {Statistiques et probabilités};}]
\vspace*{1.75cm}
\begin{minipage}[t][6.9cm][c]{\linewidth}\LARGE\centering
\begin{minipage}{0.95\linewidth}\raggedright\linespread{1.25}\selectfont Sur $120$ personnes interrogées, $30$ ont répondu «~oui~». Quelle est la fréquence, en pourcentage, des «~oui~»~?\end{minipage}
\end{minipage}

\vfill
\begin{tcolorbox}[enhanced,nobeforeafter,width=7.1cm,height=2.1cm,valign=center,colback=csta!6,colframe=csta,arc=2.5mm,boxrule=0.9pt,left=3mm]{\Large\bfseries\color{csta}a)}\quad\begin{minipage}[c]{5.3cm}\raggedright\Large $30\,\%$\end{minipage}\end{tcolorbox}\hfill\begin{tcolorbox}[enhanced,nobeforeafter,width=7.1cm,height=2.1cm,valign=center,colback=csta!6,colframe=csta,arc=2.5mm,boxrule=0.9pt,left=3mm]{\Large\bfseries\color{csta}b)}\quad\begin{minipage}[c]{5.3cm}\raggedright\Large $25\,\%$\end{minipage}\end{tcolorbox}\par\vspace{2.5mm}
\begin{tcolorbox}[enhanced,nobeforeafter,width=7.1cm,height=2.1cm,valign=center,colback=csta!6,colframe=csta,arc=2.5mm,boxrule=0.9pt,left=3mm]{\Large\bfseries\color{csta}c)}\quad\begin{minipage}[c]{5.3cm}\raggedright\Large $2{,}5\,\%$\end{minipage}\end{tcolorbox}\hfill\begin{tcolorbox}[enhanced,nobeforeafter,width=7.1cm,height=2.1cm,valign=center,colback=csta!6,colframe=csta,arc=2.5mm,boxrule=0.9pt,left=3mm]{\Large\bfseries\color{csta}d)}\quad\begin{minipage}[c]{5.3cm}\raggedright\Large $75\,\%$\end{minipage}\end{tcolorbox}
\end{tcolorbox}
\newpage
\begin{tcolorbox}[enhanced,width=16cm,height=14cm,colback=white,colframe=csta,boxrule=2pt,arc=0mm,
left=6mm,right=6mm,top=0mm,bottom=4mm,sharp corners,
overlay={\fill[csta] (frame.north west) rectangle ([yshift=-1.25cm]frame.north east);
\node[anchor=west,text=white,font=\sffamily\bfseries\LARGE] at ([xshift=5mm,yshift=-0.62cm]frame.north west) {Question 77 B};
\node[anchor=east,text=white,font=\sffamily\itshape\large] at ([xshift=-5mm,yshift=-0.62cm]frame.north east) {Statistiques et probabilités};}]
\vspace*{1.75cm}
\begin{minipage}[t][6.9cm][c]{\linewidth}\Large\centering
\begin{minipage}{0.95\linewidth}\raggedright\linespread{1.25}\selectfont Un sac contient $2$ boules rouges, $8$ boules bleues et $6$ boules vertes. On tire une boule au hasard. Quelle est la probabilité de ne pas tirer une boule rouge~?\end{minipage}
\end{minipage}

\vfill
\begin{tcolorbox}[enhanced,nobeforeafter,width=7.1cm,height=2.1cm,valign=center,colback=csta!6,colframe=csta,arc=2.5mm,boxrule=0.9pt,left=3mm]{\Large\bfseries\color{csta}a)}\quad\begin{minipage}[c]{5.3cm}\raggedright\Large $\dfrac{1}{8}$\end{minipage}\end{tcolorbox}\hfill\begin{tcolorbox}[enhanced,nobeforeafter,width=7.1cm,height=2.1cm,valign=center,colback=csta!6,colframe=csta,arc=2.5mm,boxrule=0.9pt,left=3mm]{\Large\bfseries\color{csta}b)}\quad\begin{minipage}[c]{5.3cm}\raggedright\Large $\dfrac{1}{2}$\end{minipage}\end{tcolorbox}\par\vspace{2.5mm}
\begin{tcolorbox}[enhanced,nobeforeafter,width=7.1cm,height=2.1cm,valign=center,colback=csta!6,colframe=csta,arc=2.5mm,boxrule=0.9pt,left=3mm]{\Large\bfseries\color{csta}c)}\quad\begin{minipage}[c]{5.3cm}\raggedright\Large $\dfrac{7}{8}$\end{minipage}\end{tcolorbox}\hfill\begin{tcolorbox}[enhanced,nobeforeafter,width=7.1cm,height=2.1cm,valign=center,colback=csta!6,colframe=csta,arc=2.5mm,boxrule=0.9pt,left=3mm]{\Large\bfseries\color{csta}d)}\quad\begin{minipage}[c]{5.3cm}\raggedright\Large $7$\end{minipage}\end{tcolorbox}
\end{tcolorbox}
\newpage
\begin{tcolorbox}[enhanced,width=16cm,height=14cm,colback=white,colframe=csta,boxrule=2pt,arc=0mm,
left=6mm,right=6mm,top=0mm,bottom=4mm,sharp corners,
overlay={\fill[csta] (frame.north west) rectangle ([yshift=-1.25cm]frame.north east);
\node[anchor=west,text=white,font=\sffamily\bfseries\LARGE] at ([xshift=5mm,yshift=-0.62cm]frame.north west) {Question 78 B};
\node[anchor=east,text=white,font=\sffamily\itshape\large] at ([xshift=-5mm,yshift=-0.62cm]frame.north east) {Statistiques et probabilités};}]
\vspace*{1.75cm}
\begin{minipage}[t][6.9cm][c]{\linewidth}\Large\centering
\begin{minipage}{0.95\linewidth}\raggedright\linespread{1.25}\selectfont Un sac contient $7$ boules rouges, $5$ boules bleues et $3$ boules vertes. On tire une boule au hasard. Quelle est la probabilité de ne pas tirer une boule rouge~?\end{minipage}
\end{minipage}

\vfill
\begin{tcolorbox}[enhanced,nobeforeafter,width=7.1cm,height=2.1cm,valign=center,colback=csta!6,colframe=csta,arc=2.5mm,boxrule=0.9pt,left=3mm]{\Large\bfseries\color{csta}a)}\quad\begin{minipage}[c]{5.3cm}\raggedright\Large $\dfrac{1}{3}$\end{minipage}\end{tcolorbox}\hfill\begin{tcolorbox}[enhanced,nobeforeafter,width=7.1cm,height=2.1cm,valign=center,colback=csta!6,colframe=csta,arc=2.5mm,boxrule=0.9pt,left=3mm]{\Large\bfseries\color{csta}b)}\quad\begin{minipage}[c]{5.3cm}\raggedright\Large $\dfrac{8}{15}$\end{minipage}\end{tcolorbox}\par\vspace{2.5mm}
\begin{tcolorbox}[enhanced,nobeforeafter,width=7.1cm,height=2.1cm,valign=center,colback=csta!6,colframe=csta,arc=2.5mm,boxrule=0.9pt,left=3mm]{\Large\bfseries\color{csta}c)}\quad\begin{minipage}[c]{5.3cm}\raggedright\Large $\dfrac{8}{7}$\end{minipage}\end{tcolorbox}\hfill\begin{tcolorbox}[enhanced,nobeforeafter,width=7.1cm,height=2.1cm,valign=center,colback=csta!6,colframe=csta,arc=2.5mm,boxrule=0.9pt,left=3mm]{\Large\bfseries\color{csta}d)}\quad\begin{minipage}[c]{5.3cm}\raggedright\Large $\dfrac{7}{15}$\end{minipage}\end{tcolorbox}
\end{tcolorbox}
\newpage
\begin{tcolorbox}[enhanced,width=16cm,height=14cm,colback=white,colframe=csta,boxrule=2pt,arc=0mm,
left=6mm,right=6mm,top=0mm,bottom=4mm,sharp corners,
overlay={\fill[csta] (frame.north west) rectangle ([yshift=-1.25cm]frame.north east);
\node[anchor=west,text=white,font=\sffamily\bfseries\LARGE] at ([xshift=5mm,yshift=-0.62cm]frame.north west) {Question 79 B};
\node[anchor=east,text=white,font=\sffamily\itshape\large] at ([xshift=-5mm,yshift=-0.62cm]frame.north east) {Statistiques et probabilités};}]
\vspace*{1.75cm}
\begin{minipage}[t][6.9cm][c]{\linewidth}\LARGE\centering
\begin{minipage}{0.95\linewidth}\raggedright\linespread{1.25}\selectfont La probabilité d'un événement $A$ est $\dfrac{5}{9}$. Quelle est la probabilité de l'événement contraire de $A$~?\end{minipage}
\end{minipage}

\vfill
\begin{tcolorbox}[enhanced,nobeforeafter,width=7.1cm,height=2.1cm,valign=center,colback=csta!6,colframe=csta,arc=2.5mm,boxrule=0.9pt,left=3mm]{\Large\bfseries\color{csta}a)}\quad\begin{minipage}[c]{5.3cm}\raggedright\Large $\dfrac{4}{9}$\end{minipage}\end{tcolorbox}\hfill\begin{tcolorbox}[enhanced,nobeforeafter,width=7.1cm,height=2.1cm,valign=center,colback=csta!6,colframe=csta,arc=2.5mm,boxrule=0.9pt,left=3mm]{\Large\bfseries\color{csta}b)}\quad\begin{minipage}[c]{5.3cm}\raggedright\Large $\dfrac{2}{9}$\end{minipage}\end{tcolorbox}\par\vspace{2.5mm}
\begin{tcolorbox}[enhanced,nobeforeafter,width=7.1cm,height=2.1cm,valign=center,colback=csta!6,colframe=csta,arc=2.5mm,boxrule=0.9pt,left=3mm]{\Large\bfseries\color{csta}c)}\quad\begin{minipage}[c]{5.3cm}\raggedright\Large $\dfrac{1}{18}$\end{minipage}\end{tcolorbox}\hfill\begin{tcolorbox}[enhanced,nobeforeafter,width=7.1cm,height=2.1cm,valign=center,colback=csta!6,colframe=csta,arc=2.5mm,boxrule=0.9pt,left=3mm]{\Large\bfseries\color{csta}d)}\quad\begin{minipage}[c]{5.3cm}\raggedright\Large $\dfrac{5}{9}$\end{minipage}\end{tcolorbox}
\end{tcolorbox}
\newpage
\begin{tcolorbox}[enhanced,width=16cm,height=14cm,colback=white,colframe=csta,boxrule=2pt,arc=0mm,
left=6mm,right=6mm,top=0mm,bottom=4mm,sharp corners,
overlay={\fill[csta] (frame.north west) rectangle ([yshift=-1.25cm]frame.north east);
\node[anchor=west,text=white,font=\sffamily\bfseries\LARGE] at ([xshift=5mm,yshift=-0.62cm]frame.north west) {Question 80 B};
\node[anchor=east,text=white,font=\sffamily\itshape\large] at ([xshift=-5mm,yshift=-0.62cm]frame.north east) {Statistiques et probabilités};}]
\vspace*{1.75cm}
\begin{minipage}[t][6.9cm][c]{\linewidth}\LARGE\centering
\begin{minipage}{0.95\linewidth}\raggedright\linespread{1.25}\selectfont La probabilité d'un événement $A$ est $0{,}72$. Quelle est la probabilité de l'événement contraire de $A$~?\end{minipage}
\end{minipage}

\vfill
\begin{tcolorbox}[enhanced,nobeforeafter,width=7.1cm,height=2.1cm,valign=center,colback=csta!6,colframe=csta,arc=2.5mm,boxrule=0.9pt,left=3mm]{\Large\bfseries\color{csta}a)}\quad\begin{minipage}[c]{5.3cm}\raggedright\Large $0{,}22$\end{minipage}\end{tcolorbox}\hfill\begin{tcolorbox}[enhanced,nobeforeafter,width=7.1cm,height=2.1cm,valign=center,colback=csta!6,colframe=csta,arc=2.5mm,boxrule=0.9pt,left=3mm]{\Large\bfseries\color{csta}b)}\quad\begin{minipage}[c]{5.3cm}\raggedright\Large $0{,}14$\end{minipage}\end{tcolorbox}\par\vspace{2.5mm}
\begin{tcolorbox}[enhanced,nobeforeafter,width=7.1cm,height=2.1cm,valign=center,colback=csta!6,colframe=csta,arc=2.5mm,boxrule=0.9pt,left=3mm]{\Large\bfseries\color{csta}c)}\quad\begin{minipage}[c]{5.3cm}\raggedright\Large $0{,}28$\end{minipage}\end{tcolorbox}\hfill\begin{tcolorbox}[enhanced,nobeforeafter,width=7.1cm,height=2.1cm,valign=center,colback=csta!6,colframe=csta,arc=2.5mm,boxrule=0.9pt,left=3mm]{\Large\bfseries\color{csta}d)}\quad\begin{minipage}[c]{5.3cm}\raggedright\Large $0{,}72$\end{minipage}\end{tcolorbox}
\end{tcolorbox}
\newpage
\begin{tcolorbox}[enhanced,width=16cm,height=14cm,colback=white,colframe=cgra,boxrule=2pt,arc=0mm,
left=6mm,right=6mm,top=0mm,bottom=4mm,sharp corners,
overlay={\fill[cgra] (frame.north west) rectangle ([yshift=-1.25cm]frame.north east);
\node[anchor=west,text=white,font=\sffamily\bfseries\LARGE] at ([xshift=5mm,yshift=-0.62cm]frame.north west) {Question 81 B};
\node[anchor=east,text=white,font=\sffamily\itshape\large] at ([xshift=-5mm,yshift=-0.62cm]frame.north east) {Grandeurs et proportionnalité};}]
\vspace*{1.75cm}
\begin{minipage}[t][6.9cm][c]{\linewidth}\LARGE\centering
\begin{minipage}{0.95\linewidth}\raggedright\linespread{1.25}\selectfont Une vitesse de $5$ m/s correspond à~:\end{minipage}
\end{minipage}

\vfill
\begin{tcolorbox}[enhanced,nobeforeafter,width=7.1cm,height=2.1cm,valign=center,colback=cgra!6,colframe=cgra,arc=2.5mm,boxrule=0.9pt,left=3mm]{\Large\bfseries\color{cgra}a)}\quad\begin{minipage}[c]{5.3cm}\raggedright\Large $50~\text{km/h}$\end{minipage}\end{tcolorbox}\hfill\begin{tcolorbox}[enhanced,nobeforeafter,width=7.1cm,height=2.1cm,valign=center,colback=cgra!6,colframe=cgra,arc=2.5mm,boxrule=0.9pt,left=3mm]{\Large\bfseries\color{cgra}b)}\quad\begin{minipage}[c]{5.3cm}\raggedright\Large $18~\text{km/h}$\end{minipage}\end{tcolorbox}\par\vspace{2.5mm}
\begin{tcolorbox}[enhanced,nobeforeafter,width=7.1cm,height=2.1cm,valign=center,colback=cgra!6,colframe=cgra,arc=2.5mm,boxrule=0.9pt,left=3mm]{\Large\bfseries\color{cgra}c)}\quad\begin{minipage}[c]{5.3cm}\raggedright\Large $300~\text{km/h}$\end{minipage}\end{tcolorbox}\hfill\begin{tcolorbox}[enhanced,nobeforeafter,width=7.1cm,height=2.1cm,valign=center,colback=cgra!6,colframe=cgra,arc=2.5mm,boxrule=0.9pt,left=3mm]{\Large\bfseries\color{cgra}d)}\quad\begin{minipage}[c]{5.3cm}\raggedright\Large $5\,000~\text{km/h}$\end{minipage}\end{tcolorbox}
\end{tcolorbox}
\newpage
\begin{tcolorbox}[enhanced,width=16cm,height=14cm,colback=white,colframe=cgra,boxrule=2pt,arc=0mm,
left=6mm,right=6mm,top=0mm,bottom=4mm,sharp corners,
overlay={\fill[cgra] (frame.north west) rectangle ([yshift=-1.25cm]frame.north east);
\node[anchor=west,text=white,font=\sffamily\bfseries\LARGE] at ([xshift=5mm,yshift=-0.62cm]frame.north west) {Question 82 B};
\node[anchor=east,text=white,font=\sffamily\itshape\large] at ([xshift=-5mm,yshift=-0.62cm]frame.north east) {Grandeurs et proportionnalité};}]
\vspace*{1.75cm}
\begin{minipage}[t][6.9cm][c]{\linewidth}\LARGE\centering
\begin{minipage}{0.95\linewidth}\raggedright\linespread{1.25}\selectfont Une vitesse de $15$ m/s correspond à~:\end{minipage}
\end{minipage}

\vfill
\begin{tcolorbox}[enhanced,nobeforeafter,width=7.1cm,height=2.1cm,valign=center,colback=cgra!6,colframe=cgra,arc=2.5mm,boxrule=0.9pt,left=3mm]{\Large\bfseries\color{cgra}a)}\quad\begin{minipage}[c]{5.3cm}\raggedright\Large $900~\text{km/h}$\end{minipage}\end{tcolorbox}\hfill\begin{tcolorbox}[enhanced,nobeforeafter,width=7.1cm,height=2.1cm,valign=center,colback=cgra!6,colframe=cgra,arc=2.5mm,boxrule=0.9pt,left=3mm]{\Large\bfseries\color{cgra}b)}\quad\begin{minipage}[c]{5.3cm}\raggedright\Large $54~\text{km/h}$\end{minipage}\end{tcolorbox}\par\vspace{2.5mm}
\begin{tcolorbox}[enhanced,nobeforeafter,width=7.1cm,height=2.1cm,valign=center,colback=cgra!6,colframe=cgra,arc=2.5mm,boxrule=0.9pt,left=3mm]{\Large\bfseries\color{cgra}c)}\quad\begin{minipage}[c]{5.3cm}\raggedright\Large $15\,000~\text{km/h}$\end{minipage}\end{tcolorbox}\hfill\begin{tcolorbox}[enhanced,nobeforeafter,width=7.1cm,height=2.1cm,valign=center,colback=cgra!6,colframe=cgra,arc=2.5mm,boxrule=0.9pt,left=3mm]{\Large\bfseries\color{cgra}d)}\quad\begin{minipage}[c]{5.3cm}\raggedright\Large $150~\text{km/h}$\end{minipage}\end{tcolorbox}
\end{tcolorbox}
\newpage
\begin{tcolorbox}[enhanced,width=16cm,height=14cm,colback=white,colframe=cgra,boxrule=2pt,arc=0mm,
left=6mm,right=6mm,top=0mm,bottom=4mm,sharp corners,
overlay={\fill[cgra] (frame.north west) rectangle ([yshift=-1.25cm]frame.north east);
\node[anchor=west,text=white,font=\sffamily\bfseries\LARGE] at ([xshift=5mm,yshift=-0.62cm]frame.north west) {Question 83 B};
\node[anchor=east,text=white,font=\sffamily\itshape\large] at ([xshift=-5mm,yshift=-0.62cm]frame.north east) {Grandeurs et proportionnalité};}]
\vspace*{1.75cm}
\begin{minipage}[t][6.9cm][c]{\linewidth}\LARGE\centering
\begin{minipage}{0.95\linewidth}\raggedright\linespread{1.25}\selectfont Un véhicule parcourt $60$ km en 45 min. Quelle est sa vitesse moyenne~?\end{minipage}
\end{minipage}

\vfill
\begin{tcolorbox}[enhanced,nobeforeafter,width=7.1cm,height=2.1cm,valign=center,colback=cgra!6,colframe=cgra,arc=2.5mm,boxrule=0.9pt,left=3mm]{\Large\bfseries\color{cgra}a)}\quad\begin{minipage}[c]{5.3cm}\raggedright\Large $80~\text{km/h}$\end{minipage}\end{tcolorbox}\hfill\begin{tcolorbox}[enhanced,nobeforeafter,width=7.1cm,height=2.1cm,valign=center,colback=cgra!6,colframe=cgra,arc=2.5mm,boxrule=0.9pt,left=3mm]{\Large\bfseries\color{cgra}b)}\quad\begin{minipage}[c]{5.3cm}\raggedright\Large $45~\text{km/h}$\end{minipage}\end{tcolorbox}\par\vspace{2.5mm}
\begin{tcolorbox}[enhanced,nobeforeafter,width=7.1cm,height=2.1cm,valign=center,colback=cgra!6,colframe=cgra,arc=2.5mm,boxrule=0.9pt,left=3mm]{\Large\bfseries\color{cgra}c)}\quad\begin{minipage}[c]{5.3cm}\raggedright\Large $40~\text{km/h}$\end{minipage}\end{tcolorbox}\hfill\begin{tcolorbox}[enhanced,nobeforeafter,width=7.1cm,height=2.1cm,valign=center,colback=cgra!6,colframe=cgra,arc=2.5mm,boxrule=0.9pt,left=3mm]{\Large\bfseries\color{cgra}d)}\quad\begin{minipage}[c]{5.3cm}\raggedright\Large $90~\text{km/h}$\end{minipage}\end{tcolorbox}
\end{tcolorbox}
\newpage
\begin{tcolorbox}[enhanced,width=16cm,height=14cm,colback=white,colframe=cgra,boxrule=2pt,arc=0mm,
left=6mm,right=6mm,top=0mm,bottom=4mm,sharp corners,
overlay={\fill[cgra] (frame.north west) rectangle ([yshift=-1.25cm]frame.north east);
\node[anchor=west,text=white,font=\sffamily\bfseries\LARGE] at ([xshift=5mm,yshift=-0.62cm]frame.north west) {Question 84 B};
\node[anchor=east,text=white,font=\sffamily\itshape\large] at ([xshift=-5mm,yshift=-0.62cm]frame.north east) {Grandeurs et proportionnalité};}]
\vspace*{1.75cm}
\begin{minipage}[t][6.9cm][c]{\linewidth}\LARGE\centering
\begin{minipage}{0.95\linewidth}\raggedright\linespread{1.25}\selectfont Un véhicule parcourt $420$ km en 3 h 30 min. Quelle est sa vitesse moyenne~?\end{minipage}
\end{minipage}

\vfill
\begin{tcolorbox}[enhanced,nobeforeafter,width=7.1cm,height=2.1cm,valign=center,colback=cgra!6,colframe=cgra,arc=2.5mm,boxrule=0.9pt,left=3mm]{\Large\bfseries\color{cgra}a)}\quad\begin{minipage}[c]{5.3cm}\raggedright\Large $60~\text{km/h}$\end{minipage}\end{tcolorbox}\hfill\begin{tcolorbox}[enhanced,nobeforeafter,width=7.1cm,height=2.1cm,valign=center,colback=cgra!6,colframe=cgra,arc=2.5mm,boxrule=0.9pt,left=3mm]{\Large\bfseries\color{cgra}b)}\quad\begin{minipage}[c]{5.3cm}\raggedright\Large $130~\text{km/h}$\end{minipage}\end{tcolorbox}\par\vspace{2.5mm}
\begin{tcolorbox}[enhanced,nobeforeafter,width=7.1cm,height=2.1cm,valign=center,colback=cgra!6,colframe=cgra,arc=2.5mm,boxrule=0.9pt,left=3mm]{\Large\bfseries\color{cgra}c)}\quad\begin{minipage}[c]{5.3cm}\raggedright\Large $120~\text{km/h}$\end{minipage}\end{tcolorbox}\hfill\begin{tcolorbox}[enhanced,nobeforeafter,width=7.1cm,height=2.1cm,valign=center,colback=cgra!6,colframe=cgra,arc=2.5mm,boxrule=0.9pt,left=3mm]{\Large\bfseries\color{cgra}d)}\quad\begin{minipage}[c]{5.3cm}\raggedright\Large $1\,470~\text{km/h}$\end{minipage}\end{tcolorbox}
\end{tcolorbox}
\newpage
\begin{tcolorbox}[enhanced,width=16cm,height=14cm,colback=white,colframe=cgra,boxrule=2pt,arc=0mm,
left=6mm,right=6mm,top=0mm,bottom=4mm,sharp corners,
overlay={\fill[cgra] (frame.north west) rectangle ([yshift=-1.25cm]frame.north east);
\node[anchor=west,text=white,font=\sffamily\bfseries\LARGE] at ([xshift=5mm,yshift=-0.62cm]frame.north west) {Question 85 B};
\node[anchor=east,text=white,font=\sffamily\itshape\large] at ([xshift=-5mm,yshift=-0.62cm]frame.north east) {Grandeurs et proportionnalité};}]
\vspace*{1.75cm}
\begin{minipage}[t][6.9cm][c]{\linewidth}\LARGE\centering
\begin{minipage}{0.95\linewidth}\raggedright\linespread{1.25}\selectfont Un article coûte $1\,500$ F. Son prix augmente de $12\,\%$. Quel est son nouveau prix~?\end{minipage}
\end{minipage}

\vfill
\begin{tcolorbox}[enhanced,nobeforeafter,width=7.1cm,height=2.1cm,valign=center,colback=cgra!6,colframe=cgra,arc=2.5mm,boxrule=0.9pt,left=3mm]{\Large\bfseries\color{cgra}a)}\quad\begin{minipage}[c]{5.3cm}\raggedright\Large $180~\text{F}$\end{minipage}\end{tcolorbox}\hfill\begin{tcolorbox}[enhanced,nobeforeafter,width=7.1cm,height=2.1cm,valign=center,colback=cgra!6,colframe=cgra,arc=2.5mm,boxrule=0.9pt,left=3mm]{\Large\bfseries\color{cgra}b)}\quad\begin{minipage}[c]{5.3cm}\raggedright\Large $1\,512~\text{F}$\end{minipage}\end{tcolorbox}\par\vspace{2.5mm}
\begin{tcolorbox}[enhanced,nobeforeafter,width=7.1cm,height=2.1cm,valign=center,colback=cgra!6,colframe=cgra,arc=2.5mm,boxrule=0.9pt,left=3mm]{\Large\bfseries\color{cgra}c)}\quad\begin{minipage}[c]{5.3cm}\raggedright\Large $1\,320~\text{F}$\end{minipage}\end{tcolorbox}\hfill\begin{tcolorbox}[enhanced,nobeforeafter,width=7.1cm,height=2.1cm,valign=center,colback=cgra!6,colframe=cgra,arc=2.5mm,boxrule=0.9pt,left=3mm]{\Large\bfseries\color{cgra}d)}\quad\begin{minipage}[c]{5.3cm}\raggedright\Large $1\,680~\text{F}$\end{minipage}\end{tcolorbox}
\end{tcolorbox}
\newpage
\begin{tcolorbox}[enhanced,width=16cm,height=14cm,colback=white,colframe=cgra,boxrule=2pt,arc=0mm,
left=6mm,right=6mm,top=0mm,bottom=4mm,sharp corners,
overlay={\fill[cgra] (frame.north west) rectangle ([yshift=-1.25cm]frame.north east);
\node[anchor=west,text=white,font=\sffamily\bfseries\LARGE] at ([xshift=5mm,yshift=-0.62cm]frame.north west) {Question 86 B};
\node[anchor=east,text=white,font=\sffamily\itshape\large] at ([xshift=-5mm,yshift=-0.62cm]frame.north east) {Grandeurs et proportionnalité};}]
\vspace*{1.75cm}
\begin{minipage}[t][6.9cm][c]{\linewidth}\LARGE\centering
\begin{minipage}{0.95\linewidth}\raggedright\linespread{1.25}\selectfont Un article coûte $85$ F. Son prix baisse de $20\,\%$. Quel est son nouveau prix~?\end{minipage}
\end{minipage}

\vfill
\begin{tcolorbox}[enhanced,nobeforeafter,width=7.1cm,height=2.1cm,valign=center,colback=cgra!6,colframe=cgra,arc=2.5mm,boxrule=0.9pt,left=3mm]{\Large\bfseries\color{cgra}a)}\quad\begin{minipage}[c]{5.3cm}\raggedright\Large $65~\text{F}$\end{minipage}\end{tcolorbox}\hfill\begin{tcolorbox}[enhanced,nobeforeafter,width=7.1cm,height=2.1cm,valign=center,colback=cgra!6,colframe=cgra,arc=2.5mm,boxrule=0.9pt,left=3mm]{\Large\bfseries\color{cgra}b)}\quad\begin{minipage}[c]{5.3cm}\raggedright\Large $102~\text{F}$\end{minipage}\end{tcolorbox}\par\vspace{2.5mm}
\begin{tcolorbox}[enhanced,nobeforeafter,width=7.1cm,height=2.1cm,valign=center,colback=cgra!6,colframe=cgra,arc=2.5mm,boxrule=0.9pt,left=3mm]{\Large\bfseries\color{cgra}c)}\quad\begin{minipage}[c]{5.3cm}\raggedright\Large $68~\text{F}$\end{minipage}\end{tcolorbox}\hfill\begin{tcolorbox}[enhanced,nobeforeafter,width=7.1cm,height=2.1cm,valign=center,colback=cgra!6,colframe=cgra,arc=2.5mm,boxrule=0.9pt,left=3mm]{\Large\bfseries\color{cgra}d)}\quad\begin{minipage}[c]{5.3cm}\raggedright\Large $17~\text{F}$\end{minipage}\end{tcolorbox}
\end{tcolorbox}
\newpage
\begin{tcolorbox}[enhanced,width=16cm,height=14cm,colback=white,colframe=cgra,boxrule=2pt,arc=0mm,
left=6mm,right=6mm,top=0mm,bottom=4mm,sharp corners,
overlay={\fill[cgra] (frame.north west) rectangle ([yshift=-1.25cm]frame.north east);
\node[anchor=west,text=white,font=\sffamily\bfseries\LARGE] at ([xshift=5mm,yshift=-0.62cm]frame.north west) {Question 87 B};
\node[anchor=east,text=white,font=\sffamily\itshape\large] at ([xshift=-5mm,yshift=-0.62cm]frame.north east) {Grandeurs et proportionnalité};}]
\vspace*{1.75cm}
\begin{minipage}[t][6.9cm][c]{\linewidth}\LARGE\centering
\begin{minipage}{0.95\linewidth}\raggedright\linespread{1.25}\selectfont Un prix passe de $40$ F à $46$ F. Quel est le pourcentage d'évolution~?\end{minipage}
\end{minipage}

\vfill
\begin{tcolorbox}[enhanced,nobeforeafter,width=7.1cm,height=2.1cm,valign=center,colback=cgra!6,colframe=cgra,arc=2.5mm,boxrule=0.9pt,left=3mm]{\Large\bfseries\color{cgra}a)}\quad\begin{minipage}[c]{5.3cm}\raggedright\Large $-15\,\%$\end{minipage}\end{tcolorbox}\hfill\begin{tcolorbox}[enhanced,nobeforeafter,width=7.1cm,height=2.1cm,valign=center,colback=cgra!6,colframe=cgra,arc=2.5mm,boxrule=0.9pt,left=3mm]{\Large\bfseries\color{cgra}b)}\quad\begin{minipage}[c]{5.3cm}\raggedright\Large $+6\,\%$\end{minipage}\end{tcolorbox}\par\vspace{2.5mm}
\begin{tcolorbox}[enhanced,nobeforeafter,width=7.1cm,height=2.1cm,valign=center,colback=cgra!6,colframe=cgra,arc=2.5mm,boxrule=0.9pt,left=3mm]{\Large\bfseries\color{cgra}c)}\quad\begin{minipage}[c]{5.3cm}\raggedright\Large $+15\,\%$\end{minipage}\end{tcolorbox}\hfill\begin{tcolorbox}[enhanced,nobeforeafter,width=7.1cm,height=2.1cm,valign=center,colback=cgra!6,colframe=cgra,arc=2.5mm,boxrule=0.9pt,left=3mm]{\Large\bfseries\color{cgra}d)}\quad\begin{minipage}[c]{5.3cm}\raggedright\Large $+7{,}5\,\%$\end{minipage}\end{tcolorbox}
\end{tcolorbox}
\newpage
\begin{tcolorbox}[enhanced,width=16cm,height=14cm,colback=white,colframe=cgra,boxrule=2pt,arc=0mm,
left=6mm,right=6mm,top=0mm,bottom=4mm,sharp corners,
overlay={\fill[cgra] (frame.north west) rectangle ([yshift=-1.25cm]frame.north east);
\node[anchor=west,text=white,font=\sffamily\bfseries\LARGE] at ([xshift=5mm,yshift=-0.62cm]frame.north west) {Question 88 B};
\node[anchor=east,text=white,font=\sffamily\itshape\large] at ([xshift=-5mm,yshift=-0.62cm]frame.north east) {Grandeurs et proportionnalité};}]
\vspace*{1.75cm}
\begin{minipage}[t][6.9cm][c]{\linewidth}\LARGE\centering
\begin{minipage}{0.95\linewidth}\raggedright\linespread{1.25}\selectfont Un prix passe de $120$ F à $90$ F. Quel est le pourcentage d'évolution~?\end{minipage}
\end{minipage}

\vfill
\begin{tcolorbox}[enhanced,nobeforeafter,width=7.1cm,height=2.1cm,valign=center,colback=cgra!6,colframe=cgra,arc=2.5mm,boxrule=0.9pt,left=3mm]{\Large\bfseries\color{cgra}a)}\quad\begin{minipage}[c]{5.3cm}\raggedright\Large $-30\,\%$\end{minipage}\end{tcolorbox}\hfill\begin{tcolorbox}[enhanced,nobeforeafter,width=7.1cm,height=2.1cm,valign=center,colback=cgra!6,colframe=cgra,arc=2.5mm,boxrule=0.9pt,left=3mm]{\Large\bfseries\color{cgra}b)}\quad\begin{minipage}[c]{5.3cm}\raggedright\Large $+25\,\%$\end{minipage}\end{tcolorbox}\par\vspace{2.5mm}
\begin{tcolorbox}[enhanced,nobeforeafter,width=7.1cm,height=2.1cm,valign=center,colback=cgra!6,colframe=cgra,arc=2.5mm,boxrule=0.9pt,left=3mm]{\Large\bfseries\color{cgra}c)}\quad\begin{minipage}[c]{5.3cm}\raggedright\Large $-25\,\%$\end{minipage}\end{tcolorbox}\hfill\begin{tcolorbox}[enhanced,nobeforeafter,width=7.1cm,height=2.1cm,valign=center,colback=cgra!6,colframe=cgra,arc=2.5mm,boxrule=0.9pt,left=3mm]{\Large\bfseries\color{cgra}d)}\quad\begin{minipage}[c]{5.3cm}\raggedright\Large $-12{,}5\,\%$\end{minipage}\end{tcolorbox}
\end{tcolorbox}
\newpage
\begin{tcolorbox}[enhanced,width=16cm,height=14cm,colback=white,colframe=cgra,boxrule=2pt,arc=0mm,
left=6mm,right=6mm,top=0mm,bottom=4mm,sharp corners,
overlay={\fill[cgra] (frame.north west) rectangle ([yshift=-1.25cm]frame.north east);
\node[anchor=west,text=white,font=\sffamily\bfseries\LARGE] at ([xshift=5mm,yshift=-0.62cm]frame.north west) {Question 89 B};
\node[anchor=east,text=white,font=\sffamily\itshape\large] at ([xshift=-5mm,yshift=-0.62cm]frame.north east) {Grandeurs et proportionnalité};}]
\vspace*{1.75cm}
\begin{minipage}[t][6.9cm][c]{\linewidth}\LARGE\centering
\begin{minipage}{0.95\linewidth}\raggedright\linespread{1.25}\selectfont Un robinet remplit une cuve de $250$ L en $10$ min. Quel est son débit~?\end{minipage}
\end{minipage}

\vfill
\begin{tcolorbox}[enhanced,nobeforeafter,width=7.1cm,height=2.1cm,valign=center,colback=cgra!6,colframe=cgra,arc=2.5mm,boxrule=0.9pt,left=3mm]{\Large\bfseries\color{cgra}a)}\quad\begin{minipage}[c]{5.3cm}\raggedright\Large $25~\text{L/min}$\end{minipage}\end{tcolorbox}\hfill\begin{tcolorbox}[enhanced,nobeforeafter,width=7.1cm,height=2.1cm,valign=center,colback=cgra!6,colframe=cgra,arc=2.5mm,boxrule=0.9pt,left=3mm]{\Large\bfseries\color{cgra}b)}\quad\begin{minipage}[c]{5.3cm}\raggedright\Large $2\,500~\text{L/min}$\end{minipage}\end{tcolorbox}\par\vspace{2.5mm}
\begin{tcolorbox}[enhanced,nobeforeafter,width=7.1cm,height=2.1cm,valign=center,colback=cgra!6,colframe=cgra,arc=2.5mm,boxrule=0.9pt,left=3mm]{\Large\bfseries\color{cgra}c)}\quad\begin{minipage}[c]{5.3cm}\raggedright\Large $27~\text{L/min}$\end{minipage}\end{tcolorbox}\hfill\begin{tcolorbox}[enhanced,nobeforeafter,width=7.1cm,height=2.1cm,valign=center,colback=cgra!6,colframe=cgra,arc=2.5mm,boxrule=0.9pt,left=3mm]{\Large\bfseries\color{cgra}d)}\quad\begin{minipage}[c]{5.3cm}\raggedright\Large $4~\text{L/min}$\end{minipage}\end{tcolorbox}
\end{tcolorbox}
\newpage
\begin{tcolorbox}[enhanced,width=16cm,height=14cm,colback=white,colframe=cgra,boxrule=2pt,arc=0mm,
left=6mm,right=6mm,top=0mm,bottom=4mm,sharp corners,
overlay={\fill[cgra] (frame.north west) rectangle ([yshift=-1.25cm]frame.north east);
\node[anchor=west,text=white,font=\sffamily\bfseries\LARGE] at ([xshift=5mm,yshift=-0.62cm]frame.north west) {Question 90 B};
\node[anchor=east,text=white,font=\sffamily\itshape\large] at ([xshift=-5mm,yshift=-0.62cm]frame.north east) {Grandeurs et proportionnalité};}]
\vspace*{1.75cm}
\begin{minipage}[t][6.9cm][c]{\linewidth}\LARGE\centering
\begin{minipage}{0.95\linewidth}\raggedright\linespread{1.25}\selectfont Un robinet remplit une cuve de $36$ L en $4$ min. Quel est son débit~?\end{minipage}
\end{minipage}

\vfill
\begin{tcolorbox}[enhanced,nobeforeafter,width=7.1cm,height=2.1cm,valign=center,colback=cgra!6,colframe=cgra,arc=2.5mm,boxrule=0.9pt,left=3mm]{\Large\bfseries\color{cgra}a)}\quad\begin{minipage}[c]{5.3cm}\raggedright\large $\dfrac{100}{9}~\text{L/min}$\end{minipage}\end{tcolorbox}\hfill\begin{tcolorbox}[enhanced,nobeforeafter,width=7.1cm,height=2.1cm,valign=center,colback=cgra!6,colframe=cgra,arc=2.5mm,boxrule=0.9pt,left=3mm]{\Large\bfseries\color{cgra}b)}\quad\begin{minipage}[c]{5.3cm}\raggedright\Large $144~\text{L/min}$\end{minipage}\end{tcolorbox}\par\vspace{2.5mm}
\begin{tcolorbox}[enhanced,nobeforeafter,width=7.1cm,height=2.1cm,valign=center,colback=cgra!6,colframe=cgra,arc=2.5mm,boxrule=0.9pt,left=3mm]{\Large\bfseries\color{cgra}c)}\quad\begin{minipage}[c]{5.3cm}\raggedright\Large $11~\text{L/min}$\end{minipage}\end{tcolorbox}\hfill\begin{tcolorbox}[enhanced,nobeforeafter,width=7.1cm,height=2.1cm,valign=center,colback=cgra!6,colframe=cgra,arc=2.5mm,boxrule=0.9pt,left=3mm]{\Large\bfseries\color{cgra}d)}\quad\begin{minipage}[c]{5.3cm}\raggedright\Large $9~\text{L/min}$\end{minipage}\end{tcolorbox}
\end{tcolorbox}
\newpage
\begin{tcolorbox}[enhanced,width=16cm,height=14cm,colback=white,colframe=cgra,boxrule=2pt,arc=0mm,
left=6mm,right=6mm,top=0mm,bottom=4mm,sharp corners,
overlay={\fill[cgra] (frame.north west) rectangle ([yshift=-1.25cm]frame.north east);
\node[anchor=west,text=white,font=\sffamily\bfseries\LARGE] at ([xshift=5mm,yshift=-0.62cm]frame.north west) {Question 91 B};
\node[anchor=east,text=white,font=\sffamily\itshape\large] at ([xshift=-5mm,yshift=-0.62cm]frame.north east) {Grandeurs et proportionnalité};}]
\vspace*{1.75cm}
\begin{minipage}[t][6.9cm][c]{\linewidth}\LARGE\centering
\begin{minipage}{0.95\linewidth}\raggedright\linespread{1.25}\selectfont Sur une carte à l'échelle $\dfrac{1}{25\,000}$, deux points sont distants de $2$ cm. Quelle est la distance réelle~?\end{minipage}
\end{minipage}

\vfill
\begin{tcolorbox}[enhanced,nobeforeafter,width=7.1cm,height=2.1cm,valign=center,colback=cgra!6,colframe=cgra,arc=2.5mm,boxrule=0.9pt,left=3mm]{\Large\bfseries\color{cgra}a)}\quad\begin{minipage}[c]{5.3cm}\raggedright\Large $5\,000~\text{m}$\end{minipage}\end{tcolorbox}\hfill\begin{tcolorbox}[enhanced,nobeforeafter,width=7.1cm,height=2.1cm,valign=center,colback=cgra!6,colframe=cgra,arc=2.5mm,boxrule=0.9pt,left=3mm]{\Large\bfseries\color{cgra}b)}\quad\begin{minipage}[c]{5.3cm}\raggedright\Large $50~\text{m}$\end{minipage}\end{tcolorbox}\par\vspace{2.5mm}
\begin{tcolorbox}[enhanced,nobeforeafter,width=7.1cm,height=2.1cm,valign=center,colback=cgra!6,colframe=cgra,arc=2.5mm,boxrule=0.9pt,left=3mm]{\Large\bfseries\color{cgra}c)}\quad\begin{minipage}[c]{5.3cm}\raggedright\Large $500~\text{m}$\end{minipage}\end{tcolorbox}\hfill\begin{tcolorbox}[enhanced,nobeforeafter,width=7.1cm,height=2.1cm,valign=center,colback=cgra!6,colframe=cgra,arc=2.5mm,boxrule=0.9pt,left=3mm]{\Large\bfseries\color{cgra}d)}\quad\begin{minipage}[c]{5.3cm}\raggedright\Large $50\,000~\text{m}$\end{minipage}\end{tcolorbox}
\end{tcolorbox}
\newpage
\begin{tcolorbox}[enhanced,width=16cm,height=14cm,colback=white,colframe=cgra,boxrule=2pt,arc=0mm,
left=6mm,right=6mm,top=0mm,bottom=4mm,sharp corners,
overlay={\fill[cgra] (frame.north west) rectangle ([yshift=-1.25cm]frame.north east);
\node[anchor=west,text=white,font=\sffamily\bfseries\LARGE] at ([xshift=5mm,yshift=-0.62cm]frame.north west) {Question 92 B};
\node[anchor=east,text=white,font=\sffamily\itshape\large] at ([xshift=-5mm,yshift=-0.62cm]frame.north east) {Grandeurs et proportionnalité};}]
\vspace*{1.75cm}
\begin{minipage}[t][6.9cm][c]{\linewidth}\LARGE\centering
\begin{minipage}{0.95\linewidth}\raggedright\linespread{1.25}\selectfont Sur une carte à l'échelle $\dfrac{1}{200}$, deux points sont distants de $15$ cm. Quelle est la distance réelle~?\end{minipage}
\end{minipage}

\vfill
\begin{tcolorbox}[enhanced,nobeforeafter,width=7.1cm,height=2.1cm,valign=center,colback=cgra!6,colframe=cgra,arc=2.5mm,boxrule=0.9pt,left=3mm]{\Large\bfseries\color{cgra}a)}\quad\begin{minipage}[c]{5.3cm}\raggedright\Large $30~\text{m}$\end{minipage}\end{tcolorbox}\hfill\begin{tcolorbox}[enhanced,nobeforeafter,width=7.1cm,height=2.1cm,valign=center,colback=cgra!6,colframe=cgra,arc=2.5mm,boxrule=0.9pt,left=3mm]{\Large\bfseries\color{cgra}b)}\quad\begin{minipage}[c]{5.3cm}\raggedright\Large $3~\text{m}$\end{minipage}\end{tcolorbox}\par\vspace{2.5mm}
\begin{tcolorbox}[enhanced,nobeforeafter,width=7.1cm,height=2.1cm,valign=center,colback=cgra!6,colframe=cgra,arc=2.5mm,boxrule=0.9pt,left=3mm]{\Large\bfseries\color{cgra}c)}\quad\begin{minipage}[c]{5.3cm}\raggedright\Large $300~\text{m}$\end{minipage}\end{tcolorbox}\hfill\begin{tcolorbox}[enhanced,nobeforeafter,width=7.1cm,height=2.1cm,valign=center,colback=cgra!6,colframe=cgra,arc=2.5mm,boxrule=0.9pt,left=3mm]{\Large\bfseries\color{cgra}d)}\quad\begin{minipage}[c]{5.3cm}\raggedright\Large $3\,000~\text{m}$\end{minipage}\end{tcolorbox}
\end{tcolorbox}
\newpage
\begin{tcolorbox}[enhanced,width=16cm,height=14cm,colback=white,colframe=cgra,boxrule=2pt,arc=0mm,
left=6mm,right=6mm,top=0mm,bottom=4mm,sharp corners,
overlay={\fill[cgra] (frame.north west) rectangle ([yshift=-1.25cm]frame.north east);
\node[anchor=west,text=white,font=\sffamily\bfseries\LARGE] at ([xshift=5mm,yshift=-0.62cm]frame.north west) {Question 93 B};
\node[anchor=east,text=white,font=\sffamily\itshape\large] at ([xshift=-5mm,yshift=-0.62cm]frame.north east) {Grandeurs et proportionnalité};}]
\vspace*{1.75cm}
\begin{minipage}[t][6.9cm][c]{\linewidth}\LARGE\centering
\begin{minipage}{0.95\linewidth}\raggedright\linespread{1.25}\selectfont On partage $45$ bonbons entre Léa et Noa dans le ratio $1:2$. Combien chacun en reçoit-il (Léa, puis Noa)~?\end{minipage}
\end{minipage}

\vfill
\begin{tcolorbox}[enhanced,nobeforeafter,width=7.1cm,height=2.1cm,valign=center,colback=cgra!6,colframe=cgra,arc=2.5mm,boxrule=0.9pt,left=3mm]{\Large\bfseries\color{cgra}a)}\quad\begin{minipage}[c]{5.3cm}\raggedright\Large $14\text{ et }31$\end{minipage}\end{tcolorbox}\hfill\begin{tcolorbox}[enhanced,nobeforeafter,width=7.1cm,height=2.1cm,valign=center,colback=cgra!6,colframe=cgra,arc=2.5mm,boxrule=0.9pt,left=3mm]{\Large\bfseries\color{cgra}b)}\quad\begin{minipage}[c]{5.3cm}\raggedright\Large $22{,}5\text{ et }22{,}5$\end{minipage}\end{tcolorbox}\par\vspace{2.5mm}
\begin{tcolorbox}[enhanced,nobeforeafter,width=7.1cm,height=2.1cm,valign=center,colback=cgra!6,colframe=cgra,arc=2.5mm,boxrule=0.9pt,left=3mm]{\Large\bfseries\color{cgra}c)}\quad\begin{minipage}[c]{5.3cm}\raggedright\Large $15\text{ et }30$\end{minipage}\end{tcolorbox}\hfill\begin{tcolorbox}[enhanced,nobeforeafter,width=7.1cm,height=2.1cm,valign=center,colback=cgra!6,colframe=cgra,arc=2.5mm,boxrule=0.9pt,left=3mm]{\Large\bfseries\color{cgra}d)}\quad\begin{minipage}[c]{5.3cm}\raggedright\Large $30\text{ et }15$\end{minipage}\end{tcolorbox}
\end{tcolorbox}
\newpage
\begin{tcolorbox}[enhanced,width=16cm,height=14cm,colback=white,colframe=cgra,boxrule=2pt,arc=0mm,
left=6mm,right=6mm,top=0mm,bottom=4mm,sharp corners,
overlay={\fill[cgra] (frame.north west) rectangle ([yshift=-1.25cm]frame.north east);
\node[anchor=west,text=white,font=\sffamily\bfseries\LARGE] at ([xshift=5mm,yshift=-0.62cm]frame.north west) {Question 94 B};
\node[anchor=east,text=white,font=\sffamily\itshape\large] at ([xshift=-5mm,yshift=-0.62cm]frame.north east) {Grandeurs et proportionnalité};}]
\vspace*{1.75cm}
\begin{minipage}[t][6.9cm][c]{\linewidth}\LARGE\centering
\begin{minipage}{0.95\linewidth}\raggedright\linespread{1.25}\selectfont On partage $120$ bonbons entre Léa et Noa dans le ratio $3:5$. Combien chacun en reçoit-il (Léa, puis Noa)~?\end{minipage}
\end{minipage}

\vfill
\begin{tcolorbox}[enhanced,nobeforeafter,width=7.1cm,height=2.1cm,valign=center,colback=cgra!6,colframe=cgra,arc=2.5mm,boxrule=0.9pt,left=3mm]{\Large\bfseries\color{cgra}a)}\quad\begin{minipage}[c]{5.3cm}\raggedright\Large $45\text{ et }75$\end{minipage}\end{tcolorbox}\hfill\begin{tcolorbox}[enhanced,nobeforeafter,width=7.1cm,height=2.1cm,valign=center,colback=cgra!6,colframe=cgra,arc=2.5mm,boxrule=0.9pt,left=3mm]{\Large\bfseries\color{cgra}b)}\quad\begin{minipage}[c]{5.3cm}\raggedright\Large $60\text{ et }60$\end{minipage}\end{tcolorbox}\par\vspace{2.5mm}
\begin{tcolorbox}[enhanced,nobeforeafter,width=7.1cm,height=2.1cm,valign=center,colback=cgra!6,colframe=cgra,arc=2.5mm,boxrule=0.9pt,left=3mm]{\Large\bfseries\color{cgra}c)}\quad\begin{minipage}[c]{5.3cm}\raggedright\Large $40\text{ et }80$\end{minipage}\end{tcolorbox}\hfill\begin{tcolorbox}[enhanced,nobeforeafter,width=7.1cm,height=2.1cm,valign=center,colback=cgra!6,colframe=cgra,arc=2.5mm,boxrule=0.9pt,left=3mm]{\Large\bfseries\color{cgra}d)}\quad\begin{minipage}[c]{5.3cm}\raggedright\Large $75\text{ et }45$\end{minipage}\end{tcolorbox}
\end{tcolorbox}
\newpage
\begin{tcolorbox}[enhanced,width=16cm,height=14cm,colback=white,colframe=cgra,boxrule=2pt,arc=0mm,
left=6mm,right=6mm,top=0mm,bottom=4mm,sharp corners,
overlay={\fill[cgra] (frame.north west) rectangle ([yshift=-1.25cm]frame.north east);
\node[anchor=west,text=white,font=\sffamily\bfseries\LARGE] at ([xshift=5mm,yshift=-0.62cm]frame.north west) {Question 95 B};
\node[anchor=east,text=white,font=\sffamily\itshape\large] at ([xshift=-5mm,yshift=-0.62cm]frame.north east) {Grandeurs et proportionnalité};}]
\vspace*{1.75cm}
\begin{minipage}[t][6.9cm][c]{\linewidth}\LARGE\centering
\begin{minipage}{0.95\linewidth}\raggedright\linespread{1.25}\selectfont Compléter~: $0{,}3~\text{m}^3=\;?\;\text{dm}^3$\end{minipage}
\end{minipage}

\vfill
\begin{tcolorbox}[enhanced,nobeforeafter,width=7.1cm,height=2.1cm,valign=center,colback=cgra!6,colframe=cgra,arc=2.5mm,boxrule=0.9pt,left=3mm]{\Large\bfseries\color{cgra}a)}\quad\begin{minipage}[c]{5.3cm}\raggedright\Large $300~\text{dm}^3$\end{minipage}\end{tcolorbox}\hfill\begin{tcolorbox}[enhanced,nobeforeafter,width=7.1cm,height=2.1cm,valign=center,colback=cgra!6,colframe=cgra,arc=2.5mm,boxrule=0.9pt,left=3mm]{\Large\bfseries\color{cgra}b)}\quad\begin{minipage}[c]{5.3cm}\raggedright\Large $30~\text{dm}^3$\end{minipage}\end{tcolorbox}\par\vspace{2.5mm}
\begin{tcolorbox}[enhanced,nobeforeafter,width=7.1cm,height=2.1cm,valign=center,colback=cgra!6,colframe=cgra,arc=2.5mm,boxrule=0.9pt,left=3mm]{\Large\bfseries\color{cgra}c)}\quad\begin{minipage}[c]{5.3cm}\raggedright\Large $3\,000~\text{dm}^3$\end{minipage}\end{tcolorbox}\hfill\begin{tcolorbox}[enhanced,nobeforeafter,width=7.1cm,height=2.1cm,valign=center,colback=cgra!6,colframe=cgra,arc=2.5mm,boxrule=0.9pt,left=3mm]{\Large\bfseries\color{cgra}d)}\quad\begin{minipage}[c]{5.3cm}\raggedright\Large $30\,000~\text{dm}^3$\end{minipage}\end{tcolorbox}
\end{tcolorbox}
\newpage
\begin{tcolorbox}[enhanced,width=16cm,height=14cm,colback=white,colframe=cgra,boxrule=2pt,arc=0mm,
left=6mm,right=6mm,top=0mm,bottom=4mm,sharp corners,
overlay={\fill[cgra] (frame.north west) rectangle ([yshift=-1.25cm]frame.north east);
\node[anchor=west,text=white,font=\sffamily\bfseries\LARGE] at ([xshift=5mm,yshift=-0.62cm]frame.north west) {Question 96 B};
\node[anchor=east,text=white,font=\sffamily\itshape\large] at ([xshift=-5mm,yshift=-0.62cm]frame.north east) {Grandeurs et proportionnalité};}]
\vspace*{1.75cm}
\begin{minipage}[t][6.9cm][c]{\linewidth}\LARGE\centering
\begin{minipage}{0.95\linewidth}\raggedright\linespread{1.25}\selectfont Compléter~: $4\,500~\text{L}=\;?\;\text{m}^3$\end{minipage}
\end{minipage}

\vfill
\begin{tcolorbox}[enhanced,nobeforeafter,width=7.1cm,height=2.1cm,valign=center,colback=cgra!6,colframe=cgra,arc=2.5mm,boxrule=0.9pt,left=3mm]{\Large\bfseries\color{cgra}a)}\quad\begin{minipage}[c]{5.3cm}\raggedright\Large $450~\text{m}^3$\end{minipage}\end{tcolorbox}\hfill\begin{tcolorbox}[enhanced,nobeforeafter,width=7.1cm,height=2.1cm,valign=center,colback=cgra!6,colframe=cgra,arc=2.5mm,boxrule=0.9pt,left=3mm]{\Large\bfseries\color{cgra}b)}\quad\begin{minipage}[c]{5.3cm}\raggedright\Large $45~\text{m}^3$\end{minipage}\end{tcolorbox}\par\vspace{2.5mm}
\begin{tcolorbox}[enhanced,nobeforeafter,width=7.1cm,height=2.1cm,valign=center,colback=cgra!6,colframe=cgra,arc=2.5mm,boxrule=0.9pt,left=3mm]{\Large\bfseries\color{cgra}c)}\quad\begin{minipage}[c]{5.3cm}\raggedright\Large $4{,}5~\text{m}^3$\end{minipage}\end{tcolorbox}\hfill\begin{tcolorbox}[enhanced,nobeforeafter,width=7.1cm,height=2.1cm,valign=center,colback=cgra!6,colframe=cgra,arc=2.5mm,boxrule=0.9pt,left=3mm]{\Large\bfseries\color{cgra}d)}\quad\begin{minipage}[c]{5.3cm}\raggedright\Large $0{,}45~\text{m}^3$\end{minipage}\end{tcolorbox}
\end{tcolorbox}
\newpage
\begin{tcolorbox}[enhanced,width=16cm,height=14cm,colback=white,colframe=cari,boxrule=2pt,arc=0mm,
left=6mm,right=6mm,top=0mm,bottom=4mm,sharp corners,
overlay={\fill[cari] (frame.north west) rectangle ([yshift=-1.25cm]frame.north east);
\node[anchor=west,text=white,font=\sffamily\bfseries\LARGE] at ([xshift=5mm,yshift=-0.62cm]frame.north west) {Question 1 C};
\node[anchor=east,text=white,font=\sffamily\itshape\large] at ([xshift=-5mm,yshift=-0.62cm]frame.north east) {Arithmétique et puissances};}]
\vspace*{1.75cm}
\begin{minipage}[t][6.9cm][c]{\linewidth}\LARGE\centering
\begin{minipage}{0.95\linewidth}\raggedright\linespread{1.25}\selectfont Lequel de ces nombres est premier~?\end{minipage}
\end{minipage}

\vfill
\begin{tcolorbox}[enhanced,nobeforeafter,width=7.1cm,height=2.1cm,valign=center,colback=cari!6,colframe=cari,arc=2.5mm,boxrule=0.9pt,left=3mm]{\Large\bfseries\color{cari}a)}\quad\begin{minipage}[c]{5.3cm}\raggedright\Large $105$\end{minipage}\end{tcolorbox}\hfill\begin{tcolorbox}[enhanced,nobeforeafter,width=7.1cm,height=2.1cm,valign=center,colback=cari!6,colframe=cari,arc=2.5mm,boxrule=0.9pt,left=3mm]{\Large\bfseries\color{cari}b)}\quad\begin{minipage}[c]{5.3cm}\raggedright\Large $117$\end{minipage}\end{tcolorbox}\par\vspace{2.5mm}
\begin{tcolorbox}[enhanced,nobeforeafter,width=7.1cm,height=2.1cm,valign=center,colback=cari!6,colframe=cari,arc=2.5mm,boxrule=0.9pt,left=3mm]{\Large\bfseries\color{cari}c)}\quad\begin{minipage}[c]{5.3cm}\raggedright\Large $111$\end{minipage}\end{tcolorbox}\hfill\begin{tcolorbox}[enhanced,nobeforeafter,width=7.1cm,height=2.1cm,valign=center,colback=cari!6,colframe=cari,arc=2.5mm,boxrule=0.9pt,left=3mm]{\Large\bfseries\color{cari}d)}\quad\begin{minipage}[c]{5.3cm}\raggedright\Large $103$\end{minipage}\end{tcolorbox}
\end{tcolorbox}
\newpage
\begin{tcolorbox}[enhanced,width=16cm,height=14cm,colback=white,colframe=cari,boxrule=2pt,arc=0mm,
left=6mm,right=6mm,top=0mm,bottom=4mm,sharp corners,
overlay={\fill[cari] (frame.north west) rectangle ([yshift=-1.25cm]frame.north east);
\node[anchor=west,text=white,font=\sffamily\bfseries\LARGE] at ([xshift=5mm,yshift=-0.62cm]frame.north west) {Question 2 C};
\node[anchor=east,text=white,font=\sffamily\itshape\large] at ([xshift=-5mm,yshift=-0.62cm]frame.north east) {Arithmétique et puissances};}]
\vspace*{1.75cm}
\begin{minipage}[t][6.9cm][c]{\linewidth}\LARGE\centering
\begin{minipage}{0.95\linewidth}\raggedright\linespread{1.25}\selectfont Lequel de ces nombres est premier~?\end{minipage}
\end{minipage}

\vfill
\begin{tcolorbox}[enhanced,nobeforeafter,width=7.1cm,height=2.1cm,valign=center,colback=cari!6,colframe=cari,arc=2.5mm,boxrule=0.9pt,left=3mm]{\Large\bfseries\color{cari}a)}\quad\begin{minipage}[c]{5.3cm}\raggedright\Large $49$\end{minipage}\end{tcolorbox}\hfill\begin{tcolorbox}[enhanced,nobeforeafter,width=7.1cm,height=2.1cm,valign=center,colback=cari!6,colframe=cari,arc=2.5mm,boxrule=0.9pt,left=3mm]{\Large\bfseries\color{cari}b)}\quad\begin{minipage}[c]{5.3cm}\raggedright\Large $57$\end{minipage}\end{tcolorbox}\par\vspace{2.5mm}
\begin{tcolorbox}[enhanced,nobeforeafter,width=7.1cm,height=2.1cm,valign=center,colback=cari!6,colframe=cari,arc=2.5mm,boxrule=0.9pt,left=3mm]{\Large\bfseries\color{cari}c)}\quad\begin{minipage}[c]{5.3cm}\raggedright\Large $51$\end{minipage}\end{tcolorbox}\hfill\begin{tcolorbox}[enhanced,nobeforeafter,width=7.1cm,height=2.1cm,valign=center,colback=cari!6,colframe=cari,arc=2.5mm,boxrule=0.9pt,left=3mm]{\Large\bfseries\color{cari}d)}\quad\begin{minipage}[c]{5.3cm}\raggedright\Large $47$\end{minipage}\end{tcolorbox}
\end{tcolorbox}
\newpage
\begin{tcolorbox}[enhanced,width=16cm,height=14cm,colback=white,colframe=cari,boxrule=2pt,arc=0mm,
left=6mm,right=6mm,top=0mm,bottom=4mm,sharp corners,
overlay={\fill[cari] (frame.north west) rectangle ([yshift=-1.25cm]frame.north east);
\node[anchor=west,text=white,font=\sffamily\bfseries\LARGE] at ([xshift=5mm,yshift=-0.62cm]frame.north west) {Question 3 C};
\node[anchor=east,text=white,font=\sffamily\itshape\large] at ([xshift=-5mm,yshift=-0.62cm]frame.north east) {Arithmétique et puissances};}]
\vspace*{1.75cm}
\begin{minipage}[t][6.9cm][c]{\linewidth}\LARGE\centering
\begin{minipage}{0.95\linewidth}\raggedright\linespread{1.25}\selectfont Quelle est la décomposition de $180$ en produit de facteurs premiers~?\end{minipage}
\end{minipage}

\vfill
\begin{tcolorbox}[enhanced,nobeforeafter,width=7.1cm,height=2.1cm,valign=center,colback=cari!6,colframe=cari,arc=2.5mm,boxrule=0.9pt,left=3mm]{\Large\bfseries\color{cari}a)}\quad\begin{minipage}[c]{5.3cm}\raggedright\Large $2^{2}\times3^{2}\times5$\end{minipage}\end{tcolorbox}\hfill\begin{tcolorbox}[enhanced,nobeforeafter,width=7.1cm,height=2.1cm,valign=center,colback=cari!6,colframe=cari,arc=2.5mm,boxrule=0.9pt,left=3mm]{\Large\bfseries\color{cari}b)}\quad\begin{minipage}[c]{5.3cm}\raggedright\Large $12\times15$\end{minipage}\end{tcolorbox}\par\vspace{2.5mm}
\begin{tcolorbox}[enhanced,nobeforeafter,width=7.1cm,height=2.1cm,valign=center,colback=cari!6,colframe=cari,arc=2.5mm,boxrule=0.9pt,left=3mm]{\Large\bfseries\color{cari}c)}\quad\begin{minipage}[c]{5.3cm}\raggedright\large $2^{2}\times3^{2}\times5^{2}$\end{minipage}\end{tcolorbox}\hfill\begin{tcolorbox}[enhanced,nobeforeafter,width=7.1cm,height=2.1cm,valign=center,colback=cari!6,colframe=cari,arc=2.5mm,boxrule=0.9pt,left=3mm]{\Large\bfseries\color{cari}d)}\quad\begin{minipage}[c]{5.3cm}\raggedright\Large $2^{3}\times3^{2}\times5$\end{minipage}\end{tcolorbox}
\end{tcolorbox}
\newpage
\begin{tcolorbox}[enhanced,width=16cm,height=14cm,colback=white,colframe=cari,boxrule=2pt,arc=0mm,
left=6mm,right=6mm,top=0mm,bottom=4mm,sharp corners,
overlay={\fill[cari] (frame.north west) rectangle ([yshift=-1.25cm]frame.north east);
\node[anchor=west,text=white,font=\sffamily\bfseries\LARGE] at ([xshift=5mm,yshift=-0.62cm]frame.north west) {Question 4 C};
\node[anchor=east,text=white,font=\sffamily\itshape\large] at ([xshift=-5mm,yshift=-0.62cm]frame.north east) {Arithmétique et puissances};}]
\vspace*{1.75cm}
\begin{minipage}[t][6.9cm][c]{\linewidth}\LARGE\centering
\begin{minipage}{0.95\linewidth}\raggedright\linespread{1.25}\selectfont Quelle est la décomposition de $132$ en produit de facteurs premiers~?\end{minipage}
\end{minipage}

\vfill
\begin{tcolorbox}[enhanced,nobeforeafter,width=7.1cm,height=2.1cm,valign=center,colback=cari!6,colframe=cari,arc=2.5mm,boxrule=0.9pt,left=3mm]{\Large\bfseries\color{cari}a)}\quad\begin{minipage}[c]{5.3cm}\raggedright\Large $2^{2}\times3\times11$\end{minipage}\end{tcolorbox}\hfill\begin{tcolorbox}[enhanced,nobeforeafter,width=7.1cm,height=2.1cm,valign=center,colback=cari!6,colframe=cari,arc=2.5mm,boxrule=0.9pt,left=3mm]{\Large\bfseries\color{cari}b)}\quad\begin{minipage}[c]{5.3cm}\raggedright\large $2^{2}\times3\times11^{2}$\end{minipage}\end{tcolorbox}\par\vspace{2.5mm}
\begin{tcolorbox}[enhanced,nobeforeafter,width=7.1cm,height=2.1cm,valign=center,colback=cari!6,colframe=cari,arc=2.5mm,boxrule=0.9pt,left=3mm]{\Large\bfseries\color{cari}c)}\quad\begin{minipage}[c]{5.3cm}\raggedright\Large $12\times11$\end{minipage}\end{tcolorbox}\hfill\begin{tcolorbox}[enhanced,nobeforeafter,width=7.1cm,height=2.1cm,valign=center,colback=cari!6,colframe=cari,arc=2.5mm,boxrule=0.9pt,left=3mm]{\Large\bfseries\color{cari}d)}\quad\begin{minipage}[c]{5.3cm}\raggedright\Large $2^{3}\times3\times11$\end{minipage}\end{tcolorbox}
\end{tcolorbox}
\newpage
\begin{tcolorbox}[enhanced,width=16cm,height=14cm,colback=white,colframe=cari,boxrule=2pt,arc=0mm,
left=6mm,right=6mm,top=0mm,bottom=4mm,sharp corners,
overlay={\fill[cari] (frame.north west) rectangle ([yshift=-1.25cm]frame.north east);
\node[anchor=west,text=white,font=\sffamily\bfseries\LARGE] at ([xshift=5mm,yshift=-0.62cm]frame.north west) {Question 5 C};
\node[anchor=east,text=white,font=\sffamily\itshape\large] at ([xshift=-5mm,yshift=-0.62cm]frame.north east) {Arithmétique et puissances};}]
\vspace*{1.75cm}
\begin{minipage}[t][6.9cm][c]{\linewidth}\LARGE\centering
\begin{minipage}{0.95\linewidth}\raggedright\linespread{1.25}\selectfont Quelle est la forme irréductible de $\dfrac{210}{330}$~?\end{minipage}
\end{minipage}

\vfill
\begin{tcolorbox}[enhanced,nobeforeafter,width=7.1cm,height=2.1cm,valign=center,colback=cari!6,colframe=cari,arc=2.5mm,boxrule=0.9pt,left=3mm]{\Large\bfseries\color{cari}a)}\quad\begin{minipage}[c]{5.3cm}\raggedright\Large $\dfrac{42}{66}$\end{minipage}\end{tcolorbox}\hfill\begin{tcolorbox}[enhanced,nobeforeafter,width=7.1cm,height=2.1cm,valign=center,colback=cari!6,colframe=cari,arc=2.5mm,boxrule=0.9pt,left=3mm]{\Large\bfseries\color{cari}b)}\quad\begin{minipage}[c]{5.3cm}\raggedright\Large $\dfrac{7}{11}$\end{minipage}\end{tcolorbox}\par\vspace{2.5mm}
\begin{tcolorbox}[enhanced,nobeforeafter,width=7.1cm,height=2.1cm,valign=center,colback=cari!6,colframe=cari,arc=2.5mm,boxrule=0.9pt,left=3mm]{\Large\bfseries\color{cari}c)}\quad\begin{minipage}[c]{5.3cm}\raggedright\Large $\dfrac{70}{110}$\end{minipage}\end{tcolorbox}\hfill\begin{tcolorbox}[enhanced,nobeforeafter,width=7.1cm,height=2.1cm,valign=center,colback=cari!6,colframe=cari,arc=2.5mm,boxrule=0.9pt,left=3mm]{\Large\bfseries\color{cari}d)}\quad\begin{minipage}[c]{5.3cm}\raggedright\Large $\dfrac{105}{165}$\end{minipage}\end{tcolorbox}
\end{tcolorbox}
\newpage
\begin{tcolorbox}[enhanced,width=16cm,height=14cm,colback=white,colframe=cari,boxrule=2pt,arc=0mm,
left=6mm,right=6mm,top=0mm,bottom=4mm,sharp corners,
overlay={\fill[cari] (frame.north west) rectangle ([yshift=-1.25cm]frame.north east);
\node[anchor=west,text=white,font=\sffamily\bfseries\LARGE] at ([xshift=5mm,yshift=-0.62cm]frame.north west) {Question 6 C};
\node[anchor=east,text=white,font=\sffamily\itshape\large] at ([xshift=-5mm,yshift=-0.62cm]frame.north east) {Arithmétique et puissances};}]
\vspace*{1.75cm}
\begin{minipage}[t][6.9cm][c]{\linewidth}\LARGE\centering
\begin{minipage}{0.95\linewidth}\raggedright\linespread{1.25}\selectfont Quelle est la forme irréductible de $\dfrac{132}{220}$~?\end{minipage}
\end{minipage}

\vfill
\begin{tcolorbox}[enhanced,nobeforeafter,width=7.1cm,height=2.1cm,valign=center,colback=cari!6,colframe=cari,arc=2.5mm,boxrule=0.9pt,left=3mm]{\Large\bfseries\color{cari}a)}\quad\begin{minipage}[c]{5.3cm}\raggedright\Large $\dfrac{66}{110}$\end{minipage}\end{tcolorbox}\hfill\begin{tcolorbox}[enhanced,nobeforeafter,width=7.1cm,height=2.1cm,valign=center,colback=cari!6,colframe=cari,arc=2.5mm,boxrule=0.9pt,left=3mm]{\Large\bfseries\color{cari}b)}\quad\begin{minipage}[c]{5.3cm}\raggedright\Large $\dfrac{4}{6}$\end{minipage}\end{tcolorbox}\par\vspace{2.5mm}
\begin{tcolorbox}[enhanced,nobeforeafter,width=7.1cm,height=2.1cm,valign=center,colback=cari!6,colframe=cari,arc=2.5mm,boxrule=0.9pt,left=3mm]{\Large\bfseries\color{cari}c)}\quad\begin{minipage}[c]{5.3cm}\raggedright\Large $\dfrac{3}{5}$\end{minipage}\end{tcolorbox}\hfill\begin{tcolorbox}[enhanced,nobeforeafter,width=7.1cm,height=2.1cm,valign=center,colback=cari!6,colframe=cari,arc=2.5mm,boxrule=0.9pt,left=3mm]{\Large\bfseries\color{cari}d)}\quad\begin{minipage}[c]{5.3cm}\raggedright\Large $\dfrac{5}{3}$\end{minipage}\end{tcolorbox}
\end{tcolorbox}
\newpage
\begin{tcolorbox}[enhanced,width=16cm,height=14cm,colback=white,colframe=cari,boxrule=2pt,arc=0mm,
left=6mm,right=6mm,top=0mm,bottom=4mm,sharp corners,
overlay={\fill[cari] (frame.north west) rectangle ([yshift=-1.25cm]frame.north east);
\node[anchor=west,text=white,font=\sffamily\bfseries\LARGE] at ([xshift=5mm,yshift=-0.62cm]frame.north west) {Question 7 C};
\node[anchor=east,text=white,font=\sffamily\itshape\large] at ([xshift=-5mm,yshift=-0.62cm]frame.north east) {Arithmétique et puissances};}]
\vspace*{1.75cm}
\begin{minipage}[t][6.9cm][c]{\linewidth}\LARGE\centering
\begin{minipage}{0.95\linewidth}\raggedright\linespread{1.25}\selectfont Quel est le plus grand diviseur commun à $48$ et $180$~?\end{minipage}
\end{minipage}

\vfill
\begin{tcolorbox}[enhanced,nobeforeafter,width=7.1cm,height=2.1cm,valign=center,colback=cari!6,colframe=cari,arc=2.5mm,boxrule=0.9pt,left=3mm]{\Large\bfseries\color{cari}a)}\quad\begin{minipage}[c]{5.3cm}\raggedright\Large $12$\end{minipage}\end{tcolorbox}\hfill\begin{tcolorbox}[enhanced,nobeforeafter,width=7.1cm,height=2.1cm,valign=center,colback=cari!6,colframe=cari,arc=2.5mm,boxrule=0.9pt,left=3mm]{\Large\bfseries\color{cari}b)}\quad\begin{minipage}[c]{5.3cm}\raggedright\Large $720$\end{minipage}\end{tcolorbox}\par\vspace{2.5mm}
\begin{tcolorbox}[enhanced,nobeforeafter,width=7.1cm,height=2.1cm,valign=center,colback=cari!6,colframe=cari,arc=2.5mm,boxrule=0.9pt,left=3mm]{\Large\bfseries\color{cari}c)}\quad\begin{minipage}[c]{5.3cm}\raggedright\Large $24$\end{minipage}\end{tcolorbox}\hfill\begin{tcolorbox}[enhanced,nobeforeafter,width=7.1cm,height=2.1cm,valign=center,colback=cari!6,colframe=cari,arc=2.5mm,boxrule=0.9pt,left=3mm]{\Large\bfseries\color{cari}d)}\quad\begin{minipage}[c]{5.3cm}\raggedright\Large $6$\end{minipage}\end{tcolorbox}
\end{tcolorbox}
\newpage
\begin{tcolorbox}[enhanced,width=16cm,height=14cm,colback=white,colframe=cari,boxrule=2pt,arc=0mm,
left=6mm,right=6mm,top=0mm,bottom=4mm,sharp corners,
overlay={\fill[cari] (frame.north west) rectangle ([yshift=-1.25cm]frame.north east);
\node[anchor=west,text=white,font=\sffamily\bfseries\LARGE] at ([xshift=5mm,yshift=-0.62cm]frame.north west) {Question 8 C};
\node[anchor=east,text=white,font=\sffamily\itshape\large] at ([xshift=-5mm,yshift=-0.62cm]frame.north east) {Arithmétique et puissances};}]
\vspace*{1.75cm}
\begin{minipage}[t][6.9cm][c]{\linewidth}\LARGE\centering
\begin{minipage}{0.95\linewidth}\raggedright\linespread{1.25}\selectfont Quel est le plus grand diviseur commun à $96$ et $144$~?\end{minipage}
\end{minipage}

\vfill
\begin{tcolorbox}[enhanced,nobeforeafter,width=7.1cm,height=2.1cm,valign=center,colback=cari!6,colframe=cari,arc=2.5mm,boxrule=0.9pt,left=3mm]{\Large\bfseries\color{cari}a)}\quad\begin{minipage}[c]{5.3cm}\raggedright\Large $96$\end{minipage}\end{tcolorbox}\hfill\begin{tcolorbox}[enhanced,nobeforeafter,width=7.1cm,height=2.1cm,valign=center,colback=cari!6,colframe=cari,arc=2.5mm,boxrule=0.9pt,left=3mm]{\Large\bfseries\color{cari}b)}\quad\begin{minipage}[c]{5.3cm}\raggedright\Large $24$\end{minipage}\end{tcolorbox}\par\vspace{2.5mm}
\begin{tcolorbox}[enhanced,nobeforeafter,width=7.1cm,height=2.1cm,valign=center,colback=cari!6,colframe=cari,arc=2.5mm,boxrule=0.9pt,left=3mm]{\Large\bfseries\color{cari}c)}\quad\begin{minipage}[c]{5.3cm}\raggedright\Large $48$\end{minipage}\end{tcolorbox}\hfill\begin{tcolorbox}[enhanced,nobeforeafter,width=7.1cm,height=2.1cm,valign=center,colback=cari!6,colframe=cari,arc=2.5mm,boxrule=0.9pt,left=3mm]{\Large\bfseries\color{cari}d)}\quad\begin{minipage}[c]{5.3cm}\raggedright\Large $288$\end{minipage}\end{tcolorbox}
\end{tcolorbox}
\newpage
\begin{tcolorbox}[enhanced,width=16cm,height=14cm,colback=white,colframe=cari,boxrule=2pt,arc=0mm,
left=6mm,right=6mm,top=0mm,bottom=4mm,sharp corners,
overlay={\fill[cari] (frame.north west) rectangle ([yshift=-1.25cm]frame.north east);
\node[anchor=west,text=white,font=\sffamily\bfseries\LARGE] at ([xshift=5mm,yshift=-0.62cm]frame.north west) {Question 9 C};
\node[anchor=east,text=white,font=\sffamily\itshape\large] at ([xshift=-5mm,yshift=-0.62cm]frame.north east) {Arithmétique et puissances};}]
\vspace*{1.75cm}
\begin{minipage}[t][6.9cm][c]{\linewidth}\Large\centering
\begin{minipage}{0.95\linewidth}\raggedright\linespread{1.25}\selectfont Un fleuriste a $64$ roses et $96$ tulipes. Il veut faire le plus grand nombre de bouquets identiques en utilisant toutes les fleurs. Combien de bouquets peut-il faire~?\end{minipage}
\end{minipage}

\vfill
\begin{tcolorbox}[enhanced,nobeforeafter,width=7.1cm,height=2.1cm,valign=center,colback=cari!6,colframe=cari,arc=2.5mm,boxrule=0.9pt,left=3mm]{\Large\bfseries\color{cari}a)}\quad\begin{minipage}[c]{5.3cm}\raggedright\Large $32$\end{minipage}\end{tcolorbox}\hfill\begin{tcolorbox}[enhanced,nobeforeafter,width=7.1cm,height=2.1cm,valign=center,colback=cari!6,colframe=cari,arc=2.5mm,boxrule=0.9pt,left=3mm]{\Large\bfseries\color{cari}b)}\quad\begin{minipage}[c]{5.3cm}\raggedright\Large $160$\end{minipage}\end{tcolorbox}\par\vspace{2.5mm}
\begin{tcolorbox}[enhanced,nobeforeafter,width=7.1cm,height=2.1cm,valign=center,colback=cari!6,colframe=cari,arc=2.5mm,boxrule=0.9pt,left=3mm]{\Large\bfseries\color{cari}c)}\quad\begin{minipage}[c]{5.3cm}\raggedright\Large $16$\end{minipage}\end{tcolorbox}\hfill\begin{tcolorbox}[enhanced,nobeforeafter,width=7.1cm,height=2.1cm,valign=center,colback=cari!6,colframe=cari,arc=2.5mm,boxrule=0.9pt,left=3mm]{\Large\bfseries\color{cari}d)}\quad\begin{minipage}[c]{5.3cm}\raggedright\Large $64$\end{minipage}\end{tcolorbox}
\end{tcolorbox}
\newpage
\begin{tcolorbox}[enhanced,width=16cm,height=14cm,colback=white,colframe=cari,boxrule=2pt,arc=0mm,
left=6mm,right=6mm,top=0mm,bottom=4mm,sharp corners,
overlay={\fill[cari] (frame.north west) rectangle ([yshift=-1.25cm]frame.north east);
\node[anchor=west,text=white,font=\sffamily\bfseries\LARGE] at ([xshift=5mm,yshift=-0.62cm]frame.north west) {Question 10 C};
\node[anchor=east,text=white,font=\sffamily\itshape\large] at ([xshift=-5mm,yshift=-0.62cm]frame.north east) {Arithmétique et puissances};}]
\vspace*{1.75cm}
\begin{minipage}[t][6.9cm][c]{\linewidth}\Large\centering
\begin{minipage}{0.95\linewidth}\raggedright\linespread{1.25}\selectfont Un fleuriste a $54$ roses et $81$ tulipes. Il veut faire le plus grand nombre de bouquets identiques en utilisant toutes les fleurs. Combien de bouquets peut-il faire~?\end{minipage}
\end{minipage}

\vfill
\begin{tcolorbox}[enhanced,nobeforeafter,width=7.1cm,height=2.1cm,valign=center,colback=cari!6,colframe=cari,arc=2.5mm,boxrule=0.9pt,left=3mm]{\Large\bfseries\color{cari}a)}\quad\begin{minipage}[c]{5.3cm}\raggedright\Large $54$\end{minipage}\end{tcolorbox}\hfill\begin{tcolorbox}[enhanced,nobeforeafter,width=7.1cm,height=2.1cm,valign=center,colback=cari!6,colframe=cari,arc=2.5mm,boxrule=0.9pt,left=3mm]{\Large\bfseries\color{cari}b)}\quad\begin{minipage}[c]{5.3cm}\raggedright\Large $27$\end{minipage}\end{tcolorbox}\par\vspace{2.5mm}
\begin{tcolorbox}[enhanced,nobeforeafter,width=7.1cm,height=2.1cm,valign=center,colback=cari!6,colframe=cari,arc=2.5mm,boxrule=0.9pt,left=3mm]{\Large\bfseries\color{cari}c)}\quad\begin{minipage}[c]{5.3cm}\raggedright\Large $135$\end{minipage}\end{tcolorbox}\hfill\begin{tcolorbox}[enhanced,nobeforeafter,width=7.1cm,height=2.1cm,valign=center,colback=cari!6,colframe=cari,arc=2.5mm,boxrule=0.9pt,left=3mm]{\Large\bfseries\color{cari}d)}\quad\begin{minipage}[c]{5.3cm}\raggedright\Large $81$\end{minipage}\end{tcolorbox}
\end{tcolorbox}
\newpage
\begin{tcolorbox}[enhanced,width=16cm,height=14cm,colback=white,colframe=cari,boxrule=2pt,arc=0mm,
left=6mm,right=6mm,top=0mm,bottom=4mm,sharp corners,
overlay={\fill[cari] (frame.north west) rectangle ([yshift=-1.25cm]frame.north east);
\node[anchor=west,text=white,font=\sffamily\bfseries\LARGE] at ([xshift=5mm,yshift=-0.62cm]frame.north west) {Question 11 C};
\node[anchor=east,text=white,font=\sffamily\itshape\large] at ([xshift=-5mm,yshift=-0.62cm]frame.north east) {Arithmétique et puissances};}]
\vspace*{1.75cm}
\begin{minipage}[t][6.9cm][c]{\linewidth}\Large\centering
\begin{minipage}{0.95\linewidth}\raggedright\linespread{1.25}\selectfont Un phare clignote toutes les $8$ secondes et un autre toutes les $14$ secondes. Ils viennent de clignoter ensemble. Dans combien de secondes clignoteront-ils de nouveau ensemble~?\end{minipage}
\end{minipage}

\vfill
\begin{tcolorbox}[enhanced,nobeforeafter,width=7.1cm,height=2.1cm,valign=center,colback=cari!6,colframe=cari,arc=2.5mm,boxrule=0.9pt,left=3mm]{\Large\bfseries\color{cari}a)}\quad\begin{minipage}[c]{5.3cm}\raggedright\Large $112~\text{s}$\end{minipage}\end{tcolorbox}\hfill\begin{tcolorbox}[enhanced,nobeforeafter,width=7.1cm,height=2.1cm,valign=center,colback=cari!6,colframe=cari,arc=2.5mm,boxrule=0.9pt,left=3mm]{\Large\bfseries\color{cari}b)}\quad\begin{minipage}[c]{5.3cm}\raggedright\Large $2~\text{s}$\end{minipage}\end{tcolorbox}\par\vspace{2.5mm}
\begin{tcolorbox}[enhanced,nobeforeafter,width=7.1cm,height=2.1cm,valign=center,colback=cari!6,colframe=cari,arc=2.5mm,boxrule=0.9pt,left=3mm]{\Large\bfseries\color{cari}c)}\quad\begin{minipage}[c]{5.3cm}\raggedright\Large $22~\text{s}$\end{minipage}\end{tcolorbox}\hfill\begin{tcolorbox}[enhanced,nobeforeafter,width=7.1cm,height=2.1cm,valign=center,colback=cari!6,colframe=cari,arc=2.5mm,boxrule=0.9pt,left=3mm]{\Large\bfseries\color{cari}d)}\quad\begin{minipage}[c]{5.3cm}\raggedright\Large $56~\text{s}$\end{minipage}\end{tcolorbox}
\end{tcolorbox}
\newpage
\begin{tcolorbox}[enhanced,width=16cm,height=14cm,colback=white,colframe=cari,boxrule=2pt,arc=0mm,
left=6mm,right=6mm,top=0mm,bottom=4mm,sharp corners,
overlay={\fill[cari] (frame.north west) rectangle ([yshift=-1.25cm]frame.north east);
\node[anchor=west,text=white,font=\sffamily\bfseries\LARGE] at ([xshift=5mm,yshift=-0.62cm]frame.north west) {Question 12 C};
\node[anchor=east,text=white,font=\sffamily\itshape\large] at ([xshift=-5mm,yshift=-0.62cm]frame.north east) {Arithmétique et puissances};}]
\vspace*{1.75cm}
\begin{minipage}[t][6.9cm][c]{\linewidth}\Large\centering
\begin{minipage}{0.95\linewidth}\raggedright\linespread{1.25}\selectfont Un phare clignote toutes les $12$ secondes et un autre toutes les $18$ secondes. Ils viennent de clignoter ensemble. Dans combien de secondes clignoteront-ils de nouveau ensemble~?\end{minipage}
\end{minipage}

\vfill
\begin{tcolorbox}[enhanced,nobeforeafter,width=7.1cm,height=2.1cm,valign=center,colback=cari!6,colframe=cari,arc=2.5mm,boxrule=0.9pt,left=3mm]{\Large\bfseries\color{cari}a)}\quad\begin{minipage}[c]{5.3cm}\raggedright\Large $36~\text{s}$\end{minipage}\end{tcolorbox}\hfill\begin{tcolorbox}[enhanced,nobeforeafter,width=7.1cm,height=2.1cm,valign=center,colback=cari!6,colframe=cari,arc=2.5mm,boxrule=0.9pt,left=3mm]{\Large\bfseries\color{cari}b)}\quad\begin{minipage}[c]{5.3cm}\raggedright\Large $6~\text{s}$\end{minipage}\end{tcolorbox}\par\vspace{2.5mm}
\begin{tcolorbox}[enhanced,nobeforeafter,width=7.1cm,height=2.1cm,valign=center,colback=cari!6,colframe=cari,arc=2.5mm,boxrule=0.9pt,left=3mm]{\Large\bfseries\color{cari}c)}\quad\begin{minipage}[c]{5.3cm}\raggedright\Large $30~\text{s}$\end{minipage}\end{tcolorbox}\hfill\begin{tcolorbox}[enhanced,nobeforeafter,width=7.1cm,height=2.1cm,valign=center,colback=cari!6,colframe=cari,arc=2.5mm,boxrule=0.9pt,left=3mm]{\Large\bfseries\color{cari}d)}\quad\begin{minipage}[c]{5.3cm}\raggedright\Large $216~\text{s}$\end{minipage}\end{tcolorbox}
\end{tcolorbox}
\newpage
\begin{tcolorbox}[enhanced,width=16cm,height=14cm,colback=white,colframe=cari,boxrule=2pt,arc=0mm,
left=6mm,right=6mm,top=0mm,bottom=4mm,sharp corners,
overlay={\fill[cari] (frame.north west) rectangle ([yshift=-1.25cm]frame.north east);
\node[anchor=west,text=white,font=\sffamily\bfseries\LARGE] at ([xshift=5mm,yshift=-0.62cm]frame.north west) {Question 13 C};
\node[anchor=east,text=white,font=\sffamily\itshape\large] at ([xshift=-5mm,yshift=-0.62cm]frame.north east) {Arithmétique et puissances};}]
\vspace*{1.75cm}
\begin{minipage}[t][6.9cm][c]{\linewidth}\LARGE\centering
\begin{minipage}{0.95\linewidth}\raggedright\linespread{1.25}\selectfont Que vaut $(6\times10^{-4})\times(5\times10^{-3})$, en notation scientifique~?\end{minipage}
\end{minipage}

\vfill
\begin{tcolorbox}[enhanced,nobeforeafter,width=7.1cm,height=2.1cm,valign=center,colback=cari!6,colframe=cari,arc=2.5mm,boxrule=0.9pt,left=3mm]{\Large\bfseries\color{cari}a)}\quad\begin{minipage}[c]{5.3cm}\raggedright\Large $30\times10^{-7}$\end{minipage}\end{tcolorbox}\hfill\begin{tcolorbox}[enhanced,nobeforeafter,width=7.1cm,height=2.1cm,valign=center,colback=cari!6,colframe=cari,arc=2.5mm,boxrule=0.9pt,left=3mm]{\Large\bfseries\color{cari}b)}\quad\begin{minipage}[c]{5.3cm}\raggedright\Large $3\times10^{12}$\end{minipage}\end{tcolorbox}\par\vspace{2.5mm}
\begin{tcolorbox}[enhanced,nobeforeafter,width=7.1cm,height=2.1cm,valign=center,colback=cari!6,colframe=cari,arc=2.5mm,boxrule=0.9pt,left=3mm]{\Large\bfseries\color{cari}c)}\quad\begin{minipage}[c]{5.3cm}\raggedright\Large $3\times10^{-7}$\end{minipage}\end{tcolorbox}\hfill\begin{tcolorbox}[enhanced,nobeforeafter,width=7.1cm,height=2.1cm,valign=center,colback=cari!6,colframe=cari,arc=2.5mm,boxrule=0.9pt,left=3mm]{\Large\bfseries\color{cari}d)}\quad\begin{minipage}[c]{5.3cm}\raggedright\Large $3\times10^{-6}$\end{minipage}\end{tcolorbox}
\end{tcolorbox}
\newpage
\begin{tcolorbox}[enhanced,width=16cm,height=14cm,colback=white,colframe=cari,boxrule=2pt,arc=0mm,
left=6mm,right=6mm,top=0mm,bottom=4mm,sharp corners,
overlay={\fill[cari] (frame.north west) rectangle ([yshift=-1.25cm]frame.north east);
\node[anchor=west,text=white,font=\sffamily\bfseries\LARGE] at ([xshift=5mm,yshift=-0.62cm]frame.north west) {Question 14 C};
\node[anchor=east,text=white,font=\sffamily\itshape\large] at ([xshift=-5mm,yshift=-0.62cm]frame.north east) {Arithmétique et puissances};}]
\vspace*{1.75cm}
\begin{minipage}[t][6.9cm][c]{\linewidth}\LARGE\centering
\begin{minipage}{0.95\linewidth}\raggedright\linespread{1.25}\selectfont Que vaut $(2{,}5\times10^{2})\times(4\times10^{5})$, en notation scientifique~?\end{minipage}
\end{minipage}

\vfill
\begin{tcolorbox}[enhanced,nobeforeafter,width=7.1cm,height=2.1cm,valign=center,colback=cari!6,colframe=cari,arc=2.5mm,boxrule=0.9pt,left=3mm]{\Large\bfseries\color{cari}a)}\quad\begin{minipage}[c]{5.3cm}\raggedright\Large $1\times10^{10}$\end{minipage}\end{tcolorbox}\hfill\begin{tcolorbox}[enhanced,nobeforeafter,width=7.1cm,height=2.1cm,valign=center,colback=cari!6,colframe=cari,arc=2.5mm,boxrule=0.9pt,left=3mm]{\Large\bfseries\color{cari}b)}\quad\begin{minipage}[c]{5.3cm}\raggedright\Large $1\times10^{7}$\end{minipage}\end{tcolorbox}\par\vspace{2.5mm}
\begin{tcolorbox}[enhanced,nobeforeafter,width=7.1cm,height=2.1cm,valign=center,colback=cari!6,colframe=cari,arc=2.5mm,boxrule=0.9pt,left=3mm]{\Large\bfseries\color{cari}c)}\quad\begin{minipage}[c]{5.3cm}\raggedright\Large $1\times10^{8}$\end{minipage}\end{tcolorbox}\hfill\begin{tcolorbox}[enhanced,nobeforeafter,width=7.1cm,height=2.1cm,valign=center,colback=cari!6,colframe=cari,arc=2.5mm,boxrule=0.9pt,left=3mm]{\Large\bfseries\color{cari}d)}\quad\begin{minipage}[c]{5.3cm}\raggedright\Large $10\times10^{7}$\end{minipage}\end{tcolorbox}
\end{tcolorbox}
\newpage
\begin{tcolorbox}[enhanced,width=16cm,height=14cm,colback=white,colframe=cari,boxrule=2pt,arc=0mm,
left=6mm,right=6mm,top=0mm,bottom=4mm,sharp corners,
overlay={\fill[cari] (frame.north west) rectangle ([yshift=-1.25cm]frame.north east);
\node[anchor=west,text=white,font=\sffamily\bfseries\LARGE] at ([xshift=5mm,yshift=-0.62cm]frame.north west) {Question 15 C};
\node[anchor=east,text=white,font=\sffamily\itshape\large] at ([xshift=-5mm,yshift=-0.62cm]frame.north east) {Arithmétique et puissances};}]
\vspace*{1.75cm}
\begin{minipage}[t][6.9cm][c]{\linewidth}\LARGE\centering
\begin{minipage}{0.95\linewidth}\raggedright\linespread{1.25}\selectfont Quelle est l'écriture décimale de $2^{-5}$~?\end{minipage}
\end{minipage}

\vfill
\begin{tcolorbox}[enhanced,nobeforeafter,width=7.1cm,height=2.1cm,valign=center,colback=cari!6,colframe=cari,arc=2.5mm,boxrule=0.9pt,left=3mm]{\Large\bfseries\color{cari}a)}\quad\begin{minipage}[c]{5.3cm}\raggedright\Large $0{,}1$\end{minipage}\end{tcolorbox}\hfill\begin{tcolorbox}[enhanced,nobeforeafter,width=7.1cm,height=2.1cm,valign=center,colback=cari!6,colframe=cari,arc=2.5mm,boxrule=0.9pt,left=3mm]{\Large\bfseries\color{cari}b)}\quad\begin{minipage}[c]{5.3cm}\raggedright\Large $0{,}03125$\end{minipage}\end{tcolorbox}\par\vspace{2.5mm}
\begin{tcolorbox}[enhanced,nobeforeafter,width=7.1cm,height=2.1cm,valign=center,colback=cari!6,colframe=cari,arc=2.5mm,boxrule=0.9pt,left=3mm]{\Large\bfseries\color{cari}c)}\quad\begin{minipage}[c]{5.3cm}\raggedright\Large $-10$\end{minipage}\end{tcolorbox}\hfill\begin{tcolorbox}[enhanced,nobeforeafter,width=7.1cm,height=2.1cm,valign=center,colback=cari!6,colframe=cari,arc=2.5mm,boxrule=0.9pt,left=3mm]{\Large\bfseries\color{cari}d)}\quad\begin{minipage}[c]{5.3cm}\raggedright\Large $-32$\end{minipage}\end{tcolorbox}
\end{tcolorbox}
\newpage
\begin{tcolorbox}[enhanced,width=16cm,height=14cm,colback=white,colframe=cari,boxrule=2pt,arc=0mm,
left=6mm,right=6mm,top=0mm,bottom=4mm,sharp corners,
overlay={\fill[cari] (frame.north west) rectangle ([yshift=-1.25cm]frame.north east);
\node[anchor=west,text=white,font=\sffamily\bfseries\LARGE] at ([xshift=5mm,yshift=-0.62cm]frame.north west) {Question 16 C};
\node[anchor=east,text=white,font=\sffamily\itshape\large] at ([xshift=-5mm,yshift=-0.62cm]frame.north east) {Arithmétique et puissances};}]
\vspace*{1.75cm}
\begin{minipage}[t][6.9cm][c]{\linewidth}\LARGE\centering
\begin{minipage}{0.95\linewidth}\raggedright\linespread{1.25}\selectfont Quelle est l'écriture décimale de $20^{-1}$~?\end{minipage}
\end{minipage}

\vfill
\begin{tcolorbox}[enhanced,nobeforeafter,width=7.1cm,height=2.1cm,valign=center,colback=cari!6,colframe=cari,arc=2.5mm,boxrule=0.9pt,left=3mm]{\Large\bfseries\color{cari}a)}\quad\begin{minipage}[c]{5.3cm}\raggedright\Large $0{,}005$\end{minipage}\end{tcolorbox}\hfill\begin{tcolorbox}[enhanced,nobeforeafter,width=7.1cm,height=2.1cm,valign=center,colback=cari!6,colframe=cari,arc=2.5mm,boxrule=0.9pt,left=3mm]{\Large\bfseries\color{cari}b)}\quad\begin{minipage}[c]{5.3cm}\raggedright\Large $0{,}05$\end{minipage}\end{tcolorbox}\par\vspace{2.5mm}
\begin{tcolorbox}[enhanced,nobeforeafter,width=7.1cm,height=2.1cm,valign=center,colback=cari!6,colframe=cari,arc=2.5mm,boxrule=0.9pt,left=3mm]{\Large\bfseries\color{cari}c)}\quad\begin{minipage}[c]{5.3cm}\raggedright\Large $-20$\end{minipage}\end{tcolorbox}\hfill\begin{tcolorbox}[enhanced,nobeforeafter,width=7.1cm,height=2.1cm,valign=center,colback=cari!6,colframe=cari,arc=2.5mm,boxrule=0.9pt,left=3mm]{\Large\bfseries\color{cari}d)}\quad\begin{minipage}[c]{5.3cm}\raggedright\Large $0{,}5$\end{minipage}\end{tcolorbox}
\end{tcolorbox}
\newpage
\begin{tcolorbox}[enhanced,width=16cm,height=14cm,colback=white,colframe=clit,boxrule=2pt,arc=0mm,
left=6mm,right=6mm,top=0mm,bottom=4mm,sharp corners,
overlay={\fill[clit] (frame.north west) rectangle ([yshift=-1.25cm]frame.north east);
\node[anchor=west,text=white,font=\sffamily\bfseries\LARGE] at ([xshift=5mm,yshift=-0.62cm]frame.north west) {Question 17 C};
\node[anchor=east,text=white,font=\sffamily\itshape\large] at ([xshift=-5mm,yshift=-0.62cm]frame.north east) {Calcul littéral et équations};}]
\vspace*{1.75cm}
\begin{minipage}[t][6.9cm][c]{\linewidth}\LARGE\centering
\begin{minipage}{0.95\linewidth}\raggedright\linespread{1.25}\selectfont Développer et réduire $(-2x+1)(3x+4)$.\end{minipage}
\end{minipage}

\vfill
\begin{tcolorbox}[enhanced,nobeforeafter,width=7.1cm,height=2.1cm,valign=center,colback=clit!6,colframe=clit,arc=2.5mm,boxrule=0.9pt,left=3mm]{\Large\bfseries\color{clit}a)}\quad\begin{minipage}[c]{5.3cm}\raggedright\Large $-6x^2-11x+4$\end{minipage}\end{tcolorbox}\hfill\begin{tcolorbox}[enhanced,nobeforeafter,width=7.1cm,height=2.1cm,valign=center,colback=clit!6,colframe=clit,arc=2.5mm,boxrule=0.9pt,left=3mm]{\Large\bfseries\color{clit}b)}\quad\begin{minipage}[c]{5.3cm}\raggedright\Large $-6x^2-5x+4$\end{minipage}\end{tcolorbox}\par\vspace{2.5mm}
\begin{tcolorbox}[enhanced,nobeforeafter,width=7.1cm,height=2.1cm,valign=center,colback=clit!6,colframe=clit,arc=2.5mm,boxrule=0.9pt,left=3mm]{\Large\bfseries\color{clit}c)}\quad\begin{minipage}[c]{5.3cm}\raggedright\Large $x^2-5x+4$\end{minipage}\end{tcolorbox}\hfill\begin{tcolorbox}[enhanced,nobeforeafter,width=7.1cm,height=2.1cm,valign=center,colback=clit!6,colframe=clit,arc=2.5mm,boxrule=0.9pt,left=3mm]{\Large\bfseries\color{clit}d)}\quad\begin{minipage}[c]{5.3cm}\raggedright\Large $-6x^2+4$\end{minipage}\end{tcolorbox}
\end{tcolorbox}
\newpage
\begin{tcolorbox}[enhanced,width=16cm,height=14cm,colback=white,colframe=clit,boxrule=2pt,arc=0mm,
left=6mm,right=6mm,top=0mm,bottom=4mm,sharp corners,
overlay={\fill[clit] (frame.north west) rectangle ([yshift=-1.25cm]frame.north east);
\node[anchor=west,text=white,font=\sffamily\bfseries\LARGE] at ([xshift=5mm,yshift=-0.62cm]frame.north west) {Question 18 C};
\node[anchor=east,text=white,font=\sffamily\itshape\large] at ([xshift=-5mm,yshift=-0.62cm]frame.north east) {Calcul littéral et équations};}]
\vspace*{1.75cm}
\begin{minipage}[t][6.9cm][c]{\linewidth}\LARGE\centering
\begin{minipage}{0.95\linewidth}\raggedright\linespread{1.25}\selectfont Développer et réduire $(4x+5)(2x-1)$.\end{minipage}
\end{minipage}

\vfill
\begin{tcolorbox}[enhanced,nobeforeafter,width=7.1cm,height=2.1cm,valign=center,colback=clit!6,colframe=clit,arc=2.5mm,boxrule=0.9pt,left=3mm]{\Large\bfseries\color{clit}a)}\quad\begin{minipage}[c]{5.3cm}\raggedright\Large $8x^2+6x-5$\end{minipage}\end{tcolorbox}\hfill\begin{tcolorbox}[enhanced,nobeforeafter,width=7.1cm,height=2.1cm,valign=center,colback=clit!6,colframe=clit,arc=2.5mm,boxrule=0.9pt,left=3mm]{\Large\bfseries\color{clit}b)}\quad\begin{minipage}[c]{5.3cm}\raggedright\Large $8x^2-5$\end{minipage}\end{tcolorbox}\par\vspace{2.5mm}
\begin{tcolorbox}[enhanced,nobeforeafter,width=7.1cm,height=2.1cm,valign=center,colback=clit!6,colframe=clit,arc=2.5mm,boxrule=0.9pt,left=3mm]{\Large\bfseries\color{clit}c)}\quad\begin{minipage}[c]{5.3cm}\raggedright\Large $8x^2-14x-5$\end{minipage}\end{tcolorbox}\hfill\begin{tcolorbox}[enhanced,nobeforeafter,width=7.1cm,height=2.1cm,valign=center,colback=clit!6,colframe=clit,arc=2.5mm,boxrule=0.9pt,left=3mm]{\Large\bfseries\color{clit}d)}\quad\begin{minipage}[c]{5.3cm}\raggedright\Large $6x^2+6x-5$\end{minipage}\end{tcolorbox}
\end{tcolorbox}
\newpage
\begin{tcolorbox}[enhanced,width=16cm,height=14cm,colback=white,colframe=clit,boxrule=2pt,arc=0mm,
left=6mm,right=6mm,top=0mm,bottom=4mm,sharp corners,
overlay={\fill[clit] (frame.north west) rectangle ([yshift=-1.25cm]frame.north east);
\node[anchor=west,text=white,font=\sffamily\bfseries\LARGE] at ([xshift=5mm,yshift=-0.62cm]frame.north west) {Question 19 C};
\node[anchor=east,text=white,font=\sffamily\itshape\large] at ([xshift=-5mm,yshift=-0.62cm]frame.north east) {Calcul littéral et équations};}]
\vspace*{1.75cm}
\begin{minipage}[t][6.9cm][c]{\linewidth}\LARGE\centering
\begin{minipage}{0.95\linewidth}\raggedright\linespread{1.25}\selectfont Factoriser $25x^2-9$.\end{minipage}
\end{minipage}

\vfill
\begin{tcolorbox}[enhanced,nobeforeafter,width=7.1cm,height=2.1cm,valign=center,colback=clit!6,colframe=clit,arc=2.5mm,boxrule=0.9pt,left=3mm]{\Large\bfseries\color{clit}a)}\quad\begin{minipage}[c]{5.3cm}\raggedright\Large $(5x-3)^2$\end{minipage}\end{tcolorbox}\hfill\begin{tcolorbox}[enhanced,nobeforeafter,width=7.1cm,height=2.1cm,valign=center,colback=clit!6,colframe=clit,arc=2.5mm,boxrule=0.9pt,left=3mm]{\Large\bfseries\color{clit}b)}\quad\begin{minipage}[c]{5.3cm}\raggedright\Large $(5x-3)(5x+3)$\end{minipage}\end{tcolorbox}\par\vspace{2.5mm}
\begin{tcolorbox}[enhanced,nobeforeafter,width=7.1cm,height=2.1cm,valign=center,colback=clit!6,colframe=clit,arc=2.5mm,boxrule=0.9pt,left=3mm]{\Large\bfseries\color{clit}c)}\quad\begin{minipage}[c]{5.3cm}\raggedright\Large $(5x-9)(5x+9)$\end{minipage}\end{tcolorbox}\hfill\begin{tcolorbox}[enhanced,nobeforeafter,width=7.1cm,height=2.1cm,valign=center,colback=clit!6,colframe=clit,arc=2.5mm,boxrule=0.9pt,left=3mm]{\Large\bfseries\color{clit}d)}\quad\begin{minipage}[c]{5.3cm}\raggedright\Large $(5x+3)^2$\end{minipage}\end{tcolorbox}
\end{tcolorbox}
\newpage
\begin{tcolorbox}[enhanced,width=16cm,height=14cm,colback=white,colframe=clit,boxrule=2pt,arc=0mm,
left=6mm,right=6mm,top=0mm,bottom=4mm,sharp corners,
overlay={\fill[clit] (frame.north west) rectangle ([yshift=-1.25cm]frame.north east);
\node[anchor=west,text=white,font=\sffamily\bfseries\LARGE] at ([xshift=5mm,yshift=-0.62cm]frame.north west) {Question 20 C};
\node[anchor=east,text=white,font=\sffamily\itshape\large] at ([xshift=-5mm,yshift=-0.62cm]frame.north east) {Calcul littéral et équations};}]
\vspace*{1.75cm}
\begin{minipage}[t][6.9cm][c]{\linewidth}\LARGE\centering
\begin{minipage}{0.95\linewidth}\raggedright\linespread{1.25}\selectfont Factoriser $16x^2-1$.\end{minipage}
\end{minipage}

\vfill
\begin{tcolorbox}[enhanced,nobeforeafter,width=7.1cm,height=2.1cm,valign=center,colback=clit!6,colframe=clit,arc=2.5mm,boxrule=0.9pt,left=3mm]{\Large\bfseries\color{clit}a)}\quad\begin{minipage}[c]{5.3cm}\raggedright\Large $(4x+1)^2$\end{minipage}\end{tcolorbox}\hfill\begin{tcolorbox}[enhanced,nobeforeafter,width=7.1cm,height=2.1cm,valign=center,colback=clit!6,colframe=clit,arc=2.5mm,boxrule=0.9pt,left=3mm]{\Large\bfseries\color{clit}b)}\quad\begin{minipage}[c]{5.3cm}\raggedright\Large $(4x-1)^2$\end{minipage}\end{tcolorbox}\par\vspace{2.5mm}
\begin{tcolorbox}[enhanced,nobeforeafter,width=7.1cm,height=2.1cm,valign=center,colback=clit!6,colframe=clit,arc=2.5mm,boxrule=0.9pt,left=3mm]{\Large\bfseries\color{clit}c)}\quad\begin{minipage}[c]{5.3cm}\raggedright\Large $(4x-1)(4x+1)$\end{minipage}\end{tcolorbox}\hfill\begin{tcolorbox}[enhanced,nobeforeafter,width=7.1cm,height=2.1cm,valign=center,colback=clit!6,colframe=clit,arc=2.5mm,boxrule=0.9pt,left=3mm]{\Large\bfseries\color{clit}d)}\quad\begin{minipage}[c]{5.3cm}\raggedright\Large $(16x-1)(16x+1)$\end{minipage}\end{tcolorbox}
\end{tcolorbox}
\newpage
\begin{tcolorbox}[enhanced,width=16cm,height=14cm,colback=white,colframe=clit,boxrule=2pt,arc=0mm,
left=6mm,right=6mm,top=0mm,bottom=4mm,sharp corners,
overlay={\fill[clit] (frame.north west) rectangle ([yshift=-1.25cm]frame.north east);
\node[anchor=west,text=white,font=\sffamily\bfseries\LARGE] at ([xshift=5mm,yshift=-0.62cm]frame.north west) {Question 21 C};
\node[anchor=east,text=white,font=\sffamily\itshape\large] at ([xshift=-5mm,yshift=-0.62cm]frame.north east) {Calcul littéral et équations};}]
\vspace*{1.75cm}
\begin{minipage}[t][6.9cm][c]{\linewidth}\LARGE\centering
\begin{minipage}{0.95\linewidth}\raggedright\linespread{1.25}\selectfont Quelles sont les solutions de l'équation $(x+7)(3x-2)=0$~?\end{minipage}
\end{minipage}

\vfill
\begin{tcolorbox}[enhanced,nobeforeafter,width=7.1cm,height=2.1cm,valign=center,colback=clit!6,colframe=clit,arc=2.5mm,boxrule=0.9pt,left=3mm]{\Large\bfseries\color{clit}a)}\quad\begin{minipage}[c]{5.3cm}\raggedright\Large $\{-2\,;\,7\}$\end{minipage}\end{tcolorbox}\hfill\begin{tcolorbox}[enhanced,nobeforeafter,width=7.1cm,height=2.1cm,valign=center,colback=clit!6,colframe=clit,arc=2.5mm,boxrule=0.9pt,left=3mm]{\Large\bfseries\color{clit}b)}\quad\begin{minipage}[c]{5.3cm}\raggedright\Large $\{-7\,;\,\dfrac{2}{3}\}$\end{minipage}\end{tcolorbox}\par\vspace{2.5mm}
\begin{tcolorbox}[enhanced,nobeforeafter,width=7.1cm,height=2.1cm,valign=center,colback=clit!6,colframe=clit,arc=2.5mm,boxrule=0.9pt,left=3mm]{\Large\bfseries\color{clit}c)}\quad\begin{minipage}[c]{5.3cm}\raggedright\Large $\{-7\}$\end{minipage}\end{tcolorbox}\hfill\begin{tcolorbox}[enhanced,nobeforeafter,width=7.1cm,height=2.1cm,valign=center,colback=clit!6,colframe=clit,arc=2.5mm,boxrule=0.9pt,left=3mm]{\Large\bfseries\color{clit}d)}\quad\begin{minipage}[c]{5.3cm}\raggedright\Large $\{-\dfrac{2}{3}\,;\,7\}$\end{minipage}\end{tcolorbox}
\end{tcolorbox}
\newpage
\begin{tcolorbox}[enhanced,width=16cm,height=14cm,colback=white,colframe=clit,boxrule=2pt,arc=0mm,
left=6mm,right=6mm,top=0mm,bottom=4mm,sharp corners,
overlay={\fill[clit] (frame.north west) rectangle ([yshift=-1.25cm]frame.north east);
\node[anchor=west,text=white,font=\sffamily\bfseries\LARGE] at ([xshift=5mm,yshift=-0.62cm]frame.north west) {Question 22 C};
\node[anchor=east,text=white,font=\sffamily\itshape\large] at ([xshift=-5mm,yshift=-0.62cm]frame.north east) {Calcul littéral et équations};}]
\vspace*{1.75cm}
\begin{minipage}[t][6.9cm][c]{\linewidth}\LARGE\centering
\begin{minipage}{0.95\linewidth}\raggedright\linespread{1.25}\selectfont Quelles sont les solutions de l'équation $(4x-1)(x+3)=0$~?\end{minipage}
\end{minipage}

\vfill
\begin{tcolorbox}[enhanced,nobeforeafter,width=7.1cm,height=2.1cm,valign=center,colback=clit!6,colframe=clit,arc=2.5mm,boxrule=0.9pt,left=3mm]{\Large\bfseries\color{clit}a)}\quad\begin{minipage}[c]{5.3cm}\raggedright\Large $\{-3\,;\,\dfrac{1}{4}\}$\end{minipage}\end{tcolorbox}\hfill\begin{tcolorbox}[enhanced,nobeforeafter,width=7.1cm,height=2.1cm,valign=center,colback=clit!6,colframe=clit,arc=2.5mm,boxrule=0.9pt,left=3mm]{\Large\bfseries\color{clit}b)}\quad\begin{minipage}[c]{5.3cm}\raggedright\Large $\{\dfrac{1}{4}\}$\end{minipage}\end{tcolorbox}\par\vspace{2.5mm}
\begin{tcolorbox}[enhanced,nobeforeafter,width=7.1cm,height=2.1cm,valign=center,colback=clit!6,colframe=clit,arc=2.5mm,boxrule=0.9pt,left=3mm]{\Large\bfseries\color{clit}c)}\quad\begin{minipage}[c]{5.3cm}\raggedright\Large $\{-\dfrac{1}{4}\,;\,3\}$\end{minipage}\end{tcolorbox}\hfill\begin{tcolorbox}[enhanced,nobeforeafter,width=7.1cm,height=2.1cm,valign=center,colback=clit!6,colframe=clit,arc=2.5mm,boxrule=0.9pt,left=3mm]{\Large\bfseries\color{clit}d)}\quad\begin{minipage}[c]{5.3cm}\raggedright\Large $\{-1\,;\,3\}$\end{minipage}\end{tcolorbox}
\end{tcolorbox}
\newpage
\begin{tcolorbox}[enhanced,width=16cm,height=14cm,colback=white,colframe=clit,boxrule=2pt,arc=0mm,
left=6mm,right=6mm,top=0mm,bottom=4mm,sharp corners,
overlay={\fill[clit] (frame.north west) rectangle ([yshift=-1.25cm]frame.north east);
\node[anchor=west,text=white,font=\sffamily\bfseries\LARGE] at ([xshift=5mm,yshift=-0.62cm]frame.north west) {Question 23 C};
\node[anchor=east,text=white,font=\sffamily\itshape\large] at ([xshift=-5mm,yshift=-0.62cm]frame.north east) {Calcul littéral et équations};}]
\vspace*{1.75cm}
\begin{minipage}[t][6.9cm][c]{\linewidth}\LARGE\centering
\begin{minipage}{0.95\linewidth}\raggedright\linespread{1.25}\selectfont Factoriser $(3x-1)(x+2)+(3x-1)(x+5)$.\end{minipage}
\end{minipage}

\vfill
\begin{tcolorbox}[enhanced,nobeforeafter,width=7.1cm,height=2.1cm,valign=center,colback=clit!6,colframe=clit,arc=2.5mm,boxrule=0.9pt,left=3mm]{\Large\bfseries\color{clit}a)}\quad\begin{minipage}[c]{5.3cm}\raggedright\Large $(3x-1)^2(2x+7)$\end{minipage}\end{tcolorbox}\hfill\begin{tcolorbox}[enhanced,nobeforeafter,width=7.1cm,height=2.1cm,valign=center,colback=clit!6,colframe=clit,arc=2.5mm,boxrule=0.9pt,left=3mm]{\Large\bfseries\color{clit}b)}\quad\begin{minipage}[c]{5.3cm}\raggedright\Large $(3x-1)(x+10)$\end{minipage}\end{tcolorbox}\par\vspace{2.5mm}
\begin{tcolorbox}[enhanced,nobeforeafter,width=7.1cm,height=2.1cm,valign=center,colback=clit!6,colframe=clit,arc=2.5mm,boxrule=0.9pt,left=3mm]{\Large\bfseries\color{clit}c)}\quad\begin{minipage}[c]{5.3cm}\raggedright\Large $(3x-1)(2x+7)$\end{minipage}\end{tcolorbox}\hfill\begin{tcolorbox}[enhanced,nobeforeafter,width=7.1cm,height=2.1cm,valign=center,colback=clit!6,colframe=clit,arc=2.5mm,boxrule=0.9pt,left=3mm]{\Large\bfseries\color{clit}d)}\quad\begin{minipage}[c]{5.3cm}\raggedright\Large $(3x-1)(-3)$\end{minipage}\end{tcolorbox}
\end{tcolorbox}
\newpage
\begin{tcolorbox}[enhanced,width=16cm,height=14cm,colback=white,colframe=clit,boxrule=2pt,arc=0mm,
left=6mm,right=6mm,top=0mm,bottom=4mm,sharp corners,
overlay={\fill[clit] (frame.north west) rectangle ([yshift=-1.25cm]frame.north east);
\node[anchor=west,text=white,font=\sffamily\bfseries\LARGE] at ([xshift=5mm,yshift=-0.62cm]frame.north west) {Question 24 C};
\node[anchor=east,text=white,font=\sffamily\itshape\large] at ([xshift=-5mm,yshift=-0.62cm]frame.north east) {Calcul littéral et équations};}]
\vspace*{1.75cm}
\begin{minipage}[t][6.9cm][c]{\linewidth}\LARGE\centering
\begin{minipage}{0.95\linewidth}\raggedright\linespread{1.25}\selectfont Factoriser $(x+5)(2x+7)+(x+5)(-x+1)$.\end{minipage}
\end{minipage}

\vfill
\begin{tcolorbox}[enhanced,nobeforeafter,width=7.1cm,height=2.1cm,valign=center,colback=clit!6,colframe=clit,arc=2.5mm,boxrule=0.9pt,left=3mm]{\Large\bfseries\color{clit}a)}\quad\begin{minipage}[c]{5.3cm}\raggedright\Large $(x+5)^2(x+8)$\end{minipage}\end{tcolorbox}\hfill\begin{tcolorbox}[enhanced,nobeforeafter,width=7.1cm,height=2.1cm,valign=center,colback=clit!6,colframe=clit,arc=2.5mm,boxrule=0.9pt,left=3mm]{\Large\bfseries\color{clit}b)}\quad\begin{minipage}[c]{5.3cm}\raggedright\Large $(x+5)(3x+6)$\end{minipage}\end{tcolorbox}\par\vspace{2.5mm}
\begin{tcolorbox}[enhanced,nobeforeafter,width=7.1cm,height=2.1cm,valign=center,colback=clit!6,colframe=clit,arc=2.5mm,boxrule=0.9pt,left=3mm]{\Large\bfseries\color{clit}c)}\quad\begin{minipage}[c]{5.3cm}\raggedright\Large $(x+5)(x+8)$\end{minipage}\end{tcolorbox}\hfill\begin{tcolorbox}[enhanced,nobeforeafter,width=7.1cm,height=2.1cm,valign=center,colback=clit!6,colframe=clit,arc=2.5mm,boxrule=0.9pt,left=3mm]{\Large\bfseries\color{clit}d)}\quad\begin{minipage}[c]{5.3cm}\raggedright\Large $(x+5)(-2x+7)$\end{minipage}\end{tcolorbox}
\end{tcolorbox}
\newpage
\begin{tcolorbox}[enhanced,width=16cm,height=14cm,colback=white,colframe=clit,boxrule=2pt,arc=0mm,
left=6mm,right=6mm,top=0mm,bottom=4mm,sharp corners,
overlay={\fill[clit] (frame.north west) rectangle ([yshift=-1.25cm]frame.north east);
\node[anchor=west,text=white,font=\sffamily\bfseries\LARGE] at ([xshift=5mm,yshift=-0.62cm]frame.north west) {Question 25 C};
\node[anchor=east,text=white,font=\sffamily\itshape\large] at ([xshift=-5mm,yshift=-0.62cm]frame.north east) {Calcul littéral et équations};}]
\vspace*{1.75cm}
\begin{minipage}[t][6.9cm][c]{\linewidth}\LARGE\centering
\begin{minipage}{0.95\linewidth}\raggedright\linespread{1.25}\selectfont Quelles sont les solutions de l'inéquation $2x+7<1$~?\end{minipage}
\end{minipage}

\vfill
\begin{tcolorbox}[enhanced,nobeforeafter,width=7.1cm,height=2.1cm,valign=center,colback=clit!6,colframe=clit,arc=2.5mm,boxrule=0.9pt,left=3mm]{\Large\bfseries\color{clit}a)}\quad\begin{minipage}[c]{5.3cm}\raggedright\Large $x<-3$\end{minipage}\end{tcolorbox}\hfill\begin{tcolorbox}[enhanced,nobeforeafter,width=7.1cm,height=2.1cm,valign=center,colback=clit!6,colframe=clit,arc=2.5mm,boxrule=0.9pt,left=3mm]{\Large\bfseries\color{clit}b)}\quad\begin{minipage}[c]{5.3cm}\raggedright\Large $x<4$\end{minipage}\end{tcolorbox}\par\vspace{2.5mm}
\begin{tcolorbox}[enhanced,nobeforeafter,width=7.1cm,height=2.1cm,valign=center,colback=clit!6,colframe=clit,arc=2.5mm,boxrule=0.9pt,left=3mm]{\Large\bfseries\color{clit}c)}\quad\begin{minipage}[c]{5.3cm}\raggedright\Large $x<3$\end{minipage}\end{tcolorbox}\hfill\begin{tcolorbox}[enhanced,nobeforeafter,width=7.1cm,height=2.1cm,valign=center,colback=clit!6,colframe=clit,arc=2.5mm,boxrule=0.9pt,left=3mm]{\Large\bfseries\color{clit}d)}\quad\begin{minipage}[c]{5.3cm}\raggedright\Large $x>-3$\end{minipage}\end{tcolorbox}
\end{tcolorbox}
\newpage
\begin{tcolorbox}[enhanced,width=16cm,height=14cm,colback=white,colframe=clit,boxrule=2pt,arc=0mm,
left=6mm,right=6mm,top=0mm,bottom=4mm,sharp corners,
overlay={\fill[clit] (frame.north west) rectangle ([yshift=-1.25cm]frame.north east);
\node[anchor=west,text=white,font=\sffamily\bfseries\LARGE] at ([xshift=5mm,yshift=-0.62cm]frame.north west) {Question 26 C};
\node[anchor=east,text=white,font=\sffamily\itshape\large] at ([xshift=-5mm,yshift=-0.62cm]frame.north east) {Calcul littéral et équations};}]
\vspace*{1.75cm}
\begin{minipage}[t][6.9cm][c]{\linewidth}\LARGE\centering
\begin{minipage}{0.95\linewidth}\raggedright\linespread{1.25}\selectfont Quelles sont les solutions de l'inéquation $-x+6<2$~?\end{minipage}
\end{minipage}

\vfill
\begin{tcolorbox}[enhanced,nobeforeafter,width=7.1cm,height=2.1cm,valign=center,colback=clit!6,colframe=clit,arc=2.5mm,boxrule=0.9pt,left=3mm]{\Large\bfseries\color{clit}a)}\quad\begin{minipage}[c]{5.3cm}\raggedright\Large $x>-8$\end{minipage}\end{tcolorbox}\hfill\begin{tcolorbox}[enhanced,nobeforeafter,width=7.1cm,height=2.1cm,valign=center,colback=clit!6,colframe=clit,arc=2.5mm,boxrule=0.9pt,left=3mm]{\Large\bfseries\color{clit}b)}\quad\begin{minipage}[c]{5.3cm}\raggedright\Large $x>-4$\end{minipage}\end{tcolorbox}\par\vspace{2.5mm}
\begin{tcolorbox}[enhanced,nobeforeafter,width=7.1cm,height=2.1cm,valign=center,colback=clit!6,colframe=clit,arc=2.5mm,boxrule=0.9pt,left=3mm]{\Large\bfseries\color{clit}c)}\quad\begin{minipage}[c]{5.3cm}\raggedright\Large $x<4$\end{minipage}\end{tcolorbox}\hfill\begin{tcolorbox}[enhanced,nobeforeafter,width=7.1cm,height=2.1cm,valign=center,colback=clit!6,colframe=clit,arc=2.5mm,boxrule=0.9pt,left=3mm]{\Large\bfseries\color{clit}d)}\quad\begin{minipage}[c]{5.3cm}\raggedright\Large $x>4$\end{minipage}\end{tcolorbox}
\end{tcolorbox}
\newpage
\begin{tcolorbox}[enhanced,width=16cm,height=14cm,colback=white,colframe=clit,boxrule=2pt,arc=0mm,
left=6mm,right=6mm,top=0mm,bottom=4mm,sharp corners,
overlay={\fill[clit] (frame.north west) rectangle ([yshift=-1.25cm]frame.north east);
\node[anchor=west,text=white,font=\sffamily\bfseries\LARGE] at ([xshift=5mm,yshift=-0.62cm]frame.north west) {Question 27 C};
\node[anchor=east,text=white,font=\sffamily\itshape\large] at ([xshift=-5mm,yshift=-0.62cm]frame.north east) {Calcul littéral et équations};}]
\vspace*{1.75cm}
\begin{minipage}[t][6.9cm][c]{\linewidth}\LARGE\centering
\begin{minipage}{0.95\linewidth}\raggedright\linespread{1.25}\selectfont Quelles sont les solutions de l'équation $x^2=7$~?\end{minipage}
\end{minipage}

\vfill
\begin{tcolorbox}[enhanced,nobeforeafter,width=7.1cm,height=2.1cm,valign=center,colback=clit!6,colframe=clit,arc=2.5mm,boxrule=0.9pt,left=3mm]{\Large\bfseries\color{clit}a)}\quad\begin{minipage}[c]{5.3cm}\raggedright\Large $\{\sqrt{7}\}$\end{minipage}\end{tcolorbox}\hfill\begin{tcolorbox}[enhanced,nobeforeafter,width=7.1cm,height=2.1cm,valign=center,colback=clit!6,colframe=clit,arc=2.5mm,boxrule=0.9pt,left=3mm]{\Large\bfseries\color{clit}b)}\quad\begin{minipage}[c]{5.3cm}\raggedright\Large aucune solution\end{minipage}\end{tcolorbox}\par\vspace{2.5mm}
\begin{tcolorbox}[enhanced,nobeforeafter,width=7.1cm,height=2.1cm,valign=center,colback=clit!6,colframe=clit,arc=2.5mm,boxrule=0.9pt,left=3mm]{\Large\bfseries\color{clit}c)}\quad\begin{minipage}[c]{5.3cm}\raggedright\large $\{-\sqrt{7}\,;\,\sqrt{7}\}$\end{minipage}\end{tcolorbox}\hfill\begin{tcolorbox}[enhanced,nobeforeafter,width=7.1cm,height=2.1cm,valign=center,colback=clit!6,colframe=clit,arc=2.5mm,boxrule=0.9pt,left=3mm]{\Large\bfseries\color{clit}d)}\quad\begin{minipage}[c]{5.3cm}\raggedright\Large $\{-3{,}5\,;\,3{,}5\}$\end{minipage}\end{tcolorbox}
\end{tcolorbox}
\newpage
\begin{tcolorbox}[enhanced,width=16cm,height=14cm,colback=white,colframe=clit,boxrule=2pt,arc=0mm,
left=6mm,right=6mm,top=0mm,bottom=4mm,sharp corners,
overlay={\fill[clit] (frame.north west) rectangle ([yshift=-1.25cm]frame.north east);
\node[anchor=west,text=white,font=\sffamily\bfseries\LARGE] at ([xshift=5mm,yshift=-0.62cm]frame.north west) {Question 28 C};
\node[anchor=east,text=white,font=\sffamily\itshape\large] at ([xshift=-5mm,yshift=-0.62cm]frame.north east) {Calcul littéral et équations};}]
\vspace*{1.75cm}
\begin{minipage}[t][6.9cm][c]{\linewidth}\LARGE\centering
\begin{minipage}{0.95\linewidth}\raggedright\linespread{1.25}\selectfont Quelles sont les solutions de l'équation $x^2=0$~?\end{minipage}
\end{minipage}

\vfill
\begin{tcolorbox}[enhanced,nobeforeafter,width=7.1cm,height=2.1cm,valign=center,colback=clit!6,colframe=clit,arc=2.5mm,boxrule=0.9pt,left=3mm]{\Large\bfseries\color{clit}a)}\quad\begin{minipage}[c]{5.3cm}\raggedright\Large $\{0\,;\,1\}$\end{minipage}\end{tcolorbox}\hfill\begin{tcolorbox}[enhanced,nobeforeafter,width=7.1cm,height=2.1cm,valign=center,colback=clit!6,colframe=clit,arc=2.5mm,boxrule=0.9pt,left=3mm]{\Large\bfseries\color{clit}b)}\quad\begin{minipage}[c]{5.3cm}\raggedright\Large aucune solution\end{minipage}\end{tcolorbox}\par\vspace{2.5mm}
\begin{tcolorbox}[enhanced,nobeforeafter,width=7.1cm,height=2.1cm,valign=center,colback=clit!6,colframe=clit,arc=2.5mm,boxrule=0.9pt,left=3mm]{\Large\bfseries\color{clit}c)}\quad\begin{minipage}[c]{5.3cm}\raggedright\Large $\{0\}$\end{minipage}\end{tcolorbox}\hfill\begin{tcolorbox}[enhanced,nobeforeafter,width=7.1cm,height=2.1cm,valign=center,colback=clit!6,colframe=clit,arc=2.5mm,boxrule=0.9pt,left=3mm]{\Large\bfseries\color{clit}d)}\quad\begin{minipage}[c]{5.3cm}\raggedright\Large $\{-1\,;\,1\}$\end{minipage}\end{tcolorbox}
\end{tcolorbox}
\newpage
\begin{tcolorbox}[enhanced,width=16cm,height=14cm,colback=white,colframe=clit,boxrule=2pt,arc=0mm,
left=6mm,right=6mm,top=0mm,bottom=4mm,sharp corners,
overlay={\fill[clit] (frame.north west) rectangle ([yshift=-1.25cm]frame.north east);
\node[anchor=west,text=white,font=\sffamily\bfseries\LARGE] at ([xshift=5mm,yshift=-0.62cm]frame.north west) {Question 29 C};
\node[anchor=east,text=white,font=\sffamily\itshape\large] at ([xshift=-5mm,yshift=-0.62cm]frame.north east) {Calcul littéral et équations};}]
\vspace*{1.75cm}
\begin{minipage}[t][6.9cm][c]{\linewidth}\LARGE\centering
\begin{minipage}{0.95\linewidth}\raggedright\linespread{1.25}\selectfont Le nombre $4$ est-il solution de l'équation $x^2-x-10=0$~?\end{minipage}
\end{minipage}

\vfill
\begin{tcolorbox}[enhanced,nobeforeafter,width=7.1cm,height=2.1cm,valign=center,colback=clit!6,colframe=clit,arc=2.5mm,boxrule=0.9pt,left=3mm]{\Large\bfseries\color{clit}a)}\quad\begin{minipage}[c]{5.3cm}\raggedright\Large on ne peut pas savoir\end{minipage}\end{tcolorbox}\hfill\begin{tcolorbox}[enhanced,nobeforeafter,width=7.1cm,height=2.1cm,valign=center,colback=clit!6,colframe=clit,arc=2.5mm,boxrule=0.9pt,left=3mm]{\Large\bfseries\color{clit}b)}\quad\begin{minipage}[c]{5.3cm}\raggedright\Large non\end{minipage}\end{tcolorbox}\par\vspace{2.5mm}
\begin{tcolorbox}[enhanced,nobeforeafter,width=7.1cm,height=2.1cm,valign=center,colback=clit!6,colframe=clit,arc=2.5mm,boxrule=0.9pt,left=3mm]{\Large\bfseries\color{clit}c)}\quad\begin{minipage}[c]{5.3cm}\raggedright\Large oui\end{minipage}\end{tcolorbox}\hfill\begin{tcolorbox}[enhanced,nobeforeafter,width=7.1cm,height=2.1cm,valign=center,colback=clit!6,colframe=clit,arc=2.5mm,boxrule=0.9pt,left=3mm]{\Large\bfseries\color{clit}d)}\quad\begin{minipage}[c]{5.3cm}\raggedright\large seulement si $x$ est positif\end{minipage}\end{tcolorbox}
\end{tcolorbox}
\newpage
\begin{tcolorbox}[enhanced,width=16cm,height=14cm,colback=white,colframe=clit,boxrule=2pt,arc=0mm,
left=6mm,right=6mm,top=0mm,bottom=4mm,sharp corners,
overlay={\fill[clit] (frame.north west) rectangle ([yshift=-1.25cm]frame.north east);
\node[anchor=west,text=white,font=\sffamily\bfseries\LARGE] at ([xshift=5mm,yshift=-0.62cm]frame.north west) {Question 30 C};
\node[anchor=east,text=white,font=\sffamily\itshape\large] at ([xshift=-5mm,yshift=-0.62cm]frame.north east) {Calcul littéral et équations};}]
\vspace*{1.75cm}
\begin{minipage}[t][6.9cm][c]{\linewidth}\LARGE\centering
\begin{minipage}{0.95\linewidth}\raggedright\linespread{1.25}\selectfont Le nombre $2$ est-il solution de l'équation $3x^2-2x-8=0$~?\end{minipage}
\end{minipage}

\vfill
\begin{tcolorbox}[enhanced,nobeforeafter,width=7.1cm,height=2.1cm,valign=center,colback=clit!6,colframe=clit,arc=2.5mm,boxrule=0.9pt,left=3mm]{\Large\bfseries\color{clit}a)}\quad\begin{minipage}[c]{5.3cm}\raggedright\Large non\end{minipage}\end{tcolorbox}\hfill\begin{tcolorbox}[enhanced,nobeforeafter,width=7.1cm,height=2.1cm,valign=center,colback=clit!6,colframe=clit,arc=2.5mm,boxrule=0.9pt,left=3mm]{\Large\bfseries\color{clit}b)}\quad\begin{minipage}[c]{5.3cm}\raggedright\Large on ne peut pas savoir\end{minipage}\end{tcolorbox}\par\vspace{2.5mm}
\begin{tcolorbox}[enhanced,nobeforeafter,width=7.1cm,height=2.1cm,valign=center,colback=clit!6,colframe=clit,arc=2.5mm,boxrule=0.9pt,left=3mm]{\Large\bfseries\color{clit}c)}\quad\begin{minipage}[c]{5.3cm}\raggedright\Large oui\end{minipage}\end{tcolorbox}\hfill\begin{tcolorbox}[enhanced,nobeforeafter,width=7.1cm,height=2.1cm,valign=center,colback=clit!6,colframe=clit,arc=2.5mm,boxrule=0.9pt,left=3mm]{\Large\bfseries\color{clit}d)}\quad\begin{minipage}[c]{5.3cm}\raggedright\large seulement si $x$ est positif\end{minipage}\end{tcolorbox}
\end{tcolorbox}
\newpage
\begin{tcolorbox}[enhanced,width=16cm,height=14cm,colback=white,colframe=clit,boxrule=2pt,arc=0mm,
left=6mm,right=6mm,top=0mm,bottom=4mm,sharp corners,
overlay={\fill[clit] (frame.north west) rectangle ([yshift=-1.25cm]frame.north east);
\node[anchor=west,text=white,font=\sffamily\bfseries\LARGE] at ([xshift=5mm,yshift=-0.62cm]frame.north west) {Question 31 C};
\node[anchor=east,text=white,font=\sffamily\itshape\large] at ([xshift=-5mm,yshift=-0.62cm]frame.north east) {Calcul littéral et équations};}]
\vspace*{1.75cm}
\begin{minipage}[t][6.9cm][c]{\linewidth}\Large\centering
\begin{minipage}{0.95\linewidth}\raggedright\linespread{1.25}\selectfont Programme de calcul~: choisir un nombre ; ajouter $2$ ; multiplier le résultat par le nombre choisi ; soustraire $5$. Quel résultat obtient-on en choisissant $-3$~?\end{minipage}
\end{minipage}

\vfill
\begin{tcolorbox}[enhanced,nobeforeafter,width=7.1cm,height=2.1cm,valign=center,colback=clit!6,colframe=clit,arc=2.5mm,boxrule=0.9pt,left=3mm]{\Large\bfseries\color{clit}a)}\quad\begin{minipage}[c]{5.3cm}\raggedright\Large $8$\end{minipage}\end{tcolorbox}\hfill\begin{tcolorbox}[enhanced,nobeforeafter,width=7.1cm,height=2.1cm,valign=center,colback=clit!6,colframe=clit,arc=2.5mm,boxrule=0.9pt,left=3mm]{\Large\bfseries\color{clit}b)}\quad\begin{minipage}[c]{5.3cm}\raggedright\Large $-6$\end{minipage}\end{tcolorbox}\par\vspace{2.5mm}
\begin{tcolorbox}[enhanced,nobeforeafter,width=7.1cm,height=2.1cm,valign=center,colback=clit!6,colframe=clit,arc=2.5mm,boxrule=0.9pt,left=3mm]{\Large\bfseries\color{clit}c)}\quad\begin{minipage}[c]{5.3cm}\raggedright\Large $-2$\end{minipage}\end{tcolorbox}\hfill\begin{tcolorbox}[enhanced,nobeforeafter,width=7.1cm,height=2.1cm,valign=center,colback=clit!6,colframe=clit,arc=2.5mm,boxrule=0.9pt,left=3mm]{\Large\bfseries\color{clit}d)}\quad\begin{minipage}[c]{5.3cm}\raggedright\Large $-14$\end{minipage}\end{tcolorbox}
\end{tcolorbox}
\newpage
\begin{tcolorbox}[enhanced,width=16cm,height=14cm,colback=white,colframe=clit,boxrule=2pt,arc=0mm,
left=6mm,right=6mm,top=0mm,bottom=4mm,sharp corners,
overlay={\fill[clit] (frame.north west) rectangle ([yshift=-1.25cm]frame.north east);
\node[anchor=west,text=white,font=\sffamily\bfseries\LARGE] at ([xshift=5mm,yshift=-0.62cm]frame.north west) {Question 32 C};
\node[anchor=east,text=white,font=\sffamily\itshape\large] at ([xshift=-5mm,yshift=-0.62cm]frame.north east) {Calcul littéral et équations};}]
\vspace*{1.75cm}
\begin{minipage}[t][6.9cm][c]{\linewidth}\Large\centering
\begin{minipage}{0.95\linewidth}\raggedright\linespread{1.25}\selectfont Programme de calcul~: choisir un nombre ; soustraire $3$ ; multiplier le résultat par le nombre choisi ; soustraire $2$. Quel résultat obtient-on en choisissant $7$~?\end{minipage}
\end{minipage}

\vfill
\begin{tcolorbox}[enhanced,nobeforeafter,width=7.1cm,height=2.1cm,valign=center,colback=clit!6,colframe=clit,arc=2.5mm,boxrule=0.9pt,left=3mm]{\Large\bfseries\color{clit}a)}\quad\begin{minipage}[c]{5.3cm}\raggedright\Large $-16$\end{minipage}\end{tcolorbox}\hfill\begin{tcolorbox}[enhanced,nobeforeafter,width=7.1cm,height=2.1cm,valign=center,colback=clit!6,colframe=clit,arc=2.5mm,boxrule=0.9pt,left=3mm]{\Large\bfseries\color{clit}b)}\quad\begin{minipage}[c]{5.3cm}\raggedright\Large $30$\end{minipage}\end{tcolorbox}\par\vspace{2.5mm}
\begin{tcolorbox}[enhanced,nobeforeafter,width=7.1cm,height=2.1cm,valign=center,colback=clit!6,colframe=clit,arc=2.5mm,boxrule=0.9pt,left=3mm]{\Large\bfseries\color{clit}c)}\quad\begin{minipage}[c]{5.3cm}\raggedright\Large $2$\end{minipage}\end{tcolorbox}\hfill\begin{tcolorbox}[enhanced,nobeforeafter,width=7.1cm,height=2.1cm,valign=center,colback=clit!6,colframe=clit,arc=2.5mm,boxrule=0.9pt,left=3mm]{\Large\bfseries\color{clit}d)}\quad\begin{minipage}[c]{5.3cm}\raggedright\Large $26$\end{minipage}\end{tcolorbox}
\end{tcolorbox}
\newpage
\begin{tcolorbox}[enhanced,width=16cm,height=14cm,colback=white,colframe=cfon,boxrule=2pt,arc=0mm,
left=6mm,right=6mm,top=0mm,bottom=4mm,sharp corners,
overlay={\fill[cfon] (frame.north west) rectangle ([yshift=-1.25cm]frame.north east);
\node[anchor=west,text=white,font=\sffamily\bfseries\LARGE] at ([xshift=5mm,yshift=-0.62cm]frame.north west) {Question 33 C};
\node[anchor=east,text=white,font=\sffamily\itshape\large] at ([xshift=-5mm,yshift=-0.62cm]frame.north east) {Fonctions};}]
\vspace*{1.75cm}
\begin{minipage}[t][6.9cm][c]{\linewidth}\LARGE\centering
\begin{minipage}{0.95\linewidth}\raggedright\linespread{1.25}\selectfont Soit $f(x)=6x-9$. Quelle est l'image de $2$ par $f$~?\end{minipage}
\end{minipage}

\vfill
\begin{tcolorbox}[enhanced,nobeforeafter,width=7.1cm,height=2.1cm,valign=center,colback=cfon!6,colframe=cfon,arc=2.5mm,boxrule=0.9pt,left=3mm]{\Large\bfseries\color{cfon}a)}\quad\begin{minipage}[c]{5.3cm}\raggedright\Large $21$\end{minipage}\end{tcolorbox}\hfill\begin{tcolorbox}[enhanced,nobeforeafter,width=7.1cm,height=2.1cm,valign=center,colback=cfon!6,colframe=cfon,arc=2.5mm,boxrule=0.9pt,left=3mm]{\Large\bfseries\color{cfon}b)}\quad\begin{minipage}[c]{5.3cm}\raggedright\Large $-1$\end{minipage}\end{tcolorbox}\par\vspace{2.5mm}
\begin{tcolorbox}[enhanced,nobeforeafter,width=7.1cm,height=2.1cm,valign=center,colback=cfon!6,colframe=cfon,arc=2.5mm,boxrule=0.9pt,left=3mm]{\Large\bfseries\color{cfon}c)}\quad\begin{minipage}[c]{5.3cm}\raggedright\Large $3$\end{minipage}\end{tcolorbox}\hfill\begin{tcolorbox}[enhanced,nobeforeafter,width=7.1cm,height=2.1cm,valign=center,colback=cfon!6,colframe=cfon,arc=2.5mm,boxrule=0.9pt,left=3mm]{\Large\bfseries\color{cfon}d)}\quad\begin{minipage}[c]{5.3cm}\raggedright\Large $\dfrac{11}{6}$\end{minipage}\end{tcolorbox}
\end{tcolorbox}
\newpage
\begin{tcolorbox}[enhanced,width=16cm,height=14cm,colback=white,colframe=cfon,boxrule=2pt,arc=0mm,
left=6mm,right=6mm,top=0mm,bottom=4mm,sharp corners,
overlay={\fill[cfon] (frame.north west) rectangle ([yshift=-1.25cm]frame.north east);
\node[anchor=west,text=white,font=\sffamily\bfseries\LARGE] at ([xshift=5mm,yshift=-0.62cm]frame.north west) {Question 34 C};
\node[anchor=east,text=white,font=\sffamily\itshape\large] at ([xshift=-5mm,yshift=-0.62cm]frame.north east) {Fonctions};}]
\vspace*{1.75cm}
\begin{minipage}[t][6.9cm][c]{\linewidth}\LARGE\centering
\begin{minipage}{0.95\linewidth}\raggedright\linespread{1.25}\selectfont Soit $f(x)=x+8$. Quelle est l'image de $-5$ par $f$~?\end{minipage}
\end{minipage}

\vfill
\begin{tcolorbox}[enhanced,nobeforeafter,width=7.1cm,height=2.1cm,valign=center,colback=cfon!6,colframe=cfon,arc=2.5mm,boxrule=0.9pt,left=3mm]{\Large\bfseries\color{cfon}a)}\quad\begin{minipage}[c]{5.3cm}\raggedright\Large $3$\end{minipage}\end{tcolorbox}\hfill\begin{tcolorbox}[enhanced,nobeforeafter,width=7.1cm,height=2.1cm,valign=center,colback=cfon!6,colframe=cfon,arc=2.5mm,boxrule=0.9pt,left=3mm]{\Large\bfseries\color{cfon}b)}\quad\begin{minipage}[c]{5.3cm}\raggedright\Large $-13$\end{minipage}\end{tcolorbox}\par\vspace{2.5mm}
\begin{tcolorbox}[enhanced,nobeforeafter,width=7.1cm,height=2.1cm,valign=center,colback=cfon!6,colframe=cfon,arc=2.5mm,boxrule=0.9pt,left=3mm]{\Large\bfseries\color{cfon}c)}\quad\begin{minipage}[c]{5.3cm}\raggedright\Large $2$\end{minipage}\end{tcolorbox}\hfill\begin{tcolorbox}[enhanced,nobeforeafter,width=7.1cm,height=2.1cm,valign=center,colback=cfon!6,colframe=cfon,arc=2.5mm,boxrule=0.9pt,left=3mm]{\Large\bfseries\color{cfon}d)}\quad\begin{minipage}[c]{5.3cm}\raggedright\Large $4$\end{minipage}\end{tcolorbox}
\end{tcolorbox}
\newpage
\begin{tcolorbox}[enhanced,width=16cm,height=14cm,colback=white,colframe=cfon,boxrule=2pt,arc=0mm,
left=6mm,right=6mm,top=0mm,bottom=4mm,sharp corners,
overlay={\fill[cfon] (frame.north west) rectangle ([yshift=-1.25cm]frame.north east);
\node[anchor=west,text=white,font=\sffamily\bfseries\LARGE] at ([xshift=5mm,yshift=-0.62cm]frame.north west) {Question 35 C};
\node[anchor=east,text=white,font=\sffamily\itshape\large] at ([xshift=-5mm,yshift=-0.62cm]frame.north east) {Fonctions};}]
\vspace*{1.75cm}
\begin{minipage}[t][6.9cm][c]{\linewidth}\LARGE\centering
\begin{minipage}{0.95\linewidth}\raggedright\linespread{1.25}\selectfont Soit $f(x)=-2x-6$. Quel est l'antécédent de $4$ par $f$~?\end{minipage}
\end{minipage}

\vfill
\begin{tcolorbox}[enhanced,nobeforeafter,width=7.1cm,height=2.1cm,valign=center,colback=cfon!6,colframe=cfon,arc=2.5mm,boxrule=0.9pt,left=3mm]{\Large\bfseries\color{cfon}a)}\quad\begin{minipage}[c]{5.3cm}\raggedright\Large $4$\end{minipage}\end{tcolorbox}\hfill\begin{tcolorbox}[enhanced,nobeforeafter,width=7.1cm,height=2.1cm,valign=center,colback=cfon!6,colframe=cfon,arc=2.5mm,boxrule=0.9pt,left=3mm]{\Large\bfseries\color{cfon}b)}\quad\begin{minipage}[c]{5.3cm}\raggedright\Large $-14$\end{minipage}\end{tcolorbox}\par\vspace{2.5mm}
\begin{tcolorbox}[enhanced,nobeforeafter,width=7.1cm,height=2.1cm,valign=center,colback=cfon!6,colframe=cfon,arc=2.5mm,boxrule=0.9pt,left=3mm]{\Large\bfseries\color{cfon}c)}\quad\begin{minipage}[c]{5.3cm}\raggedright\Large $1$\end{minipage}\end{tcolorbox}\hfill\begin{tcolorbox}[enhanced,nobeforeafter,width=7.1cm,height=2.1cm,valign=center,colback=cfon!6,colframe=cfon,arc=2.5mm,boxrule=0.9pt,left=3mm]{\Large\bfseries\color{cfon}d)}\quad\begin{minipage}[c]{5.3cm}\raggedright\Large $-5$\end{minipage}\end{tcolorbox}
\end{tcolorbox}
\newpage
\begin{tcolorbox}[enhanced,width=16cm,height=14cm,colback=white,colframe=cfon,boxrule=2pt,arc=0mm,
left=6mm,right=6mm,top=0mm,bottom=4mm,sharp corners,
overlay={\fill[cfon] (frame.north west) rectangle ([yshift=-1.25cm]frame.north east);
\node[anchor=west,text=white,font=\sffamily\bfseries\LARGE] at ([xshift=5mm,yshift=-0.62cm]frame.north west) {Question 36 C};
\node[anchor=east,text=white,font=\sffamily\itshape\large] at ([xshift=-5mm,yshift=-0.62cm]frame.north east) {Fonctions};}]
\vspace*{1.75cm}
\begin{minipage}[t][6.9cm][c]{\linewidth}\LARGE\centering
\begin{minipage}{0.95\linewidth}\raggedright\linespread{1.25}\selectfont Soit $f(x)=3x$. Quel est l'antécédent de $12$ par $f$~?\end{minipage}
\end{minipage}

\vfill
\begin{tcolorbox}[enhanced,nobeforeafter,width=7.1cm,height=2.1cm,valign=center,colback=cfon!6,colframe=cfon,arc=2.5mm,boxrule=0.9pt,left=3mm]{\Large\bfseries\color{cfon}a)}\quad\begin{minipage}[c]{5.3cm}\raggedright\Large $4$\end{minipage}\end{tcolorbox}\hfill\begin{tcolorbox}[enhanced,nobeforeafter,width=7.1cm,height=2.1cm,valign=center,colback=cfon!6,colframe=cfon,arc=2.5mm,boxrule=0.9pt,left=3mm]{\Large\bfseries\color{cfon}b)}\quad\begin{minipage}[c]{5.3cm}\raggedright\Large $3$\end{minipage}\end{tcolorbox}\par\vspace{2.5mm}
\begin{tcolorbox}[enhanced,nobeforeafter,width=7.1cm,height=2.1cm,valign=center,colback=cfon!6,colframe=cfon,arc=2.5mm,boxrule=0.9pt,left=3mm]{\Large\bfseries\color{cfon}c)}\quad\begin{minipage}[c]{5.3cm}\raggedright\Large $5$\end{minipage}\end{tcolorbox}\hfill\begin{tcolorbox}[enhanced,nobeforeafter,width=7.1cm,height=2.1cm,valign=center,colback=cfon!6,colframe=cfon,arc=2.5mm,boxrule=0.9pt,left=3mm]{\Large\bfseries\color{cfon}d)}\quad\begin{minipage}[c]{5.3cm}\raggedright\Large $36$\end{minipage}\end{tcolorbox}
\end{tcolorbox}
\newpage
\begin{tcolorbox}[enhanced,width=16cm,height=14cm,colback=white,colframe=cfon,boxrule=2pt,arc=0mm,
left=6mm,right=6mm,top=0mm,bottom=4mm,sharp corners,
overlay={\fill[cfon] (frame.north west) rectangle ([yshift=-1.25cm]frame.north east);
\node[anchor=west,text=white,font=\sffamily\bfseries\LARGE] at ([xshift=5mm,yshift=-0.62cm]frame.north west) {Question 37 C};
\node[anchor=east,text=white,font=\sffamily\itshape\large] at ([xshift=-5mm,yshift=-0.62cm]frame.north east) {Fonctions};}]
\vspace*{1.75cm}
\begin{minipage}[t][6.9cm][c]{\linewidth}\LARGE\centering
\begin{minipage}{0.95\linewidth}\raggedright\linespread{1.25}\selectfont $f$ est une fonction linéaire telle que $f(8)=6$. Quel est son coefficient~?\end{minipage}
\end{minipage}

\vfill
\begin{tcolorbox}[enhanced,nobeforeafter,width=7.1cm,height=2.1cm,valign=center,colback=cfon!6,colframe=cfon,arc=2.5mm,boxrule=0.9pt,left=3mm]{\Large\bfseries\color{cfon}a)}\quad\begin{minipage}[c]{5.3cm}\raggedright\Large $\dfrac{3}{4}$\end{minipage}\end{tcolorbox}\hfill\begin{tcolorbox}[enhanced,nobeforeafter,width=7.1cm,height=2.1cm,valign=center,colback=cfon!6,colframe=cfon,arc=2.5mm,boxrule=0.9pt,left=3mm]{\Large\bfseries\color{cfon}b)}\quad\begin{minipage}[c]{5.3cm}\raggedright\Large $48$\end{minipage}\end{tcolorbox}\par\vspace{2.5mm}
\begin{tcolorbox}[enhanced,nobeforeafter,width=7.1cm,height=2.1cm,valign=center,colback=cfon!6,colframe=cfon,arc=2.5mm,boxrule=0.9pt,left=3mm]{\Large\bfseries\color{cfon}c)}\quad\begin{minipage}[c]{5.3cm}\raggedright\Large $\dfrac{4}{3}$\end{minipage}\end{tcolorbox}\hfill\begin{tcolorbox}[enhanced,nobeforeafter,width=7.1cm,height=2.1cm,valign=center,colback=cfon!6,colframe=cfon,arc=2.5mm,boxrule=0.9pt,left=3mm]{\Large\bfseries\color{cfon}d)}\quad\begin{minipage}[c]{5.3cm}\raggedright\Large $-2$\end{minipage}\end{tcolorbox}
\end{tcolorbox}
\newpage
\begin{tcolorbox}[enhanced,width=16cm,height=14cm,colback=white,colframe=cfon,boxrule=2pt,arc=0mm,
left=6mm,right=6mm,top=0mm,bottom=4mm,sharp corners,
overlay={\fill[cfon] (frame.north west) rectangle ([yshift=-1.25cm]frame.north east);
\node[anchor=west,text=white,font=\sffamily\bfseries\LARGE] at ([xshift=5mm,yshift=-0.62cm]frame.north west) {Question 38 C};
\node[anchor=east,text=white,font=\sffamily\itshape\large] at ([xshift=-5mm,yshift=-0.62cm]frame.north east) {Fonctions};}]
\vspace*{1.75cm}
\begin{minipage}[t][6.9cm][c]{\linewidth}\LARGE\centering
\begin{minipage}{0.95\linewidth}\raggedright\linespread{1.25}\selectfont $f$ est une fonction linéaire telle que $f(-5)=-20$. Quel est son coefficient~?\end{minipage}
\end{minipage}

\vfill
\begin{tcolorbox}[enhanced,nobeforeafter,width=7.1cm,height=2.1cm,valign=center,colback=cfon!6,colframe=cfon,arc=2.5mm,boxrule=0.9pt,left=3mm]{\Large\bfseries\color{cfon}a)}\quad\begin{minipage}[c]{5.3cm}\raggedright\Large $100$\end{minipage}\end{tcolorbox}\hfill\begin{tcolorbox}[enhanced,nobeforeafter,width=7.1cm,height=2.1cm,valign=center,colback=cfon!6,colframe=cfon,arc=2.5mm,boxrule=0.9pt,left=3mm]{\Large\bfseries\color{cfon}b)}\quad\begin{minipage}[c]{5.3cm}\raggedright\Large $-15$\end{minipage}\end{tcolorbox}\par\vspace{2.5mm}
\begin{tcolorbox}[enhanced,nobeforeafter,width=7.1cm,height=2.1cm,valign=center,colback=cfon!6,colframe=cfon,arc=2.5mm,boxrule=0.9pt,left=3mm]{\Large\bfseries\color{cfon}c)}\quad\begin{minipage}[c]{5.3cm}\raggedright\Large $\dfrac{1}{4}$\end{minipage}\end{tcolorbox}\hfill\begin{tcolorbox}[enhanced,nobeforeafter,width=7.1cm,height=2.1cm,valign=center,colback=cfon!6,colframe=cfon,arc=2.5mm,boxrule=0.9pt,left=3mm]{\Large\bfseries\color{cfon}d)}\quad\begin{minipage}[c]{5.3cm}\raggedright\Large $4$\end{minipage}\end{tcolorbox}
\end{tcolorbox}
\newpage
\begin{tcolorbox}[enhanced,width=16cm,height=14cm,colback=white,colframe=cfon,boxrule=2pt,arc=0mm,
left=6mm,right=6mm,top=0mm,bottom=4mm,sharp corners,
overlay={\fill[cfon] (frame.north west) rectangle ([yshift=-1.25cm]frame.north east);
\node[anchor=west,text=white,font=\sffamily\bfseries\LARGE] at ([xshift=5mm,yshift=-0.62cm]frame.north west) {Question 39 C};
\node[anchor=east,text=white,font=\sffamily\itshape\large] at ([xshift=-5mm,yshift=-0.62cm]frame.north east) {Fonctions};}]
\vspace*{1.75cm}
\begin{minipage}[t][6.9cm][c]{\linewidth}\LARGE\centering
\begin{minipage}{0.95\linewidth}\raggedright\linespread{1.25}\selectfont $f$ est une fonction affine telle que $f(4)=2$ et $f(8)=4$. Quel est le coefficient $a$ de $f(x)=ax+b$~?\end{minipage}
\end{minipage}

\vfill
\begin{tcolorbox}[enhanced,nobeforeafter,width=7.1cm,height=2.1cm,valign=center,colback=cfon!6,colframe=cfon,arc=2.5mm,boxrule=0.9pt,left=3mm]{\Large\bfseries\color{cfon}a)}\quad\begin{minipage}[c]{5.3cm}\raggedright\Large $2$\end{minipage}\end{tcolorbox}\hfill\begin{tcolorbox}[enhanced,nobeforeafter,width=7.1cm,height=2.1cm,valign=center,colback=cfon!6,colframe=cfon,arc=2.5mm,boxrule=0.9pt,left=3mm]{\Large\bfseries\color{cfon}b)}\quad\begin{minipage}[c]{5.3cm}\raggedright\Large $\dfrac{3}{2}$\end{minipage}\end{tcolorbox}\par\vspace{2.5mm}
\begin{tcolorbox}[enhanced,nobeforeafter,width=7.1cm,height=2.1cm,valign=center,colback=cfon!6,colframe=cfon,arc=2.5mm,boxrule=0.9pt,left=3mm]{\Large\bfseries\color{cfon}c)}\quad\begin{minipage}[c]{5.3cm}\raggedright\Large $-\dfrac{1}{2}$\end{minipage}\end{tcolorbox}\hfill\begin{tcolorbox}[enhanced,nobeforeafter,width=7.1cm,height=2.1cm,valign=center,colback=cfon!6,colframe=cfon,arc=2.5mm,boxrule=0.9pt,left=3mm]{\Large\bfseries\color{cfon}d)}\quad\begin{minipage}[c]{5.3cm}\raggedright\Large $\dfrac{1}{2}$\end{minipage}\end{tcolorbox}
\end{tcolorbox}
\newpage
\begin{tcolorbox}[enhanced,width=16cm,height=14cm,colback=white,colframe=cfon,boxrule=2pt,arc=0mm,
left=6mm,right=6mm,top=0mm,bottom=4mm,sharp corners,
overlay={\fill[cfon] (frame.north west) rectangle ([yshift=-1.25cm]frame.north east);
\node[anchor=west,text=white,font=\sffamily\bfseries\LARGE] at ([xshift=5mm,yshift=-0.62cm]frame.north west) {Question 40 C};
\node[anchor=east,text=white,font=\sffamily\itshape\large] at ([xshift=-5mm,yshift=-0.62cm]frame.north east) {Fonctions};}]
\vspace*{1.75cm}
\begin{minipage}[t][6.9cm][c]{\linewidth}\LARGE\centering
\begin{minipage}{0.95\linewidth}\raggedright\linespread{1.25}\selectfont $f$ est une fonction affine telle que $f(-2)=-5$ et $f(3)=10$. Quel est le coefficient $a$ de $f(x)=ax+b$~?\end{minipage}
\end{minipage}

\vfill
\begin{tcolorbox}[enhanced,nobeforeafter,width=7.1cm,height=2.1cm,valign=center,colback=cfon!6,colframe=cfon,arc=2.5mm,boxrule=0.9pt,left=3mm]{\Large\bfseries\color{cfon}a)}\quad\begin{minipage}[c]{5.3cm}\raggedright\Large $3$\end{minipage}\end{tcolorbox}\hfill\begin{tcolorbox}[enhanced,nobeforeafter,width=7.1cm,height=2.1cm,valign=center,colback=cfon!6,colframe=cfon,arc=2.5mm,boxrule=0.9pt,left=3mm]{\Large\bfseries\color{cfon}b)}\quad\begin{minipage}[c]{5.3cm}\raggedright\Large $15$\end{minipage}\end{tcolorbox}\par\vspace{2.5mm}
\begin{tcolorbox}[enhanced,nobeforeafter,width=7.1cm,height=2.1cm,valign=center,colback=cfon!6,colframe=cfon,arc=2.5mm,boxrule=0.9pt,left=3mm]{\Large\bfseries\color{cfon}c)}\quad\begin{minipage}[c]{5.3cm}\raggedright\Large $-3$\end{minipage}\end{tcolorbox}\hfill\begin{tcolorbox}[enhanced,nobeforeafter,width=7.1cm,height=2.1cm,valign=center,colback=cfon!6,colframe=cfon,arc=2.5mm,boxrule=0.9pt,left=3mm]{\Large\bfseries\color{cfon}d)}\quad\begin{minipage}[c]{5.3cm}\raggedright\Large $\dfrac{1}{3}$\end{minipage}\end{tcolorbox}
\end{tcolorbox}
\newpage
\begin{tcolorbox}[enhanced,width=16cm,height=14cm,colback=white,colframe=cfon,boxrule=2pt,arc=0mm,
left=6mm,right=6mm,top=0mm,bottom=4mm,sharp corners,
overlay={\fill[cfon] (frame.north west) rectangle ([yshift=-1.25cm]frame.north east);
\node[anchor=west,text=white,font=\sffamily\bfseries\LARGE] at ([xshift=5mm,yshift=-0.62cm]frame.north west) {Question 41 C};
\node[anchor=east,text=white,font=\sffamily\itshape\large] at ([xshift=-5mm,yshift=-0.62cm]frame.north east) {Fonctions};}]
\vspace*{1.75cm}
\begin{minipage}[t][6.9cm][c]{\linewidth}\LARGE\centering
\begin{minipage}{0.95\linewidth}\raggedright\linespread{1.25}\selectfont $f$ est la fonction affine $f(x)=-x+b$ et $f(6)=0$. Que vaut $b$~?\end{minipage}
\end{minipage}

\vfill
\begin{tcolorbox}[enhanced,nobeforeafter,width=7.1cm,height=2.1cm,valign=center,colback=cfon!6,colframe=cfon,arc=2.5mm,boxrule=0.9pt,left=3mm]{\Large\bfseries\color{cfon}a)}\quad\begin{minipage}[c]{5.3cm}\raggedright\Large $7$\end{minipage}\end{tcolorbox}\hfill\begin{tcolorbox}[enhanced,nobeforeafter,width=7.1cm,height=2.1cm,valign=center,colback=cfon!6,colframe=cfon,arc=2.5mm,boxrule=0.9pt,left=3mm]{\Large\bfseries\color{cfon}b)}\quad\begin{minipage}[c]{5.3cm}\raggedright\Large $0$\end{minipage}\end{tcolorbox}\par\vspace{2.5mm}
\begin{tcolorbox}[enhanced,nobeforeafter,width=7.1cm,height=2.1cm,valign=center,colback=cfon!6,colframe=cfon,arc=2.5mm,boxrule=0.9pt,left=3mm]{\Large\bfseries\color{cfon}c)}\quad\begin{minipage}[c]{5.3cm}\raggedright\Large $6$\end{minipage}\end{tcolorbox}\hfill\begin{tcolorbox}[enhanced,nobeforeafter,width=7.1cm,height=2.1cm,valign=center,colback=cfon!6,colframe=cfon,arc=2.5mm,boxrule=0.9pt,left=3mm]{\Large\bfseries\color{cfon}d)}\quad\begin{minipage}[c]{5.3cm}\raggedright\Large $-6$\end{minipage}\end{tcolorbox}
\end{tcolorbox}
\newpage
\begin{tcolorbox}[enhanced,width=16cm,height=14cm,colback=white,colframe=cfon,boxrule=2pt,arc=0mm,
left=6mm,right=6mm,top=0mm,bottom=4mm,sharp corners,
overlay={\fill[cfon] (frame.north west) rectangle ([yshift=-1.25cm]frame.north east);
\node[anchor=west,text=white,font=\sffamily\bfseries\LARGE] at ([xshift=5mm,yshift=-0.62cm]frame.north west) {Question 42 C};
\node[anchor=east,text=white,font=\sffamily\itshape\large] at ([xshift=-5mm,yshift=-0.62cm]frame.north east) {Fonctions};}]
\vspace*{1.75cm}
\begin{minipage}[t][6.9cm][c]{\linewidth}\LARGE\centering
\begin{minipage}{0.95\linewidth}\raggedright\linespread{1.25}\selectfont $f$ est la fonction affine $f(x)=3x+b$ et $f(-1)=5$. Que vaut $b$~?\end{minipage}
\end{minipage}

\vfill
\begin{tcolorbox}[enhanced,nobeforeafter,width=7.1cm,height=2.1cm,valign=center,colback=cfon!6,colframe=cfon,arc=2.5mm,boxrule=0.9pt,left=3mm]{\Large\bfseries\color{cfon}a)}\quad\begin{minipage}[c]{5.3cm}\raggedright\Large $-5$\end{minipage}\end{tcolorbox}\hfill\begin{tcolorbox}[enhanced,nobeforeafter,width=7.1cm,height=2.1cm,valign=center,colback=cfon!6,colframe=cfon,arc=2.5mm,boxrule=0.9pt,left=3mm]{\Large\bfseries\color{cfon}b)}\quad\begin{minipage}[c]{5.3cm}\raggedright\Large $5$\end{minipage}\end{tcolorbox}\par\vspace{2.5mm}
\begin{tcolorbox}[enhanced,nobeforeafter,width=7.1cm,height=2.1cm,valign=center,colback=cfon!6,colframe=cfon,arc=2.5mm,boxrule=0.9pt,left=3mm]{\Large\bfseries\color{cfon}c)}\quad\begin{minipage}[c]{5.3cm}\raggedright\Large $2$\end{minipage}\end{tcolorbox}\hfill\begin{tcolorbox}[enhanced,nobeforeafter,width=7.1cm,height=2.1cm,valign=center,colback=cfon!6,colframe=cfon,arc=2.5mm,boxrule=0.9pt,left=3mm]{\Large\bfseries\color{cfon}d)}\quad\begin{minipage}[c]{5.3cm}\raggedright\Large $8$\end{minipage}\end{tcolorbox}
\end{tcolorbox}
\newpage
\begin{tcolorbox}[enhanced,width=16cm,height=14cm,colback=white,colframe=cfon,boxrule=2pt,arc=0mm,
left=6mm,right=6mm,top=0mm,bottom=4mm,sharp corners,
overlay={\fill[cfon] (frame.north west) rectangle ([yshift=-1.25cm]frame.north east);
\node[anchor=west,text=white,font=\sffamily\bfseries\LARGE] at ([xshift=5mm,yshift=-0.62cm]frame.north west) {Question 43 C};
\node[anchor=east,text=white,font=\sffamily\itshape\large] at ([xshift=-5mm,yshift=-0.62cm]frame.north east) {Fonctions};}]
\vspace*{1.75cm}
\begin{minipage}[t][6.9cm][c]{\linewidth}\LARGE\centering
\begin{minipage}{0.95\linewidth}\raggedright\linespread{1.25}\selectfont Augmenter une quantité de $150\,\%$ revient à la multiplier par~:\end{minipage}
\end{minipage}

\vfill
\begin{tcolorbox}[enhanced,nobeforeafter,width=7.1cm,height=2.1cm,valign=center,colback=cfon!6,colframe=cfon,arc=2.5mm,boxrule=0.9pt,left=3mm]{\Large\bfseries\color{cfon}a)}\quad\begin{minipage}[c]{5.3cm}\raggedright\Large $-0{,}5$\end{minipage}\end{tcolorbox}\hfill\begin{tcolorbox}[enhanced,nobeforeafter,width=7.1cm,height=2.1cm,valign=center,colback=cfon!6,colframe=cfon,arc=2.5mm,boxrule=0.9pt,left=3mm]{\Large\bfseries\color{cfon}b)}\quad\begin{minipage}[c]{5.3cm}\raggedright\Large $1{,}5$\end{minipage}\end{tcolorbox}\par\vspace{2.5mm}
\begin{tcolorbox}[enhanced,nobeforeafter,width=7.1cm,height=2.1cm,valign=center,colback=cfon!6,colframe=cfon,arc=2.5mm,boxrule=0.9pt,left=3mm]{\Large\bfseries\color{cfon}c)}\quad\begin{minipage}[c]{5.3cm}\raggedright\Large $2{,}5$\end{minipage}\end{tcolorbox}\hfill\begin{tcolorbox}[enhanced,nobeforeafter,width=7.1cm,height=2.1cm,valign=center,colback=cfon!6,colframe=cfon,arc=2.5mm,boxrule=0.9pt,left=3mm]{\Large\bfseries\color{cfon}d)}\quad\begin{minipage}[c]{5.3cm}\raggedright\Large $16$\end{minipage}\end{tcolorbox}
\end{tcolorbox}
\newpage
\begin{tcolorbox}[enhanced,width=16cm,height=14cm,colback=white,colframe=cfon,boxrule=2pt,arc=0mm,
left=6mm,right=6mm,top=0mm,bottom=4mm,sharp corners,
overlay={\fill[cfon] (frame.north west) rectangle ([yshift=-1.25cm]frame.north east);
\node[anchor=west,text=white,font=\sffamily\bfseries\LARGE] at ([xshift=5mm,yshift=-0.62cm]frame.north west) {Question 44 C};
\node[anchor=east,text=white,font=\sffamily\itshape\large] at ([xshift=-5mm,yshift=-0.62cm]frame.north east) {Fonctions};}]
\vspace*{1.75cm}
\begin{minipage}[t][6.9cm][c]{\linewidth}\LARGE\centering
\begin{minipage}{0.95\linewidth}\raggedright\linespread{1.25}\selectfont Diminuer une quantité de $5\,\%$ revient à la multiplier par~:\end{minipage}
\end{minipage}

\vfill
\begin{tcolorbox}[enhanced,nobeforeafter,width=7.1cm,height=2.1cm,valign=center,colback=cfon!6,colframe=cfon,arc=2.5mm,boxrule=0.9pt,left=3mm]{\Large\bfseries\color{cfon}a)}\quad\begin{minipage}[c]{5.3cm}\raggedright\Large $1{,}05$\end{minipage}\end{tcolorbox}\hfill\begin{tcolorbox}[enhanced,nobeforeafter,width=7.1cm,height=2.1cm,valign=center,colback=cfon!6,colframe=cfon,arc=2.5mm,boxrule=0.9pt,left=3mm]{\Large\bfseries\color{cfon}b)}\quad\begin{minipage}[c]{5.3cm}\raggedright\Large $0{,}95$\end{minipage}\end{tcolorbox}\par\vspace{2.5mm}
\begin{tcolorbox}[enhanced,nobeforeafter,width=7.1cm,height=2.1cm,valign=center,colback=cfon!6,colframe=cfon,arc=2.5mm,boxrule=0.9pt,left=3mm]{\Large\bfseries\color{cfon}c)}\quad\begin{minipage}[c]{5.3cm}\raggedright\Large $0{,}5$\end{minipage}\end{tcolorbox}\hfill\begin{tcolorbox}[enhanced,nobeforeafter,width=7.1cm,height=2.1cm,valign=center,colback=cfon!6,colframe=cfon,arc=2.5mm,boxrule=0.9pt,left=3mm]{\Large\bfseries\color{cfon}d)}\quad\begin{minipage}[c]{5.3cm}\raggedright\Large $0{,}05$\end{minipage}\end{tcolorbox}
\end{tcolorbox}
\newpage
\begin{tcolorbox}[enhanced,width=16cm,height=14cm,colback=white,colframe=cfon,boxrule=2pt,arc=0mm,
left=6mm,right=6mm,top=0mm,bottom=4mm,sharp corners,
overlay={\fill[cfon] (frame.north west) rectangle ([yshift=-1.25cm]frame.north east);
\node[anchor=west,text=white,font=\sffamily\bfseries\LARGE] at ([xshift=5mm,yshift=-0.62cm]frame.north west) {Question 45 C};
\node[anchor=east,text=white,font=\sffamily\itshape\large] at ([xshift=-5mm,yshift=-0.62cm]frame.north east) {Fonctions};}]
\vspace*{1.75cm}
\begin{minipage}[t][6.9cm][c]{\linewidth}\LARGE\centering
\begin{minipage}{0.95\linewidth}\raggedright\linespread{1.25}\selectfont Soit $g(x)=3x^2-4$. Que vaut $g(2)$~?\end{minipage}
\end{minipage}

\vfill
\begin{tcolorbox}[enhanced,nobeforeafter,width=7.1cm,height=2.1cm,valign=center,colback=cfon!6,colframe=cfon,arc=2.5mm,boxrule=0.9pt,left=3mm]{\Large\bfseries\color{cfon}a)}\quad\begin{minipage}[c]{5.3cm}\raggedright\Large $-16$\end{minipage}\end{tcolorbox}\hfill\begin{tcolorbox}[enhanced,nobeforeafter,width=7.1cm,height=2.1cm,valign=center,colback=cfon!6,colframe=cfon,arc=2.5mm,boxrule=0.9pt,left=3mm]{\Large\bfseries\color{cfon}b)}\quad\begin{minipage}[c]{5.3cm}\raggedright\Large $8$\end{minipage}\end{tcolorbox}\par\vspace{2.5mm}
\begin{tcolorbox}[enhanced,nobeforeafter,width=7.1cm,height=2.1cm,valign=center,colback=cfon!6,colframe=cfon,arc=2.5mm,boxrule=0.9pt,left=3mm]{\Large\bfseries\color{cfon}c)}\quad\begin{minipage}[c]{5.3cm}\raggedright\Large $10$\end{minipage}\end{tcolorbox}\hfill\begin{tcolorbox}[enhanced,nobeforeafter,width=7.1cm,height=2.1cm,valign=center,colback=cfon!6,colframe=cfon,arc=2.5mm,boxrule=0.9pt,left=3mm]{\Large\bfseries\color{cfon}d)}\quad\begin{minipage}[c]{5.3cm}\raggedright\Large $-8$\end{minipage}\end{tcolorbox}
\end{tcolorbox}
\newpage
\begin{tcolorbox}[enhanced,width=16cm,height=14cm,colback=white,colframe=cfon,boxrule=2pt,arc=0mm,
left=6mm,right=6mm,top=0mm,bottom=4mm,sharp corners,
overlay={\fill[cfon] (frame.north west) rectangle ([yshift=-1.25cm]frame.north east);
\node[anchor=west,text=white,font=\sffamily\bfseries\LARGE] at ([xshift=5mm,yshift=-0.62cm]frame.north west) {Question 46 C};
\node[anchor=east,text=white,font=\sffamily\itshape\large] at ([xshift=-5mm,yshift=-0.62cm]frame.north east) {Fonctions};}]
\vspace*{1.75cm}
\begin{minipage}[t][6.9cm][c]{\linewidth}\LARGE\centering
\begin{minipage}{0.95\linewidth}\raggedright\linespread{1.25}\selectfont Soit $g(x)=x^2-2x-8$. Que vaut $g(-3)$~?\end{minipage}
\end{minipage}

\vfill
\begin{tcolorbox}[enhanced,nobeforeafter,width=7.1cm,height=2.1cm,valign=center,colback=cfon!6,colframe=cfon,arc=2.5mm,boxrule=0.9pt,left=3mm]{\Large\bfseries\color{cfon}a)}\quad\begin{minipage}[c]{5.3cm}\raggedright\Large $7$\end{minipage}\end{tcolorbox}\hfill\begin{tcolorbox}[enhanced,nobeforeafter,width=7.1cm,height=2.1cm,valign=center,colback=cfon!6,colframe=cfon,arc=2.5mm,boxrule=0.9pt,left=3mm]{\Large\bfseries\color{cfon}b)}\quad\begin{minipage}[c]{5.3cm}\raggedright\Large $-8$\end{minipage}\end{tcolorbox}\par\vspace{2.5mm}
\begin{tcolorbox}[enhanced,nobeforeafter,width=7.1cm,height=2.1cm,valign=center,colback=cfon!6,colframe=cfon,arc=2.5mm,boxrule=0.9pt,left=3mm]{\Large\bfseries\color{cfon}c)}\quad\begin{minipage}[c]{5.3cm}\raggedright\Large $-11$\end{minipage}\end{tcolorbox}\hfill\begin{tcolorbox}[enhanced,nobeforeafter,width=7.1cm,height=2.1cm,valign=center,colback=cfon!6,colframe=cfon,arc=2.5mm,boxrule=0.9pt,left=3mm]{\Large\bfseries\color{cfon}d)}\quad\begin{minipage}[c]{5.3cm}\raggedright\Large $-7$\end{minipage}\end{tcolorbox}
\end{tcolorbox}
\newpage
\begin{tcolorbox}[enhanced,width=16cm,height=14cm,colback=white,colframe=cfon,boxrule=2pt,arc=0mm,
left=6mm,right=6mm,top=0mm,bottom=4mm,sharp corners,
overlay={\fill[cfon] (frame.north west) rectangle ([yshift=-1.25cm]frame.north east);
\node[anchor=west,text=white,font=\sffamily\bfseries\LARGE] at ([xshift=5mm,yshift=-0.62cm]frame.north west) {Question 47 C};
\node[anchor=east,text=white,font=\sffamily\itshape\large] at ([xshift=-5mm,yshift=-0.62cm]frame.north east) {Fonctions};}]
\vspace*{1.75cm}
\begin{minipage}[t][6.9cm][c]{\linewidth}\LARGE\centering
\begin{minipage}{0.95\linewidth}\raggedright\linespread{1.25}\selectfont De quel type est la fonction $f:x\mapsto -x$~?\end{minipage}
\end{minipage}

\vfill
\begin{tcolorbox}[enhanced,nobeforeafter,width=7.1cm,height=2.1cm,valign=center,colback=cfon!6,colframe=cfon,arc=2.5mm,boxrule=0.9pt,left=3mm]{\Large\bfseries\color{cfon}a)}\quad\begin{minipage}[c]{5.3cm}\raggedright\Large constante\end{minipage}\end{tcolorbox}\hfill\begin{tcolorbox}[enhanced,nobeforeafter,width=7.1cm,height=2.1cm,valign=center,colback=cfon!6,colframe=cfon,arc=2.5mm,boxrule=0.9pt,left=3mm]{\Large\bfseries\color{cfon}b)}\quad\begin{minipage}[c]{5.3cm}\raggedright\Large affine non linéaire\end{minipage}\end{tcolorbox}\par\vspace{2.5mm}
\begin{tcolorbox}[enhanced,nobeforeafter,width=7.1cm,height=2.1cm,valign=center,colback=cfon!6,colframe=cfon,arc=2.5mm,boxrule=0.9pt,left=3mm]{\Large\bfseries\color{cfon}c)}\quad\begin{minipage}[c]{5.3cm}\raggedright\Large ni affine ni linéaire\end{minipage}\end{tcolorbox}\hfill\begin{tcolorbox}[enhanced,nobeforeafter,width=7.1cm,height=2.1cm,valign=center,colback=cfon!6,colframe=cfon,arc=2.5mm,boxrule=0.9pt,left=3mm]{\Large\bfseries\color{cfon}d)}\quad\begin{minipage}[c]{5.3cm}\raggedright\Large linéaire\end{minipage}\end{tcolorbox}
\end{tcolorbox}
\newpage
\begin{tcolorbox}[enhanced,width=16cm,height=14cm,colback=white,colframe=cfon,boxrule=2pt,arc=0mm,
left=6mm,right=6mm,top=0mm,bottom=4mm,sharp corners,
overlay={\fill[cfon] (frame.north west) rectangle ([yshift=-1.25cm]frame.north east);
\node[anchor=west,text=white,font=\sffamily\bfseries\LARGE] at ([xshift=5mm,yshift=-0.62cm]frame.north west) {Question 48 C};
\node[anchor=east,text=white,font=\sffamily\itshape\large] at ([xshift=-5mm,yshift=-0.62cm]frame.north east) {Fonctions};}]
\vspace*{1.75cm}
\begin{minipage}[t][6.9cm][c]{\linewidth}\LARGE\centering
\begin{minipage}{0.95\linewidth}\raggedright\linespread{1.25}\selectfont De quel type est la fonction $f:x\mapsto 3-x$~?\end{minipage}
\end{minipage}

\vfill
\begin{tcolorbox}[enhanced,nobeforeafter,width=7.1cm,height=2.1cm,valign=center,colback=cfon!6,colframe=cfon,arc=2.5mm,boxrule=0.9pt,left=3mm]{\Large\bfseries\color{cfon}a)}\quad\begin{minipage}[c]{5.3cm}\raggedright\Large ni affine ni linéaire\end{minipage}\end{tcolorbox}\hfill\begin{tcolorbox}[enhanced,nobeforeafter,width=7.1cm,height=2.1cm,valign=center,colback=cfon!6,colframe=cfon,arc=2.5mm,boxrule=0.9pt,left=3mm]{\Large\bfseries\color{cfon}b)}\quad\begin{minipage}[c]{5.3cm}\raggedright\Large affine non linéaire\end{minipage}\end{tcolorbox}\par\vspace{2.5mm}
\begin{tcolorbox}[enhanced,nobeforeafter,width=7.1cm,height=2.1cm,valign=center,colback=cfon!6,colframe=cfon,arc=2.5mm,boxrule=0.9pt,left=3mm]{\Large\bfseries\color{cfon}c)}\quad\begin{minipage}[c]{5.3cm}\raggedright\Large constante\end{minipage}\end{tcolorbox}\hfill\begin{tcolorbox}[enhanced,nobeforeafter,width=7.1cm,height=2.1cm,valign=center,colback=cfon!6,colframe=cfon,arc=2.5mm,boxrule=0.9pt,left=3mm]{\Large\bfseries\color{cfon}d)}\quad\begin{minipage}[c]{5.3cm}\raggedright\Large linéaire\end{minipage}\end{tcolorbox}
\end{tcolorbox}
\newpage
\begin{tcolorbox}[enhanced,width=16cm,height=14cm,colback=white,colframe=cgeo,boxrule=2pt,arc=0mm,
left=6mm,right=6mm,top=0mm,bottom=4mm,sharp corners,
overlay={\fill[cgeo] (frame.north west) rectangle ([yshift=-1.25cm]frame.north east);
\node[anchor=west,text=white,font=\sffamily\bfseries\LARGE] at ([xshift=5mm,yshift=-0.62cm]frame.north west) {Question 49 C};
\node[anchor=east,text=white,font=\sffamily\itshape\large] at ([xshift=-5mm,yshift=-0.62cm]frame.north east) {Géométrie};}]
\vspace*{1.75cm}
\begin{minipage}[t][6.9cm][c]{\linewidth}\Large\centering
\begin{minipage}{0.95\linewidth}\raggedright\linespread{1.25}\selectfont Les droites $(AC)$ et $(BD)$ se coupent en $O$, et les droites $(AB)$ et $(CD)$ sont parallèles. $OA=6$ cm, $OC=9$ cm et $AB=4$ cm. Combien mesure $CD$~?\end{minipage}
\end{minipage}

\vfill
\begin{tcolorbox}[enhanced,nobeforeafter,width=7.1cm,height=2.1cm,valign=center,colback=cgeo!6,colframe=cgeo,arc=2.5mm,boxrule=0.9pt,left=3mm]{\Large\bfseries\color{cgeo}a)}\quad\begin{minipage}[c]{5.3cm}\raggedright\Large $7~\text{cm}$\end{minipage}\end{tcolorbox}\hfill\begin{tcolorbox}[enhanced,nobeforeafter,width=7.1cm,height=2.1cm,valign=center,colback=cgeo!6,colframe=cgeo,arc=2.5mm,boxrule=0.9pt,left=3mm]{\Large\bfseries\color{cgeo}b)}\quad\begin{minipage}[c]{5.3cm}\raggedright\Large $13{,}5~\text{cm}$\end{minipage}\end{tcolorbox}\par\vspace{2.5mm}
\begin{tcolorbox}[enhanced,nobeforeafter,width=7.1cm,height=2.1cm,valign=center,colback=cgeo!6,colframe=cgeo,arc=2.5mm,boxrule=0.9pt,left=3mm]{\Large\bfseries\color{cgeo}c)}\quad\begin{minipage}[c]{5.3cm}\raggedright\Large $6~\text{cm}$\end{minipage}\end{tcolorbox}\hfill\begin{tcolorbox}[enhanced,nobeforeafter,width=7.1cm,height=2.1cm,valign=center,colback=cgeo!6,colframe=cgeo,arc=2.5mm,boxrule=0.9pt,left=3mm]{\Large\bfseries\color{cgeo}d)}\quad\begin{minipage}[c]{5.3cm}\raggedright\Large $\dfrac{8}{3}~\text{cm}$\end{minipage}\end{tcolorbox}
\end{tcolorbox}
\newpage
\begin{tcolorbox}[enhanced,width=16cm,height=14cm,colback=white,colframe=cgeo,boxrule=2pt,arc=0mm,
left=6mm,right=6mm,top=0mm,bottom=4mm,sharp corners,
overlay={\fill[cgeo] (frame.north west) rectangle ([yshift=-1.25cm]frame.north east);
\node[anchor=west,text=white,font=\sffamily\bfseries\LARGE] at ([xshift=5mm,yshift=-0.62cm]frame.north west) {Question 50 C};
\node[anchor=east,text=white,font=\sffamily\itshape\large] at ([xshift=-5mm,yshift=-0.62cm]frame.north east) {Géométrie};}]
\vspace*{1.75cm}
\begin{minipage}[t][6.9cm][c]{\linewidth}\Large\centering
\begin{minipage}{0.95\linewidth}\raggedright\linespread{1.25}\selectfont Les droites $(AC)$ et $(BD)$ se coupent en $O$, et les droites $(AB)$ et $(CD)$ sont parallèles. $OA=2$ cm, $OC=7$ cm et $AB=5$ cm. Combien mesure $CD$~?\end{minipage}
\end{minipage}

\vfill
\begin{tcolorbox}[enhanced,nobeforeafter,width=7.1cm,height=2.1cm,valign=center,colback=cgeo!6,colframe=cgeo,arc=2.5mm,boxrule=0.9pt,left=3mm]{\Large\bfseries\color{cgeo}a)}\quad\begin{minipage}[c]{5.3cm}\raggedright\Large $2{,}8~\text{cm}$\end{minipage}\end{tcolorbox}\hfill\begin{tcolorbox}[enhanced,nobeforeafter,width=7.1cm,height=2.1cm,valign=center,colback=cgeo!6,colframe=cgeo,arc=2.5mm,boxrule=0.9pt,left=3mm]{\Large\bfseries\color{cgeo}b)}\quad\begin{minipage}[c]{5.3cm}\raggedright\Large $17{,}5~\text{cm}$\end{minipage}\end{tcolorbox}\par\vspace{2.5mm}
\begin{tcolorbox}[enhanced,nobeforeafter,width=7.1cm,height=2.1cm,valign=center,colback=cgeo!6,colframe=cgeo,arc=2.5mm,boxrule=0.9pt,left=3mm]{\Large\bfseries\color{cgeo}c)}\quad\begin{minipage}[c]{5.3cm}\raggedright\Large $\dfrac{10}{7}~\text{cm}$\end{minipage}\end{tcolorbox}\hfill\begin{tcolorbox}[enhanced,nobeforeafter,width=7.1cm,height=2.1cm,valign=center,colback=cgeo!6,colframe=cgeo,arc=2.5mm,boxrule=0.9pt,left=3mm]{\Large\bfseries\color{cgeo}d)}\quad\begin{minipage}[c]{5.3cm}\raggedright\Large $10~\text{cm}$\end{minipage}\end{tcolorbox}
\end{tcolorbox}
\newpage
\begin{tcolorbox}[enhanced,width=16cm,height=14cm,colback=white,colframe=cgeo,boxrule=2pt,arc=0mm,
left=6mm,right=6mm,top=0mm,bottom=4mm,sharp corners,
overlay={\fill[cgeo] (frame.north west) rectangle ([yshift=-1.25cm]frame.north east);
\node[anchor=west,text=white,font=\sffamily\bfseries\LARGE] at ([xshift=5mm,yshift=-0.62cm]frame.north west) {Question 51 C};
\node[anchor=east,text=white,font=\sffamily\itshape\large] at ([xshift=-5mm,yshift=-0.62cm]frame.north east) {Géométrie};}]
\vspace*{1.75cm}
\begin{minipage}[t][6.9cm][c]{\linewidth}\Large\centering
\begin{minipage}{0.95\linewidth}\raggedright\linespread{1.25}\selectfont Les points $A$, $M$, $B$ sont alignés dans cet ordre, ainsi que $A$, $N$, $C$. $AM=7$, $AB=21$, $AN=2$ et $AC=6$. Les droites $(MN)$ et $(BC)$ sont-elles parallèles~?\end{minipage}
\end{minipage}

\vfill
\begin{tcolorbox}[enhanced,nobeforeafter,width=7.1cm,height=2.1cm,valign=center,colback=cgeo!6,colframe=cgeo,arc=2.5mm,boxrule=0.9pt,left=3mm]{\Large\bfseries\color{cgeo}a)}\quad\begin{minipage}[c]{5.3cm}\raggedright\Large non\end{minipage}\end{tcolorbox}\hfill\begin{tcolorbox}[enhanced,nobeforeafter,width=7.1cm,height=2.1cm,valign=center,colback=cgeo!6,colframe=cgeo,arc=2.5mm,boxrule=0.9pt,left=3mm]{\Large\bfseries\color{cgeo}b)}\quad\begin{minipage}[c]{5.3cm}\raggedright\Large oui\end{minipage}\end{tcolorbox}\par\vspace{2.5mm}
\begin{tcolorbox}[enhanced,nobeforeafter,width=7.1cm,height=2.1cm,valign=center,colback=cgeo!6,colframe=cgeo,arc=2.5mm,boxrule=0.9pt,left=3mm]{\Large\bfseries\color{cgeo}c)}\quad\begin{minipage}[c]{5.3cm}\raggedright\large seulement si le triangle est rectangle\end{minipage}\end{tcolorbox}\hfill\begin{tcolorbox}[enhanced,nobeforeafter,width=7.1cm,height=2.1cm,valign=center,colback=cgeo!6,colframe=cgeo,arc=2.5mm,boxrule=0.9pt,left=3mm]{\Large\bfseries\color{cgeo}d)}\quad\begin{minipage}[c]{5.3cm}\raggedright\Large on ne peut pas savoir\end{minipage}\end{tcolorbox}
\end{tcolorbox}
\newpage
\begin{tcolorbox}[enhanced,width=16cm,height=14cm,colback=white,colframe=cgeo,boxrule=2pt,arc=0mm,
left=6mm,right=6mm,top=0mm,bottom=4mm,sharp corners,
overlay={\fill[cgeo] (frame.north west) rectangle ([yshift=-1.25cm]frame.north east);
\node[anchor=west,text=white,font=\sffamily\bfseries\LARGE] at ([xshift=5mm,yshift=-0.62cm]frame.north west) {Question 52 C};
\node[anchor=east,text=white,font=\sffamily\itshape\large] at ([xshift=-5mm,yshift=-0.62cm]frame.north east) {Géométrie};}]
\vspace*{1.75cm}
\begin{minipage}[t][6.9cm][c]{\linewidth}\Large\centering
\begin{minipage}{0.95\linewidth}\raggedright\linespread{1.25}\selectfont Les points $A$, $M$, $B$ sont alignés dans cet ordre, ainsi que $A$, $N$, $C$. $AM=3$, $AB=7$, $AN=6$ et $AC=15$. Les droites $(MN)$ et $(BC)$ sont-elles parallèles~?\end{minipage}
\end{minipage}

\vfill
\begin{tcolorbox}[enhanced,nobeforeafter,width=7.1cm,height=2.1cm,valign=center,colback=cgeo!6,colframe=cgeo,arc=2.5mm,boxrule=0.9pt,left=3mm]{\Large\bfseries\color{cgeo}a)}\quad\begin{minipage}[c]{5.3cm}\raggedright\Large non\end{minipage}\end{tcolorbox}\hfill\begin{tcolorbox}[enhanced,nobeforeafter,width=7.1cm,height=2.1cm,valign=center,colback=cgeo!6,colframe=cgeo,arc=2.5mm,boxrule=0.9pt,left=3mm]{\Large\bfseries\color{cgeo}b)}\quad\begin{minipage}[c]{5.3cm}\raggedright\large seulement si le triangle est rectangle\end{minipage}\end{tcolorbox}\par\vspace{2.5mm}
\begin{tcolorbox}[enhanced,nobeforeafter,width=7.1cm,height=2.1cm,valign=center,colback=cgeo!6,colframe=cgeo,arc=2.5mm,boxrule=0.9pt,left=3mm]{\Large\bfseries\color{cgeo}c)}\quad\begin{minipage}[c]{5.3cm}\raggedright\Large on ne peut pas savoir\end{minipage}\end{tcolorbox}\hfill\begin{tcolorbox}[enhanced,nobeforeafter,width=7.1cm,height=2.1cm,valign=center,colback=cgeo!6,colframe=cgeo,arc=2.5mm,boxrule=0.9pt,left=3mm]{\Large\bfseries\color{cgeo}d)}\quad\begin{minipage}[c]{5.3cm}\raggedright\Large oui\end{minipage}\end{tcolorbox}
\end{tcolorbox}
\newpage
\begin{tcolorbox}[enhanced,width=16cm,height=14cm,colback=white,colframe=cgeo,boxrule=2pt,arc=0mm,
left=6mm,right=6mm,top=0mm,bottom=4mm,sharp corners,
overlay={\fill[cgeo] (frame.north west) rectangle ([yshift=-1.25cm]frame.north east);
\node[anchor=west,text=white,font=\sffamily\bfseries\LARGE] at ([xshift=5mm,yshift=-0.62cm]frame.north west) {Question 53 C};
\node[anchor=east,text=white,font=\sffamily\itshape\large] at ([xshift=-5mm,yshift=-0.62cm]frame.north east) {Géométrie};}]
\vspace*{1.75cm}
\begin{minipage}[t][6.9cm][c]{\linewidth}\LARGE\centering
\begin{minipage}{0.95\linewidth}\raggedright\linespread{1.25}\selectfont Le triangle $ABC$ est rectangle en $A$, avec $BC=15$ cm et $\sin\widehat{ABC}=0{,}2$. Combien mesure $AC$~?\end{minipage}
\end{minipage}

\vfill
\begin{tcolorbox}[enhanced,nobeforeafter,width=7.1cm,height=2.1cm,valign=center,colback=cgeo!6,colframe=cgeo,arc=2.5mm,boxrule=0.9pt,left=3mm]{\Large\bfseries\color{cgeo}a)}\quad\begin{minipage}[c]{5.3cm}\raggedright\Large $6~\text{cm}$\end{minipage}\end{tcolorbox}\hfill\begin{tcolorbox}[enhanced,nobeforeafter,width=7.1cm,height=2.1cm,valign=center,colback=cgeo!6,colframe=cgeo,arc=2.5mm,boxrule=0.9pt,left=3mm]{\Large\bfseries\color{cgeo}b)}\quad\begin{minipage}[c]{5.3cm}\raggedright\Large $12~\text{cm}$\end{minipage}\end{tcolorbox}\par\vspace{2.5mm}
\begin{tcolorbox}[enhanced,nobeforeafter,width=7.1cm,height=2.1cm,valign=center,colback=cgeo!6,colframe=cgeo,arc=2.5mm,boxrule=0.9pt,left=3mm]{\Large\bfseries\color{cgeo}c)}\quad\begin{minipage}[c]{5.3cm}\raggedright\Large $3~\text{cm}$\end{minipage}\end{tcolorbox}\hfill\begin{tcolorbox}[enhanced,nobeforeafter,width=7.1cm,height=2.1cm,valign=center,colback=cgeo!6,colframe=cgeo,arc=2.5mm,boxrule=0.9pt,left=3mm]{\Large\bfseries\color{cgeo}d)}\quad\begin{minipage}[c]{5.3cm}\raggedright\Large $1{,}5~\text{cm}$\end{minipage}\end{tcolorbox}
\end{tcolorbox}
\newpage
\begin{tcolorbox}[enhanced,width=16cm,height=14cm,colback=white,colframe=cgeo,boxrule=2pt,arc=0mm,
left=6mm,right=6mm,top=0mm,bottom=4mm,sharp corners,
overlay={\fill[cgeo] (frame.north west) rectangle ([yshift=-1.25cm]frame.north east);
\node[anchor=west,text=white,font=\sffamily\bfseries\LARGE] at ([xshift=5mm,yshift=-0.62cm]frame.north west) {Question 54 C};
\node[anchor=east,text=white,font=\sffamily\itshape\large] at ([xshift=-5mm,yshift=-0.62cm]frame.north east) {Géométrie};}]
\vspace*{1.75cm}
\begin{minipage}[t][6.9cm][c]{\linewidth}\LARGE\centering
\begin{minipage}{0.95\linewidth}\raggedright\linespread{1.25}\selectfont Le triangle $ABC$ est rectangle en $A$, avec $BC=5$ cm et $\sin\widehat{ABC}=0{,}8$. Combien mesure $AC$~?\end{minipage}
\end{minipage}

\vfill
\begin{tcolorbox}[enhanced,nobeforeafter,width=7.1cm,height=2.1cm,valign=center,colback=cgeo!6,colframe=cgeo,arc=2.5mm,boxrule=0.9pt,left=3mm]{\Large\bfseries\color{cgeo}a)}\quad\begin{minipage}[c]{5.3cm}\raggedright\Large $2~\text{cm}$\end{minipage}\end{tcolorbox}\hfill\begin{tcolorbox}[enhanced,nobeforeafter,width=7.1cm,height=2.1cm,valign=center,colback=cgeo!6,colframe=cgeo,arc=2.5mm,boxrule=0.9pt,left=3mm]{\Large\bfseries\color{cgeo}b)}\quad\begin{minipage}[c]{5.3cm}\raggedright\Large $8~\text{cm}$\end{minipage}\end{tcolorbox}\par\vspace{2.5mm}
\begin{tcolorbox}[enhanced,nobeforeafter,width=7.1cm,height=2.1cm,valign=center,colback=cgeo!6,colframe=cgeo,arc=2.5mm,boxrule=0.9pt,left=3mm]{\Large\bfseries\color{cgeo}c)}\quad\begin{minipage}[c]{5.3cm}\raggedright\Large $1~\text{cm}$\end{minipage}\end{tcolorbox}\hfill\begin{tcolorbox}[enhanced,nobeforeafter,width=7.1cm,height=2.1cm,valign=center,colback=cgeo!6,colframe=cgeo,arc=2.5mm,boxrule=0.9pt,left=3mm]{\Large\bfseries\color{cgeo}d)}\quad\begin{minipage}[c]{5.3cm}\raggedright\Large $4~\text{cm}$\end{minipage}\end{tcolorbox}
\end{tcolorbox}
\newpage
\begin{tcolorbox}[enhanced,width=16cm,height=14cm,colback=white,colframe=cgeo,boxrule=2pt,arc=0mm,
left=6mm,right=6mm,top=0mm,bottom=4mm,sharp corners,
overlay={\fill[cgeo] (frame.north west) rectangle ([yshift=-1.25cm]frame.north east);
\node[anchor=west,text=white,font=\sffamily\bfseries\LARGE] at ([xshift=5mm,yshift=-0.62cm]frame.north west) {Question 55 C};
\node[anchor=east,text=white,font=\sffamily\itshape\large] at ([xshift=-5mm,yshift=-0.62cm]frame.north east) {Géométrie};}]
\vspace*{1.75cm}
\begin{minipage}[t][6.9cm][c]{\linewidth}\LARGE\centering
\begin{minipage}{0.95\linewidth}\raggedright\linespread{1.25}\selectfont Le triangle $ABC$ est rectangle en $A$, avec $AB=10$ cm et $AC=4$ cm. Que vaut $\tan\widehat{ABC}$~?\end{minipage}
\end{minipage}

\vfill
\begin{tcolorbox}[enhanced,nobeforeafter,width=7.1cm,height=2.1cm,valign=center,colback=cgeo!6,colframe=cgeo,arc=2.5mm,boxrule=0.9pt,left=3mm]{\Large\bfseries\color{cgeo}a)}\quad\begin{minipage}[c]{5.3cm}\raggedright\Large $\dfrac{7}{5}$\end{minipage}\end{tcolorbox}\hfill\begin{tcolorbox}[enhanced,nobeforeafter,width=7.1cm,height=2.1cm,valign=center,colback=cgeo!6,colframe=cgeo,arc=2.5mm,boxrule=0.9pt,left=3mm]{\Large\bfseries\color{cgeo}b)}\quad\begin{minipage}[c]{5.3cm}\raggedright\Large $\dfrac{5}{2}$\end{minipage}\end{tcolorbox}\par\vspace{2.5mm}
\begin{tcolorbox}[enhanced,nobeforeafter,width=7.1cm,height=2.1cm,valign=center,colback=cgeo!6,colframe=cgeo,arc=2.5mm,boxrule=0.9pt,left=3mm]{\Large\bfseries\color{cgeo}c)}\quad\begin{minipage}[c]{5.3cm}\raggedright\Large $\dfrac{2}{5}$\end{minipage}\end{tcolorbox}\hfill\begin{tcolorbox}[enhanced,nobeforeafter,width=7.1cm,height=2.1cm,valign=center,colback=cgeo!6,colframe=cgeo,arc=2.5mm,boxrule=0.9pt,left=3mm]{\Large\bfseries\color{cgeo}d)}\quad\begin{minipage}[c]{5.3cm}\raggedright\Large $\dfrac{5}{7}$\end{minipage}\end{tcolorbox}
\end{tcolorbox}
\newpage
\begin{tcolorbox}[enhanced,width=16cm,height=14cm,colback=white,colframe=cgeo,boxrule=2pt,arc=0mm,
left=6mm,right=6mm,top=0mm,bottom=4mm,sharp corners,
overlay={\fill[cgeo] (frame.north west) rectangle ([yshift=-1.25cm]frame.north east);
\node[anchor=west,text=white,font=\sffamily\bfseries\LARGE] at ([xshift=5mm,yshift=-0.62cm]frame.north west) {Question 56 C};
\node[anchor=east,text=white,font=\sffamily\itshape\large] at ([xshift=-5mm,yshift=-0.62cm]frame.north east) {Géométrie};}]
\vspace*{1.75cm}
\begin{minipage}[t][6.9cm][c]{\linewidth}\LARGE\centering
\begin{minipage}{0.95\linewidth}\raggedright\linespread{1.25}\selectfont Le triangle $ABC$ est rectangle en $A$, avec $AB=3$ cm et $AC=9$ cm. Que vaut $\tan\widehat{ABC}$~?\end{minipage}
\end{minipage}

\vfill
\begin{tcolorbox}[enhanced,nobeforeafter,width=7.1cm,height=2.1cm,valign=center,colback=cgeo!6,colframe=cgeo,arc=2.5mm,boxrule=0.9pt,left=3mm]{\Large\bfseries\color{cgeo}a)}\quad\begin{minipage}[c]{5.3cm}\raggedright\Large $\dfrac{1}{4}$\end{minipage}\end{tcolorbox}\hfill\begin{tcolorbox}[enhanced,nobeforeafter,width=7.1cm,height=2.1cm,valign=center,colback=cgeo!6,colframe=cgeo,arc=2.5mm,boxrule=0.9pt,left=3mm]{\Large\bfseries\color{cgeo}b)}\quad\begin{minipage}[c]{5.3cm}\raggedright\Large $4$\end{minipage}\end{tcolorbox}\par\vspace{2.5mm}
\begin{tcolorbox}[enhanced,nobeforeafter,width=7.1cm,height=2.1cm,valign=center,colback=cgeo!6,colframe=cgeo,arc=2.5mm,boxrule=0.9pt,left=3mm]{\Large\bfseries\color{cgeo}c)}\quad\begin{minipage}[c]{5.3cm}\raggedright\Large $3$\end{minipage}\end{tcolorbox}\hfill\begin{tcolorbox}[enhanced,nobeforeafter,width=7.1cm,height=2.1cm,valign=center,colback=cgeo!6,colframe=cgeo,arc=2.5mm,boxrule=0.9pt,left=3mm]{\Large\bfseries\color{cgeo}d)}\quad\begin{minipage}[c]{5.3cm}\raggedright\Large $\dfrac{1}{3}$\end{minipage}\end{tcolorbox}
\end{tcolorbox}
\newpage
\begin{tcolorbox}[enhanced,width=16cm,height=14cm,colback=white,colframe=cgeo,boxrule=2pt,arc=0mm,
left=6mm,right=6mm,top=0mm,bottom=4mm,sharp corners,
overlay={\fill[cgeo] (frame.north west) rectangle ([yshift=-1.25cm]frame.north east);
\node[anchor=west,text=white,font=\sffamily\bfseries\LARGE] at ([xshift=5mm,yshift=-0.62cm]frame.north west) {Question 57 C};
\node[anchor=east,text=white,font=\sffamily\itshape\large] at ([xshift=-5mm,yshift=-0.62cm]frame.north east) {Géométrie};}]
\vspace*{1.75cm}
\begin{minipage}[t][6.9cm][c]{\linewidth}\LARGE\centering
\begin{minipage}{0.95\linewidth}\raggedright\linespread{1.25}\selectfont Une homothétie de rapport $-\dfrac{1}{2}$ transforme le segment $[AB]$, de longueur $6$ cm, en $[A'B']$. Combien mesure $A'B'$~?\end{minipage}
\end{minipage}

\vfill
\begin{tcolorbox}[enhanced,nobeforeafter,width=7.1cm,height=2.1cm,valign=center,colback=cgeo!6,colframe=cgeo,arc=2.5mm,boxrule=0.9pt,left=3mm]{\Large\bfseries\color{cgeo}a)}\quad\begin{minipage}[c]{5.3cm}\raggedright\Large $4~\text{cm}$\end{minipage}\end{tcolorbox}\hfill\begin{tcolorbox}[enhanced,nobeforeafter,width=7.1cm,height=2.1cm,valign=center,colback=cgeo!6,colframe=cgeo,arc=2.5mm,boxrule=0.9pt,left=3mm]{\Large\bfseries\color{cgeo}b)}\quad\begin{minipage}[c]{5.3cm}\raggedright\Large $6{,}5~\text{cm}$\end{minipage}\end{tcolorbox}\par\vspace{2.5mm}
\begin{tcolorbox}[enhanced,nobeforeafter,width=7.1cm,height=2.1cm,valign=center,colback=cgeo!6,colframe=cgeo,arc=2.5mm,boxrule=0.9pt,left=3mm]{\Large\bfseries\color{cgeo}c)}\quad\begin{minipage}[c]{5.3cm}\raggedright\Large $12~\text{cm}$\end{minipage}\end{tcolorbox}\hfill\begin{tcolorbox}[enhanced,nobeforeafter,width=7.1cm,height=2.1cm,valign=center,colback=cgeo!6,colframe=cgeo,arc=2.5mm,boxrule=0.9pt,left=3mm]{\Large\bfseries\color{cgeo}d)}\quad\begin{minipage}[c]{5.3cm}\raggedright\Large $3~\text{cm}$\end{minipage}\end{tcolorbox}
\end{tcolorbox}
\newpage
\begin{tcolorbox}[enhanced,width=16cm,height=14cm,colback=white,colframe=cgeo,boxrule=2pt,arc=0mm,
left=6mm,right=6mm,top=0mm,bottom=4mm,sharp corners,
overlay={\fill[cgeo] (frame.north west) rectangle ([yshift=-1.25cm]frame.north east);
\node[anchor=west,text=white,font=\sffamily\bfseries\LARGE] at ([xshift=5mm,yshift=-0.62cm]frame.north west) {Question 58 C};
\node[anchor=east,text=white,font=\sffamily\itshape\large] at ([xshift=-5mm,yshift=-0.62cm]frame.north east) {Géométrie};}]
\vspace*{1.75cm}
\begin{minipage}[t][6.9cm][c]{\linewidth}\LARGE\centering
\begin{minipage}{0.95\linewidth}\raggedright\linespread{1.25}\selectfont Une homothétie de rapport $4$ transforme le segment $[AB]$, de longueur $1{,}5$ cm, en $[A'B']$. Combien mesure $A'B'$~?\end{minipage}
\end{minipage}

\vfill
\begin{tcolorbox}[enhanced,nobeforeafter,width=7.1cm,height=2.1cm,valign=center,colback=cgeo!6,colframe=cgeo,arc=2.5mm,boxrule=0.9pt,left=3mm]{\Large\bfseries\color{cgeo}a)}\quad\begin{minipage}[c]{5.3cm}\raggedright\Large $7~\text{cm}$\end{minipage}\end{tcolorbox}\hfill\begin{tcolorbox}[enhanced,nobeforeafter,width=7.1cm,height=2.1cm,valign=center,colback=cgeo!6,colframe=cgeo,arc=2.5mm,boxrule=0.9pt,left=3mm]{\Large\bfseries\color{cgeo}b)}\quad\begin{minipage}[c]{5.3cm}\raggedright\Large $6~\text{cm}$\end{minipage}\end{tcolorbox}\par\vspace{2.5mm}
\begin{tcolorbox}[enhanced,nobeforeafter,width=7.1cm,height=2.1cm,valign=center,colback=cgeo!6,colframe=cgeo,arc=2.5mm,boxrule=0.9pt,left=3mm]{\Large\bfseries\color{cgeo}c)}\quad\begin{minipage}[c]{5.3cm}\raggedright\Large $5{,}5~\text{cm}$\end{minipage}\end{tcolorbox}\hfill\begin{tcolorbox}[enhanced,nobeforeafter,width=7.1cm,height=2.1cm,valign=center,colback=cgeo!6,colframe=cgeo,arc=2.5mm,boxrule=0.9pt,left=3mm]{\Large\bfseries\color{cgeo}d)}\quad\begin{minipage}[c]{5.3cm}\raggedright\Large $0{,}375~\text{cm}$\end{minipage}\end{tcolorbox}
\end{tcolorbox}
\newpage
\begin{tcolorbox}[enhanced,width=16cm,height=14cm,colback=white,colframe=cgeo,boxrule=2pt,arc=0mm,
left=6mm,right=6mm,top=0mm,bottom=4mm,sharp corners,
overlay={\fill[cgeo] (frame.north west) rectangle ([yshift=-1.25cm]frame.north east);
\node[anchor=west,text=white,font=\sffamily\bfseries\LARGE] at ([xshift=5mm,yshift=-0.62cm]frame.north west) {Question 59 C};
\node[anchor=east,text=white,font=\sffamily\itshape\large] at ([xshift=-5mm,yshift=-0.62cm]frame.north east) {Géométrie};}]
\vspace*{1.75cm}
\begin{minipage}[t][6.9cm][c]{\linewidth}\LARGE\centering
\begin{minipage}{0.95\linewidth}\raggedright\linespread{1.25}\selectfont Quel est le volume exact d'une boule de rayon $2$ cm~?\end{minipage}
\end{minipage}

\vfill
\begin{tcolorbox}[enhanced,nobeforeafter,width=7.1cm,height=2.1cm,valign=center,colback=cgeo!6,colframe=cgeo,arc=2.5mm,boxrule=0.9pt,left=3mm]{\Large\bfseries\color{cgeo}a)}\quad\begin{minipage}[c]{5.3cm}\raggedright\large $\dfrac{16}{3}\pi~\text{cm}^3$\end{minipage}\end{tcolorbox}\hfill\begin{tcolorbox}[enhanced,nobeforeafter,width=7.1cm,height=2.1cm,valign=center,colback=cgeo!6,colframe=cgeo,arc=2.5mm,boxrule=0.9pt,left=3mm]{\Large\bfseries\color{cgeo}b)}\quad\begin{minipage}[c]{5.3cm}\raggedright\large $\dfrac{32}{3}\pi~\text{cm}^3$\end{minipage}\end{tcolorbox}\par\vspace{2.5mm}
\begin{tcolorbox}[enhanced,nobeforeafter,width=7.1cm,height=2.1cm,valign=center,colback=cgeo!6,colframe=cgeo,arc=2.5mm,boxrule=0.9pt,left=3mm]{\Large\bfseries\color{cgeo}c)}\quad\begin{minipage}[c]{5.3cm}\raggedright\Large $8\pi~\text{cm}^3$\end{minipage}\end{tcolorbox}\hfill\begin{tcolorbox}[enhanced,nobeforeafter,width=7.1cm,height=2.1cm,valign=center,colback=cgeo!6,colframe=cgeo,arc=2.5mm,boxrule=0.9pt,left=3mm]{\Large\bfseries\color{cgeo}d)}\quad\begin{minipage}[c]{5.3cm}\raggedright\Large $16\pi~\text{cm}^3$\end{minipage}\end{tcolorbox}
\end{tcolorbox}
\newpage
\begin{tcolorbox}[enhanced,width=16cm,height=14cm,colback=white,colframe=cgeo,boxrule=2pt,arc=0mm,
left=6mm,right=6mm,top=0mm,bottom=4mm,sharp corners,
overlay={\fill[cgeo] (frame.north west) rectangle ([yshift=-1.25cm]frame.north east);
\node[anchor=west,text=white,font=\sffamily\bfseries\LARGE] at ([xshift=5mm,yshift=-0.62cm]frame.north west) {Question 60 C};
\node[anchor=east,text=white,font=\sffamily\itshape\large] at ([xshift=-5mm,yshift=-0.62cm]frame.north east) {Géométrie};}]
\vspace*{1.75cm}
\begin{minipage}[t][6.9cm][c]{\linewidth}\LARGE\centering
\begin{minipage}{0.95\linewidth}\raggedright\linespread{1.25}\selectfont Quel est le volume exact d'une boule de rayon $1{,}5$ cm~?\end{minipage}
\end{minipage}

\vfill
\begin{tcolorbox}[enhanced,nobeforeafter,width=7.1cm,height=2.1cm,valign=center,colback=cgeo!6,colframe=cgeo,arc=2.5mm,boxrule=0.9pt,left=3mm]{\Large\bfseries\color{cgeo}a)}\quad\begin{minipage}[c]{5.3cm}\raggedright\Large $3\pi~\text{cm}^3$\end{minipage}\end{tcolorbox}\hfill\begin{tcolorbox}[enhanced,nobeforeafter,width=7.1cm,height=2.1cm,valign=center,colback=cgeo!6,colframe=cgeo,arc=2.5mm,boxrule=0.9pt,left=3mm]{\Large\bfseries\color{cgeo}b)}\quad\begin{minipage}[c]{5.3cm}\raggedright\large $\dfrac{27}{8}\pi~\text{cm}^3$\end{minipage}\end{tcolorbox}\par\vspace{2.5mm}
\begin{tcolorbox}[enhanced,nobeforeafter,width=7.1cm,height=2.1cm,valign=center,colback=cgeo!6,colframe=cgeo,arc=2.5mm,boxrule=0.9pt,left=3mm]{\Large\bfseries\color{cgeo}c)}\quad\begin{minipage}[c]{5.3cm}\raggedright\large $\dfrac{9}{2}\pi~\text{cm}^3$\end{minipage}\end{tcolorbox}\hfill\begin{tcolorbox}[enhanced,nobeforeafter,width=7.1cm,height=2.1cm,valign=center,colback=cgeo!6,colframe=cgeo,arc=2.5mm,boxrule=0.9pt,left=3mm]{\Large\bfseries\color{cgeo}d)}\quad\begin{minipage}[c]{5.3cm}\raggedright\Large $9\pi~\text{cm}^3$\end{minipage}\end{tcolorbox}
\end{tcolorbox}
\newpage
\begin{tcolorbox}[enhanced,width=16cm,height=14cm,colback=white,colframe=cgeo,boxrule=2pt,arc=0mm,
left=6mm,right=6mm,top=0mm,bottom=4mm,sharp corners,
overlay={\fill[cgeo] (frame.north west) rectangle ([yshift=-1.25cm]frame.north east);
\node[anchor=west,text=white,font=\sffamily\bfseries\LARGE] at ([xshift=5mm,yshift=-0.62cm]frame.north west) {Question 61 C};
\node[anchor=east,text=white,font=\sffamily\itshape\large] at ([xshift=-5mm,yshift=-0.62cm]frame.north east) {Géométrie};}]
\vspace*{1.75cm}
\begin{minipage}[t][6.9cm][c]{\linewidth}\LARGE\centering
\begin{minipage}{0.95\linewidth}\raggedright\linespread{1.25}\selectfont Quelle est l'aire exacte d'une sphère de rayon $7$ cm~?\end{minipage}
\end{minipage}

\vfill
\begin{tcolorbox}[enhanced,nobeforeafter,width=7.1cm,height=2.1cm,valign=center,colback=cgeo!6,colframe=cgeo,arc=2.5mm,boxrule=0.9pt,left=3mm]{\Large\bfseries\color{cgeo}a)}\quad\begin{minipage}[c]{5.3cm}\raggedright\Large $98\pi~\text{cm}^2$\end{minipage}\end{tcolorbox}\hfill\begin{tcolorbox}[enhanced,nobeforeafter,width=7.1cm,height=2.1cm,valign=center,colback=cgeo!6,colframe=cgeo,arc=2.5mm,boxrule=0.9pt,left=3mm]{\Large\bfseries\color{cgeo}b)}\quad\begin{minipage}[c]{5.3cm}\raggedright\Large $49\pi~\text{cm}^2$\end{minipage}\end{tcolorbox}\par\vspace{2.5mm}
\begin{tcolorbox}[enhanced,nobeforeafter,width=7.1cm,height=2.1cm,valign=center,colback=cgeo!6,colframe=cgeo,arc=2.5mm,boxrule=0.9pt,left=3mm]{\Large\bfseries\color{cgeo}c)}\quad\begin{minipage}[c]{5.3cm}\raggedright\Large $196\pi~\text{cm}^2$\end{minipage}\end{tcolorbox}\hfill\begin{tcolorbox}[enhanced,nobeforeafter,width=7.1cm,height=2.1cm,valign=center,colback=cgeo!6,colframe=cgeo,arc=2.5mm,boxrule=0.9pt,left=3mm]{\Large\bfseries\color{cgeo}d)}\quad\begin{minipage}[c]{5.3cm}\raggedright\large $\dfrac{1372}{3}\pi~\text{cm}^2$\end{minipage}\end{tcolorbox}
\end{tcolorbox}
\newpage
\begin{tcolorbox}[enhanced,width=16cm,height=14cm,colback=white,colframe=cgeo,boxrule=2pt,arc=0mm,
left=6mm,right=6mm,top=0mm,bottom=4mm,sharp corners,
overlay={\fill[cgeo] (frame.north west) rectangle ([yshift=-1.25cm]frame.north east);
\node[anchor=west,text=white,font=\sffamily\bfseries\LARGE] at ([xshift=5mm,yshift=-0.62cm]frame.north west) {Question 62 C};
\node[anchor=east,text=white,font=\sffamily\itshape\large] at ([xshift=-5mm,yshift=-0.62cm]frame.north east) {Géométrie};}]
\vspace*{1.75cm}
\begin{minipage}[t][6.9cm][c]{\linewidth}\LARGE\centering
\begin{minipage}{0.95\linewidth}\raggedright\linespread{1.25}\selectfont Quelle est l'aire exacte d'une sphère de rayon $0{,}5$ cm~?\end{minipage}
\end{minipage}

\vfill
\begin{tcolorbox}[enhanced,nobeforeafter,width=7.1cm,height=2.1cm,valign=center,colback=cgeo!6,colframe=cgeo,arc=2.5mm,boxrule=0.9pt,left=3mm]{\Large\bfseries\color{cgeo}a)}\quad\begin{minipage}[c]{5.3cm}\raggedright\large $\dfrac{1}{4}\pi~\text{cm}^2$\end{minipage}\end{tcolorbox}\hfill\begin{tcolorbox}[enhanced,nobeforeafter,width=7.1cm,height=2.1cm,valign=center,colback=cgeo!6,colframe=cgeo,arc=2.5mm,boxrule=0.9pt,left=3mm]{\Large\bfseries\color{cgeo}b)}\quad\begin{minipage}[c]{5.3cm}\raggedright\large $\dfrac{1}{6}\pi~\text{cm}^2$\end{minipage}\end{tcolorbox}\par\vspace{2.5mm}
\begin{tcolorbox}[enhanced,nobeforeafter,width=7.1cm,height=2.1cm,valign=center,colback=cgeo!6,colframe=cgeo,arc=2.5mm,boxrule=0.9pt,left=3mm]{\Large\bfseries\color{cgeo}c)}\quad\begin{minipage}[c]{5.3cm}\raggedright\Large $\pi~\text{cm}^2$\end{minipage}\end{tcolorbox}\hfill\begin{tcolorbox}[enhanced,nobeforeafter,width=7.1cm,height=2.1cm,valign=center,colback=cgeo!6,colframe=cgeo,arc=2.5mm,boxrule=0.9pt,left=3mm]{\Large\bfseries\color{cgeo}d)}\quad\begin{minipage}[c]{5.3cm}\raggedright\large $\dfrac{1}{2}\pi~\text{cm}^2$\end{minipage}\end{tcolorbox}
\end{tcolorbox}
\newpage
\begin{tcolorbox}[enhanced,width=16cm,height=14cm,colback=white,colframe=cgeo,boxrule=2pt,arc=0mm,
left=6mm,right=6mm,top=0mm,bottom=4mm,sharp corners,
overlay={\fill[cgeo] (frame.north west) rectangle ([yshift=-1.25cm]frame.north east);
\node[anchor=west,text=white,font=\sffamily\bfseries\LARGE] at ([xshift=5mm,yshift=-0.62cm]frame.north west) {Question 63 C};
\node[anchor=east,text=white,font=\sffamily\itshape\large] at ([xshift=-5mm,yshift=-0.62cm]frame.north east) {Géométrie};}]
\vspace*{1.75cm}
\begin{minipage}[t][6.9cm][c]{\linewidth}\LARGE\centering
\begin{minipage}{0.95\linewidth}\raggedright\linespread{1.25}\selectfont Un solide est agrandi à l'échelle $4$. Par combien son volume est-il multiplié~?\end{minipage}
\end{minipage}

\vfill
\begin{tcolorbox}[enhanced,nobeforeafter,width=7.1cm,height=2.1cm,valign=center,colback=cgeo!6,colframe=cgeo,arc=2.5mm,boxrule=0.9pt,left=3mm]{\Large\bfseries\color{cgeo}a)}\quad\begin{minipage}[c]{5.3cm}\raggedright\Large $64$\end{minipage}\end{tcolorbox}\hfill\begin{tcolorbox}[enhanced,nobeforeafter,width=7.1cm,height=2.1cm,valign=center,colback=cgeo!6,colframe=cgeo,arc=2.5mm,boxrule=0.9pt,left=3mm]{\Large\bfseries\color{cgeo}b)}\quad\begin{minipage}[c]{5.3cm}\raggedright\Large $4$\end{minipage}\end{tcolorbox}\par\vspace{2.5mm}
\begin{tcolorbox}[enhanced,nobeforeafter,width=7.1cm,height=2.1cm,valign=center,colback=cgeo!6,colframe=cgeo,arc=2.5mm,boxrule=0.9pt,left=3mm]{\Large\bfseries\color{cgeo}c)}\quad\begin{minipage}[c]{5.3cm}\raggedright\Large $16$\end{minipage}\end{tcolorbox}\hfill\begin{tcolorbox}[enhanced,nobeforeafter,width=7.1cm,height=2.1cm,valign=center,colback=cgeo!6,colframe=cgeo,arc=2.5mm,boxrule=0.9pt,left=3mm]{\Large\bfseries\color{cgeo}d)}\quad\begin{minipage}[c]{5.3cm}\raggedright\Large $12$\end{minipage}\end{tcolorbox}
\end{tcolorbox}
\newpage
\begin{tcolorbox}[enhanced,width=16cm,height=14cm,colback=white,colframe=cgeo,boxrule=2pt,arc=0mm,
left=6mm,right=6mm,top=0mm,bottom=4mm,sharp corners,
overlay={\fill[cgeo] (frame.north west) rectangle ([yshift=-1.25cm]frame.north east);
\node[anchor=west,text=white,font=\sffamily\bfseries\LARGE] at ([xshift=5mm,yshift=-0.62cm]frame.north west) {Question 64 C};
\node[anchor=east,text=white,font=\sffamily\itshape\large] at ([xshift=-5mm,yshift=-0.62cm]frame.north east) {Géométrie};}]
\vspace*{1.75cm}
\begin{minipage}[t][6.9cm][c]{\linewidth}\LARGE\centering
\begin{minipage}{0.95\linewidth}\raggedright\linespread{1.25}\selectfont Un solide est agrandi à l'échelle $\dfrac{3}{2}$. Par combien son volume est-il multiplié~?\end{minipage}
\end{minipage}

\vfill
\begin{tcolorbox}[enhanced,nobeforeafter,width=7.1cm,height=2.1cm,valign=center,colback=cgeo!6,colframe=cgeo,arc=2.5mm,boxrule=0.9pt,left=3mm]{\Large\bfseries\color{cgeo}a)}\quad\begin{minipage}[c]{5.3cm}\raggedright\Large $\dfrac{3}{2}$\end{minipage}\end{tcolorbox}\hfill\begin{tcolorbox}[enhanced,nobeforeafter,width=7.1cm,height=2.1cm,valign=center,colback=cgeo!6,colframe=cgeo,arc=2.5mm,boxrule=0.9pt,left=3mm]{\Large\bfseries\color{cgeo}b)}\quad\begin{minipage}[c]{5.3cm}\raggedright\Large $\dfrac{9}{2}$\end{minipage}\end{tcolorbox}\par\vspace{2.5mm}
\begin{tcolorbox}[enhanced,nobeforeafter,width=7.1cm,height=2.1cm,valign=center,colback=cgeo!6,colframe=cgeo,arc=2.5mm,boxrule=0.9pt,left=3mm]{\Large\bfseries\color{cgeo}c)}\quad\begin{minipage}[c]{5.3cm}\raggedright\Large $\dfrac{9}{4}$\end{minipage}\end{tcolorbox}\hfill\begin{tcolorbox}[enhanced,nobeforeafter,width=7.1cm,height=2.1cm,valign=center,colback=cgeo!6,colframe=cgeo,arc=2.5mm,boxrule=0.9pt,left=3mm]{\Large\bfseries\color{cgeo}d)}\quad\begin{minipage}[c]{5.3cm}\raggedright\Large $\dfrac{27}{8}$\end{minipage}\end{tcolorbox}
\end{tcolorbox}
\newpage
\begin{tcolorbox}[enhanced,width=16cm,height=14cm,colback=white,colframe=csta,boxrule=2pt,arc=0mm,
left=6mm,right=6mm,top=0mm,bottom=4mm,sharp corners,
overlay={\fill[csta] (frame.north west) rectangle ([yshift=-1.25cm]frame.north east);
\node[anchor=west,text=white,font=\sffamily\bfseries\LARGE] at ([xshift=5mm,yshift=-0.62cm]frame.north west) {Question 65 C};
\node[anchor=east,text=white,font=\sffamily\itshape\large] at ([xshift=-5mm,yshift=-0.62cm]frame.north east) {Statistiques et probabilités};}]
\vspace*{1.75cm}
\begin{minipage}[t][6.9cm][c]{\linewidth}\LARGE\centering
\begin{minipage}{0.95\linewidth}\raggedright\linespread{1.25}\selectfont Quelle est la médiane de la série $21\ ;\ 17\ ;\ 30\ ;\ 14\ ;\ 26\ ;\ 19\ ;\ 23$~?\end{minipage}
\end{minipage}

\vfill
\begin{tcolorbox}[enhanced,nobeforeafter,width=7.1cm,height=2.1cm,valign=center,colback=csta!6,colframe=csta,arc=2.5mm,boxrule=0.9pt,left=3mm]{\Large\bfseries\color{csta}a)}\quad\begin{minipage}[c]{5.3cm}\raggedright\Large $14$\end{minipage}\end{tcolorbox}\hfill\begin{tcolorbox}[enhanced,nobeforeafter,width=7.1cm,height=2.1cm,valign=center,colback=csta!6,colframe=csta,arc=2.5mm,boxrule=0.9pt,left=3mm]{\Large\bfseries\color{csta}b)}\quad\begin{minipage}[c]{5.3cm}\raggedright\Large $21$\end{minipage}\end{tcolorbox}\par\vspace{2.5mm}
\begin{tcolorbox}[enhanced,nobeforeafter,width=7.1cm,height=2.1cm,valign=center,colback=csta!6,colframe=csta,arc=2.5mm,boxrule=0.9pt,left=3mm]{\Large\bfseries\color{csta}c)}\quad\begin{minipage}[c]{5.3cm}\raggedright\Large $16$\end{minipage}\end{tcolorbox}\hfill\begin{tcolorbox}[enhanced,nobeforeafter,width=7.1cm,height=2.1cm,valign=center,colback=csta!6,colframe=csta,arc=2.5mm,boxrule=0.9pt,left=3mm]{\Large\bfseries\color{csta}d)}\quad\begin{minipage}[c]{5.3cm}\raggedright\Large $\dfrac{150}{7}$\end{minipage}\end{tcolorbox}
\end{tcolorbox}
\newpage
\begin{tcolorbox}[enhanced,width=16cm,height=14cm,colback=white,colframe=csta,boxrule=2pt,arc=0mm,
left=6mm,right=6mm,top=0mm,bottom=4mm,sharp corners,
overlay={\fill[csta] (frame.north west) rectangle ([yshift=-1.25cm]frame.north east);
\node[anchor=west,text=white,font=\sffamily\bfseries\LARGE] at ([xshift=5mm,yshift=-0.62cm]frame.north west) {Question 66 C};
\node[anchor=east,text=white,font=\sffamily\itshape\large] at ([xshift=-5mm,yshift=-0.62cm]frame.north east) {Statistiques et probabilités};}]
\vspace*{1.75cm}
\begin{minipage}[t][6.9cm][c]{\linewidth}\LARGE\centering
\begin{minipage}{0.95\linewidth}\raggedright\linespread{1.25}\selectfont Quelle est la médiane de la série $4\ ;\ 4\ ;\ 9\ ;\ 2\ ;\ 7\ ;\ 6$~?\end{minipage}
\end{minipage}

\vfill
\begin{tcolorbox}[enhanced,nobeforeafter,width=7.1cm,height=2.1cm,valign=center,colback=csta!6,colframe=csta,arc=2.5mm,boxrule=0.9pt,left=3mm]{\Large\bfseries\color{csta}a)}\quad\begin{minipage}[c]{5.3cm}\raggedright\Large $\dfrac{16}{3}$\end{minipage}\end{tcolorbox}\hfill\begin{tcolorbox}[enhanced,nobeforeafter,width=7.1cm,height=2.1cm,valign=center,colback=csta!6,colframe=csta,arc=2.5mm,boxrule=0.9pt,left=3mm]{\Large\bfseries\color{csta}b)}\quad\begin{minipage}[c]{5.3cm}\raggedright\Large $7$\end{minipage}\end{tcolorbox}\par\vspace{2.5mm}
\begin{tcolorbox}[enhanced,nobeforeafter,width=7.1cm,height=2.1cm,valign=center,colback=csta!6,colframe=csta,arc=2.5mm,boxrule=0.9pt,left=3mm]{\Large\bfseries\color{csta}c)}\quad\begin{minipage}[c]{5.3cm}\raggedright\Large $5$\end{minipage}\end{tcolorbox}\hfill\begin{tcolorbox}[enhanced,nobeforeafter,width=7.1cm,height=2.1cm,valign=center,colback=csta!6,colframe=csta,arc=2.5mm,boxrule=0.9pt,left=3mm]{\Large\bfseries\color{csta}d)}\quad\begin{minipage}[c]{5.3cm}\raggedright\Large $5{,}5$\end{minipage}\end{tcolorbox}
\end{tcolorbox}
\newpage
\begin{tcolorbox}[enhanced,width=16cm,height=14cm,colback=white,colframe=csta,boxrule=2pt,arc=0mm,
left=6mm,right=6mm,top=0mm,bottom=4mm,sharp corners,
overlay={\fill[csta] (frame.north west) rectangle ([yshift=-1.25cm]frame.north east);
\node[anchor=west,text=white,font=\sffamily\bfseries\LARGE] at ([xshift=5mm,yshift=-0.62cm]frame.north west) {Question 67 C};
\node[anchor=east,text=white,font=\sffamily\itshape\large] at ([xshift=-5mm,yshift=-0.62cm]frame.north east) {Statistiques et probabilités};}]
\vspace*{1.75cm}
\begin{minipage}[t][6.9cm][c]{\linewidth}\LARGE\centering
\begin{minipage}{0.95\linewidth}\raggedright\linespread{1.25}\selectfont Quelle est la moyenne de la série $25\ ;\ 30\ ;\ 35\ ;\ 18$~?\end{minipage}
\end{minipage}

\vfill
\begin{tcolorbox}[enhanced,nobeforeafter,width=7.1cm,height=2.1cm,valign=center,colback=csta!6,colframe=csta,arc=2.5mm,boxrule=0.9pt,left=3mm]{\Large\bfseries\color{csta}a)}\quad\begin{minipage}[c]{5.3cm}\raggedright\Large $27$\end{minipage}\end{tcolorbox}\hfill\begin{tcolorbox}[enhanced,nobeforeafter,width=7.1cm,height=2.1cm,valign=center,colback=csta!6,colframe=csta,arc=2.5mm,boxrule=0.9pt,left=3mm]{\Large\bfseries\color{csta}b)}\quad\begin{minipage}[c]{5.3cm}\raggedright\Large $108$\end{minipage}\end{tcolorbox}\par\vspace{2.5mm}
\begin{tcolorbox}[enhanced,nobeforeafter,width=7.1cm,height=2.1cm,valign=center,colback=csta!6,colframe=csta,arc=2.5mm,boxrule=0.9pt,left=3mm]{\Large\bfseries\color{csta}c)}\quad\begin{minipage}[c]{5.3cm}\raggedright\Large $28$\end{minipage}\end{tcolorbox}\hfill\begin{tcolorbox}[enhanced,nobeforeafter,width=7.1cm,height=2.1cm,valign=center,colback=csta!6,colframe=csta,arc=2.5mm,boxrule=0.9pt,left=3mm]{\Large\bfseries\color{csta}d)}\quad\begin{minipage}[c]{5.3cm}\raggedright\Large $27{,}5$\end{minipage}\end{tcolorbox}
\end{tcolorbox}
\newpage
\begin{tcolorbox}[enhanced,width=16cm,height=14cm,colback=white,colframe=csta,boxrule=2pt,arc=0mm,
left=6mm,right=6mm,top=0mm,bottom=4mm,sharp corners,
overlay={\fill[csta] (frame.north west) rectangle ([yshift=-1.25cm]frame.north east);
\node[anchor=west,text=white,font=\sffamily\bfseries\LARGE] at ([xshift=5mm,yshift=-0.62cm]frame.north west) {Question 68 C};
\node[anchor=east,text=white,font=\sffamily\itshape\large] at ([xshift=-5mm,yshift=-0.62cm]frame.north east) {Statistiques et probabilités};}]
\vspace*{1.75cm}
\begin{minipage}[t][6.9cm][c]{\linewidth}\LARGE\centering
\begin{minipage}{0.95\linewidth}\raggedright\linespread{1.25}\selectfont Quelle est la moyenne de la série $3\ ;\ 11\ ;\ 7\ ;\ 9\ ;\ 5$~?\end{minipage}
\end{minipage}

\vfill
\begin{tcolorbox}[enhanced,nobeforeafter,width=7.1cm,height=2.1cm,valign=center,colback=csta!6,colframe=csta,arc=2.5mm,boxrule=0.9pt,left=3mm]{\Large\bfseries\color{csta}a)}\quad\begin{minipage}[c]{5.3cm}\raggedright\Large $8$\end{minipage}\end{tcolorbox}\hfill\begin{tcolorbox}[enhanced,nobeforeafter,width=7.1cm,height=2.1cm,valign=center,colback=csta!6,colframe=csta,arc=2.5mm,boxrule=0.9pt,left=3mm]{\Large\bfseries\color{csta}b)}\quad\begin{minipage}[c]{5.3cm}\raggedright\Large $9$\end{minipage}\end{tcolorbox}\par\vspace{2.5mm}
\begin{tcolorbox}[enhanced,nobeforeafter,width=7.1cm,height=2.1cm,valign=center,colback=csta!6,colframe=csta,arc=2.5mm,boxrule=0.9pt,left=3mm]{\Large\bfseries\color{csta}c)}\quad\begin{minipage}[c]{5.3cm}\raggedright\Large $7$\end{minipage}\end{tcolorbox}\hfill\begin{tcolorbox}[enhanced,nobeforeafter,width=7.1cm,height=2.1cm,valign=center,colback=csta!6,colframe=csta,arc=2.5mm,boxrule=0.9pt,left=3mm]{\Large\bfseries\color{csta}d)}\quad\begin{minipage}[c]{5.3cm}\raggedright\Large $35$\end{minipage}\end{tcolorbox}
\end{tcolorbox}
\newpage
\begin{tcolorbox}[enhanced,width=16cm,height=14cm,colback=white,colframe=csta,boxrule=2pt,arc=0mm,
left=6mm,right=6mm,top=0mm,bottom=4mm,sharp corners,
overlay={\fill[csta] (frame.north west) rectangle ([yshift=-1.25cm]frame.north east);
\node[anchor=west,text=white,font=\sffamily\bfseries\LARGE] at ([xshift=5mm,yshift=-0.62cm]frame.north west) {Question 69 C};
\node[anchor=east,text=white,font=\sffamily\itshape\large] at ([xshift=-5mm,yshift=-0.62cm]frame.north east) {Statistiques et probabilités};}]
\vspace*{1.75cm}
\begin{minipage}[t][6.9cm][c]{\linewidth}\LARGE\centering
\begin{minipage}{0.95\linewidth}\raggedright\linespread{1.25}\selectfont Quelle est l'étendue de la série $25\ ;\ 30\ ;\ 35\ ;\ 18$~?\end{minipage}
\end{minipage}

\vfill
\begin{tcolorbox}[enhanced,nobeforeafter,width=7.1cm,height=2.1cm,valign=center,colback=csta!6,colframe=csta,arc=2.5mm,boxrule=0.9pt,left=3mm]{\Large\bfseries\color{csta}a)}\quad\begin{minipage}[c]{5.3cm}\raggedright\Large $17$\end{minipage}\end{tcolorbox}\hfill\begin{tcolorbox}[enhanced,nobeforeafter,width=7.1cm,height=2.1cm,valign=center,colback=csta!6,colframe=csta,arc=2.5mm,boxrule=0.9pt,left=3mm]{\Large\bfseries\color{csta}b)}\quad\begin{minipage}[c]{5.3cm}\raggedright\Large $35$\end{minipage}\end{tcolorbox}\par\vspace{2.5mm}
\begin{tcolorbox}[enhanced,nobeforeafter,width=7.1cm,height=2.1cm,valign=center,colback=csta!6,colframe=csta,arc=2.5mm,boxrule=0.9pt,left=3mm]{\Large\bfseries\color{csta}c)}\quad\begin{minipage}[c]{5.3cm}\raggedright\Large $53$\end{minipage}\end{tcolorbox}\hfill\begin{tcolorbox}[enhanced,nobeforeafter,width=7.1cm,height=2.1cm,valign=center,colback=csta!6,colframe=csta,arc=2.5mm,boxrule=0.9pt,left=3mm]{\Large\bfseries\color{csta}d)}\quad\begin{minipage}[c]{5.3cm}\raggedright\Large $18$\end{minipage}\end{tcolorbox}
\end{tcolorbox}
\newpage
\begin{tcolorbox}[enhanced,width=16cm,height=14cm,colback=white,colframe=csta,boxrule=2pt,arc=0mm,
left=6mm,right=6mm,top=0mm,bottom=4mm,sharp corners,
overlay={\fill[csta] (frame.north west) rectangle ([yshift=-1.25cm]frame.north east);
\node[anchor=west,text=white,font=\sffamily\bfseries\LARGE] at ([xshift=5mm,yshift=-0.62cm]frame.north west) {Question 70 C};
\node[anchor=east,text=white,font=\sffamily\itshape\large] at ([xshift=-5mm,yshift=-0.62cm]frame.north east) {Statistiques et probabilités};}]
\vspace*{1.75cm}
\begin{minipage}[t][6.9cm][c]{\linewidth}\LARGE\centering
\begin{minipage}{0.95\linewidth}\raggedright\linespread{1.25}\selectfont Quelle est l'étendue de la série $3\ ;\ 11\ ;\ 7\ ;\ 9\ ;\ 5$~?\end{minipage}
\end{minipage}

\vfill
\begin{tcolorbox}[enhanced,nobeforeafter,width=7.1cm,height=2.1cm,valign=center,colback=csta!6,colframe=csta,arc=2.5mm,boxrule=0.9pt,left=3mm]{\Large\bfseries\color{csta}a)}\quad\begin{minipage}[c]{5.3cm}\raggedright\Large $9$\end{minipage}\end{tcolorbox}\hfill\begin{tcolorbox}[enhanced,nobeforeafter,width=7.1cm,height=2.1cm,valign=center,colback=csta!6,colframe=csta,arc=2.5mm,boxrule=0.9pt,left=3mm]{\Large\bfseries\color{csta}b)}\quad\begin{minipage}[c]{5.3cm}\raggedright\Large $14$\end{minipage}\end{tcolorbox}\par\vspace{2.5mm}
\begin{tcolorbox}[enhanced,nobeforeafter,width=7.1cm,height=2.1cm,valign=center,colback=csta!6,colframe=csta,arc=2.5mm,boxrule=0.9pt,left=3mm]{\Large\bfseries\color{csta}c)}\quad\begin{minipage}[c]{5.3cm}\raggedright\Large $11$\end{minipage}\end{tcolorbox}\hfill\begin{tcolorbox}[enhanced,nobeforeafter,width=7.1cm,height=2.1cm,valign=center,colback=csta!6,colframe=csta,arc=2.5mm,boxrule=0.9pt,left=3mm]{\Large\bfseries\color{csta}d)}\quad\begin{minipage}[c]{5.3cm}\raggedright\Large $8$\end{minipage}\end{tcolorbox}
\end{tcolorbox}
\newpage
\begin{tcolorbox}[enhanced,width=16cm,height=14cm,colback=white,colframe=csta,boxrule=2pt,arc=0mm,
left=6mm,right=6mm,top=0mm,bottom=4mm,sharp corners,
overlay={\fill[csta] (frame.north west) rectangle ([yshift=-1.25cm]frame.north east);
\node[anchor=west,text=white,font=\sffamily\bfseries\LARGE] at ([xshift=5mm,yshift=-0.62cm]frame.north west) {Question 71 C};
\node[anchor=east,text=white,font=\sffamily\itshape\large] at ([xshift=-5mm,yshift=-0.62cm]frame.north east) {Statistiques et probabilités};}]
\vspace*{1.75cm}
\begin{minipage}[t][6.9cm][c]{\linewidth}\Large\centering
\begin{minipage}{0.95\linewidth}\raggedright\linespread{1.25}\selectfont Une urne contient $7$ boules dont $3$ rouges. On tire successivement deux boules, sans remise. Quelle est la probabilité de tirer deux boules rouges~?\end{minipage}
\end{minipage}

\vfill
\begin{tcolorbox}[enhanced,nobeforeafter,width=7.1cm,height=2.1cm,valign=center,colback=csta!6,colframe=csta,arc=2.5mm,boxrule=0.9pt,left=3mm]{\Large\bfseries\color{csta}a)}\quad\begin{minipage}[c]{5.3cm}\raggedright\Large $\dfrac{9}{49}$\end{minipage}\end{tcolorbox}\hfill\begin{tcolorbox}[enhanced,nobeforeafter,width=7.1cm,height=2.1cm,valign=center,colback=csta!6,colframe=csta,arc=2.5mm,boxrule=0.9pt,left=3mm]{\Large\bfseries\color{csta}b)}\quad\begin{minipage}[c]{5.3cm}\raggedright\Large $\dfrac{6}{49}$\end{minipage}\end{tcolorbox}\par\vspace{2.5mm}
\begin{tcolorbox}[enhanced,nobeforeafter,width=7.1cm,height=2.1cm,valign=center,colback=csta!6,colframe=csta,arc=2.5mm,boxrule=0.9pt,left=3mm]{\Large\bfseries\color{csta}c)}\quad\begin{minipage}[c]{5.3cm}\raggedright\Large $\dfrac{1}{7}$\end{minipage}\end{tcolorbox}\hfill\begin{tcolorbox}[enhanced,nobeforeafter,width=7.1cm,height=2.1cm,valign=center,colback=csta!6,colframe=csta,arc=2.5mm,boxrule=0.9pt,left=3mm]{\Large\bfseries\color{csta}d)}\quad\begin{minipage}[c]{5.3cm}\raggedright\Large $\dfrac{6}{7}$\end{minipage}\end{tcolorbox}
\end{tcolorbox}
\newpage
\begin{tcolorbox}[enhanced,width=16cm,height=14cm,colback=white,colframe=csta,boxrule=2pt,arc=0mm,
left=6mm,right=6mm,top=0mm,bottom=4mm,sharp corners,
overlay={\fill[csta] (frame.north west) rectangle ([yshift=-1.25cm]frame.north east);
\node[anchor=west,text=white,font=\sffamily\bfseries\LARGE] at ([xshift=5mm,yshift=-0.62cm]frame.north west) {Question 72 C};
\node[anchor=east,text=white,font=\sffamily\itshape\large] at ([xshift=-5mm,yshift=-0.62cm]frame.north east) {Statistiques et probabilités};}]
\vspace*{1.75cm}
\begin{minipage}[t][6.9cm][c]{\linewidth}\Large\centering
\begin{minipage}{0.95\linewidth}\raggedright\linespread{1.25}\selectfont Une urne contient $8$ boules dont $5$ rouges. On tire successivement deux boules, sans remise. Quelle est la probabilité de tirer deux boules rouges~?\end{minipage}
\end{minipage}

\vfill
\begin{tcolorbox}[enhanced,nobeforeafter,width=7.1cm,height=2.1cm,valign=center,colback=csta!6,colframe=csta,arc=2.5mm,boxrule=0.9pt,left=3mm]{\Large\bfseries\color{csta}a)}\quad\begin{minipage}[c]{5.3cm}\raggedright\Large $\dfrac{5}{16}$\end{minipage}\end{tcolorbox}\hfill\begin{tcolorbox}[enhanced,nobeforeafter,width=7.1cm,height=2.1cm,valign=center,colback=csta!6,colframe=csta,arc=2.5mm,boxrule=0.9pt,left=3mm]{\Large\bfseries\color{csta}b)}\quad\begin{minipage}[c]{5.3cm}\raggedright\Large $\dfrac{5}{14}$\end{minipage}\end{tcolorbox}\par\vspace{2.5mm}
\begin{tcolorbox}[enhanced,nobeforeafter,width=7.1cm,height=2.1cm,valign=center,colback=csta!6,colframe=csta,arc=2.5mm,boxrule=0.9pt,left=3mm]{\Large\bfseries\color{csta}c)}\quad\begin{minipage}[c]{5.3cm}\raggedright\Large $\dfrac{25}{64}$\end{minipage}\end{tcolorbox}\hfill\begin{tcolorbox}[enhanced,nobeforeafter,width=7.1cm,height=2.1cm,valign=center,colback=csta!6,colframe=csta,arc=2.5mm,boxrule=0.9pt,left=3mm]{\Large\bfseries\color{csta}d)}\quad\begin{minipage}[c]{5.3cm}\raggedright\Large $\dfrac{5}{8}$\end{minipage}\end{tcolorbox}
\end{tcolorbox}
\newpage
\begin{tcolorbox}[enhanced,width=16cm,height=14cm,colback=white,colframe=csta,boxrule=2pt,arc=0mm,
left=6mm,right=6mm,top=0mm,bottom=4mm,sharp corners,
overlay={\fill[csta] (frame.north west) rectangle ([yshift=-1.25cm]frame.north east);
\node[anchor=west,text=white,font=\sffamily\bfseries\LARGE] at ([xshift=5mm,yshift=-0.62cm]frame.north west) {Question 73 C};
\node[anchor=east,text=white,font=\sffamily\itshape\large] at ([xshift=-5mm,yshift=-0.62cm]frame.north east) {Statistiques et probabilités};}]
\vspace*{1.75cm}
\begin{minipage}[t][6.9cm][c]{\linewidth}\Large\centering
\begin{minipage}{0.95\linewidth}\raggedright\linespread{1.25}\selectfont On lance un dé équilibré à six faces et une pièce équilibrée (ou un deuxième dé). Quelle est la probabilité d'obtenir un multiple de 3 avec le dé et PILE avec la pièce~?\end{minipage}
\end{minipage}

\vfill
\begin{tcolorbox}[enhanced,nobeforeafter,width=7.1cm,height=2.1cm,valign=center,colback=csta!6,colframe=csta,arc=2.5mm,boxrule=0.9pt,left=3mm]{\Large\bfseries\color{csta}a)}\quad\begin{minipage}[c]{5.3cm}\raggedright\Large $\dfrac{1}{3}$\end{minipage}\end{tcolorbox}\hfill\begin{tcolorbox}[enhanced,nobeforeafter,width=7.1cm,height=2.1cm,valign=center,colback=csta!6,colframe=csta,arc=2.5mm,boxrule=0.9pt,left=3mm]{\Large\bfseries\color{csta}b)}\quad\begin{minipage}[c]{5.3cm}\raggedright\Large $\dfrac{1}{6}$\end{minipage}\end{tcolorbox}\par\vspace{2.5mm}
\begin{tcolorbox}[enhanced,nobeforeafter,width=7.1cm,height=2.1cm,valign=center,colback=csta!6,colframe=csta,arc=2.5mm,boxrule=0.9pt,left=3mm]{\Large\bfseries\color{csta}c)}\quad\begin{minipage}[c]{5.3cm}\raggedright\Large $\dfrac{1}{2}$\end{minipage}\end{tcolorbox}\hfill\begin{tcolorbox}[enhanced,nobeforeafter,width=7.1cm,height=2.1cm,valign=center,colback=csta!6,colframe=csta,arc=2.5mm,boxrule=0.9pt,left=3mm]{\Large\bfseries\color{csta}d)}\quad\begin{minipage}[c]{5.3cm}\raggedright\Large $\dfrac{5}{6}$\end{minipage}\end{tcolorbox}
\end{tcolorbox}
\newpage
\begin{tcolorbox}[enhanced,width=16cm,height=14cm,colback=white,colframe=csta,boxrule=2pt,arc=0mm,
left=6mm,right=6mm,top=0mm,bottom=4mm,sharp corners,
overlay={\fill[csta] (frame.north west) rectangle ([yshift=-1.25cm]frame.north east);
\node[anchor=west,text=white,font=\sffamily\bfseries\LARGE] at ([xshift=5mm,yshift=-0.62cm]frame.north west) {Question 74 C};
\node[anchor=east,text=white,font=\sffamily\itshape\large] at ([xshift=-5mm,yshift=-0.62cm]frame.north east) {Statistiques et probabilités};}]
\vspace*{1.75cm}
\begin{minipage}[t][6.9cm][c]{\linewidth}\Large\centering
\begin{minipage}{0.95\linewidth}\raggedright\linespread{1.25}\selectfont On lance un dé équilibré à six faces et une pièce équilibrée (ou un deuxième dé). Quelle est la probabilité d'obtenir un nombre impair avec le dé et un 6 avec un deuxième dé~?\end{minipage}
\end{minipage}

\vfill
\begin{tcolorbox}[enhanced,nobeforeafter,width=7.1cm,height=2.1cm,valign=center,colback=csta!6,colframe=csta,arc=2.5mm,boxrule=0.9pt,left=3mm]{\Large\bfseries\color{csta}a)}\quad\begin{minipage}[c]{5.3cm}\raggedright\Large $\dfrac{1}{12}$\end{minipage}\end{tcolorbox}\hfill\begin{tcolorbox}[enhanced,nobeforeafter,width=7.1cm,height=2.1cm,valign=center,colback=csta!6,colframe=csta,arc=2.5mm,boxrule=0.9pt,left=3mm]{\Large\bfseries\color{csta}b)}\quad\begin{minipage}[c]{5.3cm}\raggedright\Large $\dfrac{1}{6}$\end{minipage}\end{tcolorbox}\par\vspace{2.5mm}
\begin{tcolorbox}[enhanced,nobeforeafter,width=7.1cm,height=2.1cm,valign=center,colback=csta!6,colframe=csta,arc=2.5mm,boxrule=0.9pt,left=3mm]{\Large\bfseries\color{csta}c)}\quad\begin{minipage}[c]{5.3cm}\raggedright\Large $\dfrac{1}{2}$\end{minipage}\end{tcolorbox}\hfill\begin{tcolorbox}[enhanced,nobeforeafter,width=7.1cm,height=2.1cm,valign=center,colback=csta!6,colframe=csta,arc=2.5mm,boxrule=0.9pt,left=3mm]{\Large\bfseries\color{csta}d)}\quad\begin{minipage}[c]{5.3cm}\raggedright\Large $\dfrac{2}{3}$\end{minipage}\end{tcolorbox}
\end{tcolorbox}
\newpage
\begin{tcolorbox}[enhanced,width=16cm,height=14cm,colback=white,colframe=csta,boxrule=2pt,arc=0mm,
left=6mm,right=6mm,top=0mm,bottom=4mm,sharp corners,
overlay={\fill[csta] (frame.north west) rectangle ([yshift=-1.25cm]frame.north east);
\node[anchor=west,text=white,font=\sffamily\bfseries\LARGE] at ([xshift=5mm,yshift=-0.62cm]frame.north west) {Question 75 C};
\node[anchor=east,text=white,font=\sffamily\itshape\large] at ([xshift=-5mm,yshift=-0.62cm]frame.north east) {Statistiques et probabilités};}]
\vspace*{1.75cm}
\begin{minipage}[t][6.9cm][c]{\linewidth}\LARGE\centering
\begin{minipage}{0.95\linewidth}\raggedright\linespread{1.25}\selectfont Sur $56$ personnes interrogées, $14$ ont répondu «~oui~». Quelle est la fréquence, en pourcentage, des «~oui~»~?\end{minipage}
\end{minipage}

\vfill
\begin{tcolorbox}[enhanced,nobeforeafter,width=7.1cm,height=2.1cm,valign=center,colback=csta!6,colframe=csta,arc=2.5mm,boxrule=0.9pt,left=3mm]{\Large\bfseries\color{csta}a)}\quad\begin{minipage}[c]{5.3cm}\raggedright\Large $25\,\%$\end{minipage}\end{tcolorbox}\hfill\begin{tcolorbox}[enhanced,nobeforeafter,width=7.1cm,height=2.1cm,valign=center,colback=csta!6,colframe=csta,arc=2.5mm,boxrule=0.9pt,left=3mm]{\Large\bfseries\color{csta}b)}\quad\begin{minipage}[c]{5.3cm}\raggedright\Large $2{,}5\,\%$\end{minipage}\end{tcolorbox}\par\vspace{2.5mm}
\begin{tcolorbox}[enhanced,nobeforeafter,width=7.1cm,height=2.1cm,valign=center,colback=csta!6,colframe=csta,arc=2.5mm,boxrule=0.9pt,left=3mm]{\Large\bfseries\color{csta}c)}\quad\begin{minipage}[c]{5.3cm}\raggedright\Large $75\,\%$\end{minipage}\end{tcolorbox}\hfill\begin{tcolorbox}[enhanced,nobeforeafter,width=7.1cm,height=2.1cm,valign=center,colback=csta!6,colframe=csta,arc=2.5mm,boxrule=0.9pt,left=3mm]{\Large\bfseries\color{csta}d)}\quad\begin{minipage}[c]{5.3cm}\raggedright\Large $14\,\%$\end{minipage}\end{tcolorbox}
\end{tcolorbox}
\newpage
\begin{tcolorbox}[enhanced,width=16cm,height=14cm,colback=white,colframe=csta,boxrule=2pt,arc=0mm,
left=6mm,right=6mm,top=0mm,bottom=4mm,sharp corners,
overlay={\fill[csta] (frame.north west) rectangle ([yshift=-1.25cm]frame.north east);
\node[anchor=west,text=white,font=\sffamily\bfseries\LARGE] at ([xshift=5mm,yshift=-0.62cm]frame.north west) {Question 76 C};
\node[anchor=east,text=white,font=\sffamily\itshape\large] at ([xshift=-5mm,yshift=-0.62cm]frame.north east) {Statistiques et probabilités};}]
\vspace*{1.75cm}
\begin{minipage}[t][6.9cm][c]{\linewidth}\LARGE\centering
\begin{minipage}{0.95\linewidth}\raggedright\linespread{1.25}\selectfont Dans une classe de $30$ élèves, $21$ pratiquent un sport en club. Quelle est la fréquence, en pourcentage, de ces élèves~?\end{minipage}
\end{minipage}

\vfill
\begin{tcolorbox}[enhanced,nobeforeafter,width=7.1cm,height=2.1cm,valign=center,colback=csta!6,colframe=csta,arc=2.5mm,boxrule=0.9pt,left=3mm]{\Large\bfseries\color{csta}a)}\quad\begin{minipage}[c]{5.3cm}\raggedright\Large $21\,\%$\end{minipage}\end{tcolorbox}\hfill\begin{tcolorbox}[enhanced,nobeforeafter,width=7.1cm,height=2.1cm,valign=center,colback=csta!6,colframe=csta,arc=2.5mm,boxrule=0.9pt,left=3mm]{\Large\bfseries\color{csta}b)}\quad\begin{minipage}[c]{5.3cm}\raggedright\Large $7\,\%$\end{minipage}\end{tcolorbox}\par\vspace{2.5mm}
\begin{tcolorbox}[enhanced,nobeforeafter,width=7.1cm,height=2.1cm,valign=center,colback=csta!6,colframe=csta,arc=2.5mm,boxrule=0.9pt,left=3mm]{\Large\bfseries\color{csta}c)}\quad\begin{minipage}[c]{5.3cm}\raggedright\Large $30\,\%$\end{minipage}\end{tcolorbox}\hfill\begin{tcolorbox}[enhanced,nobeforeafter,width=7.1cm,height=2.1cm,valign=center,colback=csta!6,colframe=csta,arc=2.5mm,boxrule=0.9pt,left=3mm]{\Large\bfseries\color{csta}d)}\quad\begin{minipage}[c]{5.3cm}\raggedright\Large $70\,\%$\end{minipage}\end{tcolorbox}
\end{tcolorbox}
\newpage
\begin{tcolorbox}[enhanced,width=16cm,height=14cm,colback=white,colframe=csta,boxrule=2pt,arc=0mm,
left=6mm,right=6mm,top=0mm,bottom=4mm,sharp corners,
overlay={\fill[csta] (frame.north west) rectangle ([yshift=-1.25cm]frame.north east);
\node[anchor=west,text=white,font=\sffamily\bfseries\LARGE] at ([xshift=5mm,yshift=-0.62cm]frame.north west) {Question 77 C};
\node[anchor=east,text=white,font=\sffamily\itshape\large] at ([xshift=-5mm,yshift=-0.62cm]frame.north east) {Statistiques et probabilités};}]
\vspace*{1.75cm}
\begin{minipage}[t][6.9cm][c]{\linewidth}\Large\centering
\begin{minipage}{0.95\linewidth}\raggedright\linespread{1.25}\selectfont Un sac contient $6$ boules rouges, $6$ boules bleues et $3$ boules vertes. On tire une boule au hasard. Quelle est la probabilité de ne pas tirer une boule rouge~?\end{minipage}
\end{minipage}

\vfill
\begin{tcolorbox}[enhanced,nobeforeafter,width=7.1cm,height=2.1cm,valign=center,colback=csta!6,colframe=csta,arc=2.5mm,boxrule=0.9pt,left=3mm]{\Large\bfseries\color{csta}a)}\quad\begin{minipage}[c]{5.3cm}\raggedright\Large $\dfrac{3}{5}$\end{minipage}\end{tcolorbox}\hfill\begin{tcolorbox}[enhanced,nobeforeafter,width=7.1cm,height=2.1cm,valign=center,colback=csta!6,colframe=csta,arc=2.5mm,boxrule=0.9pt,left=3mm]{\Large\bfseries\color{csta}b)}\quad\begin{minipage}[c]{5.3cm}\raggedright\Large $\dfrac{9}{16}$\end{minipage}\end{tcolorbox}\par\vspace{2.5mm}
\begin{tcolorbox}[enhanced,nobeforeafter,width=7.1cm,height=2.1cm,valign=center,colback=csta!6,colframe=csta,arc=2.5mm,boxrule=0.9pt,left=3mm]{\Large\bfseries\color{csta}c)}\quad\begin{minipage}[c]{5.3cm}\raggedright\Large $\dfrac{2}{5}$\end{minipage}\end{tcolorbox}\hfill\begin{tcolorbox}[enhanced,nobeforeafter,width=7.1cm,height=2.1cm,valign=center,colback=csta!6,colframe=csta,arc=2.5mm,boxrule=0.9pt,left=3mm]{\Large\bfseries\color{csta}d)}\quad\begin{minipage}[c]{5.3cm}\raggedright\Large $\dfrac{3}{2}$\end{minipage}\end{tcolorbox}
\end{tcolorbox}
\newpage
\begin{tcolorbox}[enhanced,width=16cm,height=14cm,colback=white,colframe=csta,boxrule=2pt,arc=0mm,
left=6mm,right=6mm,top=0mm,bottom=4mm,sharp corners,
overlay={\fill[csta] (frame.north west) rectangle ([yshift=-1.25cm]frame.north east);
\node[anchor=west,text=white,font=\sffamily\bfseries\LARGE] at ([xshift=5mm,yshift=-0.62cm]frame.north west) {Question 78 C};
\node[anchor=east,text=white,font=\sffamily\itshape\large] at ([xshift=-5mm,yshift=-0.62cm]frame.north east) {Statistiques et probabilités};}]
\vspace*{1.75cm}
\begin{minipage}[t][6.9cm][c]{\linewidth}\Large\centering
\begin{minipage}{0.95\linewidth}\raggedright\linespread{1.25}\selectfont Un sac contient $3$ boules rouges, $9$ boules bleues et $3$ boules vertes. On tire une boule au hasard. Quelle est la probabilité de ne pas tirer une boule rouge~?\end{minipage}
\end{minipage}

\vfill
\begin{tcolorbox}[enhanced,nobeforeafter,width=7.1cm,height=2.1cm,valign=center,colback=csta!6,colframe=csta,arc=2.5mm,boxrule=0.9pt,left=3mm]{\Large\bfseries\color{csta}a)}\quad\begin{minipage}[c]{5.3cm}\raggedright\Large $\dfrac{1}{5}$\end{minipage}\end{tcolorbox}\hfill\begin{tcolorbox}[enhanced,nobeforeafter,width=7.1cm,height=2.1cm,valign=center,colback=csta!6,colframe=csta,arc=2.5mm,boxrule=0.9pt,left=3mm]{\Large\bfseries\color{csta}b)}\quad\begin{minipage}[c]{5.3cm}\raggedright\Large $4$\end{minipage}\end{tcolorbox}\par\vspace{2.5mm}
\begin{tcolorbox}[enhanced,nobeforeafter,width=7.1cm,height=2.1cm,valign=center,colback=csta!6,colframe=csta,arc=2.5mm,boxrule=0.9pt,left=3mm]{\Large\bfseries\color{csta}c)}\quad\begin{minipage}[c]{5.3cm}\raggedright\Large $\dfrac{3}{5}$\end{minipage}\end{tcolorbox}\hfill\begin{tcolorbox}[enhanced,nobeforeafter,width=7.1cm,height=2.1cm,valign=center,colback=csta!6,colframe=csta,arc=2.5mm,boxrule=0.9pt,left=3mm]{\Large\bfseries\color{csta}d)}\quad\begin{minipage}[c]{5.3cm}\raggedright\Large $\dfrac{4}{5}$\end{minipage}\end{tcolorbox}
\end{tcolorbox}
\newpage
\begin{tcolorbox}[enhanced,width=16cm,height=14cm,colback=white,colframe=csta,boxrule=2pt,arc=0mm,
left=6mm,right=6mm,top=0mm,bottom=4mm,sharp corners,
overlay={\fill[csta] (frame.north west) rectangle ([yshift=-1.25cm]frame.north east);
\node[anchor=west,text=white,font=\sffamily\bfseries\LARGE] at ([xshift=5mm,yshift=-0.62cm]frame.north west) {Question 79 C};
\node[anchor=east,text=white,font=\sffamily\itshape\large] at ([xshift=-5mm,yshift=-0.62cm]frame.north east) {Statistiques et probabilités};}]
\vspace*{1.75cm}
\begin{minipage}[t][6.9cm][c]{\linewidth}\LARGE\centering
\begin{minipage}{0.95\linewidth}\raggedright\linespread{1.25}\selectfont La probabilité d'un événement $A$ est $\dfrac{1}{8}$. Quelle est la probabilité de l'événement contraire de $A$~?\end{minipage}
\end{minipage}

\vfill
\begin{tcolorbox}[enhanced,nobeforeafter,width=7.1cm,height=2.1cm,valign=center,colback=csta!6,colframe=csta,arc=2.5mm,boxrule=0.9pt,left=3mm]{\Large\bfseries\color{csta}a)}\quad\begin{minipage}[c]{5.3cm}\raggedright\Large $\dfrac{7}{16}$\end{minipage}\end{tcolorbox}\hfill\begin{tcolorbox}[enhanced,nobeforeafter,width=7.1cm,height=2.1cm,valign=center,colback=csta!6,colframe=csta,arc=2.5mm,boxrule=0.9pt,left=3mm]{\Large\bfseries\color{csta}b)}\quad\begin{minipage}[c]{5.3cm}\raggedright\Large $\dfrac{3}{8}$\end{minipage}\end{tcolorbox}\par\vspace{2.5mm}
\begin{tcolorbox}[enhanced,nobeforeafter,width=7.1cm,height=2.1cm,valign=center,colback=csta!6,colframe=csta,arc=2.5mm,boxrule=0.9pt,left=3mm]{\Large\bfseries\color{csta}c)}\quad\begin{minipage}[c]{5.3cm}\raggedright\Large $\dfrac{7}{8}$\end{minipage}\end{tcolorbox}\hfill\begin{tcolorbox}[enhanced,nobeforeafter,width=7.1cm,height=2.1cm,valign=center,colback=csta!6,colframe=csta,arc=2.5mm,boxrule=0.9pt,left=3mm]{\Large\bfseries\color{csta}d)}\quad\begin{minipage}[c]{5.3cm}\raggedright\Large $\dfrac{1}{8}$\end{minipage}\end{tcolorbox}
\end{tcolorbox}
\newpage
\begin{tcolorbox}[enhanced,width=16cm,height=14cm,colback=white,colframe=csta,boxrule=2pt,arc=0mm,
left=6mm,right=6mm,top=0mm,bottom=4mm,sharp corners,
overlay={\fill[csta] (frame.north west) rectangle ([yshift=-1.25cm]frame.north east);
\node[anchor=west,text=white,font=\sffamily\bfseries\LARGE] at ([xshift=5mm,yshift=-0.62cm]frame.north west) {Question 80 C};
\node[anchor=east,text=white,font=\sffamily\itshape\large] at ([xshift=-5mm,yshift=-0.62cm]frame.north east) {Statistiques et probabilités};}]
\vspace*{1.75cm}
\begin{minipage}[t][6.9cm][c]{\linewidth}\LARGE\centering
\begin{minipage}{0.95\linewidth}\raggedright\linespread{1.25}\selectfont La probabilité d'un événement $A$ est $0{,}4$. Quelle est la probabilité de l'événement contraire de $A$~?\end{minipage}
\end{minipage}

\vfill
\begin{tcolorbox}[enhanced,nobeforeafter,width=7.1cm,height=2.1cm,valign=center,colback=csta!6,colframe=csta,arc=2.5mm,boxrule=0.9pt,left=3mm]{\Large\bfseries\color{csta}a)}\quad\begin{minipage}[c]{5.3cm}\raggedright\Large $0{,}1$\end{minipage}\end{tcolorbox}\hfill\begin{tcolorbox}[enhanced,nobeforeafter,width=7.1cm,height=2.1cm,valign=center,colback=csta!6,colframe=csta,arc=2.5mm,boxrule=0.9pt,left=3mm]{\Large\bfseries\color{csta}b)}\quad\begin{minipage}[c]{5.3cm}\raggedright\Large $0{,}3$\end{minipage}\end{tcolorbox}\par\vspace{2.5mm}
\begin{tcolorbox}[enhanced,nobeforeafter,width=7.1cm,height=2.1cm,valign=center,colback=csta!6,colframe=csta,arc=2.5mm,boxrule=0.9pt,left=3mm]{\Large\bfseries\color{csta}c)}\quad\begin{minipage}[c]{5.3cm}\raggedright\Large $0{,}4$\end{minipage}\end{tcolorbox}\hfill\begin{tcolorbox}[enhanced,nobeforeafter,width=7.1cm,height=2.1cm,valign=center,colback=csta!6,colframe=csta,arc=2.5mm,boxrule=0.9pt,left=3mm]{\Large\bfseries\color{csta}d)}\quad\begin{minipage}[c]{5.3cm}\raggedright\Large $0{,}6$\end{minipage}\end{tcolorbox}
\end{tcolorbox}
\newpage
\begin{tcolorbox}[enhanced,width=16cm,height=14cm,colback=white,colframe=cgra,boxrule=2pt,arc=0mm,
left=6mm,right=6mm,top=0mm,bottom=4mm,sharp corners,
overlay={\fill[cgra] (frame.north west) rectangle ([yshift=-1.25cm]frame.north east);
\node[anchor=west,text=white,font=\sffamily\bfseries\LARGE] at ([xshift=5mm,yshift=-0.62cm]frame.north west) {Question 81 C};
\node[anchor=east,text=white,font=\sffamily\itshape\large] at ([xshift=-5mm,yshift=-0.62cm]frame.north east) {Grandeurs et proportionnalité};}]
\vspace*{1.75cm}
\begin{minipage}[t][6.9cm][c]{\linewidth}\LARGE\centering
\begin{minipage}{0.95\linewidth}\raggedright\linespread{1.25}\selectfont Une vitesse de $20$ m/s correspond à~:\end{minipage}
\end{minipage}

\vfill
\begin{tcolorbox}[enhanced,nobeforeafter,width=7.1cm,height=2.1cm,valign=center,colback=cgra!6,colframe=cgra,arc=2.5mm,boxrule=0.9pt,left=3mm]{\Large\bfseries\color{cgra}a)}\quad\begin{minipage}[c]{5.3cm}\raggedright\Large $200~\text{km/h}$\end{minipage}\end{tcolorbox}\hfill\begin{tcolorbox}[enhanced,nobeforeafter,width=7.1cm,height=2.1cm,valign=center,colback=cgra!6,colframe=cgra,arc=2.5mm,boxrule=0.9pt,left=3mm]{\Large\bfseries\color{cgra}b)}\quad\begin{minipage}[c]{5.3cm}\raggedright\Large $1\,200~\text{km/h}$\end{minipage}\end{tcolorbox}\par\vspace{2.5mm}
\begin{tcolorbox}[enhanced,nobeforeafter,width=7.1cm,height=2.1cm,valign=center,colback=cgra!6,colframe=cgra,arc=2.5mm,boxrule=0.9pt,left=3mm]{\Large\bfseries\color{cgra}c)}\quad\begin{minipage}[c]{5.3cm}\raggedright\Large $20\,000~\text{km/h}$\end{minipage}\end{tcolorbox}\hfill\begin{tcolorbox}[enhanced,nobeforeafter,width=7.1cm,height=2.1cm,valign=center,colback=cgra!6,colframe=cgra,arc=2.5mm,boxrule=0.9pt,left=3mm]{\Large\bfseries\color{cgra}d)}\quad\begin{minipage}[c]{5.3cm}\raggedright\Large $72~\text{km/h}$\end{minipage}\end{tcolorbox}
\end{tcolorbox}
\newpage
\begin{tcolorbox}[enhanced,width=16cm,height=14cm,colback=white,colframe=cgra,boxrule=2pt,arc=0mm,
left=6mm,right=6mm,top=0mm,bottom=4mm,sharp corners,
overlay={\fill[cgra] (frame.north west) rectangle ([yshift=-1.25cm]frame.north east);
\node[anchor=west,text=white,font=\sffamily\bfseries\LARGE] at ([xshift=5mm,yshift=-0.62cm]frame.north west) {Question 82 C};
\node[anchor=east,text=white,font=\sffamily\itshape\large] at ([xshift=-5mm,yshift=-0.62cm]frame.north east) {Grandeurs et proportionnalité};}]
\vspace*{1.75cm}
\begin{minipage}[t][6.9cm][c]{\linewidth}\LARGE\centering
\begin{minipage}{0.95\linewidth}\raggedright\linespread{1.25}\selectfont Une vitesse de $30$ m/s correspond à~:\end{minipage}
\end{minipage}

\vfill
\begin{tcolorbox}[enhanced,nobeforeafter,width=7.1cm,height=2.1cm,valign=center,colback=cgra!6,colframe=cgra,arc=2.5mm,boxrule=0.9pt,left=3mm]{\Large\bfseries\color{cgra}a)}\quad\begin{minipage}[c]{5.3cm}\raggedright\Large $300~\text{km/h}$\end{minipage}\end{tcolorbox}\hfill\begin{tcolorbox}[enhanced,nobeforeafter,width=7.1cm,height=2.1cm,valign=center,colback=cgra!6,colframe=cgra,arc=2.5mm,boxrule=0.9pt,left=3mm]{\Large\bfseries\color{cgra}b)}\quad\begin{minipage}[c]{5.3cm}\raggedright\Large $1\,800~\text{km/h}$\end{minipage}\end{tcolorbox}\par\vspace{2.5mm}
\begin{tcolorbox}[enhanced,nobeforeafter,width=7.1cm,height=2.1cm,valign=center,colback=cgra!6,colframe=cgra,arc=2.5mm,boxrule=0.9pt,left=3mm]{\Large\bfseries\color{cgra}c)}\quad\begin{minipage}[c]{5.3cm}\raggedright\Large $30\,000~\text{km/h}$\end{minipage}\end{tcolorbox}\hfill\begin{tcolorbox}[enhanced,nobeforeafter,width=7.1cm,height=2.1cm,valign=center,colback=cgra!6,colframe=cgra,arc=2.5mm,boxrule=0.9pt,left=3mm]{\Large\bfseries\color{cgra}d)}\quad\begin{minipage}[c]{5.3cm}\raggedright\Large $108~\text{km/h}$\end{minipage}\end{tcolorbox}
\end{tcolorbox}
\newpage
\begin{tcolorbox}[enhanced,width=16cm,height=14cm,colback=white,colframe=cgra,boxrule=2pt,arc=0mm,
left=6mm,right=6mm,top=0mm,bottom=4mm,sharp corners,
overlay={\fill[cgra] (frame.north west) rectangle ([yshift=-1.25cm]frame.north east);
\node[anchor=west,text=white,font=\sffamily\bfseries\LARGE] at ([xshift=5mm,yshift=-0.62cm]frame.north west) {Question 83 C};
\node[anchor=east,text=white,font=\sffamily\itshape\large] at ([xshift=-5mm,yshift=-0.62cm]frame.north east) {Grandeurs et proportionnalité};}]
\vspace*{1.75cm}
\begin{minipage}[t][6.9cm][c]{\linewidth}\LARGE\centering
\begin{minipage}{0.95\linewidth}\raggedright\linespread{1.25}\selectfont Un véhicule parcourt $25$ km en 15 min. Quelle est sa vitesse moyenne~?\end{minipage}
\end{minipage}

\vfill
\begin{tcolorbox}[enhanced,nobeforeafter,width=7.1cm,height=2.1cm,valign=center,colback=cgra!6,colframe=cgra,arc=2.5mm,boxrule=0.9pt,left=3mm]{\Large\bfseries\color{cgra}a)}\quad\begin{minipage}[c]{5.3cm}\raggedright\Large $50~\text{km/h}$\end{minipage}\end{tcolorbox}\hfill\begin{tcolorbox}[enhanced,nobeforeafter,width=7.1cm,height=2.1cm,valign=center,colback=cgra!6,colframe=cgra,arc=2.5mm,boxrule=0.9pt,left=3mm]{\Large\bfseries\color{cgra}b)}\quad\begin{minipage}[c]{5.3cm}\raggedright\Large $6{,}25~\text{km/h}$\end{minipage}\end{tcolorbox}\par\vspace{2.5mm}
\begin{tcolorbox}[enhanced,nobeforeafter,width=7.1cm,height=2.1cm,valign=center,colback=cgra!6,colframe=cgra,arc=2.5mm,boxrule=0.9pt,left=3mm]{\Large\bfseries\color{cgra}c)}\quad\begin{minipage}[c]{5.3cm}\raggedright\Large $100~\text{km/h}$\end{minipage}\end{tcolorbox}\hfill\begin{tcolorbox}[enhanced,nobeforeafter,width=7.1cm,height=2.1cm,valign=center,colback=cgra!6,colframe=cgra,arc=2.5mm,boxrule=0.9pt,left=3mm]{\Large\bfseries\color{cgra}d)}\quad\begin{minipage}[c]{5.3cm}\raggedright\Large $110~\text{km/h}$\end{minipage}\end{tcolorbox}
\end{tcolorbox}
\newpage
\begin{tcolorbox}[enhanced,width=16cm,height=14cm,colback=white,colframe=cgra,boxrule=2pt,arc=0mm,
left=6mm,right=6mm,top=0mm,bottom=4mm,sharp corners,
overlay={\fill[cgra] (frame.north west) rectangle ([yshift=-1.25cm]frame.north east);
\node[anchor=west,text=white,font=\sffamily\bfseries\LARGE] at ([xshift=5mm,yshift=-0.62cm]frame.north west) {Question 84 C};
\node[anchor=east,text=white,font=\sffamily\itshape\large] at ([xshift=-5mm,yshift=-0.62cm]frame.north east) {Grandeurs et proportionnalité};}]
\vspace*{1.75cm}
\begin{minipage}[t][6.9cm][c]{\linewidth}\LARGE\centering
\begin{minipage}{0.95\linewidth}\raggedright\linespread{1.25}\selectfont Un véhicule parcourt $330$ km en 3 h. Quelle est sa vitesse moyenne~?\end{minipage}
\end{minipage}

\vfill
\begin{tcolorbox}[enhanced,nobeforeafter,width=7.1cm,height=2.1cm,valign=center,colback=cgra!6,colframe=cgra,arc=2.5mm,boxrule=0.9pt,left=3mm]{\Large\bfseries\color{cgra}a)}\quad\begin{minipage}[c]{5.3cm}\raggedright\Large $990~\text{km/h}$\end{minipage}\end{tcolorbox}\hfill\begin{tcolorbox}[enhanced,nobeforeafter,width=7.1cm,height=2.1cm,valign=center,colback=cgra!6,colframe=cgra,arc=2.5mm,boxrule=0.9pt,left=3mm]{\Large\bfseries\color{cgra}b)}\quad\begin{minipage}[c]{5.3cm}\raggedright\Large $110~\text{km/h}$\end{minipage}\end{tcolorbox}\par\vspace{2.5mm}
\begin{tcolorbox}[enhanced,nobeforeafter,width=7.1cm,height=2.1cm,valign=center,colback=cgra!6,colframe=cgra,arc=2.5mm,boxrule=0.9pt,left=3mm]{\Large\bfseries\color{cgra}c)}\quad\begin{minipage}[c]{5.3cm}\raggedright\Large $55~\text{km/h}$\end{minipage}\end{tcolorbox}\hfill\begin{tcolorbox}[enhanced,nobeforeafter,width=7.1cm,height=2.1cm,valign=center,colback=cgra!6,colframe=cgra,arc=2.5mm,boxrule=0.9pt,left=3mm]{\Large\bfseries\color{cgra}d)}\quad\begin{minipage}[c]{5.3cm}\raggedright\Large $120~\text{km/h}$\end{minipage}\end{tcolorbox}
\end{tcolorbox}
\newpage
\begin{tcolorbox}[enhanced,width=16cm,height=14cm,colback=white,colframe=cgra,boxrule=2pt,arc=0mm,
left=6mm,right=6mm,top=0mm,bottom=4mm,sharp corners,
overlay={\fill[cgra] (frame.north west) rectangle ([yshift=-1.25cm]frame.north east);
\node[anchor=west,text=white,font=\sffamily\bfseries\LARGE] at ([xshift=5mm,yshift=-0.62cm]frame.north west) {Question 85 C};
\node[anchor=east,text=white,font=\sffamily\itshape\large] at ([xshift=-5mm,yshift=-0.62cm]frame.north east) {Grandeurs et proportionnalité};}]
\vspace*{1.75cm}
\begin{minipage}[t][6.9cm][c]{\linewidth}\LARGE\centering
\begin{minipage}{0.95\linewidth}\raggedright\linespread{1.25}\selectfont Un article coûte $400$ F. Son prix augmente de $35\,\%$. Quel est son nouveau prix~?\end{minipage}
\end{minipage}

\vfill
\begin{tcolorbox}[enhanced,nobeforeafter,width=7.1cm,height=2.1cm,valign=center,colback=cgra!6,colframe=cgra,arc=2.5mm,boxrule=0.9pt,left=3mm]{\Large\bfseries\color{cgra}a)}\quad\begin{minipage}[c]{5.3cm}\raggedright\Large $140~\text{F}$\end{minipage}\end{tcolorbox}\hfill\begin{tcolorbox}[enhanced,nobeforeafter,width=7.1cm,height=2.1cm,valign=center,colback=cgra!6,colframe=cgra,arc=2.5mm,boxrule=0.9pt,left=3mm]{\Large\bfseries\color{cgra}b)}\quad\begin{minipage}[c]{5.3cm}\raggedright\Large $540~\text{F}$\end{minipage}\end{tcolorbox}\par\vspace{2.5mm}
\begin{tcolorbox}[enhanced,nobeforeafter,width=7.1cm,height=2.1cm,valign=center,colback=cgra!6,colframe=cgra,arc=2.5mm,boxrule=0.9pt,left=3mm]{\Large\bfseries\color{cgra}c)}\quad\begin{minipage}[c]{5.3cm}\raggedright\Large $260~\text{F}$\end{minipage}\end{tcolorbox}\hfill\begin{tcolorbox}[enhanced,nobeforeafter,width=7.1cm,height=2.1cm,valign=center,colback=cgra!6,colframe=cgra,arc=2.5mm,boxrule=0.9pt,left=3mm]{\Large\bfseries\color{cgra}d)}\quad\begin{minipage}[c]{5.3cm}\raggedright\Large $435~\text{F}$\end{minipage}\end{tcolorbox}
\end{tcolorbox}
\newpage
\begin{tcolorbox}[enhanced,width=16cm,height=14cm,colback=white,colframe=cgra,boxrule=2pt,arc=0mm,
left=6mm,right=6mm,top=0mm,bottom=4mm,sharp corners,
overlay={\fill[cgra] (frame.north west) rectangle ([yshift=-1.25cm]frame.north east);
\node[anchor=west,text=white,font=\sffamily\bfseries\LARGE] at ([xshift=5mm,yshift=-0.62cm]frame.north west) {Question 86 C};
\node[anchor=east,text=white,font=\sffamily\itshape\large] at ([xshift=-5mm,yshift=-0.62cm]frame.north east) {Grandeurs et proportionnalité};}]
\vspace*{1.75cm}
\begin{minipage}[t][6.9cm][c]{\linewidth}\LARGE\centering
\begin{minipage}{0.95\linewidth}\raggedright\linespread{1.25}\selectfont Un article coûte $2\,200$ F. Son prix baisse de $15\,\%$. Quel est son nouveau prix~?\end{minipage}
\end{minipage}

\vfill
\begin{tcolorbox}[enhanced,nobeforeafter,width=7.1cm,height=2.1cm,valign=center,colback=cgra!6,colframe=cgra,arc=2.5mm,boxrule=0.9pt,left=3mm]{\Large\bfseries\color{cgra}a)}\quad\begin{minipage}[c]{5.3cm}\raggedright\Large $330~\text{F}$\end{minipage}\end{tcolorbox}\hfill\begin{tcolorbox}[enhanced,nobeforeafter,width=7.1cm,height=2.1cm,valign=center,colback=cgra!6,colframe=cgra,arc=2.5mm,boxrule=0.9pt,left=3mm]{\Large\bfseries\color{cgra}b)}\quad\begin{minipage}[c]{5.3cm}\raggedright\Large $1\,870~\text{F}$\end{minipage}\end{tcolorbox}\par\vspace{2.5mm}
\begin{tcolorbox}[enhanced,nobeforeafter,width=7.1cm,height=2.1cm,valign=center,colback=cgra!6,colframe=cgra,arc=2.5mm,boxrule=0.9pt,left=3mm]{\Large\bfseries\color{cgra}c)}\quad\begin{minipage}[c]{5.3cm}\raggedright\Large $2\,185~\text{F}$\end{minipage}\end{tcolorbox}\hfill\begin{tcolorbox}[enhanced,nobeforeafter,width=7.1cm,height=2.1cm,valign=center,colback=cgra!6,colframe=cgra,arc=2.5mm,boxrule=0.9pt,left=3mm]{\Large\bfseries\color{cgra}d)}\quad\begin{minipage}[c]{5.3cm}\raggedright\Large $2\,530~\text{F}$\end{minipage}\end{tcolorbox}
\end{tcolorbox}
\newpage
\begin{tcolorbox}[enhanced,width=16cm,height=14cm,colback=white,colframe=cgra,boxrule=2pt,arc=0mm,
left=6mm,right=6mm,top=0mm,bottom=4mm,sharp corners,
overlay={\fill[cgra] (frame.north west) rectangle ([yshift=-1.25cm]frame.north east);
\node[anchor=west,text=white,font=\sffamily\bfseries\LARGE] at ([xshift=5mm,yshift=-0.62cm]frame.north west) {Question 87 C};
\node[anchor=east,text=white,font=\sffamily\itshape\large] at ([xshift=-5mm,yshift=-0.62cm]frame.north east) {Grandeurs et proportionnalité};}]
\vspace*{1.75cm}
\begin{minipage}[t][6.9cm][c]{\linewidth}\LARGE\centering
\begin{minipage}{0.95\linewidth}\raggedright\linespread{1.25}\selectfont Un prix passe de $500$ F à $575$ F. Quel est le pourcentage d'évolution~?\end{minipage}
\end{minipage}

\vfill
\begin{tcolorbox}[enhanced,nobeforeafter,width=7.1cm,height=2.1cm,valign=center,colback=cgra!6,colframe=cgra,arc=2.5mm,boxrule=0.9pt,left=3mm]{\Large\bfseries\color{cgra}a)}\quad\begin{minipage}[c]{5.3cm}\raggedright\Large $+15\,\%$\end{minipage}\end{tcolorbox}\hfill\begin{tcolorbox}[enhanced,nobeforeafter,width=7.1cm,height=2.1cm,valign=center,colback=cgra!6,colframe=cgra,arc=2.5mm,boxrule=0.9pt,left=3mm]{\Large\bfseries\color{cgra}b)}\quad\begin{minipage}[c]{5.3cm}\raggedright\Large $-15\,\%$\end{minipage}\end{tcolorbox}\par\vspace{2.5mm}
\begin{tcolorbox}[enhanced,nobeforeafter,width=7.1cm,height=2.1cm,valign=center,colback=cgra!6,colframe=cgra,arc=2.5mm,boxrule=0.9pt,left=3mm]{\Large\bfseries\color{cgra}c)}\quad\begin{minipage}[c]{5.3cm}\raggedright\Large $+7{,}5\,\%$\end{minipage}\end{tcolorbox}\hfill\begin{tcolorbox}[enhanced,nobeforeafter,width=7.1cm,height=2.1cm,valign=center,colback=cgra!6,colframe=cgra,arc=2.5mm,boxrule=0.9pt,left=3mm]{\Large\bfseries\color{cgra}d)}\quad\begin{minipage}[c]{5.3cm}\raggedright\Large $+75\,\%$\end{minipage}\end{tcolorbox}
\end{tcolorbox}
\newpage
\begin{tcolorbox}[enhanced,width=16cm,height=14cm,colback=white,colframe=cgra,boxrule=2pt,arc=0mm,
left=6mm,right=6mm,top=0mm,bottom=4mm,sharp corners,
overlay={\fill[cgra] (frame.north west) rectangle ([yshift=-1.25cm]frame.north east);
\node[anchor=west,text=white,font=\sffamily\bfseries\LARGE] at ([xshift=5mm,yshift=-0.62cm]frame.north west) {Question 88 C};
\node[anchor=east,text=white,font=\sffamily\itshape\large] at ([xshift=-5mm,yshift=-0.62cm]frame.north east) {Grandeurs et proportionnalité};}]
\vspace*{1.75cm}
\begin{minipage}[t][6.9cm][c]{\linewidth}\LARGE\centering
\begin{minipage}{0.95\linewidth}\raggedright\linespread{1.25}\selectfont Un prix passe de $60$ F à $51$ F. Quel est le pourcentage d'évolution~?\end{minipage}
\end{minipage}

\vfill
\begin{tcolorbox}[enhanced,nobeforeafter,width=7.1cm,height=2.1cm,valign=center,colback=cgra!6,colframe=cgra,arc=2.5mm,boxrule=0.9pt,left=3mm]{\Large\bfseries\color{cgra}a)}\quad\begin{minipage}[c]{5.3cm}\raggedright\Large $-7{,}5\,\%$\end{minipage}\end{tcolorbox}\hfill\begin{tcolorbox}[enhanced,nobeforeafter,width=7.1cm,height=2.1cm,valign=center,colback=cgra!6,colframe=cgra,arc=2.5mm,boxrule=0.9pt,left=3mm]{\Large\bfseries\color{cgra}b)}\quad\begin{minipage}[c]{5.3cm}\raggedright\Large $+15\,\%$\end{minipage}\end{tcolorbox}\par\vspace{2.5mm}
\begin{tcolorbox}[enhanced,nobeforeafter,width=7.1cm,height=2.1cm,valign=center,colback=cgra!6,colframe=cgra,arc=2.5mm,boxrule=0.9pt,left=3mm]{\Large\bfseries\color{cgra}c)}\quad\begin{minipage}[c]{5.3cm}\raggedright\Large $-9\,\%$\end{minipage}\end{tcolorbox}\hfill\begin{tcolorbox}[enhanced,nobeforeafter,width=7.1cm,height=2.1cm,valign=center,colback=cgra!6,colframe=cgra,arc=2.5mm,boxrule=0.9pt,left=3mm]{\Large\bfseries\color{cgra}d)}\quad\begin{minipage}[c]{5.3cm}\raggedright\Large $-15\,\%$\end{minipage}\end{tcolorbox}
\end{tcolorbox}
\newpage
\begin{tcolorbox}[enhanced,width=16cm,height=14cm,colback=white,colframe=cgra,boxrule=2pt,arc=0mm,
left=6mm,right=6mm,top=0mm,bottom=4mm,sharp corners,
overlay={\fill[cgra] (frame.north west) rectangle ([yshift=-1.25cm]frame.north east);
\node[anchor=west,text=white,font=\sffamily\bfseries\LARGE] at ([xshift=5mm,yshift=-0.62cm]frame.north west) {Question 89 C};
\node[anchor=east,text=white,font=\sffamily\itshape\large] at ([xshift=-5mm,yshift=-0.62cm]frame.north east) {Grandeurs et proportionnalité};}]
\vspace*{1.75cm}
\begin{minipage}[t][6.9cm][c]{\linewidth}\LARGE\centering
\begin{minipage}{0.95\linewidth}\raggedright\linespread{1.25}\selectfont Un robinet remplit une cuve de $300$ L en $12$ min. Quel est son débit~?\end{minipage}
\end{minipage}

\vfill
\begin{tcolorbox}[enhanced,nobeforeafter,width=7.1cm,height=2.1cm,valign=center,colback=cgra!6,colframe=cgra,arc=2.5mm,boxrule=0.9pt,left=3mm]{\Large\bfseries\color{cgra}a)}\quad\begin{minipage}[c]{5.3cm}\raggedright\Large $25~\text{L/min}$\end{minipage}\end{tcolorbox}\hfill\begin{tcolorbox}[enhanced,nobeforeafter,width=7.1cm,height=2.1cm,valign=center,colback=cgra!6,colframe=cgra,arc=2.5mm,boxrule=0.9pt,left=3mm]{\Large\bfseries\color{cgra}b)}\quad\begin{minipage}[c]{5.3cm}\raggedright\Large $4~\text{L/min}$\end{minipage}\end{tcolorbox}\par\vspace{2.5mm}
\begin{tcolorbox}[enhanced,nobeforeafter,width=7.1cm,height=2.1cm,valign=center,colback=cgra!6,colframe=cgra,arc=2.5mm,boxrule=0.9pt,left=3mm]{\Large\bfseries\color{cgra}c)}\quad\begin{minipage}[c]{5.3cm}\raggedright\Large $27~\text{L/min}$\end{minipage}\end{tcolorbox}\hfill\begin{tcolorbox}[enhanced,nobeforeafter,width=7.1cm,height=2.1cm,valign=center,colback=cgra!6,colframe=cgra,arc=2.5mm,boxrule=0.9pt,left=3mm]{\Large\bfseries\color{cgra}d)}\quad\begin{minipage}[c]{5.3cm}\raggedright\Large $3\,600~\text{L/min}$\end{minipage}\end{tcolorbox}
\end{tcolorbox}
\newpage
\begin{tcolorbox}[enhanced,width=16cm,height=14cm,colback=white,colframe=cgra,boxrule=2pt,arc=0mm,
left=6mm,right=6mm,top=0mm,bottom=4mm,sharp corners,
overlay={\fill[cgra] (frame.north west) rectangle ([yshift=-1.25cm]frame.north east);
\node[anchor=west,text=white,font=\sffamily\bfseries\LARGE] at ([xshift=5mm,yshift=-0.62cm]frame.north west) {Question 90 C};
\node[anchor=east,text=white,font=\sffamily\itshape\large] at ([xshift=-5mm,yshift=-0.62cm]frame.north east) {Grandeurs et proportionnalité};}]
\vspace*{1.75cm}
\begin{minipage}[t][6.9cm][c]{\linewidth}\LARGE\centering
\begin{minipage}{0.95\linewidth}\raggedright\linespread{1.25}\selectfont Un robinet remplit une cuve de $45$ L en $2$ min. Quel est son débit~?\end{minipage}
\end{minipage}

\vfill
\begin{tcolorbox}[enhanced,nobeforeafter,width=7.1cm,height=2.1cm,valign=center,colback=cgra!6,colframe=cgra,arc=2.5mm,boxrule=0.9pt,left=3mm]{\Large\bfseries\color{cgra}a)}\quad\begin{minipage}[c]{5.3cm}\raggedright\Large $24{,}5~\text{L/min}$\end{minipage}\end{tcolorbox}\hfill\begin{tcolorbox}[enhanced,nobeforeafter,width=7.1cm,height=2.1cm,valign=center,colback=cgra!6,colframe=cgra,arc=2.5mm,boxrule=0.9pt,left=3mm]{\Large\bfseries\color{cgra}b)}\quad\begin{minipage}[c]{5.3cm}\raggedright\Large $90~\text{L/min}$\end{minipage}\end{tcolorbox}\par\vspace{2.5mm}
\begin{tcolorbox}[enhanced,nobeforeafter,width=7.1cm,height=2.1cm,valign=center,colback=cgra!6,colframe=cgra,arc=2.5mm,boxrule=0.9pt,left=3mm]{\Large\bfseries\color{cgra}c)}\quad\begin{minipage}[c]{5.3cm}\raggedright\Large $22{,}5~\text{L/min}$\end{minipage}\end{tcolorbox}\hfill\begin{tcolorbox}[enhanced,nobeforeafter,width=7.1cm,height=2.1cm,valign=center,colback=cgra!6,colframe=cgra,arc=2.5mm,boxrule=0.9pt,left=3mm]{\Large\bfseries\color{cgra}d)}\quad\begin{minipage}[c]{5.3cm}\raggedright\large $\dfrac{40}{9}~\text{L/min}$\end{minipage}\end{tcolorbox}
\end{tcolorbox}
\newpage
\begin{tcolorbox}[enhanced,width=16cm,height=14cm,colback=white,colframe=cgra,boxrule=2pt,arc=0mm,
left=6mm,right=6mm,top=0mm,bottom=4mm,sharp corners,
overlay={\fill[cgra] (frame.north west) rectangle ([yshift=-1.25cm]frame.north east);
\node[anchor=west,text=white,font=\sffamily\bfseries\LARGE] at ([xshift=5mm,yshift=-0.62cm]frame.north west) {Question 91 C};
\node[anchor=east,text=white,font=\sffamily\itshape\large] at ([xshift=-5mm,yshift=-0.62cm]frame.north east) {Grandeurs et proportionnalité};}]
\vspace*{1.75cm}
\begin{minipage}[t][6.9cm][c]{\linewidth}\LARGE\centering
\begin{minipage}{0.95\linewidth}\raggedright\linespread{1.25}\selectfont Sur une carte à l'échelle $\dfrac{1}{100\,000}$, deux points sont distants de $3$ cm. Quelle est la distance réelle~?\end{minipage}
\end{minipage}

\vfill
\begin{tcolorbox}[enhanced,nobeforeafter,width=7.1cm,height=2.1cm,valign=center,colback=cgra!6,colframe=cgra,arc=2.5mm,boxrule=0.9pt,left=3mm]{\Large\bfseries\color{cgra}a)}\quad\begin{minipage}[c]{5.3cm}\raggedright\Large $3\,000~\text{m}$\end{minipage}\end{tcolorbox}\hfill\begin{tcolorbox}[enhanced,nobeforeafter,width=7.1cm,height=2.1cm,valign=center,colback=cgra!6,colframe=cgra,arc=2.5mm,boxrule=0.9pt,left=3mm]{\Large\bfseries\color{cgra}b)}\quad\begin{minipage}[c]{5.3cm}\raggedright\Large $300\,000~\text{m}$\end{minipage}\end{tcolorbox}\par\vspace{2.5mm}
\begin{tcolorbox}[enhanced,nobeforeafter,width=7.1cm,height=2.1cm,valign=center,colback=cgra!6,colframe=cgra,arc=2.5mm,boxrule=0.9pt,left=3mm]{\Large\bfseries\color{cgra}c)}\quad\begin{minipage}[c]{5.3cm}\raggedright\Large $300~\text{m}$\end{minipage}\end{tcolorbox}\hfill\begin{tcolorbox}[enhanced,nobeforeafter,width=7.1cm,height=2.1cm,valign=center,colback=cgra!6,colframe=cgra,arc=2.5mm,boxrule=0.9pt,left=3mm]{\Large\bfseries\color{cgra}d)}\quad\begin{minipage}[c]{5.3cm}\raggedright\Large $30\,000~\text{m}$\end{minipage}\end{tcolorbox}
\end{tcolorbox}
\newpage
\begin{tcolorbox}[enhanced,width=16cm,height=14cm,colback=white,colframe=cgra,boxrule=2pt,arc=0mm,
left=6mm,right=6mm,top=0mm,bottom=4mm,sharp corners,
overlay={\fill[cgra] (frame.north west) rectangle ([yshift=-1.25cm]frame.north east);
\node[anchor=west,text=white,font=\sffamily\bfseries\LARGE] at ([xshift=5mm,yshift=-0.62cm]frame.north west) {Question 92 C};
\node[anchor=east,text=white,font=\sffamily\itshape\large] at ([xshift=-5mm,yshift=-0.62cm]frame.north east) {Grandeurs et proportionnalité};}]
\vspace*{1.75cm}
\begin{minipage}[t][6.9cm][c]{\linewidth}\LARGE\centering
\begin{minipage}{0.95\linewidth}\raggedright\linespread{1.25}\selectfont Sur une carte à l'échelle $\dfrac{1}{50}$, deux points sont distants de $12$ cm. Quelle est la distance réelle~?\end{minipage}
\end{minipage}

\vfill
\begin{tcolorbox}[enhanced,nobeforeafter,width=7.1cm,height=2.1cm,valign=center,colback=cgra!6,colframe=cgra,arc=2.5mm,boxrule=0.9pt,left=3mm]{\Large\bfseries\color{cgra}a)}\quad\begin{minipage}[c]{5.3cm}\raggedright\Large $600~\text{m}$\end{minipage}\end{tcolorbox}\hfill\begin{tcolorbox}[enhanced,nobeforeafter,width=7.1cm,height=2.1cm,valign=center,colback=cgra!6,colframe=cgra,arc=2.5mm,boxrule=0.9pt,left=3mm]{\Large\bfseries\color{cgra}b)}\quad\begin{minipage}[c]{5.3cm}\raggedright\Large $60~\text{m}$\end{minipage}\end{tcolorbox}\par\vspace{2.5mm}
\begin{tcolorbox}[enhanced,nobeforeafter,width=7.1cm,height=2.1cm,valign=center,colback=cgra!6,colframe=cgra,arc=2.5mm,boxrule=0.9pt,left=3mm]{\Large\bfseries\color{cgra}c)}\quad\begin{minipage}[c]{5.3cm}\raggedright\Large $0{,}6~\text{m}$\end{minipage}\end{tcolorbox}\hfill\begin{tcolorbox}[enhanced,nobeforeafter,width=7.1cm,height=2.1cm,valign=center,colback=cgra!6,colframe=cgra,arc=2.5mm,boxrule=0.9pt,left=3mm]{\Large\bfseries\color{cgra}d)}\quad\begin{minipage}[c]{5.3cm}\raggedright\Large $6~\text{m}$\end{minipage}\end{tcolorbox}
\end{tcolorbox}
\newpage
\begin{tcolorbox}[enhanced,width=16cm,height=14cm,colback=white,colframe=cgra,boxrule=2pt,arc=0mm,
left=6mm,right=6mm,top=0mm,bottom=4mm,sharp corners,
overlay={\fill[cgra] (frame.north west) rectangle ([yshift=-1.25cm]frame.north east);
\node[anchor=west,text=white,font=\sffamily\bfseries\LARGE] at ([xshift=5mm,yshift=-0.62cm]frame.north west) {Question 93 C};
\node[anchor=east,text=white,font=\sffamily\itshape\large] at ([xshift=-5mm,yshift=-0.62cm]frame.north east) {Grandeurs et proportionnalité};}]
\vspace*{1.75cm}
\begin{minipage}[t][6.9cm][c]{\linewidth}\LARGE\centering
\begin{minipage}{0.95\linewidth}\raggedright\linespread{1.25}\selectfont On partage $90$ bonbons entre Léa et Noa dans le ratio $4:5$. Combien chacun en reçoit-il (Léa, puis Noa)~?\end{minipage}
\end{minipage}

\vfill
\begin{tcolorbox}[enhanced,nobeforeafter,width=7.1cm,height=2.1cm,valign=center,colback=cgra!6,colframe=cgra,arc=2.5mm,boxrule=0.9pt,left=3mm]{\Large\bfseries\color{cgra}a)}\quad\begin{minipage}[c]{5.3cm}\raggedright\Large $40\text{ et }50$\end{minipage}\end{tcolorbox}\hfill\begin{tcolorbox}[enhanced,nobeforeafter,width=7.1cm,height=2.1cm,valign=center,colback=cgra!6,colframe=cgra,arc=2.5mm,boxrule=0.9pt,left=3mm]{\Large\bfseries\color{cgra}b)}\quad\begin{minipage}[c]{5.3cm}\raggedright\Large $45\text{ et }45$\end{minipage}\end{tcolorbox}\par\vspace{2.5mm}
\begin{tcolorbox}[enhanced,nobeforeafter,width=7.1cm,height=2.1cm,valign=center,colback=cgra!6,colframe=cgra,arc=2.5mm,boxrule=0.9pt,left=3mm]{\Large\bfseries\color{cgra}c)}\quad\begin{minipage}[c]{5.3cm}\raggedright\Large $36\text{ et }54$\end{minipage}\end{tcolorbox}\hfill\begin{tcolorbox}[enhanced,nobeforeafter,width=7.1cm,height=2.1cm,valign=center,colback=cgra!6,colframe=cgra,arc=2.5mm,boxrule=0.9pt,left=3mm]{\Large\bfseries\color{cgra}d)}\quad\begin{minipage}[c]{5.3cm}\raggedright\Large $50\text{ et }40$\end{minipage}\end{tcolorbox}
\end{tcolorbox}
\newpage
\begin{tcolorbox}[enhanced,width=16cm,height=14cm,colback=white,colframe=cgra,boxrule=2pt,arc=0mm,
left=6mm,right=6mm,top=0mm,bottom=4mm,sharp corners,
overlay={\fill[cgra] (frame.north west) rectangle ([yshift=-1.25cm]frame.north east);
\node[anchor=west,text=white,font=\sffamily\bfseries\LARGE] at ([xshift=5mm,yshift=-0.62cm]frame.north west) {Question 94 C};
\node[anchor=east,text=white,font=\sffamily\itshape\large] at ([xshift=-5mm,yshift=-0.62cm]frame.north east) {Grandeurs et proportionnalité};}]
\vspace*{1.75cm}
\begin{minipage}[t][6.9cm][c]{\linewidth}\LARGE\centering
\begin{minipage}{0.95\linewidth}\raggedright\linespread{1.25}\selectfont On partage $56$ bonbons entre Léa et Noa dans le ratio $5:2$. Combien chacun en reçoit-il (Léa, puis Noa)~?\end{minipage}
\end{minipage}

\vfill
\begin{tcolorbox}[enhanced,nobeforeafter,width=7.1cm,height=2.1cm,valign=center,colback=cgra!6,colframe=cgra,arc=2.5mm,boxrule=0.9pt,left=3mm]{\Large\bfseries\color{cgra}a)}\quad\begin{minipage}[c]{5.3cm}\raggedright\Large $40\text{ et }16$\end{minipage}\end{tcolorbox}\hfill\begin{tcolorbox}[enhanced,nobeforeafter,width=7.1cm,height=2.1cm,valign=center,colback=cgra!6,colframe=cgra,arc=2.5mm,boxrule=0.9pt,left=3mm]{\Large\bfseries\color{cgra}b)}\quad\begin{minipage}[c]{5.3cm}\raggedright\Large $28\text{ et }28$\end{minipage}\end{tcolorbox}\par\vspace{2.5mm}
\begin{tcolorbox}[enhanced,nobeforeafter,width=7.1cm,height=2.1cm,valign=center,colback=cgra!6,colframe=cgra,arc=2.5mm,boxrule=0.9pt,left=3mm]{\Large\bfseries\color{cgra}c)}\quad\begin{minipage}[c]{5.3cm}\raggedright\Large $35\text{ et }21$\end{minipage}\end{tcolorbox}\hfill\begin{tcolorbox}[enhanced,nobeforeafter,width=7.1cm,height=2.1cm,valign=center,colback=cgra!6,colframe=cgra,arc=2.5mm,boxrule=0.9pt,left=3mm]{\Large\bfseries\color{cgra}d)}\quad\begin{minipage}[c]{5.3cm}\raggedright\Large $16\text{ et }40$\end{minipage}\end{tcolorbox}
\end{tcolorbox}
\newpage
\begin{tcolorbox}[enhanced,width=16cm,height=14cm,colback=white,colframe=cgra,boxrule=2pt,arc=0mm,
left=6mm,right=6mm,top=0mm,bottom=4mm,sharp corners,
overlay={\fill[cgra] (frame.north west) rectangle ([yshift=-1.25cm]frame.north east);
\node[anchor=west,text=white,font=\sffamily\bfseries\LARGE] at ([xshift=5mm,yshift=-0.62cm]frame.north west) {Question 95 C};
\node[anchor=east,text=white,font=\sffamily\itshape\large] at ([xshift=-5mm,yshift=-0.62cm]frame.north east) {Grandeurs et proportionnalité};}]
\vspace*{1.75cm}
\begin{minipage}[t][6.9cm][c]{\linewidth}\LARGE\centering
\begin{minipage}{0.95\linewidth}\raggedright\linespread{1.25}\selectfont Compléter~: $2~\text{dm}^3=\;?\;\text{cm}^3$\end{minipage}
\end{minipage}

\vfill
\begin{tcolorbox}[enhanced,nobeforeafter,width=7.1cm,height=2.1cm,valign=center,colback=cgra!6,colframe=cgra,arc=2.5mm,boxrule=0.9pt,left=3mm]{\Large\bfseries\color{cgra}a)}\quad\begin{minipage}[c]{5.3cm}\raggedright\Large $20\,000~\text{cm}^3$\end{minipage}\end{tcolorbox}\hfill\begin{tcolorbox}[enhanced,nobeforeafter,width=7.1cm,height=2.1cm,valign=center,colback=cgra!6,colframe=cgra,arc=2.5mm,boxrule=0.9pt,left=3mm]{\Large\bfseries\color{cgra}b)}\quad\begin{minipage}[c]{5.3cm}\raggedright\Large $2\,000~\text{cm}^3$\end{minipage}\end{tcolorbox}\par\vspace{2.5mm}
\begin{tcolorbox}[enhanced,nobeforeafter,width=7.1cm,height=2.1cm,valign=center,colback=cgra!6,colframe=cgra,arc=2.5mm,boxrule=0.9pt,left=3mm]{\Large\bfseries\color{cgra}c)}\quad\begin{minipage}[c]{5.3cm}\raggedright\Large $200~\text{cm}^3$\end{minipage}\end{tcolorbox}\hfill\begin{tcolorbox}[enhanced,nobeforeafter,width=7.1cm,height=2.1cm,valign=center,colback=cgra!6,colframe=cgra,arc=2.5mm,boxrule=0.9pt,left=3mm]{\Large\bfseries\color{cgra}d)}\quad\begin{minipage}[c]{5.3cm}\raggedright\Large $200\,000~\text{cm}^3$\end{minipage}\end{tcolorbox}
\end{tcolorbox}
\newpage
\begin{tcolorbox}[enhanced,width=16cm,height=14cm,colback=white,colframe=cgra,boxrule=2pt,arc=0mm,
left=6mm,right=6mm,top=0mm,bottom=4mm,sharp corners,
overlay={\fill[cgra] (frame.north west) rectangle ([yshift=-1.25cm]frame.north east);
\node[anchor=west,text=white,font=\sffamily\bfseries\LARGE] at ([xshift=5mm,yshift=-0.62cm]frame.north west) {Question 96 C};
\node[anchor=east,text=white,font=\sffamily\itshape\large] at ([xshift=-5mm,yshift=-0.62cm]frame.north east) {Grandeurs et proportionnalité};}]
\vspace*{1.75cm}
\begin{minipage}[t][6.9cm][c]{\linewidth}\LARGE\centering
\begin{minipage}{0.95\linewidth}\raggedright\linespread{1.25}\selectfont Compléter~: $125~\text{cm}^3=\;?\;\text{dm}^3$\end{minipage}
\end{minipage}

\vfill
\begin{tcolorbox}[enhanced,nobeforeafter,width=7.1cm,height=2.1cm,valign=center,colback=cgra!6,colframe=cgra,arc=2.5mm,boxrule=0.9pt,left=3mm]{\Large\bfseries\color{cgra}a)}\quad\begin{minipage}[c]{5.3cm}\raggedright\Large $0{,}125~\text{dm}^3$\end{minipage}\end{tcolorbox}\hfill\begin{tcolorbox}[enhanced,nobeforeafter,width=7.1cm,height=2.1cm,valign=center,colback=cgra!6,colframe=cgra,arc=2.5mm,boxrule=0.9pt,left=3mm]{\Large\bfseries\color{cgra}b)}\quad\begin{minipage}[c]{5.3cm}\raggedright\Large $12{,}5~\text{dm}^3$\end{minipage}\end{tcolorbox}\par\vspace{2.5mm}
\begin{tcolorbox}[enhanced,nobeforeafter,width=7.1cm,height=2.1cm,valign=center,colback=cgra!6,colframe=cgra,arc=2.5mm,boxrule=0.9pt,left=3mm]{\Large\bfseries\color{cgra}c)}\quad\begin{minipage}[c]{5.3cm}\raggedright\Large $1{,}25~\text{dm}^3$\end{minipage}\end{tcolorbox}\hfill\begin{tcolorbox}[enhanced,nobeforeafter,width=7.1cm,height=2.1cm,valign=center,colback=cgra!6,colframe=cgra,arc=2.5mm,boxrule=0.9pt,left=3mm]{\Large\bfseries\color{cgra}d)}\quad\begin{minipage}[c]{5.3cm}\raggedright\Large $0{,}0125~\text{dm}^3$\end{minipage}\end{tcolorbox}
\end{tcolorbox}
\newpage
\begin{tcolorbox}[enhanced,width=16cm,height=14cm,colback=white,colframe=cari,boxrule=2pt,arc=0mm,
left=6mm,right=6mm,top=0mm,bottom=4mm,sharp corners,
overlay={\fill[cari] (frame.north west) rectangle ([yshift=-1.25cm]frame.north east);
\node[anchor=west,text=white,font=\sffamily\bfseries\LARGE] at ([xshift=5mm,yshift=-0.62cm]frame.north west) {Question 1 D};
\node[anchor=east,text=white,font=\sffamily\itshape\large] at ([xshift=-5mm,yshift=-0.62cm]frame.north east) {Arithmétique et puissances};}]
\vspace*{1.75cm}
\begin{minipage}[t][6.9cm][c]{\linewidth}\LARGE\centering
\begin{minipage}{0.95\linewidth}\raggedright\linespread{1.25}\selectfont Lequel de ces nombres est premier~?\end{minipage}
\end{minipage}

\vfill
\begin{tcolorbox}[enhanced,nobeforeafter,width=7.1cm,height=2.1cm,valign=center,colback=cari!6,colframe=cari,arc=2.5mm,boxrule=0.9pt,left=3mm]{\Large\bfseries\color{cari}a)}\quad\begin{minipage}[c]{5.3cm}\raggedright\Large $109$\end{minipage}\end{tcolorbox}\hfill\begin{tcolorbox}[enhanced,nobeforeafter,width=7.1cm,height=2.1cm,valign=center,colback=cari!6,colframe=cari,arc=2.5mm,boxrule=0.9pt,left=3mm]{\Large\bfseries\color{cari}b)}\quad\begin{minipage}[c]{5.3cm}\raggedright\Large $121$\end{minipage}\end{tcolorbox}\par\vspace{2.5mm}
\begin{tcolorbox}[enhanced,nobeforeafter,width=7.1cm,height=2.1cm,valign=center,colback=cari!6,colframe=cari,arc=2.5mm,boxrule=0.9pt,left=3mm]{\Large\bfseries\color{cari}c)}\quad\begin{minipage}[c]{5.3cm}\raggedright\Large $119$\end{minipage}\end{tcolorbox}\hfill\begin{tcolorbox}[enhanced,nobeforeafter,width=7.1cm,height=2.1cm,valign=center,colback=cari!6,colframe=cari,arc=2.5mm,boxrule=0.9pt,left=3mm]{\Large\bfseries\color{cari}d)}\quad\begin{minipage}[c]{5.3cm}\raggedright\Large $99$\end{minipage}\end{tcolorbox}
\end{tcolorbox}
\newpage
\begin{tcolorbox}[enhanced,width=16cm,height=14cm,colback=white,colframe=cari,boxrule=2pt,arc=0mm,
left=6mm,right=6mm,top=0mm,bottom=4mm,sharp corners,
overlay={\fill[cari] (frame.north west) rectangle ([yshift=-1.25cm]frame.north east);
\node[anchor=west,text=white,font=\sffamily\bfseries\LARGE] at ([xshift=5mm,yshift=-0.62cm]frame.north west) {Question 2 D};
\node[anchor=east,text=white,font=\sffamily\itshape\large] at ([xshift=-5mm,yshift=-0.62cm]frame.north east) {Arithmétique et puissances};}]
\vspace*{1.75cm}
\begin{minipage}[t][6.9cm][c]{\linewidth}\LARGE\centering
\begin{minipage}{0.95\linewidth}\raggedright\linespread{1.25}\selectfont Lequel de ces nombres est premier~?\end{minipage}
\end{minipage}

\vfill
\begin{tcolorbox}[enhanced,nobeforeafter,width=7.1cm,height=2.1cm,valign=center,colback=cari!6,colframe=cari,arc=2.5mm,boxrule=0.9pt,left=3mm]{\Large\bfseries\color{cari}a)}\quad\begin{minipage}[c]{5.3cm}\raggedright\Large $39$\end{minipage}\end{tcolorbox}\hfill\begin{tcolorbox}[enhanced,nobeforeafter,width=7.1cm,height=2.1cm,valign=center,colback=cari!6,colframe=cari,arc=2.5mm,boxrule=0.9pt,left=3mm]{\Large\bfseries\color{cari}b)}\quad\begin{minipage}[c]{5.3cm}\raggedright\Large $45$\end{minipage}\end{tcolorbox}\par\vspace{2.5mm}
\begin{tcolorbox}[enhanced,nobeforeafter,width=7.1cm,height=2.1cm,valign=center,colback=cari!6,colframe=cari,arc=2.5mm,boxrule=0.9pt,left=3mm]{\Large\bfseries\color{cari}c)}\quad\begin{minipage}[c]{5.3cm}\raggedright\Large $43$\end{minipage}\end{tcolorbox}\hfill\begin{tcolorbox}[enhanced,nobeforeafter,width=7.1cm,height=2.1cm,valign=center,colback=cari!6,colframe=cari,arc=2.5mm,boxrule=0.9pt,left=3mm]{\Large\bfseries\color{cari}d)}\quad\begin{minipage}[c]{5.3cm}\raggedright\Large $51$\end{minipage}\end{tcolorbox}
\end{tcolorbox}
\newpage
\begin{tcolorbox}[enhanced,width=16cm,height=14cm,colback=white,colframe=cari,boxrule=2pt,arc=0mm,
left=6mm,right=6mm,top=0mm,bottom=4mm,sharp corners,
overlay={\fill[cari] (frame.north west) rectangle ([yshift=-1.25cm]frame.north east);
\node[anchor=west,text=white,font=\sffamily\bfseries\LARGE] at ([xshift=5mm,yshift=-0.62cm]frame.north west) {Question 3 D};
\node[anchor=east,text=white,font=\sffamily\itshape\large] at ([xshift=-5mm,yshift=-0.62cm]frame.north east) {Arithmétique et puissances};}]
\vspace*{1.75cm}
\begin{minipage}[t][6.9cm][c]{\linewidth}\LARGE\centering
\begin{minipage}{0.95\linewidth}\raggedright\linespread{1.25}\selectfont Quelle est la décomposition de $225$ en produit de facteurs premiers~?\end{minipage}
\end{minipage}

\vfill
\begin{tcolorbox}[enhanced,nobeforeafter,width=7.1cm,height=2.1cm,valign=center,colback=cari!6,colframe=cari,arc=2.5mm,boxrule=0.9pt,left=3mm]{\Large\bfseries\color{cari}a)}\quad\begin{minipage}[c]{5.3cm}\raggedright\Large $3^{3}\times5^{2}$\end{minipage}\end{tcolorbox}\hfill\begin{tcolorbox}[enhanced,nobeforeafter,width=7.1cm,height=2.1cm,valign=center,colback=cari!6,colframe=cari,arc=2.5mm,boxrule=0.9pt,left=3mm]{\Large\bfseries\color{cari}b)}\quad\begin{minipage}[c]{5.3cm}\raggedright\Large $3^{2}\times5^{2}$\end{minipage}\end{tcolorbox}\par\vspace{2.5mm}
\begin{tcolorbox}[enhanced,nobeforeafter,width=7.1cm,height=2.1cm,valign=center,colback=cari!6,colframe=cari,arc=2.5mm,boxrule=0.9pt,left=3mm]{\Large\bfseries\color{cari}c)}\quad\begin{minipage}[c]{5.3cm}\raggedright\Large $3^{2}\times5^{3}$\end{minipage}\end{tcolorbox}\hfill\begin{tcolorbox}[enhanced,nobeforeafter,width=7.1cm,height=2.1cm,valign=center,colback=cari!6,colframe=cari,arc=2.5mm,boxrule=0.9pt,left=3mm]{\Large\bfseries\color{cari}d)}\quad\begin{minipage}[c]{5.3cm}\raggedright\Large $45\times5$\end{minipage}\end{tcolorbox}
\end{tcolorbox}
\newpage
\begin{tcolorbox}[enhanced,width=16cm,height=14cm,colback=white,colframe=cari,boxrule=2pt,arc=0mm,
left=6mm,right=6mm,top=0mm,bottom=4mm,sharp corners,
overlay={\fill[cari] (frame.north west) rectangle ([yshift=-1.25cm]frame.north east);
\node[anchor=west,text=white,font=\sffamily\bfseries\LARGE] at ([xshift=5mm,yshift=-0.62cm]frame.north west) {Question 4 D};
\node[anchor=east,text=white,font=\sffamily\itshape\large] at ([xshift=-5mm,yshift=-0.62cm]frame.north east) {Arithmétique et puissances};}]
\vspace*{1.75cm}
\begin{minipage}[t][6.9cm][c]{\linewidth}\LARGE\centering
\begin{minipage}{0.95\linewidth}\raggedright\linespread{1.25}\selectfont Quelle est la décomposition de $280$ en produit de facteurs premiers~?\end{minipage}
\end{minipage}

\vfill
\begin{tcolorbox}[enhanced,nobeforeafter,width=7.1cm,height=2.1cm,valign=center,colback=cari!6,colframe=cari,arc=2.5mm,boxrule=0.9pt,left=3mm]{\Large\bfseries\color{cari}a)}\quad\begin{minipage}[c]{5.3cm}\raggedright\Large $2^{3}\times5\times7^{2}$\end{minipage}\end{tcolorbox}\hfill\begin{tcolorbox}[enhanced,nobeforeafter,width=7.1cm,height=2.1cm,valign=center,colback=cari!6,colframe=cari,arc=2.5mm,boxrule=0.9pt,left=3mm]{\Large\bfseries\color{cari}b)}\quad\begin{minipage}[c]{5.3cm}\raggedright\Large $2^{4}\times5\times7$\end{minipage}\end{tcolorbox}\par\vspace{2.5mm}
\begin{tcolorbox}[enhanced,nobeforeafter,width=7.1cm,height=2.1cm,valign=center,colback=cari!6,colframe=cari,arc=2.5mm,boxrule=0.9pt,left=3mm]{\Large\bfseries\color{cari}c)}\quad\begin{minipage}[c]{5.3cm}\raggedright\Large $2^{3}\times5\times7$\end{minipage}\end{tcolorbox}\hfill\begin{tcolorbox}[enhanced,nobeforeafter,width=7.1cm,height=2.1cm,valign=center,colback=cari!6,colframe=cari,arc=2.5mm,boxrule=0.9pt,left=3mm]{\Large\bfseries\color{cari}d)}\quad\begin{minipage}[c]{5.3cm}\raggedright\Large $40\times7$\end{minipage}\end{tcolorbox}
\end{tcolorbox}
\newpage
\begin{tcolorbox}[enhanced,width=16cm,height=14cm,colback=white,colframe=cari,boxrule=2pt,arc=0mm,
left=6mm,right=6mm,top=0mm,bottom=4mm,sharp corners,
overlay={\fill[cari] (frame.north west) rectangle ([yshift=-1.25cm]frame.north east);
\node[anchor=west,text=white,font=\sffamily\bfseries\LARGE] at ([xshift=5mm,yshift=-0.62cm]frame.north west) {Question 5 D};
\node[anchor=east,text=white,font=\sffamily\itshape\large] at ([xshift=-5mm,yshift=-0.62cm]frame.north east) {Arithmétique et puissances};}]
\vspace*{1.75cm}
\begin{minipage}[t][6.9cm][c]{\linewidth}\LARGE\centering
\begin{minipage}{0.95\linewidth}\raggedright\linespread{1.25}\selectfont Quelle est la forme irréductible de $\dfrac{72}{270}$~?\end{minipage}
\end{minipage}

\vfill
\begin{tcolorbox}[enhanced,nobeforeafter,width=7.1cm,height=2.1cm,valign=center,colback=cari!6,colframe=cari,arc=2.5mm,boxrule=0.9pt,left=3mm]{\Large\bfseries\color{cari}a)}\quad\begin{minipage}[c]{5.3cm}\raggedright\Large $\dfrac{36}{135}$\end{minipage}\end{tcolorbox}\hfill\begin{tcolorbox}[enhanced,nobeforeafter,width=7.1cm,height=2.1cm,valign=center,colback=cari!6,colframe=cari,arc=2.5mm,boxrule=0.9pt,left=3mm]{\Large\bfseries\color{cari}b)}\quad\begin{minipage}[c]{5.3cm}\raggedright\Large $\dfrac{24}{90}$\end{minipage}\end{tcolorbox}\par\vspace{2.5mm}
\begin{tcolorbox}[enhanced,nobeforeafter,width=7.1cm,height=2.1cm,valign=center,colback=cari!6,colframe=cari,arc=2.5mm,boxrule=0.9pt,left=3mm]{\Large\bfseries\color{cari}c)}\quad\begin{minipage}[c]{5.3cm}\raggedright\Large $\dfrac{4}{15}$\end{minipage}\end{tcolorbox}\hfill\begin{tcolorbox}[enhanced,nobeforeafter,width=7.1cm,height=2.1cm,valign=center,colback=cari!6,colframe=cari,arc=2.5mm,boxrule=0.9pt,left=3mm]{\Large\bfseries\color{cari}d)}\quad\begin{minipage}[c]{5.3cm}\raggedright\Large $\dfrac{15}{4}$\end{minipage}\end{tcolorbox}
\end{tcolorbox}
\newpage
\begin{tcolorbox}[enhanced,width=16cm,height=14cm,colback=white,colframe=cari,boxrule=2pt,arc=0mm,
left=6mm,right=6mm,top=0mm,bottom=4mm,sharp corners,
overlay={\fill[cari] (frame.north west) rectangle ([yshift=-1.25cm]frame.north east);
\node[anchor=west,text=white,font=\sffamily\bfseries\LARGE] at ([xshift=5mm,yshift=-0.62cm]frame.north west) {Question 6 D};
\node[anchor=east,text=white,font=\sffamily\itshape\large] at ([xshift=-5mm,yshift=-0.62cm]frame.north east) {Arithmétique et puissances};}]
\vspace*{1.75cm}
\begin{minipage}[t][6.9cm][c]{\linewidth}\LARGE\centering
\begin{minipage}{0.95\linewidth}\raggedright\linespread{1.25}\selectfont Quelle est la forme irréductible de $\dfrac{264}{360}$~?\end{minipage}
\end{minipage}

\vfill
\begin{tcolorbox}[enhanced,nobeforeafter,width=7.1cm,height=2.1cm,valign=center,colback=cari!6,colframe=cari,arc=2.5mm,boxrule=0.9pt,left=3mm]{\Large\bfseries\color{cari}a)}\quad\begin{minipage}[c]{5.3cm}\raggedright\Large $\dfrac{132}{180}$\end{minipage}\end{tcolorbox}\hfill\begin{tcolorbox}[enhanced,nobeforeafter,width=7.1cm,height=2.1cm,valign=center,colback=cari!6,colframe=cari,arc=2.5mm,boxrule=0.9pt,left=3mm]{\Large\bfseries\color{cari}b)}\quad\begin{minipage}[c]{5.3cm}\raggedright\Large $\dfrac{15}{11}$\end{minipage}\end{tcolorbox}\par\vspace{2.5mm}
\begin{tcolorbox}[enhanced,nobeforeafter,width=7.1cm,height=2.1cm,valign=center,colback=cari!6,colframe=cari,arc=2.5mm,boxrule=0.9pt,left=3mm]{\Large\bfseries\color{cari}c)}\quad\begin{minipage}[c]{5.3cm}\raggedright\Large $\dfrac{11}{15}$\end{minipage}\end{tcolorbox}\hfill\begin{tcolorbox}[enhanced,nobeforeafter,width=7.1cm,height=2.1cm,valign=center,colback=cari!6,colframe=cari,arc=2.5mm,boxrule=0.9pt,left=3mm]{\Large\bfseries\color{cari}d)}\quad\begin{minipage}[c]{5.3cm}\raggedright\Large $\dfrac{88}{120}$\end{minipage}\end{tcolorbox}
\end{tcolorbox}
\newpage
\begin{tcolorbox}[enhanced,width=16cm,height=14cm,colback=white,colframe=cari,boxrule=2pt,arc=0mm,
left=6mm,right=6mm,top=0mm,bottom=4mm,sharp corners,
overlay={\fill[cari] (frame.north west) rectangle ([yshift=-1.25cm]frame.north east);
\node[anchor=west,text=white,font=\sffamily\bfseries\LARGE] at ([xshift=5mm,yshift=-0.62cm]frame.north west) {Question 7 D};
\node[anchor=east,text=white,font=\sffamily\itshape\large] at ([xshift=-5mm,yshift=-0.62cm]frame.north east) {Arithmétique et puissances};}]
\vspace*{1.75cm}
\begin{minipage}[t][6.9cm][c]{\linewidth}\LARGE\centering
\begin{minipage}{0.95\linewidth}\raggedright\linespread{1.25}\selectfont Quel est le plus grand diviseur commun à $66$ et $110$~?\end{minipage}
\end{minipage}

\vfill
\begin{tcolorbox}[enhanced,nobeforeafter,width=7.1cm,height=2.1cm,valign=center,colback=cari!6,colframe=cari,arc=2.5mm,boxrule=0.9pt,left=3mm]{\Large\bfseries\color{cari}a)}\quad\begin{minipage}[c]{5.3cm}\raggedright\Large $330$\end{minipage}\end{tcolorbox}\hfill\begin{tcolorbox}[enhanced,nobeforeafter,width=7.1cm,height=2.1cm,valign=center,colback=cari!6,colframe=cari,arc=2.5mm,boxrule=0.9pt,left=3mm]{\Large\bfseries\color{cari}b)}\quad\begin{minipage}[c]{5.3cm}\raggedright\Large $22$\end{minipage}\end{tcolorbox}\par\vspace{2.5mm}
\begin{tcolorbox}[enhanced,nobeforeafter,width=7.1cm,height=2.1cm,valign=center,colback=cari!6,colframe=cari,arc=2.5mm,boxrule=0.9pt,left=3mm]{\Large\bfseries\color{cari}c)}\quad\begin{minipage}[c]{5.3cm}\raggedright\Large $44$\end{minipage}\end{tcolorbox}\hfill\begin{tcolorbox}[enhanced,nobeforeafter,width=7.1cm,height=2.1cm,valign=center,colback=cari!6,colframe=cari,arc=2.5mm,boxrule=0.9pt,left=3mm]{\Large\bfseries\color{cari}d)}\quad\begin{minipage}[c]{5.3cm}\raggedright\Large $11$\end{minipage}\end{tcolorbox}
\end{tcolorbox}
\newpage
\begin{tcolorbox}[enhanced,width=16cm,height=14cm,colback=white,colframe=cari,boxrule=2pt,arc=0mm,
left=6mm,right=6mm,top=0mm,bottom=4mm,sharp corners,
overlay={\fill[cari] (frame.north west) rectangle ([yshift=-1.25cm]frame.north east);
\node[anchor=west,text=white,font=\sffamily\bfseries\LARGE] at ([xshift=5mm,yshift=-0.62cm]frame.north west) {Question 8 D};
\node[anchor=east,text=white,font=\sffamily\itshape\large] at ([xshift=-5mm,yshift=-0.62cm]frame.north east) {Arithmétique et puissances};}]
\vspace*{1.75cm}
\begin{minipage}[t][6.9cm][c]{\linewidth}\LARGE\centering
\begin{minipage}{0.95\linewidth}\raggedright\linespread{1.25}\selectfont Quel est le plus grand diviseur commun à $150$ et $225$~?\end{minipage}
\end{minipage}

\vfill
\begin{tcolorbox}[enhanced,nobeforeafter,width=7.1cm,height=2.1cm,valign=center,colback=cari!6,colframe=cari,arc=2.5mm,boxrule=0.9pt,left=3mm]{\Large\bfseries\color{cari}a)}\quad\begin{minipage}[c]{5.3cm}\raggedright\Large $150$\end{minipage}\end{tcolorbox}\hfill\begin{tcolorbox}[enhanced,nobeforeafter,width=7.1cm,height=2.1cm,valign=center,colback=cari!6,colframe=cari,arc=2.5mm,boxrule=0.9pt,left=3mm]{\Large\bfseries\color{cari}b)}\quad\begin{minipage}[c]{5.3cm}\raggedright\Large $25$\end{minipage}\end{tcolorbox}\par\vspace{2.5mm}
\begin{tcolorbox}[enhanced,nobeforeafter,width=7.1cm,height=2.1cm,valign=center,colback=cari!6,colframe=cari,arc=2.5mm,boxrule=0.9pt,left=3mm]{\Large\bfseries\color{cari}c)}\quad\begin{minipage}[c]{5.3cm}\raggedright\Large $450$\end{minipage}\end{tcolorbox}\hfill\begin{tcolorbox}[enhanced,nobeforeafter,width=7.1cm,height=2.1cm,valign=center,colback=cari!6,colframe=cari,arc=2.5mm,boxrule=0.9pt,left=3mm]{\Large\bfseries\color{cari}d)}\quad\begin{minipage}[c]{5.3cm}\raggedright\Large $75$\end{minipage}\end{tcolorbox}
\end{tcolorbox}
\newpage
\begin{tcolorbox}[enhanced,width=16cm,height=14cm,colback=white,colframe=cari,boxrule=2pt,arc=0mm,
left=6mm,right=6mm,top=0mm,bottom=4mm,sharp corners,
overlay={\fill[cari] (frame.north west) rectangle ([yshift=-1.25cm]frame.north east);
\node[anchor=west,text=white,font=\sffamily\bfseries\LARGE] at ([xshift=5mm,yshift=-0.62cm]frame.north west) {Question 9 D};
\node[anchor=east,text=white,font=\sffamily\itshape\large] at ([xshift=-5mm,yshift=-0.62cm]frame.north east) {Arithmétique et puissances};}]
\vspace*{1.75cm}
\begin{minipage}[t][6.9cm][c]{\linewidth}\Large\centering
\begin{minipage}{0.95\linewidth}\raggedright\linespread{1.25}\selectfont Un fleuriste a $120$ roses et $150$ tulipes. Il veut faire le plus grand nombre de bouquets identiques en utilisant toutes les fleurs. Combien de bouquets peut-il faire~?\end{minipage}
\end{minipage}

\vfill
\begin{tcolorbox}[enhanced,nobeforeafter,width=7.1cm,height=2.1cm,valign=center,colback=cari!6,colframe=cari,arc=2.5mm,boxrule=0.9pt,left=3mm]{\Large\bfseries\color{cari}a)}\quad\begin{minipage}[c]{5.3cm}\raggedright\Large $30$\end{minipage}\end{tcolorbox}\hfill\begin{tcolorbox}[enhanced,nobeforeafter,width=7.1cm,height=2.1cm,valign=center,colback=cari!6,colframe=cari,arc=2.5mm,boxrule=0.9pt,left=3mm]{\Large\bfseries\color{cari}b)}\quad\begin{minipage}[c]{5.3cm}\raggedright\Large $60$\end{minipage}\end{tcolorbox}\par\vspace{2.5mm}
\begin{tcolorbox}[enhanced,nobeforeafter,width=7.1cm,height=2.1cm,valign=center,colback=cari!6,colframe=cari,arc=2.5mm,boxrule=0.9pt,left=3mm]{\Large\bfseries\color{cari}c)}\quad\begin{minipage}[c]{5.3cm}\raggedright\Large $270$\end{minipage}\end{tcolorbox}\hfill\begin{tcolorbox}[enhanced,nobeforeafter,width=7.1cm,height=2.1cm,valign=center,colback=cari!6,colframe=cari,arc=2.5mm,boxrule=0.9pt,left=3mm]{\Large\bfseries\color{cari}d)}\quad\begin{minipage}[c]{5.3cm}\raggedright\Large $15$\end{minipage}\end{tcolorbox}
\end{tcolorbox}
\newpage
\begin{tcolorbox}[enhanced,width=16cm,height=14cm,colback=white,colframe=cari,boxrule=2pt,arc=0mm,
left=6mm,right=6mm,top=0mm,bottom=4mm,sharp corners,
overlay={\fill[cari] (frame.north west) rectangle ([yshift=-1.25cm]frame.north east);
\node[anchor=west,text=white,font=\sffamily\bfseries\LARGE] at ([xshift=5mm,yshift=-0.62cm]frame.north west) {Question 10 D};
\node[anchor=east,text=white,font=\sffamily\itshape\large] at ([xshift=-5mm,yshift=-0.62cm]frame.north east) {Arithmétique et puissances};}]
\vspace*{1.75cm}
\begin{minipage}[t][6.9cm][c]{\linewidth}\Large\centering
\begin{minipage}{0.95\linewidth}\raggedright\linespread{1.25}\selectfont Un fleuriste a $36$ roses et $132$ tulipes. Il veut faire le plus grand nombre de bouquets identiques en utilisant toutes les fleurs. Combien de bouquets peut-il faire~?\end{minipage}
\end{minipage}

\vfill
\begin{tcolorbox}[enhanced,nobeforeafter,width=7.1cm,height=2.1cm,valign=center,colback=cari!6,colframe=cari,arc=2.5mm,boxrule=0.9pt,left=3mm]{\Large\bfseries\color{cari}a)}\quad\begin{minipage}[c]{5.3cm}\raggedright\Large $24$\end{minipage}\end{tcolorbox}\hfill\begin{tcolorbox}[enhanced,nobeforeafter,width=7.1cm,height=2.1cm,valign=center,colback=cari!6,colframe=cari,arc=2.5mm,boxrule=0.9pt,left=3mm]{\Large\bfseries\color{cari}b)}\quad\begin{minipage}[c]{5.3cm}\raggedright\Large $12$\end{minipage}\end{tcolorbox}\par\vspace{2.5mm}
\begin{tcolorbox}[enhanced,nobeforeafter,width=7.1cm,height=2.1cm,valign=center,colback=cari!6,colframe=cari,arc=2.5mm,boxrule=0.9pt,left=3mm]{\Large\bfseries\color{cari}c)}\quad\begin{minipage}[c]{5.3cm}\raggedright\Large $6$\end{minipage}\end{tcolorbox}\hfill\begin{tcolorbox}[enhanced,nobeforeafter,width=7.1cm,height=2.1cm,valign=center,colback=cari!6,colframe=cari,arc=2.5mm,boxrule=0.9pt,left=3mm]{\Large\bfseries\color{cari}d)}\quad\begin{minipage}[c]{5.3cm}\raggedright\Large $168$\end{minipage}\end{tcolorbox}
\end{tcolorbox}
\newpage
\begin{tcolorbox}[enhanced,width=16cm,height=14cm,colback=white,colframe=cari,boxrule=2pt,arc=0mm,
left=6mm,right=6mm,top=0mm,bottom=4mm,sharp corners,
overlay={\fill[cari] (frame.north west) rectangle ([yshift=-1.25cm]frame.north east);
\node[anchor=west,text=white,font=\sffamily\bfseries\LARGE] at ([xshift=5mm,yshift=-0.62cm]frame.north west) {Question 11 D};
\node[anchor=east,text=white,font=\sffamily\itshape\large] at ([xshift=-5mm,yshift=-0.62cm]frame.north east) {Arithmétique et puissances};}]
\vspace*{1.75cm}
\begin{minipage}[t][6.9cm][c]{\linewidth}\Large\centering
\begin{minipage}{0.95\linewidth}\raggedright\linespread{1.25}\selectfont Un phare clignote toutes les $10$ secondes et un autre toutes les $25$ secondes. Ils viennent de clignoter ensemble. Dans combien de secondes clignoteront-ils de nouveau ensemble~?\end{minipage}
\end{minipage}

\vfill
\begin{tcolorbox}[enhanced,nobeforeafter,width=7.1cm,height=2.1cm,valign=center,colback=cari!6,colframe=cari,arc=2.5mm,boxrule=0.9pt,left=3mm]{\Large\bfseries\color{cari}a)}\quad\begin{minipage}[c]{5.3cm}\raggedright\Large $50~\text{s}$\end{minipage}\end{tcolorbox}\hfill\begin{tcolorbox}[enhanced,nobeforeafter,width=7.1cm,height=2.1cm,valign=center,colback=cari!6,colframe=cari,arc=2.5mm,boxrule=0.9pt,left=3mm]{\Large\bfseries\color{cari}b)}\quad\begin{minipage}[c]{5.3cm}\raggedright\Large $250~\text{s}$\end{minipage}\end{tcolorbox}\par\vspace{2.5mm}
\begin{tcolorbox}[enhanced,nobeforeafter,width=7.1cm,height=2.1cm,valign=center,colback=cari!6,colframe=cari,arc=2.5mm,boxrule=0.9pt,left=3mm]{\Large\bfseries\color{cari}c)}\quad\begin{minipage}[c]{5.3cm}\raggedright\Large $35~\text{s}$\end{minipage}\end{tcolorbox}\hfill\begin{tcolorbox}[enhanced,nobeforeafter,width=7.1cm,height=2.1cm,valign=center,colback=cari!6,colframe=cari,arc=2.5mm,boxrule=0.9pt,left=3mm]{\Large\bfseries\color{cari}d)}\quad\begin{minipage}[c]{5.3cm}\raggedright\Large $5~\text{s}$\end{minipage}\end{tcolorbox}
\end{tcolorbox}
\newpage
\begin{tcolorbox}[enhanced,width=16cm,height=14cm,colback=white,colframe=cari,boxrule=2pt,arc=0mm,
left=6mm,right=6mm,top=0mm,bottom=4mm,sharp corners,
overlay={\fill[cari] (frame.north west) rectangle ([yshift=-1.25cm]frame.north east);
\node[anchor=west,text=white,font=\sffamily\bfseries\LARGE] at ([xshift=5mm,yshift=-0.62cm]frame.north west) {Question 12 D};
\node[anchor=east,text=white,font=\sffamily\itshape\large] at ([xshift=-5mm,yshift=-0.62cm]frame.north east) {Arithmétique et puissances};}]
\vspace*{1.75cm}
\begin{minipage}[t][6.9cm][c]{\linewidth}\Large\centering
\begin{minipage}{0.95\linewidth}\raggedright\linespread{1.25}\selectfont Un phare clignote toutes les $6$ secondes et un autre toutes les $21$ secondes. Ils viennent de clignoter ensemble. Dans combien de secondes clignoteront-ils de nouveau ensemble~?\end{minipage}
\end{minipage}

\vfill
\begin{tcolorbox}[enhanced,nobeforeafter,width=7.1cm,height=2.1cm,valign=center,colback=cari!6,colframe=cari,arc=2.5mm,boxrule=0.9pt,left=3mm]{\Large\bfseries\color{cari}a)}\quad\begin{minipage}[c]{5.3cm}\raggedright\Large $27~\text{s}$\end{minipage}\end{tcolorbox}\hfill\begin{tcolorbox}[enhanced,nobeforeafter,width=7.1cm,height=2.1cm,valign=center,colback=cari!6,colframe=cari,arc=2.5mm,boxrule=0.9pt,left=3mm]{\Large\bfseries\color{cari}b)}\quad\begin{minipage}[c]{5.3cm}\raggedright\Large $42~\text{s}$\end{minipage}\end{tcolorbox}\par\vspace{2.5mm}
\begin{tcolorbox}[enhanced,nobeforeafter,width=7.1cm,height=2.1cm,valign=center,colback=cari!6,colframe=cari,arc=2.5mm,boxrule=0.9pt,left=3mm]{\Large\bfseries\color{cari}c)}\quad\begin{minipage}[c]{5.3cm}\raggedright\Large $126~\text{s}$\end{minipage}\end{tcolorbox}\hfill\begin{tcolorbox}[enhanced,nobeforeafter,width=7.1cm,height=2.1cm,valign=center,colback=cari!6,colframe=cari,arc=2.5mm,boxrule=0.9pt,left=3mm]{\Large\bfseries\color{cari}d)}\quad\begin{minipage}[c]{5.3cm}\raggedright\Large $3~\text{s}$\end{minipage}\end{tcolorbox}
\end{tcolorbox}
\newpage
\begin{tcolorbox}[enhanced,width=16cm,height=14cm,colback=white,colframe=cari,boxrule=2pt,arc=0mm,
left=6mm,right=6mm,top=0mm,bottom=4mm,sharp corners,
overlay={\fill[cari] (frame.north west) rectangle ([yshift=-1.25cm]frame.north east);
\node[anchor=west,text=white,font=\sffamily\bfseries\LARGE] at ([xshift=5mm,yshift=-0.62cm]frame.north west) {Question 13 D};
\node[anchor=east,text=white,font=\sffamily\itshape\large] at ([xshift=-5mm,yshift=-0.62cm]frame.north east) {Arithmétique et puissances};}]
\vspace*{1.75cm}
\begin{minipage}[t][6.9cm][c]{\linewidth}\LARGE\centering
\begin{minipage}{0.95\linewidth}\raggedright\linespread{1.25}\selectfont Que vaut $(8\times10^{3})\times(3\times10^{-8})$, en notation scientifique~?\end{minipage}
\end{minipage}

\vfill
\begin{tcolorbox}[enhanced,nobeforeafter,width=7.1cm,height=2.1cm,valign=center,colback=cari!6,colframe=cari,arc=2.5mm,boxrule=0.9pt,left=3mm]{\Large\bfseries\color{cari}a)}\quad\begin{minipage}[c]{5.3cm}\raggedright\Large $2{,}4\times10^{-24}$\end{minipage}\end{tcolorbox}\hfill\begin{tcolorbox}[enhanced,nobeforeafter,width=7.1cm,height=2.1cm,valign=center,colback=cari!6,colframe=cari,arc=2.5mm,boxrule=0.9pt,left=3mm]{\Large\bfseries\color{cari}b)}\quad\begin{minipage}[c]{5.3cm}\raggedright\Large $24\times10^{-5}$\end{minipage}\end{tcolorbox}\par\vspace{2.5mm}
\begin{tcolorbox}[enhanced,nobeforeafter,width=7.1cm,height=2.1cm,valign=center,colback=cari!6,colframe=cari,arc=2.5mm,boxrule=0.9pt,left=3mm]{\Large\bfseries\color{cari}c)}\quad\begin{minipage}[c]{5.3cm}\raggedright\Large $2{,}4\times10^{-5}$\end{minipage}\end{tcolorbox}\hfill\begin{tcolorbox}[enhanced,nobeforeafter,width=7.1cm,height=2.1cm,valign=center,colback=cari!6,colframe=cari,arc=2.5mm,boxrule=0.9pt,left=3mm]{\Large\bfseries\color{cari}d)}\quad\begin{minipage}[c]{5.3cm}\raggedright\Large $2{,}4\times10^{-4}$\end{minipage}\end{tcolorbox}
\end{tcolorbox}
\newpage
\begin{tcolorbox}[enhanced,width=16cm,height=14cm,colback=white,colframe=cari,boxrule=2pt,arc=0mm,
left=6mm,right=6mm,top=0mm,bottom=4mm,sharp corners,
overlay={\fill[cari] (frame.north west) rectangle ([yshift=-1.25cm]frame.north east);
\node[anchor=west,text=white,font=\sffamily\bfseries\LARGE] at ([xshift=5mm,yshift=-0.62cm]frame.north west) {Question 14 D};
\node[anchor=east,text=white,font=\sffamily\itshape\large] at ([xshift=-5mm,yshift=-0.62cm]frame.north east) {Arithmétique et puissances};}]
\vspace*{1.75cm}
\begin{minipage}[t][6.9cm][c]{\linewidth}\LARGE\centering
\begin{minipage}{0.95\linewidth}\raggedright\linespread{1.25}\selectfont Que vaut $(7\times10^{-1})\times(2\times10^{9})$, en notation scientifique~?\end{minipage}
\end{minipage}

\vfill
\begin{tcolorbox}[enhanced,nobeforeafter,width=7.1cm,height=2.1cm,valign=center,colback=cari!6,colframe=cari,arc=2.5mm,boxrule=0.9pt,left=3mm]{\Large\bfseries\color{cari}a)}\quad\begin{minipage}[c]{5.3cm}\raggedright\Large $14\times10^{8}$\end{minipage}\end{tcolorbox}\hfill\begin{tcolorbox}[enhanced,nobeforeafter,width=7.1cm,height=2.1cm,valign=center,colback=cari!6,colframe=cari,arc=2.5mm,boxrule=0.9pt,left=3mm]{\Large\bfseries\color{cari}b)}\quad\begin{minipage}[c]{5.3cm}\raggedright\Large $1{,}4\times10^{8}$\end{minipage}\end{tcolorbox}\par\vspace{2.5mm}
\begin{tcolorbox}[enhanced,nobeforeafter,width=7.1cm,height=2.1cm,valign=center,colback=cari!6,colframe=cari,arc=2.5mm,boxrule=0.9pt,left=3mm]{\Large\bfseries\color{cari}c)}\quad\begin{minipage}[c]{5.3cm}\raggedright\Large $1{,}4\times10^{-9}$\end{minipage}\end{tcolorbox}\hfill\begin{tcolorbox}[enhanced,nobeforeafter,width=7.1cm,height=2.1cm,valign=center,colback=cari!6,colframe=cari,arc=2.5mm,boxrule=0.9pt,left=3mm]{\Large\bfseries\color{cari}d)}\quad\begin{minipage}[c]{5.3cm}\raggedright\Large $1{,}4\times10^{9}$\end{minipage}\end{tcolorbox}
\end{tcolorbox}
\newpage
\begin{tcolorbox}[enhanced,width=16cm,height=14cm,colback=white,colframe=cari,boxrule=2pt,arc=0mm,
left=6mm,right=6mm,top=0mm,bottom=4mm,sharp corners,
overlay={\fill[cari] (frame.north west) rectangle ([yshift=-1.25cm]frame.north east);
\node[anchor=west,text=white,font=\sffamily\bfseries\LARGE] at ([xshift=5mm,yshift=-0.62cm]frame.north west) {Question 15 D};
\node[anchor=east,text=white,font=\sffamily\itshape\large] at ([xshift=-5mm,yshift=-0.62cm]frame.north east) {Arithmétique et puissances};}]
\vspace*{1.75cm}
\begin{minipage}[t][6.9cm][c]{\linewidth}\LARGE\centering
\begin{minipage}{0.95\linewidth}\raggedright\linespread{1.25}\selectfont Quelle est l'écriture décimale de $8^{-1}$~?\end{minipage}
\end{minipage}

\vfill
\begin{tcolorbox}[enhanced,nobeforeafter,width=7.1cm,height=2.1cm,valign=center,colback=cari!6,colframe=cari,arc=2.5mm,boxrule=0.9pt,left=3mm]{\Large\bfseries\color{cari}a)}\quad\begin{minipage}[c]{5.3cm}\raggedright\Large $-8$\end{minipage}\end{tcolorbox}\hfill\begin{tcolorbox}[enhanced,nobeforeafter,width=7.1cm,height=2.1cm,valign=center,colback=cari!6,colframe=cari,arc=2.5mm,boxrule=0.9pt,left=3mm]{\Large\bfseries\color{cari}b)}\quad\begin{minipage}[c]{5.3cm}\raggedright\Large $0{,}125$\end{minipage}\end{tcolorbox}\par\vspace{2.5mm}
\begin{tcolorbox}[enhanced,nobeforeafter,width=7.1cm,height=2.1cm,valign=center,colback=cari!6,colframe=cari,arc=2.5mm,boxrule=0.9pt,left=3mm]{\Large\bfseries\color{cari}c)}\quad\begin{minipage}[c]{5.3cm}\raggedright\Large $1{,}25$\end{minipage}\end{tcolorbox}\hfill\begin{tcolorbox}[enhanced,nobeforeafter,width=7.1cm,height=2.1cm,valign=center,colback=cari!6,colframe=cari,arc=2.5mm,boxrule=0.9pt,left=3mm]{\Large\bfseries\color{cari}d)}\quad\begin{minipage}[c]{5.3cm}\raggedright\Large $0{,}0125$\end{minipage}\end{tcolorbox}
\end{tcolorbox}
\newpage
\begin{tcolorbox}[enhanced,width=16cm,height=14cm,colback=white,colframe=cari,boxrule=2pt,arc=0mm,
left=6mm,right=6mm,top=0mm,bottom=4mm,sharp corners,
overlay={\fill[cari] (frame.north west) rectangle ([yshift=-1.25cm]frame.north east);
\node[anchor=west,text=white,font=\sffamily\bfseries\LARGE] at ([xshift=5mm,yshift=-0.62cm]frame.north west) {Question 16 D};
\node[anchor=east,text=white,font=\sffamily\itshape\large] at ([xshift=-5mm,yshift=-0.62cm]frame.north east) {Arithmétique et puissances};}]
\vspace*{1.75cm}
\begin{minipage}[t][6.9cm][c]{\linewidth}\LARGE\centering
\begin{minipage}{0.95\linewidth}\raggedright\linespread{1.25}\selectfont Quelle est l'écriture décimale de $5^{-3}$~?\end{minipage}
\end{minipage}

\vfill
\begin{tcolorbox}[enhanced,nobeforeafter,width=7.1cm,height=2.1cm,valign=center,colback=cari!6,colframe=cari,arc=2.5mm,boxrule=0.9pt,left=3mm]{\Large\bfseries\color{cari}a)}\quad\begin{minipage}[c]{5.3cm}\raggedright\Large $-15$\end{minipage}\end{tcolorbox}\hfill\begin{tcolorbox}[enhanced,nobeforeafter,width=7.1cm,height=2.1cm,valign=center,colback=cari!6,colframe=cari,arc=2.5mm,boxrule=0.9pt,left=3mm]{\Large\bfseries\color{cari}b)}\quad\begin{minipage}[c]{5.3cm}\raggedright\Large $-125$\end{minipage}\end{tcolorbox}\par\vspace{2.5mm}
\begin{tcolorbox}[enhanced,nobeforeafter,width=7.1cm,height=2.1cm,valign=center,colback=cari!6,colframe=cari,arc=2.5mm,boxrule=0.9pt,left=3mm]{\Large\bfseries\color{cari}c)}\quad\begin{minipage}[c]{5.3cm}\raggedright\Large $0{,}008$\end{minipage}\end{tcolorbox}\hfill\begin{tcolorbox}[enhanced,nobeforeafter,width=7.1cm,height=2.1cm,valign=center,colback=cari!6,colframe=cari,arc=2.5mm,boxrule=0.9pt,left=3mm]{\Large\bfseries\color{cari}d)}\quad\begin{minipage}[c]{5.3cm}\raggedright\Large $\dfrac{1}{15}$\end{minipage}\end{tcolorbox}
\end{tcolorbox}
\newpage
\begin{tcolorbox}[enhanced,width=16cm,height=14cm,colback=white,colframe=clit,boxrule=2pt,arc=0mm,
left=6mm,right=6mm,top=0mm,bottom=4mm,sharp corners,
overlay={\fill[clit] (frame.north west) rectangle ([yshift=-1.25cm]frame.north east);
\node[anchor=west,text=white,font=\sffamily\bfseries\LARGE] at ([xshift=5mm,yshift=-0.62cm]frame.north west) {Question 17 D};
\node[anchor=east,text=white,font=\sffamily\itshape\large] at ([xshift=-5mm,yshift=-0.62cm]frame.north east) {Calcul littéral et équations};}]
\vspace*{1.75cm}
\begin{minipage}[t][6.9cm][c]{\linewidth}\LARGE\centering
\begin{minipage}{0.95\linewidth}\raggedright\linespread{1.25}\selectfont Développer et réduire $(x-7)(3x+2)$.\end{minipage}
\end{minipage}

\vfill
\begin{tcolorbox}[enhanced,nobeforeafter,width=7.1cm,height=2.1cm,valign=center,colback=clit!6,colframe=clit,arc=2.5mm,boxrule=0.9pt,left=3mm]{\Large\bfseries\color{clit}a)}\quad\begin{minipage}[c]{5.3cm}\raggedright\Large $3x^2-14$\end{minipage}\end{tcolorbox}\hfill\begin{tcolorbox}[enhanced,nobeforeafter,width=7.1cm,height=2.1cm,valign=center,colback=clit!6,colframe=clit,arc=2.5mm,boxrule=0.9pt,left=3mm]{\Large\bfseries\color{clit}b)}\quad\begin{minipage}[c]{5.3cm}\raggedright\Large $3x^2-19x-14$\end{minipage}\end{tcolorbox}\par\vspace{2.5mm}
\begin{tcolorbox}[enhanced,nobeforeafter,width=7.1cm,height=2.1cm,valign=center,colback=clit!6,colframe=clit,arc=2.5mm,boxrule=0.9pt,left=3mm]{\Large\bfseries\color{clit}c)}\quad\begin{minipage}[c]{5.3cm}\raggedright\Large $4x^2-19x-14$\end{minipage}\end{tcolorbox}\hfill\begin{tcolorbox}[enhanced,nobeforeafter,width=7.1cm,height=2.1cm,valign=center,colback=clit!6,colframe=clit,arc=2.5mm,boxrule=0.9pt,left=3mm]{\Large\bfseries\color{clit}d)}\quad\begin{minipage}[c]{5.3cm}\raggedright\Large $3x^2+23x-14$\end{minipage}\end{tcolorbox}
\end{tcolorbox}
\newpage
\begin{tcolorbox}[enhanced,width=16cm,height=14cm,colback=white,colframe=clit,boxrule=2pt,arc=0mm,
left=6mm,right=6mm,top=0mm,bottom=4mm,sharp corners,
overlay={\fill[clit] (frame.north west) rectangle ([yshift=-1.25cm]frame.north east);
\node[anchor=west,text=white,font=\sffamily\bfseries\LARGE] at ([xshift=5mm,yshift=-0.62cm]frame.north west) {Question 18 D};
\node[anchor=east,text=white,font=\sffamily\itshape\large] at ([xshift=-5mm,yshift=-0.62cm]frame.north east) {Calcul littéral et équations};}]
\vspace*{1.75cm}
\begin{minipage}[t][6.9cm][c]{\linewidth}\LARGE\centering
\begin{minipage}{0.95\linewidth}\raggedright\linespread{1.25}\selectfont Développer et réduire $(-3x+2)(-x+5)$.\end{minipage}
\end{minipage}

\vfill
\begin{tcolorbox}[enhanced,nobeforeafter,width=7.1cm,height=2.1cm,valign=center,colback=clit!6,colframe=clit,arc=2.5mm,boxrule=0.9pt,left=3mm]{\Large\bfseries\color{clit}a)}\quad\begin{minipage}[c]{5.3cm}\raggedright\Large $-4x^2-17x+10$\end{minipage}\end{tcolorbox}\hfill\begin{tcolorbox}[enhanced,nobeforeafter,width=7.1cm,height=2.1cm,valign=center,colback=clit!6,colframe=clit,arc=2.5mm,boxrule=0.9pt,left=3mm]{\Large\bfseries\color{clit}b)}\quad\begin{minipage}[c]{5.3cm}\raggedright\Large $3x^2-13x+10$\end{minipage}\end{tcolorbox}\par\vspace{2.5mm}
\begin{tcolorbox}[enhanced,nobeforeafter,width=7.1cm,height=2.1cm,valign=center,colback=clit!6,colframe=clit,arc=2.5mm,boxrule=0.9pt,left=3mm]{\Large\bfseries\color{clit}c)}\quad\begin{minipage}[c]{5.3cm}\raggedright\Large $3x^2+10$\end{minipage}\end{tcolorbox}\hfill\begin{tcolorbox}[enhanced,nobeforeafter,width=7.1cm,height=2.1cm,valign=center,colback=clit!6,colframe=clit,arc=2.5mm,boxrule=0.9pt,left=3mm]{\Large\bfseries\color{clit}d)}\quad\begin{minipage}[c]{5.3cm}\raggedright\Large $3x^2-17x+10$\end{minipage}\end{tcolorbox}
\end{tcolorbox}
\newpage
\begin{tcolorbox}[enhanced,width=16cm,height=14cm,colback=white,colframe=clit,boxrule=2pt,arc=0mm,
left=6mm,right=6mm,top=0mm,bottom=4mm,sharp corners,
overlay={\fill[clit] (frame.north west) rectangle ([yshift=-1.25cm]frame.north east);
\node[anchor=west,text=white,font=\sffamily\bfseries\LARGE] at ([xshift=5mm,yshift=-0.62cm]frame.north west) {Question 19 D};
\node[anchor=east,text=white,font=\sffamily\itshape\large] at ([xshift=-5mm,yshift=-0.62cm]frame.north east) {Calcul littéral et équations};}]
\vspace*{1.75cm}
\begin{minipage}[t][6.9cm][c]{\linewidth}\LARGE\centering
\begin{minipage}{0.95\linewidth}\raggedright\linespread{1.25}\selectfont Factoriser $x^2-81$.\end{minipage}
\end{minipage}

\vfill
\begin{tcolorbox}[enhanced,nobeforeafter,width=7.1cm,height=2.1cm,valign=center,colback=clit!6,colframe=clit,arc=2.5mm,boxrule=0.9pt,left=3mm]{\Large\bfseries\color{clit}a)}\quad\begin{minipage}[c]{5.3cm}\raggedright\Large $(x-81)(x+81)$\end{minipage}\end{tcolorbox}\hfill\begin{tcolorbox}[enhanced,nobeforeafter,width=7.1cm,height=2.1cm,valign=center,colback=clit!6,colframe=clit,arc=2.5mm,boxrule=0.9pt,left=3mm]{\Large\bfseries\color{clit}b)}\quad\begin{minipage}[c]{5.3cm}\raggedright\Large $(x+9)^2$\end{minipage}\end{tcolorbox}\par\vspace{2.5mm}
\begin{tcolorbox}[enhanced,nobeforeafter,width=7.1cm,height=2.1cm,valign=center,colback=clit!6,colframe=clit,arc=2.5mm,boxrule=0.9pt,left=3mm]{\Large\bfseries\color{clit}c)}\quad\begin{minipage}[c]{5.3cm}\raggedright\Large $(x-9)(x+9)$\end{minipage}\end{tcolorbox}\hfill\begin{tcolorbox}[enhanced,nobeforeafter,width=7.1cm,height=2.1cm,valign=center,colback=clit!6,colframe=clit,arc=2.5mm,boxrule=0.9pt,left=3mm]{\Large\bfseries\color{clit}d)}\quad\begin{minipage}[c]{5.3cm}\raggedright\Large $(x-9)^2$\end{minipage}\end{tcolorbox}
\end{tcolorbox}
\newpage
\begin{tcolorbox}[enhanced,width=16cm,height=14cm,colback=white,colframe=clit,boxrule=2pt,arc=0mm,
left=6mm,right=6mm,top=0mm,bottom=4mm,sharp corners,
overlay={\fill[clit] (frame.north west) rectangle ([yshift=-1.25cm]frame.north east);
\node[anchor=west,text=white,font=\sffamily\bfseries\LARGE] at ([xshift=5mm,yshift=-0.62cm]frame.north west) {Question 20 D};
\node[anchor=east,text=white,font=\sffamily\itshape\large] at ([xshift=-5mm,yshift=-0.62cm]frame.north east) {Calcul littéral et équations};}]
\vspace*{1.75cm}
\begin{minipage}[t][6.9cm][c]{\linewidth}\LARGE\centering
\begin{minipage}{0.95\linewidth}\raggedright\linespread{1.25}\selectfont Factoriser $36x^2-25$.\end{minipage}
\end{minipage}

\vfill
\begin{tcolorbox}[enhanced,nobeforeafter,width=7.1cm,height=2.1cm,valign=center,colback=clit!6,colframe=clit,arc=2.5mm,boxrule=0.9pt,left=3mm]{\Large\bfseries\color{clit}a)}\quad\begin{minipage}[c]{5.3cm}\raggedright\Large $(6x-25)(6x+25)$\end{minipage}\end{tcolorbox}\hfill\begin{tcolorbox}[enhanced,nobeforeafter,width=7.1cm,height=2.1cm,valign=center,colback=clit!6,colframe=clit,arc=2.5mm,boxrule=0.9pt,left=3mm]{\Large\bfseries\color{clit}b)}\quad\begin{minipage}[c]{5.3cm}\raggedright\Large $(6x-5)^2$\end{minipage}\end{tcolorbox}\par\vspace{2.5mm}
\begin{tcolorbox}[enhanced,nobeforeafter,width=7.1cm,height=2.1cm,valign=center,colback=clit!6,colframe=clit,arc=2.5mm,boxrule=0.9pt,left=3mm]{\Large\bfseries\color{clit}c)}\quad\begin{minipage}[c]{5.3cm}\raggedright\Large $(6x+5)^2$\end{minipage}\end{tcolorbox}\hfill\begin{tcolorbox}[enhanced,nobeforeafter,width=7.1cm,height=2.1cm,valign=center,colback=clit!6,colframe=clit,arc=2.5mm,boxrule=0.9pt,left=3mm]{\Large\bfseries\color{clit}d)}\quad\begin{minipage}[c]{5.3cm}\raggedright\Large $(6x-5)(6x+5)$\end{minipage}\end{tcolorbox}
\end{tcolorbox}
\newpage
\begin{tcolorbox}[enhanced,width=16cm,height=14cm,colback=white,colframe=clit,boxrule=2pt,arc=0mm,
left=6mm,right=6mm,top=0mm,bottom=4mm,sharp corners,
overlay={\fill[clit] (frame.north west) rectangle ([yshift=-1.25cm]frame.north east);
\node[anchor=west,text=white,font=\sffamily\bfseries\LARGE] at ([xshift=5mm,yshift=-0.62cm]frame.north west) {Question 21 D};
\node[anchor=east,text=white,font=\sffamily\itshape\large] at ([xshift=-5mm,yshift=-0.62cm]frame.north east) {Calcul littéral et équations};}]
\vspace*{1.75cm}
\begin{minipage}[t][6.9cm][c]{\linewidth}\LARGE\centering
\begin{minipage}{0.95\linewidth}\raggedright\linespread{1.25}\selectfont Quelles sont les solutions de l'équation $(5x+10)(2x-3)=0$~?\end{minipage}
\end{minipage}

\vfill
\begin{tcolorbox}[enhanced,nobeforeafter,width=7.1cm,height=2.1cm,valign=center,colback=clit!6,colframe=clit,arc=2.5mm,boxrule=0.9pt,left=3mm]{\Large\bfseries\color{clit}a)}\quad\begin{minipage}[c]{5.3cm}\raggedright\Large $\{-2\}$\end{minipage}\end{tcolorbox}\hfill\begin{tcolorbox}[enhanced,nobeforeafter,width=7.1cm,height=2.1cm,valign=center,colback=clit!6,colframe=clit,arc=2.5mm,boxrule=0.9pt,left=3mm]{\Large\bfseries\color{clit}b)}\quad\begin{minipage}[c]{5.3cm}\raggedright\Large $\{-2\,;\,\dfrac{3}{2}\}$\end{minipage}\end{tcolorbox}\par\vspace{2.5mm}
\begin{tcolorbox}[enhanced,nobeforeafter,width=7.1cm,height=2.1cm,valign=center,colback=clit!6,colframe=clit,arc=2.5mm,boxrule=0.9pt,left=3mm]{\Large\bfseries\color{clit}c)}\quad\begin{minipage}[c]{5.3cm}\raggedright\Large $\{-\dfrac{3}{2}\,;\,2\}$\end{minipage}\end{tcolorbox}\hfill\begin{tcolorbox}[enhanced,nobeforeafter,width=7.1cm,height=2.1cm,valign=center,colback=clit!6,colframe=clit,arc=2.5mm,boxrule=0.9pt,left=3mm]{\Large\bfseries\color{clit}d)}\quad\begin{minipage}[c]{5.3cm}\raggedright\Large $\{-3\,;\,10\}$\end{minipage}\end{tcolorbox}
\end{tcolorbox}
\newpage
\begin{tcolorbox}[enhanced,width=16cm,height=14cm,colback=white,colframe=clit,boxrule=2pt,arc=0mm,
left=6mm,right=6mm,top=0mm,bottom=4mm,sharp corners,
overlay={\fill[clit] (frame.north west) rectangle ([yshift=-1.25cm]frame.north east);
\node[anchor=west,text=white,font=\sffamily\bfseries\LARGE] at ([xshift=5mm,yshift=-0.62cm]frame.north west) {Question 22 D};
\node[anchor=east,text=white,font=\sffamily\itshape\large] at ([xshift=-5mm,yshift=-0.62cm]frame.north east) {Calcul littéral et équations};}]
\vspace*{1.75cm}
\begin{minipage}[t][6.9cm][c]{\linewidth}\LARGE\centering
\begin{minipage}{0.95\linewidth}\raggedright\linespread{1.25}\selectfont Quelles sont les solutions de l'équation $(x-8)(x+8)=0$~?\end{minipage}
\end{minipage}

\vfill
\begin{tcolorbox}[enhanced,nobeforeafter,width=7.1cm,height=2.1cm,valign=center,colback=clit!6,colframe=clit,arc=2.5mm,boxrule=0.9pt,left=3mm]{\Large\bfseries\color{clit}a)}\quad\begin{minipage}[c]{5.3cm}\raggedright\Large $\{8\}$\end{minipage}\end{tcolorbox}\hfill\begin{tcolorbox}[enhanced,nobeforeafter,width=7.1cm,height=2.1cm,valign=center,colback=clit!6,colframe=clit,arc=2.5mm,boxrule=0.9pt,left=3mm]{\Large\bfseries\color{clit}b)}\quad\begin{minipage}[c]{5.3cm}\raggedright\Large $\{-16\,;\,16\}$\end{minipage}\end{tcolorbox}\par\vspace{2.5mm}
\begin{tcolorbox}[enhanced,nobeforeafter,width=7.1cm,height=2.1cm,valign=center,colback=clit!6,colframe=clit,arc=2.5mm,boxrule=0.9pt,left=3mm]{\Large\bfseries\color{clit}c)}\quad\begin{minipage}[c]{5.3cm}\raggedright\Large $\{-7\,;\,8\}$\end{minipage}\end{tcolorbox}\hfill\begin{tcolorbox}[enhanced,nobeforeafter,width=7.1cm,height=2.1cm,valign=center,colback=clit!6,colframe=clit,arc=2.5mm,boxrule=0.9pt,left=3mm]{\Large\bfseries\color{clit}d)}\quad\begin{minipage}[c]{5.3cm}\raggedright\Large $\{-8\,;\,8\}$\end{minipage}\end{tcolorbox}
\end{tcolorbox}
\newpage
\begin{tcolorbox}[enhanced,width=16cm,height=14cm,colback=white,colframe=clit,boxrule=2pt,arc=0mm,
left=6mm,right=6mm,top=0mm,bottom=4mm,sharp corners,
overlay={\fill[clit] (frame.north west) rectangle ([yshift=-1.25cm]frame.north east);
\node[anchor=west,text=white,font=\sffamily\bfseries\LARGE] at ([xshift=5mm,yshift=-0.62cm]frame.north west) {Question 23 D};
\node[anchor=east,text=white,font=\sffamily\itshape\large] at ([xshift=-5mm,yshift=-0.62cm]frame.north east) {Calcul littéral et équations};}]
\vspace*{1.75cm}
\begin{minipage}[t][6.9cm][c]{\linewidth}\LARGE\centering
\begin{minipage}{0.95\linewidth}\raggedright\linespread{1.25}\selectfont Factoriser $(2x-3)(x-4)+(2x-3)(3x+2)$.\end{minipage}
\end{minipage}

\vfill
\begin{tcolorbox}[enhanced,nobeforeafter,width=7.1cm,height=2.1cm,valign=center,colback=clit!6,colframe=clit,arc=2.5mm,boxrule=0.9pt,left=3mm]{\Large\bfseries\color{clit}a)}\quad\begin{minipage}[c]{5.3cm}\raggedright\Large $(2x-3)(3x-8)$\end{minipage}\end{tcolorbox}\hfill\begin{tcolorbox}[enhanced,nobeforeafter,width=7.1cm,height=2.1cm,valign=center,colback=clit!6,colframe=clit,arc=2.5mm,boxrule=0.9pt,left=3mm]{\Large\bfseries\color{clit}b)}\quad\begin{minipage}[c]{5.3cm}\raggedright\Large $(2x-3)(4x-2)$\end{minipage}\end{tcolorbox}\par\vspace{2.5mm}
\begin{tcolorbox}[enhanced,nobeforeafter,width=7.1cm,height=2.1cm,valign=center,colback=clit!6,colframe=clit,arc=2.5mm,boxrule=0.9pt,left=3mm]{\Large\bfseries\color{clit}c)}\quad\begin{minipage}[c]{5.3cm}\raggedright\Large $(2x-3)(-2x-6)$\end{minipage}\end{tcolorbox}\hfill\begin{tcolorbox}[enhanced,nobeforeafter,width=7.1cm,height=2.1cm,valign=center,colback=clit!6,colframe=clit,arc=2.5mm,boxrule=0.9pt,left=3mm]{\Large\bfseries\color{clit}d)}\quad\begin{minipage}[c]{5.3cm}\raggedright\Large $(2x-3)^2(4x-2)$\end{minipage}\end{tcolorbox}
\end{tcolorbox}
\newpage
\begin{tcolorbox}[enhanced,width=16cm,height=14cm,colback=white,colframe=clit,boxrule=2pt,arc=0mm,
left=6mm,right=6mm,top=0mm,bottom=4mm,sharp corners,
overlay={\fill[clit] (frame.north west) rectangle ([yshift=-1.25cm]frame.north east);
\node[anchor=west,text=white,font=\sffamily\bfseries\LARGE] at ([xshift=5mm,yshift=-0.62cm]frame.north west) {Question 24 D};
\node[anchor=east,text=white,font=\sffamily\itshape\large] at ([xshift=-5mm,yshift=-0.62cm]frame.north east) {Calcul littéral et équations};}]
\vspace*{1.75cm}
\begin{minipage}[t][6.9cm][c]{\linewidth}\LARGE\centering
\begin{minipage}{0.95\linewidth}\raggedright\linespread{1.25}\selectfont Factoriser $(x-6)(5x+1)+(x-6)(2x-3)$.\end{minipage}
\end{minipage}

\vfill
\begin{tcolorbox}[enhanced,nobeforeafter,width=7.1cm,height=2.1cm,valign=center,colback=clit!6,colframe=clit,arc=2.5mm,boxrule=0.9pt,left=3mm]{\Large\bfseries\color{clit}a)}\quad\begin{minipage}[c]{5.3cm}\raggedright\Large $(x-6)(3x+4)$\end{minipage}\end{tcolorbox}\hfill\begin{tcolorbox}[enhanced,nobeforeafter,width=7.1cm,height=2.1cm,valign=center,colback=clit!6,colframe=clit,arc=2.5mm,boxrule=0.9pt,left=3mm]{\Large\bfseries\color{clit}b)}\quad\begin{minipage}[c]{5.3cm}\raggedright\Large $(x-6)(10x-3)$\end{minipage}\end{tcolorbox}\par\vspace{2.5mm}
\begin{tcolorbox}[enhanced,nobeforeafter,width=7.1cm,height=2.1cm,valign=center,colback=clit!6,colframe=clit,arc=2.5mm,boxrule=0.9pt,left=3mm]{\Large\bfseries\color{clit}c)}\quad\begin{minipage}[c]{5.3cm}\raggedright\Large $(x-6)^2(7x-2)$\end{minipage}\end{tcolorbox}\hfill\begin{tcolorbox}[enhanced,nobeforeafter,width=7.1cm,height=2.1cm,valign=center,colback=clit!6,colframe=clit,arc=2.5mm,boxrule=0.9pt,left=3mm]{\Large\bfseries\color{clit}d)}\quad\begin{minipage}[c]{5.3cm}\raggedright\Large $(x-6)(7x-2)$\end{minipage}\end{tcolorbox}
\end{tcolorbox}
\newpage
\begin{tcolorbox}[enhanced,width=16cm,height=14cm,colback=white,colframe=clit,boxrule=2pt,arc=0mm,
left=6mm,right=6mm,top=0mm,bottom=4mm,sharp corners,
overlay={\fill[clit] (frame.north west) rectangle ([yshift=-1.25cm]frame.north east);
\node[anchor=west,text=white,font=\sffamily\bfseries\LARGE] at ([xshift=5mm,yshift=-0.62cm]frame.north west) {Question 25 D};
\node[anchor=east,text=white,font=\sffamily\itshape\large] at ([xshift=-5mm,yshift=-0.62cm]frame.north east) {Calcul littéral et équations};}]
\vspace*{1.75cm}
\begin{minipage}[t][6.9cm][c]{\linewidth}\LARGE\centering
\begin{minipage}{0.95\linewidth}\raggedright\linespread{1.25}\selectfont Quelles sont les solutions de l'inéquation $6x-2>16$~?\end{minipage}
\end{minipage}

\vfill
\begin{tcolorbox}[enhanced,nobeforeafter,width=7.1cm,height=2.1cm,valign=center,colback=clit!6,colframe=clit,arc=2.5mm,boxrule=0.9pt,left=3mm]{\Large\bfseries\color{clit}a)}\quad\begin{minipage}[c]{5.3cm}\raggedright\Large $x<3$\end{minipage}\end{tcolorbox}\hfill\begin{tcolorbox}[enhanced,nobeforeafter,width=7.1cm,height=2.1cm,valign=center,colback=clit!6,colframe=clit,arc=2.5mm,boxrule=0.9pt,left=3mm]{\Large\bfseries\color{clit}b)}\quad\begin{minipage}[c]{5.3cm}\raggedright\Large $x>3$\end{minipage}\end{tcolorbox}\par\vspace{2.5mm}
\begin{tcolorbox}[enhanced,nobeforeafter,width=7.1cm,height=2.1cm,valign=center,colback=clit!6,colframe=clit,arc=2.5mm,boxrule=0.9pt,left=3mm]{\Large\bfseries\color{clit}c)}\quad\begin{minipage}[c]{5.3cm}\raggedright\Large $x>-3$\end{minipage}\end{tcolorbox}\hfill\begin{tcolorbox}[enhanced,nobeforeafter,width=7.1cm,height=2.1cm,valign=center,colback=clit!6,colframe=clit,arc=2.5mm,boxrule=0.9pt,left=3mm]{\Large\bfseries\color{clit}d)}\quad\begin{minipage}[c]{5.3cm}\raggedright\Large $x>\dfrac{7}{3}$\end{minipage}\end{tcolorbox}
\end{tcolorbox}
\newpage
\begin{tcolorbox}[enhanced,width=16cm,height=14cm,colback=white,colframe=clit,boxrule=2pt,arc=0mm,
left=6mm,right=6mm,top=0mm,bottom=4mm,sharp corners,
overlay={\fill[clit] (frame.north west) rectangle ([yshift=-1.25cm]frame.north east);
\node[anchor=west,text=white,font=\sffamily\bfseries\LARGE] at ([xshift=5mm,yshift=-0.62cm]frame.north west) {Question 26 D};
\node[anchor=east,text=white,font=\sffamily\itshape\large] at ([xshift=-5mm,yshift=-0.62cm]frame.north east) {Calcul littéral et équations};}]
\vspace*{1.75cm}
\begin{minipage}[t][6.9cm][c]{\linewidth}\LARGE\centering
\begin{minipage}{0.95\linewidth}\raggedright\linespread{1.25}\selectfont Quelles sont les solutions de l'inéquation $-3x+8>-1$~?\end{minipage}
\end{minipage}

\vfill
\begin{tcolorbox}[enhanced,nobeforeafter,width=7.1cm,height=2.1cm,valign=center,colback=clit!6,colframe=clit,arc=2.5mm,boxrule=0.9pt,left=3mm]{\Large\bfseries\color{clit}a)}\quad\begin{minipage}[c]{5.3cm}\raggedright\Large $x<3$\end{minipage}\end{tcolorbox}\hfill\begin{tcolorbox}[enhanced,nobeforeafter,width=7.1cm,height=2.1cm,valign=center,colback=clit!6,colframe=clit,arc=2.5mm,boxrule=0.9pt,left=3mm]{\Large\bfseries\color{clit}b)}\quad\begin{minipage}[c]{5.3cm}\raggedright\Large $x<-\dfrac{7}{3}$\end{minipage}\end{tcolorbox}\par\vspace{2.5mm}
\begin{tcolorbox}[enhanced,nobeforeafter,width=7.1cm,height=2.1cm,valign=center,colback=clit!6,colframe=clit,arc=2.5mm,boxrule=0.9pt,left=3mm]{\Large\bfseries\color{clit}c)}\quad\begin{minipage}[c]{5.3cm}\raggedright\Large $x>3$\end{minipage}\end{tcolorbox}\hfill\begin{tcolorbox}[enhanced,nobeforeafter,width=7.1cm,height=2.1cm,valign=center,colback=clit!6,colframe=clit,arc=2.5mm,boxrule=0.9pt,left=3mm]{\Large\bfseries\color{clit}d)}\quad\begin{minipage}[c]{5.3cm}\raggedright\Large $x<-3$\end{minipage}\end{tcolorbox}
\end{tcolorbox}
\newpage
\begin{tcolorbox}[enhanced,width=16cm,height=14cm,colback=white,colframe=clit,boxrule=2pt,arc=0mm,
left=6mm,right=6mm,top=0mm,bottom=4mm,sharp corners,
overlay={\fill[clit] (frame.north west) rectangle ([yshift=-1.25cm]frame.north east);
\node[anchor=west,text=white,font=\sffamily\bfseries\LARGE] at ([xshift=5mm,yshift=-0.62cm]frame.north west) {Question 27 D};
\node[anchor=east,text=white,font=\sffamily\itshape\large] at ([xshift=-5mm,yshift=-0.62cm]frame.north east) {Calcul littéral et équations};}]
\vspace*{1.75cm}
\begin{minipage}[t][6.9cm][c]{\linewidth}\LARGE\centering
\begin{minipage}{0.95\linewidth}\raggedright\linespread{1.25}\selectfont Quelles sont les solutions de l'équation $x^2=100$~?\end{minipage}
\end{minipage}

\vfill
\begin{tcolorbox}[enhanced,nobeforeafter,width=7.1cm,height=2.1cm,valign=center,colback=clit!6,colframe=clit,arc=2.5mm,boxrule=0.9pt,left=3mm]{\Large\bfseries\color{clit}a)}\quad\begin{minipage}[c]{5.3cm}\raggedright\Large aucune solution\end{minipage}\end{tcolorbox}\hfill\begin{tcolorbox}[enhanced,nobeforeafter,width=7.1cm,height=2.1cm,valign=center,colback=clit!6,colframe=clit,arc=2.5mm,boxrule=0.9pt,left=3mm]{\Large\bfseries\color{clit}b)}\quad\begin{minipage}[c]{5.3cm}\raggedright\Large $\{-10\,;\,10\}$\end{minipage}\end{tcolorbox}\par\vspace{2.5mm}
\begin{tcolorbox}[enhanced,nobeforeafter,width=7.1cm,height=2.1cm,valign=center,colback=clit!6,colframe=clit,arc=2.5mm,boxrule=0.9pt,left=3mm]{\Large\bfseries\color{clit}c)}\quad\begin{minipage}[c]{5.3cm}\raggedright\Large $\{10\}$\end{minipage}\end{tcolorbox}\hfill\begin{tcolorbox}[enhanced,nobeforeafter,width=7.1cm,height=2.1cm,valign=center,colback=clit!6,colframe=clit,arc=2.5mm,boxrule=0.9pt,left=3mm]{\Large\bfseries\color{clit}d)}\quad\begin{minipage}[c]{5.3cm}\raggedright\Large $\{-50\,;\,50\}$\end{minipage}\end{tcolorbox}
\end{tcolorbox}
\newpage
\begin{tcolorbox}[enhanced,width=16cm,height=14cm,colback=white,colframe=clit,boxrule=2pt,arc=0mm,
left=6mm,right=6mm,top=0mm,bottom=4mm,sharp corners,
overlay={\fill[clit] (frame.north west) rectangle ([yshift=-1.25cm]frame.north east);
\node[anchor=west,text=white,font=\sffamily\bfseries\LARGE] at ([xshift=5mm,yshift=-0.62cm]frame.north west) {Question 28 D};
\node[anchor=east,text=white,font=\sffamily\itshape\large] at ([xshift=-5mm,yshift=-0.62cm]frame.north east) {Calcul littéral et équations};}]
\vspace*{1.75cm}
\begin{minipage}[t][6.9cm][c]{\linewidth}\LARGE\centering
\begin{minipage}{0.95\linewidth}\raggedright\linespread{1.25}\selectfont Quelles sont les solutions de l'équation $x^2=-25$~?\end{minipage}
\end{minipage}

\vfill
\begin{tcolorbox}[enhanced,nobeforeafter,width=7.1cm,height=2.1cm,valign=center,colback=clit!6,colframe=clit,arc=2.5mm,boxrule=0.9pt,left=3mm]{\Large\bfseries\color{clit}a)}\quad\begin{minipage}[c]{5.3cm}\raggedright\Large $\{-5\}$\end{minipage}\end{tcolorbox}\hfill\begin{tcolorbox}[enhanced,nobeforeafter,width=7.1cm,height=2.1cm,valign=center,colback=clit!6,colframe=clit,arc=2.5mm,boxrule=0.9pt,left=3mm]{\Large\bfseries\color{clit}b)}\quad\begin{minipage}[c]{5.3cm}\raggedright\Large $\{-5\,;\,5\}$\end{minipage}\end{tcolorbox}\par\vspace{2.5mm}
\begin{tcolorbox}[enhanced,nobeforeafter,width=7.1cm,height=2.1cm,valign=center,colback=clit!6,colframe=clit,arc=2.5mm,boxrule=0.9pt,left=3mm]{\Large\bfseries\color{clit}c)}\quad\begin{minipage}[c]{5.3cm}\raggedright\Large $\{-25\}$\end{minipage}\end{tcolorbox}\hfill\begin{tcolorbox}[enhanced,nobeforeafter,width=7.1cm,height=2.1cm,valign=center,colback=clit!6,colframe=clit,arc=2.5mm,boxrule=0.9pt,left=3mm]{\Large\bfseries\color{clit}d)}\quad\begin{minipage}[c]{5.3cm}\raggedright\Large aucune solution\end{minipage}\end{tcolorbox}
\end{tcolorbox}
\newpage
\begin{tcolorbox}[enhanced,width=16cm,height=14cm,colback=white,colframe=clit,boxrule=2pt,arc=0mm,
left=6mm,right=6mm,top=0mm,bottom=4mm,sharp corners,
overlay={\fill[clit] (frame.north west) rectangle ([yshift=-1.25cm]frame.north east);
\node[anchor=west,text=white,font=\sffamily\bfseries\LARGE] at ([xshift=5mm,yshift=-0.62cm]frame.north west) {Question 29 D};
\node[anchor=east,text=white,font=\sffamily\itshape\large] at ([xshift=-5mm,yshift=-0.62cm]frame.north east) {Calcul littéral et équations};}]
\vspace*{1.75cm}
\begin{minipage}[t][6.9cm][c]{\linewidth}\LARGE\centering
\begin{minipage}{0.95\linewidth}\raggedright\linespread{1.25}\selectfont Le nombre $-3$ est-il solution de l'équation $x^2+2x-3=0$~?\end{minipage}
\end{minipage}

\vfill
\begin{tcolorbox}[enhanced,nobeforeafter,width=7.1cm,height=2.1cm,valign=center,colback=clit!6,colframe=clit,arc=2.5mm,boxrule=0.9pt,left=3mm]{\Large\bfseries\color{clit}a)}\quad\begin{minipage}[c]{5.3cm}\raggedright\Large non\end{minipage}\end{tcolorbox}\hfill\begin{tcolorbox}[enhanced,nobeforeafter,width=7.1cm,height=2.1cm,valign=center,colback=clit!6,colframe=clit,arc=2.5mm,boxrule=0.9pt,left=3mm]{\Large\bfseries\color{clit}b)}\quad\begin{minipage}[c]{5.3cm}\raggedright\large seulement si $x$ est positif\end{minipage}\end{tcolorbox}\par\vspace{2.5mm}
\begin{tcolorbox}[enhanced,nobeforeafter,width=7.1cm,height=2.1cm,valign=center,colback=clit!6,colframe=clit,arc=2.5mm,boxrule=0.9pt,left=3mm]{\Large\bfseries\color{clit}c)}\quad\begin{minipage}[c]{5.3cm}\raggedright\Large oui\end{minipage}\end{tcolorbox}\hfill\begin{tcolorbox}[enhanced,nobeforeafter,width=7.1cm,height=2.1cm,valign=center,colback=clit!6,colframe=clit,arc=2.5mm,boxrule=0.9pt,left=3mm]{\Large\bfseries\color{clit}d)}\quad\begin{minipage}[c]{5.3cm}\raggedright\Large on ne peut pas savoir\end{minipage}\end{tcolorbox}
\end{tcolorbox}
\newpage
\begin{tcolorbox}[enhanced,width=16cm,height=14cm,colback=white,colframe=clit,boxrule=2pt,arc=0mm,
left=6mm,right=6mm,top=0mm,bottom=4mm,sharp corners,
overlay={\fill[clit] (frame.north west) rectangle ([yshift=-1.25cm]frame.north east);
\node[anchor=west,text=white,font=\sffamily\bfseries\LARGE] at ([xshift=5mm,yshift=-0.62cm]frame.north west) {Question 30 D};
\node[anchor=east,text=white,font=\sffamily\itshape\large] at ([xshift=-5mm,yshift=-0.62cm]frame.north east) {Calcul littéral et équations};}]
\vspace*{1.75cm}
\begin{minipage}[t][6.9cm][c]{\linewidth}\LARGE\centering
\begin{minipage}{0.95\linewidth}\raggedright\linespread{1.25}\selectfont Le nombre $5$ est-il solution de l'équation $x^2-6x+4=0$~?\end{minipage}
\end{minipage}

\vfill
\begin{tcolorbox}[enhanced,nobeforeafter,width=7.1cm,height=2.1cm,valign=center,colback=clit!6,colframe=clit,arc=2.5mm,boxrule=0.9pt,left=3mm]{\Large\bfseries\color{clit}a)}\quad\begin{minipage}[c]{5.3cm}\raggedright\Large on ne peut pas savoir\end{minipage}\end{tcolorbox}\hfill\begin{tcolorbox}[enhanced,nobeforeafter,width=7.1cm,height=2.1cm,valign=center,colback=clit!6,colframe=clit,arc=2.5mm,boxrule=0.9pt,left=3mm]{\Large\bfseries\color{clit}b)}\quad\begin{minipage}[c]{5.3cm}\raggedright\large seulement si $x$ est positif\end{minipage}\end{tcolorbox}\par\vspace{2.5mm}
\begin{tcolorbox}[enhanced,nobeforeafter,width=7.1cm,height=2.1cm,valign=center,colback=clit!6,colframe=clit,arc=2.5mm,boxrule=0.9pt,left=3mm]{\Large\bfseries\color{clit}c)}\quad\begin{minipage}[c]{5.3cm}\raggedright\Large non\end{minipage}\end{tcolorbox}\hfill\begin{tcolorbox}[enhanced,nobeforeafter,width=7.1cm,height=2.1cm,valign=center,colback=clit!6,colframe=clit,arc=2.5mm,boxrule=0.9pt,left=3mm]{\Large\bfseries\color{clit}d)}\quad\begin{minipage}[c]{5.3cm}\raggedright\Large oui\end{minipage}\end{tcolorbox}
\end{tcolorbox}
\newpage
\begin{tcolorbox}[enhanced,width=16cm,height=14cm,colback=white,colframe=clit,boxrule=2pt,arc=0mm,
left=6mm,right=6mm,top=0mm,bottom=4mm,sharp corners,
overlay={\fill[clit] (frame.north west) rectangle ([yshift=-1.25cm]frame.north east);
\node[anchor=west,text=white,font=\sffamily\bfseries\LARGE] at ([xshift=5mm,yshift=-0.62cm]frame.north west) {Question 31 D};
\node[anchor=east,text=white,font=\sffamily\itshape\large] at ([xshift=-5mm,yshift=-0.62cm]frame.north east) {Calcul littéral et équations};}]
\vspace*{1.75cm}
\begin{minipage}[t][6.9cm][c]{\linewidth}\Large\centering
\begin{minipage}{0.95\linewidth}\raggedright\linespread{1.25}\selectfont Programme de calcul~: choisir un nombre ; ajouter $10$ ; multiplier le résultat par le nombre choisi ; ajouter $3$. Quel résultat obtient-on en choisissant $2$~?\end{minipage}
\end{minipage}

\vfill
\begin{tcolorbox}[enhanced,nobeforeafter,width=7.1cm,height=2.1cm,valign=center,colback=clit!6,colframe=clit,arc=2.5mm,boxrule=0.9pt,left=3mm]{\Large\bfseries\color{clit}a)}\quad\begin{minipage}[c]{5.3cm}\raggedright\Large $15$\end{minipage}\end{tcolorbox}\hfill\begin{tcolorbox}[enhanced,nobeforeafter,width=7.1cm,height=2.1cm,valign=center,colback=clit!6,colframe=clit,arc=2.5mm,boxrule=0.9pt,left=3mm]{\Large\bfseries\color{clit}b)}\quad\begin{minipage}[c]{5.3cm}\raggedright\Large $25$\end{minipage}\end{tcolorbox}\par\vspace{2.5mm}
\begin{tcolorbox}[enhanced,nobeforeafter,width=7.1cm,height=2.1cm,valign=center,colback=clit!6,colframe=clit,arc=2.5mm,boxrule=0.9pt,left=3mm]{\Large\bfseries\color{clit}c)}\quad\begin{minipage}[c]{5.3cm}\raggedright\Large $27$\end{minipage}\end{tcolorbox}\hfill\begin{tcolorbox}[enhanced,nobeforeafter,width=7.1cm,height=2.1cm,valign=center,colback=clit!6,colframe=clit,arc=2.5mm,boxrule=0.9pt,left=3mm]{\Large\bfseries\color{clit}d)}\quad\begin{minipage}[c]{5.3cm}\raggedright\Large $21$\end{minipage}\end{tcolorbox}
\end{tcolorbox}
\newpage
\begin{tcolorbox}[enhanced,width=16cm,height=14cm,colback=white,colframe=clit,boxrule=2pt,arc=0mm,
left=6mm,right=6mm,top=0mm,bottom=4mm,sharp corners,
overlay={\fill[clit] (frame.north west) rectangle ([yshift=-1.25cm]frame.north east);
\node[anchor=west,text=white,font=\sffamily\bfseries\LARGE] at ([xshift=5mm,yshift=-0.62cm]frame.north west) {Question 32 D};
\node[anchor=east,text=white,font=\sffamily\itshape\large] at ([xshift=-5mm,yshift=-0.62cm]frame.north east) {Calcul littéral et équations};}]
\vspace*{1.75cm}
\begin{minipage}[t][6.9cm][c]{\linewidth}\Large\centering
\begin{minipage}{0.95\linewidth}\raggedright\linespread{1.25}\selectfont Programme de calcul~: choisir un nombre ; ajouter $4$ ; multiplier le résultat par le nombre choisi ; soustraire $1$. Quel résultat obtient-on en choisissant $-5$~?\end{minipage}
\end{minipage}

\vfill
\begin{tcolorbox}[enhanced,nobeforeafter,width=7.1cm,height=2.1cm,valign=center,colback=clit!6,colframe=clit,arc=2.5mm,boxrule=0.9pt,left=3mm]{\Large\bfseries\color{clit}a)}\quad\begin{minipage}[c]{5.3cm}\raggedright\Large $6$\end{minipage}\end{tcolorbox}\hfill\begin{tcolorbox}[enhanced,nobeforeafter,width=7.1cm,height=2.1cm,valign=center,colback=clit!6,colframe=clit,arc=2.5mm,boxrule=0.9pt,left=3mm]{\Large\bfseries\color{clit}b)}\quad\begin{minipage}[c]{5.3cm}\raggedright\Large $-2$\end{minipage}\end{tcolorbox}\par\vspace{2.5mm}
\begin{tcolorbox}[enhanced,nobeforeafter,width=7.1cm,height=2.1cm,valign=center,colback=clit!6,colframe=clit,arc=2.5mm,boxrule=0.9pt,left=3mm]{\Large\bfseries\color{clit}c)}\quad\begin{minipage}[c]{5.3cm}\raggedright\Large $-26$\end{minipage}\end{tcolorbox}\hfill\begin{tcolorbox}[enhanced,nobeforeafter,width=7.1cm,height=2.1cm,valign=center,colback=clit!6,colframe=clit,arc=2.5mm,boxrule=0.9pt,left=3mm]{\Large\bfseries\color{clit}d)}\quad\begin{minipage}[c]{5.3cm}\raggedright\Large $4$\end{minipage}\end{tcolorbox}
\end{tcolorbox}
\newpage
\begin{tcolorbox}[enhanced,width=16cm,height=14cm,colback=white,colframe=cfon,boxrule=2pt,arc=0mm,
left=6mm,right=6mm,top=0mm,bottom=4mm,sharp corners,
overlay={\fill[cfon] (frame.north west) rectangle ([yshift=-1.25cm]frame.north east);
\node[anchor=west,text=white,font=\sffamily\bfseries\LARGE] at ([xshift=5mm,yshift=-0.62cm]frame.north west) {Question 33 D};
\node[anchor=east,text=white,font=\sffamily\itshape\large] at ([xshift=-5mm,yshift=-0.62cm]frame.north east) {Fonctions};}]
\vspace*{1.75cm}
\begin{minipage}[t][6.9cm][c]{\linewidth}\LARGE\centering
\begin{minipage}{0.95\linewidth}\raggedright\linespread{1.25}\selectfont Soit $f(x)=-3x+2$. Quelle est l'image de $6$ par $f$~?\end{minipage}
\end{minipage}

\vfill
\begin{tcolorbox}[enhanced,nobeforeafter,width=7.1cm,height=2.1cm,valign=center,colback=cfon!6,colframe=cfon,arc=2.5mm,boxrule=0.9pt,left=3mm]{\Large\bfseries\color{cfon}a)}\quad\begin{minipage}[c]{5.3cm}\raggedright\Large $5$\end{minipage}\end{tcolorbox}\hfill\begin{tcolorbox}[enhanced,nobeforeafter,width=7.1cm,height=2.1cm,valign=center,colback=cfon!6,colframe=cfon,arc=2.5mm,boxrule=0.9pt,left=3mm]{\Large\bfseries\color{cfon}b)}\quad\begin{minipage}[c]{5.3cm}\raggedright\Large $-\dfrac{4}{3}$\end{minipage}\end{tcolorbox}\par\vspace{2.5mm}
\begin{tcolorbox}[enhanced,nobeforeafter,width=7.1cm,height=2.1cm,valign=center,colback=cfon!6,colframe=cfon,arc=2.5mm,boxrule=0.9pt,left=3mm]{\Large\bfseries\color{cfon}c)}\quad\begin{minipage}[c]{5.3cm}\raggedright\Large $-20$\end{minipage}\end{tcolorbox}\hfill\begin{tcolorbox}[enhanced,nobeforeafter,width=7.1cm,height=2.1cm,valign=center,colback=cfon!6,colframe=cfon,arc=2.5mm,boxrule=0.9pt,left=3mm]{\Large\bfseries\color{cfon}d)}\quad\begin{minipage}[c]{5.3cm}\raggedright\Large $-16$\end{minipage}\end{tcolorbox}
\end{tcolorbox}
\newpage
\begin{tcolorbox}[enhanced,width=16cm,height=14cm,colback=white,colframe=cfon,boxrule=2pt,arc=0mm,
left=6mm,right=6mm,top=0mm,bottom=4mm,sharp corners,
overlay={\fill[cfon] (frame.north west) rectangle ([yshift=-1.25cm]frame.north east);
\node[anchor=west,text=white,font=\sffamily\bfseries\LARGE] at ([xshift=5mm,yshift=-0.62cm]frame.north west) {Question 34 D};
\node[anchor=east,text=white,font=\sffamily\itshape\large] at ([xshift=-5mm,yshift=-0.62cm]frame.north east) {Fonctions};}]
\vspace*{1.75cm}
\begin{minipage}[t][6.9cm][c]{\linewidth}\LARGE\centering
\begin{minipage}{0.95\linewidth}\raggedright\linespread{1.25}\selectfont Soit $f(x)=7x-4$. Quelle est l'image de $0$ par $f$~?\end{minipage}
\end{minipage}

\vfill
\begin{tcolorbox}[enhanced,nobeforeafter,width=7.1cm,height=2.1cm,valign=center,colback=cfon!6,colframe=cfon,arc=2.5mm,boxrule=0.9pt,left=3mm]{\Large\bfseries\color{cfon}a)}\quad\begin{minipage}[c]{5.3cm}\raggedright\Large $3$\end{minipage}\end{tcolorbox}\hfill\begin{tcolorbox}[enhanced,nobeforeafter,width=7.1cm,height=2.1cm,valign=center,colback=cfon!6,colframe=cfon,arc=2.5mm,boxrule=0.9pt,left=3mm]{\Large\bfseries\color{cfon}b)}\quad\begin{minipage}[c]{5.3cm}\raggedright\Large $4$\end{minipage}\end{tcolorbox}\par\vspace{2.5mm}
\begin{tcolorbox}[enhanced,nobeforeafter,width=7.1cm,height=2.1cm,valign=center,colback=cfon!6,colframe=cfon,arc=2.5mm,boxrule=0.9pt,left=3mm]{\Large\bfseries\color{cfon}c)}\quad\begin{minipage}[c]{5.3cm}\raggedright\Large $-4$\end{minipage}\end{tcolorbox}\hfill\begin{tcolorbox}[enhanced,nobeforeafter,width=7.1cm,height=2.1cm,valign=center,colback=cfon!6,colframe=cfon,arc=2.5mm,boxrule=0.9pt,left=3mm]{\Large\bfseries\color{cfon}d)}\quad\begin{minipage}[c]{5.3cm}\raggedright\Large $\dfrac{4}{7}$\end{minipage}\end{tcolorbox}
\end{tcolorbox}
\newpage
\begin{tcolorbox}[enhanced,width=16cm,height=14cm,colback=white,colframe=cfon,boxrule=2pt,arc=0mm,
left=6mm,right=6mm,top=0mm,bottom=4mm,sharp corners,
overlay={\fill[cfon] (frame.north west) rectangle ([yshift=-1.25cm]frame.north east);
\node[anchor=west,text=white,font=\sffamily\bfseries\LARGE] at ([xshift=5mm,yshift=-0.62cm]frame.north west) {Question 35 D};
\node[anchor=east,text=white,font=\sffamily\itshape\large] at ([xshift=-5mm,yshift=-0.62cm]frame.north east) {Fonctions};}]
\vspace*{1.75cm}
\begin{minipage}[t][6.9cm][c]{\linewidth}\LARGE\centering
\begin{minipage}{0.95\linewidth}\raggedright\linespread{1.25}\selectfont Soit $f(x)=6x-5$. Quel est l'antécédent de $13$ par $f$~?\end{minipage}
\end{minipage}

\vfill
\begin{tcolorbox}[enhanced,nobeforeafter,width=7.1cm,height=2.1cm,valign=center,colback=cfon!6,colframe=cfon,arc=2.5mm,boxrule=0.9pt,left=3mm]{\Large\bfseries\color{cfon}a)}\quad\begin{minipage}[c]{5.3cm}\raggedright\Large $73$\end{minipage}\end{tcolorbox}\hfill\begin{tcolorbox}[enhanced,nobeforeafter,width=7.1cm,height=2.1cm,valign=center,colback=cfon!6,colframe=cfon,arc=2.5mm,boxrule=0.9pt,left=3mm]{\Large\bfseries\color{cfon}b)}\quad\begin{minipage}[c]{5.3cm}\raggedright\Large $\dfrac{43}{6}$\end{minipage}\end{tcolorbox}\par\vspace{2.5mm}
\begin{tcolorbox}[enhanced,nobeforeafter,width=7.1cm,height=2.1cm,valign=center,colback=cfon!6,colframe=cfon,arc=2.5mm,boxrule=0.9pt,left=3mm]{\Large\bfseries\color{cfon}c)}\quad\begin{minipage}[c]{5.3cm}\raggedright\Large $\dfrac{4}{3}$\end{minipage}\end{tcolorbox}\hfill\begin{tcolorbox}[enhanced,nobeforeafter,width=7.1cm,height=2.1cm,valign=center,colback=cfon!6,colframe=cfon,arc=2.5mm,boxrule=0.9pt,left=3mm]{\Large\bfseries\color{cfon}d)}\quad\begin{minipage}[c]{5.3cm}\raggedright\Large $3$\end{minipage}\end{tcolorbox}
\end{tcolorbox}
\newpage
\begin{tcolorbox}[enhanced,width=16cm,height=14cm,colback=white,colframe=cfon,boxrule=2pt,arc=0mm,
left=6mm,right=6mm,top=0mm,bottom=4mm,sharp corners,
overlay={\fill[cfon] (frame.north west) rectangle ([yshift=-1.25cm]frame.north east);
\node[anchor=west,text=white,font=\sffamily\bfseries\LARGE] at ([xshift=5mm,yshift=-0.62cm]frame.north west) {Question 36 D};
\node[anchor=east,text=white,font=\sffamily\itshape\large] at ([xshift=-5mm,yshift=-0.62cm]frame.north east) {Fonctions};}]
\vspace*{1.75cm}
\begin{minipage}[t][6.9cm][c]{\linewidth}\LARGE\centering
\begin{minipage}{0.95\linewidth}\raggedright\linespread{1.25}\selectfont Soit $f(x)=-x+9$. Quel est l'antécédent de $2$ par $f$~?\end{minipage}
\end{minipage}

\vfill
\begin{tcolorbox}[enhanced,nobeforeafter,width=7.1cm,height=2.1cm,valign=center,colback=cfon!6,colframe=cfon,arc=2.5mm,boxrule=0.9pt,left=3mm]{\Large\bfseries\color{cfon}a)}\quad\begin{minipage}[c]{5.3cm}\raggedright\Large $7$\end{minipage}\end{tcolorbox}\hfill\begin{tcolorbox}[enhanced,nobeforeafter,width=7.1cm,height=2.1cm,valign=center,colback=cfon!6,colframe=cfon,arc=2.5mm,boxrule=0.9pt,left=3mm]{\Large\bfseries\color{cfon}b)}\quad\begin{minipage}[c]{5.3cm}\raggedright\Large $-11$\end{minipage}\end{tcolorbox}\par\vspace{2.5mm}
\begin{tcolorbox}[enhanced,nobeforeafter,width=7.1cm,height=2.1cm,valign=center,colback=cfon!6,colframe=cfon,arc=2.5mm,boxrule=0.9pt,left=3mm]{\Large\bfseries\color{cfon}c)}\quad\begin{minipage}[c]{5.3cm}\raggedright\Large $8$\end{minipage}\end{tcolorbox}\hfill\begin{tcolorbox}[enhanced,nobeforeafter,width=7.1cm,height=2.1cm,valign=center,colback=cfon!6,colframe=cfon,arc=2.5mm,boxrule=0.9pt,left=3mm]{\Large\bfseries\color{cfon}d)}\quad\begin{minipage}[c]{5.3cm}\raggedright\Large $6$\end{minipage}\end{tcolorbox}
\end{tcolorbox}
\newpage
\begin{tcolorbox}[enhanced,width=16cm,height=14cm,colback=white,colframe=cfon,boxrule=2pt,arc=0mm,
left=6mm,right=6mm,top=0mm,bottom=4mm,sharp corners,
overlay={\fill[cfon] (frame.north west) rectangle ([yshift=-1.25cm]frame.north east);
\node[anchor=west,text=white,font=\sffamily\bfseries\LARGE] at ([xshift=5mm,yshift=-0.62cm]frame.north west) {Question 37 D};
\node[anchor=east,text=white,font=\sffamily\itshape\large] at ([xshift=-5mm,yshift=-0.62cm]frame.north east) {Fonctions};}]
\vspace*{1.75cm}
\begin{minipage}[t][6.9cm][c]{\linewidth}\LARGE\centering
\begin{minipage}{0.95\linewidth}\raggedright\linespread{1.25}\selectfont $f$ est une fonction linéaire telle que $f(10)=4$. Quel est son coefficient~?\end{minipage}
\end{minipage}

\vfill
\begin{tcolorbox}[enhanced,nobeforeafter,width=7.1cm,height=2.1cm,valign=center,colback=cfon!6,colframe=cfon,arc=2.5mm,boxrule=0.9pt,left=3mm]{\Large\bfseries\color{cfon}a)}\quad\begin{minipage}[c]{5.3cm}\raggedright\Large $-6$\end{minipage}\end{tcolorbox}\hfill\begin{tcolorbox}[enhanced,nobeforeafter,width=7.1cm,height=2.1cm,valign=center,colback=cfon!6,colframe=cfon,arc=2.5mm,boxrule=0.9pt,left=3mm]{\Large\bfseries\color{cfon}b)}\quad\begin{minipage}[c]{5.3cm}\raggedright\Large $40$\end{minipage}\end{tcolorbox}\par\vspace{2.5mm}
\begin{tcolorbox}[enhanced,nobeforeafter,width=7.1cm,height=2.1cm,valign=center,colback=cfon!6,colframe=cfon,arc=2.5mm,boxrule=0.9pt,left=3mm]{\Large\bfseries\color{cfon}c)}\quad\begin{minipage}[c]{5.3cm}\raggedright\Large $\dfrac{2}{5}$\end{minipage}\end{tcolorbox}\hfill\begin{tcolorbox}[enhanced,nobeforeafter,width=7.1cm,height=2.1cm,valign=center,colback=cfon!6,colframe=cfon,arc=2.5mm,boxrule=0.9pt,left=3mm]{\Large\bfseries\color{cfon}d)}\quad\begin{minipage}[c]{5.3cm}\raggedright\Large $\dfrac{5}{2}$\end{minipage}\end{tcolorbox}
\end{tcolorbox}
\newpage
\begin{tcolorbox}[enhanced,width=16cm,height=14cm,colback=white,colframe=cfon,boxrule=2pt,arc=0mm,
left=6mm,right=6mm,top=0mm,bottom=4mm,sharp corners,
overlay={\fill[cfon] (frame.north west) rectangle ([yshift=-1.25cm]frame.north east);
\node[anchor=west,text=white,font=\sffamily\bfseries\LARGE] at ([xshift=5mm,yshift=-0.62cm]frame.north west) {Question 38 D};
\node[anchor=east,text=white,font=\sffamily\itshape\large] at ([xshift=-5mm,yshift=-0.62cm]frame.north east) {Fonctions};}]
\vspace*{1.75cm}
\begin{minipage}[t][6.9cm][c]{\linewidth}\LARGE\centering
\begin{minipage}{0.95\linewidth}\raggedright\linespread{1.25}\selectfont $f$ est une fonction linéaire telle que $f(6)=-9$. Quel est son coefficient~?\end{minipage}
\end{minipage}

\vfill
\begin{tcolorbox}[enhanced,nobeforeafter,width=7.1cm,height=2.1cm,valign=center,colback=cfon!6,colframe=cfon,arc=2.5mm,boxrule=0.9pt,left=3mm]{\Large\bfseries\color{cfon}a)}\quad\begin{minipage}[c]{5.3cm}\raggedright\Large $-\dfrac{2}{3}$\end{minipage}\end{tcolorbox}\hfill\begin{tcolorbox}[enhanced,nobeforeafter,width=7.1cm,height=2.1cm,valign=center,colback=cfon!6,colframe=cfon,arc=2.5mm,boxrule=0.9pt,left=3mm]{\Large\bfseries\color{cfon}b)}\quad\begin{minipage}[c]{5.3cm}\raggedright\Large $-\dfrac{3}{2}$\end{minipage}\end{tcolorbox}\par\vspace{2.5mm}
\begin{tcolorbox}[enhanced,nobeforeafter,width=7.1cm,height=2.1cm,valign=center,colback=cfon!6,colframe=cfon,arc=2.5mm,boxrule=0.9pt,left=3mm]{\Large\bfseries\color{cfon}c)}\quad\begin{minipage}[c]{5.3cm}\raggedright\Large $-54$\end{minipage}\end{tcolorbox}\hfill\begin{tcolorbox}[enhanced,nobeforeafter,width=7.1cm,height=2.1cm,valign=center,colback=cfon!6,colframe=cfon,arc=2.5mm,boxrule=0.9pt,left=3mm]{\Large\bfseries\color{cfon}d)}\quad\begin{minipage}[c]{5.3cm}\raggedright\Large $-15$\end{minipage}\end{tcolorbox}
\end{tcolorbox}
\newpage
\begin{tcolorbox}[enhanced,width=16cm,height=14cm,colback=white,colframe=cfon,boxrule=2pt,arc=0mm,
left=6mm,right=6mm,top=0mm,bottom=4mm,sharp corners,
overlay={\fill[cfon] (frame.north west) rectangle ([yshift=-1.25cm]frame.north east);
\node[anchor=west,text=white,font=\sffamily\bfseries\LARGE] at ([xshift=5mm,yshift=-0.62cm]frame.north west) {Question 39 D};
\node[anchor=east,text=white,font=\sffamily\itshape\large] at ([xshift=-5mm,yshift=-0.62cm]frame.north east) {Fonctions};}]
\vspace*{1.75cm}
\begin{minipage}[t][6.9cm][c]{\linewidth}\LARGE\centering
\begin{minipage}{0.95\linewidth}\raggedright\linespread{1.25}\selectfont $f$ est une fonction affine telle que $f(1)=0$ et $f(4)=9$. Quel est le coefficient $a$ de $f(x)=ax+b$~?\end{minipage}
\end{minipage}

\vfill
\begin{tcolorbox}[enhanced,nobeforeafter,width=7.1cm,height=2.1cm,valign=center,colback=cfon!6,colframe=cfon,arc=2.5mm,boxrule=0.9pt,left=3mm]{\Large\bfseries\color{cfon}a)}\quad\begin{minipage}[c]{5.3cm}\raggedright\Large $-3$\end{minipage}\end{tcolorbox}\hfill\begin{tcolorbox}[enhanced,nobeforeafter,width=7.1cm,height=2.1cm,valign=center,colback=cfon!6,colframe=cfon,arc=2.5mm,boxrule=0.9pt,left=3mm]{\Large\bfseries\color{cfon}b)}\quad\begin{minipage}[c]{5.3cm}\raggedright\Large $9$\end{minipage}\end{tcolorbox}\par\vspace{2.5mm}
\begin{tcolorbox}[enhanced,nobeforeafter,width=7.1cm,height=2.1cm,valign=center,colback=cfon!6,colframe=cfon,arc=2.5mm,boxrule=0.9pt,left=3mm]{\Large\bfseries\color{cfon}c)}\quad\begin{minipage}[c]{5.3cm}\raggedright\Large $\dfrac{1}{3}$\end{minipage}\end{tcolorbox}\hfill\begin{tcolorbox}[enhanced,nobeforeafter,width=7.1cm,height=2.1cm,valign=center,colback=cfon!6,colframe=cfon,arc=2.5mm,boxrule=0.9pt,left=3mm]{\Large\bfseries\color{cfon}d)}\quad\begin{minipage}[c]{5.3cm}\raggedright\Large $3$\end{minipage}\end{tcolorbox}
\end{tcolorbox}
\newpage
\begin{tcolorbox}[enhanced,width=16cm,height=14cm,colback=white,colframe=cfon,boxrule=2pt,arc=0mm,
left=6mm,right=6mm,top=0mm,bottom=4mm,sharp corners,
overlay={\fill[cfon] (frame.north west) rectangle ([yshift=-1.25cm]frame.north east);
\node[anchor=west,text=white,font=\sffamily\bfseries\LARGE] at ([xshift=5mm,yshift=-0.62cm]frame.north west) {Question 40 D};
\node[anchor=east,text=white,font=\sffamily\itshape\large] at ([xshift=-5mm,yshift=-0.62cm]frame.north east) {Fonctions};}]
\vspace*{1.75cm}
\begin{minipage}[t][6.9cm][c]{\linewidth}\LARGE\centering
\begin{minipage}{0.95\linewidth}\raggedright\linespread{1.25}\selectfont $f$ est une fonction affine telle que $f(3)=7$ et $f(5)=6$. Quel est le coefficient $a$ de $f(x)=ax+b$~?\end{minipage}
\end{minipage}

\vfill
\begin{tcolorbox}[enhanced,nobeforeafter,width=7.1cm,height=2.1cm,valign=center,colback=cfon!6,colframe=cfon,arc=2.5mm,boxrule=0.9pt,left=3mm]{\Large\bfseries\color{cfon}a)}\quad\begin{minipage}[c]{5.3cm}\raggedright\Large $-\dfrac{1}{2}$\end{minipage}\end{tcolorbox}\hfill\begin{tcolorbox}[enhanced,nobeforeafter,width=7.1cm,height=2.1cm,valign=center,colback=cfon!6,colframe=cfon,arc=2.5mm,boxrule=0.9pt,left=3mm]{\Large\bfseries\color{cfon}b)}\quad\begin{minipage}[c]{5.3cm}\raggedright\Large $-1$\end{minipage}\end{tcolorbox}\par\vspace{2.5mm}
\begin{tcolorbox}[enhanced,nobeforeafter,width=7.1cm,height=2.1cm,valign=center,colback=cfon!6,colframe=cfon,arc=2.5mm,boxrule=0.9pt,left=3mm]{\Large\bfseries\color{cfon}c)}\quad\begin{minipage}[c]{5.3cm}\raggedright\Large $-2$\end{minipage}\end{tcolorbox}\hfill\begin{tcolorbox}[enhanced,nobeforeafter,width=7.1cm,height=2.1cm,valign=center,colback=cfon!6,colframe=cfon,arc=2.5mm,boxrule=0.9pt,left=3mm]{\Large\bfseries\color{cfon}d)}\quad\begin{minipage}[c]{5.3cm}\raggedright\Large $\dfrac{1}{2}$\end{minipage}\end{tcolorbox}
\end{tcolorbox}
\newpage
\begin{tcolorbox}[enhanced,width=16cm,height=14cm,colback=white,colframe=cfon,boxrule=2pt,arc=0mm,
left=6mm,right=6mm,top=0mm,bottom=4mm,sharp corners,
overlay={\fill[cfon] (frame.north west) rectangle ([yshift=-1.25cm]frame.north east);
\node[anchor=west,text=white,font=\sffamily\bfseries\LARGE] at ([xshift=5mm,yshift=-0.62cm]frame.north west) {Question 41 D};
\node[anchor=east,text=white,font=\sffamily\itshape\large] at ([xshift=-5mm,yshift=-0.62cm]frame.north east) {Fonctions};}]
\vspace*{1.75cm}
\begin{minipage}[t][6.9cm][c]{\linewidth}\LARGE\centering
\begin{minipage}{0.95\linewidth}\raggedright\linespread{1.25}\selectfont $f$ est la fonction affine $f(x)=\dfrac{1}{2}x+b$ et $f(4)=6$. Que vaut $b$~?\end{minipage}
\end{minipage}

\vfill
\begin{tcolorbox}[enhanced,nobeforeafter,width=7.1cm,height=2.1cm,valign=center,colback=cfon!6,colframe=cfon,arc=2.5mm,boxrule=0.9pt,left=3mm]{\Large\bfseries\color{cfon}a)}\quad\begin{minipage}[c]{5.3cm}\raggedright\Large $4$\end{minipage}\end{tcolorbox}\hfill\begin{tcolorbox}[enhanced,nobeforeafter,width=7.1cm,height=2.1cm,valign=center,colback=cfon!6,colframe=cfon,arc=2.5mm,boxrule=0.9pt,left=3mm]{\Large\bfseries\color{cfon}b)}\quad\begin{minipage}[c]{5.3cm}\raggedright\Large $\dfrac{3}{2}$\end{minipage}\end{tcolorbox}\par\vspace{2.5mm}
\begin{tcolorbox}[enhanced,nobeforeafter,width=7.1cm,height=2.1cm,valign=center,colback=cfon!6,colframe=cfon,arc=2.5mm,boxrule=0.9pt,left=3mm]{\Large\bfseries\color{cfon}c)}\quad\begin{minipage}[c]{5.3cm}\raggedright\Large $8$\end{minipage}\end{tcolorbox}\hfill\begin{tcolorbox}[enhanced,nobeforeafter,width=7.1cm,height=2.1cm,valign=center,colback=cfon!6,colframe=cfon,arc=2.5mm,boxrule=0.9pt,left=3mm]{\Large\bfseries\color{cfon}d)}\quad\begin{minipage}[c]{5.3cm}\raggedright\Large $6$\end{minipage}\end{tcolorbox}
\end{tcolorbox}
\newpage
\begin{tcolorbox}[enhanced,width=16cm,height=14cm,colback=white,colframe=cfon,boxrule=2pt,arc=0mm,
left=6mm,right=6mm,top=0mm,bottom=4mm,sharp corners,
overlay={\fill[cfon] (frame.north west) rectangle ([yshift=-1.25cm]frame.north east);
\node[anchor=west,text=white,font=\sffamily\bfseries\LARGE] at ([xshift=5mm,yshift=-0.62cm]frame.north west) {Question 42 D};
\node[anchor=east,text=white,font=\sffamily\itshape\large] at ([xshift=-5mm,yshift=-0.62cm]frame.north east) {Fonctions};}]
\vspace*{1.75cm}
\begin{minipage}[t][6.9cm][c]{\linewidth}\LARGE\centering
\begin{minipage}{0.95\linewidth}\raggedright\linespread{1.25}\selectfont $f$ est la fonction affine $f(x)=-2x+b$ et $f(-3)=1$. Que vaut $b$~?\end{minipage}
\end{minipage}

\vfill
\begin{tcolorbox}[enhanced,nobeforeafter,width=7.1cm,height=2.1cm,valign=center,colback=cfon!6,colframe=cfon,arc=2.5mm,boxrule=0.9pt,left=3mm]{\Large\bfseries\color{cfon}a)}\quad\begin{minipage}[c]{5.3cm}\raggedright\Large $1$\end{minipage}\end{tcolorbox}\hfill\begin{tcolorbox}[enhanced,nobeforeafter,width=7.1cm,height=2.1cm,valign=center,colback=cfon!6,colframe=cfon,arc=2.5mm,boxrule=0.9pt,left=3mm]{\Large\bfseries\color{cfon}b)}\quad\begin{minipage}[c]{5.3cm}\raggedright\Large $7$\end{minipage}\end{tcolorbox}\par\vspace{2.5mm}
\begin{tcolorbox}[enhanced,nobeforeafter,width=7.1cm,height=2.1cm,valign=center,colback=cfon!6,colframe=cfon,arc=2.5mm,boxrule=0.9pt,left=3mm]{\Large\bfseries\color{cfon}c)}\quad\begin{minipage}[c]{5.3cm}\raggedright\Large $-\dfrac{1}{3}$\end{minipage}\end{tcolorbox}\hfill\begin{tcolorbox}[enhanced,nobeforeafter,width=7.1cm,height=2.1cm,valign=center,colback=cfon!6,colframe=cfon,arc=2.5mm,boxrule=0.9pt,left=3mm]{\Large\bfseries\color{cfon}d)}\quad\begin{minipage}[c]{5.3cm}\raggedright\Large $-5$\end{minipage}\end{tcolorbox}
\end{tcolorbox}
\newpage
\begin{tcolorbox}[enhanced,width=16cm,height=14cm,colback=white,colframe=cfon,boxrule=2pt,arc=0mm,
left=6mm,right=6mm,top=0mm,bottom=4mm,sharp corners,
overlay={\fill[cfon] (frame.north west) rectangle ([yshift=-1.25cm]frame.north east);
\node[anchor=west,text=white,font=\sffamily\bfseries\LARGE] at ([xshift=5mm,yshift=-0.62cm]frame.north west) {Question 43 D};
\node[anchor=east,text=white,font=\sffamily\itshape\large] at ([xshift=-5mm,yshift=-0.62cm]frame.north east) {Fonctions};}]
\vspace*{1.75cm}
\begin{minipage}[t][6.9cm][c]{\linewidth}\LARGE\centering
\begin{minipage}{0.95\linewidth}\raggedright\linespread{1.25}\selectfont Augmenter une quantité de $2\,\%$ revient à la multiplier par~:\end{minipage}
\end{minipage}

\vfill
\begin{tcolorbox}[enhanced,nobeforeafter,width=7.1cm,height=2.1cm,valign=center,colback=cfon!6,colframe=cfon,arc=2.5mm,boxrule=0.9pt,left=3mm]{\Large\bfseries\color{cfon}a)}\quad\begin{minipage}[c]{5.3cm}\raggedright\Large $1{,}02$\end{minipage}\end{tcolorbox}\hfill\begin{tcolorbox}[enhanced,nobeforeafter,width=7.1cm,height=2.1cm,valign=center,colback=cfon!6,colframe=cfon,arc=2.5mm,boxrule=0.9pt,left=3mm]{\Large\bfseries\color{cfon}b)}\quad\begin{minipage}[c]{5.3cm}\raggedright\Large $0{,}98$\end{minipage}\end{tcolorbox}\par\vspace{2.5mm}
\begin{tcolorbox}[enhanced,nobeforeafter,width=7.1cm,height=2.1cm,valign=center,colback=cfon!6,colframe=cfon,arc=2.5mm,boxrule=0.9pt,left=3mm]{\Large\bfseries\color{cfon}c)}\quad\begin{minipage}[c]{5.3cm}\raggedright\Large $1{,}2$\end{minipage}\end{tcolorbox}\hfill\begin{tcolorbox}[enhanced,nobeforeafter,width=7.1cm,height=2.1cm,valign=center,colback=cfon!6,colframe=cfon,arc=2.5mm,boxrule=0.9pt,left=3mm]{\Large\bfseries\color{cfon}d)}\quad\begin{minipage}[c]{5.3cm}\raggedright\Large $0{,}02$\end{minipage}\end{tcolorbox}
\end{tcolorbox}
\newpage
\begin{tcolorbox}[enhanced,width=16cm,height=14cm,colback=white,colframe=cfon,boxrule=2pt,arc=0mm,
left=6mm,right=6mm,top=0mm,bottom=4mm,sharp corners,
overlay={\fill[cfon] (frame.north west) rectangle ([yshift=-1.25cm]frame.north east);
\node[anchor=west,text=white,font=\sffamily\bfseries\LARGE] at ([xshift=5mm,yshift=-0.62cm]frame.north west) {Question 44 D};
\node[anchor=east,text=white,font=\sffamily\itshape\large] at ([xshift=-5mm,yshift=-0.62cm]frame.north east) {Fonctions};}]
\vspace*{1.75cm}
\begin{minipage}[t][6.9cm][c]{\linewidth}\LARGE\centering
\begin{minipage}{0.95\linewidth}\raggedright\linespread{1.25}\selectfont Diminuer une quantité de $60\,\%$ revient à la multiplier par~:\end{minipage}
\end{minipage}

\vfill
\begin{tcolorbox}[enhanced,nobeforeafter,width=7.1cm,height=2.1cm,valign=center,colback=cfon!6,colframe=cfon,arc=2.5mm,boxrule=0.9pt,left=3mm]{\Large\bfseries\color{cfon}a)}\quad\begin{minipage}[c]{5.3cm}\raggedright\Large $1{,}4$\end{minipage}\end{tcolorbox}\hfill\begin{tcolorbox}[enhanced,nobeforeafter,width=7.1cm,height=2.1cm,valign=center,colback=cfon!6,colframe=cfon,arc=2.5mm,boxrule=0.9pt,left=3mm]{\Large\bfseries\color{cfon}b)}\quad\begin{minipage}[c]{5.3cm}\raggedright\Large $1{,}6$\end{minipage}\end{tcolorbox}\par\vspace{2.5mm}
\begin{tcolorbox}[enhanced,nobeforeafter,width=7.1cm,height=2.1cm,valign=center,colback=cfon!6,colframe=cfon,arc=2.5mm,boxrule=0.9pt,left=3mm]{\Large\bfseries\color{cfon}c)}\quad\begin{minipage}[c]{5.3cm}\raggedright\Large $0{,}4$\end{minipage}\end{tcolorbox}\hfill\begin{tcolorbox}[enhanced,nobeforeafter,width=7.1cm,height=2.1cm,valign=center,colback=cfon!6,colframe=cfon,arc=2.5mm,boxrule=0.9pt,left=3mm]{\Large\bfseries\color{cfon}d)}\quad\begin{minipage}[c]{5.3cm}\raggedright\Large $0{,}6$\end{minipage}\end{tcolorbox}
\end{tcolorbox}
\newpage
\begin{tcolorbox}[enhanced,width=16cm,height=14cm,colback=white,colframe=cfon,boxrule=2pt,arc=0mm,
left=6mm,right=6mm,top=0mm,bottom=4mm,sharp corners,
overlay={\fill[cfon] (frame.north west) rectangle ([yshift=-1.25cm]frame.north east);
\node[anchor=west,text=white,font=\sffamily\bfseries\LARGE] at ([xshift=5mm,yshift=-0.62cm]frame.north west) {Question 45 D};
\node[anchor=east,text=white,font=\sffamily\itshape\large] at ([xshift=-5mm,yshift=-0.62cm]frame.north east) {Fonctions};}]
\vspace*{1.75cm}
\begin{minipage}[t][6.9cm][c]{\linewidth}\LARGE\centering
\begin{minipage}{0.95\linewidth}\raggedright\linespread{1.25}\selectfont Soit $g(x)=-2x^2+3x+1$. Que vaut $g(1)$~?\end{minipage}
\end{minipage}

\vfill
\begin{tcolorbox}[enhanced,nobeforeafter,width=7.1cm,height=2.1cm,valign=center,colback=cfon!6,colframe=cfon,arc=2.5mm,boxrule=0.9pt,left=3mm]{\Large\bfseries\color{cfon}a)}\quad\begin{minipage}[c]{5.3cm}\raggedright\Large $6$\end{minipage}\end{tcolorbox}\hfill\begin{tcolorbox}[enhanced,nobeforeafter,width=7.1cm,height=2.1cm,valign=center,colback=cfon!6,colframe=cfon,arc=2.5mm,boxrule=0.9pt,left=3mm]{\Large\bfseries\color{cfon}b)}\quad\begin{minipage}[c]{5.3cm}\raggedright\Large $0$\end{minipage}\end{tcolorbox}\par\vspace{2.5mm}
\begin{tcolorbox}[enhanced,nobeforeafter,width=7.1cm,height=2.1cm,valign=center,colback=cfon!6,colframe=cfon,arc=2.5mm,boxrule=0.9pt,left=3mm]{\Large\bfseries\color{cfon}c)}\quad\begin{minipage}[c]{5.3cm}\raggedright\Large $-2$\end{minipage}\end{tcolorbox}\hfill\begin{tcolorbox}[enhanced,nobeforeafter,width=7.1cm,height=2.1cm,valign=center,colback=cfon!6,colframe=cfon,arc=2.5mm,boxrule=0.9pt,left=3mm]{\Large\bfseries\color{cfon}d)}\quad\begin{minipage}[c]{5.3cm}\raggedright\Large $2$\end{minipage}\end{tcolorbox}
\end{tcolorbox}
\newpage
\begin{tcolorbox}[enhanced,width=16cm,height=14cm,colback=white,colframe=cfon,boxrule=2pt,arc=0mm,
left=6mm,right=6mm,top=0mm,bottom=4mm,sharp corners,
overlay={\fill[cfon] (frame.north west) rectangle ([yshift=-1.25cm]frame.north east);
\node[anchor=west,text=white,font=\sffamily\bfseries\LARGE] at ([xshift=5mm,yshift=-0.62cm]frame.north west) {Question 46 D};
\node[anchor=east,text=white,font=\sffamily\itshape\large] at ([xshift=-5mm,yshift=-0.62cm]frame.north east) {Fonctions};}]
\vspace*{1.75cm}
\begin{minipage}[t][6.9cm][c]{\linewidth}\LARGE\centering
\begin{minipage}{0.95\linewidth}\raggedright\linespread{1.25}\selectfont Soit $g(x)=x^2+x+1$. Que vaut $g(5)$~?\end{minipage}
\end{minipage}

\vfill
\begin{tcolorbox}[enhanced,nobeforeafter,width=7.1cm,height=2.1cm,valign=center,colback=cfon!6,colframe=cfon,arc=2.5mm,boxrule=0.9pt,left=3mm]{\Large\bfseries\color{cfon}a)}\quad\begin{minipage}[c]{5.3cm}\raggedright\Large $31$\end{minipage}\end{tcolorbox}\hfill\begin{tcolorbox}[enhanced,nobeforeafter,width=7.1cm,height=2.1cm,valign=center,colback=cfon!6,colframe=cfon,arc=2.5mm,boxrule=0.9pt,left=3mm]{\Large\bfseries\color{cfon}b)}\quad\begin{minipage}[c]{5.3cm}\raggedright\Large $16$\end{minipage}\end{tcolorbox}\par\vspace{2.5mm}
\begin{tcolorbox}[enhanced,nobeforeafter,width=7.1cm,height=2.1cm,valign=center,colback=cfon!6,colframe=cfon,arc=2.5mm,boxrule=0.9pt,left=3mm]{\Large\bfseries\color{cfon}c)}\quad\begin{minipage}[c]{5.3cm}\raggedright\Large $-19$\end{minipage}\end{tcolorbox}\hfill\begin{tcolorbox}[enhanced,nobeforeafter,width=7.1cm,height=2.1cm,valign=center,colback=cfon!6,colframe=cfon,arc=2.5mm,boxrule=0.9pt,left=3mm]{\Large\bfseries\color{cfon}d)}\quad\begin{minipage}[c]{5.3cm}\raggedright\Large $-31$\end{minipage}\end{tcolorbox}
\end{tcolorbox}
\newpage
\begin{tcolorbox}[enhanced,width=16cm,height=14cm,colback=white,colframe=cfon,boxrule=2pt,arc=0mm,
left=6mm,right=6mm,top=0mm,bottom=4mm,sharp corners,
overlay={\fill[cfon] (frame.north west) rectangle ([yshift=-1.25cm]frame.north east);
\node[anchor=west,text=white,font=\sffamily\bfseries\LARGE] at ([xshift=5mm,yshift=-0.62cm]frame.north west) {Question 47 D};
\node[anchor=east,text=white,font=\sffamily\itshape\large] at ([xshift=-5mm,yshift=-0.62cm]frame.north east) {Fonctions};}]
\vspace*{1.75cm}
\begin{minipage}[t][6.9cm][c]{\linewidth}\LARGE\centering
\begin{minipage}{0.95\linewidth}\raggedright\linespread{1.25}\selectfont De quel type est la fonction $f:x\mapsto \dfrac{1}{x}$~?\end{minipage}
\end{minipage}

\vfill
\begin{tcolorbox}[enhanced,nobeforeafter,width=7.1cm,height=2.1cm,valign=center,colback=cfon!6,colframe=cfon,arc=2.5mm,boxrule=0.9pt,left=3mm]{\Large\bfseries\color{cfon}a)}\quad\begin{minipage}[c]{5.3cm}\raggedright\Large affine non linéaire\end{minipage}\end{tcolorbox}\hfill\begin{tcolorbox}[enhanced,nobeforeafter,width=7.1cm,height=2.1cm,valign=center,colback=cfon!6,colframe=cfon,arc=2.5mm,boxrule=0.9pt,left=3mm]{\Large\bfseries\color{cfon}b)}\quad\begin{minipage}[c]{5.3cm}\raggedright\Large ni affine ni linéaire\end{minipage}\end{tcolorbox}\par\vspace{2.5mm}
\begin{tcolorbox}[enhanced,nobeforeafter,width=7.1cm,height=2.1cm,valign=center,colback=cfon!6,colframe=cfon,arc=2.5mm,boxrule=0.9pt,left=3mm]{\Large\bfseries\color{cfon}c)}\quad\begin{minipage}[c]{5.3cm}\raggedright\Large linéaire\end{minipage}\end{tcolorbox}\hfill\begin{tcolorbox}[enhanced,nobeforeafter,width=7.1cm,height=2.1cm,valign=center,colback=cfon!6,colframe=cfon,arc=2.5mm,boxrule=0.9pt,left=3mm]{\Large\bfseries\color{cfon}d)}\quad\begin{minipage}[c]{5.3cm}\raggedright\Large constante\end{minipage}\end{tcolorbox}
\end{tcolorbox}
\newpage
\begin{tcolorbox}[enhanced,width=16cm,height=14cm,colback=white,colframe=cfon,boxrule=2pt,arc=0mm,
left=6mm,right=6mm,top=0mm,bottom=4mm,sharp corners,
overlay={\fill[cfon] (frame.north west) rectangle ([yshift=-1.25cm]frame.north east);
\node[anchor=west,text=white,font=\sffamily\bfseries\LARGE] at ([xshift=5mm,yshift=-0.62cm]frame.north west) {Question 48 D};
\node[anchor=east,text=white,font=\sffamily\itshape\large] at ([xshift=-5mm,yshift=-0.62cm]frame.north east) {Fonctions};}]
\vspace*{1.75cm}
\begin{minipage}[t][6.9cm][c]{\linewidth}\LARGE\centering
\begin{minipage}{0.95\linewidth}\raggedright\linespread{1.25}\selectfont De quel type est la fonction $f:x\mapsto 2(x+1)$~?\end{minipage}
\end{minipage}

\vfill
\begin{tcolorbox}[enhanced,nobeforeafter,width=7.1cm,height=2.1cm,valign=center,colback=cfon!6,colframe=cfon,arc=2.5mm,boxrule=0.9pt,left=3mm]{\Large\bfseries\color{cfon}a)}\quad\begin{minipage}[c]{5.3cm}\raggedright\Large ni affine ni linéaire\end{minipage}\end{tcolorbox}\hfill\begin{tcolorbox}[enhanced,nobeforeafter,width=7.1cm,height=2.1cm,valign=center,colback=cfon!6,colframe=cfon,arc=2.5mm,boxrule=0.9pt,left=3mm]{\Large\bfseries\color{cfon}b)}\quad\begin{minipage}[c]{5.3cm}\raggedright\Large linéaire\end{minipage}\end{tcolorbox}\par\vspace{2.5mm}
\begin{tcolorbox}[enhanced,nobeforeafter,width=7.1cm,height=2.1cm,valign=center,colback=cfon!6,colframe=cfon,arc=2.5mm,boxrule=0.9pt,left=3mm]{\Large\bfseries\color{cfon}c)}\quad\begin{minipage}[c]{5.3cm}\raggedright\Large affine non linéaire\end{minipage}\end{tcolorbox}\hfill\begin{tcolorbox}[enhanced,nobeforeafter,width=7.1cm,height=2.1cm,valign=center,colback=cfon!6,colframe=cfon,arc=2.5mm,boxrule=0.9pt,left=3mm]{\Large\bfseries\color{cfon}d)}\quad\begin{minipage}[c]{5.3cm}\raggedright\Large constante\end{minipage}\end{tcolorbox}
\end{tcolorbox}
\newpage
\begin{tcolorbox}[enhanced,width=16cm,height=14cm,colback=white,colframe=cgeo,boxrule=2pt,arc=0mm,
left=6mm,right=6mm,top=0mm,bottom=4mm,sharp corners,
overlay={\fill[cgeo] (frame.north west) rectangle ([yshift=-1.25cm]frame.north east);
\node[anchor=west,text=white,font=\sffamily\bfseries\LARGE] at ([xshift=5mm,yshift=-0.62cm]frame.north west) {Question 49 D};
\node[anchor=east,text=white,font=\sffamily\itshape\large] at ([xshift=-5mm,yshift=-0.62cm]frame.north east) {Géométrie};}]
\vspace*{1.75cm}
\begin{minipage}[t][6.9cm][c]{\linewidth}\Large\centering
\begin{minipage}{0.95\linewidth}\raggedright\linespread{1.25}\selectfont Les droites $(AC)$ et $(BD)$ se coupent en $O$, et les droites $(AB)$ et $(CD)$ sont parallèles. $OA=8$ cm, $OC=6$ cm et $AB=10$ cm. Combien mesure $CD$~?\end{minipage}
\end{minipage}

\vfill
\begin{tcolorbox}[enhanced,nobeforeafter,width=7.1cm,height=2.1cm,valign=center,colback=cgeo!6,colframe=cgeo,arc=2.5mm,boxrule=0.9pt,left=3mm]{\Large\bfseries\color{cgeo}a)}\quad\begin{minipage}[c]{5.3cm}\raggedright\Large $4{,}8~\text{cm}$\end{minipage}\end{tcolorbox}\hfill\begin{tcolorbox}[enhanced,nobeforeafter,width=7.1cm,height=2.1cm,valign=center,colback=cgeo!6,colframe=cgeo,arc=2.5mm,boxrule=0.9pt,left=3mm]{\Large\bfseries\color{cgeo}b)}\quad\begin{minipage}[c]{5.3cm}\raggedright\Large $\dfrac{40}{3}~\text{cm}$\end{minipage}\end{tcolorbox}\par\vspace{2.5mm}
\begin{tcolorbox}[enhanced,nobeforeafter,width=7.1cm,height=2.1cm,valign=center,colback=cgeo!6,colframe=cgeo,arc=2.5mm,boxrule=0.9pt,left=3mm]{\Large\bfseries\color{cgeo}c)}\quad\begin{minipage}[c]{5.3cm}\raggedright\Large $7{,}5~\text{cm}$\end{minipage}\end{tcolorbox}\hfill\begin{tcolorbox}[enhanced,nobeforeafter,width=7.1cm,height=2.1cm,valign=center,colback=cgeo!6,colframe=cgeo,arc=2.5mm,boxrule=0.9pt,left=3mm]{\Large\bfseries\color{cgeo}d)}\quad\begin{minipage}[c]{5.3cm}\raggedright\Large $8~\text{cm}$\end{minipage}\end{tcolorbox}
\end{tcolorbox}
\newpage
\begin{tcolorbox}[enhanced,width=16cm,height=14cm,colback=white,colframe=cgeo,boxrule=2pt,arc=0mm,
left=6mm,right=6mm,top=0mm,bottom=4mm,sharp corners,
overlay={\fill[cgeo] (frame.north west) rectangle ([yshift=-1.25cm]frame.north east);
\node[anchor=west,text=white,font=\sffamily\bfseries\LARGE] at ([xshift=5mm,yshift=-0.62cm]frame.north west) {Question 50 D};
\node[anchor=east,text=white,font=\sffamily\itshape\large] at ([xshift=-5mm,yshift=-0.62cm]frame.north east) {Géométrie};}]
\vspace*{1.75cm}
\begin{minipage}[t][6.9cm][c]{\linewidth}\Large\centering
\begin{minipage}{0.95\linewidth}\raggedright\linespread{1.25}\selectfont Les droites $(AC)$ et $(BD)$ se coupent en $O$, et les droites $(AB)$ et $(CD)$ sont parallèles. $OA=3$ cm, $OC=4$ cm et $AB=9$ cm. Combien mesure $CD$~?\end{minipage}
\end{minipage}

\vfill
\begin{tcolorbox}[enhanced,nobeforeafter,width=7.1cm,height=2.1cm,valign=center,colback=cgeo!6,colframe=cgeo,arc=2.5mm,boxrule=0.9pt,left=3mm]{\Large\bfseries\color{cgeo}a)}\quad\begin{minipage}[c]{5.3cm}\raggedright\Large $10~\text{cm}$\end{minipage}\end{tcolorbox}\hfill\begin{tcolorbox}[enhanced,nobeforeafter,width=7.1cm,height=2.1cm,valign=center,colback=cgeo!6,colframe=cgeo,arc=2.5mm,boxrule=0.9pt,left=3mm]{\Large\bfseries\color{cgeo}b)}\quad\begin{minipage}[c]{5.3cm}\raggedright\Large $\dfrac{4}{3}~\text{cm}$\end{minipage}\end{tcolorbox}\par\vspace{2.5mm}
\begin{tcolorbox}[enhanced,nobeforeafter,width=7.1cm,height=2.1cm,valign=center,colback=cgeo!6,colframe=cgeo,arc=2.5mm,boxrule=0.9pt,left=3mm]{\Large\bfseries\color{cgeo}c)}\quad\begin{minipage}[c]{5.3cm}\raggedright\Large $6{,}75~\text{cm}$\end{minipage}\end{tcolorbox}\hfill\begin{tcolorbox}[enhanced,nobeforeafter,width=7.1cm,height=2.1cm,valign=center,colback=cgeo!6,colframe=cgeo,arc=2.5mm,boxrule=0.9pt,left=3mm]{\Large\bfseries\color{cgeo}d)}\quad\begin{minipage}[c]{5.3cm}\raggedright\Large $12~\text{cm}$\end{minipage}\end{tcolorbox}
\end{tcolorbox}
\newpage
\begin{tcolorbox}[enhanced,width=16cm,height=14cm,colback=white,colframe=cgeo,boxrule=2pt,arc=0mm,
left=6mm,right=6mm,top=0mm,bottom=4mm,sharp corners,
overlay={\fill[cgeo] (frame.north west) rectangle ([yshift=-1.25cm]frame.north east);
\node[anchor=west,text=white,font=\sffamily\bfseries\LARGE] at ([xshift=5mm,yshift=-0.62cm]frame.north west) {Question 51 D};
\node[anchor=east,text=white,font=\sffamily\itshape\large] at ([xshift=-5mm,yshift=-0.62cm]frame.north east) {Géométrie};}]
\vspace*{1.75cm}
\begin{minipage}[t][6.9cm][c]{\linewidth}\Large\centering
\begin{minipage}{0.95\linewidth}\raggedright\linespread{1.25}\selectfont Les points $A$, $M$, $B$ sont alignés dans cet ordre, ainsi que $A$, $N$, $C$. $AM=6$, $AB=8$, $AN=9$ et $AC=12$. Les droites $(MN)$ et $(BC)$ sont-elles parallèles~?\end{minipage}
\end{minipage}

\vfill
\begin{tcolorbox}[enhanced,nobeforeafter,width=7.1cm,height=2.1cm,valign=center,colback=cgeo!6,colframe=cgeo,arc=2.5mm,boxrule=0.9pt,left=3mm]{\Large\bfseries\color{cgeo}a)}\quad\begin{minipage}[c]{5.3cm}\raggedright\Large oui\end{minipage}\end{tcolorbox}\hfill\begin{tcolorbox}[enhanced,nobeforeafter,width=7.1cm,height=2.1cm,valign=center,colback=cgeo!6,colframe=cgeo,arc=2.5mm,boxrule=0.9pt,left=3mm]{\Large\bfseries\color{cgeo}b)}\quad\begin{minipage}[c]{5.3cm}\raggedright\Large non\end{minipage}\end{tcolorbox}\par\vspace{2.5mm}
\begin{tcolorbox}[enhanced,nobeforeafter,width=7.1cm,height=2.1cm,valign=center,colback=cgeo!6,colframe=cgeo,arc=2.5mm,boxrule=0.9pt,left=3mm]{\Large\bfseries\color{cgeo}c)}\quad\begin{minipage}[c]{5.3cm}\raggedright\large seulement si le triangle est rectangle\end{minipage}\end{tcolorbox}\hfill\begin{tcolorbox}[enhanced,nobeforeafter,width=7.1cm,height=2.1cm,valign=center,colback=cgeo!6,colframe=cgeo,arc=2.5mm,boxrule=0.9pt,left=3mm]{\Large\bfseries\color{cgeo}d)}\quad\begin{minipage}[c]{5.3cm}\raggedright\Large on ne peut pas savoir\end{minipage}\end{tcolorbox}
\end{tcolorbox}
\newpage
\begin{tcolorbox}[enhanced,width=16cm,height=14cm,colback=white,colframe=cgeo,boxrule=2pt,arc=0mm,
left=6mm,right=6mm,top=0mm,bottom=4mm,sharp corners,
overlay={\fill[cgeo] (frame.north west) rectangle ([yshift=-1.25cm]frame.north east);
\node[anchor=west,text=white,font=\sffamily\bfseries\LARGE] at ([xshift=5mm,yshift=-0.62cm]frame.north west) {Question 52 D};
\node[anchor=east,text=white,font=\sffamily\itshape\large] at ([xshift=-5mm,yshift=-0.62cm]frame.north east) {Géométrie};}]
\vspace*{1.75cm}
\begin{minipage}[t][6.9cm][c]{\linewidth}\Large\centering
\begin{minipage}{0.95\linewidth}\raggedright\linespread{1.25}\selectfont Les points $A$, $M$, $B$ sont alignés dans cet ordre, ainsi que $A$, $N$, $C$. $AM=5$, $AB=9$, $AN=4$ et $AC=8$. Les droites $(MN)$ et $(BC)$ sont-elles parallèles~?\end{minipage}
\end{minipage}

\vfill
\begin{tcolorbox}[enhanced,nobeforeafter,width=7.1cm,height=2.1cm,valign=center,colback=cgeo!6,colframe=cgeo,arc=2.5mm,boxrule=0.9pt,left=3mm]{\Large\bfseries\color{cgeo}a)}\quad\begin{minipage}[c]{5.3cm}\raggedright\large seulement si le triangle est rectangle\end{minipage}\end{tcolorbox}\hfill\begin{tcolorbox}[enhanced,nobeforeafter,width=7.1cm,height=2.1cm,valign=center,colback=cgeo!6,colframe=cgeo,arc=2.5mm,boxrule=0.9pt,left=3mm]{\Large\bfseries\color{cgeo}b)}\quad\begin{minipage}[c]{5.3cm}\raggedright\Large on ne peut pas savoir\end{minipage}\end{tcolorbox}\par\vspace{2.5mm}
\begin{tcolorbox}[enhanced,nobeforeafter,width=7.1cm,height=2.1cm,valign=center,colback=cgeo!6,colframe=cgeo,arc=2.5mm,boxrule=0.9pt,left=3mm]{\Large\bfseries\color{cgeo}c)}\quad\begin{minipage}[c]{5.3cm}\raggedright\Large oui\end{minipage}\end{tcolorbox}\hfill\begin{tcolorbox}[enhanced,nobeforeafter,width=7.1cm,height=2.1cm,valign=center,colback=cgeo!6,colframe=cgeo,arc=2.5mm,boxrule=0.9pt,left=3mm]{\Large\bfseries\color{cgeo}d)}\quad\begin{minipage}[c]{5.3cm}\raggedright\Large non\end{minipage}\end{tcolorbox}
\end{tcolorbox}
\newpage
\begin{tcolorbox}[enhanced,width=16cm,height=14cm,colback=white,colframe=cgeo,boxrule=2pt,arc=0mm,
left=6mm,right=6mm,top=0mm,bottom=4mm,sharp corners,
overlay={\fill[cgeo] (frame.north west) rectangle ([yshift=-1.25cm]frame.north east);
\node[anchor=west,text=white,font=\sffamily\bfseries\LARGE] at ([xshift=5mm,yshift=-0.62cm]frame.north west) {Question 53 D};
\node[anchor=east,text=white,font=\sffamily\itshape\large] at ([xshift=-5mm,yshift=-0.62cm]frame.north east) {Géométrie};}]
\vspace*{1.75cm}
\begin{minipage}[t][6.9cm][c]{\linewidth}\LARGE\centering
\begin{minipage}{0.95\linewidth}\raggedright\linespread{1.25}\selectfont Le triangle $ABC$ est rectangle en $A$, avec $BC=40$ cm et $\sin\widehat{ABC}=0{,}35$. Combien mesure $AC$~?\end{minipage}
\end{minipage}

\vfill
\begin{tcolorbox}[enhanced,nobeforeafter,width=7.1cm,height=2.1cm,valign=center,colback=cgeo!6,colframe=cgeo,arc=2.5mm,boxrule=0.9pt,left=3mm]{\Large\bfseries\color{cgeo}a)}\quad\begin{minipage}[c]{5.3cm}\raggedright\Large $14~\text{cm}$\end{minipage}\end{tcolorbox}\hfill\begin{tcolorbox}[enhanced,nobeforeafter,width=7.1cm,height=2.1cm,valign=center,colback=cgeo!6,colframe=cgeo,arc=2.5mm,boxrule=0.9pt,left=3mm]{\Large\bfseries\color{cgeo}b)}\quad\begin{minipage}[c]{5.3cm}\raggedright\Large $26~\text{cm}$\end{minipage}\end{tcolorbox}\par\vspace{2.5mm}
\begin{tcolorbox}[enhanced,nobeforeafter,width=7.1cm,height=2.1cm,valign=center,colback=cgeo!6,colframe=cgeo,arc=2.5mm,boxrule=0.9pt,left=3mm]{\Large\bfseries\color{cgeo}c)}\quad\begin{minipage}[c]{5.3cm}\raggedright\Large $28~\text{cm}$\end{minipage}\end{tcolorbox}\hfill\begin{tcolorbox}[enhanced,nobeforeafter,width=7.1cm,height=2.1cm,valign=center,colback=cgeo!6,colframe=cgeo,arc=2.5mm,boxrule=0.9pt,left=3mm]{\Large\bfseries\color{cgeo}d)}\quad\begin{minipage}[c]{5.3cm}\raggedright\Large $7~\text{cm}$\end{minipage}\end{tcolorbox}
\end{tcolorbox}
\newpage
\begin{tcolorbox}[enhanced,width=16cm,height=14cm,colback=white,colframe=cgeo,boxrule=2pt,arc=0mm,
left=6mm,right=6mm,top=0mm,bottom=4mm,sharp corners,
overlay={\fill[cgeo] (frame.north west) rectangle ([yshift=-1.25cm]frame.north east);
\node[anchor=west,text=white,font=\sffamily\bfseries\LARGE] at ([xshift=5mm,yshift=-0.62cm]frame.north west) {Question 54 D};
\node[anchor=east,text=white,font=\sffamily\itshape\large] at ([xshift=-5mm,yshift=-0.62cm]frame.north east) {Géométrie};}]
\vspace*{1.75cm}
\begin{minipage}[t][6.9cm][c]{\linewidth}\LARGE\centering
\begin{minipage}{0.95\linewidth}\raggedright\linespread{1.25}\selectfont Le triangle $ABC$ est rectangle en $A$, avec $BC=16$ cm et $\sin\widehat{ABC}=0{,}25$. Combien mesure $AC$~?\end{minipage}
\end{minipage}

\vfill
\begin{tcolorbox}[enhanced,nobeforeafter,width=7.1cm,height=2.1cm,valign=center,colback=cgeo!6,colframe=cgeo,arc=2.5mm,boxrule=0.9pt,left=3mm]{\Large\bfseries\color{cgeo}a)}\quad\begin{minipage}[c]{5.3cm}\raggedright\Large $12~\text{cm}$\end{minipage}\end{tcolorbox}\hfill\begin{tcolorbox}[enhanced,nobeforeafter,width=7.1cm,height=2.1cm,valign=center,colback=cgeo!6,colframe=cgeo,arc=2.5mm,boxrule=0.9pt,left=3mm]{\Large\bfseries\color{cgeo}b)}\quad\begin{minipage}[c]{5.3cm}\raggedright\Large $4~\text{cm}$\end{minipage}\end{tcolorbox}\par\vspace{2.5mm}
\begin{tcolorbox}[enhanced,nobeforeafter,width=7.1cm,height=2.1cm,valign=center,colback=cgeo!6,colframe=cgeo,arc=2.5mm,boxrule=0.9pt,left=3mm]{\Large\bfseries\color{cgeo}c)}\quad\begin{minipage}[c]{5.3cm}\raggedright\Large $8~\text{cm}$\end{minipage}\end{tcolorbox}\hfill\begin{tcolorbox}[enhanced,nobeforeafter,width=7.1cm,height=2.1cm,valign=center,colback=cgeo!6,colframe=cgeo,arc=2.5mm,boxrule=0.9pt,left=3mm]{\Large\bfseries\color{cgeo}d)}\quad\begin{minipage}[c]{5.3cm}\raggedright\Large $2~\text{cm}$\end{minipage}\end{tcolorbox}
\end{tcolorbox}
\newpage
\begin{tcolorbox}[enhanced,width=16cm,height=14cm,colback=white,colframe=cgeo,boxrule=2pt,arc=0mm,
left=6mm,right=6mm,top=0mm,bottom=4mm,sharp corners,
overlay={\fill[cgeo] (frame.north west) rectangle ([yshift=-1.25cm]frame.north east);
\node[anchor=west,text=white,font=\sffamily\bfseries\LARGE] at ([xshift=5mm,yshift=-0.62cm]frame.north west) {Question 55 D};
\node[anchor=east,text=white,font=\sffamily\itshape\large] at ([xshift=-5mm,yshift=-0.62cm]frame.north east) {Géométrie};}]
\vspace*{1.75cm}
\begin{minipage}[t][6.9cm][c]{\linewidth}\LARGE\centering
\begin{minipage}{0.95\linewidth}\raggedright\linespread{1.25}\selectfont Le triangle $ABC$ est rectangle en $A$, avec $AB=6$ cm et $AC=8$ cm. Que vaut $\tan\widehat{ABC}$~?\end{minipage}
\end{minipage}

\vfill
\begin{tcolorbox}[enhanced,nobeforeafter,width=7.1cm,height=2.1cm,valign=center,colback=cgeo!6,colframe=cgeo,arc=2.5mm,boxrule=0.9pt,left=3mm]{\Large\bfseries\color{cgeo}a)}\quad\begin{minipage}[c]{5.3cm}\raggedright\Large $\dfrac{4}{3}$\end{minipage}\end{tcolorbox}\hfill\begin{tcolorbox}[enhanced,nobeforeafter,width=7.1cm,height=2.1cm,valign=center,colback=cgeo!6,colframe=cgeo,arc=2.5mm,boxrule=0.9pt,left=3mm]{\Large\bfseries\color{cgeo}b)}\quad\begin{minipage}[c]{5.3cm}\raggedright\Large $\dfrac{7}{3}$\end{minipage}\end{tcolorbox}\par\vspace{2.5mm}
\begin{tcolorbox}[enhanced,nobeforeafter,width=7.1cm,height=2.1cm,valign=center,colback=cgeo!6,colframe=cgeo,arc=2.5mm,boxrule=0.9pt,left=3mm]{\Large\bfseries\color{cgeo}c)}\quad\begin{minipage}[c]{5.3cm}\raggedright\Large $\dfrac{3}{4}$\end{minipage}\end{tcolorbox}\hfill\begin{tcolorbox}[enhanced,nobeforeafter,width=7.1cm,height=2.1cm,valign=center,colback=cgeo!6,colframe=cgeo,arc=2.5mm,boxrule=0.9pt,left=3mm]{\Large\bfseries\color{cgeo}d)}\quad\begin{minipage}[c]{5.3cm}\raggedright\Large $\dfrac{3}{7}$\end{minipage}\end{tcolorbox}
\end{tcolorbox}
\newpage
\begin{tcolorbox}[enhanced,width=16cm,height=14cm,colback=white,colframe=cgeo,boxrule=2pt,arc=0mm,
left=6mm,right=6mm,top=0mm,bottom=4mm,sharp corners,
overlay={\fill[cgeo] (frame.north west) rectangle ([yshift=-1.25cm]frame.north east);
\node[anchor=west,text=white,font=\sffamily\bfseries\LARGE] at ([xshift=5mm,yshift=-0.62cm]frame.north west) {Question 56 D};
\node[anchor=east,text=white,font=\sffamily\itshape\large] at ([xshift=-5mm,yshift=-0.62cm]frame.north east) {Géométrie};}]
\vspace*{1.75cm}
\begin{minipage}[t][6.9cm][c]{\linewidth}\LARGE\centering
\begin{minipage}{0.95\linewidth}\raggedright\linespread{1.25}\selectfont Le triangle $ABC$ est rectangle en $A$, avec $AB=15$ cm et $AC=5$ cm. Que vaut $\tan\widehat{ABC}$~?\end{minipage}
\end{minipage}

\vfill
\begin{tcolorbox}[enhanced,nobeforeafter,width=7.1cm,height=2.1cm,valign=center,colback=cgeo!6,colframe=cgeo,arc=2.5mm,boxrule=0.9pt,left=3mm]{\Large\bfseries\color{cgeo}a)}\quad\begin{minipage}[c]{5.3cm}\raggedright\Large $\dfrac{4}{3}$\end{minipage}\end{tcolorbox}\hfill\begin{tcolorbox}[enhanced,nobeforeafter,width=7.1cm,height=2.1cm,valign=center,colback=cgeo!6,colframe=cgeo,arc=2.5mm,boxrule=0.9pt,left=3mm]{\Large\bfseries\color{cgeo}b)}\quad\begin{minipage}[c]{5.3cm}\raggedright\Large $\dfrac{1}{3}$\end{minipage}\end{tcolorbox}\par\vspace{2.5mm}
\begin{tcolorbox}[enhanced,nobeforeafter,width=7.1cm,height=2.1cm,valign=center,colback=cgeo!6,colframe=cgeo,arc=2.5mm,boxrule=0.9pt,left=3mm]{\Large\bfseries\color{cgeo}c)}\quad\begin{minipage}[c]{5.3cm}\raggedright\Large $\dfrac{3}{4}$\end{minipage}\end{tcolorbox}\hfill\begin{tcolorbox}[enhanced,nobeforeafter,width=7.1cm,height=2.1cm,valign=center,colback=cgeo!6,colframe=cgeo,arc=2.5mm,boxrule=0.9pt,left=3mm]{\Large\bfseries\color{cgeo}d)}\quad\begin{minipage}[c]{5.3cm}\raggedright\Large $3$\end{minipage}\end{tcolorbox}
\end{tcolorbox}
\newpage
\begin{tcolorbox}[enhanced,width=16cm,height=14cm,colback=white,colframe=cgeo,boxrule=2pt,arc=0mm,
left=6mm,right=6mm,top=0mm,bottom=4mm,sharp corners,
overlay={\fill[cgeo] (frame.north west) rectangle ([yshift=-1.25cm]frame.north east);
\node[anchor=west,text=white,font=\sffamily\bfseries\LARGE] at ([xshift=5mm,yshift=-0.62cm]frame.north west) {Question 57 D};
\node[anchor=east,text=white,font=\sffamily\itshape\large] at ([xshift=-5mm,yshift=-0.62cm]frame.north east) {Géométrie};}]
\vspace*{1.75cm}
\begin{minipage}[t][6.9cm][c]{\linewidth}\LARGE\centering
\begin{minipage}{0.95\linewidth}\raggedright\linespread{1.25}\selectfont Une homothétie de rapport $-1$ transforme le segment $[AB]$, de longueur $7$ cm, en $[A'B']$. Combien mesure $A'B'$~?\end{minipage}
\end{minipage}

\vfill
\begin{tcolorbox}[enhanced,nobeforeafter,width=7.1cm,height=2.1cm,valign=center,colback=cgeo!6,colframe=cgeo,arc=2.5mm,boxrule=0.9pt,left=3mm]{\Large\bfseries\color{cgeo}a)}\quad\begin{minipage}[c]{5.3cm}\raggedright\Large $14~\text{cm}$\end{minipage}\end{tcolorbox}\hfill\begin{tcolorbox}[enhanced,nobeforeafter,width=7.1cm,height=2.1cm,valign=center,colback=cgeo!6,colframe=cgeo,arc=2.5mm,boxrule=0.9pt,left=3mm]{\Large\bfseries\color{cgeo}b)}\quad\begin{minipage}[c]{5.3cm}\raggedright\Large $7~\text{cm}$\end{minipage}\end{tcolorbox}\par\vspace{2.5mm}
\begin{tcolorbox}[enhanced,nobeforeafter,width=7.1cm,height=2.1cm,valign=center,colback=cgeo!6,colframe=cgeo,arc=2.5mm,boxrule=0.9pt,left=3mm]{\Large\bfseries\color{cgeo}c)}\quad\begin{minipage}[c]{5.3cm}\raggedright\Large $9~\text{cm}$\end{minipage}\end{tcolorbox}\hfill\begin{tcolorbox}[enhanced,nobeforeafter,width=7.1cm,height=2.1cm,valign=center,colback=cgeo!6,colframe=cgeo,arc=2.5mm,boxrule=0.9pt,left=3mm]{\Large\bfseries\color{cgeo}d)}\quad\begin{minipage}[c]{5.3cm}\raggedright\Large $8~\text{cm}$\end{minipage}\end{tcolorbox}
\end{tcolorbox}
\newpage
\begin{tcolorbox}[enhanced,width=16cm,height=14cm,colback=white,colframe=cgeo,boxrule=2pt,arc=0mm,
left=6mm,right=6mm,top=0mm,bottom=4mm,sharp corners,
overlay={\fill[cgeo] (frame.north west) rectangle ([yshift=-1.25cm]frame.north east);
\node[anchor=west,text=white,font=\sffamily\bfseries\LARGE] at ([xshift=5mm,yshift=-0.62cm]frame.north west) {Question 58 D};
\node[anchor=east,text=white,font=\sffamily\itshape\large] at ([xshift=-5mm,yshift=-0.62cm]frame.north east) {Géométrie};}]
\vspace*{1.75cm}
\begin{minipage}[t][6.9cm][c]{\linewidth}\LARGE\centering
\begin{minipage}{0.95\linewidth}\raggedright\linespread{1.25}\selectfont Une homothétie de rapport $\dfrac{3}{2}$ transforme le segment $[AB]$, de longueur $10$ cm, en $[A'B']$. Combien mesure $A'B'$~?\end{minipage}
\end{minipage}

\vfill
\begin{tcolorbox}[enhanced,nobeforeafter,width=7.1cm,height=2.1cm,valign=center,colback=cgeo!6,colframe=cgeo,arc=2.5mm,boxrule=0.9pt,left=3mm]{\Large\bfseries\color{cgeo}a)}\quad\begin{minipage}[c]{5.3cm}\raggedright\Large $16~\text{cm}$\end{minipage}\end{tcolorbox}\hfill\begin{tcolorbox}[enhanced,nobeforeafter,width=7.1cm,height=2.1cm,valign=center,colback=cgeo!6,colframe=cgeo,arc=2.5mm,boxrule=0.9pt,left=3mm]{\Large\bfseries\color{cgeo}b)}\quad\begin{minipage}[c]{5.3cm}\raggedright\Large $15~\text{cm}$\end{minipage}\end{tcolorbox}\par\vspace{2.5mm}
\begin{tcolorbox}[enhanced,nobeforeafter,width=7.1cm,height=2.1cm,valign=center,colback=cgeo!6,colframe=cgeo,arc=2.5mm,boxrule=0.9pt,left=3mm]{\Large\bfseries\color{cgeo}c)}\quad\begin{minipage}[c]{5.3cm}\raggedright\Large $\dfrac{20}{3}~\text{cm}$\end{minipage}\end{tcolorbox}\hfill\begin{tcolorbox}[enhanced,nobeforeafter,width=7.1cm,height=2.1cm,valign=center,colback=cgeo!6,colframe=cgeo,arc=2.5mm,boxrule=0.9pt,left=3mm]{\Large\bfseries\color{cgeo}d)}\quad\begin{minipage}[c]{5.3cm}\raggedright\Large $11{,}5~\text{cm}$\end{minipage}\end{tcolorbox}
\end{tcolorbox}
\newpage
\begin{tcolorbox}[enhanced,width=16cm,height=14cm,colback=white,colframe=cgeo,boxrule=2pt,arc=0mm,
left=6mm,right=6mm,top=0mm,bottom=4mm,sharp corners,
overlay={\fill[cgeo] (frame.north west) rectangle ([yshift=-1.25cm]frame.north east);
\node[anchor=west,text=white,font=\sffamily\bfseries\LARGE] at ([xshift=5mm,yshift=-0.62cm]frame.north west) {Question 59 D};
\node[anchor=east,text=white,font=\sffamily\itshape\large] at ([xshift=-5mm,yshift=-0.62cm]frame.north east) {Géométrie};}]
\vspace*{1.75cm}
\begin{minipage}[t][6.9cm][c]{\linewidth}\LARGE\centering
\begin{minipage}{0.95\linewidth}\raggedright\linespread{1.25}\selectfont Quel est le volume exact d'une boule de rayon $12$ cm~?\end{minipage}
\end{minipage}

\vfill
\begin{tcolorbox}[enhanced,nobeforeafter,width=7.1cm,height=2.1cm,valign=center,colback=cgeo!6,colframe=cgeo,arc=2.5mm,boxrule=0.9pt,left=3mm]{\Large\bfseries\color{cgeo}a)}\quad\begin{minipage}[c]{5.3cm}\raggedright\Large $576\pi~\text{cm}^3$\end{minipage}\end{tcolorbox}\hfill\begin{tcolorbox}[enhanced,nobeforeafter,width=7.1cm,height=2.1cm,valign=center,colback=cgeo!6,colframe=cgeo,arc=2.5mm,boxrule=0.9pt,left=3mm]{\Large\bfseries\color{cgeo}b)}\quad\begin{minipage}[c]{5.3cm}\raggedright\Large $1\,728\pi~\text{cm}^3$\end{minipage}\end{tcolorbox}\par\vspace{2.5mm}
\begin{tcolorbox}[enhanced,nobeforeafter,width=7.1cm,height=2.1cm,valign=center,colback=cgeo!6,colframe=cgeo,arc=2.5mm,boxrule=0.9pt,left=3mm]{\Large\bfseries\color{cgeo}c)}\quad\begin{minipage}[c]{5.3cm}\raggedright\Large $192\pi~\text{cm}^3$\end{minipage}\end{tcolorbox}\hfill\begin{tcolorbox}[enhanced,nobeforeafter,width=7.1cm,height=2.1cm,valign=center,colback=cgeo!6,colframe=cgeo,arc=2.5mm,boxrule=0.9pt,left=3mm]{\Large\bfseries\color{cgeo}d)}\quad\begin{minipage}[c]{5.3cm}\raggedright\Large $2\,304\pi~\text{cm}^3$\end{minipage}\end{tcolorbox}
\end{tcolorbox}
\newpage
\begin{tcolorbox}[enhanced,width=16cm,height=14cm,colback=white,colframe=cgeo,boxrule=2pt,arc=0mm,
left=6mm,right=6mm,top=0mm,bottom=4mm,sharp corners,
overlay={\fill[cgeo] (frame.north west) rectangle ([yshift=-1.25cm]frame.north east);
\node[anchor=west,text=white,font=\sffamily\bfseries\LARGE] at ([xshift=5mm,yshift=-0.62cm]frame.north west) {Question 60 D};
\node[anchor=east,text=white,font=\sffamily\itshape\large] at ([xshift=-5mm,yshift=-0.62cm]frame.north east) {Géométrie};}]
\vspace*{1.75cm}
\begin{minipage}[t][6.9cm][c]{\linewidth}\LARGE\centering
\begin{minipage}{0.95\linewidth}\raggedright\linespread{1.25}\selectfont Quel est le volume exact d'une boule de rayon $0{,}75$ cm~?\end{minipage}
\end{minipage}

\vfill
\begin{tcolorbox}[enhanced,nobeforeafter,width=7.1cm,height=2.1cm,valign=center,colback=cgeo!6,colframe=cgeo,arc=2.5mm,boxrule=0.9pt,left=3mm]{\Large\bfseries\color{cgeo}a)}\quad\begin{minipage}[c]{5.3cm}\raggedright\large $\dfrac{9}{16}\pi~\text{cm}^3$\end{minipage}\end{tcolorbox}\hfill\begin{tcolorbox}[enhanced,nobeforeafter,width=7.1cm,height=2.1cm,valign=center,colback=cgeo!6,colframe=cgeo,arc=2.5mm,boxrule=0.9pt,left=3mm]{\Large\bfseries\color{cgeo}b)}\quad\begin{minipage}[c]{5.3cm}\raggedright\large $\dfrac{3}{4}\pi~\text{cm}^3$\end{minipage}\end{tcolorbox}\par\vspace{2.5mm}
\begin{tcolorbox}[enhanced,nobeforeafter,width=7.1cm,height=2.1cm,valign=center,colback=cgeo!6,colframe=cgeo,arc=2.5mm,boxrule=0.9pt,left=3mm]{\Large\bfseries\color{cgeo}c)}\quad\begin{minipage}[c]{5.3cm}\raggedright\large $\dfrac{27}{64}\pi~\text{cm}^3$\end{minipage}\end{tcolorbox}\hfill\begin{tcolorbox}[enhanced,nobeforeafter,width=7.1cm,height=2.1cm,valign=center,colback=cgeo!6,colframe=cgeo,arc=2.5mm,boxrule=0.9pt,left=3mm]{\Large\bfseries\color{cgeo}d)}\quad\begin{minipage}[c]{5.3cm}\raggedright\large $\dfrac{9}{4}\pi~\text{cm}^3$\end{minipage}\end{tcolorbox}
\end{tcolorbox}
\newpage
\begin{tcolorbox}[enhanced,width=16cm,height=14cm,colback=white,colframe=cgeo,boxrule=2pt,arc=0mm,
left=6mm,right=6mm,top=0mm,bottom=4mm,sharp corners,
overlay={\fill[cgeo] (frame.north west) rectangle ([yshift=-1.25cm]frame.north east);
\node[anchor=west,text=white,font=\sffamily\bfseries\LARGE] at ([xshift=5mm,yshift=-0.62cm]frame.north west) {Question 61 D};
\node[anchor=east,text=white,font=\sffamily\itshape\large] at ([xshift=-5mm,yshift=-0.62cm]frame.north east) {Géométrie};}]
\vspace*{1.75cm}
\begin{minipage}[t][6.9cm][c]{\linewidth}\LARGE\centering
\begin{minipage}{0.95\linewidth}\raggedright\linespread{1.25}\selectfont Quelle est l'aire exacte d'une sphère de rayon $6$ cm~?\end{minipage}
\end{minipage}

\vfill
\begin{tcolorbox}[enhanced,nobeforeafter,width=7.1cm,height=2.1cm,valign=center,colback=cgeo!6,colframe=cgeo,arc=2.5mm,boxrule=0.9pt,left=3mm]{\Large\bfseries\color{cgeo}a)}\quad\begin{minipage}[c]{5.3cm}\raggedright\Large $288\pi~\text{cm}^2$\end{minipage}\end{tcolorbox}\hfill\begin{tcolorbox}[enhanced,nobeforeafter,width=7.1cm,height=2.1cm,valign=center,colback=cgeo!6,colframe=cgeo,arc=2.5mm,boxrule=0.9pt,left=3mm]{\Large\bfseries\color{cgeo}b)}\quad\begin{minipage}[c]{5.3cm}\raggedright\Large $72\pi~\text{cm}^2$\end{minipage}\end{tcolorbox}\par\vspace{2.5mm}
\begin{tcolorbox}[enhanced,nobeforeafter,width=7.1cm,height=2.1cm,valign=center,colback=cgeo!6,colframe=cgeo,arc=2.5mm,boxrule=0.9pt,left=3mm]{\Large\bfseries\color{cgeo}c)}\quad\begin{minipage}[c]{5.3cm}\raggedright\Large $36\pi~\text{cm}^2$\end{minipage}\end{tcolorbox}\hfill\begin{tcolorbox}[enhanced,nobeforeafter,width=7.1cm,height=2.1cm,valign=center,colback=cgeo!6,colframe=cgeo,arc=2.5mm,boxrule=0.9pt,left=3mm]{\Large\bfseries\color{cgeo}d)}\quad\begin{minipage}[c]{5.3cm}\raggedright\Large $144\pi~\text{cm}^2$\end{minipage}\end{tcolorbox}
\end{tcolorbox}
\newpage
\begin{tcolorbox}[enhanced,width=16cm,height=14cm,colback=white,colframe=cgeo,boxrule=2pt,arc=0mm,
left=6mm,right=6mm,top=0mm,bottom=4mm,sharp corners,
overlay={\fill[cgeo] (frame.north west) rectangle ([yshift=-1.25cm]frame.north east);
\node[anchor=west,text=white,font=\sffamily\bfseries\LARGE] at ([xshift=5mm,yshift=-0.62cm]frame.north west) {Question 62 D};
\node[anchor=east,text=white,font=\sffamily\itshape\large] at ([xshift=-5mm,yshift=-0.62cm]frame.north east) {Géométrie};}]
\vspace*{1.75cm}
\begin{minipage}[t][6.9cm][c]{\linewidth}\LARGE\centering
\begin{minipage}{0.95\linewidth}\raggedright\linespread{1.25}\selectfont Quelle est l'aire exacte d'une sphère de rayon $2{,}5$ cm~?\end{minipage}
\end{minipage}

\vfill
\begin{tcolorbox}[enhanced,nobeforeafter,width=7.1cm,height=2.1cm,valign=center,colback=cgeo!6,colframe=cgeo,arc=2.5mm,boxrule=0.9pt,left=3mm]{\Large\bfseries\color{cgeo}a)}\quad\begin{minipage}[c]{5.3cm}\raggedright\large $\dfrac{25}{4}\pi~\text{cm}^2$\end{minipage}\end{tcolorbox}\hfill\begin{tcolorbox}[enhanced,nobeforeafter,width=7.1cm,height=2.1cm,valign=center,colback=cgeo!6,colframe=cgeo,arc=2.5mm,boxrule=0.9pt,left=3mm]{\Large\bfseries\color{cgeo}b)}\quad\begin{minipage}[c]{5.3cm}\raggedright\Large $25\pi~\text{cm}^2$\end{minipage}\end{tcolorbox}\par\vspace{2.5mm}
\begin{tcolorbox}[enhanced,nobeforeafter,width=7.1cm,height=2.1cm,valign=center,colback=cgeo!6,colframe=cgeo,arc=2.5mm,boxrule=0.9pt,left=3mm]{\Large\bfseries\color{cgeo}c)}\quad\begin{minipage}[c]{5.3cm}\raggedright\large $\dfrac{125}{6}\pi~\text{cm}^2$\end{minipage}\end{tcolorbox}\hfill\begin{tcolorbox}[enhanced,nobeforeafter,width=7.1cm,height=2.1cm,valign=center,colback=cgeo!6,colframe=cgeo,arc=2.5mm,boxrule=0.9pt,left=3mm]{\Large\bfseries\color{cgeo}d)}\quad\begin{minipage}[c]{5.3cm}\raggedright\large $\dfrac{25}{2}\pi~\text{cm}^2$\end{minipage}\end{tcolorbox}
\end{tcolorbox}
\newpage
\begin{tcolorbox}[enhanced,width=16cm,height=14cm,colback=white,colframe=cgeo,boxrule=2pt,arc=0mm,
left=6mm,right=6mm,top=0mm,bottom=4mm,sharp corners,
overlay={\fill[cgeo] (frame.north west) rectangle ([yshift=-1.25cm]frame.north east);
\node[anchor=west,text=white,font=\sffamily\bfseries\LARGE] at ([xshift=5mm,yshift=-0.62cm]frame.north west) {Question 63 D};
\node[anchor=east,text=white,font=\sffamily\itshape\large] at ([xshift=-5mm,yshift=-0.62cm]frame.north east) {Géométrie};}]
\vspace*{1.75cm}
\begin{minipage}[t][6.9cm][c]{\linewidth}\LARGE\centering
\begin{minipage}{0.95\linewidth}\raggedright\linespread{1.25}\selectfont Un solide est réduit à l'échelle $\dfrac{1}{3}$. Par combien son volume est-il multiplié~?\end{minipage}
\end{minipage}

\vfill
\begin{tcolorbox}[enhanced,nobeforeafter,width=7.1cm,height=2.1cm,valign=center,colback=cgeo!6,colframe=cgeo,arc=2.5mm,boxrule=0.9pt,left=3mm]{\Large\bfseries\color{cgeo}a)}\quad\begin{minipage}[c]{5.3cm}\raggedright\Large $1$\end{minipage}\end{tcolorbox}\hfill\begin{tcolorbox}[enhanced,nobeforeafter,width=7.1cm,height=2.1cm,valign=center,colback=cgeo!6,colframe=cgeo,arc=2.5mm,boxrule=0.9pt,left=3mm]{\Large\bfseries\color{cgeo}b)}\quad\begin{minipage}[c]{5.3cm}\raggedright\Large $\dfrac{1}{3}$\end{minipage}\end{tcolorbox}\par\vspace{2.5mm}
\begin{tcolorbox}[enhanced,nobeforeafter,width=7.1cm,height=2.1cm,valign=center,colback=cgeo!6,colframe=cgeo,arc=2.5mm,boxrule=0.9pt,left=3mm]{\Large\bfseries\color{cgeo}c)}\quad\begin{minipage}[c]{5.3cm}\raggedright\Large $\dfrac{1}{27}$\end{minipage}\end{tcolorbox}\hfill\begin{tcolorbox}[enhanced,nobeforeafter,width=7.1cm,height=2.1cm,valign=center,colback=cgeo!6,colframe=cgeo,arc=2.5mm,boxrule=0.9pt,left=3mm]{\Large\bfseries\color{cgeo}d)}\quad\begin{minipage}[c]{5.3cm}\raggedright\Large $\dfrac{1}{9}$\end{minipage}\end{tcolorbox}
\end{tcolorbox}
\newpage
\begin{tcolorbox}[enhanced,width=16cm,height=14cm,colback=white,colframe=cgeo,boxrule=2pt,arc=0mm,
left=6mm,right=6mm,top=0mm,bottom=4mm,sharp corners,
overlay={\fill[cgeo] (frame.north west) rectangle ([yshift=-1.25cm]frame.north east);
\node[anchor=west,text=white,font=\sffamily\bfseries\LARGE] at ([xshift=5mm,yshift=-0.62cm]frame.north west) {Question 64 D};
\node[anchor=east,text=white,font=\sffamily\itshape\large] at ([xshift=-5mm,yshift=-0.62cm]frame.north east) {Géométrie};}]
\vspace*{1.75cm}
\begin{minipage}[t][6.9cm][c]{\linewidth}\LARGE\centering
\begin{minipage}{0.95\linewidth}\raggedright\linespread{1.25}\selectfont Un solide est agrandi à l'échelle $5$. Par combien son volume est-il multiplié~?\end{minipage}
\end{minipage}

\vfill
\begin{tcolorbox}[enhanced,nobeforeafter,width=7.1cm,height=2.1cm,valign=center,colback=cgeo!6,colframe=cgeo,arc=2.5mm,boxrule=0.9pt,left=3mm]{\Large\bfseries\color{cgeo}a)}\quad\begin{minipage}[c]{5.3cm}\raggedright\Large $15$\end{minipage}\end{tcolorbox}\hfill\begin{tcolorbox}[enhanced,nobeforeafter,width=7.1cm,height=2.1cm,valign=center,colback=cgeo!6,colframe=cgeo,arc=2.5mm,boxrule=0.9pt,left=3mm]{\Large\bfseries\color{cgeo}b)}\quad\begin{minipage}[c]{5.3cm}\raggedright\Large $25$\end{minipage}\end{tcolorbox}\par\vspace{2.5mm}
\begin{tcolorbox}[enhanced,nobeforeafter,width=7.1cm,height=2.1cm,valign=center,colback=cgeo!6,colframe=cgeo,arc=2.5mm,boxrule=0.9pt,left=3mm]{\Large\bfseries\color{cgeo}c)}\quad\begin{minipage}[c]{5.3cm}\raggedright\Large $125$\end{minipage}\end{tcolorbox}\hfill\begin{tcolorbox}[enhanced,nobeforeafter,width=7.1cm,height=2.1cm,valign=center,colback=cgeo!6,colframe=cgeo,arc=2.5mm,boxrule=0.9pt,left=3mm]{\Large\bfseries\color{cgeo}d)}\quad\begin{minipage}[c]{5.3cm}\raggedright\Large $5$\end{minipage}\end{tcolorbox}
\end{tcolorbox}
\newpage
\begin{tcolorbox}[enhanced,width=16cm,height=14cm,colback=white,colframe=csta,boxrule=2pt,arc=0mm,
left=6mm,right=6mm,top=0mm,bottom=4mm,sharp corners,
overlay={\fill[csta] (frame.north west) rectangle ([yshift=-1.25cm]frame.north east);
\node[anchor=west,text=white,font=\sffamily\bfseries\LARGE] at ([xshift=5mm,yshift=-0.62cm]frame.north west) {Question 65 D};
\node[anchor=east,text=white,font=\sffamily\itshape\large] at ([xshift=-5mm,yshift=-0.62cm]frame.north east) {Statistiques et probabilités};}]
\vspace*{1.75cm}
\begin{minipage}[t][6.9cm][c]{\linewidth}\LARGE\centering
\begin{minipage}{0.95\linewidth}\raggedright\linespread{1.25}\selectfont Quelle est la médiane de la série $40\ ;\ 35\ ;\ 52\ ;\ 48\ ;\ 39$~?\end{minipage}
\end{minipage}

\vfill
\begin{tcolorbox}[enhanced,nobeforeafter,width=7.1cm,height=2.1cm,valign=center,colback=csta!6,colframe=csta,arc=2.5mm,boxrule=0.9pt,left=3mm]{\Large\bfseries\color{csta}a)}\quad\begin{minipage}[c]{5.3cm}\raggedright\Large $40$\end{minipage}\end{tcolorbox}\hfill\begin{tcolorbox}[enhanced,nobeforeafter,width=7.1cm,height=2.1cm,valign=center,colback=csta!6,colframe=csta,arc=2.5mm,boxrule=0.9pt,left=3mm]{\Large\bfseries\color{csta}b)}\quad\begin{minipage}[c]{5.3cm}\raggedright\Large $42{,}8$\end{minipage}\end{tcolorbox}\par\vspace{2.5mm}
\begin{tcolorbox}[enhanced,nobeforeafter,width=7.1cm,height=2.1cm,valign=center,colback=csta!6,colframe=csta,arc=2.5mm,boxrule=0.9pt,left=3mm]{\Large\bfseries\color{csta}c)}\quad\begin{minipage}[c]{5.3cm}\raggedright\Large $17$\end{minipage}\end{tcolorbox}\hfill\begin{tcolorbox}[enhanced,nobeforeafter,width=7.1cm,height=2.1cm,valign=center,colback=csta!6,colframe=csta,arc=2.5mm,boxrule=0.9pt,left=3mm]{\Large\bfseries\color{csta}d)}\quad\begin{minipage}[c]{5.3cm}\raggedright\Large $52$\end{minipage}\end{tcolorbox}
\end{tcolorbox}
\newpage
\begin{tcolorbox}[enhanced,width=16cm,height=14cm,colback=white,colframe=csta,boxrule=2pt,arc=0mm,
left=6mm,right=6mm,top=0mm,bottom=4mm,sharp corners,
overlay={\fill[csta] (frame.north west) rectangle ([yshift=-1.25cm]frame.north east);
\node[anchor=west,text=white,font=\sffamily\bfseries\LARGE] at ([xshift=5mm,yshift=-0.62cm]frame.north west) {Question 66 D};
\node[anchor=east,text=white,font=\sffamily\itshape\large] at ([xshift=-5mm,yshift=-0.62cm]frame.north east) {Statistiques et probabilités};}]
\vspace*{1.75cm}
\begin{minipage}[t][6.9cm][c]{\linewidth}\LARGE\centering
\begin{minipage}{0.95\linewidth}\raggedright\linespread{1.25}\selectfont Quelle est la médiane de la série $11\ ;\ 15\ ;\ 8\ ;\ 20\ ;\ 13\ ;\ 9$~?\end{minipage}
\end{minipage}

\vfill
\begin{tcolorbox}[enhanced,nobeforeafter,width=7.1cm,height=2.1cm,valign=center,colback=csta!6,colframe=csta,arc=2.5mm,boxrule=0.9pt,left=3mm]{\Large\bfseries\color{csta}a)}\quad\begin{minipage}[c]{5.3cm}\raggedright\Large $\dfrac{38}{3}$\end{minipage}\end{tcolorbox}\hfill\begin{tcolorbox}[enhanced,nobeforeafter,width=7.1cm,height=2.1cm,valign=center,colback=csta!6,colframe=csta,arc=2.5mm,boxrule=0.9pt,left=3mm]{\Large\bfseries\color{csta}b)}\quad\begin{minipage}[c]{5.3cm}\raggedright\Large $12$\end{minipage}\end{tcolorbox}\par\vspace{2.5mm}
\begin{tcolorbox}[enhanced,nobeforeafter,width=7.1cm,height=2.1cm,valign=center,colback=csta!6,colframe=csta,arc=2.5mm,boxrule=0.9pt,left=3mm]{\Large\bfseries\color{csta}c)}\quad\begin{minipage}[c]{5.3cm}\raggedright\Large $14$\end{minipage}\end{tcolorbox}\hfill\begin{tcolorbox}[enhanced,nobeforeafter,width=7.1cm,height=2.1cm,valign=center,colback=csta!6,colframe=csta,arc=2.5mm,boxrule=0.9pt,left=3mm]{\Large\bfseries\color{csta}d)}\quad\begin{minipage}[c]{5.3cm}\raggedright\Large $3{,}5$\end{minipage}\end{tcolorbox}
\end{tcolorbox}
\newpage
\begin{tcolorbox}[enhanced,width=16cm,height=14cm,colback=white,colframe=csta,boxrule=2pt,arc=0mm,
left=6mm,right=6mm,top=0mm,bottom=4mm,sharp corners,
overlay={\fill[csta] (frame.north west) rectangle ([yshift=-1.25cm]frame.north east);
\node[anchor=west,text=white,font=\sffamily\bfseries\LARGE] at ([xshift=5mm,yshift=-0.62cm]frame.north west) {Question 67 D};
\node[anchor=east,text=white,font=\sffamily\itshape\large] at ([xshift=-5mm,yshift=-0.62cm]frame.north east) {Statistiques et probabilités};}]
\vspace*{1.75cm}
\begin{minipage}[t][6.9cm][c]{\linewidth}\LARGE\centering
\begin{minipage}{0.95\linewidth}\raggedright\linespread{1.25}\selectfont Quelle est la moyenne de la série $16\ ;\ 14\ ;\ 19\ ;\ 12$~?\end{minipage}
\end{minipage}

\vfill
\begin{tcolorbox}[enhanced,nobeforeafter,width=7.1cm,height=2.1cm,valign=center,colback=csta!6,colframe=csta,arc=2.5mm,boxrule=0.9pt,left=3mm]{\Large\bfseries\color{csta}a)}\quad\begin{minipage}[c]{5.3cm}\raggedright\Large $16{,}25$\end{minipage}\end{tcolorbox}\hfill\begin{tcolorbox}[enhanced,nobeforeafter,width=7.1cm,height=2.1cm,valign=center,colback=csta!6,colframe=csta,arc=2.5mm,boxrule=0.9pt,left=3mm]{\Large\bfseries\color{csta}b)}\quad\begin{minipage}[c]{5.3cm}\raggedright\Large $61$\end{minipage}\end{tcolorbox}\par\vspace{2.5mm}
\begin{tcolorbox}[enhanced,nobeforeafter,width=7.1cm,height=2.1cm,valign=center,colback=csta!6,colframe=csta,arc=2.5mm,boxrule=0.9pt,left=3mm]{\Large\bfseries\color{csta}c)}\quad\begin{minipage}[c]{5.3cm}\raggedright\Large $15{,}25$\end{minipage}\end{tcolorbox}\hfill\begin{tcolorbox}[enhanced,nobeforeafter,width=7.1cm,height=2.1cm,valign=center,colback=csta!6,colframe=csta,arc=2.5mm,boxrule=0.9pt,left=3mm]{\Large\bfseries\color{csta}d)}\quad\begin{minipage}[c]{5.3cm}\raggedright\Large $15$\end{minipage}\end{tcolorbox}
\end{tcolorbox}
\newpage
\begin{tcolorbox}[enhanced,width=16cm,height=14cm,colback=white,colframe=csta,boxrule=2pt,arc=0mm,
left=6mm,right=6mm,top=0mm,bottom=4mm,sharp corners,
overlay={\fill[csta] (frame.north west) rectangle ([yshift=-1.25cm]frame.north east);
\node[anchor=west,text=white,font=\sffamily\bfseries\LARGE] at ([xshift=5mm,yshift=-0.62cm]frame.north west) {Question 68 D};
\node[anchor=east,text=white,font=\sffamily\itshape\large] at ([xshift=-5mm,yshift=-0.62cm]frame.north east) {Statistiques et probabilités};}]
\vspace*{1.75cm}
\begin{minipage}[t][6.9cm][c]{\linewidth}\LARGE\centering
\begin{minipage}{0.95\linewidth}\raggedright\linespread{1.25}\selectfont Quelle est la moyenne de la série $1\ ;\ 4\ ;\ 2\ ;\ 8\ ;\ 10$~?\end{minipage}
\end{minipage}

\vfill
\begin{tcolorbox}[enhanced,nobeforeafter,width=7.1cm,height=2.1cm,valign=center,colback=csta!6,colframe=csta,arc=2.5mm,boxrule=0.9pt,left=3mm]{\Large\bfseries\color{csta}a)}\quad\begin{minipage}[c]{5.3cm}\raggedright\Large $4$\end{minipage}\end{tcolorbox}\hfill\begin{tcolorbox}[enhanced,nobeforeafter,width=7.1cm,height=2.1cm,valign=center,colback=csta!6,colframe=csta,arc=2.5mm,boxrule=0.9pt,left=3mm]{\Large\bfseries\color{csta}b)}\quad\begin{minipage}[c]{5.3cm}\raggedright\Large $5$\end{minipage}\end{tcolorbox}\par\vspace{2.5mm}
\begin{tcolorbox}[enhanced,nobeforeafter,width=7.1cm,height=2.1cm,valign=center,colback=csta!6,colframe=csta,arc=2.5mm,boxrule=0.9pt,left=3mm]{\Large\bfseries\color{csta}c)}\quad\begin{minipage}[c]{5.3cm}\raggedright\Large $25$\end{minipage}\end{tcolorbox}\hfill\begin{tcolorbox}[enhanced,nobeforeafter,width=7.1cm,height=2.1cm,valign=center,colback=csta!6,colframe=csta,arc=2.5mm,boxrule=0.9pt,left=3mm]{\Large\bfseries\color{csta}d)}\quad\begin{minipage}[c]{5.3cm}\raggedright\Large $6$\end{minipage}\end{tcolorbox}
\end{tcolorbox}
\newpage
\begin{tcolorbox}[enhanced,width=16cm,height=14cm,colback=white,colframe=csta,boxrule=2pt,arc=0mm,
left=6mm,right=6mm,top=0mm,bottom=4mm,sharp corners,
overlay={\fill[csta] (frame.north west) rectangle ([yshift=-1.25cm]frame.north east);
\node[anchor=west,text=white,font=\sffamily\bfseries\LARGE] at ([xshift=5mm,yshift=-0.62cm]frame.north west) {Question 69 D};
\node[anchor=east,text=white,font=\sffamily\itshape\large] at ([xshift=-5mm,yshift=-0.62cm]frame.north east) {Statistiques et probabilités};}]
\vspace*{1.75cm}
\begin{minipage}[t][6.9cm][c]{\linewidth}\LARGE\centering
\begin{minipage}{0.95\linewidth}\raggedright\linespread{1.25}\selectfont Quelle est l'étendue de la série $16\ ;\ 14\ ;\ 19\ ;\ 12$~?\end{minipage}
\end{minipage}

\vfill
\begin{tcolorbox}[enhanced,nobeforeafter,width=7.1cm,height=2.1cm,valign=center,colback=csta!6,colframe=csta,arc=2.5mm,boxrule=0.9pt,left=3mm]{\Large\bfseries\color{csta}a)}\quad\begin{minipage}[c]{5.3cm}\raggedright\Large $8$\end{minipage}\end{tcolorbox}\hfill\begin{tcolorbox}[enhanced,nobeforeafter,width=7.1cm,height=2.1cm,valign=center,colback=csta!6,colframe=csta,arc=2.5mm,boxrule=0.9pt,left=3mm]{\Large\bfseries\color{csta}b)}\quad\begin{minipage}[c]{5.3cm}\raggedright\Large $7$\end{minipage}\end{tcolorbox}\par\vspace{2.5mm}
\begin{tcolorbox}[enhanced,nobeforeafter,width=7.1cm,height=2.1cm,valign=center,colback=csta!6,colframe=csta,arc=2.5mm,boxrule=0.9pt,left=3mm]{\Large\bfseries\color{csta}c)}\quad\begin{minipage}[c]{5.3cm}\raggedright\Large $31$\end{minipage}\end{tcolorbox}\hfill\begin{tcolorbox}[enhanced,nobeforeafter,width=7.1cm,height=2.1cm,valign=center,colback=csta!6,colframe=csta,arc=2.5mm,boxrule=0.9pt,left=3mm]{\Large\bfseries\color{csta}d)}\quad\begin{minipage}[c]{5.3cm}\raggedright\Large $19$\end{minipage}\end{tcolorbox}
\end{tcolorbox}
\newpage
\begin{tcolorbox}[enhanced,width=16cm,height=14cm,colback=white,colframe=csta,boxrule=2pt,arc=0mm,
left=6mm,right=6mm,top=0mm,bottom=4mm,sharp corners,
overlay={\fill[csta] (frame.north west) rectangle ([yshift=-1.25cm]frame.north east);
\node[anchor=west,text=white,font=\sffamily\bfseries\LARGE] at ([xshift=5mm,yshift=-0.62cm]frame.north west) {Question 70 D};
\node[anchor=east,text=white,font=\sffamily\itshape\large] at ([xshift=-5mm,yshift=-0.62cm]frame.north east) {Statistiques et probabilités};}]
\vspace*{1.75cm}
\begin{minipage}[t][6.9cm][c]{\linewidth}\LARGE\centering
\begin{minipage}{0.95\linewidth}\raggedright\linespread{1.25}\selectfont Quelle est l'étendue de la série $1\ ;\ 4\ ;\ 2\ ;\ 8\ ;\ 10$~?\end{minipage}
\end{minipage}

\vfill
\begin{tcolorbox}[enhanced,nobeforeafter,width=7.1cm,height=2.1cm,valign=center,colback=csta!6,colframe=csta,arc=2.5mm,boxrule=0.9pt,left=3mm]{\Large\bfseries\color{csta}a)}\quad\begin{minipage}[c]{5.3cm}\raggedright\Large $5$\end{minipage}\end{tcolorbox}\hfill\begin{tcolorbox}[enhanced,nobeforeafter,width=7.1cm,height=2.1cm,valign=center,colback=csta!6,colframe=csta,arc=2.5mm,boxrule=0.9pt,left=3mm]{\Large\bfseries\color{csta}b)}\quad\begin{minipage}[c]{5.3cm}\raggedright\Large $9$\end{minipage}\end{tcolorbox}\par\vspace{2.5mm}
\begin{tcolorbox}[enhanced,nobeforeafter,width=7.1cm,height=2.1cm,valign=center,colback=csta!6,colframe=csta,arc=2.5mm,boxrule=0.9pt,left=3mm]{\Large\bfseries\color{csta}c)}\quad\begin{minipage}[c]{5.3cm}\raggedright\Large $10$\end{minipage}\end{tcolorbox}\hfill\begin{tcolorbox}[enhanced,nobeforeafter,width=7.1cm,height=2.1cm,valign=center,colback=csta!6,colframe=csta,arc=2.5mm,boxrule=0.9pt,left=3mm]{\Large\bfseries\color{csta}d)}\quad\begin{minipage}[c]{5.3cm}\raggedright\Large $11$\end{minipage}\end{tcolorbox}
\end{tcolorbox}
\newpage
\begin{tcolorbox}[enhanced,width=16cm,height=14cm,colback=white,colframe=csta,boxrule=2pt,arc=0mm,
left=6mm,right=6mm,top=0mm,bottom=4mm,sharp corners,
overlay={\fill[csta] (frame.north west) rectangle ([yshift=-1.25cm]frame.north east);
\node[anchor=west,text=white,font=\sffamily\bfseries\LARGE] at ([xshift=5mm,yshift=-0.62cm]frame.north west) {Question 71 D};
\node[anchor=east,text=white,font=\sffamily\itshape\large] at ([xshift=-5mm,yshift=-0.62cm]frame.north east) {Statistiques et probabilités};}]
\vspace*{1.75cm}
\begin{minipage}[t][6.9cm][c]{\linewidth}\Large\centering
\begin{minipage}{0.95\linewidth}\raggedright\linespread{1.25}\selectfont Une urne contient $6$ boules dont $2$ rouges. On tire successivement deux boules, sans remise. Quelle est la probabilité de tirer deux boules rouges~?\end{minipage}
\end{minipage}

\vfill
\begin{tcolorbox}[enhanced,nobeforeafter,width=7.1cm,height=2.1cm,valign=center,colback=csta!6,colframe=csta,arc=2.5mm,boxrule=0.9pt,left=3mm]{\Large\bfseries\color{csta}a)}\quad\begin{minipage}[c]{5.3cm}\raggedright\Large $\dfrac{1}{9}$\end{minipage}\end{tcolorbox}\hfill\begin{tcolorbox}[enhanced,nobeforeafter,width=7.1cm,height=2.1cm,valign=center,colback=csta!6,colframe=csta,arc=2.5mm,boxrule=0.9pt,left=3mm]{\Large\bfseries\color{csta}b)}\quad\begin{minipage}[c]{5.3cm}\raggedright\Large $\dfrac{1}{18}$\end{minipage}\end{tcolorbox}\par\vspace{2.5mm}
\begin{tcolorbox}[enhanced,nobeforeafter,width=7.1cm,height=2.1cm,valign=center,colback=csta!6,colframe=csta,arc=2.5mm,boxrule=0.9pt,left=3mm]{\Large\bfseries\color{csta}c)}\quad\begin{minipage}[c]{5.3cm}\raggedright\Large $\dfrac{1}{15}$\end{minipage}\end{tcolorbox}\hfill\begin{tcolorbox}[enhanced,nobeforeafter,width=7.1cm,height=2.1cm,valign=center,colback=csta!6,colframe=csta,arc=2.5mm,boxrule=0.9pt,left=3mm]{\Large\bfseries\color{csta}d)}\quad\begin{minipage}[c]{5.3cm}\raggedright\Large $\dfrac{2}{3}$\end{minipage}\end{tcolorbox}
\end{tcolorbox}
\newpage
\begin{tcolorbox}[enhanced,width=16cm,height=14cm,colback=white,colframe=csta,boxrule=2pt,arc=0mm,
left=6mm,right=6mm,top=0mm,bottom=4mm,sharp corners,
overlay={\fill[csta] (frame.north west) rectangle ([yshift=-1.25cm]frame.north east);
\node[anchor=west,text=white,font=\sffamily\bfseries\LARGE] at ([xshift=5mm,yshift=-0.62cm]frame.north west) {Question 72 D};
\node[anchor=east,text=white,font=\sffamily\itshape\large] at ([xshift=-5mm,yshift=-0.62cm]frame.north east) {Statistiques et probabilités};}]
\vspace*{1.75cm}
\begin{minipage}[t][6.9cm][c]{\linewidth}\Large\centering
\begin{minipage}{0.95\linewidth}\raggedright\linespread{1.25}\selectfont Une urne contient $10$ boules dont $4$ rouges. On tire successivement deux boules, sans remise. Quelle est la probabilité de tirer deux boules rouges~?\end{minipage}
\end{minipage}

\vfill
\begin{tcolorbox}[enhanced,nobeforeafter,width=7.1cm,height=2.1cm,valign=center,colback=csta!6,colframe=csta,arc=2.5mm,boxrule=0.9pt,left=3mm]{\Large\bfseries\color{csta}a)}\quad\begin{minipage}[c]{5.3cm}\raggedright\Large $\dfrac{4}{25}$\end{minipage}\end{tcolorbox}\hfill\begin{tcolorbox}[enhanced,nobeforeafter,width=7.1cm,height=2.1cm,valign=center,colback=csta!6,colframe=csta,arc=2.5mm,boxrule=0.9pt,left=3mm]{\Large\bfseries\color{csta}b)}\quad\begin{minipage}[c]{5.3cm}\raggedright\Large $\dfrac{2}{15}$\end{minipage}\end{tcolorbox}\par\vspace{2.5mm}
\begin{tcolorbox}[enhanced,nobeforeafter,width=7.1cm,height=2.1cm,valign=center,colback=csta!6,colframe=csta,arc=2.5mm,boxrule=0.9pt,left=3mm]{\Large\bfseries\color{csta}c)}\quad\begin{minipage}[c]{5.3cm}\raggedright\Large $\dfrac{4}{5}$\end{minipage}\end{tcolorbox}\hfill\begin{tcolorbox}[enhanced,nobeforeafter,width=7.1cm,height=2.1cm,valign=center,colback=csta!6,colframe=csta,arc=2.5mm,boxrule=0.9pt,left=3mm]{\Large\bfseries\color{csta}d)}\quad\begin{minipage}[c]{5.3cm}\raggedright\Large $\dfrac{3}{25}$\end{minipage}\end{tcolorbox}
\end{tcolorbox}
\newpage
\begin{tcolorbox}[enhanced,width=16cm,height=14cm,colback=white,colframe=csta,boxrule=2pt,arc=0mm,
left=6mm,right=6mm,top=0mm,bottom=4mm,sharp corners,
overlay={\fill[csta] (frame.north west) rectangle ([yshift=-1.25cm]frame.north east);
\node[anchor=west,text=white,font=\sffamily\bfseries\LARGE] at ([xshift=5mm,yshift=-0.62cm]frame.north west) {Question 73 D};
\node[anchor=east,text=white,font=\sffamily\itshape\large] at ([xshift=-5mm,yshift=-0.62cm]frame.north east) {Statistiques et probabilités};}]
\vspace*{1.75cm}
\begin{minipage}[t][6.9cm][c]{\linewidth}\Large\centering
\begin{minipage}{0.95\linewidth}\raggedright\linespread{1.25}\selectfont On lance un dé équilibré à six faces et une pièce équilibrée (ou un deuxième dé). Quelle est la probabilité d'obtenir un nombre supérieur à 4 avec le dé et un nombre pair avec un deuxième dé~?\end{minipage}
\end{minipage}

\vfill
\begin{tcolorbox}[enhanced,nobeforeafter,width=7.1cm,height=2.1cm,valign=center,colback=csta!6,colframe=csta,arc=2.5mm,boxrule=0.9pt,left=3mm]{\Large\bfseries\color{csta}a)}\quad\begin{minipage}[c]{5.3cm}\raggedright\Large $\dfrac{1}{2}$\end{minipage}\end{tcolorbox}\hfill\begin{tcolorbox}[enhanced,nobeforeafter,width=7.1cm,height=2.1cm,valign=center,colback=csta!6,colframe=csta,arc=2.5mm,boxrule=0.9pt,left=3mm]{\Large\bfseries\color{csta}b)}\quad\begin{minipage}[c]{5.3cm}\raggedright\Large $\dfrac{1}{3}$\end{minipage}\end{tcolorbox}\par\vspace{2.5mm}
\begin{tcolorbox}[enhanced,nobeforeafter,width=7.1cm,height=2.1cm,valign=center,colback=csta!6,colframe=csta,arc=2.5mm,boxrule=0.9pt,left=3mm]{\Large\bfseries\color{csta}c)}\quad\begin{minipage}[c]{5.3cm}\raggedright\Large $\dfrac{1}{6}$\end{minipage}\end{tcolorbox}\hfill\begin{tcolorbox}[enhanced,nobeforeafter,width=7.1cm,height=2.1cm,valign=center,colback=csta!6,colframe=csta,arc=2.5mm,boxrule=0.9pt,left=3mm]{\Large\bfseries\color{csta}d)}\quad\begin{minipage}[c]{5.3cm}\raggedright\Large $\dfrac{5}{6}$\end{minipage}\end{tcolorbox}
\end{tcolorbox}
\newpage
\begin{tcolorbox}[enhanced,width=16cm,height=14cm,colback=white,colframe=csta,boxrule=2pt,arc=0mm,
left=6mm,right=6mm,top=0mm,bottom=4mm,sharp corners,
overlay={\fill[csta] (frame.north west) rectangle ([yshift=-1.25cm]frame.north east);
\node[anchor=west,text=white,font=\sffamily\bfseries\LARGE] at ([xshift=5mm,yshift=-0.62cm]frame.north west) {Question 74 D};
\node[anchor=east,text=white,font=\sffamily\itshape\large] at ([xshift=-5mm,yshift=-0.62cm]frame.north east) {Statistiques et probabilités};}]
\vspace*{1.75cm}
\begin{minipage}[t][6.9cm][c]{\linewidth}\Large\centering
\begin{minipage}{0.95\linewidth}\raggedright\linespread{1.25}\selectfont On lance un dé équilibré à six faces et une pièce équilibrée (ou un deuxième dé). Quelle est la probabilité d'obtenir un 3 avec le dé et FACE avec la pièce~?\end{minipage}
\end{minipage}

\vfill
\begin{tcolorbox}[enhanced,nobeforeafter,width=7.1cm,height=2.1cm,valign=center,colback=csta!6,colframe=csta,arc=2.5mm,boxrule=0.9pt,left=3mm]{\Large\bfseries\color{csta}a)}\quad\begin{minipage}[c]{5.3cm}\raggedright\Large $\dfrac{1}{12}$\end{minipage}\end{tcolorbox}\hfill\begin{tcolorbox}[enhanced,nobeforeafter,width=7.1cm,height=2.1cm,valign=center,colback=csta!6,colframe=csta,arc=2.5mm,boxrule=0.9pt,left=3mm]{\Large\bfseries\color{csta}b)}\quad\begin{minipage}[c]{5.3cm}\raggedright\Large $\dfrac{1}{6}$\end{minipage}\end{tcolorbox}\par\vspace{2.5mm}
\begin{tcolorbox}[enhanced,nobeforeafter,width=7.1cm,height=2.1cm,valign=center,colback=csta!6,colframe=csta,arc=2.5mm,boxrule=0.9pt,left=3mm]{\Large\bfseries\color{csta}c)}\quad\begin{minipage}[c]{5.3cm}\raggedright\Large $\dfrac{1}{2}$\end{minipage}\end{tcolorbox}\hfill\begin{tcolorbox}[enhanced,nobeforeafter,width=7.1cm,height=2.1cm,valign=center,colback=csta!6,colframe=csta,arc=2.5mm,boxrule=0.9pt,left=3mm]{\Large\bfseries\color{csta}d)}\quad\begin{minipage}[c]{5.3cm}\raggedright\Large $\dfrac{2}{3}$\end{minipage}\end{tcolorbox}
\end{tcolorbox}
\newpage
\begin{tcolorbox}[enhanced,width=16cm,height=14cm,colback=white,colframe=csta,boxrule=2pt,arc=0mm,
left=6mm,right=6mm,top=0mm,bottom=4mm,sharp corners,
overlay={\fill[csta] (frame.north west) rectangle ([yshift=-1.25cm]frame.north east);
\node[anchor=west,text=white,font=\sffamily\bfseries\LARGE] at ([xshift=5mm,yshift=-0.62cm]frame.north west) {Question 75 D};
\node[anchor=east,text=white,font=\sffamily\itshape\large] at ([xshift=-5mm,yshift=-0.62cm]frame.north east) {Statistiques et probabilités};}]
\vspace*{1.75cm}
\begin{minipage}[t][6.9cm][c]{\linewidth}\LARGE\centering
\begin{minipage}{0.95\linewidth}\raggedright\linespread{1.25}\selectfont Dans une classe de $35$ élèves, $7$ pratiquent un sport en club. Quelle est la fréquence, en pourcentage, de ces élèves~?\end{minipage}
\end{minipage}

\vfill
\begin{tcolorbox}[enhanced,nobeforeafter,width=7.1cm,height=2.1cm,valign=center,colback=csta!6,colframe=csta,arc=2.5mm,boxrule=0.9pt,left=3mm]{\Large\bfseries\color{csta}a)}\quad\begin{minipage}[c]{5.3cm}\raggedright\Large $20\,\%$\end{minipage}\end{tcolorbox}\hfill\begin{tcolorbox}[enhanced,nobeforeafter,width=7.1cm,height=2.1cm,valign=center,colback=csta!6,colframe=csta,arc=2.5mm,boxrule=0.9pt,left=3mm]{\Large\bfseries\color{csta}b)}\quad\begin{minipage}[c]{5.3cm}\raggedright\Large $80\,\%$\end{minipage}\end{tcolorbox}\par\vspace{2.5mm}
\begin{tcolorbox}[enhanced,nobeforeafter,width=7.1cm,height=2.1cm,valign=center,colback=csta!6,colframe=csta,arc=2.5mm,boxrule=0.9pt,left=3mm]{\Large\bfseries\color{csta}c)}\quad\begin{minipage}[c]{5.3cm}\raggedright\Large $2\,\%$\end{minipage}\end{tcolorbox}\hfill\begin{tcolorbox}[enhanced,nobeforeafter,width=7.1cm,height=2.1cm,valign=center,colback=csta!6,colframe=csta,arc=2.5mm,boxrule=0.9pt,left=3mm]{\Large\bfseries\color{csta}d)}\quad\begin{minipage}[c]{5.3cm}\raggedright\Large $7\,\%$\end{minipage}\end{tcolorbox}
\end{tcolorbox}
\newpage
\begin{tcolorbox}[enhanced,width=16cm,height=14cm,colback=white,colframe=csta,boxrule=2pt,arc=0mm,
left=6mm,right=6mm,top=0mm,bottom=4mm,sharp corners,
overlay={\fill[csta] (frame.north west) rectangle ([yshift=-1.25cm]frame.north east);
\node[anchor=west,text=white,font=\sffamily\bfseries\LARGE] at ([xshift=5mm,yshift=-0.62cm]frame.north west) {Question 76 D};
\node[anchor=east,text=white,font=\sffamily\itshape\large] at ([xshift=-5mm,yshift=-0.62cm]frame.north east) {Statistiques et probabilités};}]
\vspace*{1.75cm}
\begin{minipage}[t][6.9cm][c]{\linewidth}\LARGE\centering
\begin{minipage}{0.95\linewidth}\raggedright\linespread{1.25}\selectfont Sur $60$ personnes interrogées, $33$ ont répondu «~oui~». Quelle est la fréquence, en pourcentage, des «~oui~»~?\end{minipage}
\end{minipage}

\vfill
\begin{tcolorbox}[enhanced,nobeforeafter,width=7.1cm,height=2.1cm,valign=center,colback=csta!6,colframe=csta,arc=2.5mm,boxrule=0.9pt,left=3mm]{\Large\bfseries\color{csta}a)}\quad\begin{minipage}[c]{5.3cm}\raggedright\Large $45\,\%$\end{minipage}\end{tcolorbox}\hfill\begin{tcolorbox}[enhanced,nobeforeafter,width=7.1cm,height=2.1cm,valign=center,colback=csta!6,colframe=csta,arc=2.5mm,boxrule=0.9pt,left=3mm]{\Large\bfseries\color{csta}b)}\quad\begin{minipage}[c]{5.3cm}\raggedright\Large $5{,}5\,\%$\end{minipage}\end{tcolorbox}\par\vspace{2.5mm}
\begin{tcolorbox}[enhanced,nobeforeafter,width=7.1cm,height=2.1cm,valign=center,colback=csta!6,colframe=csta,arc=2.5mm,boxrule=0.9pt,left=3mm]{\Large\bfseries\color{csta}c)}\quad\begin{minipage}[c]{5.3cm}\raggedright\Large $33\,\%$\end{minipage}\end{tcolorbox}\hfill\begin{tcolorbox}[enhanced,nobeforeafter,width=7.1cm,height=2.1cm,valign=center,colback=csta!6,colframe=csta,arc=2.5mm,boxrule=0.9pt,left=3mm]{\Large\bfseries\color{csta}d)}\quad\begin{minipage}[c]{5.3cm}\raggedright\Large $55\,\%$\end{minipage}\end{tcolorbox}
\end{tcolorbox}
\newpage
\begin{tcolorbox}[enhanced,width=16cm,height=14cm,colback=white,colframe=csta,boxrule=2pt,arc=0mm,
left=6mm,right=6mm,top=0mm,bottom=4mm,sharp corners,
overlay={\fill[csta] (frame.north west) rectangle ([yshift=-1.25cm]frame.north east);
\node[anchor=west,text=white,font=\sffamily\bfseries\LARGE] at ([xshift=5mm,yshift=-0.62cm]frame.north west) {Question 77 D};
\node[anchor=east,text=white,font=\sffamily\itshape\large] at ([xshift=-5mm,yshift=-0.62cm]frame.north east) {Statistiques et probabilités};}]
\vspace*{1.75cm}
\begin{minipage}[t][6.9cm][c]{\linewidth}\Large\centering
\begin{minipage}{0.95\linewidth}\raggedright\linespread{1.25}\selectfont Un sac contient $10$ boules rouges, $2$ boules bleues et $4$ boules vertes. On tire une boule au hasard. Quelle est la probabilité de ne pas tirer une boule rouge~?\end{minipage}
\end{minipage}

\vfill
\begin{tcolorbox}[enhanced,nobeforeafter,width=7.1cm,height=2.1cm,valign=center,colback=csta!6,colframe=csta,arc=2.5mm,boxrule=0.9pt,left=3mm]{\Large\bfseries\color{csta}a)}\quad\begin{minipage}[c]{5.3cm}\raggedright\Large $\dfrac{5}{8}$\end{minipage}\end{tcolorbox}\hfill\begin{tcolorbox}[enhanced,nobeforeafter,width=7.1cm,height=2.1cm,valign=center,colback=csta!6,colframe=csta,arc=2.5mm,boxrule=0.9pt,left=3mm]{\Large\bfseries\color{csta}b)}\quad\begin{minipage}[c]{5.3cm}\raggedright\Large $\dfrac{3}{5}$\end{minipage}\end{tcolorbox}\par\vspace{2.5mm}
\begin{tcolorbox}[enhanced,nobeforeafter,width=7.1cm,height=2.1cm,valign=center,colback=csta!6,colframe=csta,arc=2.5mm,boxrule=0.9pt,left=3mm]{\Large\bfseries\color{csta}c)}\quad\begin{minipage}[c]{5.3cm}\raggedright\Large $\dfrac{1}{8}$\end{minipage}\end{tcolorbox}\hfill\begin{tcolorbox}[enhanced,nobeforeafter,width=7.1cm,height=2.1cm,valign=center,colback=csta!6,colframe=csta,arc=2.5mm,boxrule=0.9pt,left=3mm]{\Large\bfseries\color{csta}d)}\quad\begin{minipage}[c]{5.3cm}\raggedright\Large $\dfrac{3}{8}$\end{minipage}\end{tcolorbox}
\end{tcolorbox}
\newpage
\begin{tcolorbox}[enhanced,width=16cm,height=14cm,colback=white,colframe=csta,boxrule=2pt,arc=0mm,
left=6mm,right=6mm,top=0mm,bottom=4mm,sharp corners,
overlay={\fill[csta] (frame.north west) rectangle ([yshift=-1.25cm]frame.north east);
\node[anchor=west,text=white,font=\sffamily\bfseries\LARGE] at ([xshift=5mm,yshift=-0.62cm]frame.north west) {Question 78 D};
\node[anchor=east,text=white,font=\sffamily\itshape\large] at ([xshift=-5mm,yshift=-0.62cm]frame.north east) {Statistiques et probabilités};}]
\vspace*{1.75cm}
\begin{minipage}[t][6.9cm][c]{\linewidth}\Large\centering
\begin{minipage}{0.95\linewidth}\raggedright\linespread{1.25}\selectfont Un sac contient $1$ boule rouge, $7$ boules bleues et $4$ boules vertes. On tire une boule au hasard. Quelle est la probabilité de ne pas tirer une boule rouge~?\end{minipage}
\end{minipage}

\vfill
\begin{tcolorbox}[enhanced,nobeforeafter,width=7.1cm,height=2.1cm,valign=center,colback=csta!6,colframe=csta,arc=2.5mm,boxrule=0.9pt,left=3mm]{\Large\bfseries\color{csta}a)}\quad\begin{minipage}[c]{5.3cm}\raggedright\Large $11$\end{minipage}\end{tcolorbox}\hfill\begin{tcolorbox}[enhanced,nobeforeafter,width=7.1cm,height=2.1cm,valign=center,colback=csta!6,colframe=csta,arc=2.5mm,boxrule=0.9pt,left=3mm]{\Large\bfseries\color{csta}b)}\quad\begin{minipage}[c]{5.3cm}\raggedright\Large $\dfrac{1}{12}$\end{minipage}\end{tcolorbox}\par\vspace{2.5mm}
\begin{tcolorbox}[enhanced,nobeforeafter,width=7.1cm,height=2.1cm,valign=center,colback=csta!6,colframe=csta,arc=2.5mm,boxrule=0.9pt,left=3mm]{\Large\bfseries\color{csta}c)}\quad\begin{minipage}[c]{5.3cm}\raggedright\Large $\dfrac{11}{12}$\end{minipage}\end{tcolorbox}\hfill\begin{tcolorbox}[enhanced,nobeforeafter,width=7.1cm,height=2.1cm,valign=center,colback=csta!6,colframe=csta,arc=2.5mm,boxrule=0.9pt,left=3mm]{\Large\bfseries\color{csta}d)}\quad\begin{minipage}[c]{5.3cm}\raggedright\Large $\dfrac{7}{12}$\end{minipage}\end{tcolorbox}
\end{tcolorbox}
\newpage
\begin{tcolorbox}[enhanced,width=16cm,height=14cm,colback=white,colframe=csta,boxrule=2pt,arc=0mm,
left=6mm,right=6mm,top=0mm,bottom=4mm,sharp corners,
overlay={\fill[csta] (frame.north west) rectangle ([yshift=-1.25cm]frame.north east);
\node[anchor=west,text=white,font=\sffamily\bfseries\LARGE] at ([xshift=5mm,yshift=-0.62cm]frame.north west) {Question 79 D};
\node[anchor=east,text=white,font=\sffamily\itshape\large] at ([xshift=-5mm,yshift=-0.62cm]frame.north east) {Statistiques et probabilités};}]
\vspace*{1.75cm}
\begin{minipage}[t][6.9cm][c]{\linewidth}\LARGE\centering
\begin{minipage}{0.95\linewidth}\raggedright\linespread{1.25}\selectfont La probabilité d'un événement $A$ est $\dfrac{11}{12}$. Quelle est la probabilité de l'événement contraire de $A$~?\end{minipage}
\end{minipage}

\vfill
\begin{tcolorbox}[enhanced,nobeforeafter,width=7.1cm,height=2.1cm,valign=center,colback=csta!6,colframe=csta,arc=2.5mm,boxrule=0.9pt,left=3mm]{\Large\bfseries\color{csta}a)}\quad\begin{minipage}[c]{5.3cm}\raggedright\Large $\dfrac{5}{12}$\end{minipage}\end{tcolorbox}\hfill\begin{tcolorbox}[enhanced,nobeforeafter,width=7.1cm,height=2.1cm,valign=center,colback=csta!6,colframe=csta,arc=2.5mm,boxrule=0.9pt,left=3mm]{\Large\bfseries\color{csta}b)}\quad\begin{minipage}[c]{5.3cm}\raggedright\Large $\dfrac{11}{12}$\end{minipage}\end{tcolorbox}\par\vspace{2.5mm}
\begin{tcolorbox}[enhanced,nobeforeafter,width=7.1cm,height=2.1cm,valign=center,colback=csta!6,colframe=csta,arc=2.5mm,boxrule=0.9pt,left=3mm]{\Large\bfseries\color{csta}c)}\quad\begin{minipage}[c]{5.3cm}\raggedright\Large $\dfrac{1}{24}$\end{minipage}\end{tcolorbox}\hfill\begin{tcolorbox}[enhanced,nobeforeafter,width=7.1cm,height=2.1cm,valign=center,colback=csta!6,colframe=csta,arc=2.5mm,boxrule=0.9pt,left=3mm]{\Large\bfseries\color{csta}d)}\quad\begin{minipage}[c]{5.3cm}\raggedright\Large $\dfrac{1}{12}$\end{minipage}\end{tcolorbox}
\end{tcolorbox}
\newpage
\begin{tcolorbox}[enhanced,width=16cm,height=14cm,colback=white,colframe=csta,boxrule=2pt,arc=0mm,
left=6mm,right=6mm,top=0mm,bottom=4mm,sharp corners,
overlay={\fill[csta] (frame.north west) rectangle ([yshift=-1.25cm]frame.north east);
\node[anchor=west,text=white,font=\sffamily\bfseries\LARGE] at ([xshift=5mm,yshift=-0.62cm]frame.north west) {Question 80 D};
\node[anchor=east,text=white,font=\sffamily\itshape\large] at ([xshift=-5mm,yshift=-0.62cm]frame.north east) {Statistiques et probabilités};}]
\vspace*{1.75cm}
\begin{minipage}[t][6.9cm][c]{\linewidth}\LARGE\centering
\begin{minipage}{0.95\linewidth}\raggedright\linespread{1.25}\selectfont La probabilité d'un événement $A$ est $0{,}06$. Quelle est la probabilité de l'événement contraire de $A$~?\end{minipage}
\end{minipage}

\vfill
\begin{tcolorbox}[enhanced,nobeforeafter,width=7.1cm,height=2.1cm,valign=center,colback=csta!6,colframe=csta,arc=2.5mm,boxrule=0.9pt,left=3mm]{\Large\bfseries\color{csta}a)}\quad\begin{minipage}[c]{5.3cm}\raggedright\Large $0{,}06$\end{minipage}\end{tcolorbox}\hfill\begin{tcolorbox}[enhanced,nobeforeafter,width=7.1cm,height=2.1cm,valign=center,colback=csta!6,colframe=csta,arc=2.5mm,boxrule=0.9pt,left=3mm]{\Large\bfseries\color{csta}b)}\quad\begin{minipage}[c]{5.3cm}\raggedright\Large $0{,}47$\end{minipage}\end{tcolorbox}\par\vspace{2.5mm}
\begin{tcolorbox}[enhanced,nobeforeafter,width=7.1cm,height=2.1cm,valign=center,colback=csta!6,colframe=csta,arc=2.5mm,boxrule=0.9pt,left=3mm]{\Large\bfseries\color{csta}c)}\quad\begin{minipage}[c]{5.3cm}\raggedright\Large $0{,}94$\end{minipage}\end{tcolorbox}\hfill\begin{tcolorbox}[enhanced,nobeforeafter,width=7.1cm,height=2.1cm,valign=center,colback=csta!6,colframe=csta,arc=2.5mm,boxrule=0.9pt,left=3mm]{\Large\bfseries\color{csta}d)}\quad\begin{minipage}[c]{5.3cm}\raggedright\Large $0{,}44$\end{minipage}\end{tcolorbox}
\end{tcolorbox}
\newpage
\begin{tcolorbox}[enhanced,width=16cm,height=14cm,colback=white,colframe=cgra,boxrule=2pt,arc=0mm,
left=6mm,right=6mm,top=0mm,bottom=4mm,sharp corners,
overlay={\fill[cgra] (frame.north west) rectangle ([yshift=-1.25cm]frame.north east);
\node[anchor=west,text=white,font=\sffamily\bfseries\LARGE] at ([xshift=5mm,yshift=-0.62cm]frame.north west) {Question 81 D};
\node[anchor=east,text=white,font=\sffamily\itshape\large] at ([xshift=-5mm,yshift=-0.62cm]frame.north east) {Grandeurs et proportionnalité};}]
\vspace*{1.75cm}
\begin{minipage}[t][6.9cm][c]{\linewidth}\LARGE\centering
\begin{minipage}{0.95\linewidth}\raggedright\linespread{1.25}\selectfont Une vitesse de $2{,}5$ m/s correspond à~:\end{minipage}
\end{minipage}

\vfill
\begin{tcolorbox}[enhanced,nobeforeafter,width=7.1cm,height=2.1cm,valign=center,colback=cgra!6,colframe=cgra,arc=2.5mm,boxrule=0.9pt,left=3mm]{\Large\bfseries\color{cgra}a)}\quad\begin{minipage}[c]{5.3cm}\raggedright\Large $2\,500~\text{km/h}$\end{minipage}\end{tcolorbox}\hfill\begin{tcolorbox}[enhanced,nobeforeafter,width=7.1cm,height=2.1cm,valign=center,colback=cgra!6,colframe=cgra,arc=2.5mm,boxrule=0.9pt,left=3mm]{\Large\bfseries\color{cgra}b)}\quad\begin{minipage}[c]{5.3cm}\raggedright\Large $150~\text{km/h}$\end{minipage}\end{tcolorbox}\par\vspace{2.5mm}
\begin{tcolorbox}[enhanced,nobeforeafter,width=7.1cm,height=2.1cm,valign=center,colback=cgra!6,colframe=cgra,arc=2.5mm,boxrule=0.9pt,left=3mm]{\Large\bfseries\color{cgra}c)}\quad\begin{minipage}[c]{5.3cm}\raggedright\Large $25~\text{km/h}$\end{minipage}\end{tcolorbox}\hfill\begin{tcolorbox}[enhanced,nobeforeafter,width=7.1cm,height=2.1cm,valign=center,colback=cgra!6,colframe=cgra,arc=2.5mm,boxrule=0.9pt,left=3mm]{\Large\bfseries\color{cgra}d)}\quad\begin{minipage}[c]{5.3cm}\raggedright\Large $9~\text{km/h}$\end{minipage}\end{tcolorbox}
\end{tcolorbox}
\newpage
\begin{tcolorbox}[enhanced,width=16cm,height=14cm,colback=white,colframe=cgra,boxrule=2pt,arc=0mm,
left=6mm,right=6mm,top=0mm,bottom=4mm,sharp corners,
overlay={\fill[cgra] (frame.north west) rectangle ([yshift=-1.25cm]frame.north east);
\node[anchor=west,text=white,font=\sffamily\bfseries\LARGE] at ([xshift=5mm,yshift=-0.62cm]frame.north west) {Question 82 D};
\node[anchor=east,text=white,font=\sffamily\itshape\large] at ([xshift=-5mm,yshift=-0.62cm]frame.north east) {Grandeurs et proportionnalité};}]
\vspace*{1.75cm}
\begin{minipage}[t][6.9cm][c]{\linewidth}\LARGE\centering
\begin{minipage}{0.95\linewidth}\raggedright\linespread{1.25}\selectfont Une vitesse de $40$ m/s correspond à~:\end{minipage}
\end{minipage}

\vfill
\begin{tcolorbox}[enhanced,nobeforeafter,width=7.1cm,height=2.1cm,valign=center,colback=cgra!6,colframe=cgra,arc=2.5mm,boxrule=0.9pt,left=3mm]{\Large\bfseries\color{cgra}a)}\quad\begin{minipage}[c]{5.3cm}\raggedright\Large $40\,000~\text{km/h}$\end{minipage}\end{tcolorbox}\hfill\begin{tcolorbox}[enhanced,nobeforeafter,width=7.1cm,height=2.1cm,valign=center,colback=cgra!6,colframe=cgra,arc=2.5mm,boxrule=0.9pt,left=3mm]{\Large\bfseries\color{cgra}b)}\quad\begin{minipage}[c]{5.3cm}\raggedright\Large $400~\text{km/h}$\end{minipage}\end{tcolorbox}\par\vspace{2.5mm}
\begin{tcolorbox}[enhanced,nobeforeafter,width=7.1cm,height=2.1cm,valign=center,colback=cgra!6,colframe=cgra,arc=2.5mm,boxrule=0.9pt,left=3mm]{\Large\bfseries\color{cgra}c)}\quad\begin{minipage}[c]{5.3cm}\raggedright\Large $2\,400~\text{km/h}$\end{minipage}\end{tcolorbox}\hfill\begin{tcolorbox}[enhanced,nobeforeafter,width=7.1cm,height=2.1cm,valign=center,colback=cgra!6,colframe=cgra,arc=2.5mm,boxrule=0.9pt,left=3mm]{\Large\bfseries\color{cgra}d)}\quad\begin{minipage}[c]{5.3cm}\raggedright\Large $144~\text{km/h}$\end{minipage}\end{tcolorbox}
\end{tcolorbox}
\newpage
\begin{tcolorbox}[enhanced,width=16cm,height=14cm,colback=white,colframe=cgra,boxrule=2pt,arc=0mm,
left=6mm,right=6mm,top=0mm,bottom=4mm,sharp corners,
overlay={\fill[cgra] (frame.north west) rectangle ([yshift=-1.25cm]frame.north east);
\node[anchor=west,text=white,font=\sffamily\bfseries\LARGE] at ([xshift=5mm,yshift=-0.62cm]frame.north west) {Question 83 D};
\node[anchor=east,text=white,font=\sffamily\itshape\large] at ([xshift=-5mm,yshift=-0.62cm]frame.north east) {Grandeurs et proportionnalité};}]
\vspace*{1.75cm}
\begin{minipage}[t][6.9cm][c]{\linewidth}\LARGE\centering
\begin{minipage}{0.95\linewidth}\raggedright\linespread{1.25}\selectfont Un véhicule parcourt $99$ km en 54 min. Quelle est sa vitesse moyenne~?\end{minipage}
\end{minipage}

\vfill
\begin{tcolorbox}[enhanced,nobeforeafter,width=7.1cm,height=2.1cm,valign=center,colback=cgra!6,colframe=cgra,arc=2.5mm,boxrule=0.9pt,left=3mm]{\Large\bfseries\color{cgra}a)}\quad\begin{minipage}[c]{5.3cm}\raggedright\Large $89{,}1~\text{km/h}$\end{minipage}\end{tcolorbox}\hfill\begin{tcolorbox}[enhanced,nobeforeafter,width=7.1cm,height=2.1cm,valign=center,colback=cgra!6,colframe=cgra,arc=2.5mm,boxrule=0.9pt,left=3mm]{\Large\bfseries\color{cgra}b)}\quad\begin{minipage}[c]{5.3cm}\raggedright\Large $120~\text{km/h}$\end{minipage}\end{tcolorbox}\par\vspace{2.5mm}
\begin{tcolorbox}[enhanced,nobeforeafter,width=7.1cm,height=2.1cm,valign=center,colback=cgra!6,colframe=cgra,arc=2.5mm,boxrule=0.9pt,left=3mm]{\Large\bfseries\color{cgra}c)}\quad\begin{minipage}[c]{5.3cm}\raggedright\Large $110~\text{km/h}$\end{minipage}\end{tcolorbox}\hfill\begin{tcolorbox}[enhanced,nobeforeafter,width=7.1cm,height=2.1cm,valign=center,colback=cgra!6,colframe=cgra,arc=2.5mm,boxrule=0.9pt,left=3mm]{\Large\bfseries\color{cgra}d)}\quad\begin{minipage}[c]{5.3cm}\raggedright\Large $55~\text{km/h}$\end{minipage}\end{tcolorbox}
\end{tcolorbox}
\newpage
\begin{tcolorbox}[enhanced,width=16cm,height=14cm,colback=white,colframe=cgra,boxrule=2pt,arc=0mm,
left=6mm,right=6mm,top=0mm,bottom=4mm,sharp corners,
overlay={\fill[cgra] (frame.north west) rectangle ([yshift=-1.25cm]frame.north east);
\node[anchor=west,text=white,font=\sffamily\bfseries\LARGE] at ([xshift=5mm,yshift=-0.62cm]frame.north west) {Question 84 D};
\node[anchor=east,text=white,font=\sffamily\itshape\large] at ([xshift=-5mm,yshift=-0.62cm]frame.north east) {Grandeurs et proportionnalité};}]
\vspace*{1.75cm}
\begin{minipage}[t][6.9cm][c]{\linewidth}\LARGE\centering
\begin{minipage}{0.95\linewidth}\raggedright\linespread{1.25}\selectfont Un véhicule parcourt $70$ km en 1 h 15 min. Quelle est sa vitesse moyenne~?\end{minipage}
\end{minipage}

\vfill
\begin{tcolorbox}[enhanced,nobeforeafter,width=7.1cm,height=2.1cm,valign=center,colback=cgra!6,colframe=cgra,arc=2.5mm,boxrule=0.9pt,left=3mm]{\Large\bfseries\color{cgra}a)}\quad\begin{minipage}[c]{5.3cm}\raggedright\Large $56~\text{km/h}$\end{minipage}\end{tcolorbox}\hfill\begin{tcolorbox}[enhanced,nobeforeafter,width=7.1cm,height=2.1cm,valign=center,colback=cgra!6,colframe=cgra,arc=2.5mm,boxrule=0.9pt,left=3mm]{\Large\bfseries\color{cgra}b)}\quad\begin{minipage}[c]{5.3cm}\raggedright\Large $87{,}5~\text{km/h}$\end{minipage}\end{tcolorbox}\par\vspace{2.5mm}
\begin{tcolorbox}[enhanced,nobeforeafter,width=7.1cm,height=2.1cm,valign=center,colback=cgra!6,colframe=cgra,arc=2.5mm,boxrule=0.9pt,left=3mm]{\Large\bfseries\color{cgra}c)}\quad\begin{minipage}[c]{5.3cm}\raggedright\Large $28~\text{km/h}$\end{minipage}\end{tcolorbox}\hfill\begin{tcolorbox}[enhanced,nobeforeafter,width=7.1cm,height=2.1cm,valign=center,colback=cgra!6,colframe=cgra,arc=2.5mm,boxrule=0.9pt,left=3mm]{\Large\bfseries\color{cgra}d)}\quad\begin{minipage}[c]{5.3cm}\raggedright\Large $66~\text{km/h}$\end{minipage}\end{tcolorbox}
\end{tcolorbox}
\newpage
\begin{tcolorbox}[enhanced,width=16cm,height=14cm,colback=white,colframe=cgra,boxrule=2pt,arc=0mm,
left=6mm,right=6mm,top=0mm,bottom=4mm,sharp corners,
overlay={\fill[cgra] (frame.north west) rectangle ([yshift=-1.25cm]frame.north east);
\node[anchor=west,text=white,font=\sffamily\bfseries\LARGE] at ([xshift=5mm,yshift=-0.62cm]frame.north west) {Question 85 D};
\node[anchor=east,text=white,font=\sffamily\itshape\large] at ([xshift=-5mm,yshift=-0.62cm]frame.north east) {Grandeurs et proportionnalité};}]
\vspace*{1.75cm}
\begin{minipage}[t][6.9cm][c]{\linewidth}\LARGE\centering
\begin{minipage}{0.95\linewidth}\raggedright\linespread{1.25}\selectfont Un article coûte $75$ F. Son prix augmente de $40\,\%$. Quel est son nouveau prix~?\end{minipage}
\end{minipage}

\vfill
\begin{tcolorbox}[enhanced,nobeforeafter,width=7.1cm,height=2.1cm,valign=center,colback=cgra!6,colframe=cgra,arc=2.5mm,boxrule=0.9pt,left=3mm]{\Large\bfseries\color{cgra}a)}\quad\begin{minipage}[c]{5.3cm}\raggedright\Large $105~\text{F}$\end{minipage}\end{tcolorbox}\hfill\begin{tcolorbox}[enhanced,nobeforeafter,width=7.1cm,height=2.1cm,valign=center,colback=cgra!6,colframe=cgra,arc=2.5mm,boxrule=0.9pt,left=3mm]{\Large\bfseries\color{cgra}b)}\quad\begin{minipage}[c]{5.3cm}\raggedright\Large $115~\text{F}$\end{minipage}\end{tcolorbox}\par\vspace{2.5mm}
\begin{tcolorbox}[enhanced,nobeforeafter,width=7.1cm,height=2.1cm,valign=center,colback=cgra!6,colframe=cgra,arc=2.5mm,boxrule=0.9pt,left=3mm]{\Large\bfseries\color{cgra}c)}\quad\begin{minipage}[c]{5.3cm}\raggedright\Large $45~\text{F}$\end{minipage}\end{tcolorbox}\hfill\begin{tcolorbox}[enhanced,nobeforeafter,width=7.1cm,height=2.1cm,valign=center,colback=cgra!6,colframe=cgra,arc=2.5mm,boxrule=0.9pt,left=3mm]{\Large\bfseries\color{cgra}d)}\quad\begin{minipage}[c]{5.3cm}\raggedright\Large $30~\text{F}$\end{minipage}\end{tcolorbox}
\end{tcolorbox}
\newpage
\begin{tcolorbox}[enhanced,width=16cm,height=14cm,colback=white,colframe=cgra,boxrule=2pt,arc=0mm,
left=6mm,right=6mm,top=0mm,bottom=4mm,sharp corners,
overlay={\fill[cgra] (frame.north west) rectangle ([yshift=-1.25cm]frame.north east);
\node[anchor=west,text=white,font=\sffamily\bfseries\LARGE] at ([xshift=5mm,yshift=-0.62cm]frame.north west) {Question 86 D};
\node[anchor=east,text=white,font=\sffamily\itshape\large] at ([xshift=-5mm,yshift=-0.62cm]frame.north east) {Grandeurs et proportionnalité};}]
\vspace*{1.75cm}
\begin{minipage}[t][6.9cm][c]{\linewidth}\LARGE\centering
\begin{minipage}{0.95\linewidth}\raggedright\linespread{1.25}\selectfont Un article coûte $1\,200$ F. Son prix baisse de $5\,\%$. Quel est son nouveau prix~?\end{minipage}
\end{minipage}

\vfill
\begin{tcolorbox}[enhanced,nobeforeafter,width=7.1cm,height=2.1cm,valign=center,colback=cgra!6,colframe=cgra,arc=2.5mm,boxrule=0.9pt,left=3mm]{\Large\bfseries\color{cgra}a)}\quad\begin{minipage}[c]{5.3cm}\raggedright\Large $60~\text{F}$\end{minipage}\end{tcolorbox}\hfill\begin{tcolorbox}[enhanced,nobeforeafter,width=7.1cm,height=2.1cm,valign=center,colback=cgra!6,colframe=cgra,arc=2.5mm,boxrule=0.9pt,left=3mm]{\Large\bfseries\color{cgra}b)}\quad\begin{minipage}[c]{5.3cm}\raggedright\Large $1\,260~\text{F}$\end{minipage}\end{tcolorbox}\par\vspace{2.5mm}
\begin{tcolorbox}[enhanced,nobeforeafter,width=7.1cm,height=2.1cm,valign=center,colback=cgra!6,colframe=cgra,arc=2.5mm,boxrule=0.9pt,left=3mm]{\Large\bfseries\color{cgra}c)}\quad\begin{minipage}[c]{5.3cm}\raggedright\Large $1\,140~\text{F}$\end{minipage}\end{tcolorbox}\hfill\begin{tcolorbox}[enhanced,nobeforeafter,width=7.1cm,height=2.1cm,valign=center,colback=cgra!6,colframe=cgra,arc=2.5mm,boxrule=0.9pt,left=3mm]{\Large\bfseries\color{cgra}d)}\quad\begin{minipage}[c]{5.3cm}\raggedright\Large $1\,195~\text{F}$\end{minipage}\end{tcolorbox}
\end{tcolorbox}
\newpage
\begin{tcolorbox}[enhanced,width=16cm,height=14cm,colback=white,colframe=cgra,boxrule=2pt,arc=0mm,
left=6mm,right=6mm,top=0mm,bottom=4mm,sharp corners,
overlay={\fill[cgra] (frame.north west) rectangle ([yshift=-1.25cm]frame.north east);
\node[anchor=west,text=white,font=\sffamily\bfseries\LARGE] at ([xshift=5mm,yshift=-0.62cm]frame.north west) {Question 87 D};
\node[anchor=east,text=white,font=\sffamily\itshape\large] at ([xshift=-5mm,yshift=-0.62cm]frame.north east) {Grandeurs et proportionnalité};}]
\vspace*{1.75cm}
\begin{minipage}[t][6.9cm][c]{\linewidth}\LARGE\centering
\begin{minipage}{0.95\linewidth}\raggedright\linespread{1.25}\selectfont Un prix passe de $200$ F à $260$ F. Quel est le pourcentage d'évolution~?\end{minipage}
\end{minipage}

\vfill
\begin{tcolorbox}[enhanced,nobeforeafter,width=7.1cm,height=2.1cm,valign=center,colback=cgra!6,colframe=cgra,arc=2.5mm,boxrule=0.9pt,left=3mm]{\Large\bfseries\color{cgra}a)}\quad\begin{minipage}[c]{5.3cm}\raggedright\Large $-30\,\%$\end{minipage}\end{tcolorbox}\hfill\begin{tcolorbox}[enhanced,nobeforeafter,width=7.1cm,height=2.1cm,valign=center,colback=cgra!6,colframe=cgra,arc=2.5mm,boxrule=0.9pt,left=3mm]{\Large\bfseries\color{cgra}b)}\quad\begin{minipage}[c]{5.3cm}\raggedright\Large $+60\,\%$\end{minipage}\end{tcolorbox}\par\vspace{2.5mm}
\begin{tcolorbox}[enhanced,nobeforeafter,width=7.1cm,height=2.1cm,valign=center,colback=cgra!6,colframe=cgra,arc=2.5mm,boxrule=0.9pt,left=3mm]{\Large\bfseries\color{cgra}c)}\quad\begin{minipage}[c]{5.3cm}\raggedright\Large $+15\,\%$\end{minipage}\end{tcolorbox}\hfill\begin{tcolorbox}[enhanced,nobeforeafter,width=7.1cm,height=2.1cm,valign=center,colback=cgra!6,colframe=cgra,arc=2.5mm,boxrule=0.9pt,left=3mm]{\Large\bfseries\color{cgra}d)}\quad\begin{minipage}[c]{5.3cm}\raggedright\Large $+30\,\%$\end{minipage}\end{tcolorbox}
\end{tcolorbox}
\newpage
\begin{tcolorbox}[enhanced,width=16cm,height=14cm,colback=white,colframe=cgra,boxrule=2pt,arc=0mm,
left=6mm,right=6mm,top=0mm,bottom=4mm,sharp corners,
overlay={\fill[cgra] (frame.north west) rectangle ([yshift=-1.25cm]frame.north east);
\node[anchor=west,text=white,font=\sffamily\bfseries\LARGE] at ([xshift=5mm,yshift=-0.62cm]frame.north west) {Question 88 D};
\node[anchor=east,text=white,font=\sffamily\itshape\large] at ([xshift=-5mm,yshift=-0.62cm]frame.north east) {Grandeurs et proportionnalité};}]
\vspace*{1.75cm}
\begin{minipage}[t][6.9cm][c]{\linewidth}\LARGE\centering
\begin{minipage}{0.95\linewidth}\raggedright\linespread{1.25}\selectfont Un prix passe de $1\,000$ F à $850$ F. Quel est le pourcentage d'évolution~?\end{minipage}
\end{minipage}

\vfill
\begin{tcolorbox}[enhanced,nobeforeafter,width=7.1cm,height=2.1cm,valign=center,colback=cgra!6,colframe=cgra,arc=2.5mm,boxrule=0.9pt,left=3mm]{\Large\bfseries\color{cgra}a)}\quad\begin{minipage}[c]{5.3cm}\raggedright\Large $-7{,}5\,\%$\end{minipage}\end{tcolorbox}\hfill\begin{tcolorbox}[enhanced,nobeforeafter,width=7.1cm,height=2.1cm,valign=center,colback=cgra!6,colframe=cgra,arc=2.5mm,boxrule=0.9pt,left=3mm]{\Large\bfseries\color{cgra}b)}\quad\begin{minipage}[c]{5.3cm}\raggedright\Large $+15\,\%$\end{minipage}\end{tcolorbox}\par\vspace{2.5mm}
\begin{tcolorbox}[enhanced,nobeforeafter,width=7.1cm,height=2.1cm,valign=center,colback=cgra!6,colframe=cgra,arc=2.5mm,boxrule=0.9pt,left=3mm]{\Large\bfseries\color{cgra}c)}\quad\begin{minipage}[c]{5.3cm}\raggedright\Large $-15\,\%$\end{minipage}\end{tcolorbox}\hfill\begin{tcolorbox}[enhanced,nobeforeafter,width=7.1cm,height=2.1cm,valign=center,colback=cgra!6,colframe=cgra,arc=2.5mm,boxrule=0.9pt,left=3mm]{\Large\bfseries\color{cgra}d)}\quad\begin{minipage}[c]{5.3cm}\raggedright\Large $-150\,\%$\end{minipage}\end{tcolorbox}
\end{tcolorbox}
\newpage
\begin{tcolorbox}[enhanced,width=16cm,height=14cm,colback=white,colframe=cgra,boxrule=2pt,arc=0mm,
left=6mm,right=6mm,top=0mm,bottom=4mm,sharp corners,
overlay={\fill[cgra] (frame.north west) rectangle ([yshift=-1.25cm]frame.north east);
\node[anchor=west,text=white,font=\sffamily\bfseries\LARGE] at ([xshift=5mm,yshift=-0.62cm]frame.north west) {Question 89 D};
\node[anchor=east,text=white,font=\sffamily\itshape\large] at ([xshift=-5mm,yshift=-0.62cm]frame.north east) {Grandeurs et proportionnalité};}]
\vspace*{1.75cm}
\begin{minipage}[t][6.9cm][c]{\linewidth}\LARGE\centering
\begin{minipage}{0.95\linewidth}\raggedright\linespread{1.25}\selectfont Un robinet remplit une cuve de $84$ L en $7$ min. Quel est son débit~?\end{minipage}
\end{minipage}

\vfill
\begin{tcolorbox}[enhanced,nobeforeafter,width=7.1cm,height=2.1cm,valign=center,colback=cgra!6,colframe=cgra,arc=2.5mm,boxrule=0.9pt,left=3mm]{\Large\bfseries\color{cgra}a)}\quad\begin{minipage}[c]{5.3cm}\raggedright\Large $14~\text{L/min}$\end{minipage}\end{tcolorbox}\hfill\begin{tcolorbox}[enhanced,nobeforeafter,width=7.1cm,height=2.1cm,valign=center,colback=cgra!6,colframe=cgra,arc=2.5mm,boxrule=0.9pt,left=3mm]{\Large\bfseries\color{cgra}b)}\quad\begin{minipage}[c]{5.3cm}\raggedright\Large $588~\text{L/min}$\end{minipage}\end{tcolorbox}\par\vspace{2.5mm}
\begin{tcolorbox}[enhanced,nobeforeafter,width=7.1cm,height=2.1cm,valign=center,colback=cgra!6,colframe=cgra,arc=2.5mm,boxrule=0.9pt,left=3mm]{\Large\bfseries\color{cgra}c)}\quad\begin{minipage}[c]{5.3cm}\raggedright\Large $12~\text{L/min}$\end{minipage}\end{tcolorbox}\hfill\begin{tcolorbox}[enhanced,nobeforeafter,width=7.1cm,height=2.1cm,valign=center,colback=cgra!6,colframe=cgra,arc=2.5mm,boxrule=0.9pt,left=3mm]{\Large\bfseries\color{cgra}d)}\quad\begin{minipage}[c]{5.3cm}\raggedright\large $\dfrac{25}{3}~\text{L/min}$\end{minipage}\end{tcolorbox}
\end{tcolorbox}
\newpage
\begin{tcolorbox}[enhanced,width=16cm,height=14cm,colback=white,colframe=cgra,boxrule=2pt,arc=0mm,
left=6mm,right=6mm,top=0mm,bottom=4mm,sharp corners,
overlay={\fill[cgra] (frame.north west) rectangle ([yshift=-1.25cm]frame.north east);
\node[anchor=west,text=white,font=\sffamily\bfseries\LARGE] at ([xshift=5mm,yshift=-0.62cm]frame.north west) {Question 90 D};
\node[anchor=east,text=white,font=\sffamily\itshape\large] at ([xshift=-5mm,yshift=-0.62cm]frame.north east) {Grandeurs et proportionnalité};}]
\vspace*{1.75cm}
\begin{minipage}[t][6.9cm][c]{\linewidth}\LARGE\centering
\begin{minipage}{0.95\linewidth}\raggedright\linespread{1.25}\selectfont Un robinet remplit une cuve de $150$ L en $20$ min. Quel est son débit~?\end{minipage}
\end{minipage}

\vfill
\begin{tcolorbox}[enhanced,nobeforeafter,width=7.1cm,height=2.1cm,valign=center,colback=cgra!6,colframe=cgra,arc=2.5mm,boxrule=0.9pt,left=3mm]{\Large\bfseries\color{cgra}a)}\quad\begin{minipage}[c]{5.3cm}\raggedright\Large $3\,000~\text{L/min}$\end{minipage}\end{tcolorbox}\hfill\begin{tcolorbox}[enhanced,nobeforeafter,width=7.1cm,height=2.1cm,valign=center,colback=cgra!6,colframe=cgra,arc=2.5mm,boxrule=0.9pt,left=3mm]{\Large\bfseries\color{cgra}b)}\quad\begin{minipage}[c]{5.3cm}\raggedright\Large $7{,}5~\text{L/min}$\end{minipage}\end{tcolorbox}\par\vspace{2.5mm}
\begin{tcolorbox}[enhanced,nobeforeafter,width=7.1cm,height=2.1cm,valign=center,colback=cgra!6,colframe=cgra,arc=2.5mm,boxrule=0.9pt,left=3mm]{\Large\bfseries\color{cgra}c)}\quad\begin{minipage}[c]{5.3cm}\raggedright\Large $9{,}5~\text{L/min}$\end{minipage}\end{tcolorbox}\hfill\begin{tcolorbox}[enhanced,nobeforeafter,width=7.1cm,height=2.1cm,valign=center,colback=cgra!6,colframe=cgra,arc=2.5mm,boxrule=0.9pt,left=3mm]{\Large\bfseries\color{cgra}d)}\quad\begin{minipage}[c]{5.3cm}\raggedright\large $\dfrac{40}{3}~\text{L/min}$\end{minipage}\end{tcolorbox}
\end{tcolorbox}
\newpage
\begin{tcolorbox}[enhanced,width=16cm,height=14cm,colback=white,colframe=cgra,boxrule=2pt,arc=0mm,
left=6mm,right=6mm,top=0mm,bottom=4mm,sharp corners,
overlay={\fill[cgra] (frame.north west) rectangle ([yshift=-1.25cm]frame.north east);
\node[anchor=west,text=white,font=\sffamily\bfseries\LARGE] at ([xshift=5mm,yshift=-0.62cm]frame.north west) {Question 91 D};
\node[anchor=east,text=white,font=\sffamily\itshape\large] at ([xshift=-5mm,yshift=-0.62cm]frame.north east) {Grandeurs et proportionnalité};}]
\vspace*{1.75cm}
\begin{minipage}[t][6.9cm][c]{\linewidth}\LARGE\centering
\begin{minipage}{0.95\linewidth}\raggedright\linespread{1.25}\selectfont Sur une carte à l'échelle $\dfrac{1}{2\,000}$, deux points sont distants de $2{,}5$ cm. Quelle est la distance réelle~?\end{minipage}
\end{minipage}

\vfill
\begin{tcolorbox}[enhanced,nobeforeafter,width=7.1cm,height=2.1cm,valign=center,colback=cgra!6,colframe=cgra,arc=2.5mm,boxrule=0.9pt,left=3mm]{\Large\bfseries\color{cgra}a)}\quad\begin{minipage}[c]{5.3cm}\raggedright\Large $5~\text{m}$\end{minipage}\end{tcolorbox}\hfill\begin{tcolorbox}[enhanced,nobeforeafter,width=7.1cm,height=2.1cm,valign=center,colback=cgra!6,colframe=cgra,arc=2.5mm,boxrule=0.9pt,left=3mm]{\Large\bfseries\color{cgra}b)}\quad\begin{minipage}[c]{5.3cm}\raggedright\Large $50~\text{m}$\end{minipage}\end{tcolorbox}\par\vspace{2.5mm}
\begin{tcolorbox}[enhanced,nobeforeafter,width=7.1cm,height=2.1cm,valign=center,colback=cgra!6,colframe=cgra,arc=2.5mm,boxrule=0.9pt,left=3mm]{\Large\bfseries\color{cgra}c)}\quad\begin{minipage}[c]{5.3cm}\raggedright\Large $5\,000~\text{m}$\end{minipage}\end{tcolorbox}\hfill\begin{tcolorbox}[enhanced,nobeforeafter,width=7.1cm,height=2.1cm,valign=center,colback=cgra!6,colframe=cgra,arc=2.5mm,boxrule=0.9pt,left=3mm]{\Large\bfseries\color{cgra}d)}\quad\begin{minipage}[c]{5.3cm}\raggedright\Large $500~\text{m}$\end{minipage}\end{tcolorbox}
\end{tcolorbox}
\newpage
\begin{tcolorbox}[enhanced,width=16cm,height=14cm,colback=white,colframe=cgra,boxrule=2pt,arc=0mm,
left=6mm,right=6mm,top=0mm,bottom=4mm,sharp corners,
overlay={\fill[cgra] (frame.north west) rectangle ([yshift=-1.25cm]frame.north east);
\node[anchor=west,text=white,font=\sffamily\bfseries\LARGE] at ([xshift=5mm,yshift=-0.62cm]frame.north west) {Question 92 D};
\node[anchor=east,text=white,font=\sffamily\itshape\large] at ([xshift=-5mm,yshift=-0.62cm]frame.north east) {Grandeurs et proportionnalité};}]
\vspace*{1.75cm}
\begin{minipage}[t][6.9cm][c]{\linewidth}\LARGE\centering
\begin{minipage}{0.95\linewidth}\raggedright\linespread{1.25}\selectfont Sur une carte à l'échelle $\dfrac{1}{10\,000}$, deux points sont distants de $7$ cm. Quelle est la distance réelle~?\end{minipage}
\end{minipage}

\vfill
\begin{tcolorbox}[enhanced,nobeforeafter,width=7.1cm,height=2.1cm,valign=center,colback=cgra!6,colframe=cgra,arc=2.5mm,boxrule=0.9pt,left=3mm]{\Large\bfseries\color{cgra}a)}\quad\begin{minipage}[c]{5.3cm}\raggedright\Large $70~\text{m}$\end{minipage}\end{tcolorbox}\hfill\begin{tcolorbox}[enhanced,nobeforeafter,width=7.1cm,height=2.1cm,valign=center,colback=cgra!6,colframe=cgra,arc=2.5mm,boxrule=0.9pt,left=3mm]{\Large\bfseries\color{cgra}b)}\quad\begin{minipage}[c]{5.3cm}\raggedright\Large $700~\text{m}$\end{minipage}\end{tcolorbox}\par\vspace{2.5mm}
\begin{tcolorbox}[enhanced,nobeforeafter,width=7.1cm,height=2.1cm,valign=center,colback=cgra!6,colframe=cgra,arc=2.5mm,boxrule=0.9pt,left=3mm]{\Large\bfseries\color{cgra}c)}\quad\begin{minipage}[c]{5.3cm}\raggedright\Large $70\,000~\text{m}$\end{minipage}\end{tcolorbox}\hfill\begin{tcolorbox}[enhanced,nobeforeafter,width=7.1cm,height=2.1cm,valign=center,colback=cgra!6,colframe=cgra,arc=2.5mm,boxrule=0.9pt,left=3mm]{\Large\bfseries\color{cgra}d)}\quad\begin{minipage}[c]{5.3cm}\raggedright\Large $7\,000~\text{m}$\end{minipage}\end{tcolorbox}
\end{tcolorbox}
\newpage
\begin{tcolorbox}[enhanced,width=16cm,height=14cm,colback=white,colframe=cgra,boxrule=2pt,arc=0mm,
left=6mm,right=6mm,top=0mm,bottom=4mm,sharp corners,
overlay={\fill[cgra] (frame.north west) rectangle ([yshift=-1.25cm]frame.north east);
\node[anchor=west,text=white,font=\sffamily\bfseries\LARGE] at ([xshift=5mm,yshift=-0.62cm]frame.north west) {Question 93 D};
\node[anchor=east,text=white,font=\sffamily\itshape\large] at ([xshift=-5mm,yshift=-0.62cm]frame.north east) {Grandeurs et proportionnalité};}]
\vspace*{1.75cm}
\begin{minipage}[t][6.9cm][c]{\linewidth}\LARGE\centering
\begin{minipage}{0.95\linewidth}\raggedright\linespread{1.25}\selectfont On partage $150$ bonbons entre Léa et Noa dans le ratio $7:3$. Combien chacun en reçoit-il (Léa, puis Noa)~?\end{minipage}
\end{minipage}

\vfill
\begin{tcolorbox}[enhanced,nobeforeafter,width=7.1cm,height=2.1cm,valign=center,colback=cgra!6,colframe=cgra,arc=2.5mm,boxrule=0.9pt,left=3mm]{\Large\bfseries\color{cgra}a)}\quad\begin{minipage}[c]{5.3cm}\raggedright\Large $105\text{ et }45$\end{minipage}\end{tcolorbox}\hfill\begin{tcolorbox}[enhanced,nobeforeafter,width=7.1cm,height=2.1cm,valign=center,colback=cgra!6,colframe=cgra,arc=2.5mm,boxrule=0.9pt,left=3mm]{\Large\bfseries\color{cgra}b)}\quad\begin{minipage}[c]{5.3cm}\raggedright\Large $104\text{ et }46$\end{minipage}\end{tcolorbox}\par\vspace{2.5mm}
\begin{tcolorbox}[enhanced,nobeforeafter,width=7.1cm,height=2.1cm,valign=center,colback=cgra!6,colframe=cgra,arc=2.5mm,boxrule=0.9pt,left=3mm]{\Large\bfseries\color{cgra}c)}\quad\begin{minipage}[c]{5.3cm}\raggedright\Large $75\text{ et }75$\end{minipage}\end{tcolorbox}\hfill\begin{tcolorbox}[enhanced,nobeforeafter,width=7.1cm,height=2.1cm,valign=center,colback=cgra!6,colframe=cgra,arc=2.5mm,boxrule=0.9pt,left=3mm]{\Large\bfseries\color{cgra}d)}\quad\begin{minipage}[c]{5.3cm}\raggedright\Large $45\text{ et }105$\end{minipage}\end{tcolorbox}
\end{tcolorbox}
\newpage
\begin{tcolorbox}[enhanced,width=16cm,height=14cm,colback=white,colframe=cgra,boxrule=2pt,arc=0mm,
left=6mm,right=6mm,top=0mm,bottom=4mm,sharp corners,
overlay={\fill[cgra] (frame.north west) rectangle ([yshift=-1.25cm]frame.north east);
\node[anchor=west,text=white,font=\sffamily\bfseries\LARGE] at ([xshift=5mm,yshift=-0.62cm]frame.north west) {Question 94 D};
\node[anchor=east,text=white,font=\sffamily\itshape\large] at ([xshift=-5mm,yshift=-0.62cm]frame.north east) {Grandeurs et proportionnalité};}]
\vspace*{1.75cm}
\begin{minipage}[t][6.9cm][c]{\linewidth}\LARGE\centering
\begin{minipage}{0.95\linewidth}\raggedright\linespread{1.25}\selectfont On partage $72$ bonbons entre Léa et Noa dans le ratio $1:5$. Combien chacun en reçoit-il (Léa, puis Noa)~?\end{minipage}
\end{minipage}

\vfill
\begin{tcolorbox}[enhanced,nobeforeafter,width=7.1cm,height=2.1cm,valign=center,colback=cgra!6,colframe=cgra,arc=2.5mm,boxrule=0.9pt,left=3mm]{\Large\bfseries\color{cgra}a)}\quad\begin{minipage}[c]{5.3cm}\raggedright\Large $12\text{ et }60$\end{minipage}\end{tcolorbox}\hfill\begin{tcolorbox}[enhanced,nobeforeafter,width=7.1cm,height=2.1cm,valign=center,colback=cgra!6,colframe=cgra,arc=2.5mm,boxrule=0.9pt,left=3mm]{\Large\bfseries\color{cgra}b)}\quad\begin{minipage}[c]{5.3cm}\raggedright\Large $36\text{ et }36$\end{minipage}\end{tcolorbox}\par\vspace{2.5mm}
\begin{tcolorbox}[enhanced,nobeforeafter,width=7.1cm,height=2.1cm,valign=center,colback=cgra!6,colframe=cgra,arc=2.5mm,boxrule=0.9pt,left=3mm]{\Large\bfseries\color{cgra}c)}\quad\begin{minipage}[c]{5.3cm}\raggedright\Large $60\text{ et }12$\end{minipage}\end{tcolorbox}\hfill\begin{tcolorbox}[enhanced,nobeforeafter,width=7.1cm,height=2.1cm,valign=center,colback=cgra!6,colframe=cgra,arc=2.5mm,boxrule=0.9pt,left=3mm]{\Large\bfseries\color{cgra}d)}\quad\begin{minipage}[c]{5.3cm}\raggedright\Large $11\text{ et }61$\end{minipage}\end{tcolorbox}
\end{tcolorbox}
\newpage
\begin{tcolorbox}[enhanced,width=16cm,height=14cm,colback=white,colframe=cgra,boxrule=2pt,arc=0mm,
left=6mm,right=6mm,top=0mm,bottom=4mm,sharp corners,
overlay={\fill[cgra] (frame.north west) rectangle ([yshift=-1.25cm]frame.north east);
\node[anchor=west,text=white,font=\sffamily\bfseries\LARGE] at ([xshift=5mm,yshift=-0.62cm]frame.north west) {Question 95 D};
\node[anchor=east,text=white,font=\sffamily\itshape\large] at ([xshift=-5mm,yshift=-0.62cm]frame.north east) {Grandeurs et proportionnalité};}]
\vspace*{1.75cm}
\begin{minipage}[t][6.9cm][c]{\linewidth}\LARGE\centering
\begin{minipage}{0.95\linewidth}\raggedright\linespread{1.25}\selectfont Compléter~: $1{,}5~\text{L}=\;?\;\text{cm}^3$\end{minipage}
\end{minipage}

\vfill
\begin{tcolorbox}[enhanced,nobeforeafter,width=7.1cm,height=2.1cm,valign=center,colback=cgra!6,colframe=cgra,arc=2.5mm,boxrule=0.9pt,left=3mm]{\Large\bfseries\color{cgra}a)}\quad\begin{minipage}[c]{5.3cm}\raggedright\Large $1\,500~\text{cm}^3$\end{minipage}\end{tcolorbox}\hfill\begin{tcolorbox}[enhanced,nobeforeafter,width=7.1cm,height=2.1cm,valign=center,colback=cgra!6,colframe=cgra,arc=2.5mm,boxrule=0.9pt,left=3mm]{\Large\bfseries\color{cgra}b)}\quad\begin{minipage}[c]{5.3cm}\raggedright\Large $150~\text{cm}^3$\end{minipage}\end{tcolorbox}\par\vspace{2.5mm}
\begin{tcolorbox}[enhanced,nobeforeafter,width=7.1cm,height=2.1cm,valign=center,colback=cgra!6,colframe=cgra,arc=2.5mm,boxrule=0.9pt,left=3mm]{\Large\bfseries\color{cgra}c)}\quad\begin{minipage}[c]{5.3cm}\raggedright\Large $150\,000~\text{cm}^3$\end{minipage}\end{tcolorbox}\hfill\begin{tcolorbox}[enhanced,nobeforeafter,width=7.1cm,height=2.1cm,valign=center,colback=cgra!6,colframe=cgra,arc=2.5mm,boxrule=0.9pt,left=3mm]{\Large\bfseries\color{cgra}d)}\quad\begin{minipage}[c]{5.3cm}\raggedright\Large $15\,000~\text{cm}^3$\end{minipage}\end{tcolorbox}
\end{tcolorbox}
\newpage
\begin{tcolorbox}[enhanced,width=16cm,height=14cm,colback=white,colframe=cgra,boxrule=2pt,arc=0mm,
left=6mm,right=6mm,top=0mm,bottom=4mm,sharp corners,
overlay={\fill[cgra] (frame.north west) rectangle ([yshift=-1.25cm]frame.north east);
\node[anchor=west,text=white,font=\sffamily\bfseries\LARGE] at ([xshift=5mm,yshift=-0.62cm]frame.north west) {Question 96 D};
\node[anchor=east,text=white,font=\sffamily\itshape\large] at ([xshift=-5mm,yshift=-0.62cm]frame.north east) {Grandeurs et proportionnalité};}]
\vspace*{1.75cm}
\begin{minipage}[t][6.9cm][c]{\linewidth}\LARGE\centering
\begin{minipage}{0.95\linewidth}\raggedright\linespread{1.25}\selectfont Compléter~: $8\,000~\text{dm}^3=\;?\;\text{m}^3$\end{minipage}
\end{minipage}

\vfill
\begin{tcolorbox}[enhanced,nobeforeafter,width=7.1cm,height=2.1cm,valign=center,colback=cgra!6,colframe=cgra,arc=2.5mm,boxrule=0.9pt,left=3mm]{\Large\bfseries\color{cgra}a)}\quad\begin{minipage}[c]{5.3cm}\raggedright\Large $80~\text{m}^3$\end{minipage}\end{tcolorbox}\hfill\begin{tcolorbox}[enhanced,nobeforeafter,width=7.1cm,height=2.1cm,valign=center,colback=cgra!6,colframe=cgra,arc=2.5mm,boxrule=0.9pt,left=3mm]{\Large\bfseries\color{cgra}b)}\quad\begin{minipage}[c]{5.3cm}\raggedright\Large $8~\text{m}^3$\end{minipage}\end{tcolorbox}\par\vspace{2.5mm}
\begin{tcolorbox}[enhanced,nobeforeafter,width=7.1cm,height=2.1cm,valign=center,colback=cgra!6,colframe=cgra,arc=2.5mm,boxrule=0.9pt,left=3mm]{\Large\bfseries\color{cgra}c)}\quad\begin{minipage}[c]{5.3cm}\raggedright\Large $0{,}8~\text{m}^3$\end{minipage}\end{tcolorbox}\hfill\begin{tcolorbox}[enhanced,nobeforeafter,width=7.1cm,height=2.1cm,valign=center,colback=cgra!6,colframe=cgra,arc=2.5mm,boxrule=0.9pt,left=3mm]{\Large\bfseries\color{cgra}d)}\quad\begin{minipage}[c]{5.3cm}\raggedright\Large $800~\text{m}^3$\end{minipage}\end{tcolorbox}
\end{tcolorbox}
\newpage
\end{document}